% Framework source amendment v3.068: scoped geometry/operator additions, physical occurrence unchanged.
\documentclass[11pt,a4paper]{article}
\usepackage[T1]{fontenc}
\usepackage[utf8]{inputenc}
\usepackage{lmodern,amsmath,amssymb,amsthm,mathtools,bm,mathrsfs}
\usepackage[margin=27mm]{geometry}
\usepackage{booktabs,longtable,array,enumitem,microtype,calc}
\usepackage[numbers,sort&compress]{natbib}
\providecommand{\printbibliography}{\bibliographystyle{plainnat}\bibliography{Born-v6_59}}
\usepackage{xurl}
\usepackage[colorlinks=true,linkcolor=blue,citecolor=blue,urlcolor=blue]{hyperref}
\usepackage{bookmark}
\usepackage{alphalph}


\makeatletter
\@ifpackageloaded{xr-hyper}{}{%
  \@ifpackageloaded{xr}{}{\usepackage{xr-hyper}}%
}
\makeatother

\ExplSyntaxOn
\prop_new:N \g_bh_external_refs_prop
\cs_new_eq:NN \BHOriginalExternalDocument \externaldocument

\cs_new_protected:Npn \bh_external_document:nnnn #1#2#3#4
 {
  \tl_set:Nn \l_tmpa_tl {#3}
  \regex_replace_once:nnN
   { \A (?: xref/ )? (Paper[0-9]+[A-F]?) .* \Z }
   { \1 }
   \l_tmpa_tl
  \tl_set:Nx \l_tmpb_tl { #1 | \tl_use:N \l_tmpa_tl }
  \prop_if_in:NVTF \g_bh_external_refs_prop \l_tmpb_tl
   {
    \PackageInfo{BornHopf}
     {Skipping~duplicate~external~refs~for~#3}
   }
   {
    \prop_gput:NVn \g_bh_external_refs_prop \l_tmpb_tl {#3}
    \IfNoValueTF{#4}
     { \BHOriginalExternalDocument[#1][#2]{#3} }
     { \BHOriginalExternalDocument[#1][#2]{#3}[#4] }
   }
 }

\RenewDocumentCommand{\externaldocument}{O{} O{} m o}
 { \bh_external_document:nnnn {#1}{#2}{#3}{#4} }

\NewDocumentCommand{\BHExternalRefs}{m}
 {
  \tl_set:Nx \l_tmpa_tl { \c_sys_jobname_str }
  \regex_replace_once:nnN
   { \A (Paper[0-9]+[A-F]?) .* \Z }
   { \1 }
   \l_tmpa_tl
  \tl_set:Nn \l_tmpb_tl { #1 }
  \regex_replace_once:nnN
   { \A (Paper[0-9]+[A-F]?) .* \Z }
   { \1 }
   \l_tmpb_tl
  \str_if_eq:eeTF { \tl_use:N \l_tmpa_tl } { \tl_use:N \l_tmpb_tl }
   {
    \PackageInfo{BornHopf}
     {Skipping~external~refs~for~#1~because~this~is~the~current~document}
   }
   {
    \IfFileExists{xref/#1.aux}
     { \externaldocument[][nocite]{xref/#1}[#1.pdf] }
     {
      \IfFileExists{#1.aux}
       { \externaldocument[][nocite]{#1} }
       {
        \PackageInfo{BornHopf}
         {External~aux~file~#1.aux~not~found;~no~labels~imported}
       }
     }
   }
 }
\ExplSyntaxOff

\BHExternalRefs{Paper1A-Geometry-v5_180}
\BHExternalRefs{Paper1B-ExceptionalGeometry-v1_22}
\BHExternalRefs{Paper1D-TechnicalMaterial-v1_68}
\BHExternalRefs{Paper2F-TechnicalMaterial-v1_70}
\BHExternalRefs{Paper3B-SpectralAction-v1_191}
\BHExternalRefs{Paper3E-GravityAction-v1_203}
\BHExternalRefs{Paper3F-CompleteDynamics-v1_275}


\providecommand{\tightlist}{\setlength{\itemsep}{0pt}\setlength{\parskip}{0pt}}
\setlength{\emergencystretch}{3em}
\setcounter{tocdepth}{2}
\numberwithin{equation}{section}
\newtheorem{theorem}{Theorem}[section]
\newtheorem{proposition}[theorem]{Proposition}
\newtheorem{lemma}[theorem]{Lemma}
\newtheorem{corollary}[theorem]{Corollary}
\theoremstyle{definition}
\newtheorem{definition}[theorem]{Definition}
\theoremstyle{remark}
\newtheorem{remark}[theorem]{Remark}
\newcommand{\R}{\mathbb R}
\newcommand{\C}{\mathbb C}
\newcommand{\Z}{\mathbb Z}
\newcommand{\Q}{\mathbb Q}
\newcommand{\HH}{\mathbb H}
\newcommand{\F}{\mathbb F}
\newcommand{\T}{\mathbb T}
\newcommand{\SU}{\mathrm{SU}}
\newcommand{\SO}{\mathrm{SO}}
\newcommand{\U}{\mathrm U}
\newcommand{\Spin}{\mathrm{Spin}}
\newcommand{\Cl}{\mathrm{Cl}}
\newcommand{\id}{\mathrm{id}}
\newcommand{\Ad}{\operatorname{Ad}}
\newcommand{\End}{\operatorname{End}}
\newcommand{\Hom}{\operatorname{Hom}}
\newcommand{\tr}{\operatorname{tr}}
\newcommand{\Tr}{\operatorname{Tr}}
\newcommand{\rank}{\operatorname{rank}}
\newcommand{\im}{\operatorname{im}}
\newcommand{\diag}{\operatorname{diag}}
\newcommand{\Spec}{\operatorname{Spec}}
\newcommand{\Sym}{\operatorname{Sym}}
\newcommand{\vol}{\operatorname{vol}}
\newcommand{\sgn}{\operatorname{sgn}}
\newcommand{\supp}{\operatorname{supp}}
\newcommand{\osc}{\operatorname{osc}}
\newcommand{\cA}{\mathcal A}
\newcommand{\cE}{\mathcal E}
\newcommand{\cP}{\mathcal P}
\newcommand{\cH}{\mathcal H}
\newcommand{\dd}{\mathrm d}
\newcommand{\Stirling}[2]{\genfrac{\{}{\}}{0pt}{}{#1}{#2}}
\title{\textbf{Born--Hopf Geometry I-C:}\\[0.4em]\large Reciprocal product and $M_8$ geometry}
\author{Peter Kugelmann}
\date{Version 1.146 --- 20 November 2026 \\ Framework v3.068}


\begin{document}
\providecommand{\RR}{\mathbb R}
\providecommand{\ZZ}{\mathbb Z}
\providecommand{\Sone}{S^1}
\providecommand{\TableNotes}[1]{\par\smallskip\footnotesize #1\par}
\providecommand{\CC}{\mathbb{C}}
\providecommand{\Rmom}{R_{\mathrm{mom}}}

\maketitle
\begin{abstract}
We develop the mathematical geometry of the labelled reciprocal product
\(M_8=(SU(2)\times U(1))_x\times(SU(2)\times U(1))_p\), including
exchange and relative quotients, Cartan geometry and curvature,
bi-Hopf structures, integral topology, bundle and index invariants,
representation and Clifford operator calculi, invariant metric and
complex geometry, exact finite and harmonic carriers, and reciprocal
transport. The complete mathematical classifications retain their exact
metric, marking, representation, and operator domains; they do not by
themselves assert physical occurrence or select a particular member.

The native weak \(SU(2)\) datum is used provisionally in
Section~\ref{p1b:sec:v259-conditional-native-atlas}, in the same
way as the level-one datum in Paper~I-A.  We prove its residual
bundle, connection, exchange, index, spectral and response
consequences on their exact mathematical domains.  Papers~III-B
and III-F later establish that their generated weak sector
satisfies this premise; it is not an additional permanent axiom.

The mixed rank-two classification is stated in
Theorem~\ref{p1b:thm:complete-mixed-rank-two}. On its declared domain
the coherent totals are 72 (even) and 120 (odd), with auxiliary counts
96 and 144. The evaluated natural and inherited extension tests are
respectively 4 YES/10 NO and 4 YES/28 NO. The longer technical proofs and verification ledgers cited here are
maintained in Paper~I-D; the mathematical statements and arguments in
this compendium retain their own stated domains. Section~\ref{p1b:sec:organized-rooted} assembles the
rooted hierarchy, Clifford/triality, and finite-carrier relations;
Sections~\ref{p1b:sec:organized-boundary} and
\ref{p1b:sec:organized-interfaces} state their boundary and response
interfaces with explicit domain restrictions. 
For explicitly selected binary-icosahedral markings in the two reciprocal sphere factors, the Petersen edge modules admit an exact 225-dimensional product resolution and 120/105 exchange-even/odd tensor split, without altering the manifold or selecting a dynamical coupling.
\end{abstract}
\tableofcontents
\clearpage

\part{Reciprocal datum, symmetry and global topology}
\section{Reciprocal product, group symmetry, and Cartan geometry}
\label{p1b:sec:organized-datum}

\subsection{The mathematical datum and organization}
\label{p1b:sec:datum}
\begin{definition}[Labelled reciprocal product]
\label{p1b:def:reciprocal-product}
Fix \(K=\SU(2)\times\U(1)\), two labelled copies \(K_x,K_p\), and
isomorphisms \(\iota_x:K_x\to K\), \(\iota_p:K_p\to K\).
The reciprocal product is
\[
 M_8=K_x\times K_p\cong S_x^3\times S_p^3\times S_t^1\times S_E^1.
\]
The subscripts are labels. Write \(a=(g_x,z_t)\), \(b=(g_p,z_E)\)
in common-model coordinates. Exchange and relative comparison are
\[
 B(a,b)=(b,a),\qquad r(a,b)=ab^{-1}.
\]
All formulas in these coordinates mean their transport through
\(\iota_x,\iota_p\). A Hopf marking is a chosen maximal circle with an
orientation and a choice of direct or opposite action on the sphere.
\end{definition}
The underlying product, the labelled Lie group, the group with exchange,
and the group with Hopf or finite representation markings are different
levels of data. Each result states the level it uses. A radius,
a torus metric, a connection perturbation, a finite compression, a measure,
or a function to be differentiated is additional mathematical data when it
first occurs. \emph{Canonical} means natural relative to those data.

For applications to the native weak sector, the sole role premise
is \(\mathsf H_{\rm vert}\) of
Definition~\ref{p1b:def:conditional-residual-datum}: this role
uses the actual derived relative vertical action and inherited
exchange.  It is assumed provisionally here and proved for the
generated sector in the later chain ending at
Corollary~\ref{p3f:cor:v259-native-su2-discharge}.  The residual
bundle and parent-descended connection are independent mathematical
outputs of this paper; its proofs do not assume the later
occurrence conclusion.  Each metric, operator, action or source
input remains stated at the result that consumes it.

Paper I-A supplies the normalized one-sphere atlas \cite{Kugelmann2026a}.
Here we restate the particular one-sphere facts needed by a reciprocal
construction.
The notation \(S^3_1\) denotes a sphere together with a specified finite
representation package; it is neither an additional product factor nor a
submanifold sphere of a different dimension. Finite scalar bands and finite
spin/form carriers will be denoted separately.

The method follows Paper I-A's treatment of \(m=1\): define a carrier or
structure family, determine its free parameters and equivalences, and prove
its consequences before specializing. In particular, the metrics and
connections are studied as complete families, and the finite completion
cones include their zero and degenerate boundary strata. Mathematical
conditions such as positivity, nonzero normalization, invertibility and a
specified group action describe theorem domains. All conclusions are relative to these mathematical data.

Section~\ref{p1b:sec:all-levels} develops every pair of scalar and enriched
levels, including the mixed scalar-position/enriched-momentum family.
It tracks the actual factor types when taking exchange, invariant subspaces,
operator comparisons and limits.

The product geometry and the classification families are summarized in
Section~\ref{p1b:sec:exhaustion}. The longer spectral arguments,
obstruction proofs, marked comparisons and finite action/measure calculations
appear in the relevant mathematical sections below or, where cited,
in the technical companion Paper~I-D. Subsection~\ref{p1b:app:m8-substructures}
records the higher-rank torus-quotient construction, Hopf--Heegaard--Clifford
strata, the exact \(\Spin^c(4)\)--\(\Spin(6)\) comparison,
and the distinctions between the stated construction types. The word
\emph{complete} in a classification refers to its explicitly named class:
for example, all left-invariant metrics, all bi-invariant metric connections,
or all finite subspaces invariant under left and right translations. This
makes exhaustion a theorem with a stated domain.

Every substantive mathematical category is preceded locally by a compact result ledger. The ledger is an index, not an additional premise: each row points to the statement or labelled derivation that discharges it. The global exhaustion catalogue in Section~\ref{p1b:sec:exhaustion} is retained as the cross-category classification index.

Sections~\ref{p1b:sec:g02}--\ref{p1b:sec:g10} complete the
declared product-structure comparison with Paper~I-A. They give the
coordinate and retained-square constructions, commutants and finite
modules, contact and CR analysis, metric deformations, full bundle
spectra, evaluation measures and finite quotients. Each section states
its positive result before the detailed proof. Longer calculations are given in their corresponding subsections
or in Paper~I-D where explicitly indicated. The finite
scalar, enriched, full spin and constrained tensor bundles retain
their distinct definitions even when some dimensions coincide.


\subparagraph{Conventions.}
Matrix-valued forms use \(F_A=\dd A+A\wedge A\). The Laplacian is
nonnegative. Haar measure is normalized to probability unless \(\dd\mu_R\)
and \(V_R=2\pi^2R^3\) are written. Positive Clifford generators satisfy
\(vw+wv=2g(v,w)\), and the three-dimensional negative convention is
\(e_ie_j+e_je_i=-2\delta_{ij}\). The conventions are related in
Section~\ref{p1b:sec:clifford}. Ordered cohomology generators
are \(a_t,a_E,u_x,u_p\), of degrees \(1,1,3,3\). Set
\(\eta=a_ta_E\) and choose \(\int_{M_8}a_ta_Eu_xu_p=1\).
The scalar line with Chern class \(+\eta\) is used throughout. It is the
dual of the evaluation-character convention with class \(-\eta\) in
Paper III-A; dualization translates between the two conventions.

\subsection{Group structure, exchange and relative quotients}
\label{p1b:sec:quotients}
\begin{theorem}[Characteristic subgroups and characters]
\label{p1b:thm:group}
For the labelled Lie group,
\begin{align*}
 [M_8,M_8]&=\SU(2)_x\times\SU(2)_p,\\
 Z(M_8)&=\{\pm1\}_x\times\{\pm1\}_p\times\T^2,\\
 Z(M_8)^0&=\T^2,\qquad M_8^{\rm ab}\cong\T^2,\\
 \Hom_{\rm cont}(M_8,\U(1))&\cong\Z^2.
\end{align*}
Every continuous one-dimensional character is \(z_t^nz_E^w\).
The group automorphisms of the unmarked product form
\[
 \operatorname{Aut}(M_8)
 \cong\bigl((\SO(3)\times\SO(3))\rtimes S_2\bigr)\times GL(2,\Z).
\]
The centralizer of scalar exchange in \(GL(2,\Z)\) is
\(\{I,-I,S,-S\}\), where \(S=\left(\begin{smallmatrix}0&1\\1&0\end{smallmatrix}\right)\).
\end{theorem}
\begin{proof}
\(\SU(2)\) is perfect with center \(\{\pm1\}\), and \(\U(1)\) is abelian.
Take commutators and centers factorwise. A character kills the derived
subgroup. Both the derived subgroup and connected center are characteristic.
An automorphism permutes the two simple ideals and acts by inner
automorphisms on them. There is no nontrivial character from the semisimple
factor to the torus, or continuous homomorphism from the connected torus
to the finite semisimple center. This proves the automorphism formula.
A matrix commuting with \(S\) is
\(\left(\begin{smallmatrix}p&q\\q&p\end{smallmatrix}\right)\);
\((p-q)(p+q)=\pm1\) gives the four matrices.
\end{proof}

\begin{corollary}[Intrinsic commutator and character-carrier rigidity]
\label{p1b:cor:intrinsic-carrier-rigidity}
Put $D_6=[M_8,M_8]=\SU(2)_x\times\SU(2)_p$.  Then
\[
 D_6
 =\bigcap_{\substack{H\le M_8\\ [g,h]\in H\ \forall g,h\in M_8}} H.
\]
Thus $D_6$ is the unique least subgroup containing the complete commutator
image of $M_8$.  Dually, every continuous character of $M_8$ factors uniquely
through the abelianization
\[
 q_{\rm ab}:M_8\longrightarrow M_8^{\rm ab}\cong\T^2,
\]
and the labelled direct product identifies this quotient canonically with the
connected central scalar factor $Z(M_8)^0=\T^2_{\rm sc}$.  This last statement
is an intrinsic character-owner statement; it does not assert that an
arbitrary one-circle phase role has a unique embedding in $M_8$.
\end{corollary}
\begin{proof}
Every subgroup in the displayed intersection contains every commutator and
therefore the subgroup generated by all commutators, namely $[M_8,M_8]$.
Conversely $[M_8,M_8]$ itself is one of those subgroups.  The character
statement is the universal property of the abelianization together with
Theorem~\ref{p1b:thm:group}.
\end{proof}

\begin{theorem}[Connected closed-subgroup skeleton]
\label{p1b:thm:connected-subgroup-skeleton}
Let \(H\subset M_8\) be connected and closed, and put
\(H_{\rm ss}=[H,H]^0\).  Up to conjugacy in the two \(\SU(2)\) factors,
\[
 H_{\rm ss}\in
 \{1,\ \SU(2)_x,\ \SU(2)_p,\ \Delta\SU(2),\
       \SU(2)_x\times\SU(2)_p\}.
\]
In the one-simple-factor case whose two projections are nontrivial, the
simple factor is the graph of an automorphism of \(\SU(2)\), hence is
conjugate to \(\Delta\SU(2)\).  The connected radical of \(H\) is a
closed torus in the connected centralizer of \(H_{\rm ss}\), and
\[
\begin{array}{c|c}
H_{\rm ss}&C_{M_8}(H_{\rm ss})^0\\ \hline
1&M_8\\
\SU(2)_x&\SU(2)_p\times\T^2_{\rm sc}\\
\SU(2)_p&\SU(2)_x\times\T^2_{\rm sc}\\
\Delta\SU(2)&\T^2_{\rm sc}\\
\SU(2)_x\times\SU(2)_p&\T^2_{\rm sc}.
\end{array}
\]
Consequently every connected abelian subgroup is conjugate into a
maximal torus, and the displayed list together with rational subtori of
the corresponding centralizers gives the connected closed subgroup types.
\end{theorem}
\begin{proof}
At Lie-algebra level
\[
 \mathfrak m_8=\mathfrak{su}(2)_x\oplus
 \mathfrak{su}(2)_p\oplus\R^2_{\rm sc}.
\]
The derived algebra of a compact connected Lie algebra is semisimple.
A simple subalgebra of
\(\mathfrak{su}(2)\oplus\mathfrak{su}(2)\) is isomorphic to
\(\mathfrak{su}(2)\).  Its projection to either simple summand is zero
or injective; if both projections are nonzero it is the graph of an
automorphism.  Every automorphism of \(\SU(2)\) is inner, so this graph
is conjugate to the diagonal.  Two independent simple ideals give the
full semisimple factor.  The radical of a compact connected group is a
central torus, so it centralizes \(H_{\rm ss}\).  The displayed
centralizers follow factorwise; closed connected subtori are precisely
the rational subtori of their maximal tori.
\end{proof}

\begin{corollary}[Connected hypersubgroups and the connected normal lattice]
\label{p1b:cor:codim-one-normal}
Every connected closed seven-dimensional subgroup of \(M_8\) is
\[
 H^7_{(m,n)}=
 \SU(2)_x\times\SU(2)_p\times S^1_{(m,n)},
 \qquad
 S^1_{(m,n)}=\{(e^{im\theta},e^{in\theta})\},
 \quad \gcd(m,n)=1,
\]
and hence is diffeomorphic to \(S^3\times S^3\times S^1\).
Moreover every connected normal subgroup has Lie algebra
\[
 \mathfrak s_x\oplus\mathfrak s_p\oplus V,
 \qquad
 \mathfrak s_i\in\{0,\mathfrak{su}(2)_i\},
 \qquad V\subset\R^2_{\rm sc},
\]
with \(V\) rational when the scalar factor is required to be closed.
In particular \(\Delta\SU(2)\) is not normal.
\end{corollary}
\begin{proof}
Let \(H^7\) project to
\(D_6=\SU(2)_x\times\SU(2)_p\).  Its kernel lies in the scalar
\(\T^2\), so the image has dimension at least five.  By
Theorem~\ref{p1b:thm:connected-subgroup-skeleton}, every proper connected
subgroup of \(D_6\) has dimension at most four.  Thus the image is
\(D_6\) and the kernel is a circle.  Any residual twisting would give a
continuous character \(D_6\to\T^2/S^1\cong\U(1)\), which is trivial
because \(D_6\) is connected semisimple.  This gives the stated split
form.  The normal-subgroup assertion follows because an ideal of
\(\mathfrak m_8\) is the direct sum of a choice of the two simple ideals
and a subspace of the center; exponentiation gives the closed connected
subgroups precisely for rational central subspaces.
\end{proof}

\begin{proposition}[Circle cocharacters]
\label{p1b:result:circle-cocharacters}
Every continuous homomorphism \(\U(1)\to\SU(2)\) is conjugate to
\[
 e^{i\theta}\longmapsto
 \diag(e^{im\theta},e^{-im\theta}),\qquad m\in\Z_{\ge0}.
\]
It is nontrivial for \(m>0\) and is injective exactly when \(m=1\).
A primitive Hopf subgroup determines its axis up to orientation;
an oriented cocharacter retains the sign before Weyl conjugacy.
\end{proposition}
\begin{proof}
The compact connected abelian image lies in a maximal torus.
Its two defining characters have opposite integer weights because
the determinant is one. Weyl conjugation interchanges the weights.
For nonzero \(m\), the kernel consists of the \(m\)-th roots of unity.
\end{proof}

\begin{theorem}[Relative principal bundle and successive quotients]
\label{p1b:thm:relative}
The map \(r:M_8\to K\) is a surjective submersion with fibres the orbits
of the free right diagonal action \((a,b)\cdot h=(ah,bh)\).
It has the global trivialization
\[
 \Phi(k,h)=(kh,h),\qquad \Phi^{-1}(a,b)=(ab^{-1},b).
\]
Thus \(M_8/\Delta K\cong K\). Quotienting first by the scalar diagonal gives
\[
 X_7=M_8/\U(1)_{\rm diag}
 \cong\SU(2)_x\times\SU(2)_p\times\U(1)_{\rm rel},
 \qquad X_7/\SU(2)_{\rm diag}\cong K.
\]
Exchange descends to inversion: \(r(B(a,b))=r(a,b)^{-1}\).
Left diagonal translation gives \(r(ha,hb)=hr(a,b)h^{-1}\).
\end{theorem}
\begin{proof}
The displayed inverse proves triviality. If \(a'b'^{-1}=ab^{-1}\), take
\(h=b^{-1}b'\) to obtain \(a'=ah,b'=bh\).
The scalar quotient is \((z_t,z_E)\mapsto z_tz_E^{-1}\).
The remaining identities follow by multiplication. The quotient is a
principal-bundle orbit space; \(\Delta K\) is not normal for nonabelian \(K\).
\end{proof}

\begin{theorem}[Symmetric pair and Cartan embedding]
\label{p1b:thm:symmetric}
The differential of exchange gives
\[
 \mathfrak k\oplus\mathfrak k=\mathfrak k_+\oplus\mathfrak k_-,
 \qquad \mathfrak k_\pm=\{(X,\pm X):X\in\mathfrak k\},
\]
with brackets
\[
 [\mathfrak k_+,\mathfrak k_+]\subset\mathfrak k_+,\quad
 [\mathfrak k_+,\mathfrak k_-]\subset\mathfrak k_-,\quad
 [\mathfrak k_-,\mathfrak k_-]\subset\mathfrak k_+.
\]
The Cartan embedding is \(c(k)=(k,k^{-1})\), with \(r(c(k))=k^2\).
Its image is a subgroup exactly when \(K\) is abelian.
\end{theorem}
\begin{proof}
Compute brackets componentwise. The Cartan map
\(g\Delta K\mapsto gB(g)^{-1}\) gives the formula using \(g=(k,e)\).
The product \(c(k)c(l)=(kl,k^{-1}l^{-1})\) equals \(c(kl)\) for all
\(k,l\) exactly when they commute. The odd complement is a Lie-triple
system, with product induced by \([[X,Y],Z]\).
\end{proof}

\begin{proposition}[Exchange fixed set and quotient]
\label{p1b:prop:fixed}
The fixed set is \(\Delta K\), of dimension four. The exchange quotient
has local model \(\R^4\times(\R^4/\{\pm1\})\) along the diagonal,
is smooth off the diagonal, and is singular on it.
Exchange preserves the orientation of \(M_8\).
\end{proposition}
\begin{proof}
The tangent action is \(+I\) on the diagonal and \(-I\) on its
four-dimensional normal complement. The normal link is \(\R P^3\),
rather than \(S^3\). Interchanging two four-dimensional factors has
orientation sign \((-1)^{16}=1\).
\end{proof}

\begin{corollary}[Exact relative-fibre carrier rigidity]
\label{p1b:cor:relative-carrier-rigidity}
For the generated relative map $r(a,b)=ab^{-1}$,
\[
 r^{-1}(e)=\Delta K=\operatorname{Fix}(B).
\]
Consequently, if a subgroup $H\le M_8$ acts by right translation and its
cosets are exactly the fibres of $r$, then $H=\Delta K$.  Moreover
\[
 [\Delta K,\Delta K]=\Delta\SU(2),
\]
so $\Delta\SU(2)$ is the unique connected semisimple derived part of the
vertical carrier.
\end{corollary}
\begin{proof}
Theorem~\ref{p1b:thm:relative} identifies the fibres of $r$ with right
$\Delta K$-orbits, hence the fibre through the identity is $\Delta K$.
Proposition~\ref{p1b:prop:fixed} gives the independent fixed-set
characterization.  If right $H$-cosets are the same fibres, the identity
fibre is $H$, so $H=\Delta K$.  Finally
$\Delta K\cong\SU(2)\times\U(1)$, whose derived subgroup is
$\Delta\SU(2)$.
\end{proof}

\subsection{Cartan connection, curvature and finite defect}
\label{p1b:sec:cartan}
\begin{theorem}[Connection--solder decomposition]
\label{p1b:thm:cartan}
Let \(\theta_x=a^{-1}\dd a\), \(\theta_p=b^{-1}\dd b\), and
\[
 \cA=\tfrac12(\theta_x+\theta_p),\qquad
 \cE=\tfrac12(\theta_x-\theta_p).
\]
Then \(\cA\) is a principal connection for \(M_8\to K\), and \(\cE\)
is horizontal and adjoint-equivariant, with
\[
 D_{\cA}\cE=0,\qquad F_{\cA}=-\cE\wedge\cE.
\]
In \((a,b)=(kh,h)\),
\[
 \cA=\theta_h+\tfrac12\Ad_{h^{-1}}\theta_k,\qquad
 \cE=\tfrac12\Ad_{h^{-1}}\theta_k.
\]
\end{theorem}
\begin{proof}
A variable right diagonal transformation gives
\(\cA\mapsto h^{-1}\cA h+h^{-1}\dd h\) and
\(\cE\mapsto h^{-1}\cE h\). Their vertical values are \(X\) and zero.
Add and subtract the two Maurer--Cartan equations. Substitute \(a=kh,b=h\)
to obtain the last formulas.
\end{proof}

\begin{proposition}[Finite order-swap defect]
\label{p1b:prop:commutator}
The word \(\mu(a,b)=aba^{-1}b^{-1}\) satisfies
\[
 \mu(b,a)=\mu(a,b)^{-1},\quad
 \mu(hah^{-1},hbh^{-1})=h\mu(a,b)h^{-1},\quad
 \log\mu(e^{\epsilon X},e^{\epsilon Y})
 =\epsilon^2[X,Y]+O(\epsilon^3).
\]
Its scalar component is one. The commuting locus in \(M_8\)
has dimension six and codimension two.
\end{proposition}
\begin{proof}
Use group multiplication and the second-order exponential expansion.
A noncentral commuting pair in \(\SU(2)\) lies in a common maximal circle:
the circle conjugacy parameter has dimension two and the two entries
contribute two. Central strata have smaller dimension. Two scalar entries
give the remaining two dimensions.
\end{proof}

\begin{theorem}[Homogeneous connections, exchange and holonomy]
\label{p1b:thm:homogeneous}
On \(X_7\to K\), with structure group \(\SU(2)_{\rm diag}\), the
left-invariant principal connections are
\[
 \omega_f=f\theta_x+(1-f)\theta_p,\qquad f\in\R.
\]
The unique exchange-invariant member is \(f=1/2\). In the section
\((r,e,z)\),
\[
 A_{1/2}=\tfrac12\Theta,\qquad F_{1/2}=-\tfrac14\Theta\wedge\Theta.
\]
Its restricted holonomy is \(\SU(2)\).
The associated fundamental bundle has zero second Chern class,
and \(\tr(F\wedge F)=0\) pointwise for this connection.
\end{theorem}
\begin{proof}
An invariant connection is an equivariant linear map
\(W:\mathfrak{su}(2)^2\oplus\R\to\mathfrak{su}(2)\) restricting to the
identity on the diagonal. Irreducibility gives
\(W(X,Y,t)=fX+(1-f)Y\); no scalar-to-adjoint intertwiner exists.
Exchange sends \(f\) to \(1-f\).
The horizontal curvature values \(-[X,Y]\) span \(\mathfrak{su}(2)\);
Ambrose--Singer gives the holonomy. The global section trivializes the
bundle; curvature is pulled back from a three-sphere and its square vanishes.
\end{proof}

\begin{theorem}[Smooth Cartan deformation]
\label{p1b:thm:deformation}
For any smooth real basic \(u\),
\[
 \mathscr A_u=\cA+u\cE,\qquad
 F_u=\dd u\wedge\cE+(u^2-1)\cE\wedge\cE.
\]
On the labelled-factor section, put \(f=(1+u)/2\); then
\[
 A_f=f\Theta,\qquad F_f=\dd f\wedge\Theta+(f^2-f)\Theta\wedge\Theta.
\]
Exchange maps \(\mathscr A_u\) to \(\mathscr A_{-u\circ B}\).
\end{theorem}
\begin{proof}
A basic multiple of a horizontal equivariant form is a connection
difference. Expand curvature using \(D_{\cA}\cE=0\).
Exchange fixes \(\cA\) and negates \(\cE\).
Constant \(u=\pm1\) gives the two flat Maurer--Cartan endpoints.
\end{proof}


\subparagraph{Transgression and the cylinder functional.}
\label{p1b:sec:g11}

For the same deformation \(A_f=f\Theta\), the normalized Chern--Simons
coordinate is \(3f^2-2f^3\). On the additional round Euclidean cylinder,
the Yang--Mills functional reduces to
\[
 \frac{12\pi^2}{g_{\rm YM}^2}\int
       \{\dot f^{\,2}+4f^2(1-f)^2\}\,\dd\tau.
\]
Its unit-charge minimum is the profile
\(f=(1+\tanh(\tau-\tau_0))/2\), with action \(8\pi^2/g_{\rm YM}^2\).
Paper~I-D, Section~\ref{p1b:app:g11} proves these statements and determines the
homogeneous fluctuation spectrum. This supplies the variational
derivation of the profile used in Proposition~\ref{p1b:prop:kink}.
The infinite cylinder, its metric and the functional are additional
mathematical data; the compact product alone does not specify them.


\subsection{Scalar torus, integral eigenphases and phase lines}
\label{p1b:sec:scalar}
\begin{definition}[Primitive scalar pairing]
For a rank-one lattice \(\Lambda\subset V\), set
\(T_t=V/\Lambda\), \(T_E=V^*/\Lambda^*\). These are dual-lattice tori;
the continuous character group of a circle itself is \(\Z\).
With \(z_t=e^{i\theta_t},z_E=e^{i\theta_E}\), use
\[
 a_t=\frac{\dd\theta_t}{2\pi},\quad
 a_E=\frac{\dd\theta_E}{2\pi},\quad
 \omega_{\rm sc}=\frac{\dd\theta_t\wedge\dd\theta_E}{(2\pi)^2}.
\]
The Poincar\'e line \(\cP_{\rm sc}\) has \(c_1=\eta=[\omega_{\rm sc}]\).
\end{definition}

\begin{proposition}[Primitive phase and circle-valued moments]
\label{p1b:prop:scalar-moments}
The class \(\eta\) is primitive and \(B^*\eta=-\eta\).
For \(V_\pm=\partial_{\theta_t}\pm\partial_{\theta_E}\),
\[
 \iota_{V_+}\omega_{\rm sc}
 =-\frac{\dd(\theta_t-\theta_E)}{(2\pi)^2},\qquad
 \iota_{V_-}\omega_{\rm sc}
 =\frac{\dd(\theta_t+\theta_E)}{(2\pi)^2}.
\]
The normalized contractions have circle-valued primitives
\(q_-^{-1}\) and \(q_+\), where \(q_\pm=z_tz_E^{\pm1}\).
Neither contraction has a globally defined real primitive.
\end{proposition}
\begin{proof}
Integration over the primitive torus gives one; exchange reverses wedge
order. Contraction gives the formulas. For example
\(2\pi\iota_{V_+}\omega_{\rm sc}=(q_-^{-1})^*\dd\vartheta/(2\pi)\).
Its nonzero primitive period excludes exactness over the reals.
The Poincar\'e transition cocycle is the perfect lattice pairing
in the dual convention specified in Section~\ref{p1b:sec:datum}.
\end{proof}

On a local vector/dual evaluation chart the exchanged one-forms satisfy
\(q\,dp=d(pq)-p\,dq\). The resulting contact lift gives the local
Legendre graph, with its additive constant and nondegeneracy stated in
Proposition~\ref{p1b:prop:evaluation-legendre-contact-lift}.


\begin{theorem}[Integral eigenphase cover]
\label{p1b:thm:eigenphase}
The map
\[
 Q:\T^2\to\T^2,\quad(z_t,z_E)\mapsto(z_tz_E,z_tz_E^{-1})
\]
is a two-sheeted covering with kernel \(\{(1,1),(-1,-1)\}\).
Its angle matrix and character pullback are
\[
 A=\begin{pmatrix}1&1\\1&-1\end{pmatrix},\quad A^2=2I,\qquad
 Q^*(r,s)=(r+s,r-s).
\]
For the standard target orientation, \(Q^*\eta_\pm=-2\eta\).
Its image lattice is \(n+w\equiv0\pmod2\), giving
\[
 0\longrightarrow\Z^2\xrightarrow{Q^*}\Z^2
 \xrightarrow{(n,w)\mapsto n+w\bmod2}\Z_2\longrightarrow0.
\]
Exchange fixes \(q_+\) and inverts \(q_-\). For local eigenangles,
\[
 -i\partial_\alpha=\tfrac12(N_t+N_E),\qquad
 -i\partial_\beta=\tfrac12(N_t-N_E),\quad
 N_t=-i\partial_{\theta_t},\quad N_E=-i\partial_{\theta_E}.
\]
\end{theorem}
\begin{proof}
The matrix has determinant \(-2\). Solve
\(z_tz_E=z_tz_E^{-1}=1\) for the kernel. Its transpose gives the character
map; its inverse \(A/2\) gives the derivatives. The top class pulls back
by \(\det A\). Thus \(\eta\) cannot descend through \(Q\).
\end{proof}

\begin{corollary}[Relative characters and deck grading]
\label{p1b:cor:relative-characters}
A character descends through the diagonal scalar quotient precisely when
\(n+w=0\). Its primitive generator is
\[
 N_{\rm rel}=\tfrac12(N_t-N_E),\qquad
 N_{\rm rel}=N_t=-N_E\quad\hbox{on descended functions}.
\]
The Laurent character algebra is strongly \(\Z_2\)-graded by \(n+w\bmod2\).
Its even subalgebra is \(\C[q_+^{\pm1},q_-^{\pm1}]\).
\end{corollary}
\begin{proof}
Diagonal translation multiplies a character by \(e^{i(n+w)s}\).
Differentiate \(\psi=\theta_t-\theta_E\) while holding the mean angle fixed.
Multiplication adds lattice labels; an odd invertible character shifts
the two cosets into one another.
\end{proof}

\section{Hopf geometry, topology, and global bundles}
\label{p1b:sec:organized-topology}

\subsection{Bi-Hopf geometry and two groupoid structures}
\label{p1b:sec:bihopf}
\begin{theorem}[Marked bi-Hopf principal bundle]
\label{p1b:thm:bihopf}
Let \(P_v=S_x^3\times S_p^3\). Choose a Hopf circle in each factor;
use \((g_x,g_p)\cdot(z,w)=(g_xh_x(z),h_p(w)g_p)\).
This is a free faithful \(\T^2\)-action and
\[
 \T^2\longrightarrow P_v\xrightarrow{\pi}S_x^2\times S_p^2
\]
is a principal bundle. In primitive base classes \(h_x,h_p\), its Chern
matrix is \(\diag(\varepsilon_x,\varepsilon_p)\), with
\(\varepsilon_i\in\{\pm1\}\) fixed by orientations.
For normalized connections \(\alpha_i\),
\(\oint\alpha_i=2\pi\) and
\(\int\alpha_i\wedge\dd\alpha_i/(2\pi)^2=\varepsilon_i\).
\end{theorem}
\begin{proof}
Cancellation in each sphere proves freeness and faithfulness.
Quotients and local trivializations are products of Hopf quotients.
Each factor has unit Chern number with its chosen sign and no mixed
base dependence. Fibre integration proves the three-form integral.
On a fibre, the sums and differences of the Hopf connection forms
pull back to the sums and differences of scalar angular forms.
\end{proof}

\begin{theorem}[Primitive maximal-torus circle quotients]
\label{p1b:thm:primitive-circle-quotients}
Fix oriented primitive Hopf circles in the two sphere factors and write
\[
 T^4_{\max}=\U(1)_{H,x}\times\U(1)_{H,p}\times\T^2_{\rm sc}.
\]
A primitive cocharacter is a vector
\(w=(a,b,m,n)\in\Z^4\) with
\(\gcd(a,b,m,n)=1\).  Let \(S^1_w\) act on \(M_8\) by right
translation through the corresponding maximal-torus subgroup and set
\(Q_w=M_8/S^1_w\).  Put \(d=\gcd(m,n)\), with \(d=0\) when
\(m=n=0\).  Then the action is free and
\[
 0\longrightarrow\pi_2(Q_w)\longrightarrow\Z
 \xrightarrow{\;1\mapsto(m,n)\;}\Z^2
 \longrightarrow\pi_1(Q_w)\longrightarrow0,
\]
so
\[
\begin{array}{c|c|c}
 d&\pi_2(Q_w)&\pi_1(Q_w)\\ \hline
 0&\Z&\Z^2\\
 d>0&0&\Z\oplus\Z/d.
\end{array}
\]
Moreover \(\pi_k(Q_w)\cong\pi_k(M_8)\) for \(k\ge3\), and the
quotients have the following geometric normal forms.
\begin{enumerate}[label=(\alph*)]
\item If \(d=0\), then \(\gcd(a,b)=1\) and
\[
 Q_w\cong T^{a,b}_{\rm top}\times\T^2_{\rm sc},\qquad
 T^{a,b}_{\rm top}:=(S_x^3\times S_p^3)/S^1_{(a,b)}
 \cong S^2\times S^3.
\]
For the bi-Hopf orientation convention of
Theorem~\ref{p1b:thm:bihopf}, the residual circle bundle over
\(S_x^2\times S_p^2\) has primitive first Chern class
\(b\varepsilon_xh_x-a\varepsilon_ph_p\), up to reversal of the
residual fibre orientation.
\item If \(d=1\), then
\[
 Q_w\cong S^3\times S^3\times S^1.
\]
\item If \(d>1\), a scalar \(GL(2,\Z)\) change sends \((m,n)\) to
\((d,0)\), and
\[
 Q_w\cong
 \bigl((S^3\times S^3)/\mu_d^{(a,b)}\bigr)\times S^1,
\]
where, after the same convention choice in both Hopf factors,
\[
 (q_x,q_p)\cdot\zeta=
 (q_xh_x(\zeta)^a,q_ph_p(\zeta)^b),\qquad \zeta^d=1.
\]
The residual \(\mu_d\)-action is faithful and free.
\end{enumerate}
\end{theorem}
\begin{proof}
A primitive cocharacter is an embedded circle in the torus, hence right
translation is free.  The fibration
\(S^1_w\to M_8\to Q_w\), together with
\(\pi_2(M_8)=0\) and \(\pi_1(M_8)=\Z^2\), gives the displayed exact
sequence.  The map on fundamental groups is exactly scalar projection,
namely \(1\mapsto(m,n)\).  Its kernel and Smith normal form give the
listed \(\pi_2\) and \(\pi_1\); higher homotopy follows from the same
long exact sequence.

If \(d=0\), the scalar torus is a spectator and primitivity is
\(\gcd(a,b)=1\).  Quotienting the bi-Hopf fibre \(\T^2\) by the
primitive subcircle leaves a principal circle over
\(S^2\times S^2\); the primitive annihilating character is represented
by \((b,-a)\), which gives the stated Chern class, with the sign changed
if the second Hopf action is converted between left and right
conventions.  The homotopy sequence makes the five-manifold simply
connected, and the Gysin sequence gives \(H_2\cong\Z\).  Its tangent
bundle is the pullback of \(T(S^2\times S^2)\) plus a trivial vertical
line, so it is spin.  The Smale--Barden classification
\cite{Smale1962FiveManifolds,Barden1965FiveManifolds} therefore gives
\(T^{a,b}_{\rm top}\cong S^2\times S^3\).

For \(d=1\), choose a scalar lattice basis in which the scalar weights
are \((1,0)\).  The first scalar coordinate supplies a global gauge:
absorbing its phase into the two Hopf coordinates leaves two unrestricted
three-spheres and the untouched scalar circle.  This gives the stated
diffeomorphism.  For \(d>1\), the same gauge slice meets each orbit with
a residual \(d\)-th-root ambiguity.  The residual action is the displayed
\(\mu_d\)-action.  Since \(\gcd(a,b,d)=1\), its kernel is trivial; as a
restriction of the original free circle action it is free.
\end{proof}

\begin{corollary}[Exchange-stable primitive slopes]
\label{p1b:cor:exchange-stable-slopes}
Assume the two oriented Hopf embeddings are matched by the common-model
identifications of Definition~\ref{p1b:def:reciprocal-product}: both are
the same oriented cocharacter after transport to \(K\), and both quotient
actions use the right-translation convention of
Theorem~\ref{p1b:thm:primitive-circle-quotients}.
Born exchange acts on primitive maximal-torus weights by
\[
 \sigma(a,b,m,n)=(b,a,n,m).
\]
The circle subgroup \(S^1_w\) is setwise exchange-stable if and only if
\(\sigma(w)=\pm w\).  Hence its primitive weights are exactly the two
families
\[
 (a,a,m,m),\qquad (a,-a,m,-m),\qquad \gcd(a,m)=1.
\]
For a pure-Hopf circle, \((m,n)=(0,0)\), this forces
\(|a|=|b|=1\).  Thus, up to overall circle orientation and the choice of
orientation of either marked Hopf axis, the exchange-stable primitive
pure-Hopf quotient is the conventional \(T^{1,1}\) slope.  For a
scalar-only circle, the diagonal and relative slopes are the two
exchange-stable possibilities.
\end{corollary}
\begin{proof}
Exchange permutes the two Hopf weights and the two scalar weights.  A
primitive rank-one sublattice of \(\Z^4\) is preserved precisely when
its primitive generator is sent to itself or its negative.  Solving
\((b,a,n,m)=\pm(a,b,m,n)\) gives the two displayed eigensublattices.
For \(m=n=0\), primitivity reduces to \(|a|=1\), and hence also
\(|b|=1\).
\end{proof}

\begin{remark}[Independent Hopf markings and the exchange hypothesis]
\label{p1b:rem:exchange-matched-hopf-markings}
The preceding circle-quotient theorem permits independently marked axes;
its topology does not require the additional exchange hypothesis.  That
hypothesis is necessary for the displayed weight permutation.  For example,
in common-model coordinates let \(H_x=H\) and
\(H_p=uHu^{-1}\ne H\).  The subgroup of weight \((1,1,0,0)\) has first
projection \(H\), whereas its exchange has first projection
\(uHu^{-1}\), so it is not exchange stable.  Conjugating a marking also
transports the exchange if an exchange-adapted comparison is intended.
The scalar-only parity corollary below uses no Hopf-axis hypothesis.
\end{remark}


\begin{corollary}[Born-parity scalar-circle rigidity]
\label{p1b:cor:scalar-born-parity-rigidity}
Restrict Corollary~\ref{p1b:cor:exchange-stable-slopes} to the scalar
cocharacter lattice $\mathbb Z^2_{\rm sc}$, on which exchange sends
$(m,n)$ to $(n,m)$.  Up to reversal of circle orientation, the unique
primitive $+1$ eigendirection is
\[
 (m,n)=(1,1),\qquad U(1)_{\rm diag},
\]
and the unique primitive $-1$ eigendirection is
\[
 (m,n)=(1,-1),\qquad U(1)_{\rm rel}.
\]
Thus a scalar circle with specified Born parity has no further primitive
slope selector.
\end{corollary}
\begin{proof}
For the $+1$ eigenspace, $(n,m)=(m,n)$ gives $m=n$; primitivity gives
$(m,n)=\pm(1,1)$.  For the $-1$ eigenspace, $(n,m)=-(m,n)$ gives
$n=-m$; primitivity gives $(m,n)=\pm(1,-1)$.  The remaining sign reverses
only the circle orientation.
\end{proof}

\begin{remark}[Construction-type firewall]
\label{p1b:rem:circle-quotient-firewall}
The diffeomorphism \(Q_w\cong S^3\times S^3\times S^1\) in the
primitive-scalar regime does not identify \(Q_w\) with either the
embedded hypersubgroups of Corollary~\ref{p1b:cor:codim-one-normal} or
the scalar-diagonal quotient \(X_7\) of
Theorem~\ref{p1b:thm:relative}; these are distinct atlas constructions
with the same smooth model.
\end{remark}


\subsubsection{Contact and prequantization geometry of primitive pure-Hopf quotients}
\label{p1b:sec:g17}

\begin{theorem}[Residual connection and contact criterion]
\label{p1b:thm:g17-contact}
Let \((a,b)\in\Z^2\) be primitive and let
\[
 T^{a,b}_{\rm top}=(S_x^3\times S_p^3)/S^1_{(a,b)}
\]
be the pure-Hopf quotient of
Theorem~\ref{p1b:thm:primitive-circle-quotients}.  On the marked Hopf
fibre \(\T^2\), orient the residual circle by the primitive character
\[
 \chi_{a,b}(z_x,z_p)=z_x^{\,b}z_p^{-a}.
\]
If \(\alpha_x,\alpha_p\) are the normalized Hopf connections of
Theorem~\ref{p1b:thm:bihopf}, then
\[
 \widetilde\eta_{a,b}=b\alpha_x-a\alpha_p
\]
is basic for \(S^1_{(a,b)}\) and descends to a normalized principal-circle
connection \(\eta_{a,b}\) on
\[
 S^1\longrightarrow T^{a,b}_{\rm top}
 \longrightarrow S_x^2\times S_p^2.
\]
Its curvature and first Chern class are
\[
 \dd\eta_{a,b}=bF_x-aF_p,
 \qquad
 c_1(T^{a,b}_{\rm top})
 =b\varepsilon_xh_x-a\varepsilon_ph_p.
\]
Moreover
\[
 \boxed{\eta_{a,b}\text{ is a contact form}\iff ab\ne0.}
\]
With the displayed residual orientation,
\[
 \boxed{
 \frac1{(2\pi)^3}
 \int_{T^{a,b}_{\rm top}}
 \eta_{a,b}\wedge(\dd\eta_{a,b})^2
 =-2ab\,\varepsilon_x\varepsilon_p.}
\]
Reversing the residual fibre orientation sends
\(\eta_{a,b}\mapsto-\eta_{a,b}\) and
\(c_1\mapsto-c_1\), but leaves both the contact criterion and the
fibre-first normalized characteristic number
\(\langle c_1^2,[S^2\times S^2]\rangle\) unchanged.  If the
orientation of the total space is artificially held fixed instead,
the raw contact-form integral changes sign; this is a different
integration convention.
\end{theorem}
\begin{proof}
The infinitesimal generator of the quotient subcircle is \((a,b)\), so
\(b\alpha_x-a\alpha_p\) annihilates it and is invariant under the torus.
Thus it descends.  On any integral lift \((r,s)\) of the residual circle
with \(br-as=1\), its fibre integral is \(2\pi\), proving normalization.
The curvature and Chern class follow from
\([F_i/(2\pi)]=\varepsilon_i h_i\).
Because \(F_x^2=F_p^2=0\) on the two sphere factors,
\[
 (bF_x-aF_p)^2=-2ab\,F_x\wedge F_p.
\]
A principal-circle connection on a five-manifold is contact exactly when
its curvature is nondegenerate on the four-dimensional horizontal space,
which here is equivalent to \(ab\ne0\).  Fibre integration gives
\[
 \frac1{(2\pi)^3}\int\eta_{a,b}(\dd\eta_{a,b})^2
 =\int_{S^2\times S^2}c_1^2
 =-2ab\varepsilon_x\varepsilon_p.
\]
To check the orientation convention explicitly, write
\(\eta'=-\eta\), \(F'=-F\) and orient the total space by the
new residual fibre orientation followed by the same base orientation.
The induced total orientation is reversed, so
\[
 \int_{-T}\eta'\wedge(\dd\eta')^2
 =\int_{-T}(-\eta)\wedge(\dd\eta)^2
 =\int_T\eta\wedge(\dd\eta)^2.
\]
Equivalently \(c_1'^2=(-c_1)^2=c_1^2\).  No sign flip of this
characteristic number occurs under a consistent reorientation.
\end{proof}

\begin{corollary}[Regular prequantization branch]
\label{p1b:cor:g17-prequantization}
If \(ab\ne0\), then
\[
 \omega_{a,b}=bF_x-aF_p
\]
is an integral symplectic form on \(S_x^2\times S_p^2\), and
\((T^{a,b}_{\rm top},\eta_{a,b})\) is its regular
Boothby--Wang prequantization circle bundle.  In particular every
exchange-stable primitive pure-Hopf slope of
Corollary~\ref{p1b:cor:exchange-stable-slopes} lies in the contact branch.
For \(a=0\) or \(b=0\), the residual curvature has rank two and the
connection is not contact.
\end{corollary}
\begin{proof}
The square computed in Theorem~\ref{p1b:thm:g17-contact} is a nowhere-zero
multiple of the product area form exactly when \(ab\ne0\).  Its integral
class is the first Chern class of the residual circle bundle, so the regular
prequantization statement follows.  Exchange-stable primitive pure-Hopf
slopes have \(|a|=|b|=1\).
\end{proof}

The quotient topology, including
\(T^{a,b}_{\rm top}\cong S^2\times S^3\), remains owned by
Theorem~\ref{p1b:thm:primitive-circle-quotients}.  The longer metric
factorization and harmonic-morphism check is supplementary and is recorded
in Subsubsection~\ref{p1b:app:g17}; it shows that the quotient refinement is
compatible with the product Hopf Riemannian submersion.


\begin{proposition}[Bi-Hopf transgression]
\label{p1b:prop:bihopf-transgression}
For fibre generators \(s_x,s_p\) in the Serre sequence,
\[
 d_2s_x=\varepsilon_xh_x,\qquad d_2s_p=\varepsilon_ph_p,\qquad
 d_2(s_xs_p)=\varepsilon_xh_xs_p-\varepsilon_ps_xh_p\ne0.
\]
In particular \(H^2(P_v;\Z)=0\); the primitive two-class on the
independent scalar torus is not a two-class on \(P_v\).
\end{proposition}
\begin{proof}
Transgression sends a fibre circle generator to its Chern class.
Apply this factorwise and use the graded Leibniz rule. The two terms
lie in distinct free summands. K\"unneth directly also gives
\(H^2(S^3\times S^3;\Z)=0\).
\end{proof}

\begin{theorem}[Pair and bi-Hopf kernel-pair groupoids]
\label{p1b:thm:groupoids}
The manifold \(M_8\) has two groupoid presentations.
For \(K\times K\rightrightarrows K\),
\[
 s(a,b)=b,\quad t(a,b)=a,\quad
 (a,b)\circ(b,c)=(a,c),\quad(a,b)^{-1}=(b,a).
\]
Here \(r(a,b)=ab^{-1}\) is a multiplicative \(K\)-valued cocycle.
For the bi-Hopf bundle, the map
\[
 P_v\times\T^2\longrightarrow P_v\times_{S^2\times S^2}P_v,\qquad
 (q,z)\mapsto(q,q\cdot z)
\]
is a diffeomorphism identifying the transformation and kernel-pair
groupoids. Their arrow manifold is again \(M_8\).
\end{theorem}
\begin{proof}
The pair maps satisfy the groupoid laws and
\((ab^{-1})(bc^{-1})=ac^{-1}\). Two points of a principal fibre differ by
a unique group element, smooth in local trivializations. This proves the
second statement. The object manifolds \(K\) and \(P_v\) are distinct.
\end{proof}

\begin{theorem}[Convolution algebras and tangent deformation]
\label{p1b:thm:convolution}
Pair-groupoid smooth kernel convolution is
\[
 (f*g)(x,z)=\int_Kf(x,y)g(y,z)\,\dd\mu(y),\qquad
 f^*(x,y)=\overline{f(y,x)}.
\]
Its operator-norm completion is \(\mathcal K(L^2K)\). The bi-Hopf
groupoid algebra is Morita equivalent to \(C(S^2\times S^2)\). Thus
\[
\begin{array}{c|cc}
\text{algebra}&K_0&K_1\\\hline
C(M_8)&\Z^8&\Z^8\\
\mathcal K(L^2K)&\Z&0\\
C^*(P_v\rtimes\T^2)&\Z^4&0
\end{array}
\]
The tangent groupoid of \(K\) has positive-parameter fibres
\(\mathcal K(L^2K)\) and zero fibre \(C_0(T^*K)\).
\end{theorem}
\begin{proof}
Smooth kernels are compact and include all finite sums
\(u(x)\overline{v(y)}\) with smooth \(u,v\); these are norm dense in the
compact operators. The free proper action groupoid is equivalent to
its orbit space, giving Morita equivalence \cite{SimsWilliams}.
The two-sphere Bott classes give the K-groups of the orbit space;
Theorem~\ref{p1b:thm:k-theory} gives those of \(M_8\).
For the tangent deformation use
\(TK\times\{0\}\sqcup(K\times K)\times(0,1]\), with addition at
zero and exponential-coordinate gluing. Fibrewise Fourier transform at
zero gives \(C_0(T^*K)\) \cite{Connes1994}.
\end{proof}

\begin{theorem}[Born involution through the tangent deformation]
\label{p1b:thm:tangent-born-involution}
Pair-groupoid inversion extends to the tangent deformation of
Theorem~\ref{p1b:thm:convolution}.  On the zero fibre it is
\[
 \mathcal I_0(q,v)=(q,-v),\qquad (q,v)\in TK.
\]
After fibrewise Fourier transform this becomes cotangent fibre reversal
\[
 R(q,\xi)=(q,-\xi),\qquad (q,\xi)\in T^*K.
\]
For the canonical Liouville one-form \(\lambda\) and canonical symplectic
form \(\omega_{\rm can}=\dd\lambda\),
\[
 \boxed{R^*\lambda=-\lambda,\qquad R^*\omega_{\rm can}=-\omega_{\rm can}.}
\]
Thus Born inversion is anti-symplectic on the cotangent zero fibre.  If
\(R^*f=\epsilon_f f\) and \(R^*g=\epsilon_g g\), with
\(\epsilon_f,\epsilon_g\in\{\pm1\}\), then
\[
 R^*\{f,g\}_{\rm can}=-\epsilon_f\epsilon_g\{f,g\}_{\rm can}.
\]
The compatibility assertion refers to this cotangent zero fibre.
\end{theorem}
\begin{proof}
The tangent groupoid is the deformation to the normal cone of the diagonal
\(\Delta_K\subset K\times K\).  The swap \((a,b)\mapsto(b,a)\) preserves
\(\Delta_K\), hence extends to the deformation.  On the normal quotient,
represented by the difference class \([v_1,v_2]\mapsto v_1-v_2\), swapping
factors sends the class to its negative.  Therefore the zero-fibre action is
\((q,v)\mapsto(q,-v)\).  Fibrewise Fourier transform converts vector
reversal into covector reversal.  Finally, for
\(\lambda_{(q,\xi)}(X)=\xi(\dd\pi(X))\), one has
\((R^*\lambda)_{(q,\xi)}(X)=(-\xi)(\dd\pi(X))=-\lambda_{(q,\xi)}(X)\),
so \(R^*\omega_{\rm can}=-\omega_{\rm can}\).  The Poisson-parity identity
follows from anti-symplecticity.
\end{proof}

The simultaneous double-groupoid compatibility test is recorded in
Paper~I-D, Section~\ref{p1b:app:obstructions}.

\subsection{Integral topology and exchange lattices}
\label{p1b:sec:topology}
\begin{theorem}[Cohomology, invariant forms and rational model]
\label{p1b:thm:cohomology}
One has
\[
 H^*(M_8;\Z)=\Lambda_\Z(a_t,a_E,u_x,u_p),\quad
 P(z)=(1+z)^2(1+z^3)^2,
\]
and Betti numbers \((1,2,1,2,4,2,1,2,1)\).
The sixteen wedge monomials in normalized scalar and Cartan forms
are all bi-invariant forms. Their differential is zero and they compute
de Rham cohomology. The rational minimal model is the same exterior
algebra over \(\Q\) with zero differential. In particular \(M_8\) is formal,
and
\[
 \pi_1(M_8)=\Z^2,\qquad
 \pi_n(M_8)=\pi_n(S^3)\oplus\pi_n(S^3)\quad(n\ge2).
\]
\end{theorem}
\begin{proof}
Integral K\"unneth applies to the four factors with free cohomology.
The adjoint-invariant exterior algebra of \(\mathfrak{su}(2)\) has only
degrees zero and three; the center contributes its full exterior algebra.
Independent adjoint actions give their tensor product. Normalized forms
represent the K\"unneth basis. The product of the circle and
odd-sphere rational models gives the model displayed. Product homotopy
gives the last identities. Only degrees one and three have nonzero rational
homotopy, each of rank two; integral higher sphere homotopy is retained.
\end{proof}

\begin{theorem}[Middle-dimensional form and integral exchange defect]
\label{p1b:thm:intersection}
In the basis \((a_tu_x,a_tu_p,a_Eu_x,a_Eu_p)\), the intersection form is
\[
 Q_4=\begin{pmatrix}0&0&0&-1\\0&0&1&0\\0&1&0&0\\-1&0&0&0\end{pmatrix}
 \cong H\oplus H.
\]
For \(c=Aa_tu_x+Ba_tu_p+Ca_Eu_x+Da_Eu_p\),
\(c^2=2(BC-AD)a_ta_Eu_xu_p\). Exchange eigensublattices have ranks
and integral saturation quotients
\[
\begin{array}{c|rrrrrrrrr}
d&0&1&2&3&4&5&6&7&8\\\hline
\rank H^d_+&1&1&0&1&2&1&0&1&1\\
\rank H^d_-&0&1&1&1&2&1&1&1&0\\
\dim_{\F_2}H^d/(H^d_++H^d_-)&0&1&0&1&2&1&0&1&0
\end{array}
\]
The total quotient is \((\Z_2)^6\).
\end{theorem}
\begin{proof}
Every generator is odd. Anticommutation gives the intersection matrix and
square formula. Each two-element signed exchange orbit of basis monomials
contributes one even and one odd generator spanning an index-two sublattice.
Fixed monomials contribute their sign without an index defect.
List the sixteen monomials to obtain the table.
In degree four the even basis has form \(\diag(-2,2)\), the odd basis
\(\diag(2,-2)\), and their combined lattice has index four.
\end{proof}

\begin{proposition}[Relative primitive classes]
\label{p1b:prop:relative-cohomology}
For primitive \(a\in H^1(K;\Z)\), \(u\in H^3(K;\Z)\),
\[
 r^*a=a_t-a_E,\qquad r^*u=u_x-u_p.
\]
Diagonal pullback kills the differences and sends the sums to \(2a,2u\).
\end{proposition}
\begin{proof}
The circle and sphere classes are primitive for multiplication;
inversion acts by minus one on each. Compose these pullbacks, then
restrict to the diagonal.
\end{proof}

\begin{theorem}[Finite covers and power maps]
\label{p1b:thm:covers}
Connected finite covers correspond to finite-index sublattices
\(\Lambda\subset\Z^2\), with total space
\(S^3\times S^3\times(\R^2/\Lambda)\). Index \(d\) gives degree \(d\);
there are \(\sum_{c\mid d}c\) marked sublattices.
A based exchange lift exists precisely when scalar swap preserves
\(\Lambda\). For simultaneous powers,
\[
 P_k^*(a_t,a_E,u_x,u_p)=k(a_t,a_E,u_x,u_p),\qquad \deg P_k=k^4.
\]
Independent sphere powers \(r,s\) and a scalar integer matrix \(L\)
have degree \(rs\det L\).
\end{theorem}
\begin{proof}
The universal cover is \(S^3\times S^3\times\R^2\).
Covering theory gives classification and the lift criterion;
Hermite normal form gives the count. Primitivity of the four odd
generators gives multiplication by \(k\) under group powers,
also for negative \(k\) by inversion. Apply this to the orientation
monomial. For \(|k|>1\), sphere power maps have critical strata
and are not covering maps.
\end{proof}

\begin{theorem}[Stable splitting]
\label{p1b:thm:stable}
With identity basepoints,
\[
 \Sigma^\infty(M_8)_+\simeq
 \bigvee_{S\subset\{1,2,3,4\}}\Sigma^{d_S}\mathbb S,\qquad
 d_S=\sum_{i\in S}(1,1,3,3)_i.
\]
For every generalized cohomology theory \(E\), additively
\[
 E^q(M_8)\cong\bigoplus_{d=0}^8 E^{q-d}(\mathrm{pt})^{b_d}.
\]
\end{theorem}
\begin{proof}
Split each \(\Sigma^\infty(S^n)_+\) into the sphere spectrum and its
\(n\)-fold suspension, then distribute the four smash products.
Apply the cohomology functor. This is a stable additive splitting.
\end{proof}

\subsection{Bundles, spin structures and integral index pairings}
\label{p1b:sec:bundles}
\begin{theorem}[Framing, spin and higher lifts]
\label{p1b:thm:spin}
Left translation trivializes \(TM_8\). All positive-degree tangent
Stiefel--Whitney and Pontryagin classes vanish, as does its Euler class.
The four spin structures form an \(H^1(M_8;\Z_2)=\Z_2^2\)-torsor.
Relative to the product framing, exchange swaps the two spin bits;
two of the four classes are fixed.
Spin-c structures form a \(\Z\)-torsor and their determinant classes are
\(2n\eta\). Each fixed spin structure admits a \(\Z^2\)-torsor of string
structures. Every spin structure extends over
\(D^4\times S^3\times T^2\).
\end{theorem}
\begin{proof}
A basis at the identity gives a global left-invariant frame.
The spin obstruction \(w_2\) is zero, and spin structures are classified
by \(H^1(-;\Z_2)\). A line twist changes a Spin-c determinant by \(2c_1\);
the absence of two-torsion in \(H^2\) makes this label complete.
All spin structures induce the same Spin-c class since their differences
map to zero under the integral Bockstein.
The string obstruction \(p_1/2\) is zero: its double is zero and \(H^4\)
is torsion-free. The lifts form an \(H^3\)-torsor \cite{Redden}.
The unique spin structure of a sphere extends across the four-ball,
while the other factors and torus spin structure are unchanged.
\end{proof}

\begin{theorem}[Complex K-theory and its integral character lattice]
\label{p1b:thm:k-theory}
Normalize odd K-generators \(\xi_t,\xi_E,\xi_x,\xi_p\) on the two
circles and spheres. Then
\[
 K^*(M_8)=\Lambda_\Z(\xi_t,\xi_E,\xi_x,\xi_p),
 \qquad K^0(M_8)=\Z^8,\quad K^1(M_8)=\Z^8.
\]
They can be normalized so that their odd Chern characters are
\(a_t,a_E,u_x,u_p\). On this product the character sends each K-basis
monomial exactly to the corresponding integral cohomology monomial.
\end{theorem}
\begin{proof}
Bott periodicity gives one odd generator on each factor, with square zero.
The torsion-free K\"unneth theorem gives the exterior tensor product.
The normalized generator on \(S^3\) is the suspension of the Bott
generator on \(S^2\); its odd Chern character is the primitive three-class.
Multiplicativity proves the integral monomial claim \cite{HatcherK}.
\end{proof}

\begin{theorem}[Stable Chern congruences and index]
\label{p1b:thm:chern}
Write \(\mathcal O=a_ta_Eu_xu_p\),
\[
 c_1=n\eta,\qquad c_2=Aa_tu_x+Ba_tu_p+Ca_Eu_x+Da_Eu_p.
\]
The complete stable integrality conditions are
\[
 c_3=2m\,u_xu_p,\qquad
 c_4=(BC-AD+2nm-6s)\mathcal O,\qquad
 n,A,B,C,D,m,s\in\Z.
\]
For any chosen Riemannian metric and spin structure, the twisted spin
Dirac index is \(\operatorname{ind}D_E^+=s\).
A stable class fixed by exchange has \(n=m=0\),
\[
 c_2=A(a_tu_x+a_Eu_p)+B(a_tu_p+a_Eu_x),\quad
 c_4=(B^2-A^2-6s)\mathcal O.
\]
\end{theorem}
\begin{proof}
The ring gives \(c_1^2=c_1c_2=0\). Newton identities give
\[
 \operatorname{ch}_2=-c_2,\quad
 \operatorname{ch}_3=c_3/2,\quad
 \operatorname{ch}_4=(c_2^2+2c_1c_3-2c_4)/12.
\]
Use the square formula in Theorem~\ref{p1b:thm:intersection}
and the integral K-lattice. Conversely the virtual class with
\(\operatorname{ch}=r+n\eta-c_2+mu_xu_p+s\mathcal O\) exists by that
lattice and becomes an actual bundle after adding sufficiently many
trivial summands. Thus sufficiency is stable, not a fixed-rank assertion.
The tangent bundle is trivial, so Atiyah--Singer reduces the index to
\(\int\operatorname{ch}_4=s\) \cite{AtiyahSinger1963}.
Exchange reverses \(\eta,u_xu_p\) and has the displayed even degree-four
basis. Fixed stable classes do not by this argument acquire an equivariant
lift to a particular bundle; Theorem~\ref{p1b:thm:g21-exotic-lift}
constructs one independently for the specific nontrivial G15 pullback class.
\end{proof}


\subsubsection{A fixed-rank sector invisible to stable invariants}
\label{p1b:sec:g15}

\begin{theorem}[The two \(SU(2)\) bundles on the canonical derived core]
\label{p1b:thm:g15-unstable}
Let
\[
 D_6=[M_8,M_8]=SU(2)_x\times SU(2)_p
 \cong S^3_x\times S^3_p,
 \qquad P_v:=D_6.
\]
Principal \(SU(2)\) bundles over \(P_v\), up to bundle isomorphism,
have exactly two classes, indexed by
\(\epsilon_{\mathrm{IB15}}\in\Z/2\).  Write
\(E_0\cong\C^2\) for the zero class and \(E_{\rm ex}\) for the nonzero
class.  The exotic rank-two complex bundle obeys
\[
 c_j(E_{\rm ex})=0\quad(j>0),\qquad [E_{\rm ex}]-2=0\ \hbox{in }K^0(P_v),
 \qquad E_{\rm ex}\oplus\C\cong\C^3,\qquad
 E_{\rm ex}\not\cong\C^2.
\]
For the projection \(\pi:M_8=P_v\times\T^2\to P_v\), the bundle
\(\pi^*E_{\rm ex}\) is still nontrivial and becomes trivial after adding
one trivial complex line. Both classes are fixed as isomorphism classes
by exchange.
\end{theorem}
\begin{proof}
The \(6\)-dimensional domain sees only the \(S^4\) part of
\(BSU(2)=\mathbb H P^\infty\). Stabilization and the suspended product
splitting give
\[
 [P_v,BSU(2)]\cong[P_v,S^4]
 \cong\widetilde\pi_s^{\,4}(P_v)
 \cong\pi_6(S^4)\cong\Z/2.
\]
Paper~I-D, Section~\ref{p1b:app:g15} checks the dimension ranges and constructs
the nonzero class explicitly. Since \(H^4(P_v;\Z)=0\), \(c_2=0\);
special unitarity gives \(c_1=0\), and higher Chern classes vanish by
rank. The torsion-free group \(K^0(P_v)\cong\Z^2\) and the rational
Chern character imply \([E_{\rm ex}]-2=0\). Stable classification is already
injective for rank three over a six-dimensional CW complex, proving
\(E_{\rm ex}\oplus\C\cong\C^3\). Projection has a section, so pullback cannot
make the nontrivial bundle trivial. Exchange acts by degree \(-1\) on
the smash sphere \(S^3\wedge S^3\), which acts trivially on \(\Z/2\).
\end{proof}

\begin{proposition}[The rank boundary for stable classification]
\label{p1b:prop:g15-stable-range}
For every \(r\ge4\), stabilization identifies isomorphism classes of
rank-\(r\) complex bundles on \(M_8\) bijectively with
\(\widetilde K^0(M_8)\). Thus the stable integral character lattice,
with rank fixed to \(r\), is a complete classification in that range.
For rank two, even vanishing Chern classes and stable triviality do
not classify bundles. The pullback sector classified above is exactly
the two classes in the image of \(\pi^*\).
\end{proposition}
\begin{proof}
The inclusion \(BU(r)\to BU\) is \((2r+1)\)-connected, obtained from
the fibrations \(U(r+1)/U(r)=S^{2r+1}\).
Obstruction theory therefore makes stabilization bijective on a CW
domain of dimension at most \(2r\). Apply this first to \(M_8\), of
dimension eight, and also to rank three on \(P_v\) in the preceding
proof. Injectivity of \(\pi^*\) and the two-element source determine
the stated image.
\end{proof}
\subparagraph{Canonical-carrier and family boundary.}
The carrier \(P_v\) used here is not an independently selected subgroup:
it is the characteristic derived subgroup \(D_6=[M_8,M_8]\) of the
completed reciprocal host.  Thus this theorem classifies the complete
ordinary principal-\(SU(2)\) bundle family attached to that canonical core,
\[
 \mathfrak B_{\mathrm{IB15}}
 =\{E_0,E_{\rm ex}\}
 \cong\Z/2.
\]
The mathematical theorem does not select a value of
\(\epsilon_{\mathrm{IB15}}\).  The framework bridge retains the complete
family and introduces a member selector only for a downstream claim that
requires one actual bundle.  This bundle-sector notation is distinct from
the separately named framework G15 flavor/source obligation.

Exchange invariance of the isomorphism class alone does not construct a
coherent involutive exchange lift.  The independent equivariant construction
of Theorem~\ref{p1b:thm:g21-exotic-lift} supplies such a lift for the
nontrivial pullback class.  This section still does not classify all rank-two
or rank-three bundles on \(M_8\); it specifies the complete pullback sector
and the range in which the stable classification is complete.


\subsubsection{Actual low-rank complex bundles on \(M_8\)}
\label{p1b:sec:g20}

Write, as in Theorem~\ref{p1b:thm:chern},
\[
 c_1=n\eta,
 \qquad
 c_2=Aa_tu_x+Ba_tu_p+Ca_Eu_x+Da_Eu_p,
 \qquad
 \Delta=BC-AD.
\]
The stable theorem does not by itself determine fixed-rank desuspension or
unstable fixed-Chern multiplicity.  The results below close the rank-two
existence problem for both determinant parities, retain the independently
verified rank-three classification, and give the fixed-Chern multiplicity
in the principal-\(SU(2)\) category used by the Born-fixed count below.
This category retains a determinant trivialization and its isomorphisms
preserve that trivialization; forgetting it requires a further quotient.

\begin{theorem}[Rank-two Chern-data existence]
\label{p1b:thm:g20-rank2-existence}
\label{p1b:thm:g20-rank2-even}
There exists a complex rank-two bundle \(E\to M_8\) with the displayed
\(c_1,c_2\) if and only if
\[
 \boxed{
 \begin{cases}
  12\mid\Delta,& n\ \text{even},\\[1mm]
  6\mid\Delta,& n\ \text{odd}.
 \end{cases}}
\]
This is an existence theorem for the Chern pair; it is not a uniqueness
statement for the underlying rank-two bundle.
\end{theorem}
\begin{proof}
For even \(n\), tensor by a line bundle with first Chern class
\(-n\eta/2\).  Since \(\eta^2=0\), this reduces to determinant zero without
changing \(c_2\); Corollary 8.4 of \cite{CadekCrabbVanzura2008} then reduces
on the present exterior cohomology ring to \(12\mid\Delta\).

For odd \(n\), line twisting reduces to \(c_1=\eta\) without changing
\(c_2\).  The determinant line supplies an \(L\)-symplectic structure.
Paper~I-D, Section~\ref{p1b:app:g20} constructs the required stable
\(L\)-symplectic class from the four five-dimensional mixed generators and
a primitive twisted-hyperbolic top generator, represents it in complex rank
four, proves \(c_4=0\), and splits off the canonical
\(L\)-symplectic plane.  Necessity follows from the exact integral
K-character coefficient \(\operatorname{ch}_4(E)=\Delta\mathcal O/6\).
Thus the odd branch is exactly \(6\mid\Delta\).
\end{proof}

\begin{theorem}[Fixed-Chern multiplicity with a chosen determinant trivialization]
\label{p1b:thm:g20-rank2-fibre-even}
Assume \(c_1=0\) and \(12\mid\Delta\).  The objects counted are pairs
\((E,\tau)\), where \(E\) is a Hermitian complex rank-two bundle and
\(\tau:\det E\xrightarrow{\sim}\mathbf1\) is a unit determinant
trivialization; isomorphisms must preserve \(\tau\).  Equivalently these
are principal \(SU(2)\) bundles.  Restriction to
\(P_v=S_x^3\times S_p^3\) gives a bit
\(\epsilon\in\Z/2\), distinguishing the trivial and G15-exotic restriction
classes.  The two \(\epsilon\)-sectors have equal cardinality.

For a fixed \(\epsilon\), the number \(N_\epsilon(c_2)\) of these
determinant-preserving isomorphism classes with the prescribed \(c_2\) is
as follows.
\begin{enumerate}[label=\textup{(\roman*)}]
\item If \(c_2=0\), then
\[
 \boxed{N_\epsilon(0)=720.}
\]
\item If \(\Delta=0\) and \(c_2\ne0\), choose a primitive
\((a,b)\in\Z^2\) and integers \(m,n\), not both zero, with
\[
 (A,B)=m(a,b),\qquad (C,D)=n(a,b).
\]
Put
\[
 d=\gcd(|m|,|n|),\qquad r=\gcd(12,d).
\]
Then
\[
 \boxed{
 N_\epsilon(c_2)=
 \begin{cases}
 48r,& d\ \text{odd},\\
 72r,& d\equiv2\pmod4,\\
 60r,& 4\mid d.
 \end{cases}}
\]
\item If \(\Delta\ne0\), put
\[
 s_1=\gcd(A,B,C,D),\qquad s_2=\frac{|\Delta|}{s_1},
\]
\[
 q_i=\gcd(12,s_i).
\]
If some of \(A,B,C,D\) is odd, set \(\nu=4\).  If all four are even, set
\[
 \rho_2=\rank_{\F_2}
 \begin{pmatrix}
  B/2&A/2\\ D/2&C/2
 \end{pmatrix}\!\!\pmod2,
 \qquad
 \nu=4+2^{\rho_2}.
\]
Then
\[
 \boxed{N_\epsilon(c_2)=q_1q_2\nu.}
\]
\end{enumerate}
For even \(n\), fix a determinant line \(L\) with \(c_1(L)=n\eta\)
and count pairs \((E,\tau:\det E\xrightarrow{\sim}L)\), with
isomorphisms preserving \(\tau\).  Choose a line bundle \(J\) and an
identification \(L\otimes J^{\otimes2}\cong\mathbf1\).  Tensoring by
\(J\) is then an equivalence with the determinant-trivialized category
above; it leaves \(c_2\) unchanged because \(\eta^2=0\).  The same
cardinalities apply in this specified category.
\end{theorem}
\begin{proof}
Paper~I-D, Section~\ref{p1b:app:g20} verifies the level-twelve Heisenberg component
group
\([S^3\times S^3,S^3]\), the secondary module \((\Z/2)^3\), its
conjugation action, and the vanishing of the affine unframing class.  The
three displayed cases are then the resulting orbit counts: the zero-Chern
case, the rank-one degree-matrix normal form, and the full-rank Smith
quotient.  The two G15 restriction components have equivalent Postnikov
2-types, hence equal counts.
\end{proof}

\begin{remark}[Forgetting the determinant trivialization]
\label{p1b:rem:g20-determinant-forgetting}
Let \(\mathscr F^{SU(2)}_\epsilon(c_2)\) denote the fibre counted in
Theorem~\ref{p1b:thm:g20-rank2-fibre-even}.  The ordinary complex fibre is
\[
 \mathscr F^{\C}_\epsilon(c_2)
 \cong \mathscr F^{SU(2)}_\epsilon(c_2)/H^1(M_8;\Z),
 \qquad H^1(M_8;\Z)\cong\Z^2,
\]
where the action changes the determinant trivialization.  It preserves
\(c_2\) and \(\epsilon\), but need not act trivially.  In fact two
distinct principal-\(SU(2)\) classes on \(M_8\) with \(c_2=a_tu_x\)
and \(\epsilon=0\) have isomorphic underlying complex bundles.
Paper~I-D, Section~\ref{p1b:app:g20-determinant-forgetting} proves this by a
five-dimensional retract and the determinant Postnikov action.  Thus the
displayed \(N_\epsilon\) formulas do not enumerate ordinary complex
bundles.  The additional quotient, including even-determinant line
twisting, is computed in Theorem~\ref{p1b:thm:v2170-ordinary-counts}.
Its ordinary zero-Chern counts are 720 in the trivial sector and 180 in
the exotic sector; the determinant-preserving count above remains 720
in each sector.
\end{remark}

\begin{theorem}[Rank-three Chern-data existence]
\label{p1b:thm:g20-rank3-existence}
A complex rank-three bundle on \(M_8\) has
\[
 c_3=2m\,u_xu_p
\]
for an integer \(m\).  Conversely, prescribed data
\[
 c_1=n\eta,\qquad c_2\text{ as above},\qquad c_3=2m\,u_xu_p
\]
are realized by a rank-three bundle if and only if
\[
 \boxed{6\mid(\Delta+2nm).}
\]
In particular the actual rank-three existence condition agrees with the
stable rank-four condition after setting \(c_4=0\).
\end{theorem}
\begin{proof}
Yang's rank-three theorem identifies rank-three bundles with stable
rank-four classes having \(c_4=0\), subject to his two-primary and
three-primary Chern congruences.  On the spin, parallelizable manifold
\(M_8\), the first condition forces exactly the evenness of the coefficient
of \(u_xu_p\), while the remaining congruences reduce to
\(\Delta+2nm\equiv0\pmod2\) and
\(\Delta+2nm\equiv0\pmod3\).  The detailed specialization is recorded in
Paper~I-D, Section~\ref{p1b:app:g20}; together they are the displayed mod-six
condition.
\end{proof}

\begin{theorem}[Fixed-Chern rank-three fibre]
\label{p1b:thm:g20-rank3-fibre}
Assume the rank-three existence condition and put
\[
 g=\gcd(6,2m).
\]
The set of rank-three isomorphism classes with the prescribed Chern classes
is in bijection with
\[
 \mathfrak T
 =\frac{\Z^2}{g\Z^2+\langle(A,C),(B,D)\rangle}.
\]
Equivalently, if
\[
 d_1=\gcd(g,A,B,C,D),
\]
\[
 d_2=\gcd(g^2,gA,gB,gC,gD,AD-BC),
\]
then
\[
 \boxed{\mathfrak T\cong
 \Z/d_1\Z\oplus\Z/(d_2/d_1)\Z,}
\]
with the convention that a zero invariant factor denotes an infinite
cyclic factor only when the displayed lattice has lower rank.  Here
\(g\in\{2,6\}\), so the present quotient is finite.  For zero Chern
classes,
\[
 \boxed{\mathfrak T_0\cong(\Z/6\Z)^2,}
\]
so there are exactly thirty-six rank-three isomorphism classes with all
Chern classes zero.
\end{theorem}
\begin{proof}
Yang's fixed-Chern fibre is the product of a Bockstein ambiguity and a
quotient of \(H^7\).  The Bockstein term vanishes because the integral
cohomology of \(M_8\) is torsion free.  Bott integrality and the exterior
cohomology ring reduce the second term to the displayed quotient; see
Paper~I-D, Section~\ref{p1b:app:g20}.  The invariant factors are the Smith invariants
of
\[
 \begin{pmatrix}g&0&A&B\\0&g&C&D\end{pmatrix}.
\]
The zero-Chern case is immediate.
\end{proof}

\begin{remark}[Odd determinant: ordinary multiplicity]
\label{p1b:rem:g20-rank2-odd}
Theorem~\ref{p1b:thm:g20-rank2-existence} closes actual rank-two existence
for odd determinant.  The determinant-zero mapping-space calculation in
Theorem~\ref{p1b:thm:g20-rank2-fibre-even} is not silently promoted to a
complete fixed-Chern multiplicity formula for odd determinant.
The actual determinant-boundary and torus-filling calculation in
Theorem~\ref{p1b:thm:v2170-ordinary-counts} supplies that separate ordinary
classification: the restriction sector is \(\Delta/6\bmod2\), and the
admissible count is \(4Q_6\) if a coefficient is odd and \(8Q_6\) if
all are even.  In particular the zero-second-Chern count is 288.
\end{remark}

\subsubsection{The actual stabilization maps between the low ranks}
\label{p1b:sec:low-rank-stabilization}

\begin{theorem}[Rank-two stabilization in the edge and filling coordinates]
\label{p1b:thm:low-rank-stabilization}
Fix admissible rank-two Chern data \(c_1=n\eta,c_2\) and put
\[
 N_c=\begin{pmatrix}B&-A\\D&-C\end{pmatrix},\qquad
 T_6=\Z^2/(6\Z^2+N_c\Z^2).
\]
Then \(E\mapsto E\oplus\mathbf1\) is onto the ordinary rank-three
fixed-Chern fibre with \((c_1,c_2,c_3)=(n\eta,c_2,0)\).
Choose origins in the source and target fibres compatibly with this map,
and choose the target free edge generators as the images of the source
generators.  In the resulting torus charts its formula is
\[
 \boxed{(\varepsilon;z_t,z_E;\hbox{secondary filling labels})
 \longmapsto[(z_t,z_E)]\in T_6.}
\]
In particular the central edge coefficients reduce modulo six, and the
G15 restriction bit and all the finite secondary filling differences
disappear.  The formula is an equality of affine coordinates after the
displayed origin choices; it does not assert an origin independent of
the bundles and framings.
\end{theorem}
\begin{proof}
Use the block inclusion \(j:SU(2)\hookrightarrow SU(3)\).  Its relevant
homotopy maps are
\[
\begin{gathered}
 \pi_3SU(2)\xrightarrow{\sim}\pi_3SU(3),\qquad
 \pi_4SU(2)\longrightarrow\pi_4SU(3)=0,\\
 \pi_6SU(2)=\Z/12\twoheadrightarrow\pi_6SU(3)=\Z/6,
 \qquad\pi_7SU(3)=0.
\end{gathered}
\]
The middle surjection has kernel \(\{0,6\}\).  These are the standard
low-rank homotopy calculations for \(SU(2)\to SU(3)\to S^5\)
\cite{MimuraToda1963SU3SU4Sp2}; generators of the cyclic groups are
chosen so that the surjection is reduction modulo six.

Every rank-two restriction to \(P=S^3\times S^3\) becomes trivial in
rank three by Theorem~\ref{p1b:thm:g15-unstable}.  Choose such a
trivialization.  The central top-cell gauge insertion is changed by
\(j_*\), while the free sphere degrees are unchanged.  Naturality of
the Samelson commutator therefore changes the source level-twelve
Heisenberg group into the level-six group with the same multiplication
\[
 (a,b;z)(a',b';z')=(a+a',b+b';z+z'+ab'),
 \qquad z\in\Z/6.
\]
There is no retained finite face-filling module in the rank-three
target.  Indeed the suspension splitting
\(\Sigma^2P\simeq S^5\vee S^5\vee S^8\) gives
\[
 \pi_2\operatorname{Map}_*(P,BSU(3))
 =\pi_4SU(3)^2\oplus\pi_7SU(3)=0.
\]
Evaluation adds no term because \(\pi_2BSU(3)=0\).  Thus the two
\(\pi_4SU(2)\) differences and the top \(\pi_7SU(2)\) difference
all map to zero.  The determinant filling retains its prescribed integer
\(n\); on the now trivial slice its boundary is zero and creates no
additional finite edge coordinate.

Changing the slice trivialization changes the two central edge
coefficients by \(N_c\Z^2\), now modulo six.  The rank-two
determinant changes by coefficient six likewise vanish in the target;
their secondary changes have already vanished.  The formula therefore
descends to ordinary complex isomorphism classes, including odd
determinant.  Finally
\[
 N_c=\begin{pmatrix}A&B\\C&D\end{pmatrix}
       \begin{pmatrix}0&-1\\1&0\end{pmatrix},
\]
so \(T_6\) is exactly Yang's rank-three quotient in
Theorem~\ref{p1b:thm:g20-rank3-fibre} with \(m=0\).
Every target residue class has a lift of its two coefficients modulo
twelve, and the admitted source orbit charts provide fillings at every
such primary.  Hence stabilization is onto.
\end{proof}

\begin{corollary}[Uniform stabilization fibres and the final stable map]
\label{p1b:cor:low-rank-stabilization-fibres}
Retain the ordinary rank-two multiplicities of
Theorem~\ref{p1b:thm:v2170-ordinary-counts}, and let \(Q_6,Q_{12}\)
have its meaning.  Every rank-three target class in the preceding
theorem has the following number of ordinary rank-two preimages:
\[
\begin{array}{c|c|c}
 n\bmod2&\hbox{Chern coefficients}&\hbox{preimages}\\ \hline
 1&\hbox{some odd}&4\\
 1&\hbox{all even}&8\\
 0&\hbox{some odd}&4\hbox{ in each }\varepsilon\hbox{-sector}\\
 0&\hbox{all even}&
 \nu Q_{12}/Q_6\hbox{ in sector }0,\quad5\hbox{ in sector }1,
\end{array}
\]
where \(\nu=4+2^{\operatorname{rank}_{\F_2}(N_c/2)}\) in the
last row.  In particular the zero-Chern fibre consists of twenty-five
rank-two classes over each of the thirty-six rank-three classes: twenty
in the trivial restriction sector and five in the exotic sector.
The next map, \(E_3\mapsto E_3\oplus\mathbf1\), sends the whole
rank-three fixed-Chern fibre to its single stable rank-four class.
\end{corollary}
\begin{proof}
Reduction modulo six induces a surjection of primary quotients with
constant kernel size \(Q_{12}/Q_6\).  In the rows already normalized
modulo six, it is the identity primary quotient.  All secondary
differences are killed by the preceding theorem.  The secondary orbit
counts in Theorem~\ref{p1b:thm:v2170-ordinary-counts} are constant at
each primary and give the displayed fibre sizes.  At zero Chern data,
\(Q_{12}/Q_6=4\) and \(\nu=5\), giving \(20+5=25\).
Rank-four stabilization is classified by the integral K-character
lattice of Theorem~\ref{p1b:thm:chern}; equal Chern classes give the
same K-class, and rank four is in the injective stable range of
Proposition~\ref{p1b:prop:g15-stable-range}.  This proves the last
assertion for every rank-three fixed-Chern fibre, including \(m\ne0\).
\end{proof}


\subsubsection{The ordinary determinant quotient and its exact rank-two counts}
\label{p1b:sec:v2170-ordinary}

Retain the G15 restriction bit \(\varepsilon\), the G20 Heisenberg
mapping-space calculation and its affine-unframing lemma, and the G21 Born
action on determinant-preserving classes.  These inputs concern the actual
bundle component, not stable K-theory.  Put
\[
 P=S_x^3\times S_p^3,\quad M_8=T^2\times P,\quad
 \eta_T=a_ta_E,\quad
 N_c=\begin{pmatrix}B&-A\\D&-C\end{pmatrix},\quad
 Q_m=|\operatorname{coker}(N_c\bmod m)|.
\]
Here \(c_1=n\eta_T\), \(c_2\) and \(\Delta=BC-AD\) are the
Chern data of Section~\ref{p1b:sec:g20}.  The counts in this subsection
are ordinary complex isomorphism classes, with the determinant
trivialization forgotten.  They do not count coherent Born structures.

\begin{lemma}[The determinant boundary and the actual restriction sector]
\label{p1b:lem:v2170-determinant-boundary}
Write \(F_\varepsilon=\operatorname{Map}(P,BSU(2))_\varepsilon\) and
\(X_\varepsilon=\operatorname{Map}(P,BU(2))_\varepsilon\).  The
determinant fibration has boundary
\[
 \partial_\varepsilon:\Z=\pi_2BU(1)\longrightarrow
 \pi_1F_\varepsilon,\qquad
 \partial_\varepsilon(k)=\kappa^{6\varepsilon k},
\]
where \(\kappa\) is the central order-twelve generator of G20.  A
rank-two bundle in sector \(\varepsilon\) exists precisely when
\[
 \boxed{\Delta\equiv6\varepsilon n\pmod{12}.}
\]
\end{lemma}
\begin{proof}
Since \(H^1(P;\Z)=H^2(P;\Z)=0\), the determinant target is
\(\operatorname{Map}(P,BU(1))\simeq BU(1)\).  A generator of its
second homotopy group has clutching determinant winding one.  On the
smash sphere \(S^6\), the nonzero restriction class is represented in
the gauge fibre by \(\eta_3\eta_4\).  The determinant-circle Samelson
product is \(\eta_3\), so the boundary is
\(\eta_3\eta_4\eta_5=6\nu'\).  The top-cell collapse
\(P\to S^6\) inserts this as \(\kappa^6\) in the G20 central
subgroup.  The zero restriction has zero boundary.  Thus the central
gauge transformation \(-I\) has component \(\kappa^{6\varepsilon}\).

Use the one-vertex, two-edge, one-face torus model.  Its edge degrees are
\((A,B)\), \((C,D)\).  In the G20 group law
\[
 (a,b;z)(a',b';z')=(a+a',b+b';z+z'+ab'),
\]
their commutator is \(\kappa^{-\Delta}\).  A determinant filling of
degree \(n\) exists if and only if this equals the displayed boundary.
The right side has order at most two, so its sign is immaterial.  This is
the asserted congruence.  Equality supplies the face nullhomotopy; the
parameter torus has no higher cells.  The exponential law supplies the
actual bundle on \(T^2\times P\) with the prescribed Chern data.
\end{proof}

Let \(V_2=\mathbb F_2\{\alpha_x,\alpha_p,\beta\}\) be the G20
filling module, with coordinates \((r,s,e)\).  Its slice action is
\(\rho(a,b;z)(r,s,e)=(r,s,e+br+as)\).

\begin{lemma}[The determinant action on fillings]
\label{p1b:lem:v2170-determinant-action}
In the trivial restriction component, an integral determinant-degree-one
generator \(q\in\pi_2X_0\) can be chosen with action
\[
 (k;r,s,e)\longmapsto
 (k;r+ka,s+kb,e+br+as+kab).
\]
The complete extra edge quotient from forgetting the determinant adds
\((C,D,CD)\) and \((A,B,AB)\) to the finite filling coordinates.
In the exotic component the integral generator has determinant degree
two and is pointwise scalar; its action on the integral summand is
trivial.  A determinant half-winding translates the corresponding primary
coordinate by six.  On filling differences it has identity linear part.
\end{lemma}
\begin{proof}
In the zero component represent \(q\) by the constant gauge loop
\(\operatorname{diag}(e^{i\theta},1)\).  Conjugation by a degree map
\(f:P\to SU(2)\) displaces it by the commutator loop
\(f q f^{-1}q^{-1}\).  Its class \(\lambda(f)\) is a crossed
homomorphism: \(\lambda(gg')=\lambda(g)+\rho(g)\lambda(g')\).
The determinant-sphere Samelson value gives
\(\lambda(x)=\alpha_x\), \(\lambda(p)=\alpha_p\).
Both \(\lambda(xp)\) and \(\lambda(px)\) equal
\(\alpha_x+\alpha_p+\beta\); since \(xp=px\kappa\),
\(\lambda(\kappa)=0\).  Consequently
\(\lambda(a,b;z)=(a,b,ab)\), including the top coordinate.  The
cellular edge differential \((\Psi(g)-1)q\) gives the two stated
vectors and the full displayed action.

In the exotic component \(e^{i\theta}I\) generates the kernel of the
determinant boundary, with degree two, and commutes with every gauge.
For a degree-one change cut the torus circle and use the scalar path
\(e^{i\pi t}I\).  Its seam is \(-I\), whose actual component was
computed above.  Multiplying an edge clutching by this central element
does not change its commutator or the difference between two face
fillings.  Returning to fixed normal-form representatives may translate
the reference filling, but cannot change the linear action on filling
differences.  The two operations commute and square to the identity
because squared determinant functions have scalar lifts.  This is an
actual action on relative bundle classes.

The lower translations of a determinant winding are obtained by slanting
the first Postnikov invariant
\(Sq^2\bar c_2+\bar c_1\bar c_2\): a loop with data \((h,v)\)
has translation \(Sq^2\bar v+\bar c_1\bar v+\bar h\bar c_2\).
For the two determinant generators this gives \((C,D)\) and \((A,B)\).
All these calculations use the actual mapping components and the G20
affine-unframing statement; a homotopy module alone would not determine
an affine filling action.
\end{proof}

\begin{theorem}[All ordinary fixed-Chern rank-two multiplicities]
\label{p1b:thm:v2170-ordinary-counts}
Each entry below is zero unless
\(\Delta\equiv6\varepsilon n\pmod{12}\).  In the all-even rows put
\(\rho_2=\operatorname{rank}_{\mathbb F_2}(N_c/2)\) and
\(\nu=4+2^{\rho_2}\).
\[
\begin{array}{c|c|c|c}
 n\bmod2& (A,B,C,D)&N_{M,0}&N_{M,1}\\\hline
 0&\text{some coefficient odd}&4Q_6&4Q_6\\
 0&\text{all coefficients even}&\nu Q_{12}&5Q_6\\
 1&\text{some coefficient odd}&
 4Q_6\ (\Delta\equiv0\!\!\pmod{12})&
 4Q_6\ (\Delta\equiv6\!\!\pmod{12})\\
 1&\text{all coefficients even}&8Q_6&0
\end{array}
\]
For \(m=6,12\) the required primary count is exactly
\[
 Q_m=\begin{cases}
 m^2,&c_2=0,\\
 m\gcd(m,s_1),&\Delta=0,\ c_2\ne0,\\
 \gcd(m,s_1)\gcd(m,|\Delta|/s_1),&\Delta\ne0,
 \end{cases}\qquad s_1=\gcd(A,B,C,D).
\]
In particular \(c_2=0\) gives \(720+180=900\) ordinary classes for
even determinant and \(288\) for odd determinant.
\end{theorem}
\begin{proof}
When one degree is odd, the edge-homotopy subgroup contains the top
\(\beta\)-line.  With that line removed, the full operations are
\[
 C_v(z,l)=(z+N_cv,l+\bar n\bar v),\quad
 d_t(z,l)=(z+6\varepsilon(1,0),l+(C,D)),
\]
\[
 d_E(z,l)=(z+6\varepsilon(0,1),l+(A,B)),
 \qquad l=(r,s)\in\mathbb F_2^2.
\]
For even determinant, before determinant forgetting there are
\(4Q_{12}\) classes.  The lower determinant matrix
\(R=\left(\begin{smallmatrix}C&A\\D&B\end{smallmatrix}\right)\bmod2\)
has rank one.  Its kernel equals \(\operatorname{im}(N_c\bmod2)\);
therefore a kernel winding's primary shift \(6k\) is already a slice
shift \(N_c(6u)\).  Every effective determinant orbit has size two.
Smith reduction gives \(Q_{12}=2Q_6\), proving \(4Q_6\) in either
sector.  For odd determinant normalize \(l=0\).  The residual slice
shifts are \(2N_c\), giving \(4Q_6\) primary classes.  Renormalizing
the two determinant operations adds \((6\varepsilon-\Delta,0)\) or
\((0,6\varepsilon+\Delta)\); both vanish by the sector congruence.

For all-even coefficients in sector zero the two additional edge vectors
of the preceding lemma vanish.  For even determinant the G20 shear
count is \(\nu Q_{12}\).  For odd determinant the action preserves
\(e+rs\), and normalizing \(r=s=0\) leaves the primary quotient by
\(2N_c\).  Its size is \(4Q_6\), with the independent retained bit
\(e+rs\), giving \(8Q_6\).  All-even coefficients have
\(4\mid\Delta\), so admissible odd determinant forces
\(12\mid\Delta\) and excludes the exotic sector.

It remains to count the all-even exotic sector at even determinant.
Tensor by a fixed line to reduce to determinant zero; \(\eta_T^2=0\)
preserves \(c_2\) and the restriction sector.  Write \(N_c=2M_c\).
The combined primary quotient is
\((\Z/12)^2/(N_c\Z^2+6\mathbb F_2^2)\), of size \(Q_6\).
At a raw primary its stabilizer is
\[
 H_z=\{(v,k):N_cv+6k=0\pmod{12}\}.
\]
It acts by \(e\mapsto e+br+as+\chi_z(v,k)\), with lower coordinates
fixed.  The affine displacement is a character, since these are actual
commuting operations with identity linear part on the top direction.
The parity map \(H_z\to\mathbb F_2^2\) is onto: take
\(v=3w\), \(k=M_cw\bmod2\).  Its kernel is even \(v\) and
\(k=0\), by reduction modulo four.  Such a pure even slice change is
the identity by G20's affine-unframing lemma.  Thus \(\chi_z\) factors
through parity and can be absorbed into a translation of \((r,s)\).
The resulting full shear has two orbits at lower label zero and one at
each of the other three labels: exactly five, without requiring the
affine character to vanish.  This proves \(5Q_6\).  Finally the stated
Smith formula counts the primary cokernel in all matrix ranks.
\end{proof}

\begin{theorem}[Ordinary Born-fixed rank-two classes]
\label{p1b:thm:v2170-ordinary-born}
Born-fixed ordinary bundles require \(n=0\), \(C=B,D=A\), and
\(12\mid(B^2-A^2)\).  On these data set
\(g_m=\gcd(m,B-A)\) and
\[
 \mu=\begin{cases}
 2,&A,B\text{ odd},\\
 4,&A,B\text{ even},\ A+B\equiv2\pmod4,\\
 3,&A,B\text{ even},\ A+B\equiv0\pmod4.
 \end{cases}
\]
Then
\[
 \boxed{N^B_{M,0}=\mu g_{12},\qquad
 N^B_{M,1}=\begin{cases}2g_6,&A,B\text{ odd},\\3g_6,&A,B\text{ even}.
 \end{cases}}
\]
The count is zero outside the displayed Chern domain.  At zero Chern
data it is \(36+18=54\).  These count fixed ordinary isomorphism
classes and do not equip them with coherent involutions.
\end{theorem}
\begin{proof}
In the all-even zero sector the additional determinant quotient is
trivial, so G21 gives \(\mu g_{12}\).  In its odd branch the quotient
identifies \((r,s)\) with \((r+1,s+1)\); both resulting labels are
fixed by exchange, giving \(2g_{12}\).

For the exotic sector let \(S\) be the determinant-preserving fibre and
\(D_0=\langle d_t,d_E\rangle\) the order-four determinant group,
allowing a nonfaithful action.  Born \(b\) swaps its generators.  By
counting pairs \((s,d)\) with \(b(s)=d(s)\), without a freeness
assumption,
\[
 |\operatorname{Fix}_{S/D_0}(b)|
 =\tfrac14\sum_{d\in D_0}|\operatorname{Fix}_S(db)|
 =\tfrac12(F_b+F_{d_tb}).
\]
The primary Born action is
\(z\mapsto AB(1,1)-\operatorname{swap}(z)\).  The obstruction to
adding \(6(1,0)\) is six in
\(\operatorname{coker}(N_c\bmod12)/(1+\operatorname{swap})
\cong\Z/g_{12}\).  For odd \(A,B\), G21 gives
\(F_b=2g_{12}\); the twisted term is \(2g_{12}\) if
\(g_{12}\mid6\), otherwise zero.  The fixed-orbit identity gives
\(2g_6\).

For even \(A,B\) and \(4\mid(B-A)\), the twisted primary has no fixed
class, while \(F_b=3g_{12}\), giving \(3g_6\).  If
\(B-A\equiv2\pmod4\), \(N_c/2\) is invertible modulo two and the
determinant primary shifts are already slice shifts.  There are
\(g_{12}=g_6\) fixed primary classes.  Their secondary Born action has
the G21 form \((r,s,e)\mapsto(s,r,e+c(r+s))\).  The preceding affine
stabilizer proof gives one lower label with two ordinary classes and
three with one.  The uniquely doubled lower label must be diagonal,
because Born preserves its size and swaps the lower coordinates.  Its
two classes have fixed representatives, the other diagonal label gives
one more, and the off-diagonal labels are exchanged.  This gives three
classes per fixed primary and finishes the count.
\end{proof}


\subsubsection{Coherent Born lifts of the unstable rank-two sector}
\label{p1b:sec:g21}

Let \(C_2=\langle B\mid B^2=1\rangle\) denote the exchange group.  A
\emph{coherent Born lift} on a complex bundle \(E\to M_8\) is an
isomorphism
\[
 \Phi:B^*E\longrightarrow E,
 \qquad
 \Phi\circ B^*\Phi=\id_E.
\]
This is stronger than invariance of characteristic classes and stronger
than the existence of an unspecified ordinary isomorphism \(B^*E\cong E\).

\begin{proposition}[Gauge-extension criterion and Born norm equation]
\label{p1b:prop:g21-norm}
Suppose \(B^*E\cong E\) as an ordinary complex bundle and let
\(\mathcal G(E)\) be its gauge group.  Bundle automorphisms covering either
\(1\) or \(B\) form an extension
\[
 1\longrightarrow\mathcal G(E)
 \longrightarrow\operatorname{Sym}_B(E)
 \longrightarrow C_2\longrightarrow1.
\]
The bundle admits a coherent Born lift if and only if this extension splits.
Equivalently, choose any ordinary isomorphism \(f:B^*E\to E\), set
\[
 A_f=f\circ B^*f\in\mathcal G(E),
 \qquad
 \alpha_f(u)=f\circ B^*u\circ f^{-1},
\]
and replace \(f\) by \(u\circ f\).  Then
\[
 A_{u f}=u\,\alpha_f(u)\,A_f,
\]
so a coherent lift exists exactly when
\[
 \boxed{u\,\alpha_f(u)=A_f^{-1}}
\]
has a solution.  Once one coherent lift \(\Phi\) is fixed, the remaining
lift classes modulo gauge equivalence are the corresponding nonabelian
\(H^1(C_2;\mathcal G(E)_\Phi)\).
\end{proposition}
\begin{proof}
A splitting chooses an order-two element of
\(\operatorname{Sym}_B(E)\) above the nontrivial element of \(C_2\), which
is exactly a bundle isomorphism satisfying the cocycle equation.  Direct
composition gives the displayed transformation law for the defect.  Setting
the transformed defect equal to the identity is the norm equation.  The
last statement is the standard cocycle/coboundary description after one
splitting has been chosen.
\end{proof}

\begin{theorem}[Born-fixed principal-\(SU(2)\) multiplicity]
\label{p1b:thm:g21-born-fixed-count}
Count pairs \((E,\tau:\det E\xrightarrow{\sim}\mathbf1)\) up to
determinant-preserving isomorphism, equivalently principal \(SU(2)\)
bundles, as in Theorem~\ref{p1b:thm:g20-rank2-fibre-even}.  Born exchange
acts on the pair by pullback.  Let \(c_1=0\) and suppose the Chern data
are Born invariant:
\[
 c_2=A(a_tu_x+a_Eu_p)+B(a_tu_p+a_Eu_x).
\]
Assume ordinary rank-two existence,
\[
 12\mid(B^2-A^2).
\]
For either G15 restriction class \(\epsilon\in\Z/2\), put
\[
 g_B=\gcd(12,B-A)
\]
and
\[
 \mu(A,B)=
 \begin{cases}
 2,&A\text{ or }B\text{ odd},\\
 4,&A,B\text{ even and }A+B\equiv2\pmod4,\\
 3,&A,B\text{ even and }A+B\equiv0\pmod4.
 \end{cases}
\]
Then the number of principal-\(SU(2)\) isomorphism classes in that
restriction sector fixed by Born exchange is
\[
 \boxed{N^{(\epsilon)}_{\rm Born}(A,B)=g_B\,\mu(A,B).}
\]
The two restriction sectors have the same count, so the total is twice this
number.  In particular, for \(A=B=0\) there are \(36\) Born-fixed
principal-\(SU(2)\) classes in each restriction sector and \(72\) in
total.  These are fixed isomorphism classes, not bundles equipped with a
chosen coherent Born lift.
\end{theorem}
\begin{proof}
With the Heisenberg convention of Paper~I-D, Section~\ref{p1b:app:g20}, sphere exchange
acts by
\[
 (a,b;z)\longmapsto(b,a;ab-z).
\]
Hence the scalar central coordinates transform as
\[
 (z_t,z_E)\longmapsto(AB-z_E,AB-z_t).
\]
The primary fixed quotient has \(\gcd(12,B-A)\) classes.  Applying the
verified secondary action and its stabilizer parity gives the three values
of \(\mu\); the complete orbit calculation is recorded in
Paper~I-D, Section~\ref{p1b:app:g21}.
\end{proof}

For ordinary complex bundles, one must first take the determinant quotient
of Remark~\ref{p1b:rem:g20-determinant-forgetting}.  A quotient class is
Born fixed when its \(H^1(M_8;\Z)\)-orbit is stable under exchange; it
need not have a Born-fixed representative before that quotient.
Theorem~\ref{p1b:thm:v2170-ordinary-born} counts these exchange-stable
ordinary quotient orbits exactly; its zero-Chern total is 54.  This does
not turn a fixed ordinary isomorphism class into a chosen coherent lift.

\begin{theorem}[Equivariant Dirac obstruction]
\label{p1b:thm:g21-index-obstruction}
Let \(E\) be a complex rank-two bundle with \(c_1=0\) and
\[
 c_2=A(a_tu_x+a_Eu_p)+B(a_tu_p+a_Eu_x).
\]
Any coherent complex rank-two Born lift on \(E\) satisfies
\[
 \boxed{4\mid(A+B).}
\]
If the fixed-fibre involution has mixed eigenranks \((1,1)\), then the
stronger condition
\[
 \boxed{A+B=0}
\]
holds.  For a pure fixed-fibre action the equivariant Dirac index has virtual character
coefficients
\[
 \boxed{
 m_\pm=
 \frac{B^2-A^2}{12}\pm\frac{A+B}{4},
 }
\]
up to interchanging \(m_+\) and \(m_-\).
\end{theorem}
\begin{proof}
The exchange fixed locus is \(F\cong S^1\times S^3\), with rank-four normal
bundle on which \(B\) acts by \(-I\).  The spin fixed-point denominator has
absolute value four.  Since
\[
 i^*c_2(E)=2(A+B)au,
\]
a pure lift has nontrivial equivariant character
\(t=\pm(A+B)/2\).  Independently,
\(s=\operatorname{ind}(D_E)=(B^2-A^2)/6\).  Integrality of these virtual character
coefficients implies \(s\equiv t\pmod2\), yielding the displayed mod-four
condition; the mixed eigenline splitting has \(c_2(E|_F)=0\) because
\(H^2(F;\Z)=0\).  Paper~I-D, Section~\ref{p1b:app:g21} gives the complete replay and
sign firewall.
\end{proof}

\begin{corollary}[Quaternionic refinement for determinant-preserving lifts]
\label{p1b:cor:g21-quaternionic-index-obstruction}
In the notation of Theorem~\ref{p1b:thm:g21-index-obstruction}, a coherent
principal-\(SU(2)\) Born lift satisfies the stronger condition
\[
 \boxed{\frac{B^2-A^2}{12}\pm\frac{A+B}{4}\in2\Z.}
\]
Equivalently, on the ordinary principal rank-two existence domain,
\[
 \boxed{24\mid\bigl(B^2-A^2+3(A+B)\bigr).}
\]
This condition is necessary; it is not asserted to characterize the
coarse-quotient descent image by itself.
\end{corollary}
\begin{proof}
In real dimension eight, the two complex chiral spin representations are
complexifications of real representations.  Let \(C_S\) denote their
parallel real structure, with \(C_S^2=1\).  The defining representation
of \(SU(2)=Sp(1)\) has its invariant anti-linear quaternionic structure
\(J_E\), with \(J_E^2=-1\).  Consequently
\(J=C_S\otimes J_E\) is a well-defined anti-linear quaternionic
structure on the twisted chiral spinors.  It commutes with Clifford
multiplication and with a Dirac operator formed from an \(SU(2)\)
connection.  Average that connection and the metric over the given
coherent involution.  The connection remains \(\mathfrak{su}(2)\)-valued.

Choose either Born-fixed spin structure and either involutive spin lift,
as in the fixed-point calculation above.  The spin lift preserves
\(C_S\), and a principal \(SU(2)\) lift preserves \(J_E\).  Thus the
combined involution commutes with \(J\).  Its positive and negative
eigenspaces in each chiral Dirac kernel and cokernel are quaternionic
vector spaces.  Their complex dimensions, and hence both virtual
character multiplicities \(m_+,m_-\), are even.

Insert the formula for \(m_\pm\) from
Theorem~\ref{p1b:thm:g21-index-obstruction}.  The plus condition says
\(24\mid B^2-A^2+3(A+B)\).  Since \(12\mid B^2-A^2\), this also
forces \(4\mid A+B\); the difference between the two numerators is
\(6(A+B)\), then divisible by twenty-four.  The single congruence is
therefore equivalent to both even-multiplicity conditions.  A mixed
\(U(2)\) lift need not preserve \(J_E\), so this refinement does not
alter the mixed-lift classification.
\end{proof}

\begin{proposition}[Every pure complex Born lift reduces equivariantly to \(SU(2)\)]
\label{p1b:prop:pure-u2-determinant-reduction}
Let \(E\to M_8\) have a coherent complex rank-two Born lift whose
fixed-fibre involution is pure, of type \(+I\) or \(-I\).  Then its
equivariant determinant line admits an invariant unit trivialization.
With this trivialization the same lift is a coherent principal
\(SU(2)\) lift.  In particular every pure complex lift has G15
restriction bit zero and obeys the quaternionic mod-twenty-four
obstruction of Corollary~\ref{p1b:cor:g21-quaternionic-index-obstruction}.
\end{proposition}
\begin{proof}
Born invariance gives \(c_1(E)=B^*c_1(E)\).  Since exchange acts as
minus one on the torsion-free \(H^2(M_8;\Z)=\Z\eta\), the
determinant line is ordinarily trivial.  Average a Hermitian metric over
the involution and choose a unit trivialization of that line.  Its
equivariant action is then a function \(d:M_8\to U(1)\) satisfying
\[
 d\,(d\circ B)=1,\qquad d|_{\Delta K}=1.
\]
The second identity uses \(\det(+I)=\det(-I)=1\).

The winding class of \(d\) lies in the anti-invariant part of
\(H^1(M_8;\Z)\), hence is \(k(a_t-a_E)\).  The character
\(v=z_t^k\) has coboundary
\(v/(v\circ B)=(z_t/z_E)^k\) with exactly this winding and value
one on the fixed locus.  Divide \(d\) by that coboundary.  The
resulting \(d_0\) is nullhomotopic and has a continuous real logarithm
\(d_0=e^{i\alpha}\).  Since \(\Delta K\) is connected and
\(d_0=1\) there, subtract a constant integral multiple of \(2\pi\)
to arrange \(\alpha|_{\Delta K}=0\).  The cocycle identity makes
\(\alpha+\alpha\circ B\) a continuous integral multiple of
\(2\pi\); it vanishes on the fixed locus and therefore on the
connected manifold.  Thus
\[
 d_0=\frac{e^{i\alpha/2}}{e^{i\alpha\circ B/2}},\qquad
 d=\frac{u}{u\circ B},\quad u=v e^{i\alpha/2}.
\]
Multiplying the chosen determinant section by \(u\) gives an invariant
unit section (equivalently change the determinant coordinate by its
inverse).  The unitary frame bundle with this fixed determinant is
preserved by the original involution and has structure group \(SU(2)\).
Apply the all-Chern exotic restriction no-go and the quaternionic index
obstruction to that reduction.  This establishes existence of a
determinant-preserving reduction; it does not identify the gauge quotients
or count the different determinant trivializations.
\end{proof}

\begin{remark}[The quaternionic refinement is strict]
For \(A=0,B=12\), the ordinary rank-two existence test holds and
\(4\mid A+B\), but
\[
 \frac{B^2-A^2}{12}\pm\frac{A+B}{4}=15,9
\]
are odd. Hence no coherent principal \(SU(2)\) lift exists on this
Chern fibre. The determinant reduction above also excludes every pure
complex lift there.
\end{remark}

\begin{theorem}[Coherent Born lift of the G15 exotic complex bundle]
\label{p1b:thm:g21-exotic-lift}
Let \(E_{\rm ex}\to P_v=S_x^3\times S_p^3\) be the nontrivial rank-two
complex bundle of Theorem~\ref{p1b:thm:g15-unstable}.  Then there is an
explicit coherent Born lift
\[
 \boxed{\Phi_{\rm ex}:B^*E_{\rm ex}\longrightarrow E_{\rm ex},
 \qquad \Phi_{\rm ex}\circ B^*\Phi_{\rm ex}=1.}
\]
Its pullback
\(\widetilde E_{\rm ex}=\operatorname{pr}_{P_v}^*E_{\rm ex}\to M_8\)
therefore also has a coherent Born lift.  On the connected exchange fixed
locus the constructed fibre involution has eigenranks \((1,1)\).
\end{theorem}
\begin{proof}
Paper~I-D, Section~\ref{p1b:app:g21-exotic-construction} constructs a
\(C_2\)-equivariant map
\[
 s:P_v\longrightarrow S^5
\]
whose underlying class is \(\eta_5\circ q_\wedge\), together with the
rank-two orthogonal-complement bundle \(\xi_v=v^\perp\subset\C^3\) on
\(S^5\).  Equivariant pullback identifies \(s^*\xi\) with \(E_{\rm ex}\)
and supplies the involution.  The fixed target decomposition has one
positive and one negative complex eigenline.
\end{proof}

\begin{theorem}[No determinant-preserving Born lift of the G15 exotic bundle]
\label{p1b:thm:g21-exotic-nogo}
The bundle \(E_{\rm ex}\to S^3\times S^3\) admits no coherent
principal-\(SU(2)\) Born lift.  The same is true for its pullback
\(\widetilde E_{\rm ex}\to M_8\).
\end{theorem}
\begin{proof}
A principal fixed-fibre involution is necessarily \(\pm I\); multiplication
by the central \(-I\) reduces a hypothetical lift to trivial fixed isotropy,
which would descend through
\[
 S^3\times S^3\longrightarrow\operatorname{SP}^2(S^3).
\]
The reduced symmetric square is \(\Sigma^4\mathbb{RP}^2\), and the induced
smash-orbit map is the fourfold suspension of the double cover
\(S^2\to\mathbb{RP}^2\), hence null.  Consequently every principal bundle
pulled back from the symmetric square is trivial on \(S^3\times S^3\),
whereas \(E_{\rm ex}\) is the nonzero G15 class.  The detailed descent and
nullity proof is in Paper~I-D, Section~\ref{p1b:app:g21}.
\end{proof}

\begin{corollary}[The exotic restriction excludes every principal Born lift]
\label{p1b:cor:g21-all-chern-exotic-nogo}
Let \(E\to M_8\) be any principal \(SU(2)\) bundle, with arbitrary
admissible second Chern class.  If its G15 restriction bit is
\(\varepsilon=1\), it admits no coherent determinant-preserving Born lift.
Consequently
\[
 \operatorname{im}\bigl(q^*:[\operatorname{SP}^2(S^1\times S^3),BSU(2)]
 \longrightarrow[M_8,BSU(2)]\bigr)
 \subseteq\{\varepsilon=0\}.
\]
\end{corollary}
\begin{proof}
The scalar slice \(P_0=\{z_t=z_E=1\}\times S_x^3\times S_p^3\)
is preserved by exchange.  Restricting a coherent principal lift to this
slice preserves the cocycle equation and the determinant trivialization.
For \(\varepsilon=1\), its underlying restricted bundle is the exotic
bundle of Theorem~\ref{p1b:thm:g15-unstable}, contradicting
Theorem~\ref{p1b:thm:g21-exotic-nogo}.  A bundle pulled back from the
coarse quotient has its canonical coherent lift, proving the image
statement.  No Chern-class hypothesis beyond existence enters this
restriction argument.
\end{proof}

\begin{corollary}[Determinant character of the constructed lift]
\label{p1b:cor:g21-determinant}
For the lift of Theorem~\ref{p1b:thm:g21-exotic-lift}, the equivariant
determinant line is the sign line of \(C_2\).  Thus the construction is a
genuine coherent rank-two complex, or \(U(2)\)-level, Born lift, but it does
not preserve a chosen determinant-one trivialization.  By
Theorem~\ref{p1b:thm:g21-exotic-nogo}, this determinant-one failure cannot
be repaired by choosing a different principal-\(SU(2)\) lift of the same
ordinary exotic bundle.
\end{corollary}

\begin{theorem}[Principal-\(SU(2)\) coherent lifts as descent fibres]
\label{p1b:thm:g21-principal-descent}
Regroup
\[
 M_8=K\times K,\qquad K=S^1\times S^3,
\]
and let
\[
 q:M_8\longrightarrow Q:=\operatorname{SP}^2(K)
\]
be the swap quotient.  For a principal \(SU(2)\) bundle \(E\to M_8\), set
\[
 \operatorname{Desc}(E)
 =\{[\bar E]\in[Q,BSU(2)]:q^*\bar E\cong E\}.
\]
Then the full set of coherent determinant-preserving Born lifts modulo gauge
is
\[
 \boxed{
 \mathscr L_{SU(2)}(E)
 \cong
 \operatorname{Desc}(E)\times\{+,-\}.
 }
\]
Equivalently, after choosing one coherent lift,
\[
 \boxed{
 H^1(C_2;\mathcal G(E)_\Phi)
 \cong
 \operatorname{Desc}(E)\times\Z/2
 }
\]
as pointed sets.  In particular a principal lift exists exactly when
\([E]\in\operatorname{im}q^*\), and if the descent fibre is finite its
cardinality is half the number of principal coherent lift classes.
\end{theorem}
\begin{proof}
Every principal fixed-fibre isotropy is \(\pm I\).  Multiplying by the
central \(-I\) reduces to type \(+I\), for which the stabilizer acts
trivially and the equivariant principal bundle descends to an ordinary
principal bundle on the coarse quotient.  Conversely every descent class
pulls back with its canonical exchange action.  Choice of pullback
isomorphism changes the lift only by gauge equivalence, and the two central
fixed signs cannot be gauge conjugate.  Paper~I-D, Section~\ref{p1b:app:g21} records
the local descent proof and the nonabelian-\(H^1\) identification.
\end{proof}

\begin{remark}[Mixed \(U(2)\) lift fibre]
\label{p1b:rem:g21-mixed-fibre}
A mixed fixed-fibre involution \(\operatorname{diag}(1,-1)\) has nontrivial
isotropy and therefore is not classified by ordinary bundles on the coarse
quotient.  The correct global object is the equivariant classifying space
with the mixed local representation retained.  If
\[
 \operatorname{For}:\operatorname{Bun}^{C_2}_{U(2),\mathrm{mix}}(M_8)
 \longrightarrow\operatorname{Bun}_{U(2)}(M_8)
\]
is the forgetful map in that genuine equivariant category, then
\[
 \boxed{
 \mathscr L^{\rm mix}_{U(2)}(E)
 =\operatorname{For}^{-1}([E]).
 }
\]
This genuine equivariant fibre is evaluated by the actual mixed
gluing construction of Section~\ref{p1b:sec:v2170-mixed} and the
completed image, sphere and natural boundary calculations in
Section~\ref{p1b:sec:rank-two-closure}.
Theorem~\ref{p1b:thm:complete-mixed-rank-two} gives every nonempty
ordinary lift cardinality and fixed-row torsor group.
At \(c_1=c_2=0\), 72 coherent classes lie over 54 ordinary classes:
36 fibres are singletons and 18 are doubletons. Coherent structures
and their underlying ordinary bundles retain separate meanings.
The exotic G15 fibre is nonempty by
Theorem~\ref{p1b:thm:g21-exotic-lift}.  The use of local representations in
the equivariant classifying problem follows the standard framework of
\cite{LuckUribe2014}.
\end{remark}

\subsubsection{An explicit coarse-quotient source for principal descent}
\label{p1b:sec:principal-descent-source}

\begin{theorem}[The symmetric square as a specified mapping torus]
\label{p1b:thm:principal-descent-mapping-torus}
Put \(Y=S^3\), \(P=Y\times Y\), \(A=\operatorname{SP}^2Y\), and
let \(q_Y:P\to A\) be the unordered-pair quotient.  Form the double
mapping cylinder
\[
 F_7=\bigl(A_-\sqcup(P\times[0,\pi])\sqcup A_+\bigr)
 /\bigl((x,y,0)\sim q_Y(x,y)_-,\
         (x,y,\pi)\sim q_Y(x,y)_+\bigr).
\]
The formula
\(R[\phi,x,y]=[\pi-\phi,y,x]\) defines an involution of \(F_7\),
interchanging \(A_-\) and \(A_+\).  There is a homeomorphism
\[
 \boxed{\operatorname{SP}^2(S^1\times S^3)
 \cong (F_7\times[0,2\pi])/((f,2\pi)\sim(Rf,0)).}
\]
The projection to the mapping-torus circle is the scalar product
\(\mu\{(z,x),(w,y)\}=zw\).  The homeomorphism is specified by
\[
 [t;\phi,x,y]\longmapsto
 \{(e^{i(t/2+\phi)},x),(e^{i(t/2-\phi)},y)\}.
\]
\end{theorem}
\begin{proof}
For a fixed product \(zw=e^{it}\), divide both scalar coordinates by
the chosen square root \(e^{it/2}\).  They become
\((e^{i\phi},e^{-i\phi})\).  An unordered pair has a representative
with \(0\le\phi\le\pi\); for an interior value this representative
fixes the order of \(x,y\), while at either endpoint the two scalar
coordinates coincide and the sphere coordinates become unordered.
This is precisely the displayed double mapping cylinder.  Increasing
\(t\) by \(2\pi\) changes the square root by \(-1\).  The same
unordered pair is then represented at \(t\) by
\((\pi-\phi,y,x)\), giving the stated monodromy.  The formula is
continuous and bijective after these identifications.  Its compact
source and Hausdorff target make it a homeomorphism.
\end{proof}

\begin{corollary}[Integral groups and the actual gluing target]
\label{p1b:cor:principal-descent-source-cohomology}
The nonzero positive-degree integral cohomology groups of \(F_7\) are
\[
 H^3(F_7)=\Z,\quad H^4(F_7)=\Z,\quad
 H^6(F_7)=(\Z/2)^2,\quad H^7(F_7)=\Z.
\]
The monodromy is the identity on the free groups and interchanges the
two degree-six generators.  Consequently the nonzero positive-degree
integral groups of \(Q=\operatorname{SP}^2(S^1\times S^3)\) are
\[
\begin{array}{c|ccccccc}
 d&1&3&4&5&6&7&8\\\hline
 H^d(Q;\Z)&\Z&\Z&\Z^2&\Z&\Z/2&\Z\oplus\Z/2&\Z.
\end{array}
\]
For any CW target \(Z\), the specified quotient gives the actual
mapping-space comparison
\[
 \operatorname{Map}(F_7,Z)\simeq
 \operatorname{Map}(A,Z)
 \mathop{\times}^{h}_{\operatorname{Map}(P,Z)}
 \operatorname{Map}(A,Z),
\]
and maps from \(Q\) are the \(R\)-twisted loops in this homotopy
pullback.  Thus a numerical evaluation of principal descent must retain
the two maps \(q_Y^*\), their path data, and the monodromy on this
homotopy pullback.  The displayed integral groups alone are not such an
evaluation.
\end{corollary}
\begin{proof}
The symmetric-square cofibration used in the G21 no-go proof gives
\(H^3(A;\Z)=\Z\), \(H^6(A;\Z)=\Z/2\), and no other
positive groups.  On degree three,
\(q_Y^*(u)=u_x+u_p\); its torsion degree-six class pulls back to zero
in the free group \(H^6(P;\Z)\).  The Mayer--Vietoris sequence of
the double mapping cylinder therefore has degree-three kernel
\(\ker(\Z^2\to\Z^2)=\Z\), degree-four cokernel
\(\Z^2/\Z(1,1)=\Z\), the two independent end torsion classes in
degree six, and the connecting image of \(H^6(P;\Z)=\Z\) in degree
seven.  All other positive groups vanish.  Interchanging the two ends
acts positively on the degree-three kernel.  The interval reflection
contributes a minus sign on connecting classes: in degree four it
cancels the minus sign of sphere exchange on
\(\Z^2/\Z(1,1)\), and in degree seven it cancels the orientation
sign of sphere exchange on \(P\).  This proves the asserted action.
The mapping-torus Wang sequence gives the displayed groups; the sole
extension \(0\to\Z/2\to H^7(Q)\to\Z\to0\) splits as an
abelian-group extension because its quotient is free.

Mapping the double mapping cylinder into a fibrant CW target converts
its homotopy pushout into the indicated homotopy pullback: a map consists
of its two end maps and their actual interpolating path on \(P\).
The mapping-torus identification imposes exactly the stated twisted
loop condition.  This description retains the maps and paths, rather
than replacing them by their cohomology groups.
\end{proof}

\begin{proposition}[A coarse-quotient constructor for the ordinary top difference]
\label{p1b:prop:pure-descent-local-top-constructor}
Let \(q:M_8\to Q=M_8/C_2\) be the Born quotient.  There is a based
map \(a:Q\to S^7\), supported in an arbitrarily small quotient
neighbourhood of a point of the diagonal, such that
\(a\circ q:M_8\to S^7\) is the nonzero local top class
\(\eta_7\in\pi_8(S^7)=\Z/2\).  Composing with the fibre inclusion
\(j:S^7\to BSp(1)=BSU(2)\) of
\(S^7\to BSp(1)\to BSp(2)\) gives a principal bundle on \(Q\)
whose pullback is the nonzero ordinary local top class
\(\nu'\eta_6\in\pi_7(S^3)=\Z/2\), and this local class remains
nontrivial after extension to \(M_8\).
In particular the ordinary top difference can occur with pure principal
Born isotropy; it is not intrinsically restricted to mixed \(U(2)\)
isotropy.
\end{proposition}
\begin{proof}
An invariant ball at a diagonal point is the unit ball in
\(\R^4_{\rm triv}\oplus\R^4_{\rm sign}\).  Collapse its boundary
and the complement of the ball.  Its quotient gives the based source
\[
 (\R^4\oplus\R^4_{\rm sign})^+/C_2
 \simeq S^4\wedge\Sigma\mathbb{RP}^3
 =\Sigma^5\mathbb{RP}^3.
\]
Here passing from the unreduced suspension to the reduced suspension
collapses a contractible basepoint arc and is a based homotopy
equivalence.  The ordered-cover map on this local source is
\(\Sigma^5p:S^8\to\Sigma^5\mathbb{RP}^3\), with
\(p:S^3\to\mathbb{RP}^3\) the double cover.

The Hopf map \(\eta_2:S^3\to S^2\), written \(v\mapsto vi v^{-1}\)
in unit quaternions, is invariant under
\(v\mapsto-v\), so it descends to
\(\bar\eta_2:\mathbb{RP}^3\to S^2\), with
\(\bar\eta_2p=\eta_2\).  Define \(a\) by this local collapse
followed by \(\Sigma^5\bar\eta_2\).  Its ordered pullback on the
collapsed ball is exactly
\(\Sigma^5\eta_2=\eta_7\), which is nonzero.

The principal \(Sp(1)\) fibration \(Sp(2)\to S^7\) has clutching
generator \(\nu'\in\pi_6(S^3)=\Z/12\).  Under adjunction the
composite \(j\eta_7\) is therefore \(\nu'\eta_6\), the already
retained nonzero \(\pi_7(S^3)\) class.  To check the global assertion,
use the actual zero-Chern G20 torus chart: the local top insertion is
\((z_t,z_E;r,s,e)=(0,0;0,0,1)\).  The slice action is
\((r,s,e)\mapsto(r,s,e+br+as)\), so it cannot change \(e\) when
\(r=s=0\); the verified affine-unframing term is zero.  The ordinary
determinant operations at \(\varepsilon=0,c_2=0\) likewise have
zero primary and secondary displacement.  Thus this class is distinct
from \((0,0;0,0,0)\) even after the full ordinary quotient.  This is
the required global check; nonzero homotopy on an isolated sphere alone
would not have established it.  The bundle constructed on
\(Q\) pulls back with its canonical order-two action and trivial
fixed-fibre isotropy, proving the last assertion.  This construction
supplies a local top modification; it does not enumerate the global
principal descent fibres.
\end{proof}

\begin{theorem}[The pure principal image at zero Chern data]
\label{p1b:thm:pure-descent-zero-chern-image}
At \(c_1=c_2=0\), every one of the thirty-six Born-fixed ordinary
rank-two classes in the trivial G15 restriction sector admits a pure
coherent principal-\(SU(2)\) Born lift.  None of the exotic restriction
sector does.  Thus the pure ordinary image has exactly thirty-six
classes.  Each of those classes has at least two inequivalent pure
complex lifts, distinguished by fixed-fibre sign.  This image theorem
does not enumerate the full descent fibre over any one class.
\end{theorem}
\begin{proof}
Use the mapping-torus source of
Theorem~\ref{p1b:thm:principal-descent-mapping-torus}.  The Wang sequence is surjective onto the monodromy-fixed orientation
class. Choose
\(v\in H^7(Q;\Z)\) restricting to the orientation generator of
\(H^7(F_7;\Z)=\Z\).  On an ordered scalar-product fibre the map
\(S^1_\phi\times P\to F_7\) is an orientation-preserving degree-two
quotient: reversal of \(\phi\) and sphere exchange each reverse
orientation.  Its pullback on degree seven is therefore
\(2\,d\phi\,u_xu_p\), with the angular class normalized to integral
one.  The Born-even integral generator
\((a_t-a_E)u_xu_p\) restricts to that same class, because the scalar
product fibre has \(a_t=d\phi,a_E=-d\phi\).  The invariant part of
\(H^7(M_8;\Z)\) has rank one, so
\[
 q^*v=(a_t-a_E)u_xu_p
\]
after orientation is chosen consistently.

Every degree-seven integral class on an eight-dimensional CW complex
is represented by a map to \(S^7\).  This follows from its Postnikov
tower: the first possible obstruction beyond \(K(\Z,7)\) lies in
degree nine.  Choose \(a_v:Q\to S^7\) representing \(v\), compose
with a degree-\(z\) self-map of \(S^7\), and then with the fibre
inclusion \(j:S^7\to BSU(2)\).  Its ordered pullback has zero
Chern classes, trivial G15 restriction, zero lower filling labels, and
central primary \((z,-z)\pmod{12}\).  A possible local top term does
not affect these assertions and can be changed independently by
Proposition~\ref{p1b:prop:pure-descent-local-top-constructor}.
This realizes all twelve zero-Chern Born-fixed primary classes in the
pure quotient category.

The remaining diagonal lower label is realized without changing fixed
isotropy.  Choose two small scalar torus discs interchanged by Born,
disjoint from the diagonal and from the defining scalar slice.  Over
one disc times \(P\), insert the relative determinant-one modification
represented by the actual gauge loop \(\alpha_x\), and over the
exchanged disc insert its Born transport \(\alpha_p\).  The
modifications are paired by exchange and agree with the original bundle
on their boundaries, so the resulting quotient bundle is again
principal \(SU(2)\).  Its ordinary fifth difference is
\(a_ta_E(u_x+u_p)\pmod2\), which toggles the two lower labels
simultaneously.  It changes neither second Chern class, G15 restriction,
nor central primary.  This is the free paired-disc modification already
used in the mixed gluing construction; its gauges are determinant one,
and its support avoids the fixed locus, so it applies equally to the
pure local model.  Any top correction is again adjusted independently
by the preceding local top constructor.

At zero Chern data, the established secondary orbit chart has exactly
three Born-fixed orbits at each fixed primary: the two top labels over
zero lower label, and the single orbit over diagonal nonzero lower
label.  The constructions above realize all three, hence all
\(12\cdot3=36\) ordinary classes.  A pure lift is Born fixed as an
ordinary class, so there are no others in the trivial sector.  The
all-Chern exotic restriction no-go excludes the other sector.
Finally the two central signs \(+I\) and \(-I\) cannot be gauge
conjugate on the connected fixed locus.  Multiplying a quotient pullback
lift by \(-I\) therefore gives a second inequivalent pure lift on
each of the thirty-six underlying bundles.
\end{proof}


\begin{lemma}[The symmetric-sphere loop image and its remaining normalization]
\label{p1b:lem:pure-descent-symmetric-loop-image}
Let \(A=\operatorname{SP}^2(S^3)\), let \(i:S^3\to A\) include a
coordinate sphere, and let \(q_Y:P=S^3\times S^3\to A\) be the quotient.
Then
\[
 [A,BSU(2)]=\{0\},\qquad
 \operatorname{im}\bigl(i^*:[A,SU(2)]\to[S^3,SU(2)]=\Z\bigr)=4\Z.
\]
The image of \(q_Y^*\) in the level-twelve Heisenberg coordinates is
infinite cyclic.  For its generator \(\ell\) of positive sphere degree,
\[
 \ell=(4,4;z_0),\qquad z_0\in\{2,8\}\pmod{12},\qquad
 \ell^r=(4r,4r;rz_0+8r(r-1)).
\]
There is only one actual value of \(z_0\) once the Heisenberg coordinates
are fixed; the displayed two candidates are the restriction imposed by
exchange invariance, not two asserted geometric generators.
\end{lemma}
\begin{proof}
The Steenrod adjunction for a suspended sphere gives
\(\operatorname{SP}^2(S^3)=\Sigma\operatorname{SP}^2(S^2)\cup e^6
=\Sigma\mathbb{CP}^2\cup e^6\).  In particular its five-skeleton is
\(S^3\cup_{\eta_3}e^5\).  Together with the symmetric-square
cofibration of the G21 proof, this gives
\[
 S^3\xrightarrow{i} A\longrightarrow M_5
 \xrightarrow{\partial}S^4,
 \qquad M_5=\Sigma^4\mathbb{RP}^2=M(\Z/2,5),
\]
where the restriction of \(\partial\) to the bottom \(S^5\) is
\(\eta_4\).  The adjunction formula is stated and proved in
\cite[Section 2]{SnaithUcci1973SymmetricMaps}; the cofibration is the
James--Thomas--Toda--Whitehead symmetric-square construction.

The target \(BSU(2)\) may be replaced by \(S^4\) on this
six-dimensional Moore source.  Its Moore exact sequence has the form
\[
 0\longrightarrow\pi_6(S^4)=\Z/2
 \longrightarrow[M_5,S^4]
 \longrightarrow\pi_5(S^4)=\Z/2\longrightarrow0.
\]
For an extension \(\bar\eta_4\) of \(\eta_4\), the mod-two Moore
identity is
\(2\bar\eta_4=\eta_4\eta_5\circ c\), where
\(c:M_5\to S^6\) collapses the bottom cell.  The right side is the
nonzero element of the injected \(\Z/2\); hence the middle group is
\(\Z/4\), and every class restricting to \(\eta_4\) is a generator.
Thus \(\partial\) generates \([M_5,BSU(2)]\).  The Puppe sequence
\[
 [S^4,BSU(2)]=\Z\longrightarrow[M_5,BSU(2)]=\Z/4
 \longrightarrow[A,BSU(2)]\longrightarrow[S^3,BSU(2)]=0
\]
has its first map onto, proving the first assertion.  Its preceding
part
\[
 [A,SU(2)]\longrightarrow[S^3,SU(2)]=\Z
 \longrightarrow[M_5,BSU(2)]=\Z/4
\]
then proves the degree-four restriction assertion, including existence
at every multiple of four.

Two maps \(A\to SU(2)\) of the same coordinate-sphere degree have a
pointwise quotient whose restriction along \(i\) is nullhomotopic.
The quotient therefore factors up to homotopy through \(M_5\).  Its
pullback along \(q_Y\) is null: after the coordinate wedge is collapsed,
\(P\to A\to M_5\) is the map
\(\Sigma^4(S^2\to\mathbb{RP}^2):S^6\to M_5\), which is null by
the stable-splitting argument already used in the G21 no-go proof.
Consequently the Heisenberg image has a unique element at each sphere
degree \((4r,4r)\), and is generated by its degree-four member.

The image is fixed by sphere exchange because \(q_Y\) is its quotient.
The established Heisenberg exchange formula sends
\((4,4;z_0)\) to \((4,4;16-z_0)\).  Equality gives
\(2z_0=16\pmod{12}\), namely \(z_0=2\) or \(8\).  Finally the
Heisenberg multiplication law gives the displayed power formula.
No step of this proof chooses between those two actual-map
normalization candidates.
\end{proof}


\subsubsection{Mixed coherent lifts with their actual boundary isotropy}
\label{p1b:sec:v2170-mixed}

Regroup \(M_8=K^2\), \(K=S^1\times S^3\), and let Born exchange
the two factors.  Its fixed locus is \(F=\Delta K\).  Remove an invariant
normal disc bundle \(N\cong F\times D^4\) and put
\[
 U=M_8\setminus\operatorname{int}N,\quad W=U/C_2,\quad
 R=\partial W=F\times L,\quad L=\mathbb{RP}^3.
\]
The normal action is \(-I\).  A mixed fixed-fibre involution has trivial
positive and negative eigenlines over \(F\), since \(H^2(F;\Z)=0\).
Its descended boundary bundle is
\(E_\partial=\mathbf1\oplus\lambda_L\), where
\(\lambda_L=\gamma\otimes\C\) is the sign line.  The retained
local model, rather than a coarse-quotient descent with trivial isotropy,
is indispensable to the classification.

\begin{proposition}[Exact mixed gluing problem]
\label{p1b:prop:v2170-mixed-gluing}
Mixed coherent rank-two bundles have
\[
 c_1=0,\qquad c_2=A\mathcal D,\qquad
 \mathcal D=(a_t-a_E)(u_x-u_p),\quad A\in\Z.
\]
They are classified by compatible actual bundles \(V\to W\) with
\[
 c_1(V)=\tau,\quad c_2(V)=(A\bmod2)\tau^2,\quad
 V|_R\cong E_\partial,
\]
and boundary isomorphisms modulo restriction of both the global gauge
group of \(V\) and the local mixed gauge group on \(N\).
The equivariant determinant is the sign line, up to an equivariant
normalization.  Every \(A\) occurs in either restriction sector.
\end{proposition}
\begin{proof}
Exchange reverses \(a_ta_E\), so invariance gives \(c_1=0\).
Writing the Born-invariant second Chern class as
\(A(a_tu_x+a_Eu_p)+B(a_tu_p+a_Eu_x)\), the mixed eigenline splitting
on \(F\) gives \(A+B=0\), hence the displayed \(\mathcal D\).
The normalized determinant has negative fixed character.  On the free
quotient this is \(\tau\).  Equivariant diagonal Thom restriction gives
\(c_2(V)=(A\bmod2)\tau^2\); its ordinary pullback is zero on \(U\).
Gluing the descended bundle to the prescribed local mixed model is exactly
the stated boundary double quotient.  Conversely such a gluing pulls back
to \(U\), attaches the local involution and supplies an order-two Born
action, not merely an isomorphism with its exchange pullback.

For existence, start with either the product mixed lift or the G21 exotic
mixed lift, constant on the torus.  In the cut-torus clutching model the
opposite edge degrees are \((A,-A)\) and their Born transports.  They
commute in the actual Heisenberg group and have a face filling: their
commutator exponent is zero.  Use the sign determinant seam, whose scalar
square is the identity.  Exchange the two halves of the filling with
that seam.  This is the coherent construction on each restriction
component, and has precisely the displayed Chern class.  The full relative
clutching construction, including its determinant normalization, is
retained in the bound derivation packet (Batches 4--5).
\end{proof}

\begin{lemma}[The free quotient cohomology and restriction maps]
\label{p1b:lem:v2170-w-cohomology}
Let \(\chi\) be the cover character, \(\tau=\beta\chi\),
\(z\) the exchanged real two-plane class and \(\eta_W=\beta z\).
Use \(\eta=\eta_W\) in the gauge proofs below; it is distinct from
the torus class \(\eta_T\) in the ordinary classification.  The retained
low integral groups and their actual restrictions are
\[
 H^1(W;\Z)=\Z\{h\},\quad H^2(W;\Z)=C_2\{\tau\},\quad
 H^3(W;\Z)=\Z\{q\}\oplus C_2\{\eta_W\},
\]
\[
 H^4(W;\Z)=\Z\{c_4\}\oplus C_2\{\tau^2\},\quad
 h|_R=2a_F,\quad q|_R=2u_F,\quad
 \eta_W|_R=a_F\tau,\quad c_4|_R=2a_Fu_F.
\]
Here \(p^*q=u_x+u_p\), \(c_4=\operatorname{tr}(a_tu_x)\),
\(z|_R=\bar a_F\chi\), and
\[
 z^2=0,\quad Sq^1z=\chi z,\quad
 \bar\eta_W=\chi z,\quad\bar\tau=\chi^2,\quad
 \tau h=\tau q=\tau c_4=0.
\]
Further,
\[
 H^5(W;C_2)=C_2,\quad
 H^6(W;C_2)\xrightarrow{\cong}H^6(P;C_2),\quad
 H^6(W;C_{12})=C_2\{a_6\},
\]
where \(a_6|_P=6u_xu_p\), \(a_6|_R=6u_Fv_L\), and its mod-two
reduction is zero.  The restriction
\(H^7(W;C_{12})\to H^7(Y;C_{12})\) is an isomorphism for the
exchange mapping torus \(Y\) of \(P\);
\(H^5(Y;C_2)=0\), \(c_4|_Y\) is a primitive generator, and
\(H^8(W;C)=0\) for every coefficient group used here.
\end{lemma}
\begin{proof}
Use the equivariant diagonal Thom sequence for
\(M_{hC_2}\) and \(F\times BC_2\), with normal Euler class
\(\tau^2\), and the free-complement identification with \(W\).
The degree-one and degree-three product modules are permutation modules;
degree two and degree six have sign action, and degree four is two
permutation modules.  Their cyclic coefficient differentials are
\(T-1\) and \(1+T\).  Transfer gives \(h,q,c_4\), with the stated
factor-two restrictions and primitive coinvariant on the mapping torus.
Since \(p^*\tau=0\), the transfer projection formula gives the three
vanishing products.  The exchanged two-plane has
\(Sq^1z=\chi z\) and \(z^2=0\), giving the torsion restriction and
reductions.  The finite-coefficient Thom sequence and mapping-torus
complex yield the displayed degree-six and degree-seven restriction
maps: on the top sphere product the exchange reverses orientation, so
the invariant \(C_{12}\) values are exactly \(0,6\).  Finally \(W\)
is an oriented eight-manifold with nonempty connected boundary, yielding
the top absolute vanishing.  The cell matrices and integral-transfer
calculations are retained in Batches 6--7 and 9--11; the integral fifth
group and its nonzero Bockstein are computed explicitly in
Lemma~\ref{p1b:rt:lemma121} below.
\end{proof}

\begin{theorem}[Exactly four compatible free-region bundles]
\label{p1b:thm:v2170-compatible-four}
For each \(A\), there are exactly four compatible ordinary classes
\(V\), two in each \(\varepsilon\)-sector.  At a chosen origin they
have labels \((A\bmod2,\varepsilon,\sigma)\),
\(\sigma\in C_2\).  The origins are coordinate choices, not canonical
member selections.
\end{theorem}
\begin{proof}
An exchanged pair of small torus-disc modifications inserts the actual
gauge loop \(\alpha_x\) and its Born transport, away from the diagonal
and the antipodal slice.  On \(W\) the pair descends to a single patch
\(D^2\times P\).  Its fifth difference is
\(t_D\bar u_x\).  By excision its Poincare--Lefschetz dual is the
closed complementary sphere \(Z_p\).  Since \(\bar q|_R=0\), a
relative lift of \(\bar q\) pairs with \(Z_p\) as one.  The fifth
difference is therefore the nonzero class.  It preserves \(A\) and
\(\varepsilon\).

A loop of four-stage data has \(\ell=kh\),
\(\xi=mq+t\eta_W\).  Its fifth-stage translation is
\[
 Sq^2\bar\xi+\bar\tau\bar\xi+
 \bar\ell\,\overline{c_2(V)}=0.
\]
The free term vanishes by stable transfer and \(\tau q=0\); the
torsion terms cancel because \(Sq^2(\chi z)=\chi^3z\); and the
determinant term vanishes by \(h\tau^2=0\).  Thus the patch's bit
survives ordinary isomorphism and gives two classes in each sector.

For the upper bound there are at most two fifth choices.  At a fixed
fifth representative the sixth torsor is \(H^6(W;C_2)=C_2\), whose
restriction to \(P\) is an isomorphism; hence at most one sixth choice
has a specified \(\varepsilon\).  The actual odd-primary returning
loop pairs with the primitive degree-four transfer class and identifies
the two seventh choices (the retained degree-seven section operation of
Batch 9).  The degree-eight ambiguity is zero.  This proves exactly two
per sector.  This argument uses the actual patch and the ordinary
restriction to \(P\); it does not use a boundary top-cell detector on an
arbitrary fifth representative.
\end{proof}

\begin{theorem}[The ordinary image of mixed coherent lifts]
\label{p1b:thm:v2170-mixed-image}
At \(c_1=0,c_2=A\mathcal D\), put \(h_m=\gcd(m,A)\).
The numbers of underlying ordinary classes admitting at least one mixed
coherent lift are exactly
\[
 I_0(A)=\begin{cases}h_{12},&A\text{ odd},\\
 4h_{12},&A\equiv2\pmod4,\\3h_{12},&4\mid A,
 \end{cases}\qquad
 I_1(A)=\begin{cases}2h_6,&A\text{ odd},\\3h_6,&A\text{ even}.
 \end{cases}
\]
An ordinary class belongs to this image precisely when, in the normalized
determinant-preserving torus chart, it has a representative with
\(z_t+z_E+A^2=0\pmod{12}\) and diagonal lower labels \(r=s\),
interpreted in its actual ordinary determinant quotient.  The sector-one
reduced primary equation is modulo six; a determinant change lifts it to
the displayed representative equation.
\end{theorem}
\begin{proof}
The sign determinant seam and coherent face filling give the displayed
primary equation and diagonal lower labels.  Its primary counts are
\(h_{12}\) and \(h_6\).  The patch in the preceding proof toggles both
lower labels.  An invariant local eight-ball also realizes the retained
ordinary top direction: equivariantly suspend the quaternionic Hopf map
\(S(4\mathbf1+4\sigma)\to S(3\mathbf1+2\sigma)\), then pull back
the rank-two complement bundle in \(\C^3\).  At a trivial pole the
local source is \(\R^4+4\sigma\) and the target is
\(\R^3+2\sigma\), exactly the fixed-point normal representation.
The complement has mixed eigenranks.  Its underlying ordinary top class
is the nonzero \(\pi_7U(2)\) value; the retained Hopf pullback calculation
of Batch 5 supplies that value.  The two modifications retain Chern and
primary data and realize every admissible secondary candidate.

Apply the actual shear \(e\mapsto e+br+as\) and the determinant
quotient of the ordinary theorem to the diagonal labels.  In sector zero
their secondary counts are one for odd \(A\), four for
\(A\equiv2\pmod4\), and three for \(4\mid A\).  In sector one
they are two and three for odd and even \(A\).  Multiplying the primary
counts gives the formula.  This counts underlying bundles, not multiple
coherent structures on the same bundle.
\end{proof}

\paragraph{The actual local boundary chart}
\label{p1b:sec:v2170-boundary-chart}
Put \(H=\mathcal G_{U(2)}(\mathbf1\oplus\lambda_L)\) and
\(G=SU(2)\).  The determinant-one gauge bundle has fibre sphere
\(S(\mathbf1\oplus\operatorname{ad}E_\partial)\), with
\(\operatorname{ad}E_\partial=\R\oplus2\gamma\),
\(w_2=\chi^2\).  It is this twisted bundle that enters the chart.

\begin{proposition}[Boundary normalization and natural cover]
\label{p1b:prop:v2170-boundary-chart}
In natural generators,
\[
 \pi_1H=\Z\{q_+\}\oplus C_4\{\kappa\},\quad
 \pi_2H=0,\quad \pi_3H=\Z\{b\}\oplus C_6\{z_6\},\quad
 \pi_4H=C_2\{e\}\oplus C_2\{t\}.
\]
Here \(\kappa=q_+-q_-\), \(\zeta=2\kappa\), sphere evaluation
sends \(b\) to degree two, and the natural sixth insertion
\(j_6:C_{12}\to C_6\) has kernel \(\{0,6\}\).
Every circle class is local: local eigenline windings \((m,n)\) give
\((m+n)q_+-n\kappa\), with kernel \((4j,-4j)\).
After local normalization the additive boundary group is
\[
 \mathscr A=\Z\{b\}\oplus T,\qquad
 T=C_6\{z_6\}\oplus C_2\{e\}\oplus C_2\{t\}.
\]
Its natural finite ordered-cover map is
\[
 \boxed{(z_6,e,t)\longmapsto(2z_6,e)\in C_{12}\oplus C_2.}
\]
It has domain order 24, image order 12 and kernel \(C_2\{t\}\).
\end{proposition}
\begin{proof}
The first adjoint operation is \(Sq^2+\chi^2\).  Its Moore-cell
connecting values are \(\eta_3\), \(\eta_3\eta_4\),
\(6\nu'\); in particular the evaluation degree is even and the
sixth insertion kills coefficient six.  The determinant splitting
\(q_+\) and the two local eigenline loops give the displayed circle
group and winding map.  The even local subgroup is central.  Odd local
conjugation changes only a removable circle shear, so the normalized
slice is additive.  The circle-cofibre calculation in
Theorem~\ref{p1b:rt:theorem114} below proves the actual fourth group and
survival of both natural generators; no associated-graded extension is
assumed.  The natural Hopf pullback gives cover \(e=\mu\), whereas
the degree-two lens top cell gives cover \(t=2\mu=0\).  The sixth
insertions pull back by degree two.  The cover formula and its kernel
follow.  The full coaction and component proof are supplied below.
\end{proof}

Choose actual global gauges with low data \((h,0)\), \((0,\eta_W)\)
and \((0,2q)\); denote them by \(g_h,g_\eta,g_2\).  Their existence
uses the oriented operation \(Sq^2+\chi^2\): its obstruction vanishes
on \(q\) by transfer and on \(\eta_W\) by cancellation.  The later
extensions are checked on \(P\) and the exchange mapping torus, with the
restriction isomorphisms above; \(H^8(W;C_2)=0\) finishes them.  The
determinant-primary compatibility eliminates odd free \(q\) coefficient
for an actual gauge.  Thus the actual joint low image is
\(\Z\{h\}\times(2\Z\{q\}\oplus C_2\{\eta_W\})\).
The restrictions of \(g_2\) have sphere integer two.  The explicit
scalar \(\sigma=fI\), \([f]=h\), has joint data \((2h,0)\),
since \(h\tau=0\), and local boundary \(4q_+-2\kappa\).
After dividing by this scalar, the zero-sphere global image is a finite
2-group.  Therefore it lies in the elementary subgroup
\(T[2]=\mathbb F_2\{a,e,t\}\), \(a=3z_6\).
The detailed low extension/slant calculations are retained in Batches
7--13; the higher choices and the remaining four generators are evaluated
or isolated by the actual gauge proofs below.

\begin{theorem}[Evaluated finite mixed gluing quotients]
\label{p1b:thm:v2170-mixed-quotient}
For every compatible actual \(V\), let \(J_V\subset\mathscr A\) be
its normalized global restriction image and \(J_{0,V}=J_V\cap T\).
Write \(a=3z_6\). The actual full and determinant-one-global images are
\[
\begin{array}{c|c|c|c|c}
 A\bmod2&\varepsilon&J_{0,V}^{\det=1}&J_{0,V}&|\mathscr A/J_V|\\\hline
0&0&\langle t\rangle&\langle t\rangle&24\\
0&1&\langle t\rangle&\langle a,t\rangle&12\\
1&0&0&0&48\\
1&1&\langle t\rangle&\langle a,t\rangle&12
\end{array}
\]
Each sector has two compatible free-region rows. Consequently
\[
 \boxed{N_{\rm mix}(A)=72\ (A\text{ even}),\quad
 N_{\rm mix}(A)=120\ (A\text{ odd}).}
\]
The auxiliary quotient by determinant-one global gauges and the full
local mixed gauge group has totals 96 and 144, respectively. It is not
a coherent principal-\(SU(2)\) classification.
The gluing quotient is \(\mathscr D(V,A)=\mathscr A/J_V\); it is a
torsor until an origin is chosen. Its finite image is the actual
four-column span of Theorem~\ref{p1b:rt:theorem127}.
\end{theorem}
\begin{proof}
The positive calculations below evaluate that span:
Theorems~\ref{p1b:rt:theorem129} and
\ref{p1b:rt:theorem30-column} give \(u_\chi=0\) and
\(u_c=(1+(A\bmod2))t\); the paired-image theorem
\ref{p1b:rt:theorem130} retains both fifth-bit rows.
The even image is Theorem~\ref{p1b:rt:theorem31-even}; the odd
sector-zero image is Theorem~\ref{p1b:rt:theorem32-sector0}; the
odd exotic determinant-one column is
Theorem~\ref{p1b:rt:theorem34-kappa}. The coincident-circle
calculation \ref{p1b:rt:theorem40-e} gives \(h_e=w_e=0\), and the
complete original natural comparison gives \(h_t=1\) in
Theorem~\ref{p1b:thm:natural-rank-two-completion}. Thus the odd exotic
column is \(a+t\), which spans \(\langle a,t\rangle\) with the
already realized \(t\).
The integer restriction has image \(2\mathbb Z\), and
\(|T|=24\); hence each row has \(48/|J_{0,V}|\) classes.
Summing two rows per sector gives \(2(24+12)=72\) and
\(2(48+12)=120\). Omitting the determinant column gives
\(2(24+24)=96\) and \(2(48+24)=144\), as in
Theorem~\ref{p1b:rt:theorem35-images}.
\end{proof}

\begin{theorem}[Chern-zero forgetful map and uniform lift fibres]
\label{p1b:thm:v2170-zero-fibres}
At \(A=0\), choose origins so that the fifth free-region bit is the
ordinary diagonal bit \(r=s=\delta\).  Its four ordinary images have
sizes \(24,12,12,6\), in \((\varepsilon,\delta)\) order
\(00,01,10,11\).  The actual normalized forgetful map is
\[
 (b,z_6,e,t)\longmapsto
 (b+2z_6\bmod m_\varepsilon,\ e\text{ if }\delta=0),
 \qquad m_0=12,\quad m_1=6.
\]
Thus each column of \(J_{0,V}\) lies respectively in
\[
 \langle t\rangle,\quad\langle e,t\rangle,\quad
 \langle a,t\rangle,\quad\langle a,e,t\rangle.
\]
Writing the corresponding actual ranks as \(s_{\varepsilon\delta}\),
every nonempty ordinary fibre in that image has the same multiplicity
\[
 \boxed{f_{\varepsilon\delta}
 =2^{1+\varepsilon+\delta-s_{\varepsilon\delta}},\qquad
 N_{\rm mix}(0)=24f_{00}+12f_{01}+12f_{10}+6f_{11}.}
\]
The ranks are \(s_{00}=s_{01}=1\) and \(s_{10}=s_{11}=2\). The multiplicities are therefore \(1,2,1,2\), respectively.
This is the Chern-zero specialization of the complete fibre theorem below.
\end{theorem}
\begin{proof}
Differences in gluings pull back to the ordered boundary.  Evaluation
parts extend over \(N\); the remaining cells have gauge dimensions six
and seven, so their bundle differences start in seven and eight and
cannot change \(\delta\).  The natural cover is \((2z_6,e)\), and
the transverse sphere generator contributes one to the ordinary primary.
The G20 top shear retains the ordinary top bit for \(\delta=0\) and
identifies it for \(\delta=1\).  This gives the displayed affine map
after choosing origins.  There are twelve primary classes in sector zero
and six in sector one, with two top values when \(\delta=0\), proving
the four image sizes.  Their lower/sector labels distinguish the images.

Any actual global boundary gauge represents an ordinary isomorphism,
so its finite part is in the kernel of this map.  Intersecting with
\(T[2]\) gives exactly the four stated allowed planes, of dimensions
\(1,2,2,3\).  On a chosen sphere-parity slice the finite cover map is
a homomorphism onto the corresponding primary coset and retained top
coordinate.  Its kernel has constant size, and dividing by the actual
global image preserves uniform fibres.  Equivalently divide
\(|\mathscr D(V,0)|=48/2^s\) by its image size
\(24,12,12,6\).  This gives the formula and bound eight.
It does not extend that bound to nonzero \(A\).
\end{proof}

The evaluated sphere relations and the quotient types appear in
Theorems~\ref{p1b:rt:theorem31-sphere} and
\ref{p1b:rt:theorem41-type}. In the natural odd exotic normalization,
the full comparison of Theorem~\ref{p1b:thm:natural-rank-two-completion}
gives
\[
 j_h=a+t,\qquad j_\eta=t,\qquad
 r_A=2b+z_Az_6+t,\qquad z_A=5-2A\pmod6.
\]
The odd-determinant translation is therefore trivial in sector zero
and a free involution in sector one. The auxiliary quotient uses the
full local mixed group even when its global gauges have determinant one.
The exact individual forgetful fibres and their fixed-row torsor groups
are summarized in Theorem~\ref{p1b:thm:complete-mixed-rank-two}.


\subsubsection{The adjoint conjugation operation}
\label{p1b:sec:rtproof23}

Put \(G=SU(2)=S^3\) and let
\(q(s)=\operatorname{diag}(e^{2\pi i s},1)\) be the loop of determinant winding one in \(U(2)\). Its conjugation commutator has values in \(G\):
\[
C(s,g)=q(s)gq(s)^{-1}g^{-1}.
\]
It is the identity on the wedge, so it defines a based map
\(S^1\wedge G\to G\). Let \(A:G\to\Omega G\) be its adjoint. The Bott Samelson value identifies this smash map with \(\eta_3:S^4\to S^3\).

\begin{lemma}
\label{p1b:rt:lemma104}
For \(x\in\pi_mG\), the adjoint of \(A_*x\) is
\[
\boxed{\eta_3\circ E x\in\pi_{m+1}G.}
\tag{RT.23.2}\label{p1b:eq:rt-23-2}
\]
In particular, with integral generators,
\[
A_*\iota_3=\eta_3,\quad
A_*\eta_3=\eta_3\eta_4,\quad
A_*(\eta_3\eta_4)=6\nu',\quad
\boxed{A_*\nu'=0.}
\tag{RT.23.3}\label{p1b:eq:rt-23-3}
\]
The last value differs from
\(\nu'\eta_6=\mu\ne0\in\pi_7G=C_2\).

\end{lemma}
\begin{proof}
The commutator applied to a representative \(x:S^m\to G\) is
\(C\circ(1\wedge x)\). Under the smash identification, \(1\wedge x\) is its suspension, proving \eqref{p1b:eq:rt-23-2}. The first three values follow from the Hopf generators and the integral relation \(\eta_3\eta_4\eta_5=6\nu'\).

Here is a derivation of the fourth value that keeps the unstable composition order explicit. Choose generators so that
\[
\pi_7S^4=\mathbb Z\{\nu_4\}\oplus C_{12}\{E\nu'\},\qquad
E\nu_4=\nu_5,\quad E^2\nu'=2\nu_5,
\quad \pi_8S^5=C_{24}\{\nu_5\}.
\]
Suspension therefore sends \((a,z)\) to \(a+2z\pmod{24}\); its kernel is generated by \((2,-1)\). The EHP sequence identifies this kernel with the Whitehead-square image. The Hopf invariant of the square is two, hence, up to simultaneous orientation choices,
\[
{}[\iota_4,\iota_4]=\pm(2\nu_4-E\nu').
\tag{RT.23.4}\label{p1b:eq:rt-23-4}
\]
Postcomposition with \(\eta_3\) sends this square to
\([\eta_3,\eta_3]=0\): multiplication in the topological group \(S^3\) extends the two maps from their wedge over their product. Also \(2\eta_3\nu_4=0\), since \(\pi_7S^3=C_2\). Applying \(\eta_3\) to \eqref{p1b:eq:rt-23-4} consequently gives \(\eta_3E\nu'=0\). This is precisely \(A_*\nu'\). In contrast the retained G20 generator is \(\mu=\nu'\eta_6\); the classical relation \(\eta_3\nu_4=\nu'\eta_6\) identifies its nonzero value. 
\end{proof}


The natural pointed map \(J\) is pullback of the Chern-class-one
sphere gauge along the Euler-number-two projection \(p:L\to S^2\).
The map \(j_7\) inserts the lens three-cell coefficient.  Set
\(\mu=\nu'\eta_6\ne0\in\pi_7SU(2)=C_2\), and define
\(e=\operatorname{im}(J\nu')\),
\(t=\operatorname{im}(j_7\mu)\) after unpointed evaluation.

\subsubsection{The actual circle coaction and fourth boundary group}
\label{p1b:sec:rtproof25}

\paragraph{The circle-cofibre coaction in the nontrivial bundle component}

Let
\[
L=\mathbb {RP}^3,\qquad
H=\mathcal G_{U(2)}(\mathbf1\oplus\lambda_L),\qquad
G=SU(2)=S^3,
\]
and let \(S_L\) denote the pointed gauge group. The determinant-one adjoint bundle is \(\mathbb R\oplus2\gamma\). The following calculation computes its pointed fourth-stage extension, including the component transport. In the natural generators of Proposition~\ref{p1b:prop:v2170-boundary-chart}, the resulting unpointed group is
\[
\boxed{\pi_4H=C_2\{e\}\oplus C_2\{t\}.}
\tag{RT.25.1}\label{p1b:eq:rt-25-1}
\]
Here \(e\) and \(t\) are the natural generators specified in Batch 23. Their natural cover values remain respectively \(\mu\) and zero. The full global gauge image is a separate calculation.

Take a circle fibre \(j:S^1\to L\) of the Euler-number-two circle bundle \(L\to S^2\). It represents the generator of \(\pi_1L=C_2\). Its homotopy cofibre
\[
C=L\cup_j CS^1\simeq L/S^1\simeq S^2_c\vee S^3_v
\]
is simply connected: the cone kills the one-cell, leaving an integral two-cycle and a three-cell with zero attaching degree. Its canonical cofibre coaction is
\[
\psi:C\longrightarrow C\vee S^2_u.
\]
Choose orientations so that the connecting map on the two-sphere has degree two.

\begin{lemma}
\label{p1b:rt:lemma112}
For suitable wedge generators, the coaction has the form
\[
\psi(c)=c+2u,\qquad
\psi(v)=v+a\eta_u+d[c,u],
\qquad a=2d,\quad d\equiv1\pmod2.
\tag{RT.25.2}\label{p1b:eq:rt-25-2}
\]
Changing \(v\) by \(k\eta_c\) changes \((a,d)\) to \((a+4k,d+2k)\). In particular one can choose \((a,d)=(2,1)\).

\end{lemma}
\begin{proof}
The projection of \(\psi\) onto \(C\) is the identity. The Hilton decomposition in degree three of \(S^2_c\vee S^3_v\vee S^2_u\) therefore leaves precisely the two terms \(a\eta_u\) and \(d[c,u]\) besides \(v\). There are no further words in this dimension.

Use two copies \(u_1,u_2\) of the cone sphere. Coassociativity of the cone coaction compares \((\psi\vee1)\psi\) with \((1\vee\Delta)\psi\). On \(v\), the first composite has coefficient \(2d\) on \([u_1,u_2]\), from \(d[c+2u_1,u_2]\). The second has coefficient \(a\), from the Hopf-Hilton formula
\[
\eta_{u_1+u_2}=\eta_{u_1}+\eta_{u_2}+[u_1,u_2].
\]
These words are independent in the Hilton decomposition, so \(a=2d\).

The parity of the mixed coefficient can be read directly from the cone construction. Write \(D\) for its cone disc and choose an integral two-chain \(E\) in \(L\) with \(\partial E=2j\). The two-cycle of \(C\) is represented by \(E-2D\). Push \(j\) off to a nearby circle fibre \(j'\), transverse to \(E\). In the coaction, splitting off the extra cone transports \(D\) across this push-off. At each signed intersection of \(E\) with \(j'\), the two cone sheets give the attaching map of the boundary of a product of two discs: the mixed Whitehead word \([c,u]\). Away from these intersections the transport is a product and contributes no mixed word. Altering the chosen two-cycle by the cone-boundary ambiguity changes this count by an even integer. Thus
\[
d\equiv E\mathbin{\cdot}j'\pmod2.
\tag{RT.25.3}\label{p1b:eq:rt-25-3}
\]
This is also the usual cellular definition of the mixed functional Hopf coefficient; it is the mixed word of the actual cone coaction, rather than a differential inferred from a spectral sequence.

Here the intersection is odd. The two lifted fibres in \(S^3\) are Hopf fibres, with linking number one. Transferring \(E\) to the double cover gives a chain with boundary twice the lifted \(j\); its intersection with the lifted \(j'\) is two. Transfer doubles the downstairs intersection, so \(E\cdot j'=1\), up to orientation. Equivalently, if \(\chi\in H^1(L;\mathbb F_2)\) is the cover character, the same parity is
\(\langle\chi Sq^1\chi,[L]_2\rangle=\langle\chi^3,[L]_2\rangle=1\).
This proves \eqref{p1b:eq:rt-25-3} and the asserted parity. Finally the Hopf-Hilton formula for \(\eta_{c+2u}\) gives the change \((4k,2k)\); since \(d\) is odd, the indicated normalization exists. 
\end{proof}

\paragraph{The actual looped connecting map}

Put \(Y=BU(2)\). In the wedge coordinates on \(C\), the bundle component is represented by \((\alpha,0)\), where \(\alpha:S^2\to Y\) has first Chern class one. Its pullback to \(L\) is \(\mathbf1\oplus\lambda_L\). There is no top component ambiguity since \(\pi_3Y=0\).

The circle cofibration gives the mapping-space sequence
\[
\Omega^2Y\xrightarrow{\delta_\alpha}
\operatorname {Map}_*(C,Y)_{(\alpha,0)}
\longrightarrow\operatorname {Map}_*(L,Y)_{\alpha|_L}
\longrightarrow\Omega Y.
\tag{RT.25.4}\label{p1b:eq:rt-25-4}
\]
The homotopy fibre of \(\delta_\alpha\), based at \((\alpha,0)\), is \(S_L\). All component choices below are taken at these indicated basepoints.

\begin{lemma}
\label{p1b:rt:lemma113}
After one loop and restriction to determinant-one factors, \eqref{p1b:eq:rt-25-4} gives
\[
\boxed{\Omega S_L\simeq
\operatorname {hofib}\bigl((2,\Omega^2A):
\Omega^2G\longrightarrow\Omega^2G\times\Omega^3G\bigr).}
\tag{RT.25.5}\label{p1b:eq:rt-25-5}
\]
The first factor is induced by the two-sphere of the circle cofibre; the second is induced by the lens top cell. The operation \(A:G\to\Omega G\) is the conjugation commutator of Batch 23, whose adjoint is \(\eta_3\).

\end{lemma}
\begin{proof}
The mapping-space connecting map is the orbit map for the canonical cone coaction. Formula \eqref{p1b:eq:rt-25-2}, in its normalized coordinates, therefore evaluates it as
\[
u\longmapsto\bigl(\alpha+2u,\ u\circ2\eta_2+[\alpha,u]\bigr).
\tag{RT.25.6}\label{p1b:eq:rt-25-6}
\]
Translate the first factor from the component \(\alpha\) to the zero component using the double-loop multiplication. The first entry becomes the power-two map. This translation does not discard the mixed term in the second entry.

The ordinary Hopf term in \eqref{p1b:eq:rt-25-6} becomes null after one loop: its adjoint precomposition map is \(\Sigma(2\eta_2)=2\eta_3=0\). The mixed term has one input fixed at \(\alpha\). Looping moves only its variable input, so its adjoint is the commutator with the clutching loop of \(\alpha\), of determinant winding one. More explicitly that clutching loop is \(q(s)=\operatorname {diag}(e^{2\pi i s},1)\), and the comparison on a variable \(g\) is
\(q(s)gq(s)^{-1}g^{-1}\). Its adjoint is \(A\). Thus the remaining looped second coordinate is \(\Omega^2A\), with orientation signs immaterial for the calculations here.

The determinant factors of \(\Omega^2U(2)\) and \(\Omega^3U(2)\) are contractible. Removing them yields \eqref{p1b:eq:rt-25-5}. A change of wedge splitting changes the coordinate maps by the explicitly described Hopf shear; it gives an equivalent homotopy-fibre calculation. In particular \eqref{p1b:eq:rt-25-5} has been derived from the actual cone coaction in the specified bundle component, not reconstructed from its associated graded groups. 
\end{proof}

\paragraph{The fourth-stage relation and the natural top generator}

Recall the integral coefficient values established in Batch 23:
\[
A_*\iota_3=\eta_3,\quad
A_*\eta_3=\eta_3\eta_4,\quad
A_*(\eta_3\eta_4)=6\nu',\quad
A_*\nu'=\eta_3E\nu'=0.
\tag{RT.25.7}\label{p1b:eq:rt-25-7}
\]
They distinguish the last conjugation value from \(\nu'\eta_6=\mu\ne0\).

\begin{theorem}
\label{p1b:rt:theorem114}
The pointed extension of Batch 22 is split, and its natural lens-top insertion survives unpointed evaluation. In natural notation,
\[
\pi_4S_L=C_2^2,\qquad 2J\nu'=0,\qquad
j_7\mu\ne0,\qquad \delta_5(\eta_3\eta_4)=0,
\tag{RT.25.8}\label{p1b:eq:rt-25-8}
\]
and \eqref{p1b:eq:rt-25-1} holds. Consequently the boundary rank is exactly \(r=2\).

\end{theorem}
\begin{proof}
The relevant exact portion of \eqref{p1b:eq:rt-25-5}, expressed in the unlooped degree of \(S_L\), is
\[
C_{12}\xrightarrow{(2,0)}C_{12}\oplus C_2
\longrightarrow\pi_4S_L
\longrightarrow C_2\xrightarrow{(0,6)}C_2\oplus C_{12}.
\tag{RT.25.9}\label{p1b:eq:rt-25-9}
\]
The final arrow is injective. The preceding cokernel is
\[
(C_{12}/2C_{12})\oplus C_2=C_2\oplus C_2.
\]
The second generator is precisely the top-cell insertion \(j_7\mu\), since the second factor in \eqref{p1b:eq:rt-25-5} was obtained from that cell. It is nonzero. The sphere-pullback class \(J\nu'\) projects to the first generator; a change of extended Moore pinch can add the second generator, but cannot change its order two or its independence. This proves the pointed assertions with the natural \(J\) used in Batch 23.

For unpointed evaluation, Batch 22 identifies \(\delta_5\) as \(6J\nu'\). This is now zero. Its next connecting map \(\delta_4\) is injective on \(\pi_4U(2)=C_2\), with image the complementary evaluation element in \(\pi_3S_L\). Hence evaluation introduces neither a quotient of \(\pi_4S_L\) nor a further element in \(\pi_4H\). This proves \eqref{p1b:eq:rt-25-1}. 
\end{proof}

The same looped sequence checks the earlier positive-degree boundary groups. In degrees two and three, the consecutive coefficient maps are \((0,1):C_2\to C_2^2\) and \((0,6):C_2\to C_2\oplus C_{12}\), giving \(\pi_2S_L=C_2\) and \(\pi_3S_L=C_2\oplus C_6\). Evaluation removes the displayed complementary \(C_2\), and its primary kernel has degree two. Thus
\[
\pi_1H=\mathbb Z\oplus C_4,\quad \pi_2H=0,\quad
\pi_3H=\mathbb Z\oplus C_6,\quad \pi_4H=C_2^2,
\tag{RT.25.10}\label{p1b:eq:rt-25-10}
\]
consistent with the actual circle and sphere calculations of Batches 21–23.

\begin{corollary}
\label{p1b:rt:corollary115}
The actual normalized finite cover map is now the fixed map
\[
\boxed{T=C_6\{z\}\oplus C_2\{e\}\oplus C_2\{t\}
\longrightarrow C_{12}\oplus C_2,
\quad (z,e,t)\longmapsto(2z,e).}
\tag{RT.25.11}\label{p1b:eq:rt-25-11}
\]
It has domain order 24, image order 12, and kernel \(C_2\{t\}\).

\end{corollary}
\begin{proof}
Substitute the proved value \(r=2\) into the natural cover formula of Batch 23. Its proof uses the actual Hopf pullback and frame comparison, so the generators here are the specified natural ones. 
\end{proof}


\subsubsection{The supported sixth gauge and its whole boundary restriction}
\label{p1b:sec:rtproof26}

\paragraph{A supported gauge on the circle-pair Möbius band}

Retain the oriented eight-manifold
\[
W=((S^1\times S^3)^2\setminus\operatorname {int}N)/C_2,
\qquad R=\partial W=F\times L,
\quad F=S^1\times S^3,\quad L=\mathbb {RP}^3.
\]
For any compatible free-part bundle \(V\to W\), let \(\mathcal K_{5,V}\) denote its determinant-one global gauges fixed to the identity through stage five. Batch 25 fixes the actual boundary group
\(\pi_4H=C_2\{e\}\oplus C_2\{t\}\), with \(t\ne0\).
The following calculation computes
\[
\boxed{\operatorname {im}(\mathcal K_{5,V}\longrightarrow\mathscr A)=0
\quad\text{for every compatible }V.}
\tag{RT.26.1}\label{p1b:eq:rt-26-1}
\]
The normalized boundary group is
\(\mathscr A=\mathbb Z\{b\}\oplus C_6\{z\}\oplus C_2\{e\}\oplus C_2\{t\}\).
Equation \eqref{p1b:eq:rt-26-1} concerns the indicated high-stage kernel. Lower-stage gauges can still have nonzero restrictions in this group.

Fix a sphere point \(x_0\). The unordered distinct pairs of circle points, with both sphere coordinates equal to \(x_0\), form the proper circle-pair Möbius band \(C\subset W\) used in Batches 7, 10 and 11. Its boundary lies on the diagonal collar:
\[
\partial C=S^1_F\times\{x_0\}\times\{\ell_0\}\subset R.
\tag{RT.26.2}\label{p1b:eq:rt-26-2}
\]
Here \(\ell_0\) is the projective normal direction along the circle difference. In particular this boundary circle winds once around the fixed circle of \(F\). The band retracts to its core.

\begin{lemma}
\label{p1b:rt:lemma116}
The rank-six normal bundle of \(C\) satisfies
\[
\nu_C\cong\mathbb R^5\oplus\gamma_C,
\qquad \operatorname {Th}(\nu_C)\simeq
P^7(2)=S^6\cup_2 e^7.
\tag{RT.26.3}\label{p1b:eq:rt-26-3}
\]
Every rank-two complex bundle \(V\) is trivial on a tubular neighborhood of \(C\).

\end{lemma}
\begin{proof}
Since \(W\) is oriented, \(w_1\nu_C=w_1TC\), the nonzero character of the core circle. A real bundle of fixed positive rank over a circle is classified by this character, which proves the first isomorphism. The Thom space of the nontrivial real line over the core circle has a bottom one-cell and a two-cell attached with degree two. Suspending five times gives the displayed Moore space. Retraction of the band and its normal neighborhood onto the core preserves these bundles and their Thom homotopy type.

Complex bundles over the band are classified over a circle. Because \(U(2)\) is connected, every such bundle is trivial; a tubular neighborhood retracts to the band. Thus a unitary frame exists there for the actual \(V\), without a global splitting assumption. 
\end{proof}

\begin{lemma}
\label{p1b:rt:lemma117}
There is an actual supported determinant-one gauge with sixth difference class \(a_6\), the nonzero class in \(H^6(W;C_{12})=C_2\).

\end{lemma}
\begin{proof}
Choose a based map
\[
f:P^7(2)\longrightarrow SU(2)
\quad\text{whose bottom restriction is }6\nu'\in\pi_6SU(2)=C_{12}.
\tag{RT.26.4}\label{p1b:eq:rt-26-4}
\]
It exists because the sole attaching obstruction is
\(2(6\nu')=12\nu'=0\). Use the frame of Lemma~\ref{p1b:rt:lemma116} to put this map on the tubular neighborhood via its normal Thom collapse, and use the identity outside the neighborhood. The Thom basepoint is the identity, so the two sections paste to an actual automorphism of \(V\) of determinant one.

The normal coefficient is six; reversing a local normal orientation changes it to \(-6=6\) in \(C_{12}\). Its degree-six difference cocycle is therefore the untwisted coefficient-six Poincaré dual of the proper band. The band meets the sphere-product slice \(P=S^3_x\times S^3_p\) transversely once at an antipodal circle pair. The restriction of this cocycle is \(6u_xu_p\), which detects the nonzero class \(a_6\) by the restriction calculation of Batch 10.

The Thom space has no cells below degree six. Equivalently, a five-skeleton can be placed disjoint from the two-dimensional band by general position in dimension eight, after which the tube can be made sufficiently small. The resulting gauge has the identity lower-stage section data through degree five. This supplies a concrete realization of the nonzero sixth class for every compatible actual \(V\). 
\end{proof}

\paragraph{The boundary Thom map and its two coefficient values}

Near \(\partial C\) choose a product collar, so restriction of the tubular gauge is governed by the Thom inclusion
\[
\operatorname {Th}(\nu_C|_{\partial C})\longrightarrow
\operatorname {Th}(\nu_C).
\tag{RT.26.5}\label{p1b:eq:rt-26-5}
\]
The normal bundle is oriented on the boundary circle and is trivial there. Its boundary Thom space is consequently \(S^6\vee S^7\).

\begin{lemma}
\label{p1b:rt:lemma118}
In any such boundary trivialization the composite of \eqref{p1b:eq:rt-26-5} with \(f\) has coefficient pair
\[
\boxed{(6\nu',0)\in\pi_6SU(2)\oplus\pi_7SU(2).}
\tag{RT.26.6}\label{p1b:eq:rt-26-6}
\]
The second value is independent of the extension in \eqref{p1b:eq:rt-26-4}, the normal framing, and the chosen complex bundle frame.

\end{lemma}
\begin{proof}
On the bottom sphere the inclusion is the bottom map
\(i:S^6\to P^7(2)\), up to orientation. Its composite is thus \(\pm6\nu'=6\nu'\).

For the top sphere, the source Moore space is in the stable range in degree seven. Its cofibration gives
\[
\pi_7S^6=C_2\xrightarrow{2}C_2
\xrightarrow{i_*}\pi_7P^7(2)
\longrightarrow\pi_6S^6=\mathbb Z\xrightarrow{2}\mathbb Z.
\]
The first map is zero and the last is injective. Hence
\[
\pi_7P^7(2)=C_2\{i\eta_6\}.
\tag{RT.26.7}\label{p1b:eq:rt-26-7}
\]
Whatever value the top restriction of \eqref{p1b:eq:rt-26-5} has in this two-element group, its image under \(f\) is zero:
\[
f_*i\eta_6=6\nu'\eta_6=6\mu=0
\quad\text{in }\pi_7SU(2)=C_2.
\tag{RT.26.8}\label{p1b:eq:rt-26-8}
\]
In particular a choice of nullhomotopy extending \eqref{p1b:eq:rt-26-4} cannot supply an additional top value on the boundary; every top map into its Thom source already factors through \eqref{p1b:eq:rt-26-7}.

A change of rank-six normal frame over the boundary circle can add its stable \(J\)-value to the top sphere, namely an \(i\eta_6\) term. Equation \eqref{p1b:eq:rt-26-8} kills it. A change of the complex frame acts on the supported bottom gauge by conjugation with a loop in \(U(2)\). The only loop invariant is its determinant winding, and the resulting top comparison is a multiple of \(A_*(6\nu')\). Batch 23 gives \(A_*\nu'=0\), so this term also vanishes. These checks allow us to express \eqref{p1b:eq:rt-26-6} in the fixed product boundary coordinates of \eqref{p1b:eq:rt-26-2}, rather than in an arbitrarily transported tube frame. 
\end{proof}

\paragraph{The actual sixth-stage restriction and the entire high-stage image}

\begin{theorem}
\label{p1b:rt:theorem119}
For the actual supported gauge of Lemma~\ref{p1b:rt:lemma117}, its full boundary component is the identity. Thus the value left open in Batch 23 is
\[
\boxed{\gamma_{6,V}=0,\qquad\beta_V=0}
\tag{RT.26.9}\label{p1b:eq:rt-26-9}
\]
for all compatible bundles, in both parities of \(A\) and both ordinary sectors and fifth-stage choices.

\end{theorem}
\begin{proof}
Use product normal coordinates around the circle \eqref{p1b:eq:rt-26-2}. The pair \eqref{p1b:eq:rt-26-6} describes a circle family of gauges supported in a six-disc in \(S^3_F\times L\). Its bottom class is the constant circle family of the lens-top insertion \(j_6(6)\in\pi_3H\); its varying circle term is the displayed zero \(\pi_7SU(2)\) coefficient. Lemma~\ref{p1b:rt:lemma118} ensures that neither the normal frame nor the unitary frame introduces an additional family term.

The actual natural insertion is
\(j_6:C_{12}\to\pi_3H_{\mathrm{tors}}=C_6\), with kernel \(\{0,6\}\), as established in Batch 22. Therefore \(j_6(6)=0\). The boundary circle restriction and primary sphere degree are already zero for a gauge fixed through degree five. The zero top family term then leaves zero in \(\pi_4H\) as well. All boundary component coordinates vanish, proving \eqref{p1b:eq:rt-26-9} for this actual gauge.

Any other gauge with the same sixth difference class and the same fixed lower section differs, after choosing the common lower homotopy, by seventh-stage data. Its restriction lies in the image of \(H^7(W;C_2)\to H^7(R;C_2)\), which is zero. Indeed \(W\) is oriented and \(R\) is connected, so the boundary fundamental class comes from the relative fundamental class and every restricted top cohomology class pairs with it as zero. The boundary is seven-dimensional; later differences contribute no boundary coordinate. Thus the value of \(\gamma_{6,V}\) is the same for the earlier sixth-class realization of Batch 10. By Batch 23 it was \(\beta_Vt\), and Theorem~\ref{p1b:rt:theorem114} proves \(t\ne0\). Hence \(\beta_V=0\). 
\end{proof}

\begin{corollary}
\label{p1b:rt:corollary120}
Equation \eqref{p1b:eq:rt-26-1} holds exactly.

\end{corollary}
\begin{proof}
The two sixth difference classes are zero and \(a_6\). The zero class restricts to the identity by the same seventh-stage argument, and Theorem~\ref{p1b:rt:theorem119} proves that the nonzero class does too. Both have actual realizations, and there are no other sixth classes. This computes the whole image of \(\mathcal K_{5,V}\), not merely the value of one chosen representative. 
\end{proof}


\subsubsection{Integral fifth torsion and the stage-four-kernel image}
\label{p1b:sec:rtproof27}

\paragraph{The nonzero integral fifth class}

Retain
\[
W=((S^1\times S^3)^2\setminus\operatorname {int}N)/C_2,
\quad R=\partial W=F\times\mathbb {RP}^3,
\]
the cover character \(\chi\), the sign-line Chern class \(\tau\), and the exchanged-plane class \(z\). Write
\[
\eta=\beta z\in H^3(W;\mathbb Z),\qquad
\alpha=\chi^2z\in H^4(W;\mathbb F_2).
\]
For a compatible actual bundle \(V\), let \(\mathcal K_{j,V}\) denote determinant-one gauges with identity section data fixed through fibrewise Postnikov degree \(j\). Corollary~\ref{p1b:rt:corollary120} computes the entire boundary image of \(\mathcal K_{5,V}\) as the identity. The following argument proves the stronger statement
\[
\boxed{\operatorname {im}\bigl(\mathcal K_{4,V}\to
\pi_0\mathcal G(V|_R)\bigr)=\{1\}}
\tag{RT.27.1}\label{p1b:eq:rt-27-1}
\]
for every compatible actual \(V\). The full lower-stage global image remains a separate calculation.

\begin{lemma}
\label{p1b:rt:lemma121}
The integral fifth cohomology group and the relevant Bockstein are
\[
\boxed{H^5(W;\mathbb Z)=C_2\{\tau\eta\},\quad
\beta\alpha=\tau\eta\ne0,\quad
Sq^1\alpha=\chi^3z\ne0.}
\tag{RT.27.2}\label{p1b:eq:rt-27-2}
\]
Consequently
\[
H^4(W;\mathbb F_2)
=\mathbb F_2\{\chi^4,\bar c_4,\alpha\},
\quad Sq^1(\chi^4)=Sq^1\bar c_4=0,
\tag{RT.27.3}\label{p1b:eq:rt-27-3}
\]
and \(\chi^3z\) generates \(H^5(W;\mathbb F_2)=C_2\).

\end{lemma}
\begin{proof}
Use the equivariant diagonal Thom sequence before taking the free complement. Put \(M=(S^1\times S^3)^2\), \(B=M_{hC_2}\), and \(F_{hC_2}=F\times BC_2\). Denote the sign-line class on \(B\) by \(t\), its mod-two cover class by \(\chi_B\), and the exchanged-plane class by \(z_B\). These restrict to \(\tau,\chi,z\) on \(W\). The diagonal equivariant Thom class \(\delta=i_!(1)\) satisfies
\[
i^*\delta=t^2,\qquad
\delta|_M=(a_t-a_E)(u_x-u_p)=\mathcal D.
\tag{RT.27.4}\label{p1b:eq:rt-27-4}
\]

The integral Cartan–Leray sequence for \(B\) has the following nonzero candidates in total degree five:
\[
E_2^{0,5}=\mathbb Z\{a_ta_E(u_x-u_p)\},\qquad
E_2^{3,2}=C_2.
\tag{RT.27.5}\label{p1b:eq:rt-27-5}
\]
Indeed the degree-one and degree-three modules of \(M\) are permutation modules, its degree-four module is a sum of two permutation modules, its degree-two module is the sign module, and its degree-five module is a signed permutation module. Positive cyclic cohomology of a permutation module vanishes, while the sign module has \(C_2\) in odd positive degrees and zero in even positive degrees. All other degree-five entries vanish.

The free entry survives: the equivariant Gysin class \(i_!(a_F)\) forgets to \(\pm a_ta_E(u_x-u_p)\), a primitive invariant integral class. The torsion entry also survives. The actual characteristic class
\[
t\beta z_B\in H^5(B;\mathbb Z)
\]
has fixed-locus restriction \(a_Ft^2\ne0\), since \(z_B|_{F_{hC_2}}=\bar a_F\chi_B\) and \(\beta z_B|_{F_{hC_2}}=a_Ft\). It is a nonzero class killed by two, so it realizes the only possible torsion entry of \eqref{p1b:eq:rt-27-5}. Thus \(H^5(B;\mathbb Z)=\mathbb Z\oplus C_2\), and \(i_!(a_F)\) has primitive free coordinate.

In total degree six the only entry is the base class \(C_2\{t^3\}\). It survives by the fixed-point section. The Thom sequence contains
\[
H^1(F_{hC_2};\mathbb Z)=\mathbb Z
\xrightarrow{i_!}\mathbb Z\oplus C_2
\longrightarrow H^5(W;\mathbb Z)
\longrightarrow H^2(F_{hC_2};\mathbb Z)=C_2
\xrightarrow{i_!}C_2\{t^3\}.
\tag{RT.27.6}\label{p1b:eq:rt-27-6}
\]
The last map is injective: \(i_!(t)=t\delta\) restricts to \(t^3\ne0\) on the fixed locus. The first map is primitive in the free coordinate, regardless of its torsion coordinate. Its cokernel is therefore exactly \(C_2\), proving the group assertion. This also rules out an order-four integral fifth class, which would change the reduction argument below.

Now \(\alpha|_R=\bar a_F\chi_R^3\ne0\). Every integral fourth class is an integer combination of \(c_4\) and \(\tau^2\). Their boundary restrictions are respectively \(2a_Fu_F\) and zero, so all their mod-two reductions restrict to zero. Hence \(\alpha\) has no integral lift and \(\beta\alpha\ne0\). The Bockstein product formula, using the integral lift \(\chi^2=\bar\tau\), gives
\[
\beta(\chi^2z)=\tau\beta z=\tau\eta.
\]
Reduction on \(H^5(W;\mathbb Z)=C_2\) is injective, since its kernel is the image of multiplication by two. Thus \(Sq^1\alpha=\overline{\beta\alpha}\ne0\). The Cartan formula and \(Sq^1z=\chi z\) identify it as \(\chi^3z\).

Finally integral universal coefficients give the dimension three of \(H^4(W;\mathbb F_2)\). Its two integral-reduction directions are \(\chi^4=\overline{\tau^2}\) and \(\bar c_4\), which are independent: the latter restricts nontrivially to the exchange mapping torus \(Y\), whereas the former does not. The boundary restriction distinguishes \(\alpha\) from both. Integral reductions have zero \(Sq^1\), giving \eqref{p1b:eq:rt-27-3}. Batch 6 supplies \(H^5(W;\mathbb F_2)=C_2\), so the nonzero value in \eqref{p1b:eq:rt-27-2} is its generator. 
\end{proof}

\paragraph{Squaring an actual fourth-stage gauge}

Let \(g_{\alpha,V}\) be the actual gauge constructed in Batch 11, with primary zero and fourth difference \(\alpha\). Its boundary circle class is
\(\zeta=2\kappa\in C_4\{\kappa\}\). The symbol \(\zeta\) denotes an actual local loop; it is not an additional independent circle generator.

\begin{lemma}
\label{p1b:rt:lemma122}
If a determinant-one gauge has fixed identity primary data and fourth difference \(a\in H^4(W;\mathbb F_2)\), its square admits the induced identity homotopy through degree four whose fifth difference is
\[
\boxed{Sq^1a\in H^5(W;\mathbb F_2).}
\tag{RT.27.7}\label{p1b:eq:rt-27-7}
\]
In particular \(g_{\alpha,V}^2\) realizes the nonzero fifth difference \(\chi^3z\).

\end{lemma}
\begin{proof}
The fibre groups through degree five are
\(\pi_3S^3=\mathbb Z\), \(\pi_4S^3=C_2\{\eta_3\}\), and
\(\pi_5S^3=C_2\{\eta_3\eta_4\}\). With the primary fixed to zero, the last two rows agree with their stable rows. No product of two differences first supported in degree four can contribute in degree five.

The local relation across a degree-two cellular attaching map is computed on the Moore space
\(P^5(2)=S^4\cup_2e^5\). If \(i:S^4\to P^5(2)\) and \(p:P^5(2)\to S^5\) are its inclusion and pinch, the Moore identity is
\[
2\operatorname {id}_{P^5(2)}=i\eta_4p.
\tag{RT.27.8}\label{p1b:eq:rt-27-8}
\]
It holds unstably in this degree by the stable range. Choose an extension \(\bar\eta_3:P^5(2)\to S^3\) of \(\eta_3\); it exists since \(2\eta_3=0\). Composing \eqref{p1b:eq:rt-27-8} gives
\[
2\bar\eta_3=\eta_3\eta_4p.
\tag{RT.27.9}\label{p1b:eq:rt-27-9}
\]
Thus doubling the fourth coefficient produces exactly the next Hopf coefficient on each mod-two Bockstein incidence of the cellular cocycle. Adding these incidences yields \(\overline{\beta a}=Sq^1a\). Changing the original fifth choice adds a class of exponent two and cannot change its square.

This construction also applies to the actual automorphism sphere bundle. Its rank-three adjoint bundle is oriented, so coefficient transport adds no degree-one orientation term. Its nonzero \(w_2=\bar\tau\) enters two-degree obstruction operations, not the one-degree coefficient Bockstein in \eqref{p1b:eq:rt-27-7}. One may equivalently apply the Thom formula \(Sq^1(aU)=Sq^1(a)U\), since \(Sq^1U=w_1U=0\). The induced lower nullhomotopy is exactly the one used in the cellular doubling construction. This proves the asserted fifth difference for an actual section, without claiming that the whole eight-dimensional section problem stabilizes. Lemma~\ref{p1b:rt:lemma121} supplies the final nonzero value. 
\end{proof}

\paragraph{The entire stage-four-kernel boundary image}

\begin{theorem}
\label{p1b:rt:theorem123}
Equation \eqref{p1b:eq:rt-27-1} holds for every compatible actual \(V\). More specifically, every gauge in \(\mathcal K_{4,V}\) has the same boundary component as either the identity or \(g_{\alpha,V}^2\), and both are the identity on the boundary up to gauge homotopy.

\end{theorem}
\begin{proof}
With the lower section fixed through degree four, the first difference group is \(H^5(W;\pi_5S^3)=C_2\). Lemma~\ref{p1b:rt:lemma122} realizes its nonzero difference with the actual square \(g_{\alpha,V}^2\). Consequently a gauge with fifth class zero differs from the identity by stage-five-kernel data; one with fifth class nonzero differs from that square by stage-five-kernel data. Any identifications caused by returning lower homotopies only reduce these choices. Batch 26 proves that every such stage-five-kernel difference has identity boundary component.

It remains to evaluate the square. The gauge \(g_{\alpha,V}\) has determinant zero and free primary zero, hence boundary sphere integer \(b=0\). Its circle class is \(\zeta\). The even local circle subgroup is central in the actual boundary group, and its quotient is the additive group
\[
\mathscr A=\mathbb Z\{b\}\oplus
C_6\{z_6\}\oplus C_2\{e\}\oplus C_2\{t\}.
\]
Batch 22 proves that the normalized restriction of every zero-free-primary global gauge belongs to its elementary order-two subgroup \(T[2]\). Therefore the normalized restriction of \(g_{\alpha,V}\) squares to zero. Its actual circle restriction also squares to zero, because \(2\zeta=4\kappa=0\). A boundary class with trivial circle lies in the injectively represented zero-circle slice of \(\mathscr A\); thus the full restriction of the square is the identity.

These two actual fifth values exhaust the first difference group, and their remaining stage-five-kernel differences have zero boundary image. This proves \eqref{p1b:eq:rt-27-1} for the full boundary component group, not only after local circle normalization. 
\end{proof}

The consequence for subsequent calculations is exact: once the primary and fourth-stage section data of an actual gauge have been specified, fifth and later choices cannot change its boundary component. The fourth-stage gauge itself may still have a finite correction; Section~\ref{p1b:sec:rtproof28} computes it using a torsion-primary square.


\subsubsection{The torsion-primary square and the four remaining restriction columns}
\label{p1b:sec:rtproof28}

\paragraph{The torsion-primary square with its adjoint-bundle twist}

For every compatible actual \(V\to W\), retain the determinant-one gauge \(g_{\eta,V}\) of Batch 13 with pure primary class \(\eta=\beta z\). Its boundary circle class is an odd multiple of \(\kappa\), and hence its square has circle class \(\zeta=2\kappa\). Theorem~\ref{p1b:rt:theorem123} proves that all gauges fixed through degree four have identity boundary restriction. The following calculation computes the whole remaining boundary value of the fourth-stage gauge:
\[
\boxed{[g_{\alpha,V}|_R]=\zeta,\qquad
\alpha=\chi^2z,}
\tag{RT.28.1}\label{p1b:eq:rt-28-1}
\]
with no sphere-torsion or parameter-top correction. It then reduces the full finite global image to four specified actual restriction columns. The following adopted calculations evaluate their values without changing this four-column reduction.

\begin{lemma}
\label{p1b:rt:lemma124}
The square of a determinant-one gauge with primary \(\eta\) has primary zero and fourth difference
\[
\boxed{a_4(g_{\eta,V}^2)=Sq^2z+\bar\tau z
=\chi^2z=\alpha.}
\tag{RT.28.2}\label{p1b:eq:rt-28-2}
\]
The value is independent of the fourth and later choices in \(g_{\eta,V}\), and of the nullhomotopy chosen to remove its doubled primary.

\end{lemma}
\begin{proof}
First work only through fibrewise degree four. The primary and fourth coefficient rows are the stable rows \(\mathbb Z\) and \(C_2\{\eta_3\}\). Their oriented adjoint-bundle first operation is
\[
D=Sq^2+\bar\tau\smile-.
\tag{RT.28.3}\label{p1b:eq:rt-28-3}
\]
This is the operation already derived by clutching in Batches 7 and 13. It can also be obtained from the adjoint Thom class \(U\): its orientation gives \(Sq^1U=0\), and \(Sq^2U=w_2(\operatorname {ad}V)U=\bar\tau U\). Hence \(Sq^2(xU)=D(x)U\).

The coefficient Bockstein diagram for these two rows computes the doubled lift of an order-two primary. To make the secondary value explicit, represent \(\eta\) by an integral three-cocycle \(E\) and choose an integral two-cochain \(Z\) with \(dZ=2E\). Its reduction represents a Bockstein preimage \(z\). Doubling the primary section and using this cochain to cancel \(2E\) leaves, on each four-cell, the first Hopf obstruction applied to \(Z\bmod2\). In a trivial frame this is \(Sq^2z\), the ordinary two-stage cohomotopy extension coefficient. Transport of the oriented adjoint frame adds \(w_2(\operatorname {ad}V)z\); applying the displayed Thom identity gives precisely \(D(z)\).

Equivalently, in the relative disk filling of the doubled clutching loop, the reduced low characteristic cochain has the form \(v\smile z\), with \(v\) the relative parameter two-class. Slanting the first operation gives
\[
(Sq^2+\bar\tau)(v z)=v\,(Sq^2z+\bar\tau z),
\]
because the parameter disk has no higher squares. This is the fourth coefficient of the square, rather than the degree-five obstruction to realizing the original primary. It identifies the order-two extension coefficient with \(D\) of the Bockstein preimage.

Different cancellations change \(z\) by the reduction of an integral two-class. Here \(H^2(W;\mathbb Z)=C_2\{\tau\}\), and
\[
D(\bar\tau)=D(\chi^2)=\chi^4+\chi^4=0.
\tag{RT.28.4}\label{p1b:eq:rt-28-4}
\]
Thus the fourth value has no such ambiguity. Adding a fourth-stage choice to the original gauge adds twice an exponent-two coefficient on squaring; products of two primary differences first affect degree six, so they add no fourth term. Later choices also cannot change this value. Finally the actual exchanged-plane identity \(z^2=0\) gives \(Sq^2z=z^2=0\), proving \eqref{p1b:eq:rt-28-2}. The calculation uses only the first two stable rows, not a stabilization of the whole global gauge problem. 
\end{proof}

\paragraph{The fourth-stage gauge has a purely local boundary value}

\begin{theorem}
\label{p1b:rt:theorem125}
Equation \eqref{p1b:eq:rt-28-1} holds for every actual determinant-one gauge with fixed zero primary and fourth difference \(\alpha\). In normalized finite coordinates its restriction is
\[
\boxed{(b,z_6,e,t)=(0,0,0,0).}
\tag{RT.28.5}\label{p1b:eq:rt-28-5}
\]

\end{theorem}
\begin{proof}
Lemma~\ref{p1b:rt:lemma124} shows that \(g_{\eta,V}^2\) and \(g_{\alpha,V}\) have the same primary and fourth-stage data. After matching their lower homotopy, their quotient belongs to \(\mathcal K_{4,V}\). Theorem~\ref{p1b:rt:theorem123} computes the full boundary image of this group as the identity. Therefore their full boundary components coincide.

The normalized restriction of \(g_{\eta,V}\) is in \(T[2]\), since it has zero free primary. Its square consequently has zero normalized restriction. Before normalization its circle coordinate is \(\kappa\) or \(3\kappa\); both square to \(2\kappa=\zeta\). The even local circle subgroup is central, so squaring introduces no circle-sphere commutator. Thus the whole boundary value of the square is precisely the local circle loop \(\zeta\), proving \eqref{p1b:eq:rt-28-1} and \eqref{p1b:eq:rt-28-5}. The same argument applies to every realization of the same lower data because their differences lie in \(\mathcal K_{4,V}\). 
\end{proof}

This selects the higher correction that remained open in Batch 11. That actual gauge supplies no additional finite translation after local circle normalization.

\paragraph{Actual gauges for the other fourth-stage directions}

Use the basis of Batch 27,
\[
H^4(W;\mathbb F_2)=\mathbb F_2\{\chi^4,\bar c_4,\alpha\}.
\]
Let \(P=S^3_x\times S^3_p\) be the antipodal-circle slice and \(Y\) its exchange mapping torus. The retained restriction maps are
\[
H^6(W;C_2)\xrightarrow{\cong}H^6(P;C_2),\quad
H^7(W;C_{12})\xrightarrow{\cong}H^7(Y;C_{12}),
\quad H^5(Y;C_2)=0.
\tag{RT.28.6}\label{p1b:eq:rt-28-6}
\]

\begin{lemma}
\label{p1b:rt:lemma126}
For every compatible actual \(V\), all three fourth-stage basis directions have actual determinant-one gauge realizations with primary zero. In particular there are gauges \(g_{\chi,V}\) and \(g_{c,V}\) with fourth classes \(\chi^4\) and \(\bar c_4\).

\end{lemma}
\begin{proof}
A fourth cocycle with the identity primary gives a partial section through the five-skeleton. Its degree-six obstruction restricts to zero on \(P\): \(P^{(5)}=S^3\vee S^3\) has no independent fourth or fifth difference, and the identity is an actual extension over \(P\). The first map in \eqref{p1b:eq:rt-28-6} therefore kills this obstruction globally.

For \(\chi^4\) and \(\alpha\), the fourth restriction to \(Y\) is zero, because \(\chi_Y^2=0\). Together with \(H^5(Y;C_2)=0\), the lower section can be matched to the actual identity through degree five on \(Y\). The identity and its product with the actual sixth-class gauge of Batch 26 realize both sixth choices on \(Y\). This is the same extension argument as Batch 11 and gives \(g_{\chi,V}\) as well as the already constructed \(g_{\alpha,V}\).

For \(\bar c_4\), construct an actual gauge on \(Y\) realizing its nonzero fourth restriction. Take a small base interval \(I\) away from the mapping-torus seam and a three-disc \(D^3_x\) in the first sphere. The product box \(I\times D^3_x\times S^3_p\) retracts to \(S^3_p\). Every complex rank-two bundle is trivial there, since \(\pi_3BU(2)=0\), so choose a frame of the actual \(V\) on that box.

Put a supported gauge \(\phi(s,x)\) in this frame, independent of \(p\), by the quotient
\[
(I/\partial I)\wedge(D^3_x/\partial D^3_x)=S^4
\xrightarrow{\eta_3}SU(2).
\]
It is the identity on the boundary of the box and extends by the identity to all of \(Y\), regardless of the exchange seam or the ordinary sector. The resulting actual gauge has zero primary and fourth class the mod-two Poincaré dual of \(\{s_0,x_0\}\times S^3_p\). In mapping-torus coordinates this is the coinvariant \([u_x]\), the generator of \(H^4(Y;C_2)\). Batch 11 identifies it with \(\bar c_4|_Y\).

The partial global section with fourth class \(\bar c_4\) can now be matched to this actual \(Y\)-gauge through degree five; the only next difference at that point is its sixth choice. Multiplying the \(Y\)-gauge by the actual sixth-class gauge realizes the other choice. Thus either restricted sixth choice has an actual extension on \(Y\), and the degree-seven global obstruction restricts to zero there. The second map in \eqref{p1b:eq:rt-28-6} kills that obstruction globally. The last obstruction group is \(H^8(W;C_2)=0\). This produces the actual global \(g_{c,V}\). Products of the three actual gauges realize every fourth class, since the first differing coefficient adds. 
\end{proof}

Let
\[
u_{\chi,V}=\operatorname {res}_{\mathrm{norm}}g_{\chi,V},\qquad
u_{c,V}=\operatorname {res}_{\mathrm{norm}}g_{c,V}.
\]
These are specified actual vectors in \(T[2]=\mathbb F_2\{a,e,t\}\), where \(a\) is the order-two value \(z_6=3\). Their values remain to be evaluated. Batch 27 makes each independent of all choices in degrees five and higher. Theorem~\ref{p1b:rt:theorem125} eliminates the \(\alpha\)-column. Hence the exact normalized image of zero-primary determinant-one gauges is
\[
\boxed{\operatorname {im}(\mathcal K_{3,V}\to T)
=\langle u_{\chi,V},u_{c,V}\rangle.}
\tag{RT.28.7}\label{p1b:eq:rt-28-7}
\]
Indeed subtract the product realizing any given fourth class; the remainder lies in \(\mathcal K_{4,V}\) and has identity boundary value.

\paragraph{The full remaining image is a four-column matrix}

Choose the actual gauges \(g_{h,V}\) and \(g_{\eta,V}\) of Batches 12–13 with joint low data \((h,0)\) and \((0,\eta)\). Put
\[
j_{h,V}=\operatorname {res}_{\mathrm{norm}}g_{h,V},\qquad
j_{\eta,V}=\operatorname {res}_{\mathrm{norm}}g_{\eta,V}.
\]
The first generator has odd global determinant winding, but its boundary winding is even and local. Both vectors therefore lie in \(T[2]\).

\begin{theorem}
\label{p1b:rt:theorem127}
For every compatible actual \(V\), the full zero-sphere-degree image and its determinant-one subimage are exactly
\[
\boxed{J_{0,V}=\operatorname {im}M_V,\qquad
M_V=(j_{h,V}\ \ j_{\eta,V}\ \ u_{\chi,V}\ \ u_{c,V})
\in\operatorname {Mat}_{3\times4}(\mathbb F_2),}
\tag{RT.28.8}\label{p1b:eq:rt-28-8}
\]
\[
\boxed{J_{0,V}^{\det=1}
=\langle j_{\eta,V},u_{\chi,V},u_{c,V}\rangle.}
\tag{RT.28.9}\label{p1b:eq:rt-28-9}
\]
The odd-determinant action is redundant precisely when
\[
\boxed{j_{h,V}\in
\langle j_{\eta,V},u_{\chi,V},u_{c,V}\rangle.}
\tag{RT.28.10}\label{p1b:eq:rt-28-10}
\]
These are equalities for actual gauge images. Their four column values are evaluated in the following complete positive chain, culminating in Theorem~\ref{p1b:thm:complete-mixed-rank-two}.

\end{theorem}
\begin{proof}
A gauge contributing to \(J_{0,V}\) has free primary coefficient zero, because its normalized sphere integer is zero. Its remaining primary is zero or \(\eta\), and its determinant is an integer multiple of \(h\). Remove the corresponding powers of the two actual low generators. For determinant winding zero, a global determinant logarithm deforms the remainder to determinant one. Its primary is now zero, so \eqref{p1b:eq:rt-28-7} supplies precisely the two remaining columns.

Even determinant winding also has the actual scalar gauge \(\sigma=fI\), where \(f\) represents \(h\). Its determinant is \(2h\), its primary is \(h\tau=0\) by the transfer projection formula, and its boundary value is local. Thus only the parity of the determinant winding supplies an additional finite direction. All four restriction columns lie in the elementary group \(T[2]\), so their powers reduce to binary coefficients. Each is an actual global restriction, proving both inclusions in \eqref{p1b:eq:rt-28-8}. For determinant-one gauges the \(h\)-column is omitted; products and the same primary-zero remainder give \eqref{p1b:eq:rt-28-9}.

Changing the chosen \(g_\eta\) adds an image from \eqref{p1b:eq:rt-28-7}, and changing a chosen pure \((h,0)\) gauge also changes its column by that image. Hence the displayed spans and membership are independent of those choices. The normalized odd-determinant operation is translation by \(j_{h,V}\) on the determinant-one quotient. It is the identity exactly when that vector is in \eqref{p1b:eq:rt-28-9}, proving \eqref{p1b:eq:rt-28-10}. This expresses determinant redundancy as actual three-vector span membership. The completed image table below decides every parity/sector case; the already evaluated split-bundle test is retained. 
\end{proof}


\subsubsection{Mixed rank-two classification and technical inputs}
\label{p1b:sec:rank-two-proof-interface}
The complete mixed rank-two obstruction, gauge-image, torsor, and
forgetful-fibre proof chains are given in Paper~I-D,
Sections~\ref{p1c:sec:rank-two-full-proofs} and
\ref{p1c:sec:natural-rank-two-full-proofs}.
All the original mathematical labels stay attached to their full
statements in that technical volume, which is the unique proof owner.
The final classification, exact totals, hypotheses, and scope firewall
remain stated below in Theorem~\ref{p1b:thm:complete-mixed-rank-two}.

\subsubsection{Complete mixed counts, quotient torsors and ordinary fibres}
\label{p1b:sec:complete-mixed-rank-two}
\begin{theorem}[Complete mixed rank-two classification on the declared domain]
\label{p1b:thm:complete-mixed-rank-two}
Retain the labelled product \(M_8=K^2\), Born exchange, fixed
eigenranks \((1,1)\), actual boundary
\(E_\partial=\mathbf1\oplus\lambda_L\), and the gluing domain
\(c_1=0,c_2=A\mathcal D\). There are two compatible free-region
rows in each sector. The groups acting on each row's full and auxiliary
gluing torsors are
\[
\begin{array}{c|c|c|c|c}
 &\text{full sector0}&\text{full sector1}&
 \text{auxiliary sector0}&\text{auxiliary sector1}\\\hline
A\text{ even}&C_{12}\oplus C_2&C_6\oplus C_2&C_{12}\oplus C_2&C_{12}\oplus C_2\\
A\text{ odd}&C_{12}\oplus C_2^2&C_6\oplus C_2&C_{12}\oplus C_2^2&C_{12}\oplus C_2
\end{array}
\]
The full totals are 72 even and 120 odd; the auxiliary totals are 96 even
and 144 odd. Auxiliary means determinant-one global gauges and the full
local mixed gauge group; it is not a principal-\(SU(2)\) lift category.

Put \(h_{12}=\gcd(12,A)\), \(h_6=\gcd(6,A)\), with
\(\gcd(m,0)=m\). An ordinary bundle outside the exact mixed-image
predicate of Theorem~\ref{p1b:thm:v2170-mixed-image} has no mixed lift.
Inside that image the individual cardinalities are
\[
\begin{array}{c|c|c|c}
\text{parity/sector}&\text{row}&\text{image size}&\text{lifts}\\\hline
A\equiv2\pmod4,0&\delta=0,1\text{ separately}&2h_{12}&12/h_{12}\\
4\mid A,0&\delta=0&2h_{12}&12/h_{12}\\
4\mid A,0&\delta=1&h_{12}&24/h_{12}\\
A\text{ even},1&\text{special}&2h_6&6/h_6\\
A\text{ even},1&\text{other}&h_6&12/h_6\\
A\text{ odd},0&\text{both rows combined}&h_{12}&96/h_{12}\\
A\text{ odd},1&\text{each disjoint row}&h_6&12/h_6
\end{array}
\]
The special even exotic row is intrinsically the one on which the
ordinary top insertion survives; at \(A=0\) it is \(\delta=0\).
The groups acting simply transitively on each nonempty fixed-row fibre
are
\[
\begin{array}{c|c|c}
\text{parity/sector}&\text{row}&\text{fibre group}\\\hline
A\equiv2\pmod4,0&\delta=0,1&C_{12/h_{12}}\\
4\mid A,0&\delta=0&C_{12/h_{12}}\\
4\mid A,0&\delta=1&C_{12/h_{12}}\oplus C_2\\
A\text{ even},1&\text{special}&C_{6/h_6}\\
A\text{ even},1&\text{other}&C_{6/h_6}\oplus C_2\\
A\text{ odd},0&\text{each row}&C_{12/h_{12}}\oplus C_2^2\\
A\text{ odd},1&\text{each row}&C_{6/h_6}\oplus C_2
\end{array}
\]
Here \(C_1\) is trivial. The odd sector-zero ordinary fibre is the
disjoint union of two row torsors over the same ordinary image; the
classification does not choose a group structure on that union.
The exotic rows have disjoint ordinary images. At Chern zero the row
image sizes \(24,12,12,6\) and multiplicities \(1,2,1,2\) yield 72
coherent classes over 54 ordinary classes, with 36 singleton fibres and 18
doubleton fibres. The natural extension classification is complete:
14 natural decisions are 4 YES/10 NO and 32 inherited decisions are 4
YES/28 NO, with no unevaluated entry. These decision totals count named
extension tests, not coherent bundle classes.
\end{theorem}
\begin{proof}
Theorem~\ref{p1b:rt:theorem35-images} gives all finite images and
counts. The even sphere presentation is
Theorem~\ref{p1b:rt:theorem31-sphere}; the odd coincident-circle
calculation gives \(h_e=w_e=0\) in
Theorem~\ref{p1b:rt:theorem40-e}, so
Theorem~\ref{p1b:rt:theorem41-type} selects the split exotic full type.
The exact homogeneous ordinary fibre action and incidence are proved in
Lemma~\ref{p1b:rt:lemma33-homogeneous} and
Theorem~\ref{p1b:rt:theorem35-fibres}. The fixed-row kernel groups are
Theorem~\ref{p1b:rt:theorem37-kernels}, specialized by
Corollary~\ref{p1b:rt:corollary41-fibres}. Dividing the row orders by
their proved ordinary image sizes gives the displayed cardinalities;
at Chern zero this is the evaluated specialization of
Theorem~\ref{p1b:thm:v2170-zero-fibres}.
The complete original natural comparison and scoped extension theorem
\ref{p1b:thm:natural-rank-two-completion} decide the final coefficients
and named tests. The expanded 14- and 32-row ledgers are preserved in
Paper~I-D, Section~\ref{p1b:app:rank-two-closure-ledgers}; the main compact test
table states their exact outcome and domain.
\end{proof}
The theorem constitutes a mathematical classification on the declared
reciprocal domain. 
Those separately typed applications retain their own input obligations.


\subsubsection{Spherical Poincar\'e class and integral pairings}
\begin{theorem}[Spherical Poincar\'e class and product index]
\label{p1b:thm:spherical-pairing}
On \(P_v=S^3_x\times S^3_p\), the virtual rank-one class
\[
 \cP_v=1+\xi_x\xi_p,\qquad \operatorname{ch}\cP_v=1+u_xu_p
\]
is invertible, with
\(B^*\cP_v=\cP_v^{-1}=\cP_v^*\).
For the oriented smash map \(q:S^3\times S^3\to S^6\), a normalized
rank-three Bott bundle \(E_6\) gives
\(\cP_v=[q^*E_6]-2[\mathbf1]\).
Rank three is minimal for an actual bundle carrying \(c_3=2u_xu_p\).
The K-oriented push-pull transform
\[
 \mathscr F_v(x)=(\pi_p)_!(\pi_x^*x\,\cP_v)
\]
is a parity-changing isomorphism, with
\(\mathscr F_v(1)=\xi_p,\ \mathscr F_v(\xi_x)=1\).
On \(M_8\), set \(\cP_8=\cP_{\rm sc}\cP_v\); then
\[
 \operatorname{ch}\cP_8=(1+\eta)(1+u_xu_p),\qquad
 \operatorname{ind}D^+_{\cP_8}=1.
\]
\end{theorem}
\begin{proof}
The square of \(\xi_x\xi_p\) is zero and exchange negates it.
Duality negates its degree-six Chern character. The oriented Bott
bundle has \(c_3(E_6)=2[S^6]\), proving the representative and minimal
rank. Integrate the two K-basis elements along the first sphere to obtain
the push-pull formulas; this is the spherical Fourier--Mukai construction
\cite{BouwknegtEvslinMathai}. Finally
\([\cP_{\rm sc}]=1+\xi_t\xi_E\); the product character has top coefficient
one. Its Chern classes are \(c_1=\eta,c_2=0,c_3=2u_xu_p,c_4=-4\mathcal O\),
consistent with Theorem~\ref{p1b:thm:chern}.
\end{proof}

\begin{proposition}[Flat bundles, lines and gerbes]
\label{p1b:prop:flat}
Flat unitary rank-\(r\) bundles with connection are classified by
\[
 \Hom(\Z^2,\U(r))/\U(r)\cong\Sym^r(\U(1)^2).
\]
They are topologically trivial. For a unitary local system \(V\) with
common holonomy-fixed subspace \(V_0\),
\[
 H^k(M_8;V)\cong H^k(M_8;\C)\otimes V_0.
\]
Topological line classes are \(\Z\eta\); gerbe classes are
\(\Z u_x\oplus\Z u_p\). Flat lines with connection form
\((\R/\Z)^2\), and flat gerbes with connection form \(\R/\Z\).
Exchange-even gerbe levels are \((k,k)\) and odd levels are \((k,-k)\).
\end{proposition}
\begin{proof}
Two commuting unitary holonomies diagonalize simultaneously. Their
eigenlines have torsion first Chern class, hence zero on this space.
A nontrivial rank-one circle holonomy has cellular cochain differential
\(\rho(1)-1\), which is invertible; its contraction contracts the tensor
product complex. Only trivial eigencharacters survive. Classifications
of lines and gerbes use \(H^2(-;\Z)\) and \(H^3(-;\Z)\);
flat refinements use \(H^1(-;\R/\Z)\) and \(H^2(-;\R/\Z)\).
Universal coefficients and torsion-free homology give the groups displayed.
\end{proof}

\part{Differential, metric and analytic geometry}
\section{Reciprocal forms, representations, and invariant operators}
\label{p1b:sec:organized-operators}

\subsection{Reciprocal transgression and the nonclosed two-form}
\label{p1b:sec:transgression}
\begin{theorem}[Normalized reciprocal transgression]
\label{p1b:thm:transgression}
On the vector factors set
\[
 X=g_x^{-1}\dd g_x,\quad Y=g_p^{-1}\dd g_p,\quad
 \nu(g)=-\frac1{24\pi^2}\Tr(g^{-1}\dd g)^3,\quad
 \beta=-\frac1{8\pi^2}\Tr(X\wedge Y).
\]
Choose the sphere orientations so \(\nu\) represents the primitive class.
For \(r_v=g_xg_p^{-1}\),
\[
 r_v^*\nu=\nu_x-\nu_p+\dd\beta,\qquad B^*\beta=-\beta.
\]
The global two-form
\[
 \Omega=\beta+\omega_{\rm sc}
\]
is nondegenerate of rank eight, is exchange-odd, and obeys
\[
 \dd\Omega=r_v^*\nu-\nu_x+\nu_p\ne0.
\]
\end{theorem}
\begin{proof}
\(r_v^{-1}\dd r_v=\Ad_{g_p}(X-Y)\). Graded trace cyclicity gives
\[
 \Tr(X-Y)^3=\Tr X^3-\Tr Y^3-3\Tr(X^2Y)+3\Tr(XY^2).
\]
Maurer--Cartan gives
\(\dd\Tr(X\wedge Y)=-\Tr(X^2Y)+\Tr(XY^2)\); substitute the constants.
The fundamental trace form pairs the two three-dimensional tangent
factors nondegenerately. The scalar two-form pairs the two remaining
directions, so \(\Omega^4=4\beta^3\omega_{\rm sc}\) never vanishes.
Exchange reverses both summands. Two vector arguments in one factor
and one in the other give a nonzero bracket contraction in \(\dd\beta\).
The explicit primitive uses common-model trace data; the differential
equation alone allows addition of a closed two-form, with cohomological
ambiguity \(\R\eta\).
\end{proof}

\begin{theorem}[Exact twisted Dirac graph]
\label{p1b:thm:dirac-graph}
Set \(H=-\dd\Omega\) and define
\[
 [X+\xi,Y+\zeta]_H=[X,Y]+\mathcal L_X\zeta-\iota_Y\dd\xi
                         +\iota_Y\iota_XH.
\]
Then
\[
 L_\Omega=\{X+\iota_X\Omega:X\in TM_8\}
\]
is maximal isotropic and closed under this bracket. The transformation
\(e^{-\Omega}(X+\xi)=X+\xi-\iota_X\Omega\) identifies it with the
untwisted tangent Dirac structure. Its Lie-algebroid cohomology is
de Rham cohomology.
\end{theorem}
\begin{proof}
Skew symmetry gives isotropy and the graph has half the total rank.
Cartan's formula gives
\[
 [X+\iota_X\Omega,Y+\iota_Y\Omega]_H
 =[X,Y]+\iota_{[X,Y]}\Omega+\iota_Y\iota_X(\dd\Omega+H).
\]
The last term vanishes and \(\dd H=0\).
The displayed transform removes both the graph form and twist;
anchor projection identifies the restricted algebroid with \(TM_8\).
These are the standard twisted-Dirac conventions \cite{SeveraWeinstein2001TwistedPoisson}.
\end{proof}
The form is invariant under independent left translations and simultaneous
right translations, not independent right translations.
The obstruction to a closed symplectic form and the stronger
left-invariant locally conformally symplectic obstruction are proved
in Paper~I-D, Section~\ref{p1b:app:obstructions}.


\subsubsection{The compact double with group-valued moment map}
\label{p1b:sec:g12}

Fix a positive \(\Ad\)-invariant inner product
\(\langle\, ,\,\rangle\) on \(\mathfrak k\). Pairings of
Lie-algebra-valued forms include the exterior product. Write
\(\theta^L=g^{-1}\dd g\), \(\theta^R=\dd g\,g^{-1}\), and
\(\chi=\langle\theta^L,[\theta^L,\theta^L]\rangle/12\);
the bracket of one-forms includes both alternating terms.

\begin{theorem}[Fused compact double and its ranks]
\label{p1b:thm:g12-double}
For diagonal conjugation on \(M_8=K\times K\), put
\[
 \mu(a,b)=aba^{-1}b^{-1},
\]
\[
 \omega_D=\frac12\langle a^*\theta^L,b^*\theta^R\rangle
 +\frac12\langle a^*\theta^R,b^*\theta^L\rangle
 +\frac12\langle(ab)^*\theta^L,(a^{-1}b^{-1})^*\theta^R\rangle .
\]
For \(\xi^\#=\left.\frac{\dd}{\dd s}\right|_0
  (e^{s\xi}ae^{-s\xi},e^{s\xi}be^{-s\xi})\), one has
\[
 \dd\omega_D=-\mu^*\chi,\qquad
 \iota_{\xi^\#}\omega_D
     =\tfrac12\mu^*\langle\theta^L+\theta^R,\xi\rangle,\qquad
 \ker(\omega_D)_{(a,b)}
     =\{\xi^\#_{(a,b)}:(\Ad_{\mu(a,b)}+I)\xi=0\}.
\]
The form has rank eight except where the fundamental \(SU(2)\) trace
of \(\mu\) is zero, where it has rank six. Exchange obeys
\[
 \mu\circ B=\mu^{-1},\qquad B^*\omega_D=-\omega_D.
\]
\end{theorem}
\begin{proof}
The double and internal fusion constructions of
\cite[Proposition 3.2 and Theorem 6.1]{AlekseevMalkinMeinrenken1998} give the three
identities with precisely these forms and action. On \(\mathfrak{su}(2)\),
\(\Ad_\mu\) is a rotation through twice the class angle of \(\mu\).
It has a two-dimensional \(-1\)-eigenspace exactly at trace zero;
the central Lie-algebra direction never contributes. At trace zero the
pair cannot commute. Two noncommuting elements of \(SU(2)\) have zero
common infinitesimal centralizer, so these two fundamental vectors
are independent. Otherwise the kernel is zero. The exchange identity
follows by exchanging the first two pairings and using inversion on
the two product maps in the last pairing; Paper~I-D, Section~\ref{p1b:app:g12}
records the signs explicitly.
\end{proof}

\begin{proposition}[Commuting-pair reduction and its strata]
\label{p1b:prop:g12-reduction}
The identity-level quotient is
\[
 \mu^{-1}(1)/K\cong(\T^2/\{\pm1\})\times\T^2.
\]
Its four singular tori correspond to the four central \(SU(2)\) pairs.
On the maximal-torus cover of the smooth stratum the reduced form is
\[
 c_v\,\dd\alpha\wedge\dd\beta+
 c_0\,\dd\theta_t\wedge\dd\theta_E,
 \qquad c_v=\langle H_v,H_v\rangle,\quad
 c_0=\langle H_0,H_0\rangle,
\]
where \(H_v=i\sigma_3\) and \(H_0=i\) generate circles of period
\(2\pi\). On each singular torus the stratum form is the second
summand. The central value \((-I,1)\) instead has quotient \(\T^2\),
with that same scalar form.
\end{proposition}
\begin{proof}
Commuting \(SU(2)\) pairs lie in a common maximal torus, uniquely up to
simultaneous Weyl inversion. The two scalar circles are central and
unaffected by conjugation. Restricting \(\omega_D\) to this torus gives
the displayed form: the fusion term pairs a one-form with its negative
and is zero. At \(\mu_v=-I\), the unit quaternions \(a_v,b_v\) anticommute.
They are orthogonal purely imaginary unit quaternions, a single
conjugacy orbit. Only the scalar torus remains.
\end{proof}
The identity reduction is singular; its dimension does not follow from
a free-action regular-value formula. This construction uses the
commutator \(\mu\). Its form \(\omega_D\) and the relative-comparison
form \(\Omega\) of Theorem~\ref{p1b:thm:transgression} have different
differentials and ranks.


\subsection{Representations and invariant differential operators}
\label{p1b:sec:operators}
\begin{theorem}[Irreducibles and fusion]
\label{p1b:thm:irreducibles}
Every complex irreducible unitary representation is
\[
 V_{j_x}\otimes V_{j_p}\otimes z_t^n\otimes z_E^w,\qquad
 j_x,j_p\in\tfrac12\Z_{\ge0},\quad n,w\in\Z.
\]
Its dimension is \((2j_x+1)(2j_p+1)\); exchange sends its labels to
\((j_p,j_x,w,n)\). The sphere central signs are
\((-1)^{2j_x},(-1)^{2j_p}\).
Tensor products add scalar weights and apply, in each sphere,
\[
 V_j\otimes V_l=\bigoplus_{s=|j-l|}^{j+l}V_s
\]
with unit steps. Thus
\[
 R(M_8)=\Z[\chi_x,\chi_p,z_t^{\pm1},z_E^{\pm1}].
\]
\end{theorem}
\begin{proof}
Irreducibles of a compact product are tensor products; apply the
\(\SU(2)\) highest-weight classification and the circle character lattice.
The Clebsch--Gordan recurrence generates every \(\SU(2)\) character from
the fundamental one, and its highest weight proves independence.
A separately specified quotient \((\SU(2)\times\U(1))/\langle(-1,-1)\rangle\)
would require \(2j+n\) even for descent. This extra quotient is not part
of the product definition.
\end{proof}

\begin{theorem}[Complete bi-invariant scalar operator algebra]
\label{p1b:thm:bi-operators}
Let \(C_x,C_p\) be the nonnegative quadratic sphere Casimirs, normalized
as the unit-sphere scalar Laplacians, and let \(N_t,N_E\) be as above.
Finite-order scalar differential operators invariant under left and right
translations form
\[
 \mathcal D_{\rm bi}=\C[C_x,C_p,N_t,N_E].
\]
Set \(C_+=C_x+C_p,\ c=C_x-C_p,\ N_+=N_t+N_E,\ n=N_t-N_E\).
Then
\[
 \mathcal D_{\rm bi}^{B+}
 \cong\C[C_+,N_+,U,V,W]/(V^2-UW),\quad
 U=c^2,\ V=cn,\ W=n^2.
\]
The exact-order Hilbert series is
\[
 \frac{1-z^6}{(1-z)(1-z^2)^2(1-z^3)(1-z^4)},
\]
with coefficients \(1,1,3,4,8,10,16\) through order six.
\end{theorem}
\begin{proof}
Bi-invariant differential operators identify with the center of the
complex universal enveloping algebra. Each simple factor contributes its
quadratic Casimir; the center contributes two commuting first-order
generators. Exchange fixes the sums and negates both differences.
An even monomial in \(c,n\) has both exponents even or both odd. It reduces
uniquely to \(f(U,W)+Vg(U,W)\), proving the generators and sole relation.
The differential orders are \(2,1,4,3,2\), and the relation has order six.
The Hilbert series follows by the one homogeneous relation.
\end{proof}

\begin{corollary}[Second-order family and positivity of its symbol]
\label{p1b:cor:second-order}
A formally self-adjoint member of order at most two is
\[
 \alpha C_++\beta N_+^2+\gamma(N_t-N_E)^2+\delta N_++\epsilon,
 \qquad\alpha,\beta,\gamma,\delta,\epsilon\in\R.
\]
Its second-order principal symbol is positive definite exactly when
\(\alpha,\beta,\gamma>0\). On the relative group, imposing all left/right
translations and inversion instead gives
\(\C[C_{\rm rel},N_{\rm rel}^2]\).
For an inversion-even smooth function \(F\) on \(K\), every odd
exponential-coordinate jet at the identity vanishes.
\end{corollary}
\begin{proof}
List monomials of weighted degree at most two in the previous normal form.
The independent sphere, scalar-sum and scalar-difference cotangent
directions give the positivity test. On \(K\), inversion fixes the Casimir
and negates the circle generator. Finally
\(F(\exp X)=F(\exp(-X))\); differentiating at zero proves every odd-jet
identity. Positive principal symbol alone does not imply positivity
of the full operator with arbitrary lower-order terms.
\end{proof}

\begin{theorem}[Matched Hopf-marked operator algebra]
\label{p1b:thm:hopf-operators}
Require invariance under the transitive translation side and the chosen
Hopf circle on the other side of each sphere. Let \(Q_x,Q_p\) be their
integer charge generators. Then
\[
 \mathcal D_H=\C[C_x,C_p,Q_x,Q_p,N_t,N_E].
\]
If the markings satisfy \(BQ_xB^{-1}=Q_p\), write \(s_i,d_i\) for
the sums and differences of \((C,Q,N)\). The exchange-even subalgebra is
\[
 \C[s_1,s_2,s_3,d_id_j:1\le i\le j\le3],
\]
whose relations are the two-by-two minors of \(Z=(d_id_j)\).
Its exact-order Hilbert series is
\[
 \frac12\left\{\frac1{(1-z^2)^2(1-z)^4}
              +\frac1{(1-z^4)(1-z^2)^2}\right\},
\]
with coefficients \(1,2,7,14,31,54,97\) through order six.
\end{theorem}
\begin{proof}
A weight-zero PBW monomial in \(U(\mathfrak{sl}_2)\) reduces to powers of
the Cartan generator and \(EF\); the latter is a polynomial in the Cartan
generator and Casimir. This proves the single-sphere algebra.
Every even monomial in the three differences is a product of pairs.
The two-by-two minors relate any two pairings of the same monomial,
giving the second Veronese algebra. Averaging the identity and exchange
traces on the weighted polynomial algebra gives the Hilbert series.
The matching condition is on the generators themselves, not just their
spectra.
\end{proof}

\begin{theorem}[Bi-Hopf charge bundles and spectral floors]
\label{p1b:thm:charge-bundles}
The character \((q_x,q_p)\in\Z^2\) of the bi-Hopf fibre gives the line
\[
 L_{q_x,q_p}=L_x^{q_x}\boxtimes L_p^{q_p},\qquad
 c_1=\varepsilon_xq_xh_x+\varepsilon_pq_ph_p.
\]
Every topological line on \(S^2\times S^2\) occurs.
For round radii \(R_x,R_p\), its equivariant scalar harmonics have
\[
 \lambda_{k_x,k_p}
 =\frac{k_x(k_x+2)}{R_x^2}+\frac{k_p(k_p+2)}{R_p^2},\quad
 k_i\ge|q_i|,\quad k_i\equiv q_i\pmod2,
\]
with labelled multiplicity \((k_x+1)(k_p+1)\).
The lowest eigenvalue is obtained at \(k_i=|q_i|\).
For the horizontal sub-Laplacian subtract
\(q_x^2/R_x^2+q_p^2/R_p^2\); its floor is
\(2|q_x|/R_x^2+2|q_p|/R_p^2\).
\end{theorem}
\begin{proof}
Associated-line Chern classes add. The base has \(H^2=\Z^2\), so the
unit Chern matrix gives all lines. On a sphere, the degree-\(k\)
scalar carrier is \((k/2,k/2)\); one Hopf weight takes the values
\(-k,-k+2,\ldots,k\), each with multiplicity \(k+1\).
Take the tensor product and minimize. Distinct joint labels with the
same total eigenvalue contribute additive multiplicities.
Hopf charges \(q_x,q_p\) and independent scalar weights \(n,w\)
remain separate unless an equivariance constraint is explicitly imposed.
\end{proof}

\subsection{Quadratic and Clifford descendants}
\label{p1b:sec:clifford}
\begin{definition}[Four-generator datum]
Let \(V\) be an oriented Euclidean four-space. It may be the ambient
coordinate module of a sphere or a metrized relative tangent module,
with the identification stated in each application.
Write \(V_\C\cong(\tfrac12,\tfrac12)\) under \(\Spin(4)\).
\end{definition}
\begin{theorem}[Complete quadratic split and Clifford quotient]
\label{p1b:thm:quadratic}
There is an equivariant decomposition
\[
 V\otimes V=\R g\oplus S_0^2V\oplus\Lambda^2V,\qquad16=1+9+6.
\]
Degree-two Clifford multiplication \(c_2(v\otimes w)=vw\) gives
\[
 0\longrightarrow S_0^2V\longrightarrow V\otimes V
 \xrightarrow{c_2}\R\oplus\Lambda^2V\longrightarrow0.
\]
The full tensor is recovered as
\[
 V\otimes V\cong\Sym^2V\times_\R(\R\oplus\Lambda^2V),
\]
where the maps to \(\R\) are contraction and scalar projection.
Moreover
\[
 S_0^2V_\C\cong E_2(S^3),\qquad
 \Sym^2E_1(S^3)\cong E_0(S^3)\oplus E_2(S^3).
\]
\end{theorem}
\begin{proof}
The Clifford relation is \(vw=g(v,w)+v\wedge w\).
It kills exactly the trace-free symmetric tensor and keeps the scalar
contraction and antisymmetric component. Retaining the symmetric tensor
and matching its contraction reconstructs every component.
The polynomial \(Q_{ab}x_ax_b\) is harmonic in \(\R^4\) exactly when
\(\tr Q=0\); restriction gives the nine-dimensional \((1,1)\) module.
The trace restricts to a constant on the sphere, proving the last identity.
\end{proof}

\begin{theorem}[Protected Clifford multiplication and transvections]
\label{p1b:thm:protected}
Let \(C_1,\ldots,C_4\) be irreducible Hermitian complex Clifford matrices
with \(\{C_a,C_b\}=2\delta_{ab}\). Set \(e_i=C_iC_4\).
Then
\[
 e_i^2=-1,\quad \{e_i,e_j\}=-2\delta_{ij},\quad
 [e_i,C_j]=-2\delta_{ij}C_4,\quad [e_i,C_4]=2C_i.
\]
If a rotation-vector triple \(P_i\) has the latter commutators multiplied
by \(i\lambda/2\), then \(P_i=(i\lambda/2)e_i\).
For \(i\ne j\),
\[
 e_ie_je_i^{-1}e_j^{-1}=-1,\qquad [e_i,e_j]=2e_ie_j.
\]
With \(\omega=e_1e_2e_3\), the oriented dual of \(e_ie_j\) is
\(-e_ie_j\omega=e_k\) for cyclic \((i,j,k)\).
\end{theorem}
\begin{proof}
Apply the Clifford relations. The difference between \(P_i\) and the
displayed candidate commutes with all four \(C_a\); irreducibility makes
it scalar. A scalar triple transforming as a vector under rotations must
vanish. Since \(e_i^{-1}=-e_i\), the finite commutator is \(-1\).
Multiplication by \(\omega\) gives the final duality.
\end{proof}

\begin{proposition}[An explicit oriented seven-dimensional identification]
\label{p1b:prop:septet}
In \(\Cl(4,0)^{\rm even}\), let \(I_V=C_1C_2C_3C_4\).
Multiplication by \(I_V\) is an isometry in the coefficient trace metric
from \(\R\oplus\Lambda^2V\) onto the scalar-free part of
\(\Cl(4,0)^{\rm even}\). Under the algebra isomorphism
\(\Cl(0,3)\to\Cl(4,0)^{\rm even}\), \(e_i\mapsto C_iC_4\),
this is a concrete map to the seven-dimensional complement of the unit.
\end{proposition}
\begin{proof}
\(I_V^2=1\); it sends degree zero to degree four and degree two to
degree two by Hodge duality. These have dimensions one and six.
Products of the three \(C_iC_4\) form an eight-element basis of the
even algebra with the \(\Cl(0,3)\) multiplication law.
Thus the map uses the orientation and Clifford identification,
rather than equality of dimensions.
\end{proof}

\begin{theorem}[Tetrahedral affine incidence is exactly \(\mathcal H_0\oplus\mathcal H_2\)]
\label{p1b:thm:v285-tetrahedral-affine-incidence}
Let \(u_1,\ldots,u_4\in\operatorname{Im}\mathbb H\) be an oriented unit
tetrahedral frame,
\[
 \sum_{j=1}^4u_j=0,
 \qquad
 \sum_{j=1}^4u_j\otimes u_j=\frac43I,
 \qquad
 \langle u_i,u_j\rangle=-\frac13\quad(i\ne j).
\]
For \(q\in SU(2)\) define
\begin{equation}
\label{p1b:eq:v285-tetrahedral-affine-incidence}
 \boxed{
 b_{ij}(q)=\frac14+\frac34
 \left\langle u_j,\operatorname{Ad}_q u_i\right\rangle.}
\end{equation}
Then
\[
 \sum_{j=1}^4b_{ij}(q)=1,
 \qquad
 b_{ij}(1)=\delta_{ij},
\]
and the complete real function span is
\[
 \boxed{
 \operatorname{span}_{\mathbb R}\{b_{ij}:1\le i,j\le4\}
 =\mathcal H_0(S^3)\oplus\mathcal H_2(S^3).}
\]
Thus local finite tetrahedral incidence is an affine matrix coefficient of
the quaternionic Clifford adjoint action.
\end{theorem}

\begin{proof}
The two affine identities follow from \(\sum_j u_j=0\) and the tetrahedral
Gram matrix.  In quaternion coordinates, \(\operatorname{Ad}_q\) is quadratic
in the four coordinates of \(q\), so every nonconstant term in
\eqref{p1b:eq:v285-tetrahedral-affine-incidence} lies in the degree-two
sphere-harmonic sector.  Conversely the tensors \(u_j\otimes u_i\) span
\(\operatorname{End}(\operatorname{Im}\mathbb H)\) because the tetrahedral
vectors span the three-space.  Hence their adjoint matrix coefficients span
the full nine-dimensional spin-one matrix-coefficient space, which is
\(\mathcal H_2(S^3)\).  Adding the constant term gives the direct sum with
\(\mathcal H_0\).
\end{proof}

\begin{remark}[Incidence is not probability]
The coefficients \(b_{ij}(q)\) are exact affine/barycentric incidence
coordinates.  Without a separately constituted positive state or stochastic
owner, they are not probabilities, transition rates, or occupation numbers.
\end{remark}

\subsection{Finite restrictions, harmonics and exact sampling}
\label{p1b:sec:finite}
\begin{proposition}[Factor restrictions and exceptional derived series]
\label{p1b:prop:finite-defect}
On the inclusions \(a\mapsto(a,e)\), \(b\mapsto(e,b)\),
\[
 (\cA,\cE)=\tfrac12(\theta,\theta),\qquad
 (\cA,\cE)=\tfrac12(\theta,-\theta),
\]
respectively. On any subgroup \(\Gamma\subset\SU(2)\), the subgroup
generated by the values of the restricted defect is \([\Gamma,\Gamma]\).
For the standard binary polyhedral groups,
\[
 [2T,2T]=Q_8,\quad [Q_8,Q_8]=\{\pm1\},\quad
 [2O,2O]=2T,\quad [2I,2I]=2I.
\]
Their abelianizations are \(C_3,C_2,1\), respectively.
\end{proposition}
\begin{proof}
Pull back the Cartan forms. The defect is the group commutator by
definition. For \(2T=Q_8\rtimes C_3\), the cyclic action permutes the
three quaternion axes, and commutators generate \(Q_8\).
The projective quotients \(2T/\{\pm1\}=A_4\),
\(2O/\{\pm1\}=S_4\) give the octahedral derived subgroup;
the quaternion commutator \(-1\) supplies the central lift.
The projective icosahedral quotient \(A_5\) is perfect, and its binary
lift contains the same central commutator, hence is perfect.
\end{proof}

\begin{theorem}[Reciprocal harmonic transport]
\label{p1b:thm:harmonic-transport}
Let \(\iota:S_p^3\to S_x^3\) be the declared comparison isometry, with
transported normalized measure. Then \(Jf=f\circ\iota^{-1}\) is unitary,
\(\Delta_xJ=J\Delta_p\), and \(JE_k^p=E_k^x\).
For inverse transported actions \(T^\vee=\iota T^{-1}\iota^{-1}\),
\[
 U_x(T)J=JU_p(T^{-1}),\qquad \Pi_\Gamma^xJ=J\Pi_\Gamma^p,
 \quad \Pi_\Gamma=\frac1{|\Gamma|}\sum_{\gamma\in\Gamma}U(\gamma).
\]
Complete harmonic-band projections commute with \(\Pi_\Gamma\).
For a one-sided \(2I\)-action, write
\[
 \frac{1-t^{60}}{(1-t^{12})(1-t^{20})(1-t^{30})}
 =\sum_{k\ge0}\mu_kt^k.
\]
Then
\[
 \dim(\cP_m\cap\im\Pi_{2I})=\sum_{k=0}^m(k+1)\mu_k,\qquad
 \cP_m=\bigoplus_{k=0}^mE_k.
\]
The first nonconstant invariant subspace occurs at \(k=12\), of dimension
13; the complete invariant band there has dimension 14.
\end{theorem}
\begin{proof}
Laplacian naturality and change of variables prove unitarity and
intertwining. Substitute the transported action, then average;
inversion permutes the finite group. Every isometry commutes with the
Laplacian and its spectral projections. The one-sphere Molien series
in Paper I-B gives the invariant multiplicity in \(V_{k/2}\);
the other Peter--Weyl factor contributes \(k+1\).
The first nonconstant coefficient of that series is \(\mu_{12}=1\).
\end{proof}

\begin{theorem}[Exceptional quadrature and sharp complete bands]
\label{p1b:thm:design}
The unit-quaternion vertex sets \(2T,2O,2I\) are uniform spherical
designs of exact strengths \(5,7,11\). For a product of \(r\) profiles
in the complete scalar band \(\cP_m\), evaluation is exact when
\(rm\le t_\Gamma\). For \(r=2,3,4\) the maximal complete bands are
\[
\begin{array}{c|ccc}
 &2T&2O&2I\\\hline
 r=2&2&3&5\\
 r=3&1&2&3\\
 r=4&1&1&2
\end{array}
\]
In particular \(\cP_2\) has dimension \(1+4+9=14\), and \(2T,2O,2I\)
give the same maximal complete band for orders two, three and four,
respectively. Normalized evaluation
\[
 S_{\Gamma,m}f=(f(q)/\sqrt{|\Gamma|})_{q\in\Gamma}
\]
is an isometry whenever \(2m\le t_\Gamma\).
\end{theorem}
\begin{proof}
The round even moments are
\[
 \int_{S^3}\prod_{i=0}^3x_i^{2a_i}\,\dd\sigma
 =\frac{\prod_i(2a_i)!}
 {4^{\sum a_i}(\sum a_i+1)!\prod_i a_i!}.
\]
The explicit finite coordinate sets of Paper I-B reproduce them through
degrees \(5,7,11\); an exact algebraic verification and sharpness
witnesses are included with this paper.
Products have degree at most \(rm\). Sharp witnesses in
Paper~I-D, Section~\ref{p1b:app:finite-checks} fail in the first band above
every listed frontier. The isometry is the case of the product
\(\overline f g\). Dirichlet pairings on \(\cP_2\) are also exact under
\(2T\): in ambient polynomials,
\(\nabla_Sf\cdot\nabla_Sg=\nabla f\cdot\nabla g-
(x\cdot\nabla f)(x\cdot\nabla g)\) has degree at most four.
\end{proof}


\subsubsection{Weighted subgroup sampling and its complete domain}

\begin{lemma}[Products and weighted evaluation of complete scalar bands]
\label{p1b:lem:v233-weighted-band-products}
With probability Haar measure on the unit sphere, put
\[
 \cP_m=\bigoplus_{k=0}^m E_k,\qquad
 d_m=\frac{(m+1)(m+2)(2m+3)}6.
\]
This is the space of restricted polynomials of degree at most \(m\).
Products of \(r\) members span \(\cP_{rm}\); in particular
\(\operatorname{span}\{\overline f g:f,g\in\cP_m\}=\cP_{2m}\).
For nonnegative normalized weights on a finite set \(Z\),
\[
 S_wf=(\sqrt{w_z}f(z))_{z\in Z}
\]
is an isometry on \(\cP_m\) if and only if the rule is exact through
degree \(2m\). Its number of positive weights is at least \(d_m\).
Exact fourth powers of every real \(f\in\cP_m\) instead require
degree \(4m\) cubature.
\end{lemma}
\begin{proof}
Harmonic decomposition restricts every polynomial of degree at most \(m\)
to the displayed band. Split a coordinate monomial of degree at most
\(rm\) into \(r\) factors of degree at most \(m\); this proves the span
claim. The isometry identity is precisely the quadrature identity on
\(\overline f g\), and injectivity gives the rank bound. Real polarization
recovers every four-factor product from fourth powers.
\end{proof}

\begin{theorem}[Complete weighted finite-subgroup sampler classification]
\label{p1b:thm:v233-subgroup-sampler-classification}
The finite subgroups \(\Gamma\subset SU(2)\) admitting nonnegative
normalized weights exact for all \(\cP_m\) inner products are, up to
conjugacy, exactly
\[
\begin{array}{c|l|l}
 m&\text{admissible groups}&\text{order by subgroup inclusion}\\ \hline
 0&\text{all finite groups}&\text{trivial group least}\\
 1&\text{all noncyclic finite groups}&\text{no least class}\\
 2&2T,2O,2I&2T\text{ uniquely least}\\
 3&2O,2I&\text{incomparable}\\
 4,5&2I&\text{unique class}\\
 m\ge6&\varnothing&\text{empty domain}.
\end{array}
\]
Whenever weighted exactness is possible, uniform weights are also exact.
Consequently no independent uniformity or antipodality hypothesis is
needed for the existence classification at \(m=2\).
\end{theorem}
\begin{proof}
Left translations by \(\Gamma\) preserve Haar measure and polynomial
degree. Averaging an exact measure \(\sum w_\gamma\delta_\gamma\) over
those translations produces the uniform measure, still exact. This proves
the weight reduction, without identifying a specified instrument with its
average.

A cyclic subgroup lies in a real two-plane. A nonzero linear form
vanishing on that plane has zero sampled norm and positive Haar norm,
so cyclic groups fail for \(m\ge1\). For \(n\ge2\), the binary dihedral
group of order \(4n\) is
\[
 2D_n=\{e^{i\pi k/n},j e^{i\pi k/n}:0\le k<2n\}.
\]
These are two regular \(2n\)-gons in orthogonal planes. Their equally
weighted union has second moments \(\delta_{ab}/4\); antipodality kills
odd degrees through three. It is a 3-design.

Every cyclic or binary dihedral group preserves an unoriented conjugation
axis \(\mathbb Ru\). The degree-two harmonic
\[
 h_u(q)=\langle\operatorname{Ad}_q u,u\rangle
\]
has \(h_u^2=1\) on the group for every normalized choice of weights,
whereas \(\int h_u^2\,d\sigma=1/3\), by uniform pushforward to \(S^2\).
For \(u=i\) it is \(q_0^2+q_1^2-q_2^2-q_3^2\).
Thus these groups fail for \(m\ge2\).

The finite-subgroup classification leaves the exceptional groups.
Theorem~\ref{p1b:thm:design} gives their strengths \(5,7,11\).
Sharp failures, computed from their quaternion coordinates and the
Haar moment formula in that theorem, are
\[
\begin{array}{c|c|c|c}
 \Gamma&f&|\Gamma|^{-1}\sum_\Gamma f&\int f\,d\sigma\\ \hline
 2T&q_0^6&3/32&5/64\\
 2O&q_0^8&15/256&7/128\\
 2I&q_0^{12}&133/4096&33/1024 .
\end{array}
\]
These are squared norms from \(\cP_3,\cP_4,\cP_6\), respectively.
Weight averaging precludes repairing their failure with unequal weights.
This proves the table's entire range.

The inclusions \(2T\subset2O,2I\) give the \(m=2\) least class;
\(48\nmid120\) gives incomparability at \(m=3\).
A putative least noncyclic group at \(m=1\) would embed in both
\(Q_8\) and \(2D_3\), hence have order dividing four and be cyclic,
a contradiction. The \(m=2\) groups all contain \(-1\), proving the
last assertion.
\end{proof}

\begin{corollary}[Quadratic comparison and its exact support boundary]
\label{p1b:cor:v233-quadratic-comparison-boundary}
On the category of finite group-closed exact \(\cP_2\) samplers,
\([2T]\) is the unique least class. On nonexceptional subgroup supports,
the weighted sampling Gram error satisfies
\(\|S_w^*S_w-I\|\ge2\).
If the same role requires fourth powers of \(\cP_2\) fields, the only
admissible finite subgroup is \(2I\).
No finite-subgroup sampler serves all growing complete bands.

Group closure is essential: there is a positive, uniform, antipodal
72-point 5-design containing \(1,-1\) and containing no orthogonal
triple, hence no rotated Hurwitz cell.
Subsection~\ref{p1b:app:v233-support-countermodels} proves this exact
counterexample. Exactness and an identity anchor alone therefore do not
constitute a group-closed exceptional support.
\end{corollary}
\begin{proof}
Apply the theorem and Lemma~\ref{p1b:lem:v233-weighted-band-products}.
The normalized test \(\sqrt3 h_u\) has squared sampled norm three
and Haar norm one, so its Gram discrepancy is two.
The counterexample and its scope are proved in the cited subsection.
\end{proof}


\begin{theorem}[Cosetwise Parseval reconstruction and exact resolution transport]
\label{p1b:thm:v273-cosetwise-resolution-groupoid}
Let $G=SU(2)=S^3$ have normalized Haar measure and
$\mathcal P_m=\bigoplus_{k=0}^m E_k$.
For any finite subgroup $\Gamma\le G$ of design strength
$t_\Gamma\ge2m$ and any coset $A=g\Gamma$, define
$(S_Af)(a)=|\Gamma|^{-1/2}f(a)$ and let
$K_m(x,y)$ denote the Haar reproducing kernel of $\mathcal P_m$.
Then $S_A^*S_A=I$, and for all $f\in\mathcal P_m$,
\[
 f(x)=\frac1{|\Gamma|}\sum_{a\in A}K_m(x,a)f(a).
\]
For any two such frames $A,B$, the rectangular transition
$U_{B\leftarrow A}:=S_BS_A^*$ has entries
\[
 U_{B\leftarrow A}(b,a)=\frac{K_m(b,a)}{\sqrt{|A||B|}}
\]
and satisfies, with $P_A=S_AS_A^*$,
\[
 U_{B\leftarrow A}^*U_{B\leftarrow A}=P_A,\quad
 U_{B\leftarrow A}U_{B\leftarrow A}^*=P_B,\quad
 U_{C\leftarrow B}U_{B\leftarrow A}=U_{C\leftarrow A}.
\]
At $m=2$, $K_2(x,y)=12(x\cdot y)^2+4(x\cdot y)-2$,
and the $2T,2O,2I$ sample images have identical rank $14$.
Their redundant sample dimensions are $10,34,106$, respectively.
These transitions are partial isometries between \emph{scalar
comparison frames}, not necessarily physical unitaries on $\mathcal H_1$.
\end{theorem}
\begin{proof}
Left multiplication preserves the restricted polynomial degree and
normalized Haar measure. Degree-$2m$ quadrature therefore proves
$\langle S_Af,S_Ah\rangle=\langle f,h\rangle$ for all
$f,h\in\mathcal P_m$ and every coset. Apply $S_A^*S_A=I$ to the
reproducing vector to obtain reconstruction. The matrix entries
follow from $S_A^*e_a=|A|^{-1/2}K_m(\cdot,a)$.
All transition identities are consequences of $S_A^*S_A=I$.
The scalar $S^3$ harmonic dimension is
$\sum_{k=0}^m(k+1)^2$; at $m=2$ this equals $14$.
The stated polynomial kernel follows from the addition formula
$\sum_{k=0}^2(k+1)U_k(x\cdot y)$.
\end{proof}

\begin{proposition}[Degree-twelve discrimination of twin icosahedral cosets]
\label{p1b:prop:v273-twin-degree12-vandermonde}
Let $H\subset I_\pm$ be the two completions of
Theorem~\ref{p1:thm:v273-two-icosahedral-completions}.
The uniform functionals $\tau_\pm$ agree on $\mathcal P_{11}$.
For a standard Hurwitz-axis quaternionic frame, with branch labels
chosen as in the even/odd 600-cell permutation convention, set
$u_\lambda(q)=\sum_{a=0}^3\lambda_aq_a^2$. Then
\[
 \boxed{(\tau_+-\tau_-)(u_\lambda^6)
 =\frac{\sqrt5}{1024}\prod_{0\le a<b\le3}(\lambda_a-\lambda_b).}
\]
Thus degree $12$ is the first degree at which these two finite
cubature presentations can differ; the sign depends on the relative
branch label. This does not constitute an occurring sextic measurement.
\end{proposition}
\begin{proof}
Both completions are spherical eleven-designs. Their common Hurwitz
vertices cancel. The remaining $96$ 600-cell vertices are the signed
even versus odd permutations of
$\bigl(0,\frac12,\frac{\varphi}{2},\frac1{2\varphi}\bigr)$,
$\varphi=(1+\sqrt5)/2$. Their coordinate squares form
$s=(0,1/4,\varphi^2/4,\varphi^{-2}/4)$.
Summing signs gives
\[
 (\tau_+-\tau_-)(u_\lambda^6)
 =\frac1{15}\sum_{\pi\in S_4}\operatorname{sgn}(\pi)
      \Bigl(\sum_a\lambda_as_{\pi(a)}\Bigr)^6.
\]
The alternating multinomial identity yields the factor
$6!/(0!1!2!3!)=60$ multiplying the two Vandermonde products.
As $\prod_{a<b}(s_a-s_b)=\sqrt5/4096$ with the stated ordering,
this equals the displayed coefficient.
\end{proof}

\begin{theorem}[Three-stage finite icosahedral convolution with sharp strength nineteen]
\label{p1b:thm:v273-triple-icosahedral-nineteen-design}
There exist conjugate but differently embedded binary icosahedral groups
$I,J\le SU(2)$ such that their uniform probability measures obey
\[
 \boxed{\nu=\mu_I*\mu_J*\mu_I,
 \qquad \nu|_{\mathcal P_{19}}=\sigma|_{\mathcal P_{19}},}
\]
while $\nu$ is not an exact degree-twenty cubature rule.
Consequently its positive finite word measure is an exact isometric
quadrature for $\mathcal P_9$ (dimension $385$), although the support
is not a finite subgroup. The special relative alignment is a
mathematical existence statement and is not a physical frame selector.
\end{theorem}
\begin{proof}
Let $I$ be the standard binary icosahedral group,
$g_\theta=\cos\theta+\mathbf i\sin\theta$, and
$J=g_\theta I g_\theta^{-1}$. Set
$w=\cos^2\theta(1-\cos^2\theta)$.
The character computation for $\rho_k=\operatorname{Sym}^k(\mathbb C^2)$
shows that, for $0<k\le20$, the only nonzero $I$-invariant
spaces occur at $k=12,20$, each of rank one. If $v_k$ is a unit
$I$-fixed vector, its overlap with the rotated vector is
$F_k(w)=120^{-1}\sum_{\gamma\in I}
 U_k(\operatorname{Re}(g_\theta\gamma))$, and direct quaternionic
character expansion yields
\[
 F_{12}(w)=1-28w+196w^2-420w^3,
\]
\[
 F_{20}(w)=1-\frac{220}{3}w+\frac{4532}{3}w^2
 -\frac{36212}{3}w^3+\frac{119680}{3}w^4-44880w^5.
\]
Since $F_{12}(0)=1$ and $F_{12}(1/4)=-5/16$,
there is $w_0\in(0,1/4)$ with $F_{12}(w_0)=0$.
At that alignment the Fourier operator
$P_I^{(12)}P_J^{(12)}P_I^{(12)}$
vanishes, as do all nonconstant Fourier operators through degree
nineteen; this proves degree-nineteen exactness by Peter--Weyl.
The rational polynomial Euclidean algorithm gives
$\gcd_{\mathbb Q[w]}(F_{12},F_{20})=1$, so
$F_{20}(w_0)\ne0$. The character $U_{20}(q_0)$ has
$\nu$ average $|F_{20}(w_0)|^2\ne0$ and Haar mean zero.
The finite support contains both $I$ and $J$, while distinct
icosahedral subgroups cannot generate a finite subgroup of $SU(2)$,
and no proper positive-dimensional closed subgroup can contain $2I$;
therefore their generated subgroup is dense in $SU(2)$ and the
finite support cannot be multiplication-closed.
\end{proof}

\begin{corollary}[Positive finite coset mixtures have no fixed degree ceiling]
\label{p1b:cor:v273-positive-coset-mixture-no-ceiling}
For every finite subgroup $\Gamma\le SU(2)$ and every finite
polynomial degree $D$, a positive weighted mixture of at most
\[
 d_D=\dim\mathcal P_D=\frac{(D+1)(D+2)(2D+3)}6
\]
translated $\Gamma$-cosets agrees with Haar on $\mathcal P_D$.
For $\Gamma=2I,D=19$ at most $14$ translated cosets suffice,
with at most $1680$ weighted points. The mixture is not required to
be a subgroup or a uniform single-coset orbit.
\end{corollary}
\begin{proof}
Choose a real basis of the $d_D-1$ nonconstant modes and map each
coset $g\Gamma$ to its vector of moment discrepancies.
Haar disintegration makes zero the mean of this compact orbit in
$\mathbb R^{d_D-1}$. Carath\'eodory's theorem gives zero as a
convex combination of at most $d_D$ orbit points. For $2I$ through
degree nineteen, only the degree-twelve scalar-harmonic representation
has nonzero orbit errors; right-$I$ averaging fixes a one-dimensional
factor, leaving a real thirteen-dimensional error orbit. Therefore
at most fourteen cosets suffice. Remove zero coefficients.
\end{proof}


\begin{theorem}[Branch-symmetric finite twin convolution converges to Haar]
\label{p1b:thm:v274-branch-symmetric-twin-convolution}
Fix an embedded binary tetrahedral core $H\le SU(2)$, its
normalizer $O=N_{SU(2)}(H)$, and its unordered pair of binary
icosahedral completions $I_+,I_-$ from
Theorem~\ref{p1:thm:v273-two-icosahedral-completions}. Write
$\mu_\pm=|I_\pm|^{-1}\sum_{g\in I_\pm}\delta_g$ and define the
positive, finite, branch-independent three-stage probability measure
\[
 \xi_H=\tfrac12\bigl(
 \mu_+*\mu_-*\mu_++\mu_-*\mu_+*\mu_-\bigr),
 \qquad \xi_{H,n}=\xi_H^{*n}\quad(n\ge1).
\]
Let $\rho_k=\operatorname{Sym}^k(\mathbb C^2)$, with Fourier transform
$\widehat\nu(k)=\int\rho_k(g)\,d\nu(g)$, and let $P_H,P_O$
denote their Reynolds projectors at $k=12$. Set
$a=275/8192$ and $b=525/8192$. Then
\[
 \widehat{\xi_{H,n}}(k)=0\quad
 (1\le k\le19,\ k\ne12),\qquad
 \boxed{\widehat{\xi_{H,n}}(12)
     =a^nP_O+b^n(P_H-P_O).}
\]
In particular every $\xi_{H,n}$ integrates $\mathcal P_{11}$
exactly, but no finite $n$ integrates $\mathcal P_{12}$ exactly. The
sharp $L^2(SU(2),\sigma)$ dual norm of its complete degree-nineteen
moment error is
\[
 \boxed{\sup_{\substack{f\in\mathcal P_{19}\\\|f\|_2\le1}}
 \left|\int f\,d\xi_{H,n}-\int f\,d\sigma\right|
 =\sqrt{13\bigl(a^{2n}+b^{2n}\bigr)}.}
\]
For every fixed $m$, there is $q_m<1$ such that the Fourier error of
$\xi_{H,n}$ on $\mathcal P_m\ominus E_0$ is bounded in operator norm
by $q_m^n$. Consequently $\xi_{H,n}$ converges weakly to Haar
measure on all continuous scalar observables; the associated averaging
operators converge strongly on $L^2(SU(2))$ to the constant projector.
Every $\xi_{H,n}$ nevertheless has finite non-subgroup support.
The construction is canonical only \emph{relative to the supplied core}
$H$; it is a classical mixture of comparison averages, not a new
physical finite action, coherent quantum superposition, or selector of
an embedded $H$.
\end{theorem}
\begin{proof}
For each $k$, subgroup averaging gives orthogonal projectors
$P_\pm^{(k)}=\widehat\mu_\pm(k)$. For $1\le k\le19$ these vanish
except at $k=12$, where each has rank one. On that two-dimensional
$H$-fixed plane, Theorem~\ref{p2f:thm:v273-twin-projector-spectral-identity}
gives the overlap $|\langle v_+,v_-\rangle|=5/16$ and
\[
 \tfrac12(P_++P_-)=\tfrac{11}{32}P_O
                     +\tfrac{21}{32}(P_H-P_O).
\]
Hence $P_+P_-P_+=(25/256)P_+$, and similarly for the reversed
sandwich. Their mean is $(25/256)(P_++P_-)/2$ and thus has the two
claimed eigenvalues $a,b$. Convolution powers raise those eigenvalues
to the $n$th power. The scalar harmonics are
$E_k\simeq\rho_k\otimes\rho_k^*$, so on $E_{12}$ each left
rank-one eigenvalue has right multiplicity $13$. Peter--Weyl
orthogonality gives the norm of the moment functional as
$\sqrt{13\operatorname{Tr}(\widehat{\xi_{H,n}}(12)^*
\widehat{\xi_{H,n}}(12))}$, proving the stated \emph{equality}.

For any $k>0$, the positive contraction
$S_k=\tfrac12(P_+P_-P_++P_-P_+P_-)$ has no unit eigenvalue.
Indeed equality in $\langle v,S_kv\rangle\le\|v\|^2$ would
force $v$ to be fixed by both $I_+$ and $I_-$. Their generated
subgroup is dense in $SU(2)$ by
Theorem~\ref{p1:thm:v273-two-icosahedral-completions}, whereas a
nontrivial irreducible $\rho_k$ has no $SU(2)$-fixed vector.
Thus $\|S_k\|<1$. Taking a maximum over finitely many $k\le m$
proves the finite-band bound. Peter--Weyl density and positivity of
the probability measures prove weak convergence; spectral domination
proves strong $L^2$ convergence of the averaging operators.
Finally every three-stage word support contains both $I_+$ and
$I_-$ (insert the identities in the remaining factors). The same
is true of every convolution power. A finite subgroup containing both
would contradict their dense generation, so none of these finite
supports is a subgroup.
\end{proof}

\begin{theorem}[Unmarked sixth-moment discriminator for complete quadratic subgroup samplers]
\label{p1b:thm:v268-sixth-moment-discriminator}
Let $\Gamma\subset SU(2)\cong S^3(1)$ be a finite subgroup for which
some normalized weighted rule is exact on every $\mathcal P_2$
Hermitian inner product.  Define its \emph{uniform-support} sixth-moment
tensor and invariant by
\[
 M_6(\Gamma):=\frac1{|\Gamma|}\sum_{q\in\Gamma}q^{\otimes6},\qquad
 M_6^{\rm Haar}:=\int_{S^3}q^{\otimes6}\,d\sigma(q),\qquad
 \mathcal A_6^{\rm supp}(\Gamma)
 :=\|M_6(\Gamma)-M_6^{\rm Haar}\|^2.
\]
The induced Euclidean tensor norm makes this scalar invariant under the
full $O(4)$ action on the support; no quaternionic axis, vertex, or
member orientation is chosen.  Its exact value is
\[
 \boxed{\quad
 \mathcal A_6^{\rm supp}(\Gamma)
 =\frac1{|\Gamma|^2}\sum_{q,q'\in\Gamma}(q\cdot q')^6-\frac5{64}
 =\begin{cases}
 \frac1{64},&[\Gamma]=[2T],\\
 0,&[\Gamma]=[2O]\text{ or }[2I].
 \end{cases}\quad}
\]
In particular, \emph{among occurring group-closed supports that are
exact on the full $\mathcal P_2$ Gram},
$\mathcal A_6^{\rm supp}>0$ is equivalent to the unmarked $[2T]$ type.
This is a type test; exact Gram data alone do not force its positive value.
\end{theorem}
\begin{proof}
The tensor inner-product identity
$\langle q^{\otimes6},(q')^{\otimes6}\rangle=(q\cdot q')^6$ and
$\int_{S^3}(q\cdot v)^6d\sigma(q)=5/64$ for every unit $v$
give the displayed pair-average formula.  It is nonnegative because
it is a squared norm.  For a subgroup, the left-translation bijection
$(q,q')\mapsto(q^{-1}q',q)$ reduces the pair average to
$|\Gamma|^{-1}\sum_{h\in\Gamma}(\operatorname{Re}h)^6$.
In the standard Hurwitz coordinates, the eight axial units contribute
$2$ and the sixteen half-units contribute $16/64$, so
\[
 \frac1{24}\sum_{h\in2T}(\operatorname{Re}h)^6
 =\frac{2+16/64}{24}=\frac3{32},\qquad
 \mathcal A_6^{\rm supp}(2T)=\frac3{32}-\frac5{64}=\frac1{64}.
\]
The uniform $2O$ and $2I$ orbits are respectively spherical $7$- and
$11$-designs, and therefore integrate every homogeneous sextic
exactly: their sixth-moment tensors equal $M_6^{\rm Haar}$.
The finite-subgroup classification in
Theorem~\ref{p1b:thm:v233-subgroup-sampler-classification}
shows there are no further exact $\mathcal P_2$ subgroup types.
$O(4)$-invariance follows from equivariance of the tensor action.
\end{proof}

\begin{remark}[Exact observable scope and finite-readout no-go]
\label{p1b:rem:v268-sixth-moment-observable-scope}
The formula uses canonical equal counting on the \emph{support}; an
instrument with arbitrary physical weights need not reproduce that
sixth moment.  Measuring this invariant therefore requires access to
the occurring point support or a separately realized uniform-support
comparison.  Neither a continuum $\mathcal P_2$ Gram nor positivity,
finite support and antipodality provide group closure: the exact
$72$-node countermodel of Proposition~\ref{p1b:prop:v233-hopf72}
still applies.  Nor does the test create the physical finite role.
Unlike group-type leastness under GP3, a measured positive support
invariant would be an independent, frame-free observable selection
criterion.  The sixth-order obstruction itself is already present in
Paper~I-B's full-cell sextic and is not a new lower-order force law.
\end{remark}

\begin{proposition}[Restriction on an exact image versus restriction of support data]
\label{p1b:prop:v233-sampling-restriction}
Let \(\Gamma'\subset\Gamma\) both be uniformly exact on \(\cP_m\), and
write \(V_\Gamma=S_{\Gamma,m}\cP_m\). Then
\[
 R=\sqrt{|\Gamma|/|\Gamma'|}\,
       \operatorname{res}_{\Gamma'}|_{V_\Gamma},
 \qquad RS_{\Gamma,m}=S_{\Gamma',m}
\]
is unitary from \(V_\Gamma\) onto \(V_{\Gamma'}\).
It transports every operator already defined on \(\cP_m\).
It need not preserve extra data on the full point-support space
\(\mathbb C^\Gamma\).
\end{proposition}
\begin{proof}
The displayed identity and the two isometries prove unitarity on their
images. Conjugation preserves products, adjoints and spectra.
If the inclusion is proper, restriction on \(\mathbb C^\Gamma\)
annihilates every function supported on \(\Gamma\setminus\Gamma'\).
Also, left multiplication by the whole \(\Gamma\) does not in general
act by permutations of \(\Gamma'\). Frozen point tests or that full
permutation action can therefore be lost. Role-faithful subgroup
comparison must check those data separately.
\end{proof}

\subsubsection{Finite scalar units and the scope of reciprocal export}

\begin{theorem}[Invariant scalar line at every finite resolution]
\label{p1b:thm:v233-finite-scalar-unit}
For every \(m\ge0\),
\[
 (\cP_m)^{\Spin(4)}=\mathbb C1.
\]
The same fixed-vector assertion holds for any scalar subspace containing
\(1\), with fixed vectors understood inside the ambient scalar function
space. For a group action on the entire subspace, that subspace must also
be invariant. The reciprocal pullback \(J\) of
Theorem~\ref{p1b:thm:harmonic-transport} satisfies
\[
 J1=1,\qquad J(fg)=(Jf)(Jg),\qquad
 J\overline f=\overline{Jf}.
\]
These statements need neither large resolution nor a new scalar-line
comparison map.
\end{theorem}
\begin{proof}
Transitivity of the round action forces every fixed function to be
constant; constants are present already in \(\cP_0\).
Pullback preserves constants, products and conjugation by evaluation
of the definitions. Its unitarity and harmonic compatibility were proved
in Theorem~\ref{p1b:thm:harmonic-transport}.
\end{proof}

\begin{proposition}[Scalar export into a specified invariant target]
\label{p1b:prop:v233-unit-map}
For a complex \(G\)-module \(A\) with trivial scalar source,
\[
 \operatorname{Hom}_G(\mathbb C,A)\cong A^G,\qquad L\longmapsto L(1).
\]
If \(A\) is a nonzero unital complex star algebra, with a unit-preserving
\(G\)-action, the unique complex-linear unital map is
\[
 \eta_A(z)=zI_A .
\]
It is a covariant star homomorphism, natural under unital algebra maps.
Its image is already inside \(A\); it adds no independent internal
dimension. A new external summand is a separate direct-sum construction.
\end{proposition}
\begin{proof}
Linearity gives \(L(z)=zL(1)\), and covariance is exactly invariance
of \(L(1)\). Unitality forces \(L(1)=I_A\). Multiplication, adjoint and
naturality follow directly. Without unitality or a specified target,
the invariant-vector family remains; if \(A^G=0\), no nonzero such
map exists.
\end{proof}


\begin{theorem}[Exact compressed observables and top-block correction]
\label{p1b:thm:compiler}
Let \(P_m\) project to \(\cP_m\), and let \(Q\) be a spherical \(t\)-design.
For a polynomial symbol \(f\),
\[
 P_mM_fP_m=S_{Q,m}^*\diag(f(q))S_{Q,m}
 \quad\hbox{if }2m+\deg f\le t.
\]
At \(m=2\), guaranteed symbol degrees for \(2T,2O,2I\) are \(1,3,7\).
For a quadratic symbol on \(2T\), the error \(E_f\) satisfies
\[
 E_f=P_{E_2}E_fP_{E_2},\qquad\rank E_f\le9.
\]
For every isometry preserving the band,
\[
 U(g)(P_mM_fP_m)U(g)^*=P_mM_{f\circ g^{-1}}P_m.
\]
\end{theorem}
\begin{proof}
A matrix element integrates \(\overline u f v\) of degree at most
\(2m+\deg f\). If either test mode has degree at most one, a quadratic
symbol gives degree at most five, so only the \(E_2\)-by-\(E_2\) block
can differ for \(2T\). Isometries commute with the band projection and
conjugate multiplication as displayed. The degree-six moment defect
determines the exact correction.
\end{proof}

\begin{theorem}[Full fifteen-dimensional fibre sampling]
\label{p1b:thm:full-sampling}
On \(S_R^3\), let \(\mathbb S_\pm\) be the two lowest spinor spaces,
of complex dimension two. In
\(\cE_R=(\Lambda^*T^*S_R^3)_\C\oplus\mathbb S\), set
\[
\begin{split}
 \mathcal L_R={}&\C V_R^{-1/2}
 \oplus\operatorname{span}\left\{\frac{2\,\dd x_a}{\sqrt{3V_R}}\right\}
 \\
 &\oplus\operatorname{span}\left\{\sqrt{\frac2{V_R}}\dd x_a\wedge\dd x_b\right\}_{a<b}
 \oplus\mathbb S_+\oplus\mathbb S_-.
\end{split}
\]
This is a fifteen-dimensional orthogonal representation with sectors
\[
 (0,0),(\tfrac12,\tfrac12),(1,0),(0,1),
 (\tfrac12,0),(0,\tfrac12).
\]
For \(Q=2T\),
\[
 (\mathsf S_Rs)_q=\sqrt{V_R/24}\,s(Rq),\qquad
 \mathsf S_R^*\mathsf S_R=I.
\]
For the compressed intrinsic symbol
\(C(f)=P_{\mathcal L_R}c(\dd f)P_{\mathcal L_R}\),
exact sampling transports it unitarily to the sampled image.
\end{theorem}
\begin{proof}
The integrated one-form and two-form pairings are
\(3V_R\delta_{ab}/4\) and
\(V_R(\delta_{ac}\delta_{bd}-\delta_{ad}\delta_{bc})/2\).
Different form degrees and spinors are orthogonal.
In unit coordinates, one-form pairings have degree two; in two-form
pairings the quartic terms cancel. Lowest spinor sections in a left frame
are constant spinors or \(q^{-1}\) times constants, with pairings of
degree at most one. The \(2T\) design integrates every retained pairing.
The target fibre has complex dimension ten and the full Clifford
symbol acts on it; the retained domain has zero three-form component.
Conjugation by the isometry transports every compressed operator.
\end{proof}
For any second injective geometric evaluation \(T\), its positive-polar
normalization \(T(T^*T)^{-1/2}\) gives another isometric realization.
This identity supplies a comparison only after the evaluation itself
is specified.

\begin{corollary}[Eight-fibre native comparison]
\label{p1b:cor:v233-eight-fibre-export}
The full fifteen-dimensional carrier of
Theorem~\ref{p1b:thm:full-sampling} also has the isometry
\((T_8s)_q=\sqrt{V_R/8}s(Rq)\) over \(Q_8\).
Every already constituted operator and its image, action, tensor and
blocking transports are exact on this image. The support here is an exact sampled presentation of the specified smooth carrier.
\end{corollary}
\begin{proof}
Corollary~\ref{p1:cor:v233-eight-fibre-sampling} proves the isometry.
The image transport identities in this paper apply to any isometry.
The target remains fibre-valued, and pointwise diagonal-symbol
quadrature retains its separate degree requirement.
\end{proof}

\begin{corollary}[Charge-linked degree-two spectra and polynomial integrals]
\label{p1b:result:charge-linked-degree-two-spectra-and-polynomial-integrals}
On \(\cP_2\), explicitly represent the relative scalar generator by
the marked Hopf charge. The operator
\(A\mathsf C+BQ^2+C_0I\), with scalar coefficients \(A,B,C_0\)
and sphere Casimir \(\mathsf C\), has values
\[
 C_0,\quad3A+B+C_0,\quad8A+C_0,\quad8A+4B+C_0
\]
with multiplicities \(1,4,3,6\).
For any polynomial function of a real field in \(\cP_2\) of degree
at most four, its round integral equals the uniform \(2I\) sum,
as do all derivatives with respect to its mode coefficients.
\end{corollary}
\begin{proof}
Use the degree-zero, one and two sphere Casimirs \(0,3,8\) and
their Hopf weights. For the integral each field has coordinate
degree at most two, so every term has degree at most eight.
Exact design integration makes the two expressions identical
polynomials in the mode coefficients; differentiate that identity.
\end{proof}

\begin{theorem}[Degree-two comparison of generated action algebras]
\label{p1b:prop:degree-two}
Under one handed \(\SU(2)\), the full carrier and scalar bands restrict as
\[
 \mathcal L_R=6V_0\oplus3V_{1/2}\oplus V_1,\qquad
 \cP_m=\bigoplus_{k=0}^m(k+1)V_{k/2}.
\]
The unital complex star algebras generated by that action are
\[
 \cA_{\mathcal L}=\C\oplus M_2(\C)\oplus M_3(\C),\qquad
 \cA_m=\bigoplus_{k=0}^m M_{k+1}(\C).
\]
A generator-preserving unital star epimorphism
\(\cA_m\to\cA_{\mathcal L}\) exists exactly when \(m\ge2\).
Requiring also exact uniform \(2I\) integration of \(f^4\) for every
real \(f\in\cP_m\) therefore gives the unique comparison value \(m=2\).
At that value the algebra isomorphism induces an affine bijection of
state spaces, while the inherited normalized traces are
\[
 \tau_{\mathcal L}(a_0,a_1,a_2)
 =\frac{6a_0+3\tr_2a_1+\tr_3a_2}{15},\qquad
 \tau_2(a_0,a_1,a_2)
 =\frac{a_0+2\tr_2a_1+3\tr_3a_2}{14}.
\]
Moreover \(\dim_\C\Hom_{\SU(2)}(\cP_2,\mathcal L_R)=15\),
and the maximal rank of such an intertwiner is eight.
\end{theorem}
\begin{proof}
The Casimir separates the irreducible spins. On each spin the generators
generate the full matrix algebra, acting identically on its multiplicity
space. Projection onto spins \(0,\tfrac12,1\) is the required epimorphism;
an algebra without the \(M_3\) summand cannot surject onto it.
The degree-twelve moment of \(2I\) fails for \(f=x_0^3\), excluding every
\(m\ge3\), whereas degree-two fourth powers are exact.
A unital star isomorphism preserves positive normalized functionals.
Multiplicity counting gives the two traces and the intertwiner dimension
\(1\cdot6+2\cdot3+3\cdot1=15\). Its maximal rank is
\(1\min(1,6)+2\min(2,3)+3\min(3,1)=8\).
\end{proof}
On restriction to the marked Hopf circle, the multiplicities of weights
\((-2,-1,0,1,2)\) are \((3,2,4,2,3)\) and \((1,3,7,3,1)\).
Thus the corresponding intertwiner dimension is \(46\) and maximal
rank is \(10\). These comparisons concern representations with their
stated multiplicities; they do not identify the two carrier spaces.

\begin{corollary}[Exact conversion of the inherited degree-two traces]
\label{p1b:cor:degree-two-trace-density}
Use the generator-preserving identification of
Theorem~\ref{p1b:prop:degree-two}, and write the common algebra as
\[
 \mathcal A=\mathbb C\oplus M_2(\mathbb C)\oplus M_3(\mathbb C).
\]
There is a unique positive central element \(h\in\mathcal A\)
such that \(\tau_{\mathcal L}(a)=\tau_2(ha)\) for every
\(a\in\mathcal A\). It is invertible and equals
\begin{equation}
\label{p1b:eq:degree-two-trace-density}
 h=\left(\frac{28}{5},\frac75I_2,\frac{14}{45}I_3\right),
 \qquad
 \tau_2(h)=1,\qquad
 \tau_2(a)=\tau_{\mathcal L}(h^{-1}a).
\end{equation}
If a state has density \(\rho_{\mathcal L}\ge0\) relative to
\(\tau_{\mathcal L}\), its density relative to \(\tau_2\) is
\[
 \rho_2=h\rho_{\mathcal L},\qquad
 \tau_2(\rho_2a)=\tau_{\mathcal L}(\rho_{\mathcal L}a),
 \qquad \tau_2(\rho_2)=1.
\]
Conversely \(\rho_{\mathcal L}=h^{-1}\rho_2\).
\end{corollary}

\begin{proof}
The coefficients of the ordinary block traces in
\(\tau_{\mathcal L}\) are \((6,3,1)/15\); those in
\(\tau_2\) are \((1,2,3)/14\). Their three ratios are
\(28/5,7/5,14/45\). Equality on each matrix block proves the
formula and uniqueness. All ratios are positive. Setting \(a=I\)
gives \(\tau_2(h)=1\), and substituting \(h^{-1}a\) gives the
inverse conversion. Centrality of \(h\) makes
\(h\rho_{\mathcal L}=h^{1/2}\rho_{\mathcal L}h^{1/2}\)
positive; the trace identity proves the state formulas.
\end{proof}

The density converts two inherited normalizations on the common
action algebra. It does not identify the fifteen- and
fourteen-dimensional representing spaces or replace a parent
action trace. Since \(h\) is central, the conversion introduces
no basis rotation inside an irreducible action block. In
particular, the \(M_3\) block here is the spin-one action block;
its dimension alone does not identify it with a family carrier.

\begin{corollary}[Polynomial detectability frontier]
\label{p1b:result:polynomial-detectability-frontier}
The first nonconstant \(2I\)-invariant harmonic, of degree \(12\), belongs
to the span of products of at most \(r\) band-\(m\) factors exactly when
\(rm\ge12\). For \(r=1,2,3,4\), the thresholds are \(m=12,6,4,3\).
\end{corollary}
\begin{proof}
Below degree twelve its harmonic projection is zero. At or above that
degree, each degree-twelve monomial factors into at most \(r\) monomials
of degree at most \(m\); restricting to the sphere proves sufficiency.
\end{proof}


\subsection{Finite scalar lattices and marked exceptional incidence}
\label{p1b:sec:modular}
\begin{proposition}[Torsion lattices and symplectic frames]
\label{p1b:prop:finite-modular}
For an odd prime \(p\), the scalar \(p\)-torsion is
\(\T^2[p]\cong\F_p^2\), with alternating pairing
\(\langle(a,b),(c,d)\rangle=ad-bc\).
Its pairing-preserving automorphism group is \(SL(2,\F_p)\).
The set of ordered symplectic frames is a torsor of size \(p(p^2-1)\);
the set of cyclic order-\(p\) axes is \(\mathbb P^1(\F_p)\), of size
\(p+1\). Abstractly,
\[
 SL(2,\F_3)\cong2T,\qquad SL(2,\F_5)\cong2I.
\]
For \(p=3\), let
\[
 A_0=\begin{pmatrix}1&1\\1&2\end{pmatrix},\quad
 A_1=\begin{pmatrix}2&1\\1&1\end{pmatrix},\quad
 A_2=A_0A_1=\begin{pmatrix}0&2\\1&0\end{pmatrix}.
\]
Then \(A_i^2=-I\), \(A_0A_1=-A_1A_0\), and
\(\{\pm I,\pm A_0,\pm A_1,\pm A_2\}\cong Q_8\) is normal,
with quotient \(C_3\).
\end{proposition}
\begin{proof}
In dimension two the alternating form transforms by determinant.
Choose a nonzero first vector in \(p^2-1\) ways; its determinant-one
partners form an affine line of \(p\) vectors. Dividing nonzero vectors
by the \(p-1\) nonzero scalars gives \(p+1\) axes.
The standard binary polyhedral identifications are developed in
Paper I-B. In the displayed \(p=3\) model direct multiplication gives
the quaternion relations. The six noncentral displayed elements are
exactly the determinant-one trace-zero matrices over \(\F_3\), so
their union with the center is invariant under conjugation. Its index
is three.
\end{proof}
Identifying a symplectic-frame torsor with its group requires a chosen
frame. These finite scalar constructions and marked sphere subgroups
can be compared through specified isomorphisms.

The analogous order count at modulus four does not give
\(SL(2,\Z/4)\cong2O\). Indeed the matrices \(I+2X\), with
\(X\in M_2(\F_2)\) and \(\tr X=0\), are eight distinct
determinant-one matrices of square \(I\); seven are nonidentity.
The subgroup \(2O\subset\SU(2)\) has only one nonidentity
involution, namely \(-I\).

\begin{proposition}[Exact aligned exceptional Hopf--Clifford fingerprint]
\label{p1b:prop:exceptional-hopf-clifford}
Use quaternion coordinates
\(q=a+b\mathbf i+c\mathbf j+d\mathbf k\leftrightarrow
(z_1,z_2)=(a+ib,c+id)\), the Hopf map
\[
 h(q)=\bigl(2(ac+bd),\ 2(bc-ad),\ a^2+b^2-c^2-d^2\bigr),
\]
and the balanced Clifford torus
\[
 T^2_{\rm Cl}=\{a^2+b^2=c^2+d^2=1/2\}.
\]
For the standard aligned binary tetrahedral, octahedral and icosahedral
subgroups, exact arithmetic gives
\[
\begin{array}{c|c|c|c|c|c}
\Gamma&|\Gamma|&|\Gamma\cap T^2_{\rm Cl}|&|h(\Gamma)|
 &|\Gamma\cap h^{-1}(\{N,S\})|&N/T^2_{\rm Cl}/S\\ \hline
2T&24&16&6&8&4/16/4\\
2O&48&32&6&16&8/32/8\\
2I&120&16&30&8&52/16/52.
\end{array}
\]
Here the final column counts points with positive, zero and negative
third Hopf coordinate.  In these aligned realizations
\[
 2T\cap2O=2T,\qquad 2T\cap2I=2T,\qquad 2O\cap2I=2T.
\]
The Hopf images of \(2T\) and \(2O\) are exactly the six vertices of
the aligned octahedron.  The sixteen points
\[
 \mathcal H_{16}=\left\{
 \frac{\pm1\pm\mathbf i\pm\mathbf j\pm\mathbf k}{2}
 \right\}
\]
are exactly \(2T\cap T^2_{\rm Cl}\); this set is not a subgroup of
\(2T\).
\end{proposition}
\begin{proof}
In the standard coordinates,
\[
 2T=Q_8\cup\mathcal H_{16}.
\]
Every half-Hurwitz unit has
\(a^2+b^2=c^2+d^2=1/2\), whereas the eight points of \(Q_8\) lie on
the two core Hopf fibres.  This gives the \(2T\) Clifford and hemisphere
counts.  Direct substitution in \(h\) shows that its six Hopf images are
\(\{\pm e_1,\pm e_2,\pm e_3\}\).  Equivalently each fibre has four
points because \(2T\cap\U(1)_H=\{\pm1,\pm\mathbf i\}\).

The group \(2O\) adds the twenty-four roots with exactly two nonzero
coordinates, each equal to \(\pm1/\sqrt2\).  Sixteen have one nonzero
coordinate in each complex pair and therefore lie on
\(T^2_{\rm Cl}\); the remaining eight lie on the two core Hopf fibres.
Thus its boundary/core and hemisphere counts are \(32,16\) and
\(8/32/8\).  Its Hopf-circle stabilizer has order eight, so orbit--fibre
counting gives \(|h(2O)|=48/8=6\), and substitution gives the same six
octahedral vertices.

For \(2I\), the additional ninety-six points are the even coordinate
permutations of
\[
 \tfrac12(0,\ \pm1,\ \pm\varphi,\ \pm\varphi^{-1}),
 \qquad \varphi=\frac{1+\sqrt5}{2},
\]
with independent signs on the three nonzero entries.  Before the common
factor \(1/4\), their squared coordinate multiset is
\(\{0,1,\varphi^2,\varphi^{-2}\}\).  Since
\(\varphi^2=\varphi+1\) and \(\varphi^{-2}=2-\varphi\), none of the six
pair sums is \(2\); hence none of these ninety-six points lies on the
Clifford torus.  Each has exactly one zero coordinate, so none lies on a
core fibre.  The even pair-swap
\((a,b,c,d)\mapsto(c,d,a,b)\) preserves the set and reverses the third
Hopf coordinate; therefore the ninety-six nonboundary points split
\(48/48\) between the two open Heegaard chambers.  Adding the eight core
points of \(2T\) gives \(52/16/52\).  Also
\(2I\cap\U(1)_H=\{\pm1,\pm\mathbf i\}\), so every Hopf fibre meeting
\(2I\) contains four group points and \(|h(2I)|=120/4=30\).

Finally the \(2O\setminus2T\) roots have exactly two zero coordinates,
whereas the \(2I\setminus2T\) points have exactly one, proving the
pairwise intersections.  The assertion that \(\mathcal H_{16}\) is not
a subgroup also follows from Lagrange's theorem, since \(16\nmid24\).

\end{proof}

\begin{proposition}[Finite Weyl algebra]
\label{p1b:result:finite-weyl-algebra}
On \(\C^p\), set \(Xe_j=e_{j+1}\), \(Ze_j=\zeta^je_j\),
where \(\zeta=e^{2\pi i/p}\) and indices are modulo \(p\).
Then \(ZX=\zeta XZ\), and \(\{X^aZ^b\}_{a,b\in\F_p}\) is an
orthonormal basis of \(M_p(\C)\) for \(p^{-1}\Tr(A^*B)\).
\end{proposition}
\begin{proof}
The commutation relation follows on each basis vector. A nonidentity
shift has zero diagonal; a nontrivial diagonal character sums to zero.
Thus the \(p^2\) words are orthogonal and exhaust the matrix dimension.
\end{proof}

\begin{proposition}[Five cells and two distinct permutation sets]
\label{p1b:prop:five-cells}
Let \(G=2I\), \(H=2T\) be a marked subgroup, and \(X=G/H\).
Then \(|X|=5\), every order-five subgroup acts regularly on \(X\), and
\(\C^X=\C\oplus V_4\), where \(V_4\) is irreducible.
The diagonal \(G\)-action on \(X\times X\) has orbits of sizes \(5,20\).
The defining \(SL(2,\F_5)\)-action on \(\F_5^2\) has orbits of sizes
\(1,24\). Consequently these two \(25\)-element sets are not
equivariantly isomorphic.
\end{proposition}
\begin{proof}
The action factors through the usual \(A_5\)-action on five letters.
It is doubly transitive, so the augmentation module is irreducible
and the diagonal and off-diagonal pairs are the two orbits.
An order-five subgroup has no nonidentity stabilizer in \(H\), whose
order is \(24\), proving regularity. In the defining linear action,
zero is fixed and every nonzero vector extends to a symplectic frame.
\end{proof}

For a supplied incidence map \(B:\C^{X\times X}\to\cH\) with the
normalization
\[
 B^*B=120(P\otimes P)+30(Q\otimes Q),\qquad
 P=\tfrac15J_5,\quad Q=I-P,
\]
the orthogonal decomposition has dimensions \(1,4,4,16\).
Hence \(\rank B=17\), \(\dim\ker B=8\), and its normalized polar
map is isometric on \(\C\oplus(V_4\otimes V_4)\).
This conclusion is immediate from the displayed positive Gram matrix.
The particular incidence construction and its counting normalization
are the finite datum of Papers V-B and III-F; an arbitrary
\(25\)-column map need not have this Gram matrix.

\subsubsection{Petersen--Hopf packets on the two reciprocal factors}
\label{p1b:sec:petersen-reciprocal-packets}

The complete finite $5\to10\to15$ construction of Paper~I-B,
Section~\ref{p1:sec:petersen-finite-resolution}, may be attached to
either labelled $SU(2)$ factor of $M_8$ under a \emph{chosen} embedded
binary-icosahedral subgroup and its five tetrahedral cores. The
underlying unmarked Lie group $M_8$ selects neither the position nor
the momentum packet. In particular, the following is a construction
on marked finite representation modules, not a new factor of $M_8$.

\begin{theorem}[Factorwise finite packet and Born-exchange covariance]
\label{p1b:thm:petersen-factor-exchange}
Choose binary-icosahedral packet markings $I_x\subset SU(2)_x$ and
$I_p\subset SU(2)_p$, with compatible identification of the standard
five-label constructions. Let $W_f=C_1(\mathcal P_f;\mathbb R)$ be
the oriented real edge space ($f=x,p$), with the four Paper~I-B
projectors $\Pi_{4,f},\Pi_{5,f},\Pi_{3,f},\Pi_{3',f}$, of ranks
$r_4=4,r_5=5,r_3=r_{3'}=3$. Then
\[
 W_f=\mathbf4_f\oplus\mathbf5_f\oplus\mathbf3_f\oplus\mathbf3'_f,
 \qquad\dim W_f=15,
\]
and on the reciprocal tensor product
\[
 \boxed{W_x\otimes W_p
 =\bigoplus_{\alpha,\beta\in\{4,5,3,3'\}}
  (\Pi_{\alpha,x}\otimes\Pi_{\beta,p})(W_x\otimes W_p),}
 \qquad\dim(W_x\otimes W_p)=225.
\]
The sixteen displayed projectors are pairwise orthogonal, sum to the
identity, have ranks $r_\alpha r_\beta$, and commute with the
independent $A_{5,x}\times A_{5,p}$ action. Relative to transported
and matched packet markings, Born exchange $B$ satisfies
\[
 B(\Pi_{\alpha,x}\otimes\Pi_{\beta,p})B^{-1}
   =\Pi_{\beta,x}\otimes\Pi_{\alpha,p}.
\]
Neither the unmarked product nor Born exchange alone determines the
markings or a mixed physical source. The sixteen-projector
statement concerns the independent product representation; the
commutant of the \emph{diagonal} $A_5$ action may be larger.
\end{theorem}
\begin{proof}
The factorwise projector identities are transported from the exact
Paper~I-B real edge resolution. Tensor products of orthogonal
projections give sixteen orthogonal projections, their ranks
multiply, and $\sum_\alpha\Pi_{\alpha,f}=I_{15}$ in each factor.
Exchange is the usual factor-flip after using the specified common
marking identification; it transports each tensor product projector
to the one with reversed labels.
\end{proof}

\begin{corollary}[Exchange-even and exchange-odd finite edge tensors]
\label{p1b:cor:petersen-born-even-odd}
If the two factor markings and coefficient inner products are
identified so that $W_x=W_p=W$, then the exchange flip $\tau$ on
$W\otimes W$ has exact eigenspaces
\[
 \dim\operatorname{Sym}^2(W)=\tfrac12(15^2+15)=120,
 \qquad\dim\bigwedge^2W=\tfrac12(15^2-15)=105.
\]
For each diagonal sector $\alpha\otimes\alpha$, exchange-even and
exchange-odd dimensions are respectively
$r_\alpha(r_\alpha+1)/2$ and $r_\alpha(r_\alpha-1)/2$.
For every unordered pair $\alpha\ne\beta$, the combined
$(\alpha,\beta)\oplus(\beta,\alpha)$ block contributes
$r_\alpha r_\beta$ to each exchange parity. The number $120$ here
is the dimension of a symmetric tensor square, not an identification
with the $120$ distinct elements of $2I$.
\end{corollary}
\begin{proof}
In a basis of fifteen edge channels, the tensor flip fixes the
fifteen diagonal basis tensors and permutes the other $210$ in
pairs; hence $\operatorname{tr}\tau=15$, giving the dimensions.
Apply the same argument to each invariant summand, and pair
distinct summands under the flip.
\end{proof}

\begin{proposition}[Marking orbit and absence of a new $M_8$ invariant]
\label{p1b:prop:petersen-marking-functor}
An unlabelled embedded $2I$ packet in either $SU(2)$ has conjugacy
orbit $SU(2)/2I$, since $N_{SU(2)}(2I)=2I$. A jointly marked pair
$(I,H\cong2T)$, with $H\subset I$, has orbit $SU(2)/2T$.
Consequently the independent two-factor family of such marked pairs
is parametrized by
\[
 (SU(2)_x/2T_x)\times(SU(2)_p/2T_p),
\]
with Born exchange interchanging the factors under the chosen label
identification. Forgetting an embedded completion retains the
already proved two-sheeted normalizer fork over a tetrahedral core,
not a preferred global branch. In particular:
\begin{enumerate}
\item no additional dimension or topological class is forced on
$M_8\cong S^3\times S^3\times T^2$;
\item the scalar $U(1)$ factors and their integral character lattice
are unchanged by the projective Petersen module;
\item an invariant physical bilinear, link holonomy, kinetic spectrum,
or finite interaction requires an independently typed state/action
attachment rather than an equality of graph and dynamical operators.
\end{enumerate}
\end{proposition}
\begin{proof}
The normalizers of the selected exceptional groups are the ones
proved in Paper~I-B: $N_{SU(2)}(2I)=2I$ and
$N_{2I}(2T)=2T$, giving the indicated stabilizers. The reciprocal
product retains both $SU(2)$ factors and both central circles as
separate direct-product factors. A marked finite representation
module is auxiliary data and does not change the topology or
character group of the ambient manifold.
\end{proof}

\begin{remark}[Precise reciprocal and physical scope]
Factor exchange is a geometrically available symmetry of the labelled
product, but invariance of a chosen metric or physical operator under
exchange must be checked in its own theorem domain. The signs of the
Petersen $\mathbb Z_2$ covering cannot be identified with an arbitrary
$U(1)$-scalar holonomy, the center $\mathbb Z_3$ of color $SU(3)$, or
a hadronic mass coefficient. A diagonal subgroup $\Delta2I$ is a
different marked action from independent $2I_x\times2I_p$ and may
have additional invariant tensors; none is silently selected here.
\end{remark}

\begin{proposition}[Recentered two-plane pencil]
\label{p1b:prop:two-plane}
Let \(u_0,u_{\rm out}\) be orthonormal. Set
\[
 u_I=\frac{u_0+2u_{\rm out}}{\sqrt5},\quad
 v_I=\frac{2u_0-u_{\rm out}}{\sqrt5},\quad
 u_\pm=\frac{u_0\pm u_{\rm out}}{\sqrt2}.
\]
For \(P_I=|v_I\rangle\langle v_I|\), \(P_O=|u_-\rangle\langle u_-|\),
\[
 |\langle v_I,u_-\rangle|^2=\frac9{10},\qquad
 \Spec(aP_I+bP_O)=
 \left\{\frac{a+b\pm\sqrt{(a-b)^2+18ab/5}}2\right\}.
\]
In particular
\(\Spec(\frac54P_I+2P_O)=\{(13\pm3\sqrt{17})/8\}\).
\end{proposition}
\begin{proof}
\(u_- =(-u_I+3v_I)/\sqrt{10}\).
The pencil has trace \(a+b\) and determinant \(ab/10\); solve its
quadratic characteristic polynomial.
\end{proof}

\subsection{Joint comparison tensors and their global lines}
\label{p1b:sec:joint-lines}
\begin{theorem}[Even relative comparison and its quadratic covariant]
\label{p1b:thm:joint-tensor}
Let \(c:\HH\to\Hom(\C_R^2,\C_L^2)\) be quaternionic Clifford
evaluation and \(u\) a retained real unit. Then
\[
 A(u,q)=c(u)^*c(q)=\rho(\bar u q)
\]
is independent of opposite chiral rephasings and transforms by
right-spin conjugation under common two-input spin transport.
With \(\varepsilon(s,t)=s^T\left(\begin{smallmatrix}0&1\\-1&0\end{smallmatrix}\right)t\),
the polynomial
\[
 \mathcal T_u(q;s)=\Tr A(u,q)\,\varepsilon(s,A(u,q)s)
\]
has degree two in each of \(q,s\). For \(u=1\), \(q=r+x\mathbf i+
y\mathbf j+z\mathbf k\), it is \(s^TS(q)s\), where
\[
 S(q)=2r\begin{pmatrix}-y+iz&-ix\\-ix&-y-iz\end{pmatrix}.
\]
The space of such homogeneous quadratic vector covariants under full
\(\SU(2)\) conjugation is one-dimensional. On the positive Hopf line
with axis \(n\) and real axial line \(\C h\), its restriction \(m\)
satisfies
\[
 \|m\|^2=4r^2(|v|^2-(n\cdot v)^2),\qquad \bar u h=r+v.
\]
For \(n=\mathbf i\), \(h=(1+\mathbf i+\mathbf j+\mathbf k)/2\),
its coefficient is \((-1+i)/2\) and \(\|m\|^2=1/2\).
\end{theorem}
\begin{proof}
Opposite chiral phases cancel in \(c(u)^*c(q)\). Under transport
\(c(q)\mapsto g_Lc(q)g_R^{-1}\), the same cancellation leaves
conjugation by \(g_R\). Trace and \(\varepsilon\) are invariant.
Matrix multiplication gives \(S\). A quadratic rotation-equivariant
map on \(\R\oplus\R^3\) to \(\R^3\) is a multiple of \(rv\):
there is no invariant vector, and a half-turn reversing \(v\) rules
out a quadratic \(v\)-only term. Complexification preserves this
classification. Restriction in the \(n=\mathbf i\) frame gives
coefficient \(2r(-y+iz)\); covariance yields its stated norm.
\end{proof}
Finite Hurwitz conjugation by itself permits an additional covariant
\((yz,zx,xy)\). The continuous symmetry in the uniqueness statement
is therefore an essential hypothesis. Also
\(2T\cap\{\cos t+n\sin t\}\) is \(C_4\) for coordinate axes,
\(C_6\) for body diagonals, and \(C_2\) otherwise, as follows by
examining the imaginary directions of the Hurwitz units.


\subsubsection{Root-normalized restrictions to the two Hopf lines}
\label{p1b:sec:v135-hopf-sector-phase}
\begin{theorem}[Conjugate sector phases from the joint quadratic]
\label{p1b:thm:v135-hopf-sector-phase}
In the Cartan frame of Theorem~\ref{p1b:thm:joint-tensor}, retain the
ordered rooted trimer
\[
 q_0=\tfrac12(1-\mathbf i+\mathbf j+\mathbf k),\quad
 q_1=\tfrac12(1+\mathbf i-\mathbf j+\mathbf k),\quad
 q_2=\tfrac12(1+\mathbf i+\mathbf j-\mathbf k).
\]
Let $s_+=(1,0)^T$, $s_-=(0,1)^T$ represent the two Hopf lines and put
$m_\sigma(q)=s_\sigma^TS(q)s_\sigma$.
Normalizing each restriction at the same root $q_0$ gives
\[
 \left(\frac{m_+(q_j)}{m_+(q_0)}\right)_{j=0}^2=(1,-i,i),
 \qquad
 \left(\frac{m_-(q_j)}{m_-(q_0)}\right)_{j=0}^2=(1,i,-i).
\]
Thus the structural restrictions are $R_G=\operatorname{diag}(1,-i,i)$
and $R_G^\dagger$.  They do not depend on the unit frames of the Hopf lines.
\end{theorem}
\begin{proof}
The displayed joint tensor gives $m_\pm(q)=2r(-y\pm iz)$.
Its values on the positive line are
$\tfrac12(-1+i,1+i,-1-i)$; the negative-line values are their conjugates.
All root coefficients are nonzero, and division gives the two triples.
A line-frame change $s_\sigma\mapsto e^{i\theta_\sigma}s_\sigma$
multiplies numerator and denominator by $e^{2i\theta_\sigma}$, which cancels.
\end{proof}

\begin{corollary}[One joint structural phase operator]
\label{p1b:cor:v135-joint-phase-operator}
For Hopf projectors $P_\pm$ and ordered trimer projectors $P_0,P_1,P_2$,
\[
 \mathscr R=P_+\otimes R_G+P_-\otimes R_G^\dagger
 =I\otimes P_0-i(P_+-P_-)\otimes(P_1-P_2)
\]
is unitary. With the retained normalization $P_+-P_-=2T_R^3$, its
two restrictions give precisely the sector phases above.
\end{corollary}
\begin{proof}
The Hopf summands are orthogonal, and each diagonal restriction has unit
modulus. Expanding the two diagonal matrices proves the second formula.
\end{proof}
This is a relation of the joint geometric parent. Its typed descent supplies
structural sector data; an operator on another carrier still requires an explicit intertwiner. An internal interchange of the Hopf lines does not reverse the
retained trimer cycle. Reversal of the whole rooted tuple also reverses that
cycle and its oriented Hermitian sign line.

\begin{theorem}[Joint orbit, quadratic descent and determinant]
\label{p1b:thm:joint-orbit}
The common conjugation orbit of the complete tuple
\[
 (1,2T,\mathbf i,(1+\mathbf i+\mathbf j+\mathbf k)/2)
\]
is \(\mathcal O=\SU(2)/\{\pm1\}\cong\SO(3)\).
Its positive-spin-line tensor axial-line bundle is
\(N=\SU(2)\times_{\{\pm1\}}\C_{\rm sign}\). Hence
\[
 \dim_\C\Gamma_{\rm eq}(\mathcal O,N^{*\otimes p})=
 \begin{cases}0&p\text{ odd},\\1&p\text{ even}.\end{cases}
\]
The preceding tensor supplies a nowhere-zero section \(m\) of
\(N^{*2}\), of constant norm, and the unique unitary connection
preserving \(m\) is flat with holonomy \(-1\) on the nontrivial loop.
The class \(c_1(N)\) generates \(H^2(\SO(3);\Z)=\Z/2\).

For any line \(L\) and ordered rank-four split \(Q=\ell\oplus C\)
with \(\rank C=3\), there is a natural isomorphism
\[
 (L\otimes\ell)^{*2}\cong
 \det(L^{*2}\oplus\Lambda^2C)\otimes(\det Q)^{-2}.
\]
\end{theorem}
\begin{proof}
The Hopf axis and the axial imaginary vector are nonparallel; fixing
both fixes a frame. The stabilizer is therefore the two spin lifts
of the identity rotation. It acts by sign on the spin line and
trivially on the axial line. Equivariant sections are precisely
isotropy-fixed fibre vectors, giving the parity rule. The computed
nonzero coefficient descends because two signs cancel.
Local unit frames in which \(m\) has its constant coefficient have
locally constant sign transitions. Declaring them parallel defines
the unique preserving unitary connection.
The Bockstein of the nontrivial double cover generates integral
\(H^2(\mathbb{RP}^3)\), giving \(c_1(N)\).
Finally \(\det(\Lambda^2C)=(\det C)^2\), while
\(\det Q=\ell\otimes\det C\); cancel the factors.
\end{proof}

\begin{theorem}[Quadratic lines on \(M_8\)]
\label{p1b:thm:global-quadratic-line}
Suppose a Hermitian line \(N\to M_8\) and an everywhere nonzero
\(m\in\Gamma(N^{*2})\) are given. Then \(c_1(N)=0\).
The normalized real square-root line has a sign class
\(\alpha\in H^1(M_8;\Z/2)=(\Z/2)^2\), and admits a nonzero global
real frame exactly when \(\alpha=0\).
A preserving unitary connection exists exactly when \(\|m\|\) is
constant, and is then unique and flat, with scalar-cycle holonomies
\((-1)^{\alpha_t},(-1)^{\alpha_E}\).

If \(c_1(N)=d\eta\) without the nonzero-section hypothesis, a
nowhere-zero section of \(N^{*2}\otimes\mathcal P^k\), with
\(c_1(\mathcal P)=\eta\), is possible precisely when \(k=2d\).
For spinor bundles, tensoring by the flat sign line of class
\(\alpha\) changes the spin structure from \(s\) to \(s+\alpha\).
\end{theorem}
\begin{proof}
Nonzero \(m\) implies \(2c_1(N)=0\), and \(H^2(M_8;\Z)\) is
torsion-free. In a global unit frame write \(m=\sigma f(e^*)^2\)
with \(f:M_8\to\U(1)\). Its winding class modulo two is unchanged
by changing \(e\); a square root exists precisely when both windings
are even. Local square-root frames have sign transitions. The
connection argument is as above, while unitary parallelism forces
constant norm. The Chern class of the tensor line is
\((k-2d)\eta\); complex lines are classified by this class.
Moving the same locally constant signs between spin transition
functions and the tensor line proves the spin-twist assertion.
\end{proof}
The parameter orbit \(\mathcal O\) and the base \(M_8\) have different
second cohomology. A pullback between them requires an actual map.

\subsection{Clifford block geometry and positive trace calculus}
\label{p1b:sec:block-calculus}
\begin{proposition}[Quaternionic rank blocks and off-diagonal edges]
\label{p1b:prop:block-edges}
Give \(\Cl_{0,3}\) its standard real coefficient inner product.
Its even algebra is \(\HH\). Choose an imaginary unit \(u\), and set
\(V=u^\perp\subset\operatorname{Im}\HH\).
Let \(Q_0,Q_4,Q_5,Q_7\) project respectively onto the odd algebra,
\(\R u,V,\R1\), of ranks \(4,1,2,1\). For \(v\in V\), put
\[
 T(v)=Q_5L_vQ_4=|vu\rangle\langle u|,\quad
 S(v)=T(v)+T(v)^*,\quad J_v=\tfrac12(L_v-R_v).
\]
Then
\[
 T(v)^*T(w)=\langle v,w\rangle Q_4,\quad
 [J_v,Q_4]=S(v),\quad
 \tfrac18\Tr(S(v)^*S(w))=\tfrac14\langle v,w\rangle.
\]
For positive \(H=\sum_i h_iQ_i\), the positive square root of the
superoperator \(X\mapsto HXH\) is \(X\mapsto H^{1/2}XH^{1/2}\).
Its \(i,j\) block weight is \(\sqrt{h_ih_j}\), and
\([J_v,H]=(h_4-h_5)S(v)\).
\end{proposition}
\begin{proof}
Quaternion multiplication preserves the coefficient norm and
\(vu=-uv\) for \(v\perp u\). Apply each map to \(u\) and to its
orthogonal complement. The trace of the two mutually adjoint edge
products is \(2\langle v,w\rangle\).
The orthogonal spaces \(\Hom(\im Q_j,\im Q_i)\) diagonalize the
positive superoperator with eigenvalues \(h_ih_j\).
\end{proof}

For a supplied depth operator
\(\mathsf Q=0Q_0+4Q_4+5Q_5+7Q_7\) and \(\alpha>1\), the
positive functional calculus gives \(H=E_0\alpha^{-\mathsf Q/4}\).
Off-diagonal exponents are arithmetic means of the endpoint depths.
On the even endomorphism algebra the labels \(4(q_i+q_j)\) are
\(32,36,40,44,48,56\), with multiplicities \(1,4,4,2,4,1\).
This is a consequence of the declared depth assignment.

\begin{proposition}[Common-rate trace coordinate]
\label{p1b:prop:trace-coordinate}
Let \(\R^8=\R1\oplus W\), \(\dim W=7\), and suppose a
self-adjoint semigroup fixes \(1\), commutes with the full orthogonal
group of \(W\), and has a single positive decay rate normalized to one.
It is \(T_\tau=P_{\R1}+e^{-\tau}P_W\), and
\[
 q=8-\Tr T_\tau=7(1-e^{-\tau}),\qquad 0\le q<7.
\]
A continuous positive multiplicative scalar response of this
additive coordinate has form \(h(q)=e^{-Aq}\).
For \(\alpha=e^A\), its unique positive fourth root is
\(\alpha^{-q/4}\).
\end{proposition}
\begin{proof}
The commutant of the full orthogonal action on \(W\) consists of
scalars. The semigroup law and continuity give an exponential.
Taking the trace gives \(q\). For the scalar response,
\(\log h(q_1+q_2)=\log h(q_1)+\log h(q_2)\) wherever the sum is
defined; continuity solves the additive equation.
\end{proof}
An ordered sequence of rank-one transitions requires an additional
flag; the trace coordinate alone has no preferred ordering.

\section{Metrics, connections, complex structures, and spectra}
\label{p1b:sec:organized-analytic}

\subsection{Complete invariant metric and connection families}
\label{p1b:sec:geometric-families}

We now treat variable geometric structures in the same way as a finite
coordinate-completion family. A member is defined by its parameters; the
classification identifies every member with the stated invariance.
Write
\[
 \mathfrak g=\mathfrak s_x\oplus\mathfrak s_p\oplus\mathfrak z,
 \qquad \mathfrak s_i\cong\mathfrak{su}(2),\quad\mathfrak z\cong\R^2.
\]
For matched oriented Hopf circles put
\(\mathfrak s_i=U_i\oplus\R H_i\), where \(U_i\) is the rotating
two-plane. Set
\(\mathfrak t=\R H_x\oplus\R H_p\oplus\mathfrak z\).
The two Hopf rotations act independently; the scalar circles are central.

\begin{theorem}[Metric cones and their exchange-fixed parts]
\label{p1b:thm:metric-cones}
Fix bases for the displayed real spaces. The following are parameter spaces
of metrics, before quotienting by changes of marking:
\[
\begin{array}{c|c|c}
\text{invariance}&\text{parameter space}&\text{dimension}\\ \hline
\text{left translations}&\Sym^+_8&36\\
\text{left translations and }B&\Sym^+_4\times\Sym^+_4&20\\
\text{left translations, right Hopf rotations}
 &(\R_{>0})^2\times\Sym^+_4&12\\
\text{left, right Hopf rotations and }B
 &\R_{>0}\times\Sym^+_2\times\Sym^+_2&7\\
\text{left and right translations}
 &(\R_{>0})^2\times\Sym^+_2&5\\
\text{left and right translations and }B
 &\R_{>0}\times\R_{>0}\times\R_{>0}&3
\end{array}
\]
Here \(\Sym^+_d\) is the cone of positive definite real symmetric
\(d\)-by-\(d\) matrices. These are exhaustive descriptions of the
indicated classes. In particular they are not assertions that every
left-invariant metric is bi-invariant.
\end{theorem}
\begin{proof}
A left-invariant metric is precisely its positive inner product at the
identity. Exchange on \(\mathfrak g\) has eigenspaces of dimensions four
and four. Invariance makes these orthogonal, giving the second row.
For the Hopf action the two rotating planes \(U_x,U_p\) are inequivalent
irreducible real representations and \(\mathfrak t\) is a four-dimensional
trivial representation. An invariant symmetric form is a positive scalar
on each rotating plane and an arbitrary positive form on \(\mathfrak t\).
There are no cross terms. Exchange equates the two scalar coefficients;
on \(\mathfrak t\) its eigenspaces both have dimension two. This proves
the third and fourth rows. For the full adjoint action each simple ideal
has its scalar Killing multiple, with no mixed terms, and the center
has an arbitrary positive form. In the exchange-fixed case the simple
coefficients agree and the central form is diagonal in the two exchange
eigenlines. Positivity gives all displayed cones and dimensions.
\end{proof}

\begin{theorem}[All left-invariant metric connections]
\label{p1b:thm:left-metric-connections}
Fix any left-invariant metric \(g\). Every left-invariant metric connection
is uniquely
\[
 \nabla_XY=\nabla^{\rm LC}_XY+K(X)Y,\qquad
 K\in\Hom_{\R}(\mathfrak g,\mathfrak{so}(\mathfrak g,g)).
\]
This is an affine space of dimension \(8\binom82=224\).
Its connections with totally alternating metric torsion form the
affine subspace
\[
 g(\nabla_XY,Z)=g(\nabla^{\rm LC}_XY,Z)+\tfrac12 T(X,Y,Z),
 \qquad T\in\Lambda^3\mathfrak g^*,
\]
of dimension \(56\). If \(g\) is exchange invariant, the
exchange-invariant connections have dimension \(112\), and their
alternating-torsion subfamily has dimension \(28\).
\end{theorem}
\begin{proof}
The difference of two connections is a tensor; metric compatibility
makes its last two indices skew. Left invariance identifies it with
the displayed linear map. The Levi-Civita term is fixed by
\[
 2g(\nabla^{\rm LC}_XY,Z)
 =g([X,Y],Z)-g([Y,Z],X)+g([Z,X],Y).
\]
If metric torsion is a three-form, the torsion equation and metric
compatibility solve uniquely for \(K=T/2\).
Under exchange, \(\mathfrak g=V_+\oplus V_-\) with dimensions \(4,4\).
The even and odd dimensions of \(\Lambda^2\mathfrak g\) are \(12,16\).
Thus the invariant Hom space has dimension \(4\cdot12+4\cdot16=112\).
The even part of \(\Lambda^3\mathfrak g^*\) is
\(\Lambda^3V_+^*\oplus(V_+^*\otimes\Lambda^2V_-^*)\), of dimension
\(4+4\cdot6=28\).
\end{proof}

\begin{theorem}[All bi-invariant metric connections]
\label{p1b:thm:bi-metric-connections}
Fix a bi-invariant metric and an orientation on its central plane.
Let \(J_{\mathfrak z}\) be the oriented orthogonal complex structure on
that plane. Every connection invariant under both left and right
translations and compatible with the metric is uniquely given, on
left-invariant fields \(X=(X_x,X_p,z)\), \(Y=(Y_x,Y_p,w)\), by
\begin{align*}
 (\nabla_XY)_i&=s_i[X_i,Y_i]+\ell_i(w)X_i,\qquad i=x,p,\\
 (\nabla_XY)_{\mathfrak z}
 &=\kappa(z)J_{\mathfrak z}w
     -\sum_{i=x,p}g(X_i,Y_i)\ell_i^\sharp ,
\end{align*}
where \(s_x,s_p\in\R\) and
\(\ell_x,\ell_p,\kappa\in\mathfrak z^*\).
The family has dimension eight. Its totally alternating-torsion part
is exactly \(\ell_x=\ell_p=\kappa=0\), with
\[
 T(X,Y)=\sum_i(2s_i-1)[X_i,Y_i],\qquad
 R(X,Y)Z=\sum_i s_i(s_i-1)[[X_i,Y_i],Z_i].
\]
For an exchange-invariant metric, write \(S=B|_{\mathfrak z}\).
The full exchange-invariant connection family has dimension four and is
defined by
\[
 s_x=s_p,\qquad \ell_p=\ell_x\circ S,\qquad
 \kappa\circ S=-\kappa.
\]
Its alternating-torsion part is the one-parameter family \(s_x=s_p\).
\end{theorem}
\begin{proof}
In a left frame, right invariance is adjoint equivariance.
Under the independent adjoint rotations,
\(\Lambda^2\mathfrak g\) contains one copy of \(\mathfrak s_i\)
from \(\Lambda^2\mathfrak s_i\), two copies from
\(\mathfrak s_i\otimes\mathfrak z\), and one trivial copy from
\(\Lambda^2\mathfrak z\). The mixed
\(\mathfrak s_x\otimes\mathfrak s_p\) has no matching input in
\(\mathfrak g\). Consequently
\(\Hom_{\Ad}(\mathfrak g,\Lambda^2\mathfrak g)\) has dimension
\(3+3+2=8\). The displayed maps are independent equivariant maps
and are skew in their last two metric indices, so they exhaust this space.
For \(X\in\mathfrak s_i\), \(z\in\mathfrak z\),
\(g(T(X,z),X)=\ell_i(z)g(X,X)\); alternating torsion forces
\(\ell_i=0\). On the central two-plane an alternating three-form
vanishes, forcing \(\kappa=0\). The remaining bracket terms have
alternating torsion; Jacobi gives the curvature formula.
Exchange swaps the simple ideals and is an orientation-reversing
orthogonal map on \(\mathfrak z\), so
\(SJ_{\mathfrak z}=-J_{\mathfrak z}S\).
Substitution proves exactly the three exchange conditions.
\end{proof}

\begin{theorem}[The full principal-connection family]
\label{p1b:thm:all-principal-connections}
For either trivial principal bundle
\(M_8\to K\) with group \(K\), or \(X_7\to K\) with group
\(\SU(2)\), fix its stated global section and let \(G\) denote its
structure group. Every principal connection is represented uniquely in
that section by
\[
 A\in\Omega^1(K;\mathfrak g_G),\qquad F_A=\dd A+A\wedge A.
\]
Changing the section by \(g:K\to G\) gives
\[
 A^g=g^{-1}Ag+g^{-1}\dd g,\qquad F_{A^g}=g^{-1}F_Ag.
\]
Equivalently, on the total space the connections form the affine space
\(\omega_0+\Omega^1_{\rm hor,Ad}(P;\mathfrak g_G)\).
If \(\omega_0\) is exchange invariant, all exchange-invariant
connections are exactly \(\omega_0+a\) with \(B^*a=a\).
Every connection has a unique even and odd difference from \(\omega_0\).
\end{theorem}
\begin{proof}
On the trivialization \(P=K\times G\), the connection reconstructed from
\(A\) is
\(\omega=\Ad(h^{-1})A+h^{-1}\dd h\).
It has the requisite vertical normalization and equivariance, and
pullback to \(h=e\) recovers \(A\). This also proves the section-change
formulas. A difference of connections is horizontal and adjoint equivariant,
and any such difference can be added to a connection.
Exchange is a principal bundle automorphism covering \(r\mapsto r^{-1}\);
it preserves this vector space of differences and squares to one.
The projections \((a+B^*a)/2,(a-B^*a)/2\) prove the final assertions.
\end{proof}

For the \(\SU(2)\) bundle, Theorem~\ref{p1b:thm:homogeneous}
classifies the entire homogeneous subfamily. The scalar deformations
\(\mathscr A_u\) of Theorem~\ref{p1b:thm:deformation} define a
useful function-parametrized slice of the full affine family.


\subsubsection{Finite operation calculus for invariant metric connections}
\label{p1b:sec:g19}

Theorem~\ref{p1b:thm:left-metric-connections} classifies the full
\(224\)-dimensional affine family.  The statements below close the basic
operations on that entire family, rather than only on the special
bi-invariant one-parameter torsion branch.

\begin{proposition}[Torsion--contorsion isomorphism]
\label{p1b:prop:g19-torsion}
Fix a left-invariant metric \(g\) and write
\(\nabla=\nabla^{\rm LC}+K\) as in
Theorem~\ref{p1b:thm:left-metric-connections}.  In a left-invariant frame,
put
\[
 K_{abc}=g(K_{e_a}e_b,e_c),\qquad
 T_{abc}=g(T(e_a,e_b),e_c).
\]
Then
\[
 T_{abc}=K_{abc}-K_{bac},\qquad
 \boxed{K_{abc}=\tfrac12(T_{abc}-T_{bca}+T_{cab}).}
\]
Hence torsion is a linear isomorphism
\[
 \mathfrak g^*\otimes\mathfrak{so}(\mathfrak g,g)
 \cong \Lambda^2\mathfrak g^*\otimes\mathfrak g.
\]
Under the orthogonal group, the \(224\)-dimensional torsion space splits
into its vectorial, alternating-three-form and traceless Cartan parts with
dimensions
\[
 \boxed{224=8+56+160.}
\]
If \(g\) is exchange invariant, the corresponding exchange-even dimensions
are
\[
 \boxed{112=4+28+80.}
\]
\end{proposition}
\begin{proof}
Metric compatibility gives \(K_{abc}=-K_{acb}\), while the
Levi--Civita torsion vanishes.  Thus
\(T_{abc}=K_{abc}+K_{bca}\).  Cyclically permuting this identity and taking
\(T_{abc}-T_{bca}+T_{cab}\) gives \(2K_{abc}\), proving invertibility.
The standard orthogonal decomposition is obtained from the trace part, the
full alternation and their common traceless complement.  Its dimensions are
\(8\), \(\binom83=56\), and \(224-8-56=160\).  Under exchange the full
fixed torsion space has dimension \(112\) by
Theorem~\ref{p1b:thm:left-metric-connections}; the fixed covector space has
dimension four and the fixed three-form space dimension twenty-eight, so
the remaining fixed Cartan dimension is eighty.
\end{proof}

\begin{theorem}[Invariant curvature calculus]
\label{p1b:thm:g19-curvature}
Let \([e_a,e_b]=c_{ab}{}^m e_m\) and write
\[
 \nabla_{e_a}e_b=A_{ab}{}^c e_c,
 \qquad
 (A_a)^c{}_{b}=A_{ab}{}^c.
\]
For every left-invariant metric connection,
\[
 \boxed{R_{ab}=[A_a,A_b]-c_{ab}{}^mA_m.}
\]
Equivalently,
\[
 R_{abc}{}^d
 =A_{bc}{}^mA_{am}{}^d
  -A_{ac}{}^mA_{bm}{}^d
  -c_{ab}{}^mA_{mc}{}^d.
\]
All covariant derivatives \(\nabla^rT\) and \(\nabla^rR\) have constant
left-frame components obtained recursively by finite tensor contraction
with the matrices \(A_a\) and the structure constants.  In particular the
infinitesimal holonomy algebra is determined by finite-dimensional linear
algebra: it is the Lie algebra generated by the curvature endomorphisms and
their iterated covariant derivatives at the identity.
\end{theorem}
\begin{proof}
The coefficients \(A_{ab}{}^c\) are constant in a left-invariant frame.
Expanding
\(R(e_a,e_b)e_c=\nabla_a\nabla_b e_c-\nabla_b\nabla_a e_c-
\nabla_{[e_a,e_b]}e_c\) gives the component formula and hence the matrix
identity.  Covariant differentiation of an invariant tensor differentiates
no scalar coefficient; only the connection action on its tensor indices and
the fixed frame brackets remain.  Repetition therefore gives a finite
algebraic recursion.  The connection and metric are real analytic in the
left-invariant frame, so the infinitesimal holonomy theorem identifies the
holonomy algebra with the Lie closure of curvature and its iterated
covariant derivatives.
\end{proof}

\begin{corollary}[Geodesic invariance under alternating torsion]
\label{p1b:cor:g19-geodesics}
Every metric connection whose torsion is a three-form has exactly the same
affinely parametrized geodesics as the Levi--Civita connection.
\end{corollary}
\begin{proof}
For alternating torsion, Proposition~\ref{p1b:prop:g19-torsion} gives
\(K_XY=\tfrac12T(X,Y,\cdot)^\sharp\), hence \(K_XX=0\).  Therefore
\(\nabla_{\dot\gamma}\dot\gamma=0\) is identical to the
Levi--Civita geodesic equation.
\end{proof}

\begin{theorem}[Finite flatness and full-holonomy criteria for every invariant connection]
\label{p1b:thm:v214-finite-full-holonomy}
Fix the metric and a global oriented orthonormal left frame.  Let
\(A_X\in\mathfrak{so}(8)\) be the connection matrix, linear in
\(X\in\mathfrak g\), and set
\[
 R_{ab}=[A_{e_a},A_{e_b}]-A_{[e_a,e_b]},\qquad
 \mathfrak h_0=\operatorname{span}_{\R}\{R_{ab}\}.
\]
Recursively form the real vector spaces
\[
 \mathfrak h_{r+1}=\mathfrak h_r+
 \operatorname{span}[\mathfrak h_r,\mathfrak h_r]+
 \sum_{a=1}^8[A_{e_a},\mathfrak h_r].
\]
This sequence stabilizes after at most twenty-eight strict enlargements.
Its stable space \(\mathfrak h\) is the restricted holonomy Lie algebra.
Consequently \(R_{ab}=0\) for all \(a,b\) is a necessary and sufficient
finite quadratic criterion for flatness on the entire
\(224\)-parameter family.  The displayed rank and commutator recursion
likewise determines every holonomy-algebra stratum; it does not require
an unbounded search through covariant derivatives.

Let \(Z_t,Z_E\) be the central coordinate generators of period
\(2\pi\), and let \(H_{\rm res}\subset SO(8)\) be the connected analytic
subgroup with Lie algebra \(\mathfrak h\).  Put
\[
 M_j=\exp(-2\pi A_{Z_j}),\qquad j=t,E.
\]
Then the full tangent holonomy is exactly
\[
 \boxed{\operatorname{Hol}(\nabla)
 =\bigcup_{(m,n)\in\Z^2}H_{\rm res}M_t^mM_E^n.}
\]
The two matrices normalize \(H_{\rm res}\) and commute modulo it.
Thus their cosets give the complete monodromy homomorphism
\(\Z^2\to\operatorname{Hol}/H_{\rm res}\); its image, equivalently
\(\Z^2\) modulo its kernel, gives all holonomy components in the
natural holonomy Lie-group structure.  No closure of a possibly
nonclosed holonomy subgroup is implicit in this formula.

For a spin structure with central transition signs
\(\epsilon_j=(-1)^{2\sigma_j}\), replace \(A_X\) by
\[
 \widehat A_X=\frac14\sum_{b,c}
       g(A_Xe_b,e_c)c(e_b)c(e_c).
\]
The same construction in \(\Spin(8)\), with the connected lifted
restricted group and generators
\(\epsilon_j\exp(-2\pi\widehat A_{Z_j})\), gives its full spin
holonomy.  Parallel tensors or spinors are precisely the vectors killed
by the corresponding infinitesimal action of \(\mathfrak h\) and fixed
by the two monodromy matrices.
\end{theorem}
\begin{proof}
The curvature identity is Theorem~\ref{p1b:thm:g19-curvature}.
Every term in the recursion belongs to \(\mathfrak{so}(8)\), of
dimension twenty-eight.  Once two successive spaces agree, that space
is a Lie algebra stable under all \(\operatorname{ad}A_X\), so the
recursion remains stationary.  This proves finite termination.

For an invariant curvature tensor,
\[
 (\nabla_XR)(Y,Z)=[A_X,R(Y,Z)]
                  -R(A_XY,Z)-R(Y,A_XZ)
\]
in the left frame.  The last two terms are linear combinations of the
curvature matrices.  More generally, when taking any successive
covariant derivative, the connection action on the argument indices
only gives combinations of derivative matrices already present; the
endomorphism index gives commutators with \(A_X\).  Induction in both
directions therefore identifies the Lie closure of all curvature
covariant derivatives with the smallest Lie algebra containing the
\(R_{ab}\) and stable under \(\operatorname{ad}A_X\), exactly the
stable recursion above.  The analytic infinitesimal-holonomy theorem
used in Theorem~\ref{p1b:thm:g19-curvature} identifies this algebra
with restricted holonomy.  A connection is flat exactly when its
curvature vanishes, proving the quadratic criterion, including all
zero and lower-rank strata.

Along the scalar coordinate loop \(\exp(tZ_j)\), \(0\le t\le2\pi\),
parallel transport in the global left frame solves
\(\dot P=-A_{Z_j}P\), giving \(M_j\).  These loops generate
\(\pi_1(M_8)=\Z^2\).  Any loop differs from the corresponding word
in them by a contractible loop, whose holonomy belongs to
\(H_{\rm res}\).  Conversely every displayed product is holonomy of
such a concatenation.  Restricted holonomy is normal, and the scalar
loop classes commute, proving the normalization, coset and component
claims.  The spin connection formula is the Lie-algebra lift of the
orthogonal connection.  Lifting a scalar loop multiplies its lifted
transport by its prescribed spin transition sign.  Finally, a parallel
section is uniquely determined by its value at a base point, and that
value extends consistently exactly when all holonomy fixes it.  A
connected group fixes a vector exactly when its Lie algebra kills it.
\end{proof}

\begin{theorem}[Complete flat strata of the bi-invariant metric-connection family]
\label{p1b:thm:v214-flat-bi-connections}
In Theorem~\ref{p1b:thm:bi-metric-connections}, write
\(v_i=\ell_i^\sharp\), and normalize orthonormal sphere generators
by \([e_{i,a},e_{i,b}]=c_i\epsilon_{abc}e_{i,c}\),
\(c_i=2/R_i\).  Flatness is equivalent to
\[
 \boxed{(2s_i-1)v_i=0,\qquad
 s_i(s_i-1)+\frac{|v_i|^2}{c_i^2}=0\quad(i=x,p),}
\]
\[
 \boxed{\langle v_x,v_p\rangle=0,\qquad
 \kappa=0\ \text{if }v_x\ne0\text{ or }v_p\ne0.}
\]
In particular each sphere independently has precisely these alternatives:
\[
 \begin{array}{c|c}
 v_i=0&s_i=0\ \text{or }1\\
 v_i\ne0&s_i=\tfrac12,\quad |v_i|=1/R_i .
 \end{array}
\]
The two nonzero central vectors must be orthogonal.  When both vanish,
\(\kappa\in\mathfrak z^*\) is arbitrary; these statements exhaust
all eight parameters and every boundary case.

All these connections have trivial restricted holonomy.  If at least
one \(v_i\ne0\), their full tangent holonomy is trivial.  If both
vanish, it is generated by the two rotations
\[
 \exp[-2\pi\kappa(Z_j)J_{\mathfrak z}],\qquad j=t,E,
\]
acting as the identity on the sphere directions.  A global parallel
oriented orthonormal tangent frame exists exactly when these rotations
are both the identity.  In the latter branch this is
\(\kappa(Z_t),\kappa(Z_E)\in\Z\).

In an oriented orthonormal central basis \(z_1,z_2\), the spin
monodromy in that branch is
\[
 \epsilon_j\exp[-\pi\kappa(Z_j)c(z_1)c(z_2)].
\]
Thus a global parallel spin frame exists exactly when
\(\kappa(Z_j)\in\Z\) and
\(\epsilon_j=(-1)^{\kappa(Z_j)}\) for both circles.
In the branch with a nonzero \(v_i\), the corresponding condition is
\(\epsilon_t=\epsilon_E=1\).
\end{theorem}
\begin{proof}
For \(X\in\mathfrak s_i\), the matrix \(A_X\) acts by
\(s_i\operatorname{ad}_X\) on its sphere block, maps
\(w\in\mathfrak z\) to \(\ell_i(w)X\), and maps
\(Y\in\mathfrak s_i\) to \(-g(X,Y)v_i\) in the center.
For \(z\in\mathfrak z\), \(A_z\) is
\(\kappa(z)J_{\mathfrak z}\) on the center and zero elsewhere.
Insert these blocks in \(R(X,Y)=[A_X,A_Y]-A_{[X,Y]}\).
On a single sphere block the result is
\[
 R(X,Y)|_{\mathfrak s_i}
 =\left(s_i(s_i-1)+|v_i|^2/c_i^2\right)
                     \operatorname{ad}_{[X,Y]}.
\]
On a central input its off-diagonal block is
\(R(X,Y)w=(2s_i-1)\ell_i(w)[X,Y]\); the reverse block
is determined by skew-adjointness.  These formulas follow also from
\((X\times Y)\times U=Y\langle X,U\rangle-X\langle Y,U\rangle\).
For \(X\in\mathfrak s_x,Y\in\mathfrak s_p\), the mixed curvature
is
\[
 \langle v_x,v_p\rangle
       (|Y\rangle\langle X|-|X\rangle\langle Y|).
\]
Finally \(R(X,z)=[A_X,A_z]\) vanishes for every \(X,z\) exactly
when \(\kappa(z)v_i=0\) for every \(z,i\), since
\(J_{\mathfrak z}\) is invertible; \(R(z,w)=0\) always.
The independent matrix blocks prove necessity and sufficiency of all
four displayed conditions.  Solving the two scalar sphere conditions
gives exactly the stated alternatives.

For flat connections Theorem~\ref{p1b:thm:v214-finite-full-holonomy}
reduces full holonomy to the two central loop transports.
Here \(A_{Z_j}\) is precisely the central rotation matrix just
written.  In the nonzero-vector branches it is zero.  A flat tangent
bundle has a global parallel frame precisely when its full holonomy
is trivial, proving the tangent assertions.  The Lie-algebra spin
lift of \(J_{\mathfrak z}\) is
\(\tfrac12c(z_1)c(z_2)\); its exponential at an integral full
rotation is the scalar \((-1)^{\kappa(Z_j)}\).  This proves all spin
conditions, including the nonperiodic spin structures.
\end{proof}

\begin{theorem}[Alternating-torsion spin Dirac spectra at every parameter]
\label{p1b:thm:v214-torsion-dirac}
Use the alternating-torsion subfamily
\(\nabla_XY=\sum_i s_i[X_i,Y_i]\), with the flat central
connection, and let \(D^{\boldsymbol s}\) be its spin Dirac operator.
Then
\[
 D^{\boldsymbol s}=D^{\rm LC}+\frac34c(T^{\boldsymbol s}),\qquad
 T^{\boldsymbol s}=\sum_i(2s_i-1)H_{g,i}.
\]
On one sphere use the Clifford convention of Section~\ref{p1b:der:g08:s61}:
\(c_a=-i\sigma_a\) in its oriented right-invariant frame, so
\(c(e_1)c(e_2)c(e_3)=-I\).  The complete signed sphere values are
\[
 \boxed{\lambda_{\ell,+}^{s}=\frac{\ell+3-3s}{R},\qquad
 \lambda_{\ell,-}^{s}=-\frac{\ell+3s}{R},\qquad
 \ell\ge0,}
\]
each with labelled multiplicity
\(d_\ell=(\ell+1)(\ell+2)\).  These are labels continuing the two
Levi--Civita branches, not assertions that their values retain their
sign for every real \(s\).  Equal numerical values have additive
multiplicities.  In the opposite volume-Clifford convention both
branches are obtained instead by adding
\(3(2s-1)/(2R)\) to the corresponding Levi--Civita values.

For arbitrary \(s\in\R\), put
\[
 h(s)=\begin{cases}
 (r+1)(r+2),&s=-r/3\text{ or }s=1+r/3,
                      \quad r\in\Z_{\ge0},\\
 0,&\text{otherwise}.
 \end{cases}
\]
This is the complete one-sphere kernel dimension.  In particular
\(0<s<1\) gives no kernel, while \(s=0,1\) each gives two zero
modes.

For the eight-dimensional product, choose either labelled sphere
value in each factor and a scalar Fourier label \(n\in\Z^2\).
Set
\[
 E^2=(\lambda_{\ell_x,\delta_x}^{s_x})^2+
      (\lambda_{\ell_p,\delta_p}^{s_p})^2+
      (n+\sigma)^TG^{-1}(n+\sigma).
\]
If \(E>0\), the product values are \(\pm E\), each with
multiplicity \(2d_{\ell_x}d_{\ell_p}\) for this full label.
At \(E=0\), that label contributes
\(4d_{\ell_x}d_{\ell_p}\) zero modes.  Consequently
\[
 \boxed{\dim\ker D^{\boldsymbol s}=
 \begin{cases}
 4h(s_x)h(s_p),&\sigma=(0,0),\\
 0,&\sigma\ne(0,0).
 \end{cases}}
\]
It splits equally between product chiralities.  Thus either choice of
endpoint connection on each sphere has sixteen zero modes for the
periodic torus spin structure.  The positive Levi--Civita product gap
is specifically the midpoint \(s_x=s_p=\tfrac12\), not a
connection-independent lower bound.
\end{theorem}
\begin{proof}
For a metric connection with three-form torsion, the connection
difference on spinors is
\(\tfrac18\sum_{b,c}T(X,e_b,e_c)c(e_b)c(e_c)\).
Clifford contraction in \(X\) gives
\(D^{\boldsymbol s}-D^{\rm LC}=\tfrac34c(T^{\boldsymbol s})\),
using \(c(T)=\tfrac16\sum T_{abc}c_ac_bc_c\).
On the radius-\(R\) sphere the Cartan form is
\(H_g=(2/R)e^{123}\).  In the stated right-frame convention its
Clifford action is \(-2I/R\).  Hence the torsion shift is
\(-3(2s-1)/(2R)\).  Adding it to the two complete round
Levi--Civita spectra \(\pm(\ell+3/2)/R\) proves the two formulas
and their multiplicities.  Their zeros occur exactly at the two
arithmetic families displayed in \(h(s)\), which are disjoint.
The convention-independent description is the same calculation with
\(c(e^{123})=+I\).

On the product, use the four auxiliary anticommuting matrices of
(G08.7).  The spin Dirac operator of the product connection is
\[
 D^{\boldsymbol s}=D_x^{s_x}\otimes\alpha_1+
          D_p^{s_p}\otimes\alpha_2+
          \kappa_1\beta_1+\kappa_2\beta_2,
 \qquad \kappa=G^{-1/2}(n+\sigma).
\]
Each connection term acts in its own factor, so the square is the
sum of the three displayed nonnegative squares.  On fixed sphere
branches its auxiliary four-by-four matrix is traceless Clifford
multiplication by a vector of length \(E\): its eigenvalues are
\(\pm E\), each twice, or it is zero with multiplicity four.
Completeness follows by tensoring the complete shifted sphere and
Fourier bases.  A zero requires both sphere values and the torus
frequency to vanish; the last is possible exactly for
\(\sigma=(0,0),n=0\).  This proves the entire kernel formula.
The auxiliary chirality has two eigenvectors of each chirality,
proving the equal splitting and the endpoint count.
\end{proof}


\subsubsection{Reciprocal double cone and hyper-K\"ahler ambient geometry}
\label{p1b:sec:g18}

The compact manifold \(M_8\) is not itself a cone.  The vector carrier
nevertheless has a canonical noncompact cone host, and the distinction is
useful enough to state explicitly.

\begin{theorem}[Reciprocal double-cone host]
\label{p1b:thm:g18-double-cone}
Let
\[
 \mathcal C^\times
 =C_x^\times\times C_p^\times
 =(0,\infty)\times S_x^3\times(0,\infty)\times S_p^3
 \cong(\HH\setminus\{0\})^2
\]
with
\[
 G_{\mathcal C}
 =\dd r_x^2+r_x^2g_{S^3,1;x}
  +\dd r_p^2+r_p^2g_{S^3,1;p}.
\]
Then \((\mathcal C^\times,G_{\mathcal C})\) is flat.  If
\(I_1,I_2,I_3\) are the diagonal left-quaternionic complex structures on
\(\HH\oplus\HH\), then
\[
 \boxed{(\mathcal C^\times,G_{\mathcal C},I_1,I_2,I_3)
 \text{ is hyper-K\"ahler}.}
\]
Its metric completion is Euclidean quaternionic two-space,
\[
 \boxed{\overline{\mathcal C}=\HH^2.}
\]
\end{theorem}
\begin{proof}
Polar coordinates identify
\(\dd r^2+r^2g_{S^3,1}\) with the Euclidean metric on
\(\HH\setminus\{0\}\).  Taking the product gives the Euclidean metric on
\((\HH\setminus\{0\})^2\).  Left multiplication by the three imaginary
quaternionic units is orthogonal, parallel, and satisfies the quaternion
relations on each factor, hence diagonally on the product.  Adding the
missing radial loci gives the Euclidean completion \(\HH^2\).
\end{proof}

\begin{proposition}[Born extension and fixed diagonal]
\label{p1b:prop:g18-born}
For matched quaternionic markings the exchange
\[
 \mathcal B_{\mathcal C}(q_x,q_p)=(q_p,q_x)
\]
is an involutive triholomorphic isometry of the completed cone and of the
punctured double cone.  Its fixed set in \(\HH^2\) is the quaternionic
linear diagonal
\[
 \Delta_{\HH}=\{(q,q):q\in\HH\}\cong\HH,
\]
and its fixed set in \(\mathcal C^\times\) is
\(\Delta_{\HH}\setminus\{0\}\cong\HH^\times\).
\end{proposition}
\begin{proof}
Factor exchange preserves the product Euclidean metric and commutes with
simultaneous left multiplication by each imaginary quaternionic unit.
The fixed-point equation is \(q_x=q_p\).
\end{proof}

\begin{corollary}[Fixed-radius leaf and type firewall]
\label{p1b:cor:g18-firewall}
For \(R_x,R_p>0\),
\[
 P_v(R_x,R_p)=S_x^3(R_x)\times S_p^3(R_p)
 =\{r_x=R_x,\ r_p=R_p\}\subset\mathcal C^\times
\]
is a codimension-two compact leaf with the round product metric.  The
full reciprocal compact carrier is instead
\[
 \boxed{M_8=P_v(R_x,R_p)\times\T^2_{\rm sc}.}
\]
Thus \(\mathcal C^\times\) and \(M_8\) have the same dimension but are
not the same geometric object: the former is noncompact and has no scalar
torus, whereas the latter is compact with \(\pi_1(M_8)\cong\Z^2\).
A scalar-preserving ambient host for the full compact carrier is the
ten-dimensional product
\(\mathcal C^\times\times\T^2_{\rm sc}\), in which \(M_8\) is again a
codimension-two fixed-radius leaf.
\end{corollary}

The independent radial dilations and the second fundamental geometry of
the fixed-radius leaves are recorded in Paper~I-D, Section~\ref{p1b:app:g18}.


\subsubsection{Isometry groups and exact scaling weights}
\label{p1b:sec:g24}

The radius and flat-torus parameters already appear throughout the paper.
The following result places their symmetry enhancement and scaling laws in
one atlas and separates group automorphisms from metric isometries.

\begin{theorem}[Full isometry group and Born enhancement]
\label{p1b:thm:g24-isometry}
Let
\[
 g=R_x^2g_{S^3,1;x}+R_p^2g_{S^3,1;p}+G_{\rm sc}
\]
with \(G_{\rm sc}\) a flat metric on the marked scalar torus.  Put
\[
 O(G_{\rm sc};\Z)=\{A\in GL(2,\Z):A^TG_{\rm sc}A=G_{\rm sc}\}.
\]
If \(R_x\ne R_p\),
\[
 \operatorname{Isom}(M_8,g)
 \cong O(4)_x\times O(4)_p\times
 \bigl(\T^2\rtimes O(G_{\rm sc};\Z)\bigr).
\]
If \(R_x=R_p\), the two round three-sphere factors may additionally be
interchanged, giving
\[
 \operatorname{Isom}(M_8,g)
 \cong (O(4)^2\rtimes C_2)\times
 \bigl(\T^2\rtimes O(G_{\rm sc};\Z)\bigr).
\]
In either case the identity component is
\[
 SO(4)_x\times SO(4)_p\times\T^2,
\]
of dimension fourteen, in agreement with the Killing-algebra result of the
round product atlas.

Let \(S=\begin{psmallmatrix}0&1\\1&0\end{psmallmatrix}\) be the scalar-circle
swap in the marked basis.  The paper's Born exchange is a metric isometry
if and only if
\[
 \boxed{R_x=R_p,\qquad S^TG_{\rm sc}S=G_{\rm sc}.}
\]
This is a metric condition, not a statement about the full Lie-group
automorphism group.
\end{theorem}
\begin{proof}
The universal cover is the Riemannian product
\(S^3(R_x)\times S^3(R_p)\times\R^2\).  The de Rham factors are preserved
by every isometry, except that isometric sphere factors may be permuted.
The full isometry group of a round three-sphere is \(O(4)\); the full
isometry group of a flat torus is its translation torus semidirect its
integral orthogonal point group.  This proves the two cases.  The Born
criterion simply requires that both the sphere swap and the marked torus
swap preserve the corresponding metrics.
\end{proof}

\begin{theorem}[Global homothety weights in dimension eight]
\label{p1b:thm:g24-scaling}
Under \(g\mapsto g_\lambda=\lambda^2g\), \(\lambda>0\), the basic weights
are
\[
\begin{array}{c|c}
\text{quantity}&\text{weight}\\\hline
\text{length, diameter, injectivity radius}&\lambda\\
\vol_8&\lambda^8\\
\nabla^{\rm LC}\text{ as an affine connection}&1\\
R^i{}_{jkl}&1\\
\text{sectional curvature and scalar curvature}&\lambda^{-2}\\
\operatorname{Ric}\text{ as a }(0,2)\text{ tensor}&1\\
\Delta\text{ and its eigenvalues}&\lambda^{-2}\\
D\text{ and Dirac eigenvalues}&\lambda^{-1}\\
*|_{\Omega^k}&\lambda^{8-2k}\\
|\alpha|\text{ for a fixed }k\text{-form}&\lambda^{-k}\\
\|\alpha\|_{L^2}^2&\lambda^{8-2k}\\
|F_A|\text{ for a fixed connection curvature}&\lambda^{-2}\\
\int|F_A|^2\vol_8&\lambda^4.
\end{array}
\]
Characteristic classes are unchanged.  If
\(\operatorname{Tr}(e^{-t\Delta_g})\) is a scalar heat trace, then
\[
 \boxed{\operatorname{Tr}(e^{-t\Delta_{g_\lambda}})
 =\operatorname{Tr}(e^{-(t/\lambda^2)\Delta_g}).}
\]
Consequently the integrated heat coefficient of order \(2j\) scales as
\(\lambda^{8-2j}\).
\end{theorem}
\begin{proof}
All entries follow from constant metric rescaling and the defining tensor
variance.  In particular the Christoffel symbols are unchanged, the
\((1,3)\) curvature is unchanged, contraction with the inverse metric
produces the \(\lambda^{-2}\) curvature weights, and Clifford
multiplication contributes one inverse length to the Dirac operator.  The
Hodge-star factor is \(\lambda^{n-2k}\) in dimension \(n=8\).  Functional
calculus of \(\Delta_{g_\lambda}=\lambda^{-2}\Delta_g\) gives the heat
identity and hence the coefficient weights.
\end{proof}

Independent factor rescaling and the distinction between automorphisms and
isometries are expanded in Paper~I-D, Appendix~\ref{p1b:app:g24}.


\begin{corollary}[Relative metric closure and surviving scale moduli]
\label{p1b:cor:v233-relative-scale-boundary}
For the labelled product metric with sphere radii \(R_x,R_p>0\)
and positive central block \(G\), Born exchange is an isometry
of this same metric exactly when
\[
 R_x=R_p,\qquad
 G=\begin{pmatrix}a&b\\b&a\end{pmatrix},\qquad a+b>0,\quad a-b>0.
\]
The common radius and the two central eigenvalues remain free.
Matched exchange in the raw coordinate atlas similarly equates the
two triples of edge magnitudes; it does not fix their common
intrinsic point. That point comes from the specified native
Levi--Civita symbol construction.
\end{corollary}
\begin{proof}
Pullback by exchange swaps the orthogonal sphere coefficients and
the two central diagonal entries. Comparing those entries gives the
metric assertion of Theorem~\ref{p1b:thm:g24-isometry}.
The coordinate assertion is Theorem~\ref{p1b:thm:product-completion}.
Under a global homothety \(g\mapsto\lambda^2g\), the unchanged
Levi--Civita connection gives
\(\Delta\mapsto\lambda^{-2}\Delta\),
\(D\mapsto\lambda^{-1}D\) under the canonical isometric
bundle identification, and
\(\operatorname{vol}\mapsto\lambda^8\operatorname{vol}\),
as in Theorem~\ref{p1b:thm:g24-scaling}.
Hence
\(\operatorname{Tr}e^{-t\Delta_{\lambda^2g}}
=\operatorname{Tr}e^{-(t/\lambda^2)\Delta_g}\).
Topology and an uncalibrated exchange therefore do not fix an
absolute dimensionful standard.
\end{proof}

\subsection{Complex structures and flat equivariant bundles}
\label{p1b:sec:complex-flat}

The root-space construction below is the compact-group construction
usually associated with Samelson and Wang; see \cite{LoebManjarinNicolau}
for its broader setting. Here the classification and integrability are
proved directly for the marked product.

\begin{theorem}[Complete Hopf-invariant complex-structure family]
\label{p1b:thm:complex-family}
Choose matched oriented Hopf axes. All left-invariant almost complex
structures on \(M_8\) invariant under the two independent right Hopf
circles are
\[
 J|_{U_x}=\epsilon_xJ_x,\quad J|_{U_p}=\epsilon_pJ_p,\quad
 J|_{\mathfrak t}=J_0,\qquad
 \epsilon_x,\epsilon_p\in\{\pm1\},\quad J_0^2=-I.
\]
Every member is integrable. Their parameter space is
\[
 \{\pm1\}^2\times GL(4,\R)/GL(2,\C),
\]
of continuous real dimension eight. For a fixed Hopf-invariant metric,
its orthogonal members replace the final factor by \(O(4)/U(2)\).
Exchange-invariant members are exactly
\[
 \epsilon_x=\epsilon_p,\qquad J_0B=BJ_0;
\]
their parameter space is
\(\{\pm1\}\times(GL(2,\R)/GL(1,\C))^2\), of continuous
dimension four. For a fixed exchange-invariant Hopf metric there are
exactly eight orthogonal exchange-invariant members.
\end{theorem}
\begin{proof}
The Hopf-isotypic decomposition is \(U_x\oplus U_p\oplus\mathfrak t\).
Any commuting endomorphism preserves it. On a rotating real two-plane
it has the form \(aI+bJ_i\); its square equals \(-I\) precisely for
\(a=0,b=\pm1\). The trivial summand permits every complex structure
\(J_0\). Such structures form \(GL(4,\R)/GL(2,\C)\).
To prove integrability, complexify the Lie algebra. The \(+i\)
eigenspace consists of a complex two-plane
\(L\subset\mathfrak t_{\C}\), with \(L\oplus\overline L=\mathfrak t_\C\),
and one root line from each \(\mathfrak{sl}_2\) factor. The Cartan
subalgebra is abelian, brackets with \(L\) preserve each chosen root
line, each root line brackets trivially with itself, and the two factors
commute. The eigenspace is thus a Lie subalgebra, exactly the
vanishing-Nijenhuis integrability condition.
Orthogonality on the Cartan space gives \(O(4)/U(2)\).
Exchange interchanges the two rotating planes with their matched
orientations. On \(\mathfrak t\), its real eigenspaces each have
dimension two, so a commuting \(J_0\) is an arbitrary complex
structure on each. For a fixed invariant metric each such two-plane
has two orthogonal complex structures. The two root signs and the
two independent Cartan signs give \(2\cdot2\cdot2=8\).
\end{proof}

\begin{corollary}[Tangent Chern classes and the Kähler obstruction]
\label{p1b:cor:complex-topology}
Every complex structure in Theorem~\ref{p1b:thm:complex-family}
has a smoothly trivial complex tangent bundle and therefore
\(c_j(TM_8,J)=0\) for \(j>0\). None admits a Kähler metric.
\end{corollary}
\begin{proof}
Choose a complex basis of \((\mathfrak g,J)\) and translate it
on the left. This is a global smooth complex frame; holomorphic
triviality is not asserted. By Theorem~\ref{p1b:thm:cohomology},
every degree-two real class is a multiple of \(\eta\), with
\(\eta^2=0\). A symplectic form would have nonzero integral of its
fourth power, which this ring makes impossible. A Kähler form would
be symplectic.
\end{proof}


\subsubsection{Automorphism saturation of the left-invariant complex family}
\label{p1b:sec:g22}

The marked family of Theorem~\ref{p1b:thm:complex-family} was proved
there under right-Hopf invariance.  For the Lie algebra of \(M_8\), that
restriction does not lose any left-invariant integrable complex structure
up to the natural compact-group automorphisms.

\begin{theorem}[Global normal form for left-invariant complex structures]
\label{p1b:thm:g22-global-complex}
Let \(J\) be any left-invariant integrable complex structure on
\[
 M_8=(\SU(2)\times\U(1))_x\times(\SU(2)\times\U(1))_p.
\]
There exists an inner automorphism of the two \(\SU(2)\) factors whose
pullback carries \(J\) into the family of
Theorem~\ref{p1b:thm:complex-family}.  Equivalently, after choosing two
Hopf axes, every left-invariant integrable \(J\) has the normal form
\[
 J|_{U_x}=\epsilon_xJ_x,
 \qquad
 J|_{U_p}=\epsilon_pJ_p,
 \qquad
 J|_{\mathfrak t}=J_0,
 \qquad J_0^2=-I.
\]
Thus the already displayed marked parameter space is an
automorphism-saturating normal-form family for \emph{all} left-invariant
integrable complex structures on \(M_8\).  It is not asserted to be the
orbit space modulo all automorphisms; different marked parameters may be
automorphism equivalent.
\end{theorem}
\begin{proof}
The Lie algebra is
\(\mathfrak u(2)\oplus\mathfrak u(2)\).  Magnin's explicit
classification \cite{Magnin2008ComplexSU2SU2} shows that every integrable
complex structure on this algebra is conjugate by
\((SO(3)\times\R_+^*)^2\) to a normal form \(K(M)\), where the four-dimensional
Cartan--central block is an arbitrary real complex structure
\(M^2=-I\).  The two \(SO(3)\) factors are inner automorphisms of the two
\(\SU(2)\) factors.  The positive rescalings act only in the central
coordinates; conjugating \(M\) by those rescalings merely replaces it by
another arbitrary complex structure on the same four-dimensional block.
That change is already absorbed by the unrestricted parameter \(J_0\) of
Theorem~\ref{p1b:thm:complex-family}.  Reversing the marked Hopf axis
accounts for the two root-plane signs.  The resulting normal form is
therefore exactly the displayed family.
\end{proof}

\begin{corollary}[Transfer of the invariant Hermitian and Dolbeault atlas]
\label{p1b:cor:g22-transfer}
Every left-invariant integrable complex structure on \(M_8\) admits
left-invariant Hermitian metrics.  After pullback by a compact-group
automorphism, the conclusions of Corollary~\ref{p1b:cor:complex-topology}
and the full Dolbeault and Fr\"olicher calculations of
Section~\ref{p1b:sec:g16} apply to every such complex structure.  In
particular every left-invariant complex structure has smoothly trivial
complex tangent bundle and none is K\"ahler.

The flat-Bismut/pluriclosed conclusions of Section~\ref{p1b:sec:g14}
remain exactly the bi-invariant-metric-compatible subfamily stated there;
Theorem~\ref{p1b:thm:g22-global-complex} does not enlarge that metric
hypothesis.
\end{corollary}
\begin{proof}
Any Hermitian inner product on \((\mathfrak g,J)\) extends by left
translation to a left-invariant Hermitian metric.  Theorem~\ref{p1b:thm:g22-global-complex}
reduces an arbitrary invariant complex structure to the marked normal form.
Pullback by a diffeomorphism preserves integrability, Chern classes and
Dolbeault cohomology, so the topological and Dolbeault results transport
back to the original structure.  The final sentence simply retains the
additional bi-invariance and orthogonality assumptions used in G14.
\end{proof}

The precise Lie-algebra normal form used in the reduction is summarized in
Paper~I-D, Appendix~\ref{p1b:app:g22}.


\subsubsection{Hermitian structures with closed Cartan torsion}
\label{p1b:sec:g14}

Choose a bi-invariant metric \(g\) and a \(g\)-orthogonal member \(J_L\)
of Theorem~\ref{p1b:thm:complex-family}. Write
\(\omega_L(X,Y)=g(J_LX,Y)\). In this subsection
\(d^c=i(\bar\partial-\partial)\), and torsion means
\(c(X,Y,Z)=g(T(X,Y),Z)\).

\begin{theorem}[Flat Bismut connection and pluriclosed metric]
\label{p1b:thm:g14-bismut}
The Bismut connection is the connection \(\nabla^0\) making the global
left-invariant frame parallel. With
\[
 H_g(X,Y,Z)=g([X,Y],Z)
\]
on left-invariant fields, it satisfies
\[
 T^0(X,Y)=-[X,Y],\qquad c_L=-H_g,\qquad
 R^{\rm Bismut}=0,\qquad
 \nabla^{\rm Bismut}c_L=0,\qquad \dd c_L=0.
\]
Its holonomy is trivial. The Hermitian form is pluriclosed:
\(\partial\bar\partial\omega_L=0\).
For the metric
\(g=R_x^2g_{S^3_x}+R_p^2g_{S^3_p}+G\), the periods are
\[
 \int_{S^3_x}H_g=4\pi^2R_x^2,\qquad
 \int_{S^3_p}H_g=4\pi^2R_p^2.
\]
In particular this closed torsion is nonexact.
\end{theorem}
\begin{proof}
The metric and complex structure have constant coefficients in the
left frame, so \(\nabla^0\) preserves both. Ad-invariance makes
\(-g([X,Y],Z)\) alternating. The unique Hermitian connection with
alternating torsion is the Bismut connection, proving the identification
as in \cite[Lemma 6]{WangYangZheng2020BismutFlat}. The frame is global and parallel;
its constant bracket coefficients prove flatness, trivial holonomy and
parallel torsion. The Chevalley--Eilenberg differential of \(H_g\)
vanishes by ad-invariance and Jacobi. With the stated convention
\(c_L=-d^c\omega_L\), hence
\(0=\dd c_L=-2i\partial\bar\partial\omega_L\).
On the unit-sphere coframe,
\(H_g=2R_x^2e_x^{123}+2R_p^2e_p^{123}\); integration proves the
periods.
\end{proof}

\begin{theorem}[Twisted generalized Kähler family]
\label{p1b:thm:g14-generalized}
Choose independently a right-invariant orthogonal integrable complex
structure \(J_R\) for the same \(g\), by right-translating any admissible
complex structure at the identity. Put
\[
 J_+=J_R,\quad J_-=J_L,\qquad
 \omega_\pm(X,Y)=g(J_\pm X,Y).
\]
Using the \(H_g\)-twisted bracket of
Theorem~\ref{p1b:thm:dirac-graph}, the endomorphisms of \(TM_8\oplus T^*M_8\)
\[
 \mathcal J_{1,2}=\frac12
 \begin{pmatrix}
  J_+\pm J_-&-(\omega_+^{-1}\mp\omega_-^{-1})\\
  \omega_+\mp\omega_-&-(J_+^*\pm J_-^*)
 \end{pmatrix}
\]
are commuting integrable generalized complex structures.
Here two-forms are maps \(X\mapsto\iota_X\omega\) and \(J^*\)
is the dual map. The generalized metric is
\[
 -\mathcal J_1\mathcal J_2=
       \begin{pmatrix}0&g^{-1}\\g&0\end{pmatrix}.
\]
The Bismut torsions of \(J_+\) and \(J_-\) are respectively \(H_g\)
and \(-H_g\).
\end{theorem}
\begin{proof}
The connections parallelizing the right and left frames are
\(\nabla^{\rm LC}+\tfrac12g^{-1}H_g\) and
\(\nabla^{\rm LC}-\tfrac12g^{-1}H_g\), respectively.
The matrix formula restricts to \(J_+\) on the positive metric graph
\(C_+=\{X+gX\}\), and to \(J_-\) on \(C_-=\{X-gX\}\);
the second matrix reverses the latter sign. These restrictions prove
the algebraic identities and positivity. The bracket computation in
Paper~I-D, Section~\ref{p1b:app:g14} proves integrability with exactly this twist
sign. It is the compact-group construction of \cite{GualtieriThesis},
with the left/right order fixed here by the bracket already used in
this paper.
\end{proof}
The central two-torus contributes no Lie bracket, so arbitrary positive
central metric data are permitted. These statements concern
bi-invariant metrics compatible with the chosen complex structures.
The nonzero periods distinguish \(H_g\) from the exact form
\(-\dd\Omega\); a global change by a two-form cannot remove \(H_g\).


\subsubsection{Full Dolbeault cohomology of the marked complex family}
\label{p1b:sec:g16}

This subsection concerns the complex structure on all of \(TM_8\) in
Theorem~\ref{p1b:thm:complex-family}. Let
\[
 W=\operatorname{span}_{\R}\{X_x,X_p\}\subset\mathfrak t
\]
be the plane generated by the two Hopf-circle directions. The scalar
circle directions complete \(\mathfrak t\) to a four-dimensional real
space. The root signs determine complex structures on the base
\(\mathbb{CP}^1\times\mathbb{CP}^1\); changing these signs does not
change the following dimensions.

\begin{theorem}[Finite model for the full Dolbeault complex]
\label{p1b:thm:g16-model}
Choose a basis \(\ell_1,\ell_2\) of
\((\mathfrak t^*)^{1,0}\) and the marked principal \(\T^4\) connection.
Normalize its two nonzero curvature components to real base forms
\(x,y\) of type \((1,1)\), with \(x^2=y^2=0\).
Let \(C\) be the matrix of those components in the complex connection
forms \(u_1,u_2\):
\(\dd u_i=C_{ix}x+C_{ip}y\).
The double complex
\[
 \mathcal A_J=
 \C[x,y]/(x^2,y^2)\otimes
       \Lambda(u_1,u_2,\bar u_1,\bar u_2)
\]
has degrees \(|x|=|y|=(1,1)\), \(|u_i|=(1,0)\),
\(|\bar u_i|=(0,1)\), and differentials
\[
 \bar\partial u_i=C_{ix}x+C_{ip}y,\qquad
 \partial\bar u_i=\overline{C_{ix}}x+\overline{C_{ip}}y,
\]
with all other generator differentials zero. Its inclusion in the
smooth forms on \(M_8\) induces the full Dolbeault cohomology algebra
and all pages of the Frölicher spectral sequence.
The curvature rank \(\rho=\rank_\C C\) is either one or two, and
\[
 \rho=1\quad\Longleftrightarrow\quad J_0W=W.
\]
\end{theorem}
\begin{proof}
The projection to the two Hopf bases is holomorphic, with complex
torus fiber \((\mathfrak t,J_0)/\Lambda\). The right \(\T^4\) action is
holomorphic. Averaging over it is a Dolbeault quasi-isomorphism:
every nonzero Fourier character is contracted by a vertical
\((0,1)\) vector field with nonzero character derivative.
The character derivatives have a uniform lower bound proportional
to the character norm, so this contraction acts on smooth Fourier
series. Invariant forms are the base forms with the four connection
generators adjoined. The base Dolbeault cohomology is represented,
as an algebra, by \(1,x,y,xy\). Filtering by the number of vertical
generators proves that its inclusion remains a quasi-isomorphism
after adjoining the connection forms. Paper~I-D, Section~\ref{p1b:app:g16}
provides the contraction and filtered argument.

The matrix \(C\) is restriction of the complex linear covectors to
the real curvature plane \(W\), up to nonzero real normalizations.
A real vector annihilated by every \((1,0)\) covector is zero, so
\(\rho\ne0\). A nonzero complex linear covector has a real
two-dimensional \(J_0\)-invariant kernel. Its annihilation of \(W\)
is therefore equivalent to \(J_0W=W\), proving the rank criterion.
\end{proof}

\begin{theorem}[Complete Hodge algebras and the two parameter strata]
\label{p1b:thm:g16-hodge}
Write
\(P_J(z,w)=\sum_{p,q=0}^4h^{p,q}_{\bar\partial}(M_8,J)z^pw^q\).
If \(\rho=2\), then
\[
 H_{\bar\partial}^{*,*}(M_8,J)
  \cong\Lambda(a_x,a_p,\gamma_x,\gamma_p),\qquad
 |a_i|=(0,1),\quad|\gamma_i|=(2,1),
\]
\[
 P_J(z,w)=(1+w)^2(1+z^2w)^2 .
\]
If \(\rho=1\), then
\[
 H_{\bar\partial}^{*,*}(M_8,J)
 \cong \C[\zeta]/(\zeta^2)\otimes\Lambda(s,a,b,\gamma),
\]
\[
 |\zeta|=(1,1),\quad |s|=(1,0),\quad
 |a|=|b|=(0,1),\quad|\gamma|=(2,1),
\]
\[
 P_J(z,w)=(1+z)(1+w)^2(1+zw)(1+z^2w).
\]
Exterior signs use total degree \(p+q\). There are no further
Hodge-number strata within the stated complex-structure family.
\end{theorem}
\begin{proof}
For rank two change the complex vertical basis so that
\(\bar\partial u_x=x,\ \bar\partial u_p=y\).
For each factor
\(\C[x]/x^2\otimes\Lambda(u_x,\bar u_x)\), the only classes before
adjoining \(\bar u_x\) are \(1,u_xx\). Tensoring the two factors gives
the first algebra.

For rank one choose vertical forms \(u,v\) with
\[
 \bar\partial u=x+\tau y,\qquad \bar\partial v=0,\qquad
 \Im\tau\ne0.
\]
The ratio is nonreal because the two real curvature vectors are
independent. The cokernel of multiplication by \(x+\tau y\) in
\(\C[x,y]/(x^2,y^2)\) has basis \(1,y\); its kernel has basis
\(x-\tau y,xy\). Thus representatives are
\[
 \zeta=[y],\quad \gamma=[u(x-\tau y)],\quad
 a=[\bar u],\quad s=[v],\quad b=[\bar v],
 \qquad \zeta\gamma=[uxy]\ne0.
\]
They give precisely the second algebra and no additional relations.
These normal forms exhaust both possible ranks.
\end{proof}

\begin{proposition}[Frölicher differentials and exchange sectors]
\label{p1b:prop:g16-frolicher}
For \(\rho=2\), the Frölicher sequence degenerates at \(E_1\).
For \(\rho=1\), its \(d_1\) is the derivation
\[
 d_1a=(\bar\tau-\tau)\zeta,\qquad
 d_1s=d_1b=d_1\gamma=d_1\zeta=0.
\]
It is nonzero and the sequence degenerates at \(E_2\), where
\[
 E_2\cong\Lambda(s,b,\gamma,\delta),\qquad
 |\delta|=(1,2),\quad \delta=[a\zeta].
\]
Every exchange-invariant complex structure in
Theorem~\ref{p1b:thm:complex-family} lies in the rank-two stratum.
For matched base generators its exchange trace polynomial is
\[
 \sum_{p,q}\tr(B^*|H^{p,q}_{\bar\partial})z^pw^q
       =(1-w^2)(1-z^4w^2).
\]
The dimensions of the two exchange sectors are the coefficients of
\(\frac12\{P_J(z,w)\pm(1-w^2)(1-z^4w^2)\}\).
\end{proposition}
\begin{proof}
In rank two, \(\partial\bar u_i=x_i=\bar\partial u_i\), so \(d_1\)
vanishes on the algebra generators. The total-degree dimensions at
\(E_1\) already equal
\[
 (1+t)^2(1+t^3)^2
   =\sum_k b_k(M_8)t^k,
\]
which leaves no room for a later differential.
In rank one, reducing
\(\partial\bar u=x+\bar\tau y\) modulo \(x+\tau y\) gives the stated
coefficient. The complex generated by \(a,\zeta\), with \(\zeta^2=0\),
has cohomology \(1,a\zeta\); the other three generators survive.
The resulting total-degree dimensions again equal the Betti
polynomial, proving degeneration at \(E_2\).

On \(W\), exchange has real eigenspaces of dimensions one and one.
A commuting complex structure cannot preserve these lines. Thus
\(J_0B=BJ_0\) rules out \(J_0W=W\).
In rank two the normalized vertical forms with derivatives \(x,y\)
are unique and exchange interchanges them. It therefore exchanges
both the pair \(a_x,a_p\) and the pair \(\gamma_x,\gamma_p\).
The trace on their exterior algebras is respectively \(1-w^2\) and
\(1-z^4w^2\), proving the result.
\end{proof}

\begin{corollary}[Derived-subgroup complex locus and exchange incompatibility]
\label{p1b:cor:g16-derived-complex}
Let
\[
 D_6=[M_8,M_8]=\SU(2)_x\times\SU(2)_p\cong S^3\times S^3
\]
and let \(W=\operatorname{span}_{\R}\{X_x,X_p\}\) be the Hopf plane
used above.  For every complex structure in
Theorem~\ref{p1b:thm:complex-family},
\[
 D_6\text{ is a complex submanifold of }(M_8,J)
 \quad\Longleftrightarrow\quad J_0W=W
 \quad\Longleftrightarrow\quad \rho=1.
\]
Consequently every exchange-invariant member of the declared complex
family lies in the complementary rank-two stratum and does not preserve
\(D_6\) as a complex submanifold.
\end{corollary}
\begin{proof}
The tangent algebra of the derived subgroup is
\[
 TD_6=U_x\oplus U_p\oplus W,
\]
where the two root planes \(U_x,U_p\) are preserved by the complex
structures of Theorem~\ref{p1b:thm:complex-family}.  Therefore
\(TD_6\) is \(J\)-invariant exactly when \(W\) is \(J_0\)-invariant.
Theorem~\ref{p1b:thm:g16-model} identifies this condition with
\(\rho=1\).  Proposition~\ref{p1b:prop:g16-frolicher} proves that every
exchange-invariant complex structure has \(\rho=2\), giving the final
assertion.
\end{proof}

In particular \(h^{1,0}\) jumps from zero to one precisely when the
Hopf plane becomes \(J_0\)-invariant. This computation concerns the
full complex fourfold; the CR and basic cohomology constructions in
Section~\ref{p1b:sec:g06} retain their distinct complexes.


\begin{theorem}[The complete pluriclosed existence locus of the complex family]
\label{p1b:thm:v214-full-skt-locus}
For the invariant complex structures of
Theorem~\ref{p1b:thm:complex-family}, the existence of a Hermitian
metric with \(\partial\bar\partial\omega=0\) is as follows.
Every rank-two member of Theorem~\ref{p1b:thm:g16-model} admits one.
In its rank-one normal form
\[
 \dd u=x+\tau y,\qquad \dd v=0,\qquad\Im\tau\ne0,
\]
one exists if and only if \(\Re\tau=0\).
This is an existence criterion for all Hermitian metrics, not only
for invariant metrics or the bi-invariant Bismut-flat subfamily.

More precisely, average over left translations and the holomorphic
right four-torus.  The resulting invariant Hermitian forms are exactly
\[
 \omega= a\omega_x+b\omega_p+
       \frac i2\sum_{r,s=1}^2h_{r\bar s}u_r\wedge\bar u_s,
 \qquad a,b>0,\quad H=(h_{r\bar s})=H^*>0,
\]
where \(\omega_i\) are the positive base area forms and the
\(u_r\) are the indicated vertical complex connection forms.
In rank two, choose \(\dd u_1=x,\dd u_2=y\); this metric is
pluriclosed exactly when \(\Re h_{1\bar2}=0\).
In rank one, put \((u_1,u_2)=(u,v)\); it is pluriclosed exactly
when \(\Re\tau=0\), independently of the positive matrix \(H\).
\end{theorem}
\begin{proof}
Left translations and the right four-torus act holomorphically on
the complex manifold, as proved in
Theorem~\ref{p1b:thm:g16-model} and its appendix.  Pullback commutes
with \(\partial,\bar\partial\), and averaging positive Hermitian
forms preserves positivity.  Thus any pluriclosed metric yields
one invariant under both actions.  Conversely any invariant example
is already a Hermitian metric on the full manifold.

The two independent Hopf rotations have inequivalent rotating root
planes and a trivial four-dimensional vertical space.  They remove
all root--vertical and mixed-root metric terms.  On each root plane
an invariant Hermitian metric is a positive scalar multiple of the
base metric; on the vertical space it is an arbitrary positive
Hermitian matrix.  This proves the claimed exhaustive averaged form.
The base area forms are closed and of type \((1,1)\), so contribute
nothing to \(\partial\bar\partial\omega\).

Set \(L_r=\bar\partial u_r=\dd u_r\).
The finite-model identities also hold as identities of global
forms, and give
\[
 \partial\bar\partial(u_r\wedge\bar u_s)
       =L_r\wedge\bar L_s.
\]
In rank two this makes
\(\partial\bar\partial\omega=i\Re(h_{1\bar2})xy\),
because \(x^2=y^2=0\).  A positive diagonal \(H\) supplies an
example for every such complex structure.  In rank one the only
nonzero term is
\[
 \partial\bar\partial\omega
   =\frac i2h_{1\bar1}(x+\tau y)(x+\bar\tau y)
   =i h_{1\bar1}\Re\tau\,xy.
\]
Here \(h_{1\bar1}>0\) and the global base four-form \(xy\) is
not zero as a differential form, even though its pullback may be
exact.  Vanishing is therefore equivalent to \(\Re\tau=0\).
The averaging argument proves necessity even for noninvariant
metrics, completing the existence classification.  The rank-one
criterion is unchanged by a nonzero real normalization of either
curvature generator or by interchanging the two generators.
\end{proof}

\begin{corollary}[Holomorphic canonical bundle and all plurigenera]
\label{p1b:cor:v214-canonical-kodaira}
For every invariant complex structure in
Theorem~\ref{p1b:thm:complex-family}, let
\(\pi:M_8\to\mathbb{CP}^1\times\mathbb{CP}^1\) be its
holomorphic torus fibration, with the base complex structures
specified by the two root signs.  Then
\[
 \boxed{K_{M_8}\cong\pi^*\mathcal O(-2,-2),\qquad
 H^0(M_8,K_{M_8}^{\otimes m})=0\quad(m\ge1).}
\]
Its Kodaira dimension is \(-\infty\).  The canonical line is
smoothly trivial but is of infinite order in the holomorphic Picard
group; in particular it has no holomorphic volume form.  These
conclusions transport to all left-invariant integrable complex
structures by Theorem~\ref{p1b:thm:g22-global-complex}.
\end{corollary}
\begin{proof}
The vertical complex torus acts holomorphically and freely.  Its two
complex infinitesimal translation fields are a global holomorphic
frame of the vertical tangent bundle.  The differential of the
holomorphic submersion therefore gives an exact sequence
\[
 0\longrightarrow\mathcal O_{M_8}^{\oplus2}
 \longrightarrow T^{1,0}M_8
 \longrightarrow\pi^*T^{1,0}
        (\mathbb{CP}^1\times\mathbb{CP}^1)\longrightarrow0.
\]
Taking dual determinant lines proves the canonical-bundle formula.
Each fiber is a compact connected complex torus, so every
holomorphic function on a fiber is constant.  In local holomorphic
submersion coordinates this also proves holomorphic descent of a
function constant on the fibers; hence
\(\pi_*\mathcal O_{M_8}=\mathcal O_{\mathbb{CP}^1\times\mathbb{CP}^1}\).
Apply this descent to local coefficients of sections of a pulled-back
line.  It identifies the displayed pluricanonical section space
with \(H^0(\mathbb{CP}^1\times\mathbb{CP}^1,
\mathcal O(-2m,-2m))\), which is zero: restriction to either
projective-line fiber has negative degree and therefore no nonzero
holomorphic section.  All positive plurigenera vanish, proving
Kodaira dimension \(-\infty\).  Smooth triviality follows from
Corollary~\ref{p1b:cor:complex-topology}.  If any positive tensor
power of the canonical line were holomorphically trivial, it would
have the nonzero constant section, contradicting the section
calculation.  Pullback under the automorphisms in the global normal
form preserves every assertion.
\end{proof}

\subsubsection{Flat special-unitary and exchange-equivariant bundles}
\begin{theorem}[Flat special-unitary and exchange-equivariant families]
\label{p1b:thm:equivariant-flat}
Flat \(\SU(r)\) connections up to gauge on \(M_8\) are classified by
\[
 \left\{((z_1,w_1),\ldots,(z_r,w_r))\in(\T^2)^r:
       \prod_jz_j=\prod_jw_j=1\right\}/S_r.
\]
All underlying bundles are trivial. For \(r=2\) the parameter space
is \(\T^2/((z,w)\sim(z^{-1},w^{-1}))\), with four fixed points
of the involution.

A flat unitary bundle with a parallel exchange lift squaring to the
identity is, up to equivariant unitary isomorphism, a unique orthogonal
sum of the following irreducibles:
\begin{enumerate}[label=(\roman*),nosep]
\item for each \(z\in\T\), a one-dimensional character
      \(U_t=U_E=z\), with exchange \(J=+1\) or \(J=-1\);
\item for each unordered pair \(\{z,w\}\), \(z\ne w\), a
      two-dimensional block
\[
 U_t=\begin{pmatrix}z&0\\0&w\end{pmatrix},\quad
 U_E=\begin{pmatrix}w&0\\0&z\end{pmatrix},\quad
 J=\begin{pmatrix}0&1\\1&0\end{pmatrix}.
\]
\end{enumerate}
Multiplicities are arbitrary finite nonnegative integers. Thus a flat
\(\U(r)\) bundle has such an exchange lift exactly when its joint
holonomy multiplicities are invariant under \((z,w)\mapsto(w,z)\).
On a diagonal character, the positive and negative lift multiplicities
are the remaining independent data.
\end{theorem}
\begin{proof}
Since \(\pi_1(M_8)=\Z^2\), flat bundles are classified by two
commuting holonomies. Commuting unitary matrices diagonalize
simultaneously; special unitarity imposes the two products.
The representation factors through the projection to \(\T^2\).
Every principal \(\SU(r)\) bundle over that two-dimensional CW
complex is trivial, since \(\SU(r)\) is connected and simply
connected for \(r\ge2\). Pullback proves the assertion, including
the immediate \(r=1\) case.
The \(\SU(2)\) Weyl group acts by simultaneous inversion and has
the four fixed points \((\pm1,\pm1)\).

The identity is fixed by exchange, so a parallel lift is a
representation of \(\Z^2\rtimes\Z/2\) with
\(J^2=I\), \(JU_tJ^{-1}=U_E\).
On the joint eigenspace of \((z,w)\), \(J\) maps unitarily to the
space for \((w,z)\). For a two-point orbit choose bases in one
space and their images in the other; the displayed blocks result.
For a fixed character diagonalize the unitary involution into
its \(+1\) and \(-1\) eigenspaces. This proves existence,
irreducibility of the listed blocks, uniqueness and exhaustion.
\end{proof}


\subsection{Metrics, invariant connections and covariant scalar spectra}
\label{p1b:app:metrics}
\begin{theorem}[Bi-invariant metrics and the product Dirac gap]
\label{p1b:thm:metrics}
Every bi-invariant Riemannian metric on \(M_8\) is
\[
 g=R_x^2g_{S^3}+R_p^2g_{S^3}
       +\sum_{i,j=t,E}G_{ij}\dd\theta_i\dd\theta_j,
 \qquad R_x,R_p>0,\quad G=G^T>0.
\]
It is exchange invariant exactly when \(R_x=R_p\) and \(G_{tt}=G_{EE}\).
All sixteen bi-invariant forms are parallel and harmonic, and their
wedge products are harmonic. Its scalar curvature is
\(6/R_x^2+6/R_p^2\).
For scalar-cycle spin shifts \(\sigma_i\in\{0,\frac12\}\), the
untwisted product spin Dirac square has spectral values
\[
 \frac{(\ell_x+3/2)^2}{R_x^2}
 +\frac{(\ell_p+3/2)^2}{R_p^2}
 +(n+\sigma)^TG^{-1}(n+\sigma),
 \qquad \ell_x,\ell_p\ge0,\quad n\in\Z^2.
\]
In particular it is invertible with squared gap at least
\(\frac94(R_x^{-2}+R_p^{-2})\).
\end{theorem}
\begin{proof}
An adjoint-invariant inner product has scalar Killing multiples on
the two independent simple ideals and an arbitrary positive form on
the center; mixed terms vanish. Exchange gives the stated equalities.
For left-invariant fields the Levi-Civita connection is
\(\nabla_XY=[X,Y]/2\), so adjoint-invariant forms are parallel.
Their cohomology dimensions exhaust the Betti numbers.
The product curvature and Clifford-anticommuting factor Dirac
operators give the curvature and squared spectrum. The round sphere
spin eigenvalues are \(\pm(\ell+3/2)/R\).
\end{proof}

\begin{proposition}[A natural torsion family]
\label{p1b:result:a-natural-torsion-family}
For \(\nabla^s_XY=s[X,Y]\) on left-invariant fields,
\[
 T^s(X,Y)=(2s-1)[X,Y],\qquad
 R^s(X,Y)Z=s(s-1)[[X,Y],Z].
\]
Every member is metric for a bi-invariant metric, the endpoints
\(s=0,1\) are flat, \(s=\frac12\) is Levi-Civita, and inversion
interchanges \(s\) with \(1-s\). Their affinely parametrized
geodesics through the identity are one-parameter subgroups.
\end{proposition}
\begin{proof}
Insert the connection into the definitions of torsion and curvature
and use Jacobi. Adjoint operators are skew for a bi-invariant metric.
Inversion exchanges the parallel left and right frames, and the
antisymmetric connection term vanishes on a repeated velocity.
\end{proof}

\begin{proposition}[Parallel maps under the midpoint connection]
\label{p1b:result:parallel-maps-under-the-midpoint-connection}
For representations
\(V=\bigoplus_jV_j\otimes\C^{m_j}\) and
\(W=\bigoplus_jV_j\otimes\C^{n_j}\), the parallel maps of associated
bundles for the midpoint connection are
\[
 \Hom_\nabla(E_V,E_W)
 \cong\bigoplus_j I_{V_j}\otimes M_{n_j\times m_j}(\C).
\]
\end{proposition}
\begin{proof}
A parallel map is determined by its value at one point and that value
must commute with holonomy. The full holonomy of
Theorem~\ref{p1b:thm:homogeneous} and Schur's lemma give the answer;
holonomy invariance also makes parallel extension path independent.
\end{proof}

\begin{theorem}[Covariant scalar spectrum]
\label{p1b:thm:covariant-spectrum}
On \(S_R^3\), supply the connection \(A_f=f\Theta\), with constant
\(f\), in the spin-\(\iota\) representation. Use an orthonormal
unit-sphere coframe satisfying
\(\dd e^a=-\epsilon_{abc}e^b\wedge e^c\), write
\(\Theta=-2iI_ae^a\), and let \(J=iL/2\).
Then
\[
 \Delta_{f,\iota}=\frac4{R^2}(J+fI)^2.
\]
On the coupling \(V_j\otimes V_\iota\to V_k\), its eigenvalue is
\[
 R^2\lambda_{jk}
 =4\{(1-f)j(j+1)+fk(k+1)+(f^2-f)\iota(\iota+1)\},
\]
where \(j\in\frac12\Z_{\ge0}\),
\(|j-\iota|\le k\le j+\iota\) in unit steps, with multiplicity
\((2j+1)(2k+1)\).
\end{theorem}
\begin{proof}
The covariant derivatives in the chosen frame are proportional to
\(J_a+fI_a\). The frame fields have zero divergence. Expand the square
and use \(2J\cdot I=(J+I)^2-J^2-I^2\).
Peter--Weyl supplies the spectator multiplicity \(2j+1\), and
Clebsch--Gordan supplies the coupled multiplicity \(2k+1\).
\end{proof}

\begin{corollary}[Embedded adjoint gap and transverse solution]
\label{p1b:cor:adjoint-gap}
For the embedding \(\SU(2)\to\SU(3)\), \(g\mapsto\diag(g,1)\),
the complexified adjoint restricts as
\(V_1\oplus2V_{1/2}\oplus V_0\).
For \(0<f<1\), the kernel of the adjoint scalar covariant Laplacian
is the constant centralizer line, and its nonzero spectral gap is
\[
 \lambda_*(f)=\frac3{R^2}\min\{f^2,(1-f)^2\}.
\]
Let \(P_c\) project onto that line and
\(G=(\Delta|_{P_c^\perp})^{-1}(I-P_c)\). Then
\(\|G\|=\lambda_*^{-1}\).
For adjoint one-forms define \(P_T=I-\dd_A G\dd_A^*\).
The equation \(\dd_A^*E=j\) is solvable exactly when \(P_cj=0\);
all solutions are
\[
 E=\dd_A Gj+E_T,\qquad \dd_A^*E_T=0,
\]
and the minimal squared norm is \(\langle j,Gj\rangle\).
\end{corollary}
\begin{proof}
For \(f\ge0\), the smallest \(k=|j-\iota|\) minimizes the
displayed eigenvalue. For \(\iota=\frac12\), its minimum comes from
\((j,k)=(0,\frac12),(\frac12,0)\), giving \(3f^2\) and
\(3(1-f)^2\). For \(\iota=1\), the candidates \(j=0,\frac12,1\)
give \(8f^2,\ 3-8f+8f^2,\ 8(1-f)^2\), all at least the doublet
minimum on \([0,1]\); larger \(j\) increase it.
The trivial representation has its constant kernel and next
eigenvalue \(3/R^2\). This proves the gap and kernel.
Functional calculus gives \(G\), and
\(\dd_A^*\dd_A G=I-P_c\). Orthogonal range--kernel decomposition
then proves the solution and minimal-norm formulas.
No assertion that \(\dd_A^2=0\) is needed.
\end{proof}

\begin{proposition}[A controlled local gauge slice]
\label{p1b:result:a-controlled-local-gauge-slice}
Let \(\dd_B\) denote a fixed midpoint covariant derivative and
let \(a\) be a smooth perturbing one-form. On the centralizer
complement suppose
\[
 \|[a,\xi]\|_{L^2}\le\beta\|\xi\|_{L^2},
 \qquad \beta<\sqrt{\lambda_*}.
\]
The projected operator
\((I-P_c)\dd_B^*\dd_{B+a}(I-P_c)\) is invertible, with
\[
 \|M_a^{-1}\|_{L^2\to L^2}
 \le(\lambda_*-\beta\sqrt{\lambda_*})^{-1}.
\]
\end{proposition}
\begin{proof}
Factor it in quadratic-form sense as
\(\Delta_B^{1/2}(I+T)\Delta_B^{1/2}\).
The bound on \(a\) gives \(\|T\|\le\beta/\sqrt{\lambda_*}<1\).
The Neumann inverse and \(\|\Delta_B^{-1/2}\|\le\lambda_*^{-1/2}\)
prove the estimate. Elliptic regularity gives the smooth
operator's inverse on the usual \(H^2\)-to-\(L^2\) domain.
\end{proof}


\subsubsection{Nontrivial scalar-line operator families}
\label{p1b:sec:g13}

Let \(L_q=\mathcal P_{\rm sc}^{\otimes q}\) on the oriented scalar
torus, where \(q\in\Z\). Fix the flat metric \(G\) from
Theorem~\ref{p1b:thm:metrics}, with area
\(A=(2\pi)^2\sqrt{\det G}\), and choose a unitary connection satisfying
\[
 \frac{iF_\nabla}{2\pi}=q\,\frac{\dd V}{A}.
\]
Such connections exist. For fixed curvature their gauge classes form
the flat-connection torsor \(H^1(\T^2;\R)/H^1(\T^2;\Z)\).

\begin{theorem}[Exact scalar-line and product spectra]
\label{p1b:thm:g13-landau}
Suppose \(q\ne0\), put \(b=2\pi|q|/A\), and take the nonnegative
Bochner Laplacian \(\Delta_\nabla=\nabla^*\nabla\) with its
self-adjoint domain \(H^2(L_q)\). Its complete spectrum is
\[
 \lambda_r=b(2r+1),\qquad
 \operatorname{mult}(\lambda_r)=|q|,\qquad r=0,1,\ldots.
\]
The spectrum is independent of the flat part of the connection. Its
heat trace is \(|q|/(2\sinh bt)\), \(t>0\).
For the pullback line on the round product \(M_8\), the scalar spectrum is
\[
 \lambda_{k_x,k_p,r}
 =\frac{k_x(k_x+2)}{R_x^2}
  +\frac{k_p(k_p+2)}{R_p^2}+b(2r+1),
\]
with labelled multiplicity \((k_x+1)^2(k_p+1)^2|q|\).
Coincident values from distinct labels have their multiplicities added.
\end{theorem}
\begin{proof}
In an oriented parallel orthonormal frame set
\(\pi_j=-i\nabla_j\). For \(q>0\),
\([\pi_1,\pi_2]=ib\) and
\[
 a=\frac{\pi_1+i\pi_2}{\sqrt{2b}},\qquad
 [a,a^*]=1,\qquad \Delta_\nabla=b(2a^*a+1).
\]
Paper~I-D, Section~\ref{p1b:app:g13} constructs exactly \(q\) independent ground
states, proves completeness by lowering and raising, and gives unitary
transport for every flat part. Complex conjugation handles \(q<0\).
Summing the geometric series proves the heat trace. Tensoring with the
complete scalar sphere bases gives the product statement. The
normalization agrees with the flat-torus specialization of
\cite[Appendix 7]{CharlesLandau}, after doubling that reference's
half-Bochner Laplacian.
\end{proof}

\begin{proposition}[Twisted spin Dirac spectrum and product gap]
\label{p1b:prop:g13-dirac}
For any torus spin structure, the Dirac operator twisted by \(L_q\)
has index \(q\), \(|q|\) zero modes of the chirality \(\sgn q\), and
nonzero eigenvalues
\[
 \pm\sqrt{2br},\qquad r=1,2,\ldots,
\]
each of multiplicity \(|q|\). Chirality is fixed by the oriented
Pauli representation used in Paper~I-D, Section~\ref{p1b:app:g13}.
Writing \(m_\ell=(\ell+1)(\ell+2)\), the corresponding eight-dimensional
product Dirac eigenvalues are
\[
 \pm\left\{\frac{(\ell_x+3/2)^2}{R_x^2}
          +\frac{(\ell_p+3/2)^2}{R_p^2}+2br\right\}^{1/2},
 \qquad r\ge0,
\]
with multiplicity per sign
\[
 \begin{cases}
 4m_{\ell_x}m_{\ell_p}|q|,&r=0,\\
 8m_{\ell_x}m_{\ell_p}|q|,&r\ge1.
 \end{cases}
\]
It is invertible, has squared gap
\(\frac94(R_x^{-2}+R_p^{-2})\), and has chiral index zero.
\end{proposition}
\begin{proof}
In the positive-degree branch,
\[
 D_{\T^2}=\sqrt{2b}\begin{pmatrix}0&a^*\\a&0\end{pmatrix},
 \qquad D_{\T^2}^2=\begin{pmatrix}\Delta_\nabla-b&0\\
                                      0&\Delta_\nabla+b\end{pmatrix}.
\]
The oscillator decomposition gives the zero modes and paired spectrum.
A spin change is a flat line twist, which leaves these counts unchanged.
For the product of the two odd-dimensional sphere spin bundles one
includes the standard two-dimensional doubling factor. Each sphere
Dirac square has multiplicity \(2m_\ell\); their product square
therefore has multiplicity \(8m_{\ell_x}m_{\ell_p}\).
The torus square has multiplicity \(|q|\) at zero and \(2|q|\) above
zero. Product chirality pairs the two nonzero signs, giving the stated
halves. The least sphere levels yield the gap. Independently,
\(\widehat A(TM_8)=1\) and \(\operatorname{ch}(L_q)=e^{q\eta}=1+q\eta\),
so the degree-eight index density is zero.
\end{proof}
The \(q=0\) branch is the flat-twisted Fourier family already stated
above. Exchange transports \((L_q,\nabla,G)\) to
\((B^*L_q,B^*\nabla,B^*G)\), of degree \(-q\), and gives the
corresponding unitary scalar spectral equivalence. An identification
with fixed connection data or an involutive antiunitary lift requires
that additional identification to be specified.


\subsection{Product spectra and exchange heat sectors}
\label{p1b:sec:product-heat}

\begin{theorem}[Complete scalar and Hodge heat traces]
\label{p1b:thm:product-hodge}
For the bi-invariant metric of Theorem~\ref{p1b:thm:metrics}, the scalar
eigenvalues, with multiplicities combined when values coincide, are
\[
 \lambda=\frac{k_x(k_x+2)}{R_x^2}
          +\frac{k_p(k_p+2)}{R_p^2}+n^TG^{-1}n,\qquad
 (k_x,k_p,n)\in\Z_{\ge0}^2\times\Z^2,
\]
with block multiplicity \((k_x+1)^2(k_p+1)^2\).
For \(t>0\) define the absolutely convergent series
\begin{align*}
 A_R(t)&=\sum_{k=0}^\infty(k+1)^2e^{-tk(k+2)/R^2},\\
 C_R(t)&=\sum_{k=1}^\infty(k+1)^2e^{-tk(k+2)/R^2}
       +2\sum_{k=1}^\infty k(k+2)e^{-t(k+1)^2/R^2},\\
 \Theta_G(t)&=\sum_{n\in\Z^2}e^{-tn^TG^{-1}n}.
\end{align*}
Then the generating polynomial of all Hodge heat traces on \(M_8\) is
\begin{align*}
 \mathcal Z(z,t)&:=\sum_{p=0}^8 z^p\Tr(e^{-t\Delta_p})\\
 &=\bigl[A_{R_x}(t)(1+z^3)+C_{R_x}(t)(z+z^2)\bigr]\\
 &\quad\cdot\bigl[A_{R_p}(t)(1+z^3)+C_{R_p}(t)(z+z^2)\bigr]
       (1+z)^2\Theta_G(t).
\end{align*}
Thus coefficient extraction gives every degree, eigenvalue and
multiplicity. Its zero-eigenvalue polynomial is
\((1+z)^2(1+z^3)^2\).
\end{theorem}
\begin{proof}
The product Laplacian is the sum of the factor Laplacians.
Scalar harmonics of degree \(k\) on \(S_R^3\) have eigenvalue
\(k(k+2)/R^2\) and multiplicity \((k+1)^2\).
Nonconstant scalar eigenfunctions give the exact one-forms with the
same eigenvalues. The coexact one-forms have eigenvalues
\((k+1)^2/R^2\) and multiplicities \(2k(k+2)\), \(k\ge1\);
these are the two \(\Spin(4)\) summands
\(((k+1)/2,(k-1)/2)\) and its interchange.
Hodge duality gives the two- and three-form terms.
The flat torus has the same Fourier eigenvalue on each constant
exterior form, producing \((1+z)^2\Theta_G\).
Tensor the complete orthogonal eigenbases; eigenvalues add and
multiplicities multiply. Polynomial multiplicity growth and positive
quadratic eigenvalue growth prove convergence and permit these products.
Taking \(t\to\infty\) retains exactly the harmonic forms.
\end{proof}

\begin{theorem}[Exact exchange-even and exchange-odd heat traces]
\label{p1b:thm:exchange-heat}
For an exchange-invariant bi-invariant metric put \(R_x=R_p=R\) and
\[
 \Theta_{\rm diag}(t)=
 \sum_{n\in\Z}\exp[-t(n,n)G^{-1}(n,n)^T].
\]
The scalar heat traces in the two exchange sectors are
\[
 Z_{0,\pm}(t)=\tfrac12\bigl[
 A_R(t)^2\Theta_G(t)\ \pm\ A_R(2t)\Theta_{\rm diag}(t)\bigr].
\]
In all form degrees, put
\[
 \mathcal Z_B(z,t)=
 \bigl[A_R(2t)(1-z^6)+C_R(2t)(-z^2+z^4)\bigr]
 (1-z^2)\Theta_{\rm diag}(t).
\]
The \(p\)-form heat traces in the exchange sectors are the coefficients
of \((\mathcal Z\pm\mathcal Z_B)/2\).
In particular the exchange trace on harmonic forms is
\[
 \lim_{t\to\infty}\mathcal Z_B(z,t)=1-z^2-z^6+z^8.
\]
\end{theorem}
\begin{proof}
For a trace-class operator \(T\), the flip on a tensor square has
\(\Tr(\mathrm{flip}(T\otimes T))=\Tr(T^2)\).
In a graded tensor square the equal-degree \(p\) summand adds
the Koszul sign \((-1)^{p^2}=(-1)^p\) and degree \(2p\).
This gives the displayed sphere expression. A torus Fourier mode
contributes to the exchange trace only when its two charges agree.
On its constant exterior forms, the exchange trace polynomial is
\(\det(I+zS)=1-z^2\).
Finally \((I\pm B^*)/2\) are the two orthogonal projectors commuting
with the Hodge Laplacian. Taking their traces proves every formula.
\end{proof}

\part{Finite carriers, level structures and the complete product atlas}
\section{Finite carriers, level structures, and exact transport}
\label{p1b:sec:organized-carriers}

\subsection{Finite-band and level-one completion atlases}
\label{p1b:sec:finite-families}

\begin{definition}[Scalar and enriched product families]
\label{p1b:def:product-bands}
Let \(E_k\) be the degree-\(k\) scalar harmonic space on \(S^3\) and
\[
 P_m=\bigoplus_{k=0}^m E_k,\qquad
 D_m=\dim P_m=\frac{(m+1)(m+2)(2m+3)}6.
\]
For a finite charge set \(F\subset\Z^2\), let
\(\C[F]=\operatorname{span}\{z_t^nz_E^w:(n,w)\in F\}\).
Define the scalar product carrier
\[
 W_{m_x,m_p,F}=P_{m_x}\otimes P_{m_p}\otimes\C[F].
\]
Separately, let
\[
 H_1=(0,0)\oplus(\tfrac12,\tfrac12)\oplus(1,0)\oplus(0,1)
          \oplus(\tfrac12,0)\oplus(0,\tfrac12)
\]
be Paper I-A's fifteen-dimensional form/spin carrier. The enriched
level-one product is
\(\mathbb H_{1,1,F}=H_{1,x}\otimes H_{1,p}\otimes\C[F]\).
These are definitions for every displayed integer and finite set.
\end{definition}

\begin{theorem}[Exhaustion of finite regular bands]
\label{p1b:thm:finite-band-classification}
Every finite-dimensional scalar subspace of \(L^2(M_8)\) invariant
under all left and right translations is exactly the sum of a finite
set of distinct Peter--Weyl blocks
\[
 E_{k_x}\otimes E_{k_p}\otimes\C z_t^nz_E^w.
\]
It is exchange invariant precisely when that finite index set is
closed under \((k_x,k_p,n,w)\mapsto(k_p,k_x,w,n)\).
In particular,
\[
 \dim W_{m_x,m_p,F}=D_{m_x}D_{m_p}|F|,\qquad
 \dim\mathbb H_{1,1,F}=225|F|.
\]
For a fixed Hopf charge \(q\) and retained degrees
\(|q|,|q|+2,\ldots,|q|+2L\), the one-sphere dimension is
\[
 \sum_{a=0}^L(|q|+2a+1)=(L+1)(L+|q|+1).
\]
Thus arbitrary finite degrees, charges, matched exchange orbits and
fixed-charge bands are described without selecting a preferred cutoff.
\end{theorem}
\begin{proof}
Peter--Weyl decomposes the regular bimodule as
\(\widehat{\bigoplus}_{\pi}\pi\otimes\pi^*\), with multiplicity one
for the product of left and right groups. A finite-dimensional
invariant subspace is a finite direct sum of these inequivalent
irreducibles. The irreducibles of \(M_8\) are classified in
Theorem~\ref{p1b:thm:irreducibles}, giving exactly the displayed
blocks. Pullback by exchange permutes their indices as stated.
Sum \((k+1)^2\) for \(D_m\), and multiply tensor dimensions.
In \(E_k\cong V_{k/2}\otimes V_{k/2}^*\), each allowed right Hopf
weight has multiplicity \(k+1\); the weights are
\(-k,-k+2,\ldots,k\). This proves the final sum.
\end{proof}

\begin{theorem}[Factorwise level-one completion cones]
\label{p1b:thm:product-completion}
On \(H_1\), take Hermitian \(\Spin(4)\)-equivariant coordinate
quadruples transforming as \((\tfrac12,\tfrac12)\), with the
scalar-to-vector edge fixed to normalized scalar multiplication.
Modulo sector rephasings their full completion cone is
\[
 \mathfrak M_1\cong\R_{\ge0}^3,\qquad
 (r_{1,0},r_{0,1},\varrho_{\rm sp}).
\]
On \(\mathbb H_{1,1,F}\), define factorwise coordinates to act as
\(C_{x,a}\otimes I\otimes I\) and
\(I\otimes C_{p,a}\otimes I\). Their completion cone is
\(\mathfrak M_1\times\mathfrak M_1\cong\R_{\ge0}^6\).
For an exchange-invariant \(F\), the completions compatible with the
fixed matched exchange are the diagonal cone \(\R_{\ge0}^3\).
All these cones include the faces where any reduced magnitude vanishes.
\end{theorem}
\begin{proof}
Tensoring the six inequivalent sectors of \(H_1\) with
\((\tfrac12,\tfrac12)\) leaves exactly four unoriented internal
edges: the scalar--vector edge, the vector--\((1,0)\) edge,
the vector--\((0,1)\) edge, and the edge between the two chiral
spinors. Each intertwiner space is one dimensional. Hermiticity
determines the reverse edge by adjoint. The graph is a three-edge
star and a disjoint one-edge component, hence a forest.
Sector phases remove all nonzero edge phases, including on every
boundary stratum; fixing the scalar edge leaves exactly three
nonnegative magnitudes. The factorwise definition makes the two
choices independent, so their cones multiply. Exchange interchanges
the two tuples, making its fixed cone diagonal.
No claim about completions with additional mixed-factor coordinate
operators is included in this factorwise classification.
\end{proof}

\begin{theorem}[The full equivariant product-coordinate atlas]
\label{p1b:thm:full-product-coordinate-atlas}
On \(H_{1,x}\otimes H_{1,p}\), consider all Hermitian coordinate
octuples transforming as the sum of the two quaternionic vector
representations of \(\Spin(4)_x\times\Spin(4)_p\).
No factorwise restriction or scalar-edge normalization is imposed in
this larger class. Let \(G_1\) be the six-vertex, four-edge graph from
the proof of Theorem~\ref{p1b:thm:product-completion}, and let
\(G_\square=G_1\square G_1\) be its Cartesian product.
Then the coordinate space is \(\C^{48}\) in chosen normalized
intertwiner bases, with sector rephasing gauge \(\U(1)^{36}\).

For any support \(E'\) of nonzero edges, let \(c(E')\) count the
connected components of its spanning subgraph, including isolated
vertices. The corresponding quotient stratum is
\[
 (\R_{>0})^{|E'|}\times
 \T^{\,|E'|-36+c(E')}.
\]
Consequently the stratum with all edges nonzero is
\[
 (\R_{>0})^{48}\times\T^{16}.
\]
For the fixed matched exchange lift, the exchange-compatible
all-nonzero stratum, modulo exchange-commuting sector rephasings, is
\[
 (\R_{>0})^{24}\times\T^6.
\]
These are classification spaces for the specified equivariance and
gauge. The factorwise cone is the subfamily obtained by making each
edge coefficient independent of its spectator sector and fixing its
two scalar-edge normalizations.
\end{theorem}
\begin{proof}
Each of the six sphere sectors occurs with multiplicity one.
An \(x\)-vector operator must preserve the \(p\)-sector and connect
an edge of \(G_1\) in the \(x\)-factor; its intertwiner is unique
up to one complex scalar. There are \(4\cdot6=24\) such edges.
There are another \(24\) for the \(p\)-vector. Hermiticity fixes
the reverse coefficients. Product sectors give \(6\cdot6=36\)
independent phases, acting on an edge coefficient by the phase
difference of its endpoints.
For any nonzero support, fix a spanning forest and gauge its
coefficients positive. The remaining phases are the independent
cycle phases, numbering \(|E'|-36+c(E')\); all edge magnitudes
remain free. This proves every stratum, including vanishing edges.

The product graph has four components, the products of the two
components of \(G_1\), so the all-nonzero cycle rank is
\(48-36+4=16\). Exchange swaps \(x\)- and \(p\)-edges, giving
24 edge orbits and equal complex coefficients within each orbit.
It has \((36+6)/2=21\) vertex orbits. Its commuting sector gauge
therefore has dimension 21, with kernel the three constant phases
on the component orbits: the star-square, the paired star--edge
components, and the edge-square. Its phase action has rank 18,
leaving \(24-18=6\) independent phase coordinates.
Equivalently this is the cycle rank of the quotient graph, which
has 24 edges, 21 vertices and three components. Spanning forests
on that graph give the stated torus and establish uniqueness.
\end{proof}

\begin{corollary}[Product protected Clifford module]
\label{p1b:cor:product-clifford}
On the interior where the two protected magnitudes are nonzero, normalize
each protected quadruple to positive Clifford generators
\(\Gamma_{x,a},\Gamma_{p,a}\), \(a=1,\ldots,4\).
With \(\Gamma_{x,5}=\Gamma_{x,1}\Gamma_{x,2}
\Gamma_{x,3}\Gamma_{x,4}\), the eight matrices
\[
 \Gamma_{x,a}\otimes I,\qquad
 \Gamma_{x,5}\otimes\Gamma_{p,a}
\]
give the irreducible complex \(\Cl_{8,0}\) module of dimension \(16\)
on the product protected block. If a protected magnitude vanishes,
its unnormalized four coordinates vanish; that boundary member
still belongs to the completion cone.
\end{corollary}
\begin{proof}
The fourfold volume element squares to \(I\) and anticommutes with
each first-factor generator. This gives the eight Clifford
anticommutators. The generators contain \(M_4(\C)\otimes I\);
its volume element recovers \(I\otimes\Gamma_{p,a}\) from the
second four matrices. They therefore generate
\(M_4(\C)\otimes M_4(\C)=M_{16}(\C)\).
The zero-magnitude statement is the corresponding zero intertwiner.
\end{proof}

\begin{theorem}[Product cubature and exact finite multiplication]
\label{p1b:thm:product-cubature}
Let \(Q_x,Q_p\) be positive weighted spherical designs on the two
unit three-spheres, of strengths \(t_x,t_p\). Let the scalar circles
carry uniform grids of \(L_t,L_E\) points. Their product rule
integrates exactly every tensor product with sphere degrees at most
\(t_x,t_p\) and scalar frequencies satisfying
\(|n|<L_t,\ |w|<L_E\).
For
\[
 W=P_{m_x}\otimes P_{m_p}\otimes
 \operatorname{span}\{z_t^nz_E^w:|n|\le N_t,\ |w|\le N_E\},
\]
weighted evaluation is an isometry onto its image whenever
\[
 2m_x\le t_x,\quad2m_p\le t_p,\quad
 2N_t<L_t,\quad2N_E<L_E.
\]
Products of \(r\) retained modes and their conjugates are integrated
exactly under
\[
 rm_x\le t_x,\quad rm_p\le t_p,\quad
 rN_t<L_t,\quad rN_E<L_E.
\]
More generally, the matrix of compressed scalar multiplication by
a symbol with sphere degrees \(d_x,d_p\) and frequency bounds
\(D_t,D_E\) is computed exactly if
\[
 2m_i+d_i\le t_i,\qquad 2N_j+D_j<L_j.
\]
For matched designs, grids and bands, evaluation also intertwines exchange.
\end{theorem}
\begin{proof}
For a uniform \(L\)-point circle grid,
\(L^{-1}\sum_{\ell=0}^{L-1}e^{2\pi i n\ell/L}\) is one if
\(L\mid n\), and zero otherwise. Within \(|n|<L\) this agrees
with the circle integral. The design identities and the scalar
identities multiply under product measure. Products add degree
and frequency bounds, and complex conjugation preserves their
absolute bounds. Apply the exact rule to \(\bar uv\), to the
\(r\)-fold products, and to \(\bar u f v\), respectively.
The square roots of the positive quadrature weights define the
isometric evaluation. Exchange permutes matched sample factors
and their equal weights. Factorwise counterexamples to a stronger
design strength or the character of frequency \(L_j\) also show
why the degree and no-alias conditions cannot simply be discarded.
\end{proof}

\begin{proposition}[Finite function-algebra closure]
\label{p1b:prop:finite-function-algebra}
A finite-dimensional unital subalgebra of \(C^\infty(M_8;\C)\)
consists only of constants. For a finite orthogonal projection \(P\),
compressed multiplication instead obeys the exact identity
\[
 PM_fP\,PM_gP-PM_{fg}P=-PM_f(I-P)M_gP.
\]
\end{proposition}
\begin{proof}
For any \(f\) in a finite-dimensional unital algebra, its powers
are linearly dependent, so \(p(f)=0\) for a nonzero polynomial \(p\).
Its range is a finite subset of the roots of \(p\). Continuity and
connectedness of \(M_8\) force \(f\) to be constant. The second
formula inserts \(I=P+(I-P)\) between \(M_f\) and \(M_g\).
\end{proof}

\subsection{Arbitrary levels, mixed carriers and fixed-factor limits}
\label{p1b:sec:all-levels}

The product construction has independent levels.  We develop both the full
Paper I-A carrier family and its scalar subfamily, including the mixed product
used when one side is a scalar harmonic band and the other is the enriched
level-one carrier.  The results cover every nonnegative integer pair, their
boundary cases, exact comparison maps, and increasing-level limits.
All spaces in this section use normalized invariant inner products.

\subsubsection{The three product families and their exact sizes}

\begin{definition}[Independent scalar and enriched levels]
\label{p1b:def:all-level-products}
Write $U_{a,b}=V_{a/2}\otimes V_{b/2}$ and, for $m,\ell\ge0$, set
\[
 \Lambda_m=\{(a,b)\in\Z_{\ge0}^2:a+b\le2m\},\qquad
 H_m=\bigoplus_{(a,b)\in\Lambda_m}U_{a,b},\qquad
 P_\ell=\bigoplus_{k=0}^{\ell}U_{k,k}.
\]
Here $H_m$ is precisely Paper I-A's $\mathcal H_m$ and $P_\ell$ is the
scalar band already used above.  Put
\[
 h_m=\binom{2m+4}{4},\qquad
 d_\ell=\frac{(\ell+1)(\ell+2)(2\ell+3)}6,\qquad
 v_m=\binom{2m+2}{2}.
\]
For every finite scalar charge set $F\subset\Z^2$, define
\[
 \mathbb H_{m_x,m_p,F}=H_{m_x}\otimes H_{m_p}\otimes\C[F],\quad
 \mathbb V_{\ell_x,m_p,F}=P_{\ell_x}\otimes H_{m_p}\otimes\C[F],
\]
and retain the scalar--scalar family $W_{\ell_x,\ell_p,F}$.
The reversed mixed family is obtained by exchanging the factor types as
well as their labels.  All support comparisons use fixed unitary models $U_{a,b}$; their
intersections are taken in that common abstract model.  A specified
scalar, form or spinor realization uses its transported intertwiner.
The scalar degree $\ell$ and enriched level $m$
are independent labels; numerically equal labels need not give equal carriers.
The notation $m_x$ for scalar reciprocal-position resolution in Paper II-A
corresponds to $\ell_x$ in $\mathbb V$ here.  This is a notation translation,
with no change to that paper's Born-derived reciprocal-position spectral window.
The spaces $H_{m_x}$ and $P_{\ell_x}$ are the specified comparison families with their displayed representation actions.
\end{definition}

\begin{remark}[Reciprocal-position comparison firewall]
\label{p1b:rem:reciprocal-position-comparison-firewall}
The all-level families in this section are mathematical comparison objects.
The scalar comparison space $P_{\ell_x}$ is the finite harmonic window
$\mathcal P_x(\ell_x)=\bigoplus_{k=0}^{\ell_x}E_k^x$ on the already smooth
Born-derived position sphere.  It is not a site set, not a microscopic
position lattice, and not a second copy of the compact momentum carrier
$\mathcal H_1$.  The mixed family $P_{\ell_x}\otimes H_{m_p}$ therefore
compares differently typed spectral-resolution and momentum-carrier data; it
does not identify their finite dimensions or microscopic ontologies.
\end{remark}

\begin{theorem}[Dimensions, shells and coordinate parity at every pair]
\label{p1b:thm:mixed-dimensions}
The dimensions of the three products are, respectively,
\[
 h_{m_x}h_{m_p}|F|,\qquad d_{\ell_x}h_{m_p}|F|,\qquad
 d_{\ell_x}d_{\ell_p}|F|.
\]
For $S=U_{1,0}\oplus U_{0,1}$ one has
\[
 H_m\cong\bigoplus_{r=0}^{2m}\Sym^r S,\qquad
 \dim\Sym^rS=\binom{r+3}{3}.
\]
Thus the increment from $H_{m-1}$ to $H_m$ consists exactly of shells
$r=2m-1,2m$, for $m\ge1$; the increment from $P_{\ell-1}$ to $P_\ell$
is $U_{\ell,\ell}$ of dimension $(\ell+1)^2$.
The central sign $(-1,-1)\in\Spin(4)$ gives the two dimensions
\[
 e_m=\frac{(m+1)(m+2)(2m^2+6m+3)}6,\qquad
 o_m=\frac{m(m+1)^2(m+2)}3.
\]
Vector coordinates preserve this sign. The four product parity dimensions,
each multiplied by $|F|$, are
\[
\begin{gathered}
 e_{m_x}e_{m_p},\qquad e_{m_x}o_{m_p},\\
 o_{m_x}e_{m_p},\qquad o_{m_x}o_{m_p}.
\end{gathered}
\]
The mixed scalar--$H_1$ product has parity dimensions
$11d_\ell|F|$ and $4d_\ell|F|$.
In particular
\[
\begin{array}{c|c|c}
 \text{auxiliary level}&H_m\otimes H_1&P_\ell\otimes H_1\ \ (\ell=m)\\\hline
0&15&15\\
1&225&75\\
2&1050&210\\
3&3150&450
\end{array}
\]
before the independent scalar factor.  With momentum fixed at one,
$\dim(H_m\otimes H_1)\sim10m^4$, whereas
$\dim(P_\ell\otimes H_1)\sim5\ell^3$.
\end{theorem}
\begin{proof}
The direct-sum symmetric-power identity gives
$\Sym^r(S_+\oplus S_-)=\bigoplus_{a+b=r}U_{a,b}$.
Summing its dimensions gives $h_m$; summing $(k+1)^2$ gives $d_\ell$.
Separate the shell sum into even and odd $r$ for $e_m,o_m$.
Tensoring with $U_{1,1}$ changes $(a,b)$ by $(\pm1,\pm1)$, hence
preserves $a+b$ modulo two.  Scalar sectors have $a=b$ and even parity.
Tensor dimensions, substitution, and the leading polynomial terms give
the remaining assertions.
\end{proof}


\begin{theorem}[Operator algebras and intertwiners of the product families]
\label{p1b:thm:mixed-algebras}
Let $A=\bigoplus_{s\in I}U_s$ and $B=\bigoplus_{t\in J}U_t$ be either
of the specified multiplicity-one sphere carriers.  For the independent
$\Spin(4)_x\times\Spin(4)_p\times\T^2$ action on
$A\otimes B\otimes\C[F]$, the generated complex star algebra, its
commutant, and the full endomorphism algebra are
\[
 \bigoplus_{s,t,\chi\in F}M_{\dim U_s\dim U_t}(\C),\qquad
 \C^{|I||J||F|},\qquad M_{\dim A\dim B|F|}(\C).
\]
Consequently the first algebra has dimension
$a(A)a(B)|F|$, where
\[
 a(H_m)=\frac{(n+1)(n+2)(n+3)(n+4)(2n^2+10n+15)}{360},\quad n=2m,
\]
\[
 a(P_\ell)=\sum_{k=0}^{\ell}(k+1)^4
 =\frac{(\ell+1)(\ell+2)(2\ell+3)(3\ell^2+9\ell+5)}{30}.
\]
Every invariant Hermitian operator is an arbitrary real scalar on each
displayed irreducible block; positivity and invertibility are exactly the
corresponding conditions on those scalars.

If $A,A'$ are embedded by these sector labels in the same multiplicity-one
Hilbert direct sum, every $\Spin(4)$ intertwiner is a scalar on each common
sector.  Hence its vector-space dimension is the number of common labels
and its maximal rank is the sum of their representation dimensions.
For $H_m,H_n$ these numbers are $v_{\min(m,n)},h_{\min(m,n)}$; for
$P_\ell,H_m$ they are $\min(\ell,m)+1,d_{\min(\ell,m)}$.
The analogous product statement uses the intersection of the two joint
label sets, including scalar characters.
\end{theorem}
\begin{proof}
The joint sectors are inequivalent irreducibles of the independent product
group.  Schur's lemma and the finite-dimensional double commutant theorem
give the three algebras.  Sum their matrix dimensions; the single-factor
formulas follow by summing $(a+1)^2(b+1)^2$ on the triangle, and fourth
powers on the diagonal.  An intertwiner vanishes between inequivalent
sectors and is scalar on an equal sector, proving every count and rank.
\end{proof}

\begin{proposition}[Nesting, fusion and all finite comparison kernels]
\label{p1b:prop:mixed-nesting}
All level inclusions are the isometric sector inclusions.  Their orthogonal
projections commute, and intersection is obtained by intersection of support
sets.  In particular $H_m\cap H_n=H_{\min(m,n)}$ and
$P_\ell\cap H_m=P_{\min(\ell,m)}$.
Under a common diagonal $\Spin(4)$ action, the tensor product $H_m\otimes H_n$
has irreducible support exactly $\Lambda_{m+n}$.  Its multiplicities are
\[
 \mu_{A,B}=
 \sum_{a+b\le2m}\sum_{c+d\le2n}N_{a,c}^{A}N_{b,d}^{B},\qquad
 N_{a,c}^{A}=\mathbf1_{\{|a-c|\le A\le a+c,\ A\equiv a+c\ (2)\}}.
\]
Pointwise products of scalar bands span $P_{\ell+k}$.
Each generator-preserving restriction of the full group-action algebras
to a smaller support is the coordinate projection of their matrix-block
direct sums, with kernel exactly the omitted blocks.
\end{proposition}
\begin{proof}
The inclusions and intersections are diagonal in the common sector
decomposition.  Clebsch--Gordan gives the multiplicity formula and the
upper support bound.  Every $(A,B)$ with $A+B\le2(m+n)$ can be partitioned
as $(a,b)+(c,d)$ with the two required total bounds; the highest coupling
then realizes it.  A polynomial monomial of degree at most $\ell+k$ factors
into two monomials of degrees at most $\ell,k$, proving the scalar span
after restriction to the sphere.  Direct sums of full matrix algebras give
the final homomorphism and its kernel.
\end{proof}

\subsubsection{Exchange at unequal levels and exact invariant comparisons}

\begin{theorem}[Exchange transport, common part and minimal completion]
\label{p1b:thm:unequal-exchange}
In a common labelled Hilbert direct sum let $A,B$ be two of the sphere
carriers, $I=A\cap B$, and put $a=\dim A,b=\dim B,c=\dim I$.
Write $F^\top=\{(w,n):(n,w)\in F\}$,
$f=|F|$, $f_\cap=|F\cap F^\top|$, and
$f_\Delta=|\{n:(n,n)\in F\}|$.  Exchange is a unitary map
\[
 A_x\otimes B_p\otimes\C[F]\longrightarrow
 B_x\otimes A_p\otimes\C[F^\top].
\]
The intersection of these two carriers is
$I\otimes I\otimes\C[F\cap F^\top]$.
Their sum is the least exchange-invariant subspace containing the first,
and has dimension $2abf-c^2f_\cap$.  Its exchange sectors have dimensions
\[
 \frac{2abf-c^2f_\cap\ \pm\ cf_\Delta}{2}.
\]
If $P$ projects to the first carrier, the compressed exchange $PBP$ is
self-adjoint there and
\[
 (PBP)^2=P_I\otimes P_I\otimes P_{F\cap F^\top}.
\]
It is an involution on that carrier exactly when $A=B$ and $F=F^\top$
(for nonzero carriers).  In the enriched case $c=h_{\min(m_x,m_p)}$;
in the mixed scalar--$H_{m_p}$ case $c=d_{\min(\ell_x,m_p)}$.
At $(H_1,H_1)$ with the zero scalar character, the exchange ranks are
$120,105$.  At $(P_1,H_1)$ the carrier has dimension $75$, its exchange
completion has dimension $125$, and the completion's ranks are $65,60$.
More generally, for $P_\ell\otimes H_1$, $\ell\ge1$, and the zero
scalar character, the common part is always $P_1\otimes P_1$.
The completion has dimension $30d_\ell-25$ and exchange ranks
$15d_\ell-10,15d_\ell-15$; compressed exchange on the original
carrier has positive and negative ranks $15,10$ and kernel dimension
$15d_\ell-25$.
\end{theorem}
\begin{proof}
The flip permutes orthogonal sector-pair and character labels.  Both
orientations are retained precisely on the displayed intersection.
Every exchanged two-element orbit in the sum has trace zero.  Only equal
basis vectors in $I\otimes I$ with a diagonal scalar character contribute
to the trace, giving $cf_\Delta$.  Dimension and trace determine the two
eigenspace dimensions.  Compression kills the labels whose swapped partner
is absent and is the ordinary flip on the intersection, giving its square.
Substitute the intersection formulas and dimensions for the examples.
\end{proof}

For the corresponding labelled projectors, exchange also transports every
defined compressed operator:
\[
 B(P_{m,n}DP_{m,n})B^{-1}
 =P_{n,m}(BDB^{-1})P_{n,m}.
\]
This follows by conjugating the three factors, with the domain transported
when \(D\) is unbounded. Thus exchange is available before truncation as
a map between the two labelled presentations. In the enriched
\(H_m\otimes H_1\) family, \(m\ge1\), with no additional scalar-character
band, the minimal invariant comparison has dimension
\(30h_m-225\), where \(h_m=\binom{2m+4}{4}\).
The mixed \(P_\ell\otimes H_1\) dimensions remain the distinct values
in the theorem.


\begin{theorem}[Right-diagonal singlets and the scalar relative quotient]
\label{p1b:thm:mixed-diagonal}
Impose the common right $\SU(2)$ action on the second spin label of both
factors.  On $U_{a,b}\otimes U_{c,d}$ its invariant space vanishes unless
$b=d$, and otherwise is $V_{a/2}\otimes V_{c/2}$.
Thus, before an independent scalar factor, the invariant dimensions are
\[
 \dim(H_M\otimes H_m)^{\SU(2)_{R,\rm diag}}
 =\sum_{b=0}^{2\min(M,m)}
       \binom{2M-b+2}{2}\binom{2m-b+2}{2},
\]
\[
 \dim(P_\ell\otimes H_m)^{\SU(2)_{R,\rm diag}}
 =\sum_{k=0}^{\min(\ell,2m)}(k+1)\binom{2m-k+2}{2}.
\]
At $m=1$, the first is $20M^2+20M+6$ for $M\ge1$; the second is
$6,12,15$ for $\ell=0,1,\ell\ge2$, respectively.
If the whole diagonal $\Spin(4)$ is imposed instead, the invariant
dimensions are $v_{\min(M,m)}$ and $\min(\ell,m)+1$.

For scalar--scalar products the right-diagonal projection is exactly the
relative-quotient construction:
\[
 (P_{\ell_x}\otimes P_{\ell_p}\otimes\C[F])^{K_{\rm diag}}
 =r^*\bigl(P_{\min(\ell_x,\ell_p)}\otimes\C[F_{\rm rel}]\bigr),
\]
where $F_{\rm rel}=\{n:(n,-n)\in F\}$.  Pullback is isometric.
For the bi-invariant metric, its scalar spectral values are
\[
 k(k+2)(R_x^{-2}+R_p^{-2})
       +n^2(1,-1)G^{-1}(1,-1)^T,
\]
with multiplicity $(k+1)^2$, for $k\le\min(\ell_x,\ell_p)$ and
$n\in F_{\rm rel}$.  An independent diagonal scalar constraint in the
enriched or mixed formulas likewise multiplies their dimensions by
$|F_{\rm rel}|$.
\end{theorem}
\begin{proof}
The singlet multiplicity in $V_{b/2}\otimes V_{d/2}$ is $\delta_{bd}$.
The first labels survive as spectator representations.  Sum their
dimensions to obtain the two formulas.  Invariance under both diagonal
chiral groups additionally requires $a=c$, leaving one singlet per common
sector.  In the scalar case $a=b=k,c=d=l$, so only $k=l$ survives;
Peter--Weyl matrix coefficients of $g_xg_p^{-1}$ span exactly this space.
The scalar diagonal requires $n+w=0$.  Haar invariance gives
$\int|f(ab^{-1})|^2\,da\,db=\int|f(k)|^2\,dk$.
Restrict the product Laplacian eigenvalues to the matching labels.
\end{proof}

\begin{theorem}[One-handed action comparison and inherited traces]
\label{p1b:thm:all-level-chiral}
Under one handed $\SU(2)$,
\[
 H_m=\bigoplus_{a=0}^{2m}b_m(a)V_{a/2},\quad
 b_m(a)=\binom{2m-a+2}{2},\qquad
 P_\ell=\bigoplus_{a=0}^{\ell}(a+1)V_{a/2}.
\]
Their generated algebras are $\bigoplus_{a=0}^{2m}M_{a+1}(\C)$ and
$\bigoplus_{a=0}^{\ell}M_{a+1}(\C)$.
A generator-preserving unital star epimorphism from the latter to the
former exists exactly when $\ell\ge2m$, and is the projection onto
$a\le2m$.  A full $\Spin(4)$ action comparison instead uses the full
two-label supports of Theorem~\ref{p1b:thm:mixed-algebras}.
The inherited normalized traces are
\[
 \tau_{H_m}(A)=h_m^{-1}\sum_{a=0}^{2m}b_m(a)\tr_{a+1}A_a,\qquad
 \tau_{P_\ell}(A)=d_\ell^{-1}\sum_{a=0}^{\ell}(a+1)\tr_{a+1}A_a.
\]
The intertwiner dimension and maximal rank for this one-handed comparison
are, respectively,
\[
 \sum_{a=0}^{\min(\ell,2m)}(a+1)b_m(a),\qquad
 \sum_{a=0}^{\min(\ell,2m)}(a+1)\min\{a+1,b_m(a)\}.
\]
For $m=1$, the pairs (intertwiner dimension, maximal rank) are
$(6,1)$ at $\ell=0$, $(12,5)$ at $\ell=1$, and $(15,8)$ for every
$\ell\ge2$.  In particular the level-one scalar comparison band $P_1$ has
no generator-preserving epimorphism onto the full one-handed action
algebra of $H_1$, although the corresponding mixed auxiliary tensor-product
comparison family exists.
The comparison dimensions remain $15$ and $8$ as $\ell$ grows above two.
All spins above one lie in the kernel of the action-algebra projection.
Their omission from this comparison does not remove them from $P_\ell$.
The corner $a\le2$ has normalized trace weight $14/d_\ell$ in $P_\ell$;
in $H_M$ it has weight $(12M^2+2M+1)/h_M$, for $M\ge1$.
\end{theorem}
\begin{proof}
For fixed $a$, summing the dimensions $b+1$ of the opposite chiral
factor gives $b_m(a)$.  Distinct Casimirs separate the spins; on each spin
the action generates its full matrix algebra.  This proves the quotient
criterion and its kernel.  Tracing with the displayed multiplicities gives
the functionals.  An intertwiner on a spin is the identity tensored with an
arbitrary multiplicity-space map, proving both formulas.  Substitution at
$m=1$ gives the stated constants.  Sum the first three isotypic dimensions
for the two corner weights.
\end{proof}

\begin{proposition}[All Hopf weights and fixed-momentum charge matching]
\label{p1b:prop:all-level-hopf}
For one marked chiral Hopf circle, the weight-$q$ multiplicities are
\[
 h_m(q)=\sum_{\substack{|q|\le a\le2m\\a\equiv q\ (2)}}b_m(a),\qquad
 p_\ell(q)=(L+1)(L+|q|+1),\quad
 L=\left\lfloor\frac{\ell-|q|}{2}\right\rfloor,
\]
with value zero when the upper bound is negative.  For both chiral Hopf
circles on $H_m$, the joint weight $(q_L,q_R)$ has multiplicity
\[
 \binom{L_0+2}{2},\qquad
 L_0=\left\lfloor\frac{2m-|q_L|-|q_R|}{2}\right\rfloor\ge0,
\]
and zero otherwise.  Tensor-product weight multiplicities multiply;
diagonal-circle invariant dimensions are sums of opposite-weight products.
In particular $h_1(0)=7,h_1(\pm1)=3,h_1(\pm2)=1$, and
\[
 \dim(P_\ell\otimes H_1)^{\U(1)_{H,\rm diag}}
 =7p_\ell(0)+6p_\ell(1)+2p_\ell(2)
 =\begin{cases}15r^2+24r+7,&\ell=2r,\\
                 15r^2+36r+19,&\ell=2r+1.
   \end{cases}
\]
Thus this charge-matched dimension grows as $15\ell^2/4$, while the
full mixed dimension grows as $5\ell^3$ and the common right-$\SU(2)$
singlet dimension stabilizes at $15$.  These are three specified spaces.
Independent scalar characters remain arbitrary elements of $F$.  If an
additional equivariance identifies them with Hopf charges, their support
is the corresponding finite weight support; if both scalar diagonal and
Hopf identifications are imposed, the same opposite-weight sums apply.
\end{proposition}
\begin{proof}
The weights of $V_{a/2}$ are $-a,-a+2,\ldots,a$, each once.  Sum the
spectator multiplicities.  For two weights write
$a=|q_L|+2i,b=|q_R|+2j$; the allowed pairs satisfy $i+j\le L_0$.
Opposite charges are precisely the singlets of a diagonal circle.
Substitution into the arithmetic sums gives the two polynomials.
No equation has related the separate scalar characters to these weights
until the stated additional equivariance is imposed.
\end{proof}


\subsubsection{Complete coordinate atlases at arbitrary levels}

The following atlas concerns every Hermitian vector-coordinate tuple with
the stated equivariance.  Scalar multiplication is a particular member on
$P_\ell$.  Its coefficients are fixed by that multiplication; the larger
atlas also describes every other equivariant tuple.

\begin{theorem}[The all-level one-factor coordinate graph]
\label{p1b:thm:all-level-graph}
For $H_m$, let $G_m$ have vertex set $\Lambda_m$ and an edge between
$(a,b)$ and $(a',b')$ exactly when $a'=a\pm1,b'=b\pm1$ and both
vertices are retained.  Each vector-coordinate intertwiner space on an
edge is complex one-dimensional.  The numbers of vertices, edges,
connected components and independent cycles are
\[
 v_m=\binom{2m+2}{2},\qquad e_m^{\rm gr}=4m^2,\qquad
 c_m=\begin{cases}1&m=0,\\2&m\ge1,\end{cases}\qquad
 \beta_m=e_m^{\rm gr}-v_m+c_m.
\]
For $m\ge1$, $\beta_m=(2m-1)(m-1)$; the two connected components
are exactly the two central parities.  For $P_\ell$ the corresponding
graph is the path with $\ell+1$ vertices, $\ell$ edges and one component.

Choose normalized intertwiner bases and an orientation of every edge.
Every Hermitian coordinate tuple is one complex coefficient per edge,
with reverse coefficients determined by adjoints.  On any support $E'$
of nonzero edges, modulo sector rephasings, its exact parameter space is
\[
 (\R_{>0})^{|E'|}\times\T^{|E'|-v+c(E')},
\]
where isolated vertices count in $c(E')$.  These strata include every
vanishing-edge case.  Fixing the nonzero scalar-to-vector edge to its
normalized scalar-multiplication magnitude removes one magnitude
parameter whenever that edge exists.  It leaves $4m^2-1$ magnitudes and
$\beta_m$ phases in the all-nonzero $H_m$ stratum, recovering
$\R_{\ge0}^3$ with its boundary strata at $m=1$.
\end{theorem}
\begin{proof}
Clebsch--Gordan for tensoring by $U_{1,1}$ gives exactly the four adjacent
possibilities, each once.  Orient the edges either as $(a,b)\to(a+1,b+1)$
or $(a,b)\to(a+1,b-1)$.  With $n=2m$ there are $n(n-1)/2$ of the first
type and $n(n+1)/2$ of the second, hence $n^2$ edges.
Along a fixed shell the second type connects all its vertices, and the
first type joins each shell to the next shell of the same parity.
This proves the components and cycle count.  The diagonal scalar support
retains exactly the first-type path.  Hermiticity determines each reverse
edge.  Vertex phases act by endpoint phase differences.  Make a spanning
forest's coefficients positive; its remaining edge phases give exactly
the cycle torus, independently on every support.  The fixed scalar edge
has a fixed positive magnitude and belongs to a spanning forest.
\end{proof}

\begin{theorem}[All mixed-level product-coordinate strata]
\label{p1b:thm:all-level-product-graph}
Let $G_A,G_B$ be the preceding scalar or enriched graphs, with counts
$(v_A,e_A,c_A)$ and $(v_B,e_B,c_B)$.  For independent product-group
equivariance, all Hermitian coordinate octuples on $A\otimes B$ have
the Cartesian product graph $G_A\square G_B$.  Its counts are
\[
 v=v_Av_B,\qquad e=e_Av_B+v_Ae_B,\qquad c=c_Ac_B.
\]
Every nonzero-support stratum has the preceding magnitude-and-cycle
description using this graph.  The all-nonzero stratum has $e$ positive
magnitudes and $e-v+c$ phases.  In particular
\[
\begin{array}{c|c|c|c}
 \text{carrier}&v&e&\text{cycle phases}\\\hline
 H_M\otimes H_1\ (M\ge1)&6v_M&32M^2+12M+4&20M^2-6M+2\\
 P_\ell\otimes H_1&6(\ell+1)&10\ell+4&4\ell\\
 P_\ell\otimes P_k&(\ell+1)(k+1)&2\ell k+\ell+k&\ell k
\end{array}
\]
Thus $(H_1,H_1)$ has $48$ magnitudes and $16$ phases before any
scalar-edge normalization, while $(P_1,H_1)$ has $14$ magnitudes and
four phases in the corresponding larger atlas.

For equal factor types and equal levels with fixed exchange, the
all-nonzero exchange-compatible atlas, modulo exchange-commuting sector
phases, has
\[
 ev\ \text{magnitudes},\qquad
 ev-\frac{v(v+1)}2+\frac{c(c+1)}2\ \text{phases},
\]
where now $(v,e,c)$ are the single-factor counts.  At $H_1$ this is
$24$ magnitudes and six phases; at $P_\ell$ it is $\ell(\ell+1)$
magnitudes and $\ell(\ell-1)/2$ phases.  Every exchange-invariant
zero-edge support is likewise classified by the cycle space of its
quotient graph, retaining all vertex orbits.

Neutral coordinates which commute with both independent scalar-circle
actions preserve each character in $F$.  Allowing arbitrary character
dependence repeats these atlases independently over $F$; requiring
character-independent coordinates selects their diagonal subfamily.
Exchange acts on this family by the sphere flip and $F\to F^\top$.
\end{theorem}
\begin{proof}
An $x$-vector preserves the entire $p$-sector and joins one edge in its
$x$-sector; each such intertwiner occurs once.  The reverse statement
holds for a $p$-vector.  This gives the Cartesian graph and its counts.
Apply the spanning-forest proof to every support and substitute the
single-factor formulas.  For equal factors exchange swaps the two edge
directions, leaving $ev$ edge orbits, and has $v(v+1)/2$ vertex orbits.
The component orbits are the $c(c+1)/2$ unordered pairs of single-factor
components.  Its commuting gauge is the vertex-phase group on these
orbits; its kernel is one constant phase per component orbit.  The same
argument on an invariant support proves the boundary statement.
Scalar characters are inequivalent and neutral operators cannot mix them.
\end{proof}

\begin{theorem}[Mixed commutators and the tensor-factorwise subfamilies]
\label{p1b:thm:coordinate-square-defects}
For oriented single-factor edges $e:s\to s'$ and $f:t\to t'$, denote
the product coefficients by $x_e(t)$ and $p_f(s)$.
All cross commutators vanish exactly when every rectangle obeys the
following equation for all independent edge orientations, including
adjoint reversal of either edge:
\[
 x_e(t')p_f(s)-p_f(s')x_e(t)=0.
\]
On the all-nonzero stratum, these equations are equivalent, on each
connected component pair, to a sector-rephasing of tensor-factorwise
coordinates.  In particular the $x$-edge magnitudes are independent of
the spectator vertex within each $p$ component, the $p$-edge magnitudes
are independent within each $x$ component, and every mixed rectangle
has trivial phase holonomy.  Dependence on the spectator component may
remain.  Global factorwise coordinates impose equality across those
components as well.

If the scalar edge in each factor is normalized, the globally factorwise
enriched family has, on its all-nonzero stratum,
\[
 (4m_x^2-1)+(4m_p^2-1)\ \text{magnitudes},\qquad
 \beta_{m_x}+\beta_{m_p}\ \text{phases}
\]
for $m_x,m_p\ge1$.  At equal levels the matched exchange subfamily
has the single-factor counts.  If the position coordinates are fixed
to actual scalar multiplication on $P_\ell$, only the chosen momentum
coordinate family remains free in the globally factorwise product.
\end{theorem}
\begin{proof}
The two orders of traversing a rectangle have the same nonzero tensor
intertwiner and the two displayed scalar products as coefficients.
These are the only two-step paths contributing to that block of a cross
commutator.  Reversing $f$ and using Hermitian adjoints, together with the
forward equation, gives $|x_e(t')|=|x_e(t)|$ and
$|p_f(s')|=|p_f(s)|$ whenever all coefficients are nonzero.
The remaining equations set every rectangle phase product to one.
Choose base vertices and spanning trees in the two connected factors.
Use sector phases to transport the edge phases from the two base slices.
The rectangle identities make this transport consistent; cycles remaining
within either factor retain precisely that factor's phases.  The result
is tensor-factorwise on this component pair.  Conversely such coordinates
commute, and conjugation preserves commutators.  The parameter counts
follow from the one-factor atlas.  For zero supports the displayed
rectangle equations remain necessary and sufficient, without requiring
the componentwise factorization conclusion.
\end{proof}

\begin{theorem}[Generated operation algebras and protected momentum sectors]
\label{p1b:thm:all-level-generated-coordinates}
For a specified nonzero-support graph, the algebra generated jointly by
the independent group action and the coordinate tuple is
\[
 \bigoplus_{C\in\pi_0(G)}\End\left(\bigoplus_{v\in C}U_v\right).
\]
Its commutant is one scalar per graph component.  With all allowed edges
nonzero, $H_{m_x}\otimes H_{m_p}$ has one full matrix block on each
nonempty product parity space (four when $m_x,m_p\ge1$); $P_\ell\otimes H_1$ has exactly
$M_{11d_\ell}(\C)\oplus M_{4d_\ell}(\C)$.

For any level of the other factor, $A\otimes S_p$, with
$S_p=U_{1,0}\oplus U_{0,1}\subset H_{1,p}$, is reducing for every
coordinate tuple in the atlas.  At an enriched level $m\ge2$, the
four-dimensional subspace $S\subset H_m$ is reducing exactly when the
two allowed exterior-edge coefficients
\[
 U_{1,0}\leftrightarrow U_{2,1},\qquad
 U_{0,1}\leftrightarrow U_{1,2}
\]
vanish.  At $m=1$ these edges are absent from the graph.  On a product,
this criterion holds separately for every retained spectator sector.

On globally factorwise coordinates with nonzero protected magnitudes and
both protected subspaces reducing, normalizing the two protected quadruples
gives the dimension-$16$ complex $\Cl_{8,0}$ module of
Corollary~\ref{p1b:cor:product-clifford}.  The general mixed atlas is
governed by the preceding rectangle defects as well as the exit edges;
the factorwise Clifford construction occupies its specified subfamily.
\end{theorem}
\begin{proof}
The group action contains the full endomorphism algebra on each irreducible
vertex.  Multiplying a nonzero edge map on both sides by these matrix
algebras spans the full rectangular Hom space between its vertices.
Products along paths give all maps inside a graph component.  There are
no maps between distinct components, proving the algebra and commutant.
The all-nonzero component counts follow from parity.  The momentum odd
component at level one is exactly $S_p$; position coordinates preserve
momentum sectors and momentum coordinates preserve this component.
Clebsch--Gordan lists exactly the two exits from $S$ displayed above.
Hermiticity makes absence of the exit maps equivalent to reducingness.
For factorwise protected coordinates the earlier Clifford proof applies
without change to the two four-dimensional reducing blocks.
\end{proof}

\subsubsection{Exact spectra, boundary effects and large position resolution}

\begin{theorem}[Finite joint spectra and fixed-momentum towers]
\label{p1b:thm:mixed-spectra}
On $U_{a,b}$ let $K_m$ have the representation Casimir value
\[
 \kappa_{a,b}=\frac{a(a+2)+b(b+2)}4.
\]
For arbitrary positive $\gamma_x,\gamma_p$ and torus metric $G$, the
product operator $\gamma_x K_{m_x}+\gamma_p K_{m_p}+N^TG^{-1}N$
has, for every retained joint label, eigenvalue
\[
 \gamma_x\kappa_{a,b}+\gamma_p\kappa_{c,d}+\chi^TG^{-1}\chi
\]
and block multiplicity $(a+1)(b+1)(c+1)(d+1)$; multiplicities are
summed whenever different joint labels give the same numerical eigenvalue.
Its heat trace is the product
of the two finite Casimir sums and $\sum_{\chi\in F}e^{-t\chi^TG^{-1}\chi}$.
These statements concern the representation Casimir; a geometric bundle
Laplacian uses its separately specified connection.

On the scalar-position product choose instead
$L=\Delta_x+\gamma_p K_{m_p}+N^TG^{-1}N$.
For $m_p=1$ this gives the four towers
\[
 \frac{k(k+2)}{R_x^2}+\gamma_p\kappa+\chi^TG^{-1}\chi,
 \quad k=0,\ldots,\ell_x,\quad
 (\kappa,d_\kappa)=(0,1),(3/4,4),(3/2,4),(2,6),
\]
of block multiplicities $(k+1)^2d_\kappa$.  For any other fixed
Hermitian momentum endomorphism, replace the four Casimir entries by
its eigenvalues and their multiplicities; the same additive tower
formula holds for $\Delta_x$ plus that endomorphism.
The scalar position spectral ceiling is $\ell_x(\ell_x+2)/R_x^2$;
its first omitted level is $(\ell_x+1)(\ell_x+3)/R_x^2$.
The enriched Casimir ceiling is $m_x(m_x+1)$ and its first positive
level is $3/4$ when $m_x\ge1$.

For equal carriers, matched spectral operators, an exchange-invariant
torus metric $G$, and $F=F^\top$, if $Z_A(t)$
is the finite one-factor heat trace, the exchange heat traces are
\[
 \tfrac12\left[Z_A(t)^2\Theta_F(t)
                   \ \pm\ Z_A(2t)\Theta_{F,\Delta}(t)\right],
\]
where the second scalar sum retains $(n,n)\in F$.  At unequal levels,
the flip transports spectra between the swapped parameter-and-level
members.  On the minimal exchange completion with a common parent
one-factor operator, its ordinary heat trace is
$2Z_AZ_B\Theta_F-Z_I^2\Theta_{F\cap F^\top}$, and its exchange trace
is $Z_I(2t)\Theta_{F,\Delta}$, assuming $G$ is exchange invariant.
\end{theorem}
\begin{proof}
All specified operators are diagonal in the independent joint sector
basis.  Eigenvalues add, multiplicities multiply, and exponentials
factor.  Substitute the four $H_1$ Casimir values.  On the enriched
triangle the maximum of the convex sum $j_L(j_L+1)+j_R(j_R+1)$ is
at $(m,0)$ or $(0,m)$, giving its ceiling.
For exchange, only diagonal basis pairs contribute to its trace; their
weights square.  In the completion use inclusion--exclusion on the two
coordinate subspaces and their common part.  Matched parent spectra
ensure that all these projections commute with the operator.
\end{proof}

\begin{theorem}[Scalar coordinate radius and commutator at every cutoff]
\label{p1b:thm:scalar-cutoff-identities}
On the unit sphere let $P_\ell$ also denote its scalar band projector,
$\Pi_\ell$ the projector to the top harmonic degree, and
$C_a=P_\ell M_{x_a}P_\ell$, $a=1,\ldots,4$.
With $\mathcal L_{ab}=x_a\partial_b-x_b\partial_a$, one has
\[
 \sum_{a=1}^4 C_a^2
 =I-\frac{\ell+2}{2(\ell+1)}\Pi_\ell,
 \qquad
 [C_a,C_b]=\frac1{2(\ell+1)}\mathcal L_{ab}\Pi_\ell.
\]
For $\tau_\ell=d_\ell^{-1}\Tr$,
\[
 \tau_\ell\left(I-\sum C_a^2\right)=\frac3{2\ell+3},\qquad
 \tau_\ell\left(\sum_{a<b}[C_a,C_b]^*[C_a,C_b]\right)
 =\frac{\ell(\ell+2)}{4d_\ell}.
\]
The radius-defect norm tends to $1/2$ and, for each $a\ne b$, the
commutator norm is $\ell/[2(\ell+1)]$.  Both defects nevertheless
vanish strongly on the increasing scalar Hilbert space, and their stated
normalized traces tend to zero.  Tensoring with any fixed momentum
carrier preserves these norm and normalized-trace statements.  Radius coordinates $R_xC_a$ multiply the two operator identities by
$R_x^2$ in their quadratic terms.
\end{theorem}
\begin{proof}
For a homogeneous harmonic polynomial $f_k$ in four variables,
\[
 x_af_k=H_{k+1,a}+\frac{|x|^2}{2(k+1)}\partial_af_k
\]
is its raising/lowering decomposition.  The total squared norm of its
lowering pieces is $k/[2(k+1)]$ times $\|f_k\|^2$: integrate the
ambient gradient norm $k(2k+2)\|f_k\|^2$ and divide by $4(k+1)^2$.
Since $\sum x_a^2=1$, the raising norm is $(k+2)/[2(k+1)]$.
Only the top raising piece is removed by the band projection, proving
the radius identity.  Antisymmetrizing two top-level lowering/raising
operations cancels the second derivatives and yields the stated rotation
operator.  All lower blocks have zero commutator by the exact compression
identity.  Trace the top projector of rank $(\ell+1)^2$ and use
$\sum_{a<b}\mathcal L_{ab}^*\mathcal L_{ab}=\ell(\ell+2)I$ there.
The largest rotation weight in degree $\ell$ is $\ell$, attained by
$(x_a+ix_b)^\ell$, giving the norm.  Each fixed finite harmonic band is
eventually below the top boundary; the defects are uniformly bounded,
so density proves strong convergence.  Tensoring with a fixed identity
preserves norms and normalized traces.
\end{proof}

\begin{theorem}[Compatibility of words and all finite boundary corrections]
\label{p1b:thm:all-level-boundaries}
Fix a common vector-coordinate coefficient array on the algebraic sum
of all enriched sectors, and compress it at each level.  A word of length
$r$ applied to $H_{m_0}$ is unchanged by enriched cutoff once
$m\ge m_0+r$.  For scalar coordinate words on $P_{\ell_0}$ the analogous
bound is $\ell\ge\ell_0+r$.  For product words, count the operations
in each factor separately.  These are sharp sufficient degree bounds.

For an orthogonal projector $P$, $Q=I-P$, and operators defined on the
finite sector core, all two-step corrections are
\[
 PAP\,PBP-PABP=-PAQBP,\qquad
 [PAP,PBP]-P[A,B]P=-PAQBP+PBQAP.
\]
Longer corrections follow by inserting $P+Q$ between every adjacent
pair.  For enriched vector coordinates each two-step correction is
supported on the top two shells $a+b=2m-1,2m$; for scalar multiplication
it is supported on the top harmonic degree.  In a globally factorwise
product all cross-factor commutators vanish at every cutoff.  Increasing
only the position level removes position boundary effects on fixed
position words, while momentum truncation corrections remain those of
the fixed momentum carrier.
\end{theorem}
\begin{proof}
One enriched coordinate changes $a+b$ by at most two; one scalar
coordinate changes harmonic degree by at most one.  Iteration gives the
word bounds, and repeated raising paths attain them when their
coefficients are nonzero.  Insert $I=P+Q$ for the algebraic identities.
An enriched raising edge can leave the cutoff only from the two displayed
shells; Hermiticity gives the same support on return.  The scalar path
has just its top exit.  Tensor-factorwise compressions commute across
factors, and their omitted spaces remain factorwise as claimed.
\end{proof}

\begin{theorem}[One-factor and two-factor limits, with radius kept explicit]
\label{p1b:thm:all-level-limits}
The scalar inclusions have Hilbert completion $L^2(S^3)$; the enriched
inclusions have completion
\[
 H_\infty=\widehat{\bigoplus}_{a,b\ge0}U_{a,b}.
\]
With momentum fixed at one the completions are, respectively,
$L^2(S_x^3)\otimes H_1\otimes\C[F]$ and
$H_\infty\otimes H_1\otimes\C[F]$.  They retain their finite momentum
factor.  Letting both enriched levels grow gives
$H_\infty\widehat\otimes H_\infty\otimes\C[F]$; any cofinal sequence
of pairs gives the same completion.  Letting both levels of a mixed
family grow gives $L^2(S_x^3)\widehat\otimes H_\infty\otimes\C[F]$,
and the reversed mixed family has the swapped completion.

For a bounded parent operator $A$, its finite compressions converge
strongly to $A$ on the relevant completion, and fixed finite products
of the compressions converge strongly to the corresponding parent
products.  For the diagonal positive Casimir or scalar Laplacian,
the inclusions preserve every retained eigenvalue and the compressed
resolvents, extended by zero, converge strongly to the parent resolvent.
Their heat traces increase to the parent trace for each $t>0$.

For a fixed sphere radius and fixed finite scalar set $F$, the iterated
limit (position level first, then $t\downarrow0$) on
$P_\ell\otimes H_1$ has heat growth
\[
 15|F|\frac{\sqrt\pi R_x^3}{4t^{3/2}}
\]
for $\Delta_x$ plus any fixed finite momentum operator.  On
$H_m\otimes H_1$, with $\gamma_xK_m$ instead, it is
$60|F|/(\gamma_x^2t^2)$.
For both enriched factors growing it is
$16|F|/(\gamma_x^2\gamma_p^2t^4)$; for both scalar factors growing
it is $|F|\pi R_x^3R_p^3/(16t^3)$.  For both levels of the mixed
scalar--enriched family growing, it is
$|F|\sqrt\pi R_x^3/(\gamma_p^2t^{7/2})$, with the swapped formula
for the reversed mixed family.
At every fixed finite pair the heat trace instead tends to the finite
carrier dimension as $t\downarrow0$.

Finally, for a scalar position radius depending on its level,
$R_\ell/\ell\to L>0$, one has
\[
 \lambda_{\max}\to L^{-2},\qquad
 \lambda_1=3/R_\ell^2\to0,\qquad
 \frac{d_\ell}{2\pi^2R_\ell^3}\longrightarrow\frac1{6\pi^2L^3}.
\]
At fixed $R_x$ the same level increase instead sends the spectral
ceiling to infinity.  Thus a level and a geometric radius remain two
explicit mathematical parameters throughout the family.
\end{theorem}
\begin{proof}
Peter--Weyl proves scalar density.  The enriched completion is the
definition of the Hilbert direct sum of its exhaustive sector union.
Tensor inclusions give the fixed-factor and cofinal two-factor statements.
If $P_m\to I$ strongly, boundedness proves
$P_mAP_m\to A$ strongly; uniform operator bounds permit finite products.
The diagonal resolvent assertion follows on every eigenbasis vector and
then by its uniform bound.  Positive heat sums increase to their full
absolutely convergent sums.

For the scalar sum write $j=k+1$; then
$\sum j^2e^{-tj^2/R_x^2}\sim\sqrt\pi R_x^3/(4t^{3/2})$, by its
Gaussian Riemann sum.  For the enriched full quadrant,
\[
 \Tr e^{-t\gamma K_\infty}
 =\left[\sum_{a\ge0}(a+1)e^{-t\gamma a(a+2)/4}\right]^2
 \sim\frac4{\gamma^2t^2}.
\]
The same Riemann-sum argument gives the last asymptotic, and finite-factor
heat traces tend to their dimensions.  Products give all displayed
constants.  Finite traces and the different order of limits are immediate.
The radius assertions follow from the exact scalar ceiling, gap, and
dimension formula, without asserting any additional local geometric limit.
\end{proof}


\subsubsection{Exceptional restrictions and exact presentations at every level}

\begin{theorem}[Finite-subgroup restrictions across all levels]
\label{p1b:thm:all-level-finite-restriction}
Let $\Gamma\subset\SU(2)$ act through one chosen chiral factor and let
$\mu_a=\dim V_{a/2}^{\Gamma}$.  Its exact invariant ranks are
\[
 r_\Gamma(P_\ell)=\sum_{a=0}^{\ell}(a+1)\mu_a,\qquad
 r_\Gamma(H_m)=\sum_{a=0}^{2m}b_m(a)\mu_a.
\]
The finite averaging projector commutes with every level projection and
every Casimir function.  Independent finite actions on the two factors
multiply the ranks.  A common finite action instead has invariant rank
$|\Gamma|^{-1}\sum_{g\in\Gamma}\chi_A(g)\chi_B(g)$; its intertwiner
dimension is $|\Gamma|^{-1}\sum\overline{\chi_A(g)}\chi_B(g)$.
These formulas determine every finite-level restriction, including mixed
and equal levels, from the single-factor characters.

For $2I$ the coefficients are supplied by
\[
 \sum_{a\ge0}\mu_at^a
 =\frac{1-t^{60}}{(1-t^{12})(1-t^{20})(1-t^{30})}.
\]
Thus $P_\ell$ has only its constant invariant through $\ell=11$ and gains
13 dimensions at $\ell=12$.  For $H_m$ the acted-spin-zero contribution
is always $b_m(0)=v_m$; the first additional acted-spin contribution is
one-dimensional at $m=6$.  For each of $2T,2O,2I$ one has
$r_\Gamma(H_1)=6$.
For every fixed finite subgroup,
\[
 \frac{r_\Gamma(P_\ell)}{d_\ell}\longrightarrow\frac1{|\Gamma|},
 \qquad
 \frac{r_\Gamma(H_m)}{h_m}\longrightarrow\frac1{|\Gamma|}.
\]
Consequently independent exceptional averaging on $P_\ell\otimes H_1$
has asymptotic retained fraction $2/(5|\Gamma|)$ when the same exceptional
group acts on both sides and only $\ell$ grows.
\end{theorem}
\begin{proof}
Apply the single-chiral branching formula and take invariant subspaces.
The Haar sum is a projector and commutes with the spectral and sector
projections.  Tensor characters give the remaining exact rank formulas.
The supplied Molien series has first positive coefficient at $a=12$;
substitution gives both onsets.  Each exceptional group has no invariant
vector in spins $1/2$ or $1$, leaving only the six copies of spin zero in
$H_1$.  This also follows from the absence of degrees one and two in
their one-sphere invariant series.
For the asymptotic, average characters.  At the identity the contribution
is the full dimension.  For a noncentral group element with angle $\theta$,
$\chi_{a/2}(g)=\sin((a+1)\theta)/\sin\theta$ is bounded.  Its scalar
sum is $O(\ell^2)$ and enriched sum is $O(m^3)$.  At $g=-1$ the character
is $(-1)^a(a+1)$; pairing consecutive terms gives the same respective
orders.  Divide by $d_\ell\sim\ell^3/3$ and $h_m\sim2m^4/3$.
The fixed momentum fraction is $6/15=2/5$.
\end{proof}

\begin{theorem}[All-level bundle realization and positive exact sampling]
\label{p1b:thm:all-level-sampling}
Every $H_m$ has an explicit $\Spin(4)$-equivariant isometric realization
in smooth sections of a finite-rank homogeneous bundle over $S^3$.
One such coherent family uses
\[
 \mathcal E_m=\bigoplus_{s=0,1/2,\ldots,m}
       \Spin(4)\times_{\SU(2)_{\rm diag}}V_s,
 \qquad \rank_\C\mathcal E_m=v_m.
\]
In the global section $(q,1)$, every retained section in this realization
has component polynomials of degree at most $2m$ in unit quaternion
coordinates.  A positive spherical design of strength $t\ge4m$ therefore
gives isometric evaluation.  A multilinear expression in $r$ fields and
their conjugates, contracted with polynomial coefficient tensors of degree
at most $d$, is integrated exactly when $2rm+d\le t$.

These sufficient bounds extend factorwise to arbitrary products.  For a
scalar factor use its degree bound $\ell$ in place of $2m$; for scalar
Fourier grids use the separate no-alias bounds of
Theorem~\ref{p1b:thm:product-cubature}.  At the fixed momentum level one
the original fifteen-dimensional smooth model and its exact $2T$ sampling
may be retained unchanged while the position sampling is refined.

More generally, for any specified finite-dimensional smooth section
realization on a compact manifold and any finite collection of continuous
real pairing, operator, or interaction coefficient functions, there is a
positive finite quadrature exact on their real span and on constants.
If their real span together with $1$ has dimension $D$, at most $D$
sampling points suffice.
Exact pairing quadrature makes weighted evaluation an isometry; all
operator, action, tensor and measure transport theorems of
Section~\ref{p1b:sec:transport} then apply at every level pair.
The construction and the quadrature proof are given in
Paper~I-D, Section~\ref{p1b:app:all-level-sampling}.
\end{theorem}

\begin{corollary}[Sharp exceptional comparison ranges and fixed-factor refinement]
\label{p1b:cor:all-level-design-comparison}
For $r=2,3,4$, impose both a one-handed generator-preserving quotient
$\mathcal B(P_\ell)\to\mathcal B(H_m)$ and exact uniform $\Gamma$
integration of all $r$-fold scalar-band products.  Writing $L_\Gamma(r)$
for the sharp scalar frontier in Theorem~\ref{p1b:thm:design}, the exact
range is
\[
 2m\le\ell\le L_\Gamma(r).
\]
Its largest enriched level is
\[
\begin{array}{c|ccc}
 &2T&2O&2I\\\hline
 r=2&1&1&2\\
 r=3&0&1&1\\
 r=4&0&0&1
\end{array}
\]
with $m=0$ allowed.  At $m\ge1$, quartic $2I$ comparison uniquely gives
$(m,\ell)=(1,2)$.  For $m_p=1$ and arbitrarily larger position degree,
the quotient onto the same comparison algebra persists, while the full
position field requires its larger exact integration rule.

On a product, a fixed momentum quadrature never repairs a failed position
moment: test a position polynomial tensored with the scalar unit in every
other factor.  Thus a fixed $2I$ position rule ceases to be isometric for
all scalar fields at $\ell=6$ and ceases to integrate all quartics at
$\ell=3$.  Positive exact quadrature at every larger finite level is
available by Theorem~\ref{p1b:thm:all-level-sampling}.
\end{corollary}
\begin{proof}
Combine the necessary and sufficient quotient condition $\ell\ge2m$
with the already proved sharp scalar frontiers, then take the integer
part of $L_\Gamma(r)/2$.  Testing a function constant in the other factors
reduces product exactness to the failed one-factor identity.  The $2I$
frontiers for quadratic and quartic products are five and two.
\end{proof}

The families developed here specify which conclusions survive a fixed
momentum factor, which grow with position resolution, and which require a
stated projection or exchange completion.  The all-level formulas retain
every nonnegative pair, arbitrary finite scalar character sets, and every
coordinate support stratum.  The count, graph, character, spectral and
sampling descriptions provide exact rules for any further member of these
named families; they introduce no preferred cutoff or metric scale.


\subsection{Exact transport of the finite mathematical families}
\label{p1b:sec:transport}

\begin{theorem}[Algebras, spectra and constraints on an isometric image]
\label{p1b:thm:image-transport}
Let \(T:W\to V\) be an isometry of finite-dimensional Hermitian
spaces and \(P=TT^*\). On \(V_0=\im T\), the map
\[
 \mathcal T(A)=TAT^*
\]
is a unital star-algebra isomorphism
\(\End(W)\to\End(V_0)\). It preserves products, adjoints,
commutators, positivity, rank, norm, trace, determinant, spectrum
and spectral functional calculus. For a unital star-subalgebra
\(\mathfrak A\subset\End(W)\) and an operator family \(A_j\),
\[
 \mathcal T(\mathfrak A')=\mathcal T(\mathfrak A)'_{\End(V_0)},
 \qquad
 T\bigcap_j\ker A_j=\bigcap_j\ker(\mathcal T(A_j)|_{V_0}).
\]
In the larger host \(\End(V)\), the unit of the image algebra is
\(P\), and its commutant is
\[
 \mathcal T(\mathfrak A')\oplus\End(V_0^\perp).
\]
\end{theorem}
\begin{proof}
Insert \(T^*T=I\) between products; on \(V_0\), \(TT^*=I\).
The inverse is conjugation by \(T^*\). Transport an orthonormal
basis and a spectral resolution for the listed invariants.
The commutant and joint-kernel equations follow by conjugation.
An ambient commuting operator commutes with \(P\), because
\(I_W\in\mathfrak A\), hence is block diagonal for
\(V_0\oplus V_0^\perp\). Its first block is the transported
commutant and its second block is arbitrary.
\end{proof}

\begin{theorem}[Actions, reductions and moving presentations]
\label{p1b:thm:action-transport}
For a \(C^k\) real-valued function \(S\) on an open subset of the
underlying real space of \(W\), put \(\widetilde S=S\circ T^*\)
on its image. Then
\[
 D^j\widetilde S(Tq)[Tu_1,\ldots,Tu_j]
       =D^jS(q)[u_1,\ldots,u_j],\qquad 0\le j\le k.
\]
Stationarity, Hessians in the transported metric, symmetry identities,
and Schur complements with an invertible interior block transport
exactly. For \(H=H^*\),
\[
 T e^{-itH}T^*=e^{-it\mathcal T(H)}
\]
as operators on \(V_0\). For smooth isometries \(T(s)\) onto
varying images, define on their image bundle
\[
 \nabla_s^T v=T(s)\partial_s(T(s)^*v).
\]
Then \(\nabla_s^T(T\psi)=T\partial_s\psi\); replacing ordinary
derivatives by this connection gives the exact moving-image law.
\end{theorem}
\begin{proof}
The linear chain rule gives all derivatives. It transports an
invariance equation before differentiation, hence its associated
differential identities. Transport both subspaces of a Hessian
block decomposition and cancel \(T_I^*T_I\) in the inverse block
to obtain the Schur identity. The matrix exponential follows
from the convergent power series on \(V_0\).
For the moving formula differentiate \(T^*T=I\) in
\(T\partial_s(T^*T\psi)\).
\end{proof}

\begin{corollary}[Tensor, state and invariant-subspace transport]
\label{p1b:cor:tensor-transport}
Isometric image transport preserves the pairing
\(\Tr(\rho A)\), positive trace-one matrices, and effects
\(0\le A\le I\), with the identity interpreted on the image.
It extends to every finite tensor, symmetric-power and exterior-power
carrier. If \(U(g)\) is a unitary representation of a compact group,
then its Haar projector \(P_G=\int U(g)\,\dd g\), its invariant
subspace, and every finite sum of those tensor carriers transport
exactly.
\end{corollary}
\begin{proof}
Use cyclicity of trace and \(T^*T=I\).
The tensor power \(T^{\otimes n}\) is an isometry and commutes
with the permutation action, so it restricts to symmetric and
alternating subspaces. Transport the group matrices under the
Haar integral. Its range is the invariant subspace, giving the
last assertion.
\end{proof}

\begin{theorem}[Measure and blocking transport]
\label{p1b:thm:blocking-transport}
Let \((X_n,\nu_n)\) be standard Borel probability spaces with
measurable coarse maps \(\pi_{n,m}\), with their reference
measures compatible under pushforward. Let
\(F_n:X_n\to\widetilde X_n\) be bimeasurable bijections.
Transport
\[
 \widetilde\nu_n=F_{n*}\nu_n,\quad
 \widetilde S_n=S_n\circ F_n^{-1},\quad
 \widetilde\pi_{n,m}=F_m\pi_{n,m}F_n^{-1}.
\]
For positive integrable weights \(e^{-S_n}\), conditional blocking
\(\mathcal B_{n\to m}S=-\log
 \mathbb E_{\nu_n}(e^{-S}\mid\pi_{n,m})\) satisfies
\[
 \widetilde{\mathcal B}_{n\to m}\widetilde S_n
       =(\mathcal B_{n\to m}S_n)\circ F_m^{-1}
\]
almost everywhere. Normalized integrals, total variation distances
and essential action oscillations are preserved. Consequently
exact presentation transport contributes zero blocking discrepancy.
\end{theorem}
\begin{proof}
Disintegrate \(\nu_n\) over \(\pi_{n,m}\) and push each fibre
kernel forward by \(F_n\). Its integral of the transported weight
is exactly the original conditional integral evaluated at
\(F_m^{-1}\widetilde x\). Negative logarithms give the identity.
The change-of-variables formula preserves normalized integrals.
Bimeasurable bijections preserve measurable test sets and essential
ranges, proving the variation-distance and oscillation assertions.
\end{proof}


\section{Complete product-geometry and harmonic atlas}
\label{p1b:sec:organized-atlas}

\hypertarget{p1b:sec:g02}{%
\subsection{Coordinate ancestry and retained geometric squares}\label{p1b:sec:g02}}


\begin{proposition}[Coordinate ancestry and the retained product square]
\label{p1b:prop:g02-product}
For the normalized fifteen-dimensional retained sphere model, the
intrinsic tangential and ambient cone constructions define two coordinate
classes. Under the grammar that applies either construction independently
to each labelled factor, their product has four ancestry classes.
For radii \(R_x,R_p>0\), sector projections \(\Pi_s\), and the flat scalar
torus metric \(G>0\), the compressed geometric square is
\[
 K_x\otimes I+I\otimes K_p+\Delta_{T_G^2},\qquad
 K_i=R_i^{-2}\sum_s\lambda_s\Pi_s,\quad
 (\lambda_0,\lambda_v,\lambda_L,\lambda_R,\lambda_+,\lambda_-)
 =(0,3,4,4,9/4,9/4).
\]
The torus factor is omitted when absent. Compression of the square here
means \(P D^2P\), with its stated smooth realization; it is not defined
by squaring \(PDP\). The scalar band and the enriched carrier have their
separate realizations. Isometric sampling transports the displayed
operators and their products to the sampled image.
\end{proposition}

The following normalization, coordinate-block and spectral calculations give the proof, including the four factor choices and their exchange.

\hypertarget{p1b:der:g02:s2}{%
\subsubsection{One-factor carrier and geometric normalization}\label{p1b:der:g02:s2}}

Fix an oriented round sphere of radius \(R>0\). Use the complex multiplicity-one carrier

\[
H_1=U_0\oplus U_v\oplus U_L\oplus U_R\oplus S_+\oplus S_-,
\]

with sector labels and dimensions

\[
\begin{array}{c|cccccc}
s&0&v&L&R&+&-\\\hline
(j_L,j_R)&(0,0)&(\tfrac12,\tfrac12)&(1,0)&(0,1)&(\tfrac12,0)&(0,\tfrac12)\\
d_s&1&4&3&3&2&2.
\end{array}
\]

The symbols \(+\) and \(-\) in this table label the two lowest spinor sectors. Their particular assignment to signed Dirac eigenvalues can be reversed without affecting any square or fingerprint below.

Write \(V_R=\operatorname{Vol}(S_R^3)\) and \(q_a=x_a/R\). The retained smooth model is

\[
e_0=V_R^{-1/2},\qquad
u_a=\frac{2}{\sqrt{3V_R}}\,dx_a,\qquad
w_{ab}=\sqrt{\frac2{V_R}}\,dx_a\wedge dx_b\quad(a<b),
\]

together with the two lowest spinor eigenspaces. Let \(J_R:H_1\to\mathcal L_R\) be the equivariant unitary identification, defined up to phases on the six sectors, and let \(P_R\) project onto \(\mathcal L_R\).

To verify the normalization directly, use

\[
\langle dx_a,dx_b\rangle=\delta_{ab}-q_aq_b,
\qquad \int q_aq_b\,d\mu_H=\frac14\delta_{ab}.
\]

The determinant formula for the two-form pairing has cancelling quartic terms. Therefore

\[
\int\langle dx_a,dx_b\rangle d\mathrm{vol}
=\frac{3V_R}{4}\delta_{ab},
\] \[
\int\langle dx_a\wedge dx_b,dx_c\wedge dx_d\rangle d\mathrm{vol}
=\frac{V_R}{2}(\delta_{ac}\delta_{bd}-\delta_{ad}\delta_{bc}).
\]

Different form degrees and the spinor bundle are orthogonal. The displayed modes give the required \(1+4+6+4=15\) orthonormal carrier.

The normalized scalar-to-vector coordinate edge is \(R/2\). Indeed, on the scalar harmonic realization, the normalized first mode is \(2x_a/(R\sqrt{V_R})\), and multiplication by \(x_a\) sends the normalized constant to \(R/2\) times that mode. This fixes the common coordinate scale when comparing the form/spin and ambient realizations.

\textbf{Source input:} Paper I, \path{p1:thm:canonical-retained-smooth-model}; Paper I-C, result~\ref{p1b:thm:full-sampling}.

\hypertarget{p1b:der:g02:s3}{%
\subsubsection{Deriving the two one-factor coordinate points}\label{p1b:der:g02:s3}}

\hypertarget{p1b:der:g02:s3:sub1}{%
\paragraph{G02.1. Intrinsic tangential coordinates}\label{p1b:der:g02:s3:sub1}}

Let

\[
c(\xi)=\bigl(\varepsilon(\xi)+\iota_{\xi^\sharp}\bigr)\oplus c_S(\xi)
\]

be the Hermitian metric Clifford symbol. Define its retained intrinsic coordinates by

\[
X_a^D=\frac R{\sqrt3}\,J_R^*P_Rc(dx_a)P_RJ_R.
\]

These formulas select a representative of the source's phase class. The harmless first-order sign convention is reconciled in Paper~I-D, Section~\ref{p1b:app:g02:1}.

On the form modes,

\[
c(dx_a)e_0=\frac{\sqrt3}{2}u_a,
\] \[
P_Rc(dx_a)u_b
=\frac{\sqrt3}{2}\delta_{ab}e_0
 +\sqrt{\frac23}\,w_{ab},\qquad w_{ba}=-w_{ab}.
\]

The first term in the second identity follows by averaging \(\delta_{ab}-q_aq_b\); the second follows by inserting the normalization of \(w_{ab}\). Hermiticity determines the reverse edges. Consequently the scalar edge is \(a=R/2\) and each raw exterior edge has coefficient

\[
b_D=\frac{\sqrt2R}{3}.
\]

The spin coefficient can be obtained without assuming a reduced Clebsch--Gordan normalization. In a left spin frame, use normalized modes represented by constant two-spinors and by \(q^{-1}\) times constant two-spinors. Write

\[
q=\sum_{a=1}^4q_a\tau_a,\qquad
\tau_1=i\sigma_1,\quad\tau_2=i\sigma_2,\quad
\tau_3=i\sigma_3,\quad\tau_4=I_2.
\]

Let \(P_{\rm const}^{\rm sp}\) and \(P_q^{\rm sp}\) denote these two spinor-space projections in the chosen frame. The retained cross block of scalar multiplication is

\[
P_q^{\rm sp}M_{x_a}P_{\rm const}^{\rm sp}=\frac R4\tau_a,
\]

up to the sector phase convention, since \(\int q_aq_b\,d\mu_H=\delta_{ab}/4\). The gap between the two lowest signed Dirac eigenvalues is \(3/R\). Thus the retained commutator block has magnitude \(3/4\), and the common \(R/\sqrt3\) coordinate normalization gives

\[
\varrho_D=\frac{\sqrt3R}{4}.
\]

Equivalently, on \(S_+\oplus S_-\), \(X_a^D=\varrho_D\gamma_a\) for Hermitian matrices satisfying \(\{\gamma_a,\gamma_b\}=2\delta_{ab}I\).

Now sum the outgoing squared norms. The scalar sees four edges of magnitude \(a\); each vector sees one scalar edge and three exterior edges; a two-form has two contraction channels; each spinor sees four Clifford directions. Hence

\[
Q_D:=\sum_a(X_a^D)^2
=\operatorname{diag}_s\left(4a^2,\ a^2+3b_D^2,\ 2b_D^2,
 2b_D^2,\ 4\varrho_D^2,\ 4\varrho_D^2\right),
\]

or

\[
\boxed{R^{-2}Q_D=(1,\;11/12,\;4/9,\;4/9,\;3/4,\;3/4).}
\]

The raw completion tuple is therefore

\[
\boxed{(r_{1,0},r_{0,1},\varrho_{\rm sp})_D
=R(\sqrt2/3,\sqrt2/3,\sqrt3/4).}
\]

For the raw wedge map \(T(v\otimes u)=v\wedge u\), \(TT^*=2I\). Replacing \(T\) by its coisometric normalization \(T/\sqrt2\) changes the written coefficient to \(2R/3\), while the coordinate operators and \(Q_D\) remain the same.

\hypertarget{p1b:der:g02:s3:sub2}{%
\paragraph{G02.2. Ambient cone coordinates}\label{p1b:der:g02:s3:sub2}}

At fixed radius, an ambient constant covector decomposes as

\[
e^a=dx_a+q_a\,dr.
\]

The intrinsic construction removes its normal part. The cone construction retains the full parallel ambient Clifford symbol, on

\[
H_1\cong\Lambda^{\le2}_{\mathbb C}(\mathbb R^4)^*\oplus S.
\]

With \(P_{\le2}\) the exterior-degree projection, set

\[
X_a^C=\frac R2\bigl(P_{\le2}(\varepsilon_a+\iota_a)P_{\le2}
                    \oplus\gamma_a\bigr).
\]

The scalar-edge condition fixes the multiplier \(R/2\). The common ambient metric Clifford relation fixes the same multiplier on the exterior and spin summands.

On degrees zero and one all four outgoing Clifford channels survive. On degree two, only the two contractions survive compression. All four spin channels survive. Therefore

\[
\boxed{R^{-2}Q_C=(1,\;1,\;1/2,\;1/2,\;1,\;1),}
\] \[
\boxed{(r_{1,0},r_{0,1},\varrho_{\rm sp})_C=R(1/2,1/2,1/2).}
\]

These fingerprints prove that the intrinsic and cone points are inequivalent under sector rephasings. Their scalar edges agree, so an extra overall rescaling cannot identify the normalized points either.

\textbf{Scope of the two-point result.} The source's primitive grammar admits exactly the tangent-symbol anchor and the parallel ambient-cone anchor, followed by the stated compression and scalar normalization. At a sphere point, the stabilizer decomposition \(V^*=W^*\oplus\mathbb R\) leaves a one-dimensional equivariant map to \(T^*S^3\cong W^*\). On the parallel ambient branch, a constant equivariant endomorphism of the irreducible real vector representation is scalar. These observations give one normalized class per anchor. The differing fingerprints make the two images distinct. This proves exhaustion of that declared grammar; enlarged grammars and the continuous completion cone retain their larger domains.

\textbf{Source inputs:} Paper I, \path{p1:thm:intrinsic-tangential-coordinate-class}, \path{p1:rem:intrinsic-edge-normalization-conversion}, \path{p1:thm:metric-cone-clifford-coordinate-class}, \path{p1:cor:metric-cone-quadratic-fingerprint}, \path{p1:def:primitive-round-symbol-ancestry-grammar}, \path{p1:thm:primitive-round-symbol-ancestry-exhaustion}.

\hypertarget{p1b:der:g02:s4}{%
\subsubsection{The four product ancestry choices and exchange}\label{p1b:der:g02:s4}}

\hypertarget{p1b:der:g02:s4:sub1}{%
\paragraph{G02.3. Factorwise product theorem}\label{p1b:der:g02:s4:sub1}}

Take \(R_x,R_p>0\), a nonempty finite scalar-character set \(F\subset\mathbb Z^2\), and

\[
H=H_{1,x}\otimes H_{1,p}\otimes\mathbb C[F],
\qquad\dim H=225|F|.
\]

The torus characters here are independent of the Hopf weights. For \((\eta_x,\eta_p)\in\{D,C\}^2\), define

\[
\mathsf X_a=X_a^{\eta_x}(R_x)\otimes I\otimes I_F,
\qquad
\mathsf P_b=I\otimes X_b^{\eta_p}(R_p)\otimes I_F.
\]

Then:

\begin{enumerate}
\def\labelenumi{\arabic{enumi}.}
\tightlist
\item
  All cross commutators vanish: \([\mathsf X_a,\mathsf P_b]=0\).
\item
  Their factor quadratic operators are \[
  \mathsf Q_x=Q_{\eta_x}(R_x)\otimes I\otimes I_F,
  \qquad
  \mathsf Q_p=I\otimes Q_{\eta_p}(R_p)\otimes I_F.
  \]
\item
  The joint fingerprint on the sector \(U_s\otimes U_t\otimes\mathbb C\chi\) is \[
  \bigl(R_x^{-2}\mathsf Q_x,R_p^{-2}\mathsf Q_p\bigr)
  =\bigl(q_{\eta_x}(s),q_{\eta_p}(t)\bigr),
  \] with sector multiplicity \(d_sd_t\). Equal numerical pairs acquire the sum of the corresponding sector multiplicities.
\item
  The factorwise primitive product grammar has exactly four inequivalent labelled outputs: \[
  \boxed{(D,D),\quad(D,C),\quad(C,D),\quad(C,C).}
  \]
\end{enumerate}

\textbf{Proof.} Tensor factors commute, proving the first two assertions. The one-factor sector tables give the third. The product of the two two-element primitive images gives at most four outputs. On the joint vector--scalar sector, \(\mathsf Q_x/R_x^2\) distinguishes the first ancestry label; on scalar--vector, \(\mathsf Q_p/R_p^2\) distinguishes the second. Every product-group-equivariant sector rephasing commutes with these block-scalar fingerprints, even when the rephasing does not factor into separate one-sphere phases. Hence none of the four outputs is duplicated. \(\square\)

All allowed one-factor edges are nonzero at these four points. Their placement in the preceding atlas's larger product atlas is the globally factorwise, character-independent subfamily. The four-point theorem does not classify the 48-edge atlas, arbitrary spectator-dependent coefficients, or independently chosen ancestry on each scalar-character block.

\hypertarget{p1b:der:g02:s4:sub2}{%
\paragraph{Exchange}\label{p1b:der:g02:s4:sub2}}

Let \(F^\top=\{(w,n):(n,w)\in F\}\). Exchange transports

\[
(\eta_x,\eta_p;R_x,R_p;F)
\longmapsto(\eta_p,\eta_x;R_p,R_x;F^\top),
\]

interchanging the coordinate quadruples. This is transport between the corresponding parameter-labelled objects. When radii, common-model data, and scalar support match, \(R_x=R_p\) and \(F=F^\top\), it is an involution on the fixed carrier and family.

Then \((D,D)\) and \((C,C)\) are fixed ancestry points, while \((D,C)\) and \((C,D)\) are exchanged. Forgetting the factor labels gives three orbits: two singletons and one mixed pair. The mixed pair is not an additional fixed point of the labelled family.

\textbf{Product inputs:} Paper I-C, result~\ref{p1b:thm:all-level-product-graph}, result~\ref{p1b:thm:coordinate-square-defects}, and result~\ref{p1b:thm:unequal-exchange}.

\hypertarget{p1b:der:g02:s5}{%
\subsubsection{Retained Levi--Civita square and its product spectrum}\label{p1b:der:g02:s5}}

\hypertarget{p1b:der:g02:s5:sub1}{%
\paragraph{G02.4. The one-factor square and scalar projector}\label{p1b:der:g02:s5:sub1}}

Use the specified Levi--Civita first-order representative on forms plus spinors. Choose the exterior sign convention \(-i(d-d^\dagger)\) so that \(i[\mathscr D,M_f]=c(df)\); its square agrees with that of the \(+i(d-d^\dagger)\) exterior convention in Paper I. Define

\[
K_R=J_R^*\mathscr D_{\rm LC}^{\,2}J_R,
\qquad A_R=R^2K_R.
\]

On forms, this square is the Hodge Laplacian. Constants have eigenvalue zero. Since \(\Delta x_a=3x_a/R^2\) and \(d\) commutes with the Hodge Laplacian, \(dx_a\) has eigenvalue \(3/R^2\). For the two-forms, the round Hessian identity \(\nabla^2x_a=-x_ag/R^2\) gives

\[
d^\dagger(dx_a\wedge dx_b)=\frac2{R^2}(x_a\,dx_b-x_b\,dx_a).
\]

They are closed, and applying \(d\) to the right side gives \(4\,dx_a\wedge dx_b/R^2\). This directly proves their eigenvalue \(4/R^2\). The retained spinors have signed Dirac eigenvalues \(\pm3/(2R)\). Thus

\[
\boxed{A_R=(0,\;3,\;4,\;4,\;9/4,\;9/4).}
\]

For reference, the complete one-factor data are:

\begin{longtable}[]{@{}lrrrr@{}}
\toprule\noalign{}
Sector & Dimension & \(Q_D/R^2\) & \(Q_C/R^2\) & \(R^2K_R\) \\
\midrule\noalign{}
\endhead
\bottomrule\noalign{}
\endlastfoot
\((0,0)\) & 1 & 1 & 1 & 0 \\
\((1/2,1/2)\) & 4 & \(11/12\) & 1 & 3 \\
\((1,0)\) & 3 & \(4/9\) & \(1/2\) & 4 \\
\((0,1)\) & 3 & \(4/9\) & \(1/2\) & 4 \\
\((1/2,0)\) & 2 & \(3/4\) & 1 & \(9/4\) \\
\((0,1/2)\) & 2 & \(3/4\) & 1 & \(9/4\) \\
\end{longtable}

The scalar line is exactly the kernel. An explicit polynomial for its geometric projector is

\[
\boxed{\Pi_{\rm sc}
=\left(I-\frac49A_R\right)
 \left(I-\frac13A_R\right)
 \left(I-\frac14A_R\right).}
\]

Each nonzero eigenvalue is a root of one factor; the polynomial equals one at zero. This proves the formula and the rank-one assertion.

This square belongs to the explicit Levi--Civita representative. Changing its zero-order term can change the square while leaving the symbol-derived coordinates unchanged. The cone coordinate class by itself likewise does not specify a first-order operator whose square is \(K_R\). All four ancestry choices can be equipped with this common, explicitly declared intrinsic LC square on the same retained carrier.

\hypertarget{p1b:der:g02:s5:sub2}{%
\paragraph{Relation to the representation Casimir}\label{p1b:der:g02:s5:sub2}}

Let \(\mathcal C=C_L+C_{\rm right}\) denote the representation Casimir, with sector eigenvalues \(0,3/2,2,2,3/4,3/4\). If \(\Pi_{\rm sp}\) projects onto \(S_+\oplus S_-\), then inspection of the six sectors gives

\[
\boxed{A_R=2\mathcal C+\frac34\Pi_{\rm sp}.}
\]

This bridges the joint Casimir spectrum to the geometric spectrum without identifying two different operators. On a two-factor carrier, \(A_x+A_p\) cannot be recovered from \(\mathcal C_x+\mathcal C_p\) alone: scalar--vector and spin--spin both have summed Casimir \(3/2\), but their summed LC values are \(3\) and \(9/2\), respectively. The joint sector labels, or the displayed spin-projector correction, retain the necessary information.

On this retained one-factor spectrum the spin projector, and
hence \(A_R\), is a cubic polynomial in \(\mathcal C\);
Paper~I-D, Section~\ref{p1b:app:g02:casimir-polynomial} gives the explicit
coefficients.

\hypertarget{p1b:der:g02:s5:sub3}{%
\paragraph{G02.5. Product square, including the scalar torus}\label{p1b:der:g02:s5:sub3}}

For a positive-definite scalar-torus metric \(G\), with angles of period \(2\pi\), put

\[
T_{G,F}e_\chi=(\chi^TG^{-1}\chi)e_\chi.
\]

Define the positive second-order product operator

\[
\mathcal K=K_{R_x}\otimes I\otimes I_F
 +I\otimes K_{R_p}\otimes I_F
 +I\otimes I\otimes T_{G,F}.
\]

This is the retained tensor sum of the specified one-factor squares and the scalar Laplacian. The definition does not assume an ungraded tensor sum of first-order Dirac operators has this square.

On each retained \((s,t,\chi)\), its eigenvalue and multiplicity are

\[
\boxed{\lambda_{s,t,\chi}
=\frac{a_s}{R_x^2}+\frac{a_t}{R_p^2}+\chi^TG^{-1}\chi,
\qquad \operatorname{mult}_{s,t,\chi}=d_sd_t,}
\]

where \(a_s=(0,3,4,4,9/4,9/4)\). The joint projectors are

\[
\Pi_s\otimes\Pi_t\otimes|e_\chi\rangle\langle e_\chi|;
\]

numerically coincident eigenvalues use their sums.

In particular,

\[
\boxed{\Pi_{\ker\mathcal K}
=\Pi_{{\rm sc},x}\otimes\Pi_{{\rm sc},p}\otimes\Pi_{\chi=0}.}
\]

The kernel has dimension one if \(0\in F\), and zero otherwise. If the torus kinetic term is omitted, the kernel instead has dimension \(|F|\). These facts follow because every summand of the eigenvalue is nonnegative, and vanishing forces every summand to vanish.

The heat trace is also exact:

\[
\operatorname{Tr}e^{-t\mathcal K}
=Z_{R_x}(t)Z_{R_p}(t)\sum_{\chi\in F}e^{-t\chi^TG^{-1}\chi},
\] \[
Z_R(t)=1+4e^{-9t/(4R^2)}+4e^{-3t/R^2}+6e^{-4t/R^2}.
\]

At equal radii, \(F=\{0\}\), the eigenvalues of \(R^2\mathcal K\) and their full multiplicities are

\begin{center}
\begin{tabular}{@{}lrrrrr@{}}
\toprule
Eigenvalue & $0$ & $9/4$ & $3$ & $4$ & $9/2$ \\
Multiplicity & 1 & 8 & 8 & 12 & 16 \\
\midrule
Eigenvalue & $21/4$ & $6$ & $25/4$ & $7$ & $8$ \\
Multiplicity & 32 & 16 & 48 & 48 & 36 \\
\bottomrule
\end{tabular}
\end{center}

They sum to 225. The lowest positive eigenvalue on the sphere product is \(\min\{9/(4R_x^2),9/(4R_p^2)\}\).

Exchange transports this operator with its radii, torus metric, and scalar support. It commutes with exchange on the fixed product when the sphere data match, \(F=F^\top\), and \(J^TGJ=G\), where \(J\) is the matrix swapping the two torus coordinates.

These identities use the round one-factor Dirac, Hodge and harmonic
representation results of Paper~I-A (Corollary~\ref{p1:cor:canonical-retained-dirac-square}
and Theorems~\ref{p1:thm:hodge-one-form-spectrum},
\ref{p1:thm:representation-harmonic-package}) and the reciprocal tensor
spectra of Theorem~\ref{p1b:thm:mixed-spectra}.

\hypertarget{p1b:der:g02:s6}{%
\subsubsection{Hopf comparison and its exact product kernels}\label{p1b:der:g02:s6}}

\hypertarget{p1b:der:g02:s6:sub1}{%
\paragraph{G02.6. One-factor comparison}\label{p1b:der:g02:s6:sub1}}

Let \(N_R\) be the integer-weight right Hopf generator and set

\[
E_R=R\sqrt{K_R}=\sqrt{A_R}.
\]

The right weights in \(V_{j_L}\otimes V_{j_R}\) are \(-2j_R,-2j_R+2,\ldots,2j_R\), each with multiplicity \(2j_L+1\). In the retained sector order:

\begin{longtable}[]{@{}lrl@{}}
\toprule\noalign{}
Sector & \(E_R\) & Right weights, with their multiplicities \\
\midrule\noalign{}
\endhead
\bottomrule\noalign{}
\endlastfoot
scalar & 0 & \(0\), once \\
vector & \(\sqrt3\) & \(-1,+1\), twice each \\
left bivector & 2 & \(0\), three times \\
right bivector & 2 & \(-2,0,+2\), once each \\
left spinor & \(3/2\) & \(0\), twice \\
right spinor & \(3/2\) & \(-1,+1\), once each \\
\end{longtable}

Since \(K_R\) is Spin(4)-invariant, \(E_R\) commutes with \(N_R\). The table proves \(E_R\ge|N_R|\). Thus, for \(\epsilon\in\{+1,-1\}\),

\[
B_{R,\epsilon}=\frac{E_R+\epsilon N_R}{\sqrt2}\ge0,
\] \[
\boxed{\ker B_{R,\epsilon}
=U_0\oplus(U_R)_{q=-2\epsilon},\qquad
\dim\ker B_{R,\epsilon}=2.}
\]

The notation \(B_{R,\epsilon}\) is used here to distinguish these positive operators from the exchange operator and from the paper's other uses of \(N_\pm\).

\hypertarget{p1b:der:g02:s6:sub2}{%
\paragraph{Product theorem}\label{p1b:der:g02:s6:sub2}}

On the two-factor carrier, suppress identity tensor factors and define

\[
E_\Sigma=E_x+E_p,\qquad N_\sigma=N_x+\sigma N_p,
\qquad \sigma\in\{+1,-1\}.
\]

Then

\[
\boxed{E_\Sigma\ge|N_\sigma|,\qquad
B_{\sigma,\epsilon}:=\frac{E_\Sigma+\epsilon N_\sigma}{\sqrt2}\ge0.}
\]

More precisely,

\[
B_{\sigma,\epsilon}
=B_{x,\epsilon}\otimes I+I\otimes B_{p,\epsilon\sigma}.
\]

For positive operators \(A,B\) on finite-dimensional factors,

\[
\ker(A\otimes I+I\otimes B)=\ker A\otimes\ker B.
\]

To prove this, choose separate orthonormal eigenbases: a joint eigenvalue is \(\alpha+\beta\) with \(\alpha,\beta\ge0\), and it vanishes exactly when both vanish. It follows that

\[
\boxed{\ker B_{\sigma,\epsilon}
=(U_{0,x}\oplus(U_{R,x})_{q=-2\epsilon})
 \otimes(U_{0,p}\oplus(U_{R,p})_{q=-2\epsilon\sigma})
 \otimes\mathbb C[F],}
\] \[
\boxed{\dim\ker B_{\sigma,\epsilon}=4|F|.}
\]

This derives the product kernel count, including its actual subspace. The inequality can alternatively be read directly in the joint eigenbasis from

\[
e_x+e_p\ge|q_x|+|q_p|\ge|q_x+\sigma q_p|.
\]

At matched exchange data, exchange preserves the \(\sigma=+1\) operators and swaps the two \(\epsilon\) values when \(\sigma=-1\).

The absolute comparison has

\[
\dim\ker(E_\Sigma-|N_\sigma|)=7|F|.
\]

Indeed, its zero space is the sum of the two four-dimensional signed zero spaces. Their intersection is precisely the scalar--scalar line: simultaneous equations require \(E_\Sigma=0\). Therefore its dimension is \(4+4-1=7\) before the scalar factor.

These are unconstrained carrier dimensions. A subsequent Hopf matching condition or quotient takes intersections with the specified invariant space; it does not inherit a kernel dimension without that restriction calculation.

\hypertarget{p1b:der:g02:s6:sub3}{%
\paragraph{The square-root distinction}\label{p1b:der:g02:s6:sub3}}

The additive energy just used is not \(\sqrt{A_x+A_p}\). The latter obeys the different valid inequality

\[
\boxed{\sqrt{A_x+A_p}\ge\sqrt{N_x^2+N_p^2}.}
\]

All operators commute, and the inequality follows from \(A_i\ge N_i^2\). The equality space has dimension \(9|F|\): each factor has three equality modes, namely the scalar and both extreme right-bivector weights.

It is essential to use \(E_\Sigma\) for the preceding summed-charge inequality. The minimal witness explaining this distinction is in Paper~I-D, Section~\ref{p1b:app:g02:3}.

\textbf{Source input:} Paper I, \path{p1:cor:v203-geometric-positive-polar}. The product inequalities and product kernel counts above are derived here.

\hypertarget{p1b:der:g02:s7}{%
\subsubsection{Scalar position at arbitrary level, enriched momentum fixed at one}\label{p1b:der:g02:s7}}

\hypertarget{p1b:der:g02:s7:sub1}{%
\paragraph{G02.7. Hybrid specialization}\label{p1b:der:g02:s7:sub1}}

The independent-level family used for scalar position and enriched momentum is

\[
\mathcal V_{\ell,1,F}=P_{\ell,x}\otimes H_{1,p}\otimes\mathbb C[F],
\qquad\ell\ge0.
\]

For the position sphere use the actual scalar Laplacian and its actual compressed scalar coordinates. For momentum choose either of the two level-one coordinate ancestors above, together with the specified retained LC square. This gives two momentum ancestry choices at every scalar position level. It does not replace \(P_\ell\) by the fifteen-dimensional enriched carrier when the level labels happen to agree.

The full second-order operator is

\[
\mathcal K^{\rm mix}
=\Delta_x|_{P_\ell}\otimes I+I\otimes K_{R_p}+T_{G,F}.
\]

Its exact joint spectrum is

\[
\boxed{\frac{k(k+2)}{R_x^2}+\frac{a_s}{R_p^2}+\chi^TG^{-1}\chi,
\quad \text{multiplicity }(k+1)^2d_s,\quad0\le k\le\ell.}
\]

This is a specific geometric specialization of the preceding atlas's fixed-momentum towers. The zero-space projector is the product of the position constant projector, the momentum scalar projector, and the zero-character projector. Thus its dimension is one when \(0\in F\), otherwise zero.

For the Hopf comparison, define \(E_x^{\rm sc}=R_x\sqrt{\Delta_x|_{P_\ell}}\). On degree \(k\), its eigenvalue is \(\sqrt{k(k+2)}\) and the right weights satisfy \(|q|\le k\). For \(k>0\),

\[
\sqrt{k(k+2)}>k\ge|q|.
\]

Therefore each signed scalar-position operator \((E_x^{\rm sc}+\epsilon N_x)/\sqrt2\) has exactly the constant line as kernel. Tensoring with the two-dimensional momentum kernel yields

\[
\boxed{\dim\ker\frac{E_x^{\rm sc}+E_p+\epsilon(N_x+\sigma N_p)}{\sqrt2}
=2|F|,\qquad\ell\ge0.}
\]

This proves the fixed-momentum result for every scalar position level, including \(\ell=0\). The different \(4|F|\) count in Section~\ref{p1b:der:g02:s6} belongs to the enriched--enriched carrier.

No intrinsic all-level LC connection or symbol grammar on every abstract \(H_m\) is inferred from the level-one smooth model. Such an extension would require specifying its smooth realization and operator. The arbitrary-level abstract carrier and Casimir family in the preceding sections already have their own stated definitions.

\hypertarget{p1b:der:g02:s8}{%
\subsubsection{Exact product sampling transport}\label{p1b:der:g02:s8}}

\hypertarget{p1b:der:g02:s8:sub1}{%
\paragraph{G02.8. Isometric transport theorem}\label{p1b:der:g02:s8:sub1}}

Let \(\mathsf S_R:\mathcal L_R\to\bigoplus_{q\in2T}\mathcal E_{Rq}\) be the full fibre sampling map

\[
(\mathsf S_Rs)_q=\sqrt{V_R/24}\,s(Rq),
\qquad\mathsf S_R^*\mathsf S_R=I.
\]

On the product define

\[
\mathsf T=(\mathsf S_{R_x}J_{R_x})\otimes
            (\mathsf S_{R_p}J_{R_p})\otimes I_F.
\]

Then \(\mathsf T^*\mathsf T=I_H\), and \(\mathsf T\) is unitary from \(H\) onto its image \(\mathcal V\), of dimension \(225|F|\). Each sphere has a ten-dimensional intrinsic fibre and 24 nodes. The full coordinate host has dimension \(24^2\cdot10^2|F|=57{,}600|F|\); its retained image is still \(225|F|\)-dimensional.

For every retained operator \(A\), put \(\widetilde A=\mathsf T A\mathsf T^*|_{\mathcal V}\). Then

\[
\widetilde{AB}=\widetilde A\widetilde B,\qquad
\widetilde{A^*}=\widetilde A^*,\qquad
\widetilde{f(A)}=f(\widetilde A)
\]

for spectral functions of a normal \(A\). The proof is insertion of \(\mathsf T^*\mathsf T=I\), followed by transport of the finite spectral resolution. Consequently every coordinate fingerprint, joint spectral projector, LC scalar projector, Hopf inequality, and kernel dimension proved above is preserved on \(\mathcal V\).

The relevant identity on the ambient coordinate host is the image projector \(\mathsf T\mathsf T^*\). Kernel and spectrum statements in this theorem are made on \(\mathcal V\); extending an operator by zero to the entire host adds unrelated zero modes.

For intrinsic coordinates there is also a pointwise symbol compiler:

\[
X_a^D=J_R^*\mathsf S_R^*\operatorname{diag}_{q\in2T}
 \left(\frac R{\sqrt3}c(dx_a)(Rq)\right)\mathsf S_RJ_R.
\]

The retained symbol matrix elements have polynomial degree at most five in a left frame, so the Hurwitz 5-design integrates them exactly. The product compiler follows by tensoring the two one-factor identities. The cone coordinates have exact operator transport under \(\mathsf T\); a pointwise compiler using the intrinsic symbol is specific to the intrinsic class and is not assigned to the cone class.

For the hybrid family, replace the position sampling factor by any isometry from an exact scalar-band quadrature, as supplied by the preceding sections. The same proof transports its geometric towers and \(2|F|\) Hopf kernel counts. The scalar-character factor in these constructions is retained in its Fourier-label model; point evaluation on the two independent scalar circles would require specifying the corresponding scalar quadrature as well.

\textbf{Source inputs:} Paper I, \path{p1:thm:v203-full-carrier-sampling}, \path{p1:cor:v203-full-symbol-transport}, \path{p1:thm:v204-isometric-algebra-transport}; Paper I-C, result~\ref{p1b:thm:full-sampling} and result~\ref{p1b:thm:all-level-sampling}.

Paper~I-D, Section~\ref{p1b:app:g02} proves the convention-comparison identities and extended-case formulas used in the results above.


\hypertarget{p1b:sec:g03}{%
\subsection{Product commutants, Clifford algebras and Casimir filtrations}\label{p1b:sec:g03}}


\begin{proposition}[Independent and diagonal operator algebras]
\label{p1b:prop:g03-product}
On \(H=H_{1,x}\otimes H_{1,p}\), let the two copies of
\(\mathrm{Spin}(4)\) act independently. Each \(H_1\) contains the six
inequivalent sectors of dimensions \(1,4,3,3,2,2\) once. Then
\[
 \operatorname{End}_{\mathrm{Spin}(4)^2}(H)\cong\mathbb C^{36},
 \qquad
 \dim_{\mathbb C}\operatorname{alg}^{*}\rho(\mathrm{Spin}(4)^2)=43^2.
\]
Under diagonal restriction, if
\(C_i=\sum_a L_{i,a}^{\,2}+\sum_a R_{i,a}^{\,2}\), the diagonal
Casimir is
\[
 C_{\rm diag}=C_x+C_p
 +2\sum_aL_{x,a}L_{p,a}+2\sum_aR_{x,a}R_{p,a}.
\]
The joint, summed and diagonal Casimir decompositions therefore use
different operators. Adjoining coordinates, ordinary or graded
exchange, or the even or full Clifford action yields the distinct
commutants calculated below on their stated reducing subspaces.
\end{proposition}

Schur decomposition proves the independent assertion. Expanding the diagonal Lie algebra generators proves the second identity. The following block decompositions give the remaining commutants and spectral projectors explicitly.

\hypertarget{p1b:der:g03:s10}{%
\subsubsection{Scope and notation}\label{p1b:der:g03:s10}}

Unless a scalar factor is explicitly included, work on

\[
H=H_{1,x}\otimes H_{1,p},\qquad \dim_{\mathbb C}H=225,
\qquad G=\operatorname{Spin}(4)_x\times\operatorname{Spin}(4)_p.
\]

All algebras and commutants are complex. A commutant is taken inside the full endomorphism algebra of the stated carrier. Write

\[
\mathcal A_G=\operatorname{alg}_{\mathbb C}^*\rho(G).
\]

The six one-factor sector dimensions are

\[
(d_0,d_v,d_L,d_R,d_+,d_-)=(1,4,3,3,2,2).
\]

Two different sign operators will be needed:

\begin{itemize}
\tightlist
\item
  On all of \(H_1\), the central representation sign is \(Z=I-2\Pi_{\rm sp}\). Its eigenspaces have dimensions \(11\) and \(4\), and vector coordinates preserve them.
\item
  On the protected four-dimensional space \(\mathcal K=S_+\oplus S_-\), the chiral sign is \(\Gamma=P_+-P_-\). Its eigenspaces have dimensions \(2\) and \(2\), and the protected vector coordinates interchange them.
\end{itemize}

In particular, \(Z|_{\mathcal K}=-I\); it is not the chiral sign \(\Gamma\). Equal commutant dimensions below do not identify these differently acting operators.

The preceding sections supply the general product-group algebra theorem, support-component theorem, and joint Casimir eigenvalue formula. The calculations below specialize those results, prove the remaining distinctions and projectors, and extend the explicit dictionary to exchange and diagonal restriction.

\textbf{Source anchors:} p1:thm:s3one-commutant-line; p1:cor:level-one-operator-algebra-dimensions; p1:thm:s3one-casimir-scale-space; p1b:thm:mixed-algebras; p1b:thm:all-level-generated-coordinates; p1b:thm:mixed-spectra.

\hypertarget{p1b:der:g03:s11}{%
\subsubsection{Full-carrier algebras and commutants}\label{p1b:der:g03:s11}}

\hypertarget{p1b:der:g03:s11:sub1}{%
\paragraph{The independent product group}\label{p1b:der:g03:s11:sub1}}

Each \(U_s\otimes U_t\) is an irreducible representation of \(G\), and all 36 joint labels are inequivalent. Schur's lemma gives

\[
\boxed{\mathcal A_G'
=\bigoplus_{s,t}\mathbb C(\Pi_s\otimes\Pi_t)
\cong\mathbb C^{36}.}
\]

The finite-dimensional double commutant theorem then gives

\[
\mathcal A_G=\bigoplus_{s,t}M_{d_sd_t}(\mathbb C),
\qquad
\dim_{\mathbb C}\mathcal A_G
=\left(\sum_s d_s^2\right)^2=43^2=1849.
\]

Every invariant Hermitian operator is

\[
A=\sum_{s,t}a_{st}\Pi_s\otimes\Pi_t,\qquad a_{st}\in\mathbb R.
\]

It is positive exactly when every \(a_{st}\ge0\), positive definite exactly when every \(a_{st}>0\), and has rank \(\sum_{a_{st}\ne0}d_sd_t\). Its normalized trace is \(225^{-1}\sum_{s,t}d_sd_ta_{st}\).

\hypertarget{p1b:der:g03:s11:sub2}{%
\paragraph{Adjoining the coordinate tuple}\label{p1b:der:g03:s11:sub2}}

Suppose every allowed edge of the level-one product coordinate graph is nonzero. An \(x\)-coordinate changes the \(x\) sector along a one-factor edge and retains the \(p\) sector; a \(p\)-coordinate does the reverse. The four connected components have dimensions

\[
11\cdot11=121,\qquad11\cdot4=44,\qquad4\cdot11=44,\qquad4\cdot4=16.
\]

To derive the generated algebra, sandwich any nonzero edge map between the full matrix algebras on its two irreducible endpoints. This spans the full rectangular Hom space between those endpoints. Multiplication along graph paths gives every matrix between vertices in a connected component. No path joins distinct components. Therefore

\[
\boxed{\mathcal A_{\rm op}:=
\operatorname{alg}_{\mathbb C}^*(\rho(G),X_1,\ldots,X_4,P_1,\ldots,P_4)
=M_{121}\oplus M_{44}\oplus M_{44}\oplus M_{16}.}
\]

Here and below the matrix fields are \(\mathbb C\). Consequently

\[
\boxed{\mathcal A_{\rm op}'
=\operatorname{span}_{\mathbb C}\{I,Z_x,Z_p,Z_xZ_p\}
\cong\mathbb C^4,}
\]

and \(\dim_{\mathbb C}\mathcal A_{\rm op}=18769\).

This conclusion depends on the nonzero support, not on edge phases or on vanishing cross commutators. It therefore applies throughout the full nonzero 48-edge atlas, including the four factorwise ancestry points derived in G02.

\hypertarget{p1b:der:g03:s11:sub3}{%
\paragraph{Adjoining exchange}\label{p1b:der:g03:s11:sub3}}

Let \(B(v\otimes w)=w\otimes v\) in the specified common model. Conjugation by \(B\) exchanges sector labels \((s,t)\leftrightarrow(t,s)\). An operator in \(\mathcal A_G'\) commutes with \(B\) precisely when \(a_{st}=a_{ts}\). Thus

\[
\boxed{\operatorname{alg}^*(\mathcal A_G,B)'
\cong\mathbb C^{21}.}
\]

The six diagonal sector orbits have carrier dimensions \(d_s^2\). Each of the 15 off-diagonal orbits has dimension \(2d_sd_t\). The group supplies all matrices on each original sector, and \(B\) supplies a nonzero map between the two members of an off-diagonal orbit. Hence the generated algebra is the full matrix algebra on each orbit, of total dimension

\[
\sum_s d_s^4+4\sum_{s<t}d_s^2d_t^2
=2(43)^2-\sum_s d_s^4=3247.
\]

For the coordinate-generated algebra, exchange fixes the \(121\)- and \(16\)-dimensional components and joins the two \(44\)-dimensional components. Therefore

\[
\boxed{\operatorname{alg}^*(\mathcal A_{\rm op},B)
=M_{121}\oplus M_{88}\oplus M_{16},}
\] \[
\boxed{\operatorname{alg}^*(\mathcal A_{\rm op},B)'
=\operatorname{span}_{\mathbb C}\{I,Z_x+Z_p,Z_xZ_p\}
\cong\mathbb C^3.}
\]

Its algebra dimension is \(22641\). Adjoining \(B\) here is an explicitly specified algebra construction. A particular coordinate tuple being fixed by exchange is an additional symmetry statement about that tuple.

The full-carrier dictionary is:

\begin{longtable}[]{@{}
  >{\raggedright\arraybackslash}p{(\columnwidth - 4\tabcolsep) * \real{0.3000}}
  >{\raggedleft\arraybackslash}p{(\columnwidth - 4\tabcolsep) * \real{0.4000}}
  >{\raggedright\arraybackslash}p{(\columnwidth - 4\tabcolsep) * \real{0.3000}}@{}}
\toprule\noalign{}
\begin{minipage}[b]{\linewidth}\raggedright
Acting algebra on \(H\)
\end{minipage} & \begin{minipage}[b]{\linewidth}\raggedleft
Complex dimension of algebra
\end{minipage} & \begin{minipage}[b]{\linewidth}\raggedright
Its commutant
\end{minipage} \\
\midrule\noalign{}
\endhead
\bottomrule\noalign{}
\endlastfoot
Independent group action & 1,849 & \(\mathbb C^{36}\) \\
Independent group action and exchange & 3,247 & \(\mathbb C^{21}\) \\
Group action and all nonzero coordinate edges & 18,769 & \(\mathbb C^4\), generated by the central signs \\
Those operations and exchange & 22,641 & \(\mathbb C^3\) \\
\(\operatorname{End}(H)=M_{225}\) & 50,625 & \(\mathbb C I\) \\
\end{longtable}

The algebra of all operators commuting with exchange alone is a different object:

\[
\{B\}'=M_{120}\oplus M_{105},
\qquad\dim_{\mathbb C}\{B\}'=25425.
\]

Indeed, the flip has trace \(15\) and dimension \(225\), giving its \(120\) and \(105\) eigenspace ranks. This exchange commutant does not impose \(G\)-invariance.

\hypertarget{p1b:der:g03:s12}{%
\subsubsection{Protected commutant, positivity, and exchange}\label{p1b:der:g03:s12}}

On

\[
\mathcal K_{xp}=(S_{x,+}\oplus S_{x,-})
\otimes(S_{p,+}\oplus S_{p,-}),\qquad\dim\mathcal K_{xp}=16,
\]

the four independent-group irreducibles have dimensions four. Write \(\Gamma_x=\Gamma\otimes I\) and \(\Gamma_p=I\otimes\Gamma\), and let

\[
P_{\epsilon\eta}
=\frac14(I+\epsilon\Gamma_x)(I+\eta\Gamma_p),
\qquad\epsilon,\eta\in\{+1,-1\}.
\]

These are mutually orthogonal rank-four projectors summing to the identity. Schur's lemma gives

\[
\boxed{\operatorname{End}_G(\mathcal K_{xp})
=\bigoplus_{\epsilon,\eta}\mathbb C P_{\epsilon\eta}
=\operatorname{span}_{\mathbb C}\{I,\Gamma_x,\Gamma_p,\Gamma_x\Gamma_p\}
\cong\mathbb C^4.}
\]

Each of the three nonidentity basis elements has trace zero. Thus the traceless invariant space is three-dimensional. The one-factor traceless commutant line becomes a three-dimensional space on this product.

For a Hermitian invariant operator

\[
A=aI+b\Gamma_x+c\Gamma_p+d\Gamma_x\Gamma_p,
\qquad a,b,c,d\in\mathbb R,
\]

the four eigenvalues, each of multiplicity four, are

\[
\boxed{\lambda_{\epsilon\eta}=a+\epsilon b+\eta c+\epsilon\eta d.}
\]

This completely specifies positivity, invertibility, and rank: require all four values nonnegative for positivity, none zero for invertibility, and count four times the number of nonzero values for rank. Its trace is \(16a\).

With matched exchange data,

\[
B\Gamma_xB=\Gamma_p,\qquad
B\Gamma_x\Gamma_pB=\Gamma_x\Gamma_p.
\]

Therefore the exchange-even invariant algebra is

\[
\boxed{\operatorname{End}_{G\rtimes\langle B\rangle}(\mathcal K_{xp})
=\operatorname{span}_{\mathbb C}\{I,\Gamma_x+\Gamma_p,\Gamma_x\Gamma_p\}
\cong\mathbb C^3.}
\]

Its three joint blocks have ranks \(4,8,4\), namely \(P_{++}\), \(P_{+-}+P_{-+}\), and \(P_{--}\). The exchange-even traceless invariant space is two-dimensional, while the exchange-odd invariant space is the line \(\mathbb C(\Gamma_x-\Gamma_p)\).

For an exchange-even Hermitian invariant operator \(c=b\), positivity reduces to

\[
a+2b+d\ge0,\qquad a-d\ge0,\qquad a-2b+d\ge0.
\]

Trace-one positive invariant operators form a tetrahedron in the four block weights; imposing exchange equality on the two mixed weights gives a triangle. This is a description of a mathematical positive cone and its trace-one section.

The protected group image is \(M_4^{\oplus4}\), of dimension 64. Adjoining \(B\) joins the two mixed blocks and gives

\[
M_4\oplus M_8\oplus M_4,
\qquad \dim_{\mathbb C}=96.
\]

\hypertarget{p1b:der:g03:s13}{%
\subsubsection{Even and full Clifford algebras}\label{p1b:der:g03:s13}}

Assume both protected coordinate magnitudes are positive and both protected subspaces are reducing for the coordinates. This holds at both level-one ancestry points in G02. Normalize the two quadruples to Hermitian matrices

\[
\{\gamma_{i,a},\gamma_{i,b}\}=2\delta_{ab}I_4.
\]

Choose compatible oriented Clifford conventions so that their fourfold volume elements are the chiral signs \(\Gamma_i\). The eight product generators are

\[
G_a=\gamma_{x,a}\otimes I,\qquad
G_{4+a}=\Gamma_x\otimes\gamma_{p,a},\qquad a=1,\ldots,4.
\]

The inserted chiral sign anticommutes with every \(\gamma_{x,a}\). Therefore all eight satisfy \(\{G_i,G_j\}=2\delta_{ij}I_{16}\).

The first four generate \(M_4\otimes I\). Their volume element supplies \(\Gamma_x\otimes I\); multiplying it into the second four recovers \(I\otimes\gamma_{p,a}\). Hence

\[
\boxed{\operatorname{alg}_{\mathbb C}^*(G_1,\ldots,G_8)=M_{16},
\qquad \operatorname{alg}_{\mathbb C}^*(G_1,\ldots,G_8)'=\mathbb C I.}
\]

The ordinary factorwise coordinate quadruples \(\gamma_{x,a}\otimes I\) and \(I\otimes\gamma_{p,a}\) also generate \(M_{16}\). Their cross commutators vanish; their combined eight-tuple does not satisfy the eight-generator Clifford anticommutators. Equality of the generated algebras does not identify those tuples.

For the even Clifford algebra, the total volume element is

\[
\Omega=G_1G_2\cdots G_8=\Gamma_x\Gamma_p.
\]

It has two eigenspaces of dimension eight. Every even Clifford word commutes with \(\Omega\). There are 128 independent even Clifford words in the faithful complex matrix representation, and the full block algebra preserving its two eigenspaces also has dimension \(8^2+8^2=128\). Consequently

\[
\boxed{\operatorname{Cl}_{8}^{0}(\mathbb C)\big|_{\mathcal K_{xp}}
=M_8\oplus M_8,\qquad
\left(\operatorname{Cl}_{8}^{0}(\mathbb C)\big|_{\mathcal K_{xp}}\right)'
=\mathbb C I\oplus\mathbb C\Omega.}
\]

This acting algebra has a one-dimensional traceless commutant. It couples \(++\) to \(--\) and \(+-\) to \(-+\), keeping only total Clifford parity fixed. The independent group action preserves each of the four chiral sectors separately.

\begin{longtable}[]{@{}
  >{\raggedright\arraybackslash}p{(\columnwidth - 4\tabcolsep) * \real{0.2727}}
  >{\raggedleft\arraybackslash}p{(\columnwidth - 4\tabcolsep) * \real{0.3636}}
  >{\raggedleft\arraybackslash}p{(\columnwidth - 4\tabcolsep) * \real{0.3636}}@{}}
\toprule\noalign{}
\begin{minipage}[b]{\linewidth}\raggedright
Acting algebra on the protected product
\end{minipage} & \begin{minipage}[b]{\linewidth}\raggedleft
Algebra dimension
\end{minipage} & \begin{minipage}[b]{\linewidth}\raggedleft
Commutant dimension
\end{minipage} \\
\midrule\noalign{}
\endhead
\bottomrule\noalign{}
\endlastfoot
Independent \(\operatorname{Spin}(4)^2\) & 64 & 4 \\
Independent group and ordinary exchange & 96 & 3 \\
Even complex \(\operatorname{Cl}_8\) image & 128 & 2 \\
Full complex \(\operatorname{Cl}_8\) image & 256 & 1 \\
\end{longtable}

The first row's two chiral signs are not the restrictions of the central signs generating the full-carrier coordinate commutant in Section~\ref{p1b:der:g03:s11}. On the protected product those central signs are scalar, and the full coordinate algebra has scalar commutant, consistently with the last row here.

\hypertarget{p1b:der:g03:s13:sub1}{%
\paragraph{The Clifford exchange operator}\label{p1b:der:g03:s13:sub1}}

Ordinary exchange \(B\) swaps tensor factors, but conjugates \(\gamma_{x,a}\otimes I\) to \(I\otimes\gamma_{p,a}\). The normalized Clifford tuple uses \(\Gamma_x\otimes\gamma_{p,a}\) in that position.

Define instead

\[
\boxed{\mathcal B_{\rm Cl}
=B(I-2P_{--}),\qquad
P_{--}=\frac14(I-\Gamma_x)(I-\Gamma_p).}
\]

On homogeneous chiral vectors, this is

\[
\mathcal B_{\rm Cl}(v_a\otimes w_b)
=(-1)^{ab}w_b\otimes v_a,\qquad a,b\in\{0,1\},
\]

where \(a=0\) means chirality \(+1\). Applying an odd first-factor generator changes \(a\) by one. The two exchange signs therefore leave precisely a factor \((-1)^b\), which is the inserted first-factor chirality after exchanging. This proves

\[
\boxed{\mathcal B_{\rm Cl}G_a\mathcal B_{\rm Cl}=G_{4+a},\qquad
\mathcal B_{\rm Cl}G_{4+a}\mathcal B_{\rm Cl}=G_a.}
\]

Also \(\mathcal B_{\rm Cl}^2=I\), \(\mathcal B_{\rm Cl}^*=\mathcal B_{\rm Cl}\), and \([\mathcal B_{\rm Cl},\Omega]=0\).

For ordinary exchange, only equal tensor basis vectors contribute to the trace, giving \(\operatorname{Tr}B=4\). For \(\mathcal B_{\rm Cl}\), the two positive-chirality equal pairs contribute \(+1\) and the two negative-chirality equal pairs contribute \(-1\). Thus

\[
\boxed{\operatorname{rank}P_{B,\pm}=(10,6),\qquad
\operatorname{rank}P_{\mathcal B_{\rm Cl},\pm}=(8,8).}
\]

On total Clifford parity \(+1\), the ordinary ranks are \((6,2)\) and the Clifford-exchange ranks are \((4,4)\). On total parity \(-1\), both have ranks \((4,4)\).

Since \(I-2P_{--}\) belongs to the independent group image, adjoining either \(B\) or \(\mathcal B_{\rm Cl}\) to that group image gives the same algebra and the same three-dimensional commutant. Their different eigenspace ranks concern the individual exchange operators.

This is exchange of the finite Clifford module with its declared chiral grading. It does not identify that grading with the central representation sign or replace the smooth spin-bundle exchange calculation scheduled under G08.

\textbf{Source input:} p1b:cor:product-clifford. The explicit even-image commutant and Clifford-exchange formula are derived above.

\hypertarget{p1b:der:g03:s14}{%
\subsubsection{Changing to the diagonal group action}\label{p1b:der:g03:s14}}

Independent and common diagonal actions give different commutants. For any diagonal restriction

\[
H\big|_{G_{\rm diag}}\cong\bigoplus_\alpha
V_\alpha\otimes\mathbb C^{\mu_\alpha},
\]

Schur's lemma gives

\[
\boxed{\operatorname{End}_{G_{\rm diag}}(H)
\cong\bigoplus_\alpha M_{\mu_\alpha}(\mathbb C).}
\]

The matrix blocks act on multiplicity spaces; the group acts on the \(V_\alpha\) factors. This template applies to a specified subgroup restriction, with the corresponding actual multiplicities.

\hypertarget{p1b:der:g03:s14:sub1}{%
\paragraph{Protected product}\label{p1b:der:g03:s14:sub1}}

For the full diagonal \(\operatorname{Spin}(4)\),

\[
\mathcal K_{xp}\big|_{G_{\rm diag}}
\cong2(0,0)\oplus(1,0)\oplus(0,1)
\oplus2(\tfrac12,\tfrac12).
\]

Indeed, \(S_+\otimes S_+\) and \(S_-\otimes S_-\) each give a scalar plus a triplet, while the two mixed products give two copies of the vector. Therefore

\[
\boxed{\operatorname{End}_{G_{\rm diag}}(\mathcal K_{xp})
\cong M_2\oplus\mathbb C\oplus\mathbb C\oplus M_2,
\qquad\dim_{\mathbb C}=10.}
\]

Ordinary exchange is \(-1\) on both scalar singlets, \(+1\) on the two triplets, and swaps the two vector copies. Requiring commutation with it leaves the full \(M_2\) on the scalars and splits the vector multiplicity space into its two exchange lines:

\[
\dim_{\mathbb C}\operatorname{End}_{G_{\rm diag},B}(\mathcal K_{xp})=8.
\]

For \(\mathcal B_{\rm Cl}\), the two scalar singlets have opposite signs, as do the two triplets. The vector multiplicity exchange is unchanged. Thus every surviving joint multiplicity is one and

\[
\dim_{\mathbb C}\operatorname{End}_{G_{\rm diag},\mathcal B_{\rm Cl}}
(\mathcal K_{xp})=6.
\]

These claims follow explicitly by using the antisymmetric and symmetric projectors on \(P_{++}\) and \(P_{--}\), and the exchange projectors on \(P_{+-}+P_{-+}\).

\hypertarget{p1b:der:g03:s14:sub2}{%
\paragraph{Full level-one product}\label{p1b:der:g03:s14:sub2}}

Write \(U_{a,b}=V_{a/2}\otimes V_{b/2}\) and let

\[
\Lambda_1=\{(0,0),(1,0),(0,1),(2,0),(1,1),(0,2)\}.
\]

The diagonal multiplicity of \(U_{A,B}\) in \(H_1\otimes H_1\) is

\[
\mu_{A,B}
=\sum_{(a,b),(c,d)\in\Lambda_1}N_{a,c}^{A}N_{b,d}^{B},
\] \[
N_{a,c}^{A}
=\mathbf1_{\{|a-c|\le A\le a+c,\ A\equiv a+c\pmod2\}}.
\]

This is the independent-level fusion formula with its commutant consequence made explicit. It gives all labels \(A+B\le4\). The multiplicities can be arranged as follows; omitted entries have multiplicity zero:

\begin{longtable}[]{@{}
  >{\raggedleft\arraybackslash}p{(\columnwidth - 10\tabcolsep) * \real{0.1667}}
  >{\raggedleft\arraybackslash}p{(\columnwidth - 10\tabcolsep) * \real{0.1667}}
  >{\raggedleft\arraybackslash}p{(\columnwidth - 10\tabcolsep) * \real{0.1667}}
  >{\raggedleft\arraybackslash}p{(\columnwidth - 10\tabcolsep) * \real{0.1667}}
  >{\raggedleft\arraybackslash}p{(\columnwidth - 10\tabcolsep) * \real{0.1667}}
  >{\raggedleft\arraybackslash}p{(\columnwidth - 10\tabcolsep) * \real{0.1667}}@{}}
\toprule\noalign{}
\begin{minipage}[b]{\linewidth}\raggedleft
\(A\)
\end{minipage} & \begin{minipage}[b]{\linewidth}\raggedleft
\(\mu_{A,0}\)
\end{minipage} & \begin{minipage}[b]{\linewidth}\raggedleft
\(\mu_{A,1}\)
\end{minipage} & \begin{minipage}[b]{\linewidth}\raggedleft
\(\mu_{A,2}\)
\end{minipage} & \begin{minipage}[b]{\linewidth}\raggedleft
\(\mu_{A,3}\)
\end{minipage} & \begin{minipage}[b]{\linewidth}\raggedleft
\(\mu_{A,4}\)
\end{minipage} \\
\midrule\noalign{}
\endhead
\bottomrule\noalign{}
\endlastfoot
0 & 6 & 6 & 5 & 2 & 1 \\
1 & 6 & 8 & 4 & 2 & --- \\
2 & 5 & 4 & 3 & --- & --- \\
3 & 2 & 2 & --- & --- & --- \\
4 & 1 & --- & --- & --- & --- \\
\end{longtable}

They satisfy

\[
\sum_{A,B}(A+1)(B+1)\mu_{A,B}=225,
\qquad
\boxed{\dim_{\mathbb C}\operatorname{End}_{G_{\rm diag}}(H)=281.}
\]

The represented diagonal-group algebra has dimension \(\sum_{A+B\le4}(A+1)^2(B+1)^2=406\). Its scalar invariant subspace has dimension six, agreeing with the existing full-diagonal singlet count.

Exchange on a self-pair \(U_{a,b}\otimes U_{a,b}\) has sign \((-1)^{a+b-(A+B)/2}\) on its \(U_{A,B}\) constituent. This is the product of the two usual symmetric/antisymmetric Clebsch--Gordan signs; on an off-diagonal pair exchange interchanges the two copies and has zero multiplicity trace. Consequently

\[
t_{A,B}
=\sum_{\substack{(a,b)\in\Lambda_1\\
 A\in\{0,2,\ldots,2a\}\\ B\in\{0,2,\ldots,2b\}}}
(-1)^{a+b-(A+B)/2},
\qquad
\mu_{A,B}^{\pm}=\frac{\mu_{A,B}\pm t_{A,B}}2.
\]

This derives the full exchange restriction without choosing multiplicity bases:

\[
\operatorname{End}_{G_{\rm diag},B}(H)
\cong\bigoplus_{A,B}\left(M_{\mu_{A,B}^+}\oplus M_{\mu_{A,B}^-}\right),
\qquad
\boxed{\dim_{\mathbb C}=145.}
\]

Zero-size blocks are omitted. The six diagonal singlets split into four exchange-even and two exchange-odd modes.

Finally, the common diagonal Casimir is

\[
C_{\rm diag}
=C_x+C_p
+2\sum_i L_{x,i}L_{p,i}
+2\sum_i R_{x,i}R_{p,i}.
\]

Its eigenvalues use the coupled labels \((A,B)\). The tensor sum \(C_x+C_p\) studied in the following filtration retains its own definition and eigenvalues.

\textbf{Source inputs:} p1b:prop:mixed-nesting and p1b:thm:mixed-diagonal.

\hypertarget{p1b:der:g03:s15}{%
\subsubsection{Explicit Casimir projectors and the full invariant algebra}\label{p1b:der:g03:s15}}

On one \(H_1\) let \(C=C_L+C_R\) and \(D=C_L-C_R\). The distinct values and ranks of \(C\) are

\[
(c_0,c_s,c_v,c_b)=(0,\tfrac34,\tfrac32,2),
\qquad(r_0,r_s,r_v,r_b)=(1,4,4,6).
\]

Lagrange interpolation gives the four exact spectral projectors:

\[
\begin{aligned}
p_0(C)&=(I-\tfrac43C)(I-\tfrac23C)(I-\tfrac12C),\\
p_s(C)&=\tfrac{64}{45}C(C-\tfrac32I)(C-2I),\\
p_v(C)&=-\tfrac{16}{9}C(C-\tfrac34I)(C-2I),\\
p_b(C)&=\tfrac45C(C-\tfrac34I)(C-\tfrac32I).
\end{aligned}
\]

Each polynomial takes value one on its named eigenvalue and zero on the other three. Thus the projectors are orthogonal, sum to the identity, and have ranks \((1,4,4,6)\).

The difference \(D\) splits the chiral pairs:

\[
\boxed{
\begin{aligned}
\Pi_{S_\pm}&=\tfrac12p_s(C)(I\pm\tfrac43D),\\
\Pi_{U_L},\Pi_{U_R}&=\tfrac12p_b(C)(I\pm\tfrac12D),\\
\Pi_{U_0}&=p_0(C),\qquad \Pi_{U_v}=p_v(C).
\end{aligned}}
\]

The two signs in the second line give left and right respectively. The formula follows because \(D=\pm3/4\) on the spin pair and \(D=\pm2\) on the bivector pair. Consequently

\[
\boxed{\mathbb C[C,D]=\operatorname{End}_{\operatorname{Spin}(4)}(H_1)
\cong\mathbb C^6.}
\]

Tensoring the six projectors yields every independent product-group invariant:

\[
\boxed{\mathbb C[C_x,D_x,C_p,D_p]=\mathcal A_G'\cong\mathbb C^{36}.}
\]

In contrast, \(\mathbb C[C_x,C_p]\cong\mathbb C^{16}\) resolves the four aggregate Casimir sectors in each factor. It does not split the two chiral spin or the two chiral bivector sectors.

For the general enriched carrier \(H_m\), the one-factor sector projector can be constructed from the two chiral Casimirs. Put \(\lambda_a=a(a+2)/4\) and

\[
L_a(T)=\prod_{\substack{0\le c\le2m\\c\ne a}}
\frac{T-\lambda_cI}{\lambda_a-\lambda_c}.
\]

Since the \(\lambda_a\) are strictly increasing, on every retained \((a,b)\),

\[
\Pi_{a,b}=L_a(C_L)L_b(C_R).
\]

The same construction works on a scalar band using its actual retained labels. This proves that the full chiral-Casimir family generates the independent-group commutant at every finite pair of levels. With a finite scalar-character set \(F\) and the independent torus action included, append its two charge generators to obtain the individual character projectors as well.

The degree-three formulas \(p_0,p_s,p_v,p_b\) displayed above are specific to the four-point spectrum of \(H_1\). At higher levels use the actual spectral projector or the interpolation polynomial for the enlarged spectrum.

\hypertarget{p1b:der:g03:s16}{%
\subsubsection{Joint and summed Casimir filtrations}\label{p1b:der:g03:s16}}

\hypertarget{p1b:der:g03:s16:sub1}{%
\paragraph{Joint filtration}\label{p1b:der:g03:s16:sub1}}

Define

\[
\Pi_{u,v}^{\rm rect}
=\mathbf1_{[0,u]}(C_x)\otimes\mathbf1_{[0,v]}(C_p).
\]

The projectors multiply by taking the minimum cutoff separately in each factor. Their ranks at the one-factor thresholds are:

\begin{longtable}[]{@{}lrrrr@{}}
\toprule\noalign{}
Position cutoff / momentum cutoff & \(0\) & \(3/4\) & \(3/2\) & \(2\) \\
\midrule\noalign{}
\endhead
\bottomrule\noalign{}
\endlastfoot
\(0\) & 1 & 5 & 9 & 15 \\
\(3/4\) & 5 & 25 & 45 & 75 \\
\(3/2\) & 9 & 45 & 81 & 135 \\
\(2\) & 15 & 75 & 135 & 225 \\
\end{longtable}

The proof is multiplication of the one-factor cumulative ranks \((1,5,9,15)\). Exact joint eigenspace projectors are \(p_i(C_x)p_j(C_p)\), of ranks \(r_ir_j\).

The associated two-parameter contraction semigroup is

\[
S_{t_x,t_p}=e^{-t_xC_x}\otimes e^{-t_pC_p},\qquad t_x,t_p\ge0.
\]

Its trace is \(Z(t_x)Z(t_p)\), where

\[
Z(t)=1+4e^{-3t/4}+4e^{-3t/2}+6e^{-2t}.
\]

Both parameters tending to infinity leaves precisely the scalar--scalar projector. At finite parameters the exponential has full rank; changes of rank belong to the sharp cutoff projectors.

\hypertarget{p1b:der:g03:s16:sub2}{%
\paragraph{Weighted sum and all positive parameter ratios}\label{p1b:der:g03:s16:sub2}}

For \(\alpha,\beta>0\) put \(K_{\alpha,\beta}=\alpha C_x+\beta C_p\). Its spectral projector at a numerical value \(\lambda\) is

\[
\boxed{P_\lambda(K_{\alpha,\beta})
=\sum_{\alpha c_i+\beta c_j=\lambda}p_i(C_x)p_j(C_p).}
\]

Its multiplicity is the same sum of \(r_ir_j\), and the cumulative cutoff rank is

\[
\boxed{N_{\alpha,\beta}(\Lambda)
=\sum_{i,j}r_ir_j\mathbf1_{\{\alpha c_i+\beta c_j\le\Lambda\}}.}
\]

These formulas describe every parameter value, including all eigenvalue collisions. The heat operator factors, and

\[
\left\|e^{-tK_{\alpha,\beta}}-P_{00}\right\|
=e^{-\frac34t\min(\alpha,\beta)}.
\]

For a complete ordering description, let \(\rho=\alpha/\beta\) and \(q_i=4c_i\in\{0,3,6,8\}\). Two distinct joint values coincide only if both labels change in opposite directions, and then

\[
\rho=\frac{q_l-q_j}{q_i-q_k}>0.
\]

The positive one-factor differences are \(\{2,3,5,6,8\}\). Thus the complete collision set is

\[
\boxed{\mathcal W=\{d/e:d,e\in\{2,3,5,6,8\}\}.}
\]

There are 21 distinct ratios: the ten values

\[
\tfrac14,\ \tfrac13,\ \tfrac38,\ \tfrac25,\ \tfrac12,\
\tfrac35,\ \tfrac58,\ \tfrac23,\ \tfrac34,\ \tfrac56,
\]

together with \(1\) and those ten reciprocals. Every ratio is attained by actual grid differences. The 22 complementary positive intervals have 16 distinct numerical eigenvalues each and constant ordering within each interval. The projector formula above gives every collision multiplicity on the boundaries.

At \(\alpha=0,\beta>0\), the four eigenspace ranks are \(15r_j\) and the zero space has dimension 15. The reversed boundary has the analogous result. At \(\alpha=\beta=0\), the entire carrier is the zero space. These boundary cases concern the weighted-Casimir family and do not change the representation carrier.

\hypertarget{p1b:der:g03:s16:sub3}{%
\paragraph{Equal weights and exchange}\label{p1b:der:g03:s16:sub3}}

For \(K=C_x+C_p\), the complete sharp filtration is:

\begin{longtable}[]{@{}
  >{\raggedleft\arraybackslash}p{(\columnwidth - 8\tabcolsep) * \real{0.2000}}
  >{\raggedleft\arraybackslash}p{(\columnwidth - 8\tabcolsep) * \real{0.2000}}
  >{\raggedleft\arraybackslash}p{(\columnwidth - 8\tabcolsep) * \real{0.2000}}
  >{\raggedleft\arraybackslash}p{(\columnwidth - 8\tabcolsep) * \real{0.2000}}
  >{\raggedleft\arraybackslash}p{(\columnwidth - 8\tabcolsep) * \real{0.2000}}@{}}
\toprule\noalign{}
\begin{minipage}[b]{\linewidth}\raggedleft
Eigenvalue
\end{minipage} & \begin{minipage}[b]{\linewidth}\raggedleft
Eigenspace rank
\end{minipage} & \begin{minipage}[b]{\linewidth}\raggedleft
Cumulative rank
\end{minipage} & \begin{minipage}[b]{\linewidth}\raggedleft
Exchange \(+\) rank
\end{minipage} & \begin{minipage}[b]{\linewidth}\raggedleft
Exchange \(-\) rank
\end{minipage} \\
\midrule\noalign{}
\endhead
\bottomrule\noalign{}
\endlastfoot
\(0\) & 1 & 1 & 1 & 0 \\
\(3/4\) & 8 & 9 & 4 & 4 \\
\(3/2\) & 24 & 33 & 14 & 10 \\
\(2\) & 12 & 45 & 6 & 6 \\
\(9/4\) & 32 & 77 & 16 & 16 \\
\(11/4\) & 48 & 125 & 24 & 24 \\
\(3\) & 16 & 141 & 10 & 6 \\
\(7/2\) & 48 & 189 & 24 & 24 \\
\(4\) & 36 & 225 & 21 & 15 \\
\end{longtable}

The eigenspace ranks are the convolution of \((c_i,r_i)\). For exchange, a distinct ordered pair \((i,j)\) is swapped with \((j,i)\) and has trace zero; an equal pair contributes trace \(r_i\) at energy \(2c_i\). Taking half the dimension plus or minus that trace proves the two exchange columns. They sum to 120 and 105.

Equivalently the exchange heat traces are

\[
Z_\pm(t)=\tfrac12\bigl(Z(t)^2\pm Z(2t)\bigr).
\]

For unequal weights, exchange transports \(K_{\alpha,\beta}\) to \(K_{\beta,\alpha}\); it does not preserve the two exchange eigenspaces as spectral subspaces of a fixed unequal-weight operator.

\hypertarget{p1b:der:g03:s16:sub4}{%
\paragraph{What the sum forgets}\label{p1b:der:g03:s16:sub4}}

The protected product is the rank-16 exact joint block

\[
\boxed{\Pi_{\rm prot}=p_s(C_x)p_s(C_p).}
\]

At total Casimir \(3/2\), it shares the eigenspace with \(U_0\otimes U_v\) and \(U_v\otimes U_0\), of combined dimension eight. Hence the rank-24 projector for \(K=3/2\) is not the protected projector. No function of \(K\) alone can split that eigenspace.

An exchange-even refinement is

\[
\Pi_{\rm prot}
=P_{3/2}(K)\,
\mathbf1_{\{0\}}\bigl((C_x-C_p)^2\bigr).
\]

The scalar--vector modes have \((C_x-C_p)^2=9/4\), while the protected modes have value zero.

The associated algebra dimensions make the loss of information exact:

\[
\mathbb C[K]\cong\mathbb C^9
\ \subset\
\mathbb C[C_x,C_p]\cong\mathbb C^{16}
\ \subset\
\mathcal A_G'\cong\mathbb C^{36}.
\]

The corresponding spectral commutants are

\[
\{C_x,C_p\}'=\bigoplus_{i,j}M_{r_ir_j},
\qquad \dim_{\mathbb C}=69^2=4761,
\] \[
\{K\}'=\bigoplus_\lambda M_{\operatorname{rank}P_\lambda(K)},
\qquad \dim_{\mathbb C}=7969.
\]

These are algebras of all operators preserving the stated spectral blocks. They allow additional mixing inside degenerate blocks and differ from the independent-group invariant algebra.

Away from the 21 collision ratios, every numerical eigenvalue of \(K_{\alpha,\beta}\) identifies one joint pair. Finite interpolation then gives \(\mathbb C[K_{\alpha,\beta}]=\mathbb C[C_x,C_p]\), and its commutant has dimension 4761. On a collision ratio, replace each merged group by the full matrix algebra on its combined eigenspace; the commutant dimension is the sum of those combined ranks squared. With exactly one zero weight it is \(15^2\sum_jr_j^2=15525\); with both weights zero it is \(50625\). This completes the commutant description throughout the nonnegative weighted-Casimir family.

For a coordinate block between sectors \(s,t\), direct subtraction gives

\[
\Pi_s[C,X_a]\Pi_t
=(c_s-c_t)\Pi_sX_a\Pi_t.
\]

At level one, the scalar--vector Casimir gap is \(3/2\), the vector--bivector gap is \(1/2\), and the protected chiral gap is zero. Thus a sharp spectral projector need not commute with the coordinates. A projector formed from whole representation sectors is reducing for the group-plus-coordinate algebra exactly when its selected vertex set is a union of active graph components. This follows from the nonzero edge blocks and supplies the precise compatibility criterion between the filtration and the operations.

The G02 geometric square obeys \(R^2K_{\rm LC}=2C+\frac34\Pi_{\rm sp}\) on each factor. The joint one-factor projectors therefore also identify its sectors, but the summed LC spectrum retains the different values and multiplicities already derived in G02.

\hypertarget{p1b:der:g03:s17}{%
\subsubsection{Arbitrary levels and fixed momentum one}\label{p1b:der:g03:s17}}

For a scalar or enriched position carrier \(A\), let \(C_A\) be its representation Casimir and define

\[
N_A(u)=\operatorname{rank}\mathbf1_{[0,u]}(C_A),
\qquad N_A(u)=0\text{ for }u<0.
\]

With enriched momentum fixed at one and positive weights, the complete counting formula is

\[
\boxed{
N_{A,1}(\Lambda)
=\sum_{j\in\{0,s,v,b\}}r_j\,
N_A\!\left(\frac{\Lambda-\beta c_j}{\alpha}\right).
}
\]

Its proof is the orthogonal decomposition into the four momentum Casimir sectors: each sector translates the position threshold and repeats it \(r_j\) times. The corresponding projector formula retains the position spectral projector tensored with \(p_j(C_p)\), so it also gives the full subspace and not only a dimension.

For \(A=P_\ell\), the position Casimir has value \(k(k+2)/2\) on degree \(k\), of dimension \((k+1)^2\). For \(u\ge0\) put

\[
L=\min\{\ell,\lfloor\sqrt{1+2u}-1\rfloor\}.
\]

Then

\[
\boxed{N_{P_\ell}(u)=\frac{(L+1)(L+2)(2L+3)}6.}
\]

The scalar geometric Laplacian is obtained by the separately specified weight \(2/R_x^2\) on this Casimir.

For the enriched carrier,

\[
\boxed{
N_{H_m}(u)
=\sum_{\substack{a,b\ge0,\ a+b\le2m\\
(a+1)^2+(b+1)^2\le4u+2}}(a+1)(b+1).
}
\]

This is an exact finite lattice sum, including every Casimir degeneracy. It and the four-tower formula apply at every finite position level, including \(m=0\).

The level inclusions preserve each sector eigenvalue. Every fixed Casimir window contains only finitely many labels, so increasing the position level eventually stabilizes its projector exactly. Momentum remains the same finite factor. This proves the filtration statement for arbitrarily large position levels.

\hypertarget{p1b:der:g03:s17:sub1}{%
\paragraph{Level truncation and spectral truncation}\label{p1b:der:g03:s17:sub1}}

The enriched cutoff \(a+b\le2m\) is not generally a sharp sublevel of the Casimir. At \(m=2\), \((a,b)=(4,0)\) is retained with Casimir 6, while \((2,3)\) is omitted with the smaller value \(23/4\). At \(m=3\), a stronger ambiguity occurs: retained \((6,0)\) and omitted \((4,4)\) both have Casimir 12. Thus at that level even an arbitrary function of the unrestricted total Casimir cannot select the level carrier.

This distinguishes the level family from its internal Casimir filtration. The full pair of chiral Casimirs always resolves the sector labels.

The four-dimensional spin pair remains the spectral subspace \(\mathbf1_{\{3/4\}}(C_{H_m})H_m\) for every \(m\ge1\). At \(m\ge2\), however, being a spectral subspace does not make it reducing for all coordinate tuples. The two exit edges

\[
U_{1,0}\leftrightarrow U_{2,1},\qquad
U_{0,1}\leftrightarrow U_{1,2}
\]

must vanish for that reducing property, exactly as specified in the preceding sections. The protected Clifford conclusions in Section~\ref{p1b:der:g03:s13} retain this domain.

For scalar position and enriched momentum, the protected momentum space \(P_\ell\otimes\mathcal K_p\) is always reducing. Its independent-group commutant is \(\mathbb C^{2(\ell+1)}\). Adding the actual scalar coordinates and nonzero protected momentum coordinates joins every position path vertex and both momentum chiral vertices. The represented algebra is then \(M_{4d_\ell}\), with scalar commutant. On the whole \(P_\ell\otimes H_1\), nonzero momentum exterior edges instead give \(M_{11d_\ell}\oplus M_{4d_\ell}\), with commutant \(\mathbb C^2\).

\textbf{Source inputs:} p1b:thm:mixed-dimensions; p1b:thm:mixed-spectra; p1b:thm:all-level-generated-coordinates; p1b:thm:all-level-limits.

\hypertarget{p1b:der:g03:s18}{%
\subsubsection{Scalar characters and coordinate boundary strata}\label{p1b:der:g03:s18}}

\hypertarget{p1b:der:g03:s18:sub1}{%
\paragraph{Independent scalar-character factor}\label{p1b:der:g03:s18:sub1}}

Let \(f=|F|>0\). If the independent torus action on \(\mathbb C[F]\) is included, its distinct characters give separate one-dimensional representation blocks. Therefore:

\begin{itemize}
\tightlist
\item
  The full independent-group algebra and commutant dimensions become \(1849f\) and \(36f\).
\item
  For neutral coordinates with the full nonzero support in every character block, the operation-algebra dimension becomes \(18769f\) and the commutant is \(\mathbb C^{4f}\).
\item
  On the protected product, the group commutant is \(\mathbb C^{4f}\). The full Clifford image together with the torus action has commutant \(\mathbb C^f\).
\item
  The full endomorphism algebra has dimension \(50625f^2\).
\end{itemize}

If only the Clifford algebra acting as \(M_{16}\otimes I_F\) is specified, its commutant is instead \(I_{16}\otimes M_f\). Including the torus action reduces that multiplicity algebra to the diagonal character algebra. The acting algebra must therefore be named in any scalar-factor commutant claim.

When \(F=F^\top\), exchange acts on joint labels by \((s,t,\chi)\mapsto(t,s,\chi^\top)\). Put

\[
f_\Delta=|\{(n,n)\in F\}|.
\]

The numbers of invariant scalar blocks are the numbers of these two-element orbits and fixed points:

\[
\dim\operatorname{End}_{G\times T^2,B}(H\otimes\mathbb C[F])
=\frac{36f+6f_\Delta}{2}=18f+3f_\Delta,
\] \[
\dim\operatorname{End}_{G\times T^2,B}(\mathcal K_{xp}\otimes\mathbb C[F])
=\frac{4f+2f_\Delta}{2}=2f+f_\Delta.
\]

The analogous commutant after all nonzero coordinate operations and exchange also has dimension \(2f+f_\Delta\), from the four central-sign components. If \(F\) is not exchange-invariant, exchange transports to the swapped carrier, and a fixed-carrier exchange commutant is not assigned without choosing an invariant completion.

\hypertarget{p1b:der:g03:s18:sub2}{%
\paragraph{Zero coordinate magnitudes}\label{p1b:der:g03:s18:sub2}}

The group commutants are independent of coordinate magnitudes. The normalized Clifford construction requires the stated positive magnitudes.

On the protected product, the boundary algebras are:

\begin{longtable}[]{@{}
  >{\raggedright\arraybackslash}p{(\columnwidth - 6\tabcolsep) * \real{0.2500}}
  >{\raggedright\arraybackslash}p{(\columnwidth - 6\tabcolsep) * \real{0.2500}}
  >{\raggedright\arraybackslash}p{(\columnwidth - 6\tabcolsep) * \real{0.2500}}
  >{\raggedright\arraybackslash}p{(\columnwidth - 6\tabcolsep) * \real{0.2500}}@{}}
\toprule\noalign{}
\begin{minipage}[b]{\linewidth}\raggedright
Active protected coordinate factors
\end{minipage} & \begin{minipage}[b]{\linewidth}\raggedright
Coordinate algebra alone
\end{minipage} & \begin{minipage}[b]{\linewidth}\raggedright
Its commutant
\end{minipage} & \begin{minipage}[b]{\linewidth}\raggedright
Commutant after including \(G\)
\end{minipage} \\
\midrule\noalign{}
\endhead
\bottomrule\noalign{}
\endlastfoot
Both & \(M_{16}\) & \(\mathbb C\) & \(\mathbb C\) \\
Position only & \(M_4\otimes I_4\) & \(I_4\otimes M_4\) & \(\operatorname{span}\{I,\Gamma_p\}\) \\
Momentum only & \(I_4\otimes M_4\) & \(M_4\otimes I_4\) & \(\operatorname{span}\{I,\Gamma_x\}\) \\
Neither & \(\mathbb C I_{16}\) & \(M_{16}\) & \(\mathbb C^4\) \\
\end{longtable}

For the complete one-factor level-one coordinate family with its nonzero normalized scalar edge, let \(l,r,s\in\{0,1\}\) indicate whether the left exterior, right exterior, and protected spin edge are nonzero. Its component dimensions are

\[
5+3(l+r),\quad
3\text{ for each missing exterior edge},\quad
\begin{cases}4&s=1,\\2,2&s=0.\end{cases}
\]

Thus its group-plus-coordinate algebra dimension is

\[
\boxed{[5+3(l+r)]^2+9(2-l-r)+16s+8(1-s),}
\]

and its commutant has dimension \(5-l-r-s\). Tensoring two globally factorwise supports multiplies their algebra and commutant dimensions. For a general product support, the connected-component proof of Section~\ref{p1b:der:g03:s11} applies directly. This covers the coordinate boundary strata without extending a normalization through a zero magnitude.


\hypertarget{p1b:sec:g04}{%
\subsection{Hurwitz, trimer, Walsh and resolution structures}\label{p1b:sec:g04}}


\begin{proposition}[The scalar Hurwitz bridge and typed finite products]
\label{p1b:prop:g04-product}
Let \(A_+\) be the Hurwitz adjacency operator, and equip its twenty-four
nodes with normalized counting measure. Evaluation is isometric on
the scalar band \(H_1^{\rm sc}=\mathbb C1\oplus E_1^{\rm sc}\) of
dimension five. Its product image has dimension twenty-five, and
\[
 \operatorname{ev}_{\times}(\Delta_x+\Delta_p)=
 \left[
 \frac{3}{4R_x^2}(8I-A_+)\otimes I+
 \frac{3}{4R_p^2}I\otimes(8I-A_+)
 \right]\operatorname{ev}_{\times}.
\]
This identity is restricted to that scalar product image.
Hurwitz relation modules, trimer corners, real Walsh--Clifford products
and resolution modules are defined by their actions and explicit
embeddings below; equality of dimensions does not identify these
objects with the enriched fifteen-dimensional representation.
\end{proposition}

The adjacency eigenspaces, polar trimer embeddings and twisted multiplication laws below prove these assertions and give the corresponding product and exchange constructions.

\hypertarget{p1b:der:g04:s20}{%
\subsubsection{Finite supports, actions, and source conventions}\label{p1b:der:g04:s20}}

This derivation covers the two sphere factors of \(M_8=S_x^3\times S_p^3\times T^2\): the Hurwitz relation algebra, oriented trimers and their shell embeddings, trimer operator spaces, twisted Walsh algebras, and resolution-data modules. Section~\ref{p1b:der:g04:s26} gives the scalar-character extension. Exchange statements assume an identification of the two factor packages, including their chosen frames, axes, and trimer orderings. Exchange is a symmetry of the round metric only when the corresponding radii agree; the algebraic flip exists independently.

All finite function spaces in Section~\ref{p1b:der:g04:s21}--\ref{p1b:der:g04:s23} have complex scalars and counting inner products, with coordinate delta functions orthonormal. Clifford and resolution spaces in Section~\ref{p1b:der:g04:s24}--\ref{p1b:der:g04:s25} have real scalars. Scalar evaluation in Section~\ref{p1b:der:g04:s26} uses normalized counting measure. These inner products differ by explicit constant factors, not by their vector-space dimensions.

The relevant one-factor objects have different types:

\begin{longtable}[]{@{}
  >{\raggedright\arraybackslash}p{(\columnwidth - 4\tabcolsep) * \real{0.3333}}
  >{\raggedright\arraybackslash}p{(\columnwidth - 4\tabcolsep) * \real{0.3333}}
  >{\raggedright\arraybackslash}p{(\columnwidth - 4\tabcolsep) * \real{0.3333}}@{}}
\toprule\noalign{}
\begin{minipage}[b]{\linewidth}\raggedright
Object
\end{minipage} & \begin{minipage}[b]{\linewidth}\raggedright
Definition
\end{minipage} & \begin{minipage}[b]{\linewidth}\raggedright
Dimension and field
\end{minipage} \\
\midrule\noalign{}
\endhead
\bottomrule\noalign{}
\endlastfoot
Full Hurwitz support & \(\mathscr H=\mathbb C^{2T}\) & \(24\), complex \\
Coarse shell spaces & \(\mathscr E_9=\ker A_+\), \(\mathscr E_{15}=\mathscr E_9^\perp\) & \(9,15\), complex \\
Oriented trimer & \(E_\sigma=\mathbb C^{T_\sigma}\), \(\sigma\in\{+,-\}\) & \(3\), complex \\
Trimer endomorphisms & \(\operatorname{End}_{\mathbb C}(E_\sigma)\) & \(9\), complex \\
Clifford carrier & \(A=\mathrm{Cl}_{0,3}\) & \(8\), real \\
Scalar-free carrier & \(W=A/\mathbb R1\) & \(7\), real \\
Resolution data & \(D=\operatorname{Hom}_{\mathbb R}(W,A)\) & \(56\), real \\
Enriched carrier from G02--G03 & \(H_1\) & \(15\), complex \\
\end{longtable}

The maps proved below determine the relationships between these objects. Equal dimensions alone do not supply an intertwiner. In particular, the natural scalar evaluation bridge has domain \(H_1^{\mathrm{sc}}\) of dimension five.

\textbf{One-factor source inputs from Paper I:}

\begin{longtable}[]{@{}
  >{\raggedright\arraybackslash}p{(\columnwidth - 2\tabcolsep) * \real{0.5000}}
  >{\raggedright\arraybackslash}p{(\columnwidth - 2\tabcolsep) * \real{0.5000}}@{}}
\toprule\noalign{}
\begin{minipage}[b]{\linewidth}\raggedright
Construction
\end{minipage} & \begin{minipage}[b]{\linewidth}\raggedright
Stable source labels
\end{minipage} \\
\midrule\noalign{}
\endhead
\bottomrule\noalign{}
\endlastfoot
Clifford carrier and resolution module & p1:thm:hopf-clifford-finite-carrier; p1:thm:septet-octet-resolution-module \\
Cube link, Walsh twist, 56-entry support & p1:prop:hurwitz-link-cube; p1:thm:hurwitz-walsh-clifford; p1:cor:walsh-evaluation-support-56 \\
Axial trimer decomposition & p1:thm:hurwitz-diagonal-trimer \\
Rooted vector and bivector overlaps &
p1:thm:rooted-trimer-vector-overlap;
p1:cor:rooted-trimer-bivector-overlap \\
Trimer grammar and Weyl basis & p1:thm:cyclic-trimer-operator-grammar; p1:prop:math-trimer-weyl-a2 \\
Zero extension and shell embedding & p1:def:s3one-trimer-zero-extension; p1:prop:s3one-trimer-adjacency-compression; p1:thm:s3one-trimer-bose-mesner-interface; p1:cor:s3one-trimer-operator-corner \\
Shell relation algebra & p1:thm:s3one-numerical-data; p1:prop:math-twentyfourcell-bose-mesner \\
Scalar evaluation and natural-action comparison & p1:prop:h1-scalar-hurwitz-evaluation-bridge; p1:prop:h1-e15-left-2t-obstruction \\
Protected Clifford transport & p1:prop:s3one-even-clifford-sign-transport \\
\end{longtable}

The one-factor inputs are replayed where needed; the tensor, exchange, quotient, and restriction formulas below are derived from them. The scope is this explicit list of constructions and operations.

\hypertarget{p1b:der:g04:s21}{%
\subsubsection{Hurwitz relation algebra on the product}\label{p1b:der:g04:s21}}

Write the unit Hurwitz shell as \[
\mathsf H=\{\pm e_0,\pm e_1,\pm e_2,\pm e_3\}
\ \cup\
\left\{\tfrac12(\epsilon_0,\epsilon_1,\epsilon_2,\epsilon_3):
       \epsilon_i=\pm1\right\}\subset S^3.
\] For \(t\in\{1,\tfrac12,0,-\tfrac12,-1\}\), let \[
(A_t)_{uv}=\mathbf1_{\{\langle u,v\rangle=t\}}.
\] Denote these matrices by \(I,A_+,A_0,A_-,A_\pi\), respectively. Their supports partition all ordered pairs. Their valencies are \((1,8,6,8,1)\). Counting intermediate nodes gives \[
\begin{aligned}
A_+^2&=8I+3A_++4A_0+A_-,&
A_+A_0&=3A_++3A_-,\\
A_+A_-&=A_++4A_0+3A_-+8A_\pi,&
A_+A_\pi&=A_-,\\
A_0^2&=6I+4A_0+6A_\pi,&
A_0A_-&=3A_++3A_-,\\
A_0A_\pi&=A_0,&
A_-^2&=A_+^2,\\
A_-A_\pi&=A_+,&
A_\pi^2&=I .
\end{aligned}
\] These identities show that the five-dimensional span is a commutative \(*\)-algebra \(\mathcal B\). Every displayed identity is an equality of the defined 24-node rational relation matrices.

The adjacency matrix \(A_+\) has eigenvalues and multiplicities \[
\begin{array}{c|rrrrr}
\lambda_i&8&4&0&-2&-4\\ \hline
r_i&1&4&9&8&2.
\end{array}
\] Consequently its five nonzero spectral projections \[
P_i=\prod_{j\ne i}\frac{A_+-\lambda_jI}{\lambda_i-\lambda_j}
\] are the primitive idempotents of \(\mathcal B\), and \(\mathcal B=\mathbb C[A_+]\cong\mathbb C^5\). In particular, \[
P_9=\frac{3I-A_0+3A_\pi}{8},\qquad
P_{15}=I-P_9=\frac{5I+A_0-3A_\pi}{8}.
\] Here \(P_9=P_{\lambda=0}\) is primitive, whereas \(P_{15}\) is the sum of the other four primitive projections.

On \(\mathscr H_x\otimes\mathscr H_p=\mathbb C^{\mathsf H\times\mathsf H}\), the product relation algebra is \[
\boxed{\mathcal B_x\otimes\mathcal B_p\cong\mathbb C^{25}.}
\] Its primitive projections are \(P_i^x\otimes P_j^p\), of rank \(r_ir_j\). This follows because the tensor projections are nonzero, mutually orthogonal, sum to the identity, and their scalar spans exhaust the 25 independent tensor relations.

For the factor flip \(B(u\otimes v)=v\otimes u\), \[
B(P_i\otimes P_j)B=P_j\otimes P_i.
\] Thus the subalgebra fixed under conjugation by \(B\) is \[
(\mathcal B\otimes\mathcal B)^B\cong\mathbb C^{15},
\] with primitive projections \[
P_i\otimes P_i,\qquad
P_i\otimes P_j+P_j\otimes P_i\quad(i<j),
\] of ranks \(r_i^2\) and \(2r_ir_j\). The dimension is \(5+\binom52=15\). This is the fixed part of the relation algebra, not the entire algebra of endomorphisms commuting with \(B\).

The coarse \(9/15\) decomposition gives \[
\mathscr H_x\otimes\mathscr H_p
=
(\mathscr E_9\otimes\mathscr E_9)
\oplus(\mathscr E_9\otimes\mathscr E_{15})
\oplus(\mathscr E_{15}\otimes\mathscr E_9)
\oplus(\mathscr E_{15}\otimes\mathscr E_{15}),
\] with dimensions \(81,135,135,225\). Exchange acts as follows:

\begin{longtable}[]{@{}
  >{\raggedright\arraybackslash}p{(\columnwidth - 6\tabcolsep) * \real{0.2000}}
  >{\raggedleft\arraybackslash}p{(\columnwidth - 6\tabcolsep) * \real{0.2667}}
  >{\raggedleft\arraybackslash}p{(\columnwidth - 6\tabcolsep) * \real{0.2667}}
  >{\raggedleft\arraybackslash}p{(\columnwidth - 6\tabcolsep) * \real{0.2667}}@{}}
\toprule\noalign{}
\begin{minipage}[b]{\linewidth}\raggedright
Exchange-stable space
\end{minipage} & \begin{minipage}[b]{\linewidth}\raggedleft
Dimension
\end{minipage} & \begin{minipage}[b]{\linewidth}\raggedleft
\(B=+1\) rank
\end{minipage} & \begin{minipage}[b]{\linewidth}\raggedleft
\(B=-1\) rank
\end{minipage} \\
\midrule\noalign{}
\endhead
\bottomrule\noalign{}
\endlastfoot
\(\mathscr E_9\otimes\mathscr E_9\) & 81 & 45 & 36 \\
\((\mathscr E_9\otimes\mathscr E_{15})\oplus(\mathscr E_{15}\otimes\mathscr E_9)\) & 270 & 135 & 135 \\
\(\mathscr E_{15}\otimes\mathscr E_{15}\) & 225 & 120 & 105 \\
Full product support & 576 & 300 & 276 \\
\end{longtable}

For a \(d\)-dimensional factor, the flip has trace \(d\) on its tensor square: only basis vectors \(e_i\otimes e_i\) contribute. Its two ranks are \(d(d\pm1)/2\). On two swapped summands of equal dimension the trace is zero. These facts prove the table.

The product shell projections retain both factor labels. A chosen product adjacency can coarsen them: the additive adjacency \(A_+\otimes I+I\otimes A_+\) has kernel dimension 97, whereas \(A_+\otimes A_+\) has kernel dimension 351. Paper~I-D, Section~\ref{p1b:app:g04:1} gives their spectral derivation.

\hypertarget{p1b:der:g04:s22}{%
\subsubsection{Trimers, polar embeddings, and operator corners}\label{p1b:der:g04:s22}}

At the vertex \(1\), the eight neighbors are \(\tfrac12(1,\epsilon_1,\epsilon_2,\epsilon_3)\). Choose the oriented body diagonal \((+++)\to(---)\) and order \[
T_+=\{(-++),(+-+),(++-)\},\qquad
T_-=\{(+--),(-+-),(--+)\}.
\] The notation lists the three signs; the corresponding actual nodes have real part \(1/2\) and imaginary coordinates equal to half the signs. The two remaining sign vectors are the axial pair. Cyclic permutation of the three imaginary coordinates advances each ordered trimer.

Let \(i_\sigma:E_\sigma\to\mathscr H\) be zero extension and \(i_W=i_+\oplus i_-\). The maps are isometries for counting measure. Put \[
J=\mathbf1\mathbf1^{\mathsf T},\quad
Q_0=J/3,\quad Q_\perp=I-Q_0,\quad C=J-I.
\] Direct node inner products give \[
i_W^*A_+i_W=
\begin{pmatrix}0&C\\ C&0\end{pmatrix},
\] the adjacency matrix of a six-cycle. Each of its nodes has two neighbors inside this support and six outside, so this is an adjacency compression.

\begin{proposition}[Rooted transport and ambient overlap compatibility]
\label{p1b:prop:trimer-transport-overlap}
For the ordered trimers above, let
\[
 I_0:E_+\longrightarrow E_-,\quad I_0e_{p_j}=e_{m_j},
 \qquad
 U_\pm e_{p_j}=e_{m_{j\pm1}}.
\]
Thus \(I_0\) is induced by the root involution, whereas \(U_\pm\)
are the two short-edge matchings. In the aligned coordinates put
\[
 X=I_0^{-1}U_+,\qquad X^{-1}=I_0^{-1}U_-,\qquad
 C_1=X+X^{-1}=C,\qquad C_2=-i(X-X^{-1}).
\]
Complexify the point-linearization maps \(L_\pm\) and the
real-Hodge exterior maps \(B_\pm\) of Paper~I-A,
Theorem~\ref{p1:thm:rooted-trimer-vector-overlap} and
Corollary~\ref{p1:cor:rooted-trimer-bivector-overlap}.
Regarding their domains as \(E_\pm\), one has
\[
 L_+^*L_-I_0=C_1^{-1},\qquad
 B_+^*B_-I_0=\tfrac12C_1.
\]
These are fixed functions of the same even transport port.
On the real augmentation plane, the imaginary-coordinate
restriction of \(L_+\) is
\[
 \Psi(a)=\sum_{j=0}^2a_j\operatorname{Im}p_j=-a
 \quad\left(\sum_j a_j=0\right).
\]
It is an isometry onto \(u^\perp\) and intertwines the trimer
cycle with the corresponding cyclic rotation of that plane.
\end{proposition}

\begin{proof}
The two edge maps advance or reverse the aligned cyclic labels,
so their pulled-back sum is the previously computed incidence
\(C=J-I\). The two overlap identities are the cited one-factor
results. The formula for \(\Psi\) is the restriction of \(L_+\)
to zero-sum coefficients. Cyclic permutation of the imaginary
coordinates carries \(p_j\) to \(p_{j+1}\), proving equivariance.
\end{proof}

\begin{corollary}[Products of the rooted vector and bivector overlaps]
\label{p1b:cor:product-trimer-overlap}
Give each factor \(f\in\{x,p\}\) the unit rooted Hurwitz data of
Proposition~\ref{p1b:prop:trimer-transport-overlap}.
Use the aligned counting space \(G_f\) for both trimer domains,
and write \(C_f=X_f+X_f^{-1}\).
For \(s\in\{V,\Lambda^2\}\), put
\[
 T^f_{V,\pm}=L^f_\pm,\qquad
 T^f_{\Lambda^2,\pm}=B^f_\pm,
 \qquad
 M_{s,t,\pm}=T^x_{s,\pm}\otimes T^p_{t,\pm}.
\]
These are isometries from \(G_x\otimes G_p\) into the indicated
tensor products of vector or bivector spaces. Their overlaps are
\[
\begin{aligned}
 \mathbb K_{V,V}&=C_x^{-1}\otimes C_p^{-1},&
 \mathbb K_{V,\Lambda^2}&=\tfrac12C_x^{-1}\otimes C_p,\\
 \mathbb K_{\Lambda^2,V}&=\tfrac12C_x\otimes C_p^{-1},&
 \mathbb K_{\Lambda^2,\Lambda^2}&=\tfrac14C_x\otimes C_p,
\end{aligned}
\qquad
 \mathbb K_{s,t}:=M_{s,t,+}^*M_{s,t,-}.
\]
In particular,
\[
 \mathbb K_{V,V}\mathbb K_{\Lambda^2,\Lambda^2}=\tfrac14I_9,
 \qquad
 \mathbb H_{\Lambda^2,\Lambda^2}
 =\tfrac1{16}\mathbb H_{V,V}^{-1},
 \quad \mathbb H_{s,t}:=\mathbb K_{s,t}^*\mathbb K_{s,t}.
\]
The factor flip transports these constructions to the corresponding
ones with \(x,p\) and \(s,t\) exchanged. It is an endomorphism
of an identified factor package when the rooted data are matched.
\end{corollary}

\begin{proof}
For isometries \(A_\pm,D_\pm\),
\[
 (A_+\otimes D_+)^*(A_-\otimes D_-)
 =(A_+^*A_-)\otimes(D_+^*D_-).
\]
Apply the one-factor overlap formulas in each slot. The reciprocal
identity follows by multiplication. The displayed overlaps are
self-adjoint, so multiplying their squares gives the stated
positive-return identity.
Naturality of the tensor flip proves the exchange statement.
\end{proof}

These tensor targets include
\(\Lambda^2V_x\otimes\Lambda^2V_p\); they are not being
identified with \(\Lambda^2(V_x\otimes V_p)\).
The statements concern the declared rooted mathematical factor
packages. Their \(x\)-labels alone do not constitute an independent
microscopic position carrier.

The shell Gram matrices are \[
i_\sigma^*P_ri_\tau=\delta_{\sigma\tau}B_r,\qquad
B_9=\frac{4I-J}{8}=\frac18Q_0+\frac12Q_\perp,\qquad
B_{15}=\frac{4I+J}{8}=\frac78Q_0+\frac12Q_\perp.
\] For example, the three nodes in either trimer are mutually in relation \(A_0\); cross-trimer nodes lie in neither \(A_0\) nor \(A_\pi\). Substitution into the formulas for \(P_r\) proves all blocks. Both \(B_r\) are strictly positive. Define \[
U_{9,\sigma}
=P_9i_\sigma(\sqrt8\,Q_0+\sqrt2\,Q_\perp),\qquad
U_{15,\sigma}
=P_{15}i_\sigma(\sqrt{8/7}\,Q_0+\sqrt2\,Q_\perp).
\] Then \[
U_{r,\sigma}^*U_{r,\tau}=\delta_{\sigma\tau}I_3.
\] Indeed, these are the polar normalizations \(P_ri_\sigma B_r^{-1/2}\). The cyclic symmetry commutes with \(P_r,Q_0,Q_\perp\), so the embeddings are \(C_3\)-equivariant.

For two sphere factors, set \[
i_{\sigma\tau}=i_\sigma^x\otimes i_\tau^p:
E_\sigma^x\otimes E_\tau^p\longrightarrow\mathscr H_x\otimes\mathscr H_p.
\] The domain has dimension nine. For \(r,s\in\{9,15\}\), \[
i_{\sigma\tau}^*(P_r\otimes P_s)i_{\sigma'\tau'}
=\delta_{\sigma\sigma'}\delta_{\tau\tau'}(B_r\otimes B_s).
\] Therefore the product polar embeddings are exactly \[
\boxed{
U_{rs,\sigma\tau}
=U_{r,\sigma}\otimes U_{s,\tau}.
}
\] All four sign-pair embeddings have mutually orthogonal ranges in each \(\mathscr E_r\otimes\mathscr E_s\). Each coarse block consequently contains four specified nine-dimensional copies, of total dimension 36. Their construction is determined by the chosen axis and trimer order.

The eigenvalues of the unnormalized product Gram matrix are:

\begin{longtable}[]{@{}
  >{\raggedright\arraybackslash}p{(\columnwidth - 2\tabcolsep) * \real{0.5000}}
  >{\raggedright\arraybackslash}p{(\columnwidth - 2\tabcolsep) * \real{0.5000}}@{}}
\toprule\noalign{}
\begin{minipage}[b]{\linewidth}\raggedright
\((r,s)\)
\end{minipage} & \begin{minipage}[b]{\linewidth}\raggedright
Eigenvalues, with multiplicities
\end{minipage} \\
\midrule\noalign{}
\endhead
\bottomrule\noalign{}
\endlastfoot
\((9,9)\) & \(1/64\) (1), \(1/16\) (4), \(1/4\) (4) \\
\((9,15)\) or \((15,9)\) & \(7/64\) (1), \(1/16\) (2), \(7/16\) (2), \(1/4\) (4) \\
\((15,15)\) & \(49/64\) (1), \(7/16\) (4), \(1/4\) (4) \\
\end{longtable}

These follow by multiplying the eigenvalues on \(Q_0\otimes Q_0,Q_0\otimes Q_\perp,Q_\perp\otimes Q_0,Q_\perp\otimes Q_\perp\). All are positive and less than one. Thus projection of the 36-dimensional supported product space into any one coarse block is injective, while no nonzero vector of that support already lies in one block. The polar embeddings use projections and normalization; they do not identify the original support with a shell eigenspace.

For every isometry \(V:\mathbb C^d\to\mathbb C^N\), define \[
j_V(a)=VaV^*.
\] Then \(j_V\) is an injective \(*\)-homomorphism onto the corner with identity \(VV^*\), and \[
\tau_N(j_V(a))=\frac dN\,\tau_d(a).
\] The product trimer zero-extension corner consequently satisfies \[
\tau_{576}\!\left(i_{\sigma\tau}a\,i_{\sigma\tau}^*\right)
=\frac1{64}\tau_9(a).
\] The polar corner in \(\mathscr E_r\otimes\mathscr E_s\) satisfies \[
\tau_{rs}\!\left(U_{rs,\sigma\tau}a\,U_{rs,\sigma\tau}^*\right)
=\frac9{rs}\tau_9(a).
\] For all four sign pairs together, the corner has domain dimension 36 and normalized-trace factor \(36/576=1/16\). Cross maps \(V_\alpha aV_\beta^*\) have source corner \(\beta\) and target corner \(\alpha\). Orthogonality gives the matrix-unit composition rule \[
(V_\alpha aV_\beta^*)(V_\gamma bV_\delta^*)
=\delta_{\beta\gamma}V_\alpha abV_\delta^*.
\]

Exchange carries the product embedding with indices \((r,s,\sigma,\tau)\) to the embedding with indices \((s,r,\tau,\sigma)\), composed with the domain flip. For \(r=s\) and \(\sigma=\tau\), its rank split on the embedded trimer square is \(6+3\). On the full paired-trimer product, \((E_+\oplus E_-)^{\otimes2}\), the split is \(21+15\).

For completeness, the finite cyclic restrictions can also be made explicit. The imaginary-coordinate cycle on all 24 nodes has character \((24,6,6)\). Taking traces after \(P_9,P_{15}\) gives \[
\chi_{\mathscr E_9}=(9,0,0),\qquad
\chi_{\mathscr E_{15}}=(15,6,6).
\] The character multiplicities for charges \(a=0,1,2\) are therefore \[
m_9=(3,3,3),\qquad m_{15}=(9,3,3).
\] This follows either from the three-term Fourier transform or from \(\mathscr E_9\cong3\mathbb C[C_3]\) and \(\mathscr E_{15}\cong6\mathbf1\oplus3\mathbb C[C_3]\). On a product, independent \(C_3\times C_3\) multiplicities are \(m_r(a)m_s(b)\). Hence \[
\dim(\mathscr E_r\otimes\mathscr E_s)^{C_3\times C_3}
=m_r(0)m_s(0),
\quad
\dim(\mathscr E_r\otimes\mathscr E_s)^{C_3^{\rm diag}}
=\sum_{a=0}^2m_r(a)m_s(-a).
\]

\begin{longtable}[]{@{}lrr@{}}
\toprule\noalign{}
\((r,s)\) & Independent invariant rank & Common-cycle invariant rank \\
\midrule\noalign{}
\endhead
\bottomrule\noalign{}
\endlastfoot
\((9,9)\) & 9 & 27 \\
\((9,15)\), \((15,9)\) & 27 & 45 \\
\((15,15)\) & 81 & 99 \\
\end{longtable}

On the common-cycle invariant space for \(r=s\), the trace of exchange is \(m_r(0)\): the only charge satisfying \(a=-a\) in \(\mathbb Z_3\) is zero, and the flip trace there is the dimension of the zero-charge factor. The corresponding exchange ranks are \(15,12\) for \(r=9\) and \(54,45\) for \(r=15\).

\hypertarget{p1b:der:g04:s23}{%
\subsubsection{Product trimer operator spaces and defect geometry}\label{p1b:der:g04:s23}}

On an oriented trimer let \(Xe_j=e_{j+1}\), \(Ze_j=\omega^je_j\), \(\omega=e^{2\pi i/3}\). Then \[
X^3=Z^3=I,\qquad ZX=\omega XZ,\qquad
\operatorname{Tr}((X^aZ^b)^*X^{a'}Z^{b'})
=3\delta_{aa'}\delta_{bb'}.
\] For the trace identity, a nonzero net shift has no diagonal entries; with zero net shift, the diagonal sum is either \(3\) or \(1+\omega+\omega^2=0\). Thus the nine words form a complex basis of \(M_3(\mathbb C)\).

The product trimer has the 81 trace-orthogonal words \[
(X_x^aZ_x^b)\otimes(X_p^cZ_p^d),\qquad
(a,b,c,d)\in\mathbb Z_3^4,
\] with squared Hilbert--Schmidt norm \(9\). They form a basis of \(\operatorname{End}_{\mathbb C}(E_x\otimes E_p)=M_9(\mathbb C)\). The two factor Weyl algebras commute. Exchange swaps their two index pairs.

The source's primitive Hermitian trimer slice is \[
\mathcal T=\operatorname{span}_{\mathbb R}
\left\{I,D_1,D_2,X+X^*,-i(X-X^*)\right\},
\] where \[
D_1=\operatorname{diag}(2/3,-1/3,-1/3),\qquad
D_2=\operatorname{diag}(0,1,-1).
\] Its dimension is five: three independent diagonal matrices and two independent off-diagonal Hermitian matrices. With tensor operations specified explicitly, it has the following product descendants: \[
\mathcal T_{\rm loc}
=(\mathcal T_x\otimes I)+(I\otimes\mathcal T_p),
\qquad
\mathcal T_{\rm tens}
=\operatorname{span}_{\mathbb R}
\{a\otimes b:a\in\mathcal T_x,\ b\in\mathcal T_p\}.
\] The two local summands intersect exactly in the scalar line, so \[
\dim_{\mathbb R}\mathcal T_{\rm loc}=5+5-1=9,\qquad
\dim_{\mathbb R}\mathcal T_{\rm tens}=25.
\] Their exchange-even dimensions are respectively \(5\) and \(\dim\operatorname{Sym}^2(\mathbb R^5)=15\); the exchange-odd dimensions are \(4\) and \(10\). These are linear slices. The complex \(*\)-algebra generated by either product slice is \(M_9(\mathbb C)\): each factor slice recovers \(X\) and all diagonal projectors, which generate every matrix unit; tensoring the two factor algebras supplies all product matrix units. The full Hermitian space of \(M_9(\mathbb C)\) has real dimension 81. Algebra generation does not change the definition of the primitive slices.

For \(P_j=|e_j\rangle\langle e_j|\), the centered diagonal defects \(D_j=P_j-I/3\) satisfy \[
\operatorname{Tr}(D_iD_j)=\delta_{ij}-\frac13,\qquad
\sum_jD_j=0.
\] The simple differences \(\alpha_1=P_0-P_1\), \(\alpha_2=P_1-P_2\) have Gram matrix \[
\begin{pmatrix}2&-1\\-1&2\end{pmatrix}.
\] On the product, the factor-root planes are spanned by \(\alpha_i\otimes I\) and \(I\otimes\alpha_i\). They are orthogonal, since the \(\alpha_i\) have zero trace, and each Gram matrix acquires the factor \(\operatorname{Tr}I_3=3\). They give a four-dimensional \(A_2\oplus A_2\) factor-root space (with the original Gram normalization after division by \(\sqrt3\)).

The entire product diagonal trace-zero space is larger. Writing \(\mathrm{Diag}_3=\mathbb RI\oplus\mathfrak d_0\), \(\dim\mathfrak d_0=2\), gives \[
(\mathrm{Diag}_3\otimes\mathrm{Diag}_3)_0
=(\mathfrak d_0\otimes I)
\oplus(I\otimes\mathfrak d_0)
\oplus(\mathfrak d_0\otimes\mathfrak d_0),
\] of dimension \(2+2+4=8\). If the nine product rank-one projectors are listed as \(Q_1,\ldots,Q_9\), then \[
\operatorname{Tr}\bigl((Q_i-I/9)(Q_j-I/9)\bigr)
=\delta_{ij}-\frac19.
\] The successive differences \(Q_i-Q_{i+1}\) have the \(A_8\) Cartan Gram matrix. Thus the factor \(A_2\oplus A_2\) geometry and the full nine-point defect geometry occupy explicitly different subspaces.

\hypertarget{p1b:der:g04:s24}{%
\subsubsection{Walsh carrier and the two Clifford products}\label{p1b:der:g04:s24}}

Let \(G=\mathbb F_2^3\) label the sign cube, with \[
\chi_a(s)=(-1)^{a\cdot s},\qquad
H_{sa}=\chi_a(s),\qquad a,s\in G.
\] For \(a\ne b\), summing \(\chi_{a+b}\) over \(G\) gives zero by pairing each point with a translate on which the character changes sign. Consequently \[
H^{\mathsf T}H=8I,\qquad
\left(\frac H{\sqrt8}\right)^{\mathsf T}
\left(\frac H{\sqrt8}\right)=I.
\] Deleting the constant column gives a matrix \(H_*\) with \[
H_*^{\mathsf T}H_*=8I_7,\qquad
\frac18H_*H_*^{\mathsf T}=I_8-\frac18\mathbf1\mathbf1^{\mathsf T}.
\] The normalized nonconstant Walsh columns therefore identify the seven-dimensional character-label space with the mean-zero function subspace on the cube.

To obtain the Clifford product, keep the ordered three-coordinate frame and specify the cocycle \[
\sigma_-(a,b)=(-1)^{\beta(a,b)},\qquad
\beta(a,b)=\sum_i a_ib_i+\sum_{i>j}a_ib_j
\quad\text{in }\mathbb F_2.
\] Define the product on the Walsh basis by \[
\chi_a\star\chi_b=\sigma_-(a,b)\chi_{a+b}.
\] The bilinearity of \(\beta\) proves \[
\beta(a,b)+\beta(a+b,c)=\beta(b,c)+\beta(a,b+c);
\] hence the product is associative and \(\chi_0\) is its unit. The three generators \(e_i=\chi_{\mathbf e_i}\) satisfy \[
e_i^2=-1,\qquad e_ie_j=-e_je_i\quad(i\ne j).
\] Their ordered monomials give all eight basis vectors. The universal Clifford relations thus give an isomorphism \[
(\mathbb R^G,\star)\cong A=\mathrm{Cl}_{0,3}.
\] The grading by \(|a|\bmod2\) has even and odd dimensions four each. The finer exterior-degree dimensions are \(1,3,3,1\). The Clifford product here is the displayed twist of the Walsh basis; pointwise multiplication of cube functions instead gives \(\chi_a\chi_b=\chi_{a+b}\).

The explicit quaternion model \[
e_1\mapsto(i,i),\qquad e_2\mapsto(j,j),\qquad
e_3\mapsto(k,-k)
\] identifies \(A\cong\mathbb H\oplus\mathbb H\). The volume element \(\omega=e_1e_2e_3\) maps to \((-1,1)\), is central, and squares to one.

For two factors, use labels \((a,b)\in G_x\times G_p\). There are two specified associative products on the same 64-dimensional underlying real vector space.

The \textbf{ordinary tensor product} uses \[
\sigma_{\rm ord}((a,b),(c,d))
=\sigma_-(a,c)\sigma_-(b,d).
\] The factor generators commute across the two factors. Therefore \[
A_{\rm ord}=A_x\otimes_{\mathbb R}A_p.
\] Its center has dimension four, with basis \[
1,\quad\omega_x\otimes1,\quad1\otimes\omega_p,\quad
\omega_x\otimes\omega_p.
\] For example, tensoring the two quaternion direct sums gives four copies of \(\mathbb H\otimes_{\mathbb R}\mathbb H\). The map \(a\otimes b:v\mapsto av\bar b\) identifies each of these with \(\operatorname{End}_{\mathbb R}(\mathbb H)\cong M_4(\mathbb R)\): the multiplication law is respected. For the orthonormal quaternion basis \(1,i,j,k\), the 16 corresponding operators \(T_{ab}\) obey \(\operatorname{Tr}(T_{ab}^{\mathsf T}T_{cd}) =4\delta_{ac}\delta_{bd}\), by direct quaternion multiplication (equivalently, use \(\operatorname{Tr}(v\mapsto av\bar b) =4\operatorname{Re}(a)\operatorname{Re}(b)\)). They are therefore linearly independent. Thus \[
A_{\rm ord}\cong M_4(\mathbb R)^{\oplus4}.
\]

The \textbf{graded tensor product} uses the native Clifford parities: \[
\sigma_{\rm gr}((a,b),(c,d))
=(-1)^{|b||c|}\sigma_-(a,c)\sigma_-(b,d).
\] Its six generators anticommute pairwise and square to \(-1\), so \[
A_{\rm gr}=A_x\widehat\otimes A_p\cong\mathrm{Cl}_{0,6}.
\] The additional exponent is bilinear and therefore also satisfies the cocycle identity. The center of this algebra is the scalar line: for an ordered Clifford monomial of degree \(k\), commuting with a generator outside its support requires \(k\) even, while commuting with one inside its support requires \(k\) odd. A nonempty proper support cannot meet both conditions; the full six-generator support also fails. Only the empty support survives. This center calculation already distinguishes the two products without appealing to a classification table.

Ordinary exchange \[
\tau_{\rm ord}(u\otimes v)=v\otimes u
\] is an involutive algebra automorphism of \(A_{\rm ord}\). Its trace is eight, and its rank split on the 64-dimensional carrier is \(36+28\). Graded exchange \[
\tau_{\rm gr}(u\otimes v)=(-1)^{|u||v|}v\otimes u
\] is an involutive algebra automorphism of \(A_{\rm gr}\). Its trace is \(4-4=0\), so its rank split is \(32+32\). The signs enter because equal odd basis monomials acquire a minus sign under graded exchange. These exchanges act on real Clifford carriers; the protected complex spinor exchanges of G03 act on different spaces.

The product Walsh matrix is \(H_x\otimes H_p\). Dividing it by eight gives an orthogonal 64-character transform. Its 63 nonconstant columns split, using the declared factorization, into \[
\underbrace{(\widehat G_x\setminus\{0\})\times\{0\}}_{7}
\ \sqcup\
\underbrace{\{0\}\times(\widehat G_p\setminus\{0\})}_{7}
\ \sqcup\
\underbrace{(\widehat G_x\setminus\{0\})\times
(\widehat G_p\setminus\{0\})}_{49}.
\] Exchange swaps the two seven-column sets and preserves the 49-column set. This is the marginal/mixed input decomposition used next.

\hypertarget{p1b:der:g04:s25}{%
\subsubsection{Resolution modules, restrictions, and exchange}\label{p1b:der:g04:s25}}

For any finite-dimensional unital real algebra \(A\), the quotient \(W=A/\mathbb R1\) here is a vector-space quotient. The scalar line is not an ideal: multiplying \(1\) by a nonscalar element leaves that line. For \(A=\mathrm{Cl}_{0,3}\), define \[
D=\operatorname{Hom}_{\mathbb R}(W,A),\qquad
\dim D=(\dim A)(\dim W)=8\cdot7=56.
\] If \(\pi:A\to W\) is the quotient map, composition with \(\pi\) gives the canonical vector-space identification \[
D\cong\{T\in\operatorname{End}_{\mathbb R}(A):T(1)=0\}.
\] Every map on the right factors uniquely through \(W\). No Leibniz rule is imposed: \(D\) is a space of linear resolution data.

For \(Q\subset W\), \(\dim Q=q\), the restriction map fits into the canonical exact sequence \[
0\longrightarrow
\operatorname{Hom}_{\mathbb R}(W/Q,A)
\longrightarrow D
\xrightarrow{\mathrm{res}_Q}
\operatorname{Hom}_{\mathbb R}(Q,A)
\longrightarrow0.
\] Its surjectivity follows by extending a basis of \(Q\) to a basis of \(W\). Thus \(\dim D(Q)=8q\). A flag gives the dimension ladder \(0,8,\ldots,56\), together with restriction surjections in the reverse direction. To realize the small modules as nested subspaces of \(D\), specify a splitting. The Chevalley/Walsh inner product supplies one: \(f\mapsto f\circ P_Q\), using the orthogonal projector onto \(Q\). This makes the relationship between a flag and embedded resolution subspaces explicit.

In Walsh coordinates, \(D\) is the space of all real \(8\times7\) matrices indexed by cube nodes and nonconstant characters. The particular Walsh evaluation table is one specified element of this 56-dimensional matrix space; the dimension counts all independent linear maps, not independent parameters of the fixed Walsh matrix.

Now set \[
A_\times=A_x\otimes A_p,\qquad
W_\times=A_\times/\mathbb R(1\otimes1).
\] This paragraph uses the ordinary product and its ordinary exchange. The vector-space dimension calculations also apply to the underlying graded product carrier, but its exchange signs must be recomputed. Here \[
\dim A_\times=64,\qquad \dim W_\times=63.
\] There is a canonical exact sequence of input spaces \[
0\longrightarrow W_x\oplus W_p
\longrightarrow W_\times
\xrightarrow{\ \overline{\pi_x\otimes\pi_p}\ }
W_x\otimes W_p
\longrightarrow0.
\] The first map sends \(([a],[b])\) to \([a\otimes1+1\otimes b]\); changing representatives changes the result only by a scalar unit. The last map is induced by \(\pi_x\otimes\pi_p\). Its kernel is exactly the marginal input space. One can verify this using any decompositions \(A_i=\mathbb R1\oplus W_i\); the exact sequence itself does not require selecting them.

Under the source's Chevalley scalar complements the sequence splits as \[
W_\times\cong W_x\oplus W_p\oplus(W_x\otimes W_p),
\qquad 63=7+7+49.
\] Consequently the \textbf{complete product resolution module} is \[
\boxed{
D_{\rm full}
=\operatorname{Hom}_{\mathbb R}(W_\times,A_\times),
\qquad \dim D_{\rm full}=64\cdot63=4032.
}
\] It decomposes, with these complements, into two marginal-input summands of dimension \(64\cdot7=448\) each and one mixed-input summand of dimension \(64\cdot49=3136\).

The \textbf{mixed-input resolution module} is intrinsically \[
\boxed{
D_{\rm mix}
=\operatorname{Hom}_{\mathbb R}(W_x\otimes W_p,A_x\otimes A_p)
\cong D_x\otimes D_p,
\qquad \dim D_{\rm mix}=56^2=3136.
}
\] The Hom-tensor identification sends \(F\otimes G\) to \([a]\otimes[b]\mapsto F([a])\otimes G([b])\). It is an isomorphism because tensor products of the two matrix-unit bases give a basis of the full target Hom space. Precomposition with the quotient \(W_\times\to W_x\otimes W_p\) canonically embeds \(D_{\rm mix}\) into \(D_{\rm full}\). Equivalently, it consists of maps that vanish on both marginal input spaces. Applying contravariant Hom to the previous exact sequence gives \[
0\longrightarrow D_{\rm mix}
\longrightarrow D_{\rm full}
\longrightarrow
\operatorname{Hom}_{\mathbb R}(W_x\oplus W_p,A_\times)
\longrightarrow0,
\] with dimensions \(3136,4032,896\).

There is also a canonical \textbf{factor-local tensor-sum module}. For \(F_x\in D_x,F_p\in D_p\), let \(\widetilde F_i=F_i\pi_i\), so \(\widetilde F_i(1)=0\). Then \[
\widetilde F_x\otimes I+I\otimes\widetilde F_p
\] kills \(1\otimes1\), and therefore defines an element of \(D_{\rm full}\). Evaluation on \(a\otimes1\) and \(1\otimes b\) proves that this map from \(D_x\oplus D_p\) is injective. Its image \(D_{\rm loc}\) has dimension \[
\boxed{\dim D_{\rm loc}=56+56=112.}
\] The same evaluations show \(D_{\rm loc}\cap D_{\rm mix}=0\). Their sum has dimension 3248, still smaller than \(D_{\rm full}\). The tensor-sum formula is a linear construction; its summands need not be algebra derivations.

For partial inputs \(Q_x\subset W_x\), \(Q_p\subset W_p\), write \(\dim Q_i=q_i\). The corresponding mixed module has dimension \[
64q_xq_p=(8q_x)(8q_p).
\] With the Chevalley complements, the full partial input \(Q_x\oplus Q_p\oplus(Q_x\otimes Q_p)\) has dimension \(q_x+q_p+q_xq_p\). Hence its resolution module has dimension \[
64(q_x+q_p+q_xq_p).
\] These formulas include zero-dimensional and full-input cases. They describe restrictions to the specified input spaces; embedded partial modules additionally use the declared orthogonal extensions.

To obtain exchange ranks, let \(\tau_A\) be the ordinary flip on \(A_\times\), and let \(\tau_W\) be its induced map on \(W_\times\). Conjugation on resolution data is \[
\mathscr B(F)=\tau_A F\tau_W^{-1}.
\] The trace of \(\tau_A\) is eight. Removing its fixed scalar line leaves \(\operatorname{Tr}\tau_W=7\). Since \(\operatorname{Hom}(W_\times,A_\times)\cong W_\times^*\otimes A_\times\), \[
\operatorname{Tr}\mathscr B=7\cdot8=56.
\] It follows that \[
\boxed{
\dim D_{\rm full}^{+}=2044,\qquad
\dim D_{\rm full}^{-}=1988.
}
\] Explicitly, \(A_\times\) has exchange dimensions \(36,28\) and \(W_\times\) has dimensions \(35,28\). Even maps preserve those eigenspaces and odd maps interchange them, giving \[
35\cdot36+28\cdot28=2044,\qquad
35\cdot28+28\cdot36=1988.
\]

On \(D_{\rm mix}\cong D\otimes D\), exchange is the usual flip. Its trace is 56 and \[
\boxed{
\dim D_{\rm mix}^{+}=\frac{56\cdot57}{2}=1596,\qquad
\dim D_{\rm mix}^{-}=\frac{56\cdot55}{2}=1540.
}
\] On \(D_{\rm loc}\cong D\oplus D\), exchange swaps the two summands and has ranks \(56,56\). For matched partial inputs \(Q_x=Q_p=Q\), \(\dim Q=q\), the mixed exchange ranks are \((8q)(8q\pm1)/2\). On the full partial module they are \[
\frac{64(q^2+2q)\pm8q}{2};
\] the partial input flip has trace \(q\), and the output flip has trace eight.

Under Paper I's inverse-dimension definition, the weights of \(D\), \(D_{\rm loc}\), \(D_{\rm mix}\), and \(D_{\rm full}\) are, respectively, \[
\frac1{56},\qquad \frac1{112},\qquad
\frac1{3136},\qquad \frac1{4032}.
\] Each number refers to the specified real vector space. The full and mixed product modules have different input domains and therefore different weights.

\hypertarget{p1b:der:g04:s26}{%
\subsubsection{Scalar evaluation and the scalar torus factor}\label{p1b:der:g04:s26}}

The scalar subcarrier \[
H_1^{\rm sc}
=(0,0)\oplus(\tfrac12,\tfrac12)
=\operatorname{span}_{\mathbb C}\{1,q_0,q_1,q_2,q_3\}
\] has dimension five, where \(q_a=x_a/R\) are unit-sphere coordinates. For normalized round measure, \(\{1,2q_0,2q_1,2q_2,2q_3\}\) is orthonormal. Directly from the 24 Hurwitz nodes, \[
\sum_{q\in\mathsf H}q=0,\qquad
\sum_{q\in\mathsf H}qq^{\mathsf T}=6I_4.
\] These two identities prove the same orthonormality for normalized counting measure \(24^{-1}\sum\), without needing the full design theorem. Thus evaluation \[
\mathrm{ev}:H_1^{\rm sc}\to\mathbb C^{\mathsf H}
\] is isometric with that normalization. Its image is \(\ker(A_+-8I)\oplus\ker(A_+-4I)\subset\mathscr E_{15}\): constants have adjacency eigenvalue eight, and the sum of the eight neighbors of \(q\) is \(4q\), so coordinates have eigenvalue four. The five independent columns exhaust those two eigenspaces.

For the Casimir \(K_1=C_L+C_R\) from G03, \[
\mathrm{ev}\,K_1=\frac38(8I-A_+)\mathrm{ev},\qquad
\mathrm{ev}\,\Delta_{S_R^3}
=\frac{3}{4R^2}(8I-A_+)\mathrm{ev}.
\] The constants have Casimir zero and the four-coordinate sector has Casimir \(3/2\); the displayed operators agree on both sectors.

Tensoring these isometries gives \[
\boxed{
\mathrm{ev}_x\otimes\mathrm{ev}_p:
H_{1,x}^{\rm sc}\otimes H_{1,p}^{\rm sc}
\hookrightarrow
\mathscr E_{15,x}\otimes\mathscr E_{15,p},
\qquad \dim\operatorname{im}=25.
}
\] The target function inner product is \(576^{-1}\sum\). In the counting convention of Section~\ref{p1b:der:g04:s22}, the corresponding isometry is \((\mathrm{ev}_x\otimes\mathrm{ev}_p)/24\). The product identities are \[
\mathrm{ev}_\times(K_x+K_p)
=\frac38(16I-A_+\otimes I-I\otimes A_+)\mathrm{ev}_\times,
\] \[
\mathrm{ev}_\times(\Delta_x+\Delta_p)
=
\left[
\frac3{4R_x^2}(8I-A_+)\otimes I
+\frac3{4R_p^2}I\otimes(8I-A_+)
\right]\mathrm{ev}_\times.
\] These follow by applying the one-factor identity in each tensor leg. The product image is invariant under factor exchange with matched factor identifications; its exchange ranks are \(15,10\).

This bridge has precisely the stated scalar domain. Under the independent natural left-\(2T_x\times2T_p\) action, 25 is also the largest possible rank of any map from \(H_{1,x}\otimes H_{1,p}\) into \(\mathscr E_{15,x}\otimes\mathscr E_{15,p}\). Paper~I-D, Section~\ref{p1b:app:g04:2} proves that sharp rank statement and distinguishes general product intertwiners from simple tensor maps.

For the finite scalar-character carrier of the existing \(M_8\) family, let \[
\mathscr T_F=\operatorname{span}_{\mathbb C}
\{e^{i(n_1\theta_1+n_2\theta_2)}:n\in F\subset\mathbb Z^2\},
\qquad f=|F|,
\] with normalized torus measure. Tensor every complex finite-support space, projection, and isometry in Section~\ref{p1b:der:g04:s21}--\ref{p1b:der:g04:s23} and this section with \(\mathscr T_F\), leaving the displayed sphere operators tensored with the identity. Ranks multiply by \(f\); the corner normalized-trace ratios remain unchanged: \[
\frac{9f}{576f}=\frac1{64},\qquad
\frac{36f}{576f}=\frac1{16}.
\] The scalar evaluation domain has dimension \(25f\), with image in the \(225f\)-dimensional coarse shell carrier. If the scalar Laplacian has eigenvalues \(\lambda_T(n)\), its product intertwiner adds \(I\otimes\operatorname{diag}_{n\in F}\lambda_T(n)\). No metric on the scalar torus beyond the chosen source branch is needed for the finite vector-space constructions.

If an exchange action on \(\mathscr T_F\) is specified and has involutive matrix \(B_F\), write \(t_F=\operatorname{Tr}B_F\). For a sphere-product space of dimension \(d\) with exchange trace \(t\), the total exchange ranks are \[
\frac{df\pm t\,t_F}{2}.
\] For coordinate swap \(n\mapsto(n_2,n_1)\), the set \(F\) must be closed under that swap and \(t_F\) is the number of fixed labels. If exchange acts trivially on the torus carrier, \(t_F=f\). This specifies when the previously computed exchange ranks simply multiply by \(f\).

Paper~I-D, Section~\ref{p1b:app:g04} proves the convention-comparison identities and extended-case formulas used in the results above.


\hypertarget{p1b:sec:g05}{%
\subsection{Contact pairs and product control geometry}\label{p1b:sec:g05}}


\begin{proposition}[Contact pair and product control distance]
\label{p1b:prop:g05-product}
Fix the marked round sphere factors and a positive flat torus metric.
The two sphere contact forms define a contact pair of type \((1,1)\)
on \(P_v=S_x^3\times S_p^3\). Its horizontal distribution has growth
vector \((4,6)\) and homogeneous dimension eight. On
\(M_8=P_v\times T_G^2\), adjoin the two torus directions to the
horizontal distribution. The growth vector becomes \((6,8)\), and
the homogeneous dimension is ten. With the orthogonal horizontal
metric, the control distance satisfies
\[
 d_{\mathcal H}^{\,2}
 =d_{H,x}^{\,2}+d_{H,p}^{\,2}+d_{T_G^2}^{\,2}.
\]
For the product distribution an abnormal curve has at least one sphere
factor stationary throughout its interval. Every length minimizer
also admits a normal lift. The contact-pair assertion concerns \(P_v\);
on \(M_8\) the extra torus directions are included in the distribution.
\end{proposition}

The bracket calculation and separated energy prove the growth and distance assertions. The multiplier equations and factor geodesics below prove the statements about abnormality and minimizers.

\hypertarget{p1b:der:g05:s28}{%
\subsubsection{Marked Hopf data and the control metric}\label{p1b:der:g05:s28}}

This derivation covers contact-pair geometry, reachability, normal and abnormal curves, exact product distances and cut behavior, nilpotent tangents, metric-measure estimates, horizontal scalar analysis, and exchange. It uses the smooth round sphere factors, independently of finite enriched carrier levels.

Let \[
P_v=S^3(R_x)\times S^3(R_p),\qquad
M_8=P_v\times T^2,\qquad R_x,R_p>0,
\] and let the scalar angles have period \(2\pi\). For the full control metric choose any positive definite constant matrix \(G\) on \(T^2\), as in the metric family of this paper. Write \[
g=g_x\oplus g_p\oplus G,\qquad
\alpha_i=R_i\eta_{0,i},\qquad H_i=\ker\alpha_i .
\] The sphere forms have the normalization of Paper I, p1:prop:contact-normalization: \[
\eta_{0,i}(X_{H,i})=1,\quad
\xi_i=R_i^{-1}X_{H,i},\quad
\alpha_i(\xi_i)=1,\quad
\operatorname{period}(X_{H,i})=2\pi.
\]

For a consistent explicit frame, represent \(S^3(R)\) by unit quaternions \(q\) with metric \(R^2g_{\rm unit}\), and use the source's common-phase action \(q\mapsto e^{i\theta}q\). Set \[
X(q)=\frac{jq}{R},\qquad Y(q)=\frac{kq}{R},\qquad
Z(q)=\frac{iq}{R}=\xi(q).
\] These are right-invariant vector fields and are orthonormal. For such fields the Lie bracket is the negative of the quaternion Lie-algebra commutator: \[
[X,Y]=-\frac2R Z,\qquad
[Y,Z]=-\frac2R X,\qquad
[Z,X]=-\frac2R Y.
\] It follows that \[
d\alpha(X,Y)=-\alpha([X,Y])=\frac2R,\qquad
\alpha\wedge d\alpha=\frac2R\,d\mathrm{vol}_{g}.
\] The frame fixes the positive Hopf orientation directly.  Paper I-A,
Proposition~\ref{p1:prop:marked-hopf-inversion-transport}, uses this same
right-invariant common-phase structure with its second horizontal vector
reversed.  Its opposite left-invariant copy is separately defined by
inversion; the associated circle parameter and orientation are transported
there.  Right translation preserves the primary model used here. In particular, \[
d_{CC,R}(q_0,q_1)=d_{CC,R}(1,q_1q_0^{-1}).
\] All distance and dimension formulas are unchanged under a consistent exchange of handedness.

\begin{corollary}[Same-marking frame and opposite-marking transport]
\label{p1b:cor:marked-hopf-frame-transport}
Let \(E_1,E_2,\xi_\eta\) be the metric-normalized primary frame of
Paper I-A, Proposition~\ref{p1:prop:marked-hopf-inversion-transport}.
On the same marked radius-\(R\) sphere,
\[
 (E_1,E_2,\xi_\eta)=(X,-Y,Z),\qquad X_H=RZ.
\]
Thus the two primary frames define the identical horizontal metric,
positive scalar operator \(-X^2-Y^2\), principal Hopf charge and contact
form.  If the transverse complex structure is written \(JX=Y\) in the
present frame, the same tensor is \(JE_1=-E_2\) in Paper I-A's frame.
No Hopf charge reversal accompanies this horizontal basis change.
For the opposite marked copy, inversion pulls back the contact form and
orientation and sends the common-phase orbit \(e^{i\alpha}q\) to
\(q^{-1}e^{-i\alpha}\).  Charge reverses only if that latter circle is
reparametrized as a positive right orbit.
\end{corollary}
\begin{proof}
Insert the explicit quaternion fields in the two frames.  The sign in
the second horizontal direction disappears from its squared norm and
squared differential operator, while the common vertical field is \(Z\).
The displayed expression for \(J\) is a change of basis of the same tensor.
The inversion identity and its orientation/charge consequences are proved
in the cited Paper I-A proposition, so no invariant-frame label is
identified merely from the sign of its brackets.
\end{proof}

\textbf{Source inputs.} The principal Paper I locators are p1:sec:sub-riemannian, p1:subsec:cc-global, p1:prop:nilpotent-approx, p1:cor:hausdorff-dim, p1:thm:sr-poincare, p1:thm:sr-sobolev, and p1:subsec:sub-laplacian-spectrum. Paper 1B already supplies the horizontal charge floors in p1b:thm:charge-bundles; Section~\ref{p1b:der:g05:s35} below derives their full product operator and scalar-torus extension.

The general inputs used below are stated explicitly: the one-factor global normal synthesis; the equiregular tangent-cone and ball-box theorems; and the local Hörmander Poincaré/Sobolev results already imported in Paper I. These inputs are not purportedly proved by finite matrix checks. The one-factor synthesis and its normalization were also checked against \href{https://arxiv.org/html/0709.3997v2}{Boscain--Rossi, §§3.1 and 5.1}.

\hypertarget{p1b:der:g05:s29}{%
\subsubsection{Contact pair, commuting Reeb fields, and reachability}\label{p1b:der:g05:s29}}

Pull the two sphere contact forms to \(P_v\). They satisfy \[
(d\alpha_x)^2=(d\alpha_p)^2=0,\qquad
\Omega_{\rm pair}
=\alpha_x\wedge d\alpha_x\wedge\alpha_p\wedge d\alpha_p
=\frac4{R_xR_p}\,d\mathrm{vol}_{P_v}\ne0.
\] Thus they define a contact pair of type \((1,1)\) on \(P_v\). This follows directly from the product forms and agrees with the definition and product example in \href{https://arxiv.org/html/math/0305381v2}{Bande--Hadjar, §§2.1--2.2}. Its two Reeb fields are the factor lifts \(Z_x,Z_p\), since \[
\alpha_i(Z_j)=\delta_{ij},\qquad
\iota_{Z_j}d\alpha_i=0,\qquad [Z_x,Z_p]=0.
\] The first two equations determine them uniquely on \(P_v\): the common kernel of \(d\alpha_x,d\alpha_p\) is exactly their two-plane. The characteristic distributions \(\ker\alpha_x\cap\ker d\alpha_x\) and \(\ker\alpha_p\cap\ker d\alpha_p\) are respectively the other sphere's tangent bundle. Their leaves are the two sphere-factor foliations. The contact-pair volume integral is \[
\int_{P_v}\Omega_{\rm pair}=16\pi^4R_x^2R_p^2.
\]

The horizontal distribution on \(P_v\) is \[
\mathcal D_v=H_x\oplus H_p
=\ker\alpha_x\cap\ker\alpha_p,\qquad
g_{\mathcal D_v}=g_{H_x}\oplus g_{H_p}.
\] In the factor frame, cross brackets vanish, while \[
[X_x,Y_x]=-\frac2{R_x}Z_x,\qquad
[X_p,Y_p]=-\frac2{R_p}Z_p.
\] These brackets produce exactly the two missing directions. Hence \[
\boxed{\operatorname{rank}\mathcal D_v=4,\qquad
\mathcal D_v+[\mathcal D_v,\mathcal D_v]=TP_v,\qquad
\text{growth }(4,6).}
\] Here a bracket span means the span of values of brackets of local horizontal sections. The displayed global frame suffices to calculate it.

On all of \(M_8\), the full control distribution is \[
\boxed{
\mathcal D=H_x\oplus H_p\oplus TT^2
=\ker\alpha_x\cap\ker\alpha_p,\qquad
g_{\mathcal D}=g_{H_x}\oplus g_{H_p}\oplus G.
}
\] The forms in this equality are pulled back from the spheres, so their common kernel includes the scalar tangent directions. Since constant torus fields commute with every sphere field, \[
\boxed{
\operatorname{rank}\mathcal D=6,\qquad
\mathcal D+[\mathcal D,\mathcal D]=TM_8,\qquad
\text{growth }(6,8).
}
\] Both structures are equiregular of step two and connect any two points of their respective connected manifolds by horizontal curves. Their control distances are finite and induce the manifold topology. Compactness gives completeness and existence of minimizing curves.

If instead one embeds only \(H_x\oplus H_p\) into \(TM_8\), omitting torus velocities, its accessible leaves are exactly \[
P_v\times\{\theta\}.
\] Every iterated bracket has zero scalar component; every two points within one leaf are reachable by the preceding sphere-product argument. The resulting extended distance is finite precisely for equal scalar coordinates. This is the leafwise construction, whereas \(\mathcal D\) above gives a metric on all of \(M_8\). The contact pair itself lives on the six-dimensional factor \(P_v\).

For the Reeb two-torus, a vertical flow \(aZ_x+bZ_p\) has principal angular frequencies \(a/R_x,b/R_p\). It is periodic exactly when its nonzero frequencies are commensurable; with an irrational nonzero ratio its orbit is dense in the fiber two-torus. For a primitive integer pair \((r,s)\), \(\gcd(|r|,|s|)=1\), the vector field \[
rX_{H,x}+sX_{H,p}
\] has period \(2\pi\) and constant norm \(\sqrt{r^2R_x^2+s^2R_p^2}\). Its unit-speed orbit has period \[
2\pi\sqrt{r^2R_x^2+s^2R_p^2}.
\] These statements follow from the free product circle action. The Reeb fiber torus and the independent scalar torus are separate factors.

\hypertarget{p1b:der:g05:s30}{%
\subsubsection{Nilpotent tangent, ball-box scaling, and dimension}\label{p1b:der:g05:s30}}

At a point in a sphere factor, write the vertical bracket coefficient as \(\kappa_i=-2/R_i\) in the chosen unit-Reeb frame. In privileged coordinates \((x_i,y_i,t_i)\), the leading horizontal fields are \[
\widehat X_i=\partial_{x_i}
-\frac{\kappa_i}{2}y_i\partial_{t_i},\qquad
\widehat Y_i=\partial_{y_i}
+\frac{\kappa_i}{2}x_i\partial_{t_i}.
\] The signed rescaling \(z_i=t_i/\kappa_i\) gives \[
\widehat X_i=\partial_{x_i}-\frac12y_i\partial_{z_i},\qquad
\widehat Y_i=\partial_{y_i}+\frac12x_i\partial_{z_i},\qquad
[\widehat X_i,\widehat Y_i]=\partial_{z_i}.
\] Cross-factor brackets remain zero. Therefore the graded tangent algebras and their horizontal metrics are \[
\operatorname{gr}(TP_v)=\mathfrak h_1\oplus\mathfrak h_1,\qquad
\operatorname{gr}(TM_8)=
\mathfrak h_1\oplus\mathfrak h_1\oplus\mathbb R^2_G.
\] The two Euclidean scalar directions have degree one. The two Reeb quotient directions have degree two. Integrating these graded algebras gives the metric tangents \[
\boxed{
\operatorname{Tan}(P_v,d_v)=\mathbb H^1\times\mathbb H^1,\qquad
\operatorname{Tan}(M_8,d_{\mathcal D})
=\mathbb H^1\times\mathbb H^1\times\mathbb R^2_G,
}
\] with the orthogonal product horizontal metric. Precisely, at each base point the pointed spaces with distance \(\varepsilon^{-1}d\) converge to the stated tangent as \(\varepsilon\downarrow0\). Equiregularity and the ball-box theorem supply the metric convergence, not merely formal convergence of frames.

In scalar orthonormal coordinates \(u\in\mathbb R^2\), the full dilations are \[
\delta_\varepsilon(x_x,y_x,z_x,x_p,y_p,z_p,u)
=(\varepsilon x_x,\varepsilon y_x,\varepsilon^2z_x,
  \varepsilon x_p,\varepsilon y_p,\varepsilon^2z_p,
  \varepsilon u).
\] They scale the tangent control distance by \(\varepsilon\). The sum of the dilation weights gives \[
\boxed{
\dim_H(P_v,d_v)=4+2\cdot2=8,\qquad
\dim_H(M_8,d_{\mathcal D})=6+2\cdot2=10.
}
\] These are Hausdorff dimensions for the specified control metrics; the smooth manifold dimensions remain six and eight.

The local ball-box estimate has the form \[
d_{\mathcal D}(o,z)\asymp
|x_x|+|y_x|+|x_p|+|y_p|+\|u\|
+\sqrt{|z_x|}+\sqrt{|z_p|}.
\] The constants are uniform on the compact homogeneous geometry. It explains directly why two missing velocity directions still become reachable through horizontal loops, while carrying weight two in the distance and volume scaling.

\hypertarget{p1b:der:g05:s31}{%
\subsubsection{Exact product distance and minimizers}\label{p1b:der:g05:s31}}

Let \(d_x,d_p\) denote the two sphere control distances, and let \(d_T\) be the flat torus distance associated with \(G\). Then \[
\boxed{
d_v((x,p),(x',p'))^2=d_x(x,x')^2+d_p(p,p')^2,
}
\] \[
\boxed{
d_{\mathcal D}((x,p,\theta),(x',p',\theta'))^2
=d_x(x,x')^2+d_p(p,p')^2+d_T(\theta,\theta')^2.
}
\] To prove this, use the energy without a factor \(1/2\) on a fixed interval \([0,T]\): \[
E(\gamma)=\int_0^T
\bigl(|\dot\gamma_x|_{H_x}^2+
      |\dot\gamma_p|_{H_p}^2+
      |\dot\gamma_T|_G^2\bigr)\,dt.
\] Each factor contributes at least its squared endpoint distance divided by \(T\), by Cauchy--Schwarz and the definition of distance. Conversely, choose a minimizing path in each factor and parametrize it with constant speed \(d_i/T\) on the common interval. Their product attains the sum of those three bounds. Energy minimization at fixed time is equivalent to length minimization after constant-speed parametrization, proving the distance identities.

The equality conditions also classify all minimizers: after a common constant-speed parametrization, every factor is a constant-speed minimizer, with speed \(d_i/T\). A factor whose endpoints coincide is constant. This classification will be used for both abnormality and cut behavior.

The torus contribution is fully explicit: \[
d_T(\theta,\theta')^2
=\min_{n\in\mathbb Z^2}
(\theta'-\theta+2\pi n)^{\mathsf T}
G(\theta'-\theta+2\pi n).
\] Each shortest lattice representative supplies one minimizing torus segment. The sphere factors obey \(d_i\ge d_{g_i}\), since horizontal curves are an admissible subset of round-metric curves with the same length. Consequently \[
d_{\mathcal D}\ge d_g .
\] Homothety gives \(d_{CC,R}=R\,d_{CC,1}\) in unit-quaternion coordinates.

The distance formulas apply only to the displayed orthogonal product control metric. Arbitrary additional cross terms would define a different energy minimization problem.

\hypertarget{p1b:der:g05:s32}{%
\subsubsection{Normal synthesis, cut times, and global distance data}\label{p1b:der:g05:s32}}

In a sphere factor take initial horizontal speed \(v>0\), angle \(\phi\), and vertical covector component \(c\in\mathbb R\). Put \[
U_\phi=j\cos\phi+k\sin\phi,\qquad \lambda=\sqrt{v^2+c^2}.
\] In the common-phase convention of Section~\ref{p1b:der:g05:s28} the normal curve from \(q_0\) is \[
\boxed{
q(t)=
\exp\!\left(-\frac{ci t}{R}\right)
\exp\!\left(\frac{(vU_\phi+ci)t}{R}\right)q_0.
}
\] Differentiating gives \[
\dot q\,q^{-1}
=\frac vR
\operatorname{Ad}_{\exp(-ci t/R)}U_\phi .
\] This is horizontal and has round-metric norm \(v\). Its controls rotate through angle \(-2ct/R\). For the normal Hamiltonian \[
\mathsf H_R=\frac12(h_X^2+h_Y^2),
\] the right-invariant bracket convention gives \[
\dot h_X=\frac{2c}{R}h_Y,\qquad
\dot h_Y=-\frac{2c}{R}h_X,\qquad \dot c=0.
\] Thus the exponential product solves the normal Hamiltonian equations, with the specified initial data.

For \(v>0\), the first conjugate and cut time of this factor curve is \[
\boxed{t_{\rm cut}=t_{\rm conj,1}=\frac{\pi R}{\sqrt{v^2+c^2}}.}
\] The global optimality assertion is the one-factor synthesis imported in Paper I and checked in Boscain--Rossi. Their orthonormal quaternion generators have size \(1/2\); replacing them by size \(1/R\) and reversing handedness gives the normalization here. The direct endpoint and Jacobian calculations below independently check the constants.

At \(t_{\rm cut}\), the second exponential is \(-1\), so \[
q(t_{\rm cut})q_0^{-1}
=\exp(i\vartheta),\qquad
\vartheta=\pi\left(1-\frac c{\sqrt{v^2+c^2}}\right)\in(0,2\pi).
\] The horizontal angle disappears. Eliminating \(c/\lambda\) gives the exact vertical-fiber distance \[
\boxed{
d_{CC,R}(q_0,e^{i\vartheta}q_0)
=R\sqrt{\vartheta(2\pi-\vartheta)},
\qquad 0\le\vartheta\le2\pi .
}
\] The two endpoint values refer to the same base point and give zero. Every other vertical endpoint has a circle of minimizing curves. The cut locus of the factor, with distinct minimizing paths counted, is its Hopf circle through \(q_0\) with \(q_0\) removed.

For an arbitrary relative endpoint \[
q_1q_0^{-1}=z_1+z_2j,\qquad |z_1|^2+|z_2|^2=1,\quad z_2\ne0,
\] one obtains the distance by setting \(v=1\) and solving \[
z_1=e^{-ict/R}
\left(\cos\frac{\sqrt{1+c^2}t}{R}
 +i\frac c{\sqrt{1+c^2}}\sin\frac{\sqrt{1+c^2}t}{R}\right),
\] \[
|z_2|=\frac1{\sqrt{1+c^2}}
\sin\frac{\sqrt{1+c^2}t}{R},\qquad
0<t<\frac{\pi R}{\sqrt{1+c^2}}.
\] The unique pre-cut solution gives \(d_{CC,R}=t\); the phase of \(z_2\) recovers \(\phi\). The complex equation retains its argument branch and avoids using an ambiguous scalar arctangent.

The factor diameter is \(\pi R\), attained at \(-q_0\). Indeed every minimizing nonconstant path has length \(vt\le \pi Rv/\sqrt{v^2+c^2}\le\pi R\), while the vertical formula at \(\vartheta=\pi\) attains equality. Hence \[
\boxed{\operatorname{diam}(P_v,d_v)
=\pi\sqrt{R_x^2+R_p^2},}
\] \[
\boxed{\operatorname{diam}(M_8,d_{\mathcal D})
=\sqrt{\pi^2R_x^2+\pi^2R_p^2+\operatorname{diam}(T^2,G)^2}.}
\] Both follow by maximizing independent factor distances. For a rectangular torus \(G=\operatorname{diag}(a^2,b^2)\), its diameter is \(\pi\sqrt{a^2+b^2}\). For general \(G\), it is the covering radius of the corresponding \(2\pi\mathbb Z^2\) lattice.

The full normal Hamiltonian is \[
\mathsf H=\mathsf H_x+\mathsf H_p+
\frac12\zeta^{\mathsf T}G^{-1}\zeta .
\] It splits into two independent sphere systems and the flat torus motion \(\theta(t)=\theta_0+tG^{-1}\zeta\). The three factor speeds are constant. Along such a curve, \[
\boxed{t_{\rm cut}^{\rm product}
=\min\{t_{\rm cut}^x,t_{\rm cut}^p,t_{\rm cut}^T\}.}
\] For a stationary factor its cut time is defined to be \(+\infty\). In particular one must not substitute \(v=0\) into the nonconstant sphere cut formula: any vertical covector then produces the same constant path. The minimum rule follows from Section~\ref{p1b:der:g05:s31}: a product segment minimizes exactly as long as each factor segment minimizes.

Let \(K_x,K_p\) be the two factor cut sets through a fixed starting point, and \(K_T\) the translated boundary of a flat-torus Voronoi cell, equivalently the points with more than one shortest lattice representative. The locus of endpoints with multiple distinct minimizing paths is exactly \[
\boxed{
(K_x\times S_p^3\times T^2)
\ \cup\ (S_x^3\times K_p\times T^2)
\ \cup\ (S_x^3\times S_p^3\times K_T).
}
\] Every factor choice combines independently with the other minimizing paths. Distinct covector lifts of the same constant factor are not counted as distinct paths. This distinction is needed for the abnormal cases treated next.

\hypertarget{p1b:der:g05:s33}{%
\subsubsection{Abnormal curves and the regular conjugate-time formula}\label{p1b:der:g05:s33}}

For a nonzero annihilating covector \(\lambda=a\alpha_x+b\alpha_p\), define the Levi form by \(\mathcal L_\lambda(V,W)=\lambda([V,W]\bmod\mathcal D_v)\). On \(\mathcal D_v=H_x\oplus H_p\) it has two blocks \[
\mathcal L_\lambda
=
\begin{pmatrix}0&-2a/R_x\\2a/R_x&0\end{pmatrix}
\oplus
\begin{pmatrix}0&-2b/R_p\\2b/R_p&0\end{pmatrix}.
\] Its determinant is \[
\det\mathcal L_\lambda=\frac{16a^2b^2}{R_x^2R_p^2}.
\] It has rank four when \(ab\ne0\), and rank two when exactly one coefficient is nonzero. On \(\mathcal D\), a zero two-by-two block for the scalar tangent directions is added. Thus the product is bracket-generating but is not fat, although each sphere factor is contact and fat.

For a horizontal absolutely continuous curve on a fixed interval, the abnormal multiplier equations give \[
\lambda(t)=a(t)\alpha_x+b(t)\alpha_p,\qquad
\dot a=\dot b=0,
\] \[
a\,\dot\gamma_x=0,\qquad b\,\dot\gamma_p=0
\quad\text{almost everywhere}.
\] To see this directly, abnormal stationarity requires the covector to annihilate every horizontal control field, including the torus fields. Differentiating its pairing with \(X_x,Y_x\) gives \(a u_{x,1}=a u_{x,2}=0\) using the nonzero bracket coefficient; the \(p\) equations are identical. Pairing the adjoint equation with \(Z_x,Z_p\) shows \(a,b\) are constant, because their brackets with control fields are horizontal and hence already annihilated.

At least one of \(a,b\) is nonzero. Conversely, a constant sphere factor allows a nonzero constant multiple of its contact form as an abnormal multiplier, for arbitrary motion in the remaining factors. Therefore \[
\boxed{
\gamma\text{ admits an abnormal lift}
\quad\Longleftrightarrow\quad
\gamma_x\text{ or }\gamma_p
\text{ is constant throughout the interval}.
}
\] The condition concerns the whole interval. Merely switching which factor is momentarily stationary does not supply a single abnormal lift.

Nevertheless, every minimizing product curve admits a normal lift. Section~\ref{p1b:der:g05:s31} shows that it is a product of minimizing factor curves with constant speeds. Each nonconstant contact-factor minimizer is normal; a constant factor is produced by a zero normal horizontal covector; a torus minimizing segment is a flat normal geodesic. The sum Hamiltonian combines these lifts. Hence \[
\boxed{\text{There are no strictly abnormal minimizing curves.}}
\] Minimizers with a stationary sphere factor are both normal and abnormal. This is the precise product replacement for the one-factor absence of nontrivial abnormal curves.

From a fixed starting point \((x_0,p_0,\theta_0)\), the set of abnormal endpoints is exactly \[
\mathcal A=
(\{x_0\}\times S_p^3\times T^2)
\ \cup\
(S_x^3\times\{p_0\}\times T^2).
\] Reachability in each remaining factor proves equality. These are submanifolds of dimension five in the eight-dimensional ambient manifold, so \(\mathcal A\) has smooth-volume measure zero. On \(P_v\), the corresponding two pieces have dimension three.

For the regular normal exponential, assume both sphere speeds are nonzero. In a factor put \[
s=\frac{\lambda t}{R},\qquad \beta=\frac c\lambda,\qquad
|\beta|<1.
\] The parameterization in Section~\ref{p1b:der:g05:s32} has complex components \[
z_1=e^{-i\beta s}(\cos s+i\beta\sin s),\qquad
z_2=\sqrt{1-\beta^2}\sin s\,e^{i(\phi-\beta s)}.
\] An endpoint-volume Jacobian in parameters \((s,\beta,\phi)\), up to orientation, is \[
(1-\beta^2)\sin s\,(\sin s-s\cos s).
\] Paper~I-D, Section~\ref{p1b:app:g05:1} derives it explicitly. Thus all positive factor conjugate times are determined by \[
s=n\pi\quad(n\ge1),\qquad\text{or}\qquad
\tan s=s,\quad s>0.
\] The first is \(s=\pi\), since \(\sin s-s\cos s\) is positive on \((0,\pi)\): its derivative is \(s\sin s>0\). The torus exponential has derivative \(tG^{-1}\), which is invertible for \(t>0\). The regular product exponential has block-diagonal differential, so its conjugate times are the union of the two sphere lists, with kernel dimensions adding at simultaneous singular times.

For a stationary sphere factor, its endpoint does not depend on its vertical covector. The full normal exponential is then singular at every positive time. This persistent degeneracy is separated from the regular first-conjugate-time statement and from loss of minimizing optimality.

\hypertarget{p1b:der:g05:s34}{%
\subsubsection{Volumes, doubling, and horizontal inequalities}\label{p1b:der:g05:s34}}

Use the product of the round sphere volumes and the flat torus volume, \(\mu=\mu_x\otimes\mu_p\otimes\mu_T\). All these measures and the corresponding control metrics are homogeneous, so ball volume is independent of the center. Let the unit-sphere small-ball coefficient for round volume be \(C_*>0\). The one-factor expansion imported in Paper I and homothety give, for each fixed \(R>0\), \[
\mu_R(B_{CC}(q,r))
=C_Rr^4(1+O(r^2)),\qquad C_R=\frac{C_*}{R}.
\] Indeed \(\mu_R=R^3\mu_1\) and \(d_R=Rd_1\). The \(O(r^2)\) constant may depend on \(R\).

The exact squared-distance identity determines the product ball: in the variables \(s=d_x/r\) and \(t=d_p/r\), its radial domain is \(s^2+t^2<1\). Passing to rescaled radial volume measures gives \[
\mu_{P_v}(B_v(r))
=\frac{C_xC_p}{6}\,r^8(1+o(1)).
\] The coefficient follows from \[
\int_0^1 4s^3(1-s^2)^2\,ds=\frac16.
\] For \(M_8\), sufficiently small torus balls have exact volume \(\pi\rho^2\). Integrating the remaining torus radius gives \[
\boxed{
\mu(B_{\mathcal D}(r))
=\frac{\pi C_xC_p}{30}\,r^{10}(1+o(1)).
}
\] For an explicit coefficient check set \(u=s^2,v=t^2\); the limiting radial measure is \(4uv\,du\,dv\), and \[
\int_{\substack{u,v\ge0\\u+v\le1}}4uv\,du\,dv=\frac16,\qquad
\int_{\substack{u,v\ge0\\u+v\le1}}4\pi uv(1-u-v)\,du\,dv=\frac\pi{30}.
\] The limit follows from the one-factor radial asymptotics, by weak convergence of finite rescaled radial measures on a compact interval; the boundary has zero limiting radial measure. No differentiation of an unspecified error term is required. The displayed claims use \(1+o(1)\), which is sufficient for the following estimates.

For normalized probability measure divide these leading constants by the total product volume. For any other smooth invariant positive volume, including Popp volume, multiply by its fixed density ratio. The metric exponents do not depend on that normalization.

For \(Q=8\) on \(P_v\), or \(Q=10\) on \(M_8\), there are constants \(0<c\le C<\infty\), depending on the fixed geometric parameters, such that \[
cr^Q\le\mu(B(z,r))\le Cr^Q,\qquad
0<r\le\operatorname{diam}.
\] The small-radius bounds follow from the preceding asymptotic. For larger radii, monotonicity bounds ball volume between a fixed small-ball volume and total volume; adjusting constants proves the all-scale statement. This also proves doubling. No continuity of ball volume is needed for that extension.

The global smooth horizontal frame satisfies the step-two Hörmander condition. The local Hörmander Poincaré theorem used in Paper I therefore applies here, initially with a fixed enlargement of the energy ball. Compactness makes its local constants uniform. Together with the doubling and lower-volume bounds just proved, the same metric Sobolev and John-domain chaining argument used in p1:thm:sr-poincare yields \[
\int_B|f-f_B|^2\,d\mu
\le C r^2\int_B|\nabla_{\mathcal D}f|^2\,d\mu.
\] For completeness, the ball-domain hypothesis has a direct verification: a minimizing path from a point in \(B(z,r)\) toward \(z\), after travelling length \(s\), is at distance at least \(s\) from the complement of the ball. Thus these geodesic balls have a uniform John-domain constant. The horizontal gradient also has the truncation property required by the cited metric theorem.

The critical horizontal \(W^{1,2}\) exponent is \(2Q/(Q-2)\). Consequently the compact-space Sobolev inequalities are \[
\boxed{
\|f\|_{L^{8/3}(P_v)}
\le C\bigl(\|\nabla_{\mathcal D_v}f\|_{L^2}
           +\|f\|_{L^2}\bigr),
}
\] \[
\boxed{
\|f\|_{L^{5/2}(M_8)}
\le C\bigl(\|\nabla_{\mathcal D}f\|_{L^2}
           +\|f\|_{L^2}\bigr).
}
\] They hold for the horizontal Sobolev completions by density. The inhomogeneous term accounts for constants on the compact space. The exact best global \(L^2\) Poincaré constant is obtained spectrally in the next section; the local constants here are not claimed sharp.

\hypertarget{p1b:der:g05:s35}{%
\subsubsection{Horizontal scalar operator and finite spectral exports}\label{p1b:der:g05:s35}}

The invariant horizontal fields are divergence-free for the product volume. Their nonnegative scalar operator is therefore \[
\boxed{
L=-X_x^2-Y_x^2-X_p^2-Y_p^2
-\sum_{a,b=1}^2(G^{-1})^{ab}\partial_{\theta_a}\partial_{\theta_b}.
}
\] Its quadratic form is \(\int|\nabla_{\mathcal D}f|^2d\mu\). The operator is hypoelliptic by the verified bracket condition. On a sphere, \[
L_i=\Delta_{S_i^3}+Z_i^2
=\Delta_{S_i^3}-\frac{Q_i^2}{R_i^2},
\qquad Q_i=-iX_{H,i}.
\] The geometric Hopf weights are the integers \(q_i\), rather than the half-normalized protected weights used elsewhere in the finite spinor carrier.

The scalar Peter--Weyl basis gives the complete labeled spectrum \[
\boxed{
\lambda=
\frac{k_x(k_x+2)-q_x^2}{R_x^2}
+\frac{k_p(k_p+2)-q_p^2}{R_p^2}
+n^{\mathsf T}G^{-1}n ,
}
\] \[
k_i\in\mathbb Z_{\ge0},\quad
q_i=-k_i,-k_i+2,\ldots,k_i,\qquad n\in\mathbb Z^2,
\] with labeled multiplicity \((k_x+1)(k_p+1)\). Distinct labels with equal eigenvalues contribute additive multiplicities. The eigenbasis is complete by the scalar Peter--Weyl and torus Fourier decompositions; since \[
k(k+2)-q^2\ge2k,
\] only finitely many labels lie below any finite spectral bound. This also proves compact resolvent for the self-adjoint spectral realization of \(L\).

For fixed nonzero or zero bi-Hopf charges, the sphere contribution has the exact floor \[
\frac{2|q_x|}{R_x^2}+\frac{2|q_p|}{R_p^2},
\] attained at \(k_i=|q_i|\). This recovers the existing Paper 1B charge-floor theorem and shows its operator origin. The full kernel consists of constants. Put \[
\lambda_T=\min_{n\in\mathbb Z^2\setminus\{0\}}
n^{\mathsf T}G^{-1}n>0.
\] Then \[
\boxed{
\lambda_1(L)=
\min\left\{\frac2{R_x^2},\frac2{R_p^2},\lambda_T\right\}.
}
\] Each sphere minimum is attained by \(k=1,q=\pm1\), with total multiplicity four. Tensoring with constants in the other factors attains the product minimum. All other modes have at least that energy. Spectral expansion now gives the sharp global estimate \[
\int|f-\bar f|^2\,d\mu
\le\lambda_1(L)^{-1}\int|\nabla_{\mathcal D}f|^2\,d\mu.
\] For \(P_v\), omit the torus term in the minimum and in the operator.

The commuting factor operators also give the exact semigroup identity \[
e^{-tL}=e^{-tL_x}\otimes e^{-tL_p}\otimes e^{-t\Delta_T},
\qquad t\ge0.
\] This is an export for the later analytic derivation; no additional heat asymptotic is needed here.

For arbitrary scalar harmonic cutoffs use the notation from G03: \[
\mathcal P_{\ell_x}\otimes\mathcal P_{\ell_p}\otimes\mathscr T_F,
\qquad
\mathcal P_\ell=\bigoplus_{k=0}^{\ell}U_{kk}.
\] Restriction simply imposes \(k_i\le\ell_i\), \(n\in F\) in the displayed spectrum. These scalar spaces reduce \(L\), and their dimension is \[
\frac{(\ell_x+1)(\ell_x+2)(2\ell_x+3)}6
\frac{(\ell_p+1)(\ell_p+2)(2\ell_p+3)}6\,|F|.
\] The finite kernel has dimension one if \(0\in F\), and zero otherwise. When \(0\in F\), the first positive eigenvalue is the minimum of the present candidates \[
\left\{\frac2{R_i^2}:\ell_i\ge1\right\}
\ \cup\
\{n^{\mathsf T}G^{-1}n:n\in F\setminus\{0\}\}.
\] If that set is empty the finite space contains only the zero mode. If \(0\notin F\), its lowest eigenvalue is \(\min_{n\in F}n^{\mathsf T}G^{-1}n\), provided \(F\ne\varnothing\). Finite cutoff changes affect retained scalar modes, while the smooth control distribution and its growth vector are the same. This scalar-function export does not define a horizontal connection on the additional form and spinor summands of the enriched carriers.

\hypertarget{p1b:der:g05:s36}{%
\subsubsection{Exchange and the fixed geometry}\label{p1b:der:g05:s36}}

Let the sphere markings be identified and let \[
B(x,p,\theta)=(p,x,S\theta),
\] where \(S\in GL(2,\mathbb Z)\), \(S^2=I\). The full control distribution is preserved by this map. It is an isometry of the displayed control metric when \[
R_x=R_p,\qquad S^{\mathsf T}GS=G.
\] It then commutes with \(L\), preserves the control distance and positive volume measure, exchanges the Reeb fields and the two Heisenberg tangent factors, and carries all factorwise cut and abnormal strata to their exchanged partners. These assertions follow by applying \(B\) to the explicit frame, metric, Hamiltonian, and distance formulas.

For plain sphere exchange on \(P_v\), the fixed diagonal is a copy of \(S^3\), with induced horizontal metric \(2g_H\). Its intrinsic control distance and its ambient restricted distance coincide: \[
d_v((q,q),(q',q'))=\sqrt2\,d_{CC,R}(q,q').
\] The lower bound is the product formula and the upper bound is attained by taking the same minimizing curve in both factors. Its intrinsic growth and Hausdorff dimension are \((2,3)\) and four.

For scalar coordinate swap \(S(\theta_1,\theta_2)=(\theta_2,\theta_1)\), the fixed set in \(M_8\) is \[
\operatorname{Fix}B=\{(q,q,s,s)\}\cong S^3\times S^1.
\] An exchange-invariant torus metric has the form \[
G=\begin{pmatrix}a&b\\b&a\end{pmatrix},\qquad a>|b|.
\] The induced distribution and its metric on this fixed set are \[
\mathcal D_{\rm fix}\cong H\oplus TS^1,\qquad
g_{\rm fix}=2g_H\oplus2(a+b)\,ds^2.
\] It follows by the same bracket computation that \[
\boxed{
\text{intrinsic fixed-set growth }(3,4),\quad
\operatorname{Tan}_{\rm intrinsic}
=\mathbb H^1\times\mathbb R,\quad
Q_{\rm intrinsic}=5.
}
\] The word intrinsic specifies paths constrained to the fixed set. Its circle can have a longer constrained distance than its distance through the ambient torus; Paper~I-D, Section~\ref{p1b:app:g05:2} gives an exact example.

On \(P_v\), swapping the two three-dimensional sphere tangent blocks reverses orientation. On \(M_8\), combining that swap with scalar coordinate swap reverses orientation twice and hence preserves it. Equivalently, the full tangent involution has four positive and four negative eigenvectors; the horizontal subspace has three of each and the Reeb quotient one of each. Positive volume densities are preserved in either case, independently of oriented-form signs.

Paper~I-D, Section~\ref{p1b:app:g05} proves the convention-comparison identities and extended-case formulas used in the results above.


\hypertarget{p1b:sec:g06}{%
\subsection{CR complexes, Hardy spaces and bi-Hopf quantization}\label{p1b:sec:g06}}


\begin{proposition}[Graded CR products and integer bi-Hopf twists]
\label{p1b:prop:g06-product}
Use the specified common-phase Hopf markings and compatible CR
structures on the two spheres. On the graded tensor complex,
\[
 \bar\partial_b
 =\bar\partial_{b,x}\otimes I
   +(-1)^{N_x}\otimes\bar\partial_{b,p},
 \qquad
 \Box_b=\Box_{b,x}\otimes I+I\otimes\Box_{b,p}.
\]
On the quotient \(\mathbb{CP}^1_x\times\mathbb{CP}^1_p\), every
integer pair \((q,r)\) defines the line
\(\mathcal O(q)\boxtimes\mathcal O(r)\). Its Dolbeault cohomology is
the graded product of the one-factor groups
\[
 h^0(\mathcal O(q))=\max(q+1,0),\qquad
 h^1(\mathcal O(q))=\max(-q-1,0).
\]
The charge of the degree-one CR coefficient is shifted by minus two.
Including a scalar elliptic curve requires a chosen compatible complex
structure on the torus. Exchange is a chain symmetry only when it
respects these complex data and the graded signs.
The explicitly defined product Paneitz operators below retain their
definitions; the one-factor \(P'\), \(Q'\) and rigidity theorems remain
one-factor references.
\end{proposition}

The local CR fields, adjoints, graded differential and harmonic representatives below prove these formulas and determine the Hardy, Szeg\H{o} and quantization constructions.

\hypertarget{p1b:der:g06:s38}{%
\subsubsection{Normalization}\label{p1b:der:g06:s38}}

This derivation supplies the CR, basic Dolbeault, Hardy/Szegő, Paneitz, and bi-Hopf quantization package for the marked orthogonal product \[
P_v=S^3(R_x)\times S^3(R_p),\qquad
M_8=P_v\times T^2_G,\qquad R_x,R_p>0,\quad G>0.
\] The scalar angles have period \(2\pi\). The Hopf circles on the spheres are distinct from these scalar circles.

The principal quotient is \[
B=P_v/(S^1_{H,x}\times S^1_{H,p})
 =\mathbb {CP}^1_x\times\mathbb {CP}^1_p,
\qquad
M_8/(S^1_{H,x}\times S^1_{H,p})=B\times T^2_G.
\] Each quotient sphere has radius \(R_i/2\). When the scalar torus is included in a Dolbeault complex, choose its orientation and the compatible flat complex structure \(J_T\); call the resulting elliptic curve \(E\). This choice is explicitly tracked below.

The controlling one-factor source inputs are Paper I, the section labeled p1:subsec:cr-structure, p1:thm:hopf-sasakian-dolbeault-hodge, p1:thm:hopf-basic-dolbeault-index, p1:thm:hardy-szego-projection, p1:subsec:cr-paneitz, p1:thm:borel-weil-cp1, and p1:cor:cp1-geometric-quantization. Paper I-C, p1b:thm:charge-bundles supplies the associated-bundle setting; p1b:thm:complex-family classifies a different object, an integrable complex structure on all of \(TM_8\).

We retain G05's common-phase Hopf convention. On one factor use unit complex coordinates \(w=(w_1,w_2)\), \(|w|=1\), and the round metric \(R^2g_{\rm unit}\). Its real orthonormal frame obeys \[
[X,Y]=-\frac2R Z,\qquad [Z,X]=-\frac2R Y,\qquad
[Z,Y]=\frac2R X,\qquad JX=Y.
\] Here \[
\alpha=R\eta_0,\quad X_H=RZ,\quad
Q_H=-iX_H,\quad Z=\frac{i}{R}Q_H,\quad
L_H=-(X^2+Y^2)=\Delta_{S^3}-\frac{Q_H^2}{R^2}.
\] \(Q_H\) denotes integer geometric charge, not the nilpotent Dolbeault differential. Normalized round measure is used for scalar inner products; constant rescaling of that measure does not change any displayed adjoint.

For the pseudohermitian formulas it is convenient to make the horizontal Levi metric equal to this round horizontal metric: \[
\vartheta=\frac{R^2}{2}\eta_0=\frac R2\alpha,\qquad
T_\vartheta=\frac2R Z,\qquad
Z_1=\frac{X-iY}{\sqrt2},\quad
Z_{\bar1}=\frac{X+iY}{\sqrt2}.
\] The complex scalar \(i\) here denotes complexification. With \(\theta^1=(X^*+iY^*)/\sqrt2\) and \(\theta^{\bar1}=\overline{\theta^1}\), \[
d\vartheta=i\theta^1\wedge\theta^{\bar1},\qquad
\omega_1{}^1=-\frac{4i}{R^2}\vartheta,\qquad
d\theta^1=\theta^1\wedge\omega_1{}^1,\qquad
d\omega_1{}^1=\frac4{R^2}\theta^1\wedge\theta^{\bar1}.
\] These identities follow by evaluating both sides on \(X,Y,Z\). Thus the Tanaka--Webster torsion vanishes and, in the convention \(d\omega_1{}^1=\mathrm{Scal}_\vartheta \theta^1\wedge\theta^{\bar1}\), \(\mathrm{Scal}_\vartheta=4/R^2\). For the dimensionless form \(\eta_0\) the scalar is \(2\); for G05's form \(\alpha\) it is \(2/R\). The different powers reflect the stated contact-form scaling. They do not replace the round Riemannian scalar \(6/R^2\).

The structure-equation and Paneitz conventions used here are those in equations (2.5)--(2.13) and (3.4)--(3.5) of \href{https://arxiv.org/html/2506.19215v1}{Takeuchi, CR Paneitz operator and embeddability}. All product claims below are derived for the displayed operators.

\hypertarget{p1b:der:g06:s39}{%
\subsubsection{One-factor Kohn complex and the charge shift}\label{p1b:der:g06:s39}}

In the chosen coordinates the tangent CR fields are \[
Z_1=\frac{\sqrt2}{R}
 \left(-\bar w_2\partial_{w_1}+\bar w_1\partial_{w_2}\right),
\qquad
Z_{\bar1}=\frac{\sqrt2}{R}
 \left(-w_2\partial_{\bar w_1}+w_1\partial_{\bar w_2}\right).
\] They annihilate \(|w|^2\), so these are intrinsic tangential operators. The complex is \[
C^\infty(S^3)\xrightarrow{\ \bar\partial_b\ }
C^\infty(S^3)\theta^{\bar1}\longrightarrow0,
\qquad
\bar\partial_bf=(Z_{\bar1}f)\theta^{\bar1}.
\] Divergence-free real invariant fields give \[
\bar\partial_b^\dagger(g\theta^{\bar1})=-Z_1g.
\] Set \[
A=\Box_b^{(0)}=-Z_1Z_{\bar1}
 =\frac12L_H-\frac{Q_H}{R^2},
\qquad
B_c=-Z_{\bar1}Z_1
 =\frac12L_H+\frac{Q_H}{R^2}.
\tag{G06.1}\label{p1b:der:g06:eq-G06-1}
\] The subscript in \(B_c\) distinguishes the conjugate Kohn operator from the quotient manifold \(B\). The degree-one Laplacian is \[
\Box_b^{(1)}(g\theta^{\bar1})=(B_cg)\theta^{\bar1}.
\] Both operators are nonnegative, commute, and sum to \(L_H\).

Let \(\mathcal H^{r,s}\) be the restrictions of complex harmonic polynomials bihomogeneous of degrees \(r,s\) in \(w,\bar w\), for integers \(r,s\ge0\). The complex Euclidean Laplacian maps the \((r,s)\) polynomial space onto the \((r-1,s-1)\) space. For completeness, with \(t=|w|^2\) and bihomogeneous \(f\) of degrees \(a,b\), \[
\sum_j\partial_{w_j}\partial_{\bar w_j}(tf)
 =(a+b+2)f+t\sum_j\partial_{w_j}\partial_{\bar w_j}f.
\] Successive subtraction of radial multiples terminates and proves surjectivity and the Fischer decomposition. Consequently \[
\dim_{\mathbb C}\mathcal H^{r,s}
 =(r+1)(s+1)-rs=r+s+1.
\] The scalar spherical degree and charge are \[
k=r+s,\qquad q=r-s.
\] The ordinary spherical harmonic decomposition, refined by the commuting circle action, gives the complete orthogonal sum \[
L^2(S^3)=\widehat{\bigoplus}_{r,s\ge0}\mathcal H^{r,s}.
\] Substituting \(k,q\) in (G06.1), or directly differentiating the polynomials, yields \[
\boxed{\quad
A|_{\mathcal H^{r,s}}=\frac{2s(r+1)}{R^2},\qquad
B_c|_{\mathcal H^{r,s}}=\frac{2r(s+1)}{R^2}.
\quad}
\tag{G06.2}\label{p1b:der:g06:eq-G06-2}
\] In particular the first positive eigenvalue of either operator is \(2/R^2\). Its zero eigenspace is infinite dimensional.

There is an essential form-charge correction: \[
[X_H,Z_{\bar1}]=2iZ_{\bar1},\qquad
\mathcal L_{X_H}\theta^{\bar1}=-2i\theta^{\bar1}.
\tag{G06.3}\label{p1b:der:g06:eq-G06-3}
\] Thus \(g\theta^{\bar1}\) has total charge \(q_g-2\). The differential raises the coefficient charge by two and preserves the total charge: \[
\bar\partial_b:\mathcal H^{r,s}
\longrightarrow\mathcal H^{r+1,s-1}\theta^{\bar1}
\quad(s>0).
\] The map is an isomorphism onto this partner space because the two spaces have dimension \(r+s+1\), and \[
\|\bar\partial_bf\|^2
=\frac{2s(r+1)}{R^2}\|f\|^2.
\] Its adjoint gives the inverse up to this eigenvalue.

Writing \[
\mathscr H^+=\widehat{\bigoplus}_{r\ge0}\mathcal H^{r,0},
\qquad
\mathscr H^-=\widehat{\bigoplus}_{s\ge0}\mathcal H^{0,s},
\] the harmonic cohomology is therefore \[
H_b^{0,0}=\mathscr H^+,\qquad
H_b^{0,1}=\mathscr H^-\theta^{\bar1}.
\tag{G06.4}\label{p1b:der:g06:eq-G06-4}
\] These are Hilbert cohomology spaces; their smooth parts have the usual rapid spherical-harmonic coefficient decay. The positive gap and the explicit orthogonal paired decomposition prove closed range. No finite Dolbeault index is assigned to this unrestricted CR complex.

\hypertarget{p1b:der:g06:s40}{%
\subsubsection{Product CR complexes and Hodge operators}\label{p1b:der:g06:s40}}

On \(P_v\) take \[
T^{1,0}_{CR}P_v
=T^{1,0}S_x^3\oplus T^{1,0}S_p^3.
\] Cross-factor brackets vanish; the one-factor complex lines are involutive. This proves formal integrability and gives CR dimension two and real CR codimension two. For a conormal covector \(a\alpha_x+b\alpha_p\), its Levi form in the marked unit CR frames is \[
\operatorname{diag}\left(\frac{2a}{R_x},\frac{2b}{R_p}\right).
\] It is positive definite when \(a,b>0\). This is a CR structure of codimension two, so a theorem stated for strictly pseudoconvex real hypersurfaces cannot be applied by merely changing the dimension label.

Use the graded tensor differential \[
d_{CR}=\bar\partial_{b,x}\otimes I+
(-1)^{N_x}\otimes\bar\partial_{b,p}.
\tag{G06.5}\label{p1b:der:g06:eq-G06-5}
\] On a homogeneous tensor the sign is \((-1)^{\deg_x}\). Its mixed terms cancel, proving \[
d_{CR}^2=0,\qquad
\Box_{CR}=\Box_{b,x}\otimes I+I\otimes\Box_{b,p}.
\tag{G06.6}\label{p1b:der:g06:eq-G06-6}
\] For \(\mathcal Q=\sqrt2\,d_{CR}\), \[
\mathcal Q^2=(\mathcal Q^\dagger)^2=0,\qquad
\{\mathcal Q,\mathcal Q^\dagger\}=2\Box_{CR},\qquad
[2\Box_{CR},\mathcal Q]=0.
\] These identities hold on smooth forms and extend to the closed Hilbert complexes with their spectral domains.

The explicit one-factor paired decomposition proves the product cohomology without a formal dimension argument: \[
\begin{aligned}
H_{CR}^{0,0}(P_v)
 &=\mathscr H_x^+\widehat\otimes\mathscr H_p^+,\\
H_{CR}^{0,1}(P_v)
 &=(\mathscr H_x^-\theta_x^{\bar1}
        \widehat\otimes\mathscr H_p^+)
 \oplus(\mathscr H_x^+\widehat\otimes
        \mathscr H_p^-\theta_p^{\bar1}),\\
H_{CR}^{0,2}(P_v)
 &=\mathscr H_x^-\theta_x^{\bar1}
       \widehat\otimes\mathscr H_p^-\theta_p^{\bar1}.
\end{aligned}
\tag{G06.7}\label{p1b:der:g06:eq-G06-7}
\] Indeed every nonzero factor eigenvalue remains positive in the sum, while simultaneous factor zero modes give precisely these tensors. The least positive eigenvalue of \(2\Box_{CR}\) is \(\min(4/R_x^2,4/R_p^2)\). Multiplicities in this unrestricted product can be infinite because one factor can remain in its infinite harmonic space.

There are two useful extensions to \(M_8\), and they have different differentials.

\begin{longtable}[]{@{}
  >{\raggedright\arraybackslash}p{(\columnwidth - 4\tabcolsep) * \real{0.3333}}
  >{\raggedright\arraybackslash}p{(\columnwidth - 4\tabcolsep) * \real{0.3333}}
  >{\raggedright\arraybackslash}p{(\columnwidth - 4\tabcolsep) * \real{0.3333}}@{}}
\toprule\noalign{}
\begin{minipage}[b]{\linewidth}\raggedright
Complex on \(M_8\)
\end{minipage} & \begin{minipage}[b]{\linewidth}\raggedright
CR dimension / codimension
\end{minipage} & \begin{minipage}[b]{\linewidth}\raggedright
Differential and scalar kernel
\end{minipage} \\
\midrule\noalign{}
\endhead
\bottomrule\noalign{}
\endlastfoot
Lift of the two sphere CR structures; scalar torus as parameter & \(2/4\) & \(d_{CR}\otimes I\); kernel \(\mathscr H_x^+\widehat\otimes\mathscr H_p^+\widehat\otimes L^2(T^2)\) \\
Product with the elliptic curve \(E\) & \(3/2\) & Graded sum of \(d_{CR}\) and \(\bar\partial_E\); kernel \(\mathscr H_x^+\widehat\otimes\mathscr H_p^+\otimes\mathbb C\) \\
\end{longtable}

For the second row the underlying real distribution is exactly G05's \(\mathcal D=\mathcal H_x\oplus\mathcal H_p\oplus TT^2\). Its conormal Levi form has the additional zero elliptic direction. Since holomorphic functions on a compact connected elliptic curve are constant, the stated scalar kernel follows. On all degrees its CR cohomology is (G06.7) tensored with \(H^{0,*}(E)\); this remains infinite dimensional.

These CR complexes and their quotients do not assert a Dolbeault-cohomology formula for any full complex structure on \(TM_8\). Their relation to Paper I-C's complex family is treated in Section~\ref{p1b:der:g06:s46}.
The full complex-fourfold Dolbeault algebra and its two parameter
strata are computed separately in Theorems~\ref{p1b:thm:g16-model}
and~\ref{p1b:thm:g16-hodge}, with the Frölicher differentials in
Proposition~\ref{p1b:prop:g16-frolicher}.

\hypertarget{p1b:der:g06:s41}{%
\subsubsection{The basic quotient and the torus contribution}\label{p1b:der:g06:s41}}

A form is bi-Hopf basic when contraction and Lie derivative with respect to each \(X_{H,i}\) vanish. The free circle actions and their horizontal connections identify basic forms with quotient forms: \[
\Omega_B^{p,q}(P_v)\cong\pi^*\Omega^{p,q}(B),
\qquad
\Omega_B^{p,q}(M_8)\cong\pi^*\Omega^{p,q}(B\times E).
\tag{G06.8}\label{p1b:der:g06:eq-G06-8}
\] The second identity uses the chosen \(J_T\). One can prove the identification in a principal-bundle trivialization: horizontality removes fiber differentials and invariance makes the remaining coefficients descend. Conversely pullbacks satisfy both conditions. The complex structures, metrics, and differentials descend with them. Constant fiber volumes make the identification unitary after a constant normalization.

On \(B\), the graded Dolbeault differential and its adjoint give \[
H_B=2\Box_B,\qquad
H_B|_{\Omega^{0,0}(B)}=\Delta_B.
\] Equation (G06.3) must be used for basic forms: zero coefficient charge on a one-form is not the basic condition. At total charge zero the one-factor harmonic calculation gives \(H^{0,0}(\mathbb {CP}^1)=\mathbb C\) and \(H^{0,1}(\mathbb {CP}^1)=0\). The tensor complex therefore yields \[
\boxed{\quad
h^{0,q}(B)=(1,0,0),\qquad \chi(B,\mathcal O)=1.
\quad}
\tag{G06.9}\label{p1b:der:g06:eq-G06-9}
\] The scalar eigenvalues and multiplicities are \[
\lambda_{\ell_x,\ell_p}
=\frac{4\ell_x(\ell_x+1)}{R_x^2}
 +\frac{4\ell_p(\ell_p+1)}{R_p^2},
\qquad
\mathrm{mult}=(2\ell_x+1)(2\ell_p+1),
\quad \ell_x,\ell_p\ge0.
\tag{G06.10}\label{p1b:der:g06:eq-G06-10}
\] These labels are base spherical degrees; their basic lifts have three-sphere degrees \(2\ell_i\). The first positive eigenvalue is \(\min(8/R_x^2,8/R_p^2)\); multiplicities add if the two minima coincide.

For completeness the basic de Rham ring and Hodge polynomial are \[
H^*_{\rm basic}(P_v;\mathbb R)
=\mathbb R[h_x,h_p]/(h_x^2,h_p^2),\quad |h_i|=2,
\qquad
\sum h^{p,q}(B)u^pv^q=(1+uv)^2.
\tag{G06.11}\label{p1b:der:g06:eq-G06-11}
\] The normalized curvature forms represent the generators. This uses the ordinary two-sphere cohomology and its Kähler Hodge types, followed by the product complex.

Write \(G=\begin{pmatrix}a&b\\b&c\end{pmatrix}\) and \(\delta=\sqrt{ac-b^2}>0\). In the chosen orientation, \[
J_T=\frac1\delta
\begin{pmatrix}-b&-c\\a&b\end{pmatrix},
\quad J_T^2=-I,\quad J_T^{\mathsf T}GJ_T=G,
\quad
z_T=\theta_1+\tau\theta_2,\quad
\tau=\frac{b+i\delta}{a}.
\tag{G06.12}\label{p1b:der:g06:eq-G06-12}
\] Indeed \(dz_T\circ J_T=i\,dz_T\). Thus \(E=\mathbb C/[2\pi(\mathbb Z+\tau\mathbb Z)]\). Fourier modes show that \[
2\Box_E=\Delta_T,\qquad
\lambda_n=n^{\mathsf T}G^{-1}n,\quad n\in\mathbb Z^2,
\] on functions and in a parallel unit \((0,1)\) coframe. Only \(n=0\) is harmonic, once in degree zero and once in degree one. Hence \[
h^{0,q}(E)=(1,1),\qquad \chi(E,\mathcal O)=0.
\] Tensoring gives the full bi-Hopf basic result \[
\boxed{\quad
h_B^{0,q}(M_8)=(1,1,0,0),\qquad
\chi_B(M_8,\mathcal O)=0.
\quad}
\tag{G06.13}\label{p1b:der:g06:eq-G06-13}
\] There are two harmonic states in the \((0,*)\) complex: the scalar constant and the constant torus \((0,1)\)-form. The Hodge polynomial is \[
(1+uv)^2(1+u)(1+v).
\] Equivalently:

\begin{longtable}[]{@{}lrrrr@{}}
\toprule\noalign{}
\(p\backslash q\) & 0 & 1 & 2 & 3 \\
\midrule\noalign{}
\endhead
\bottomrule\noalign{}
\endlastfoot
0 & 1 & 1 & 0 & 0 \\
1 & 1 & 3 & 2 & 0 \\
2 & 0 & 2 & 3 & 1 \\
3 & 0 & 0 & 1 & 1 \\
\end{longtable}

The basic scalar spectrum is (G06.10) plus \(\lambda_n\). Writing \(\lambda_{T,1}=\min_{n\ne0}n^{\mathsf T}G^{-1}n\), its sharp gap is \[
\min\left(\frac8{R_x^2},\frac8{R_p^2},\lambda_{T,1}\right).
\tag{G06.14}\label{p1b:der:g06:eq-G06-14}
\] The Dolbeault index of this quotient, the finite scalar kernel, and the infinite CR harmonic spaces in Section~\ref{p1b:der:g06:s40} are distinct results for distinct complexes.

\hypertarget{p1b:der:g06:s42}{%
\subsubsection{All integer twists and the bi-Hopf index}\label{p1b:der:g06:s42}}

For each factor define \[
\omega_i=\frac2{R_i^2}\,\mathrm{vol}_{S^2(R_i/2)},
\qquad d\eta_{0,i}=\pi_i^*\omega_i,\qquad
\int_{\mathbb {CP}^1_i}\omega_i=2\pi.
\] Use the positive associated-line convention of Paper I: charge-\(q\) functions obey \[
f(e^{i\phi}w)=e^{iq\phi}f(w)
\] and represent sections of \(\mathcal O(q)\). The induced unitary curvature is \[
F_{\nabla_q}=-iq\,\omega,\qquad
c_1(\mathcal O(q))=\left[\frac{iF_{\nabla_q}}{2\pi}\right]
 =q\left[\frac{\omega}{2\pi}\right].
\tag{G06.15}\label{p1b:der:g06:eq-G06-15}
\] The invariant horizontal \(\bar\partial_b\) is precisely the pullback of this Chern connection's Dolbeault operator. In a local section of the Hopf bundle this follows by subtracting the vertical connection term; its \((0,1)\) part is the associated-line \(\bar\partial\).

Equation (G06.2) gives the complete one-factor harmonic result without invoking an index to guess dimensions: \[
h^0(q)=\max(q+1,0),\qquad
h^1(q)=\max(-q-1,0).
\tag{G06.16}\label{p1b:der:g06:eq-G06-16}
\] To prove it, scalar zero modes require \(s=0\) and \(r=q\ge0\). Degree-one zero modes have coefficient \(r=0\), and their total charge is \(-s-2=q\), requiring \(q\le-2\) and giving dimension \(-q-1\). At \(q=-1\) neither case exists. In particular negative-charge harmonic one-forms are not obtained by assigning the same charge to an antiholomorphic scalar coefficient.

There is also a compact complete spectral formula. For \(H_q=2\Box_{\mathcal O(q)}\), the zero eigenspace has dimension \(|q+1|\) if \(q\ne-1\), in degree zero for \(q\ge0\) and degree one for \(q\le-2\). All positive levels are \[
\boxed{\quad
\nu_{q,n}=\frac{4n(n+|q+1|)}{R^2},\qquad
n=1,2,\ldots,\qquad
\mathrm{mult}_{(0,0)}=\mathrm{mult}_{(0,1)}
=2n+|q+1|.
\quad}
\tag{G06.17}\label{p1b:der:g06:eq-G06-17}
\] For \(q\ge0\), put \(s=n,\ r=n+q\) in the scalar sector. For \(q\le-2\), put \(r=n-1,\ s=n-q-1\). For \(q=-1\), put \(r=n-1,s=n\). The one-form partner from Section~\ref{p1b:der:g06:s39} has the same eigenvalue and dimension in every case. These exhaust the charge sector. This proves both the spectrum and the differential pairing. At \(q=0\) the first level is \(8/R^2\) with multiplicity three in each degree; at \(q=-1\) it is \(4/R^2\) with multiplicity two in each degree.

For integers \(a,b\), let \[
L_{a,b}=\mathcal O(a)\boxtimes\mathcal O(b)\longrightarrow B,
\qquad
F_{a,b}=-i(a\omega_x+b\omega_p).
\] Tensor the two paired complexes. Their only cohomology comes from harmonic tensors, so \[
\begin{aligned}
h^0(a,b)&=h^0(a)h^0(b),\\
h^1(a,b)&=h^1(a)h^0(b)+h^0(a)h^1(b),\\
h^2(a,b)&=h^1(a)h^1(b),\\
\chi(B,L_{a,b})&=(a+1)(b+1).
\end{aligned}
\tag{G06.18}\label{p1b:der:g06:eq-G06-18}
\] If either charge equals \(-1\), the complex is acyclic. Otherwise the cohomology is concentrated in degree equal to the number of charges at most \(-2\), with dimension \(|a+1||b+1|\). All positive eigenvalues and degree multiplicities follow by adding the two spectra (G06.17), including their zero levels, and combining coincident sums.

On \(B\times E\), with no additional twist on \(E\), the harmonic degree polynomial is multiplied by \(1+t\). Thus the total harmonic dimension is \(2|a+1||b+1|\) unless a charge is \(-1\), in which case it is zero; the Dolbeault index is always zero. Torus Fourier restriction preserves these zero modes exactly when the retained set contains \(n=0\).

For \(a,b\ge0\) the holomorphic quantization space is \[
\boxed{\quad
\mathscr Q_{a,b}
=H^0(B,L_{a,b})
\cong\operatorname{Sym}^a(\mathbb C^2)^*
 \otimes\operatorname{Sym}^b(\mathbb C^2)^*,
\qquad
\dim\mathscr Q_{a,b}=(a+1)(b+1).
\quad}
\tag{G06.19}\label{p1b:der:g06:eq-G06-19}
\] Homogeneous polynomial sections prove this realization directly. The unitary polynomial action is the source's dual convention; the two spins are \(a/2,b/2\). For \(a,b>0\), the form \(a\omega_x+b\omega_p\) is symplectic and this is its ordinary positive holomorphic quantization. At a zero charge the corresponding line is trivial and the holomorphic space remains well defined; the two-factor curvature form is then degenerate on that factor, so no nondegenerate four-dimensional symplectic form is claimed there.

\hypertarget{p1b:der:g06:s43}{%
\subsubsection{Hardy spaces, Szegő kernels, and finite restrictions}\label{p1b:der:g06:s43}}

Take normalized Haar probability measure \(d\mu\) on each unit coordinate sphere. Phase integration and the uniform distribution of \(|w_1|^2\) on \([0,1]\) give \[
\int w^\alpha\overline{w^\beta}\,d\mu
=\delta_{\alpha\beta}\frac{\alpha!}{(|\alpha|+1)!}.
\tag{G06.20}\label{p1b:der:g06:eq-G06-20}
\] Indeed the phase integral vanishes unless the two exponent vectors agree; in that case the remaining integral is \(\int_0^1t^{\alpha_1}(1-t)^{\alpha_2}dt =\alpha_1!\alpha_2!/(|\alpha|+1)!\). For \(|\alpha|=n\), an orthonormal holomorphic basis is therefore \[
e_\alpha(w)=\sqrt{\frac{(n+1)!}{\alpha!}}\,w^\alpha.
\] The binomial theorem supplies the degree-\(n\) projector kernel: \[
K_n(w,v)=\sum_{|\alpha|=n}e_\alpha(w)\overline{e_\alpha(v)}
 =(n+1)\langle w,v\rangle^n,
\quad \langle w,v\rangle=\sum_jw_j\bar v_j.
\tag{G06.21}\label{p1b:der:g06:eq-G06-21}
\] Thus \(\mathscr H^+\) in Section~\ref{p1b:der:g06:s39} is exactly the Hardy space of the complex two-ball: a holomorphic polynomial has only \(s=0\) components, and the Hilbert closure of these components is \(\ker\bar\partial_b\). Conversely a square-summable sequence of such homogeneous components defines a holomorphic function in the ball, with bounded radial \(L^2\) norms and the stated boundary limit.

If \(\Pi_n\) denotes the orthogonal degree-\(n\) projection, the Szegő projector and its Abel kernels are \[
\Pi_{\rm Sz}=\sum_{n\ge0}\Pi_n,\qquad
\sum_{n\ge0}\rho^nK_n(w,v)
=\frac1{(1-\rho\langle w,v\rangle)^2},
\quad 0\le\rho<1.
\tag{G06.22}\label{p1b:der:g06:eq-G06-22}
\] The operators on the left converge strongly to \(\Pi_{\rm Sz}\) as \(\rho\uparrow1\), by dominated convergence of the squared Fourier coefficients. At \(\rho=1\) the displayed expression is understood as a boundary distribution or this Abel limit. It is not an ordinary bounded kernel at the diagonal.

On \(P_v\), \[
\mathscr H^+_{xp}
=\mathscr H_x^+\widehat\otimes\mathscr H_p^+,\qquad
\Pi_{{\rm Sz},xp}=\Pi_{{\rm Sz},x}\otimes\Pi_{{\rm Sz},p}.
\tag{G06.23}\label{p1b:der:g06:eq-G06-23}
\] The Abel kernel is the product \[
\frac1{(1-\rho_x\langle w_x,v_x\rangle)^2
          (1-\rho_p\langle w_p,v_p\rangle)^2}.
\] Separate radial limits and the orthogonal coefficient sum prove the tensor identification. Geometrically this is the Hardy space on the distinguished boundary \(S^3\times S^3\) of \(B^2\times B^2\). The full topological boundary of that product domain has additional faces; \(S^3\times S^3\) is not the boundary sphere of a ball in \(\mathbb C^4\).

The fixed-charge \((a,b)\) block of (G06.23), for \(a,b\ge0\), is exactly \(\mathscr Q_{a,b}\), with kernel \[
K_{a,b}(w,v)
=(a+1)(b+1)
\langle w_x,v_x\rangle^a\langle w_p,v_p\rangle^b.
\tag{G06.24}\label{p1b:der:g06:eq-G06-24}
\] This proves agreement between homogeneous boundary quantization and the positive associated-line Dolbeault construction.

Let \(\mathcal P_\ell\) be the scalar spherical band of degrees \(0,\ldots,\ell\), using G03's notation. Its dimension and its Hardy rank are \[
d_\ell=\frac{(\ell+1)(\ell+2)(2\ell+3)}6,\qquad
b_\ell=\frac{(\ell+1)(\ell+2)}2.
\tag{G06.25}\label{p1b:der:g06:eq-G06-25}
\] Both the band projector and the Szegő projector are diagonal in \(\mathcal H^{r,s}\), so they commute. The Hardy restriction has rank \(b_\ell\) because it retains just \(\mathcal H^{n,0}\), \(0\le n\le\ell\). Consequently independent scalar cutoffs have Hardy rank \(b_{\ell_x}b_{\ell_p}\). For the different restriction \(a+b\le L\), the rank is \[
\sum_{a+b\le L}(a+1)(b+1)=\binom{L+4}{4}.
\tag{G06.26}\label{p1b:der:g06:eq-G06-26}
\] The shell \(a+b=N\) has dimension \(\binom{N+3}{3}\), as follows either by convolution or by counting four-variable monomials.

For a finite scalar-torus Fourier set \(F\), the sphere-only CR extension has finite Hardy rank \(b_{\ell_x}b_{\ell_p}|F|\). The elliptic CR extension has rank \(b_{\ell_x}b_{\ell_p}\) if \(0\in F\), and zero otherwise. These are scalar ranks. The enriched carrier \(H_m\), its bundles, and the protected \(K_{16}\) from G02--G04 are different spaces. In particular \(\dim\mathscr Q_{1,1}=4\) and \(\dim(\mathcal P_1\otimes\mathcal P_1)=25\); neither is 16 or the full enriched dimension \(15^2\).

\hypertarget{p1b:der:g06:s44}{%
\subsubsection{Factor Paneitz operators and the mixed terms}\label{p1b:der:g06:s44}}

For the contact form \(\vartheta_i\) in Section~\ref{p1b:der:g06:s38} the torsion vanishes. The three-dimensional Paneitz formula therefore reduces to the commuting Kohn product. Fix its normalization by \[
\mathsf P_i=A_iB_{c,i}
=\frac14\left(L_{H,i}^2-\frac{4Q_{H,i}^2}{R_i^4}\right).
\tag{G06.27}\label{p1b:der:g06:eq-G06-27}
\] This is the torsion-free specialization of the Paneitz definition in \href{https://arxiv.org/html/2506.19215v1}{Takeuchi, equations (3.4)--(3.5)}. For a convention that doubles each Kohn operator the corresponding factor \(1/4\) must also be included. Under the constant change \(\alpha_i=(2/R_i)\vartheta_i\), the Paneitz operator for \(\alpha_i\) is \(R_i^2\mathsf P_i/4\); for \(\eta_{0,i}\) it is \(R_i^4\mathsf P_i/4\). All product formulas below use (G06.27), whose metric units are length to the power \(-4\).

Equations (G06.2) give \[
\mathsf P_i|_{\mathcal H^{r,s}}
=\frac{4rs(r+1)(s+1)}{R_i^4}.
\tag{G06.28}\label{p1b:der:g06:eq-G06-28}
\] This is nonnegative, and it vanishes exactly when \(r=0\) or \(s=0\). Hence \[
\ker\mathsf P_i
=\mathscr P_i^{\mathbb C}
:=\mathscr H_i^++\mathscr H_i^-,
\qquad
\mathscr H_i^+\cap\mathscr H_i^-=\mathbb C.
\tag{G06.29}\label{p1b:der:g06:eq-G06-29}
\] For smooth real functions the kernel is the real part of smooth CR-holomorphic functions, as the conjugate harmonic coefficients pair uniquely apart from the constant. This directly proves the standard-sphere positivity and kernel statement in the source. The first positive eigenvalue is \(16/R_i^4\), at \(r=s=1\).

Several different product constructions are possible. On scalar functions on \(P_v\) define \[
\mathsf P_\Sigma=\mathsf P_x\otimes I+I\otimes\mathsf P_p,
\qquad
\mathsf P_\otimes=\mathsf P_x\otimes\mathsf P_p,
\tag{G06.30}\label{p1b:der:g06:eq-G06-30}
\] and \[
\begin{aligned}
\mathsf P_J
 &=(A_x+A_p)(B_{c,x}+B_{c,p})\\
 &=\mathsf P_\Sigma+A_xB_{c,p}+A_pB_{c,x}.
\end{aligned}
\tag{G06.31}\label{p1b:der:g06:eq-G06-31}
\] The tensor identities are implicit in the second display. The sum and joint operators are fourth order; the tensor operator is eighth order. They are all nonnegative and self-adjoint on their spectral domains, since all four factor Kohn operators strongly commute on the complete product harmonic basis.

Their kernels follow from products and sums of nonnegative eigenvalues: \[
\begin{aligned}
\ker\mathsf P_\Sigma
 &=\mathscr P_x^{\mathbb C}\widehat\otimes
   \mathscr P_p^{\mathbb C},\\
\ker\mathsf P_\otimes
 &=(\mathscr P_x^{\mathbb C}\widehat\otimes L^2(S_p^3))
  +(L^2(S_x^3)\widehat\otimes\mathscr P_p^{\mathbb C}),\\
\ker\mathsf P_J
 &=(\mathscr H_x^+\widehat\otimes\mathscr H_p^+)
   +(\mathscr H_x^-\widehat\otimes\mathscr H_p^-).
\end{aligned}
\tag{G06.32}\label{p1b:der:g06:eq-G06-32}
\] For the last identity, a joint eigenvalue \((a_x+a_p)(b_x+b_p)\) vanishes precisely when \(a_x=a_p=0\) or \(b_x=b_p=0\). The two summands intersect in the constants and are otherwise orthogonal. Thus the real smooth joint kernel is exactly the real parts of globally CR-holomorphic functions for the product CR structure. Here that description defines the global product CR-pluriharmonic space; no unproved local-to-global theorem for general codimension-two CR manifolds is needed.

The added mixed terms are necessary for this specific joint kernel. For example \[
f=w_{x,1}\bar w_{p,1}
\] lies in \(\ker\mathsf P_\Sigma\), but \[
\mathsf P_Jf=\frac4{R_x^2R_p^2}f.
\tag{G06.33}\label{p1b:der:g06:eq-G06-33}
\] The sum remains a valid factorwise operator; it has a larger kernel than the joint construction. Neither its name nor (G06.31) asserts a universally defined Paneitz invariant on arbitrary CR manifolds of codimension two. The product metric, marked splitting, and normalization are part of their definitions.

At independent scalar cutoffs set \[
d_i=d_{\ell_i},\qquad b_i=b_{\ell_i},\qquad c_i=2b_i-1.
\] Counting the same harmonic axes proves \[
\begin{aligned}
\dim\ker\mathsf P_\Sigma&=c_xc_p,\\
\dim\ker\mathsf P_J&=2b_xb_p-1,\\
\dim\ker\mathsf P_\otimes
 &=d_xd_p-(d_x-c_x)(d_p-c_p),\\
\dim\ker\mathsf P_\Sigma-\dim\ker\mathsf P_J
 &=2(b_x-1)(b_p-1).
\end{aligned}
\tag{G06.34}\label{p1b:der:g06:eq-G06-34}
\] At \(\ell_x=\ell_p=1\), \(d_i=5,\ b_i=3,\ c_i=5\). Therefore \[
\boxed{\quad
\mathsf P_\Sigma=0\ \text{on all 25 scalar modes},\qquad
\operatorname{Spec}\mathsf P_J
=\left\{0^{[17]},
\left(\frac4{R_x^2R_p^2}\right)^{[8]}\right\}.
\quad}
\tag{G06.35}\label{p1b:der:g06:eq-G06-35}
\] Four positive modes are holomorphic-linear in \(x\) and antiholomorphic-linear in \(p\); their four conjugates supply the other orientation. Every other retained sector is joint holomorphic or joint antiholomorphic.

The global positive gaps are \[
\operatorname{gap}\mathsf P_\Sigma
=\min\left(\frac{16}{R_x^4},\frac{16}{R_p^4}\right),
\] \[
\operatorname{gap}\mathsf P_J
=\min\left(\frac{16}{R_x^4},\frac{16}{R_p^4},
           \frac4{R_x^2R_p^2}\right).
\tag{G06.36}\label{p1b:der:g06:eq-G06-36}
\] To prove the second formula, if a factor has \(r,s\ge1\) its own Paneitz term gives the corresponding first bound. Otherwise every factor is on a harmonic axis. A positive joint eigenvalue then requires two opposite nonconstant axes and is at least the mixed value. Each candidate is attained. This also proves a bounded inverse on the orthogonal complement of the respective kernel. The inverse is not defined on a nonzero kernel vector; its Moore--Penrose extension assigns zero there, including the constant.

On \(M_8\), the sphere-only operators can be tensored with the identity on \(L^2(T^2)\). For the elliptic CR structure there is also the explicitly defined joint operator \[
C=\Box_E^{(0)}=\frac12\Delta_T,\qquad
\mathsf P_{J,E}=(A_x+A_p+C)(B_{c,x}+B_{c,p}+C).
\tag{G06.37}\label{p1b:der:g06:eq-G06-37}
\] Its kernel is the last line of (G06.32) tensored with the torus constants: a nonzero Fourier mode has \(C>0\) in both factors. Its gap is \[
\min\left(\operatorname{gap}\mathsf P_J,
          \frac{\lambda_{T,1}^2}{4}\right).
\] For \(n\ne0\) the eigenvalue is at least \(C^2\), attained with both sphere functions constant. This proves the assertion.

The one-factor \(P'\), \(Q'\), and conditional CR sphere-rigidity results in Paper I remain one-factor reference theorems. The construction of \(P'\) on pluriharmonic functions and its associated curvature uses the particular three-dimensional pseudohermitian setting described by \href{https://arxiv.org/abs/1309.2528}{Case--Yang}. No product \(P'\), total \(Q'\), or rigidity conclusion on \(M_8\) is implied by (G06.30)--(G06.37). Keeping those reference statements in Paper I, with their stated hypotheses, is the declared disposition of that part of G06.

\hypertarget{p1b:der:g06:s45}{%
\subsubsection{Finite quantization operators and extremal pairs}\label{p1b:der:g06:s45}}

On \(\mathscr H^+\), multiplication by a holomorphic coordinate preserves the Hardy space. Equation (G06.20) gives the exact weighted shifts \[
M_j e_\alpha
=\sqrt{\frac{\alpha_j+1}{n+2}}\ e_{\alpha+e_j},
\qquad |\alpha|=n,
\tag{G06.38}\label{p1b:der:g06:eq-G06-38}
\] and \[
M_j^\dagger e_\alpha
=\sqrt{\frac{\alpha_j}{n+1}}\ e_{\alpha-e_j},
\] with zero when \(\alpha_j=0\). Consequently, with \(N\) the holomorphic degree operator, \[
\sum_{j=1}^2M_j^\dagger M_j=I,\qquad
\sum_{j=1}^2M_jM_j^\dagger=\frac{N}{N+1}.
\tag{G06.39}\label{p1b:der:g06:eq-G06-39}
\] These identities hold on each homogeneous block and hence on the whole Hardy space.

For bounded scalar \(f\), the degree-\(n\) Toeplitz operator is \[
T_f^{(n)}=\Pi_n M_f\Pi_n.
\] Compression gives \[
T_1^{(n)}=I,\quad
(T_f^{(n)})^\dagger=T_{\bar f}^{(n)},\quad
\|T_f^{(n)}\|\le\|f\|_\infty,\quad f\ge0\Rightarrow T_f^{(n)}\ge0.
\] For a Hopf-invariant \(f\), this is an operator on \(H^0(\mathbb {CP}^1,\mathcal O(n))\). The integral of \(\bar e_\gamma w_i\bar w_j e_\alpha\) vanishes unless the exponents balance. In the surviving case (G06.20) gives, in occupation-number notation, \[
\boxed{\quad
T_{w_i\bar w_j}^{(n)}
=\frac{a_i^\dagger a_j+\delta_{ij}}{n+2}.
\quad}
\tag{G06.40}\label{p1b:der:g06:eq-G06-40}
\] For \(i=j\) the coefficient is \((\alpha_i+1)/(n+2)\); for \(i\ne j\) it is \(\sqrt{\alpha_j(\alpha_i+1)}/(n+2)\), proving the formula. The creation and annihilation operators here are algebraic operators on symmetric powers, with \([a_i,a_j^\dagger]=\delta_{ij}\).

To keep the source's dual convention explicit, define the real coordinate functions \[
\beta_a(w)=\sum_{i,j}(\sigma_a)_{ij}w_i\bar w_j
=w^{\mathsf T}\sigma_a\bar w,\qquad
J_a=\frac12\sum_{i,j}(\sigma_a)_{ij}a_i^\dagger a_j,
\] where \(\sigma_a\) are the ordinary Pauli matrices. Then \[
\sum_a\beta_a^2=1,\qquad
[J_a,J_b]=i\epsilon_{abc}J_c,\qquad
\sum_aJ_a^2=\frac n2\left(\frac n2+1\right)I.
\] It follows that \[
T_{\beta_a}^{(n)}=\frac{2J_a}{n+2},
\quad
[T_{\beta_a}^{(n)},T_{\beta_b}^{(n)}]
=\frac{2i}{n+2}\epsilon_{abc}T_{\beta_c}^{(n)},
\] \[
\boxed{\quad
\sum_a(T_{\beta_a}^{(n)})^2=\frac n{n+2}I,\qquad
\sum_aT_{\beta_a^2}^{(n)}=I.
\quad}
\tag{G06.41}\label{p1b:der:g06:eq-G06-41}
\] The compression defect is \(2I/(n+2)\). These exact identities hold also at \(n=0\).

Paper I writes the ordinary Bloch vector as \(n_a=w^\dagger\sigma_aw\). With these fixed coordinate and matrix conventions, \[
(\beta_1,\beta_2,\beta_3)=(n_1,-n_2,n_3).
\tag{G06.42}\label{p1b:der:g06:eq-G06-42}
\] Thus translating (G06.41) to those Bloch components requires the displayed dual-axis change; it is not a silent orientation-preserving rotation. Poisson-bracket and orientation-sensitive sign comparisons require that translation.

On \(\mathscr Q_{a,b}\), factorized symbols obey \[
T_{f_xf_p}^{(a,b)}=T_{f_x}^{(a)}\otimes T_{f_p}^{(b)},\qquad
[T_{\beta_{x,j}}^{(a)}\otimes I,\,
 I\otimes T_{\beta_{p,k}}^{(b)}]=0.
\tag{G06.43}\label{p1b:der:g06:eq-G06-43}
\] For a positive level, the three \(J_a\) generate the full matrix algebra: \(J_3\) has distinct eigenvalues, its spectral projections are polynomials in \(J_3\), and the nonzero ladder coefficients of \(J_\pm\) join neighboring weight spaces. Products therefore give every matrix unit. For level zero the algebra is \(\mathbb C\). The two factor algebras consequently generate \[
\operatorname{End}(\mathscr Q_{a,b})
\cong M_{(a+1)(b+1)}(\mathbb C).
\] This is a statement about these quantization spaces and compressed symbol operators.

The same realization gives the product of the Borel--Weil extremal-pair construction. For a unit axis \(\mathbf u\in S^2\), \(\mathbf u\cdot\mathbf J\) is unitarily conjugate to \(J_3\), with eigenvalues \(-n/2,-n/2+1,\ldots,n/2\). For \(n\ge1\) its two extremal eigenlines are distinct and orthogonal. If \(v_\pm\) are normalized eigenvectors of \(\mathbf u\cdot\boldsymbol\sigma\), these lines are the symmetric powers of \(v_\pm\) in the stated polynomial identification. The corresponding rank-one projectors have a rank-two sum. For two positive levels the two commuting axis operators give four orthogonal tensor extremal lines, with joint signs \((+,+),(+,-),(-,+),(-,-)\), and a rank-four sum. If one level is zero the number of distinct lines is two, and if both are zero it is one. Exchange swaps the factor axes and the two signs. This derives the product representation geometry associated with Paper I, p1:prop:v191-bw-extremal-pair; the resulting extremal subspace depends on the specified axis.

The reproducing kernels needed for later normalized measures are now available. The actual operator-valued measures and the framework's sampling moments remain the separate G09 item.

\hypertarget{p1b:der:g06:s46}{%
\subsubsection{Exchange, graded signs, and complex compatibility}\label{p1b:der:g06:s46}}

First consider sphere exchange with the scalar torus held fixed. On forms its natural action is the graded flip \[
\mathcal B(\xi\otimes\zeta)
=(-1)^{\deg\xi\,\deg\zeta}\zeta\otimes\xi.
\tag{G06.44}\label{p1b:der:g06:eq-G06-44}
\] This is the sign obtained by expressing the pulled-back wedge in the chosen factor order. Direct substitution in (G06.5) proves that it is a chain map to the complex with the factors interchanged. It is a unitary symmetry at equal radii, with matched Hopf markings. On \(L_{a,b}\) it transports the complex to \(L_{b,a}\). It is an endomorphism of this single line-bundle complex when \(a=b\); a paired sum can also be used for unequal charges. The line-bundle exchange used here is the canonical external tensor flip, without an extra scalar phase.

For \(a=b=n\ge0\), put \(d=n+1\). The harmonic space is in degree zero and exchange is the ordinary flip. Thus \[
\dim\mathscr Q_{n,n}^{+}=\frac{d(d+1)}2,\qquad
\dim\mathscr Q_{n,n}^{-}=\frac{d(d-1)}2.
\tag{G06.45}\label{p1b:der:g06:eq-G06-45}
\] For \(n\le-2\), put \(d=-n-1\). Each factor cohomology is in degree one, so exchange on their degree-two tensor product is minus the ordinary flip: \[
\dim H^{0,2}(B,L_{n,n})^{+}=\frac{d(d-1)}2,\qquad
\dim H^{0,2}(B,L_{n,n})^{-}=\frac{d(d+1)}2.
\tag{G06.46}\label{p1b:der:g06:eq-G06-46}
\] At \(n=-1\) there is no harmonic space. For every integer \(n\), these calculations give \[
\chi(B,L_{n,n})=(n+1)^2,\qquad
\operatorname{Str}\bigl(\mathcal B\mid H^{0,*}(B,L_{n,n})\bigr)
=n+1.
\tag{G06.47}\label{p1b:der:g06:eq-G06-47}
\] The negative sign in the second formula for \(n\le-2\) is the exterior-form exchange sign, even though the final harmonic degree is two.

The full Born exchange in G05 also swaps the scalar angles: \[
\mathcal B_{\rm full}(w_x,w_p,\theta)
=(w_p,w_x,S\theta),\qquad
S=\begin{pmatrix}0&1\\1&0\end{pmatrix}.
\tag{G06.48}\label{p1b:der:g06:eq-G06-48}
\] It preserves the product metric when \(R_x=R_p\) and \(S^{\mathsf T}GS=G\). But \(\det S=-1\). In oriented dimension two an orientation reversing isometry anticommutes with the compatible complex structure: \[
SJ_T=-J_TS.
\tag{G06.49}\label{p1b:der:g06:eq-G06-49}
\] This can also be verified immediately in (G06.12), where metric invariance forces \(a=c\). Therefore the scalar swap is antiholomorphic on \(E\). The sphere exchange is holomorphic on \(B\), so their product preserves neither the holomorphic nor the antiholomorphic splitting on all of \(B\times E\). In particular it sends the parallel torus \((0,1)\)-form to a \((1,0)\)-form and is not an endomorphism of the chosen \((0,*)\) basic Dolbeault complex. No Dolbeault equivariant supertrace for (G06.48) on that single complex is asserted.

It does transport the complex with \(J_T\) to the one with \(-J_T\). It also commutes with the scalar product Hardy projector, the scalar Laplacians, and the operators of Section~\ref{p1b:der:g06:s44} under the metric conditions above. Thus this qualification concerns the differential complex, not the existence of those scalar symmetries.

For comparison, a specified holomorphic scalar-torus involution does give a chain action on \(B\times E\). With the identity on \(E\), the untwisted equivariant index is zero: the two torus harmonic degrees have the same trace. With \(-I\) on \(T^2\), it is two: constants have trace one and the constant \((0,1)\)-form has trace minus one. These are different declared scalar actions from (G06.48).

There is also a precise reason why the full invariant complex family in Paper I-C does not supply the missing basic complex. In that source its residual real four-space is \[
W=\operatorname{span}(Z_x,Z_p)
 \oplus\operatorname{span}(\partial_{\theta_1},
                           \partial_{\theta_2}),
\] and full exchange acts as \(S\oplus S\). A full exchange-invariant complex structure has \([J_0,S\oplus S]=0\). If \(J_0\) preserved either displayed real two-plane, its restriction there would commute with the reflection \(S\). It would therefore preserve the two real one-dimensional eigenspaces of \(S\), which is incompatible with \(J_0^2=-I\). Thus an exchange-invariant \(J_0\) must mix Reeb and scalar directions. It neither preserves the torus plane in \(\mathcal D\) nor preserves the vertical two-plane needed for this complex quotient. This proves the domain distinction by linear algebra.

The chosen CR geometry has the natural compact product action \(U(2)_x\times U(2)_p\), with the matched exchange extension when appropriate. The unweighted Szegő projectors commute with these unitary actions. The larger one-factor group \(\mathrm{PU}(2,1)\) in Paper I concerns CR automorphisms; an arbitrary such map does not preserve the chosen round pseudohermitian metric or its unweighted \(L^2\) inner product. No exhaustive automorphism classification of the product is needed for these exports.

Paper~I-D, Section~\ref{p1b:app:g06} proves the convention-comparison identities and extended-case formulas used in the results above.


\hypertarget{p1b:sec:g07}{%
\subsection{Global metric geometry and deformations}\label{p1b:sec:g07}}


\begin{proposition}[Orthogonal product geometry and metric parameters]
\label{p1b:prop:g07-product}
Let \(g=g_x\oplus g_p\oplus G\), where the sphere metrics are the
round or stated Berger metrics and \(G\) is any positive symmetric
matrix on \(\mathbb R^2/(2\pi\mathbb Z)^2\). Then
\[
 d_g^2=d_{g_x}^2+d_{g_p}^2+d_G^2,\qquad
 \operatorname{diam}(g)^2
 =\operatorname{diam}(g_x)^2+\operatorname{diam}(g_p)^2+
   \operatorname{diam}(G)^2,
\]
\[
 \operatorname{inj}(g)=
 \min\{\operatorname{inj}(g_x),\operatorname{inj}(g_p),
                  \operatorname{inj}(G)\}.
\]
The torus distance minimizes the quadratic form \(G\) over its full
period lattice. The Berger cut, curvature and spectral formulas below
hold on this orthogonal subfamily. The larger invariant metric cones
are treated by their full algebraic metric tensors. Collapsing the
vertical directions and penalizing them to obtain the control metric
are distinct limits. Factor-normalized and full
eight-dimensional volume-normalized Ricci flows obey the separately
stated evolution equations.
\end{proposition}

Energy separation proves the product assertions. The subsequent geodesic, lattice, curvature and evolution calculations prove the quantitative formulas and their precise parameter ranges.

\hypertarget{p1b:der:g07:s48}{%
\subsubsection{Metric domain, and theorem inputs}\label{p1b:der:g07:s48}}

G07 concerns global Riemannian geometry and metric deformations. The starting geometry is the specified bi-invariant metric \[
g_0=R_x^2g_{\rm unit,x}\oplus R_p^2g_{\rm unit,p}\oplus G,
\qquad
M_8=S^3_x\times S^3_p\times(\mathbb R^2/2\pi\mathbb Z^2),
\quad R_i>0,\quad G=G^{\mathsf T}>0.
\tag{G07.1}\label{p1b:der:g07:eq-G07-1}
\] The six-dimensional leaf is \(P_v=S^3_x\times S^3_p\). All torus lattice statements retain the marked periods \(2\pi\); an off-diagonal \(G\) is not silently replaced by a rectangular torus.

The marked orthogonal Berger extension is \[
g_{\boldsymbol\varepsilon}
=g_{x,\varepsilon_x}\oplus g_{p,\varepsilon_p}\oplus G,\qquad
g_{i,\varepsilon_i}
=R_i^2\bigl((\omega_i^1)^2+(\omega_i^2)^2
                    +\varepsilon_i^2(\omega_i^3)^2\bigr),
\quad \varepsilon_i>0.
\tag{G07.2}\label{p1b:der:g07:eq-G07-2}
\] Here the unit right-invariant coframe is that of G05, with \(\omega_i^3=\eta_{0,i}\). Its dual frame has cyclic brackets \([\xi_1,\xi_2]=-2\xi_3\), and cyclic permutations. Thus \(\varepsilon_i\) scales vertical lengths and \(\varepsilon_i^2\) scales the vertical metric coefficient. Paper I-A's primary frame is right-invariant as well:
Corollary~\ref{p1b:cor:marked-hopf-frame-transport} gives its exact
horizontal sign change and preserves every squared metric term here.
Its separately marked left-invariant copy is the inverse pullback of that
primary structure, with the circle and orientation transport of
Proposition~\ref{p1:prop:marked-hopf-inversion-transport}.

The marked Berger product has seven parameters: the two pairs \((R_i,\varepsilon_i)\) and the three independent entries of \(G\). Its exchange-fixed part has four: the two sphere parameters agree and \(S^{\mathsf T}GS=G\), where \(S\) swaps the scalar angles. It is a specified subfamily of Paper I-C's twelve-dimensional Hopf-invariant metric cone, which also permits mixing among Reeb and scalar directions. Exact product-distance identities below require the orthogonal product splitting. Paper~I-D, Section~\ref{p1b:app:g07:4} supplies algebraic curvature, flow, and scalar-operator formulas for the larger invariant cones.

The source inputs are Paper I, p1:sec:geodesics, p1:sec:berger-collapse-rigorous, p1:thm:berger-sectional, p1:thm:berger-spectrum, p1:thm:hamilton-ricci-flow-convergence, and p1:thm:lowest-dirac-pair-stability; and Paper I-C, p1b:thm:metric-cones and p1b:thm:metrics. The one-factor global Berger cut theorem is imported explicitly from the SU(2) result in \href{https://arxiv.org/abs/1504.05472}{Podobryaev--Sachkov, §9, Theorem 5}. Its normalization is checked against the directly derived exponential map. The diameter comparison uses \href{https://arxiv.org/abs/1710.02945}{Podobryaev, Theorem 1}. General small product-factor Ricci-flow convergence uses the Hamilton theorem already stated in Paper I. The particular Dirac first-variation coefficient is checked in \href{https://arxiv.org/abs/1605.08589}{Fang--Levitin--Vassiliev, Theorem 2.1}. All remaining product, curvature, spectral-limit, and homogeneous flow formulas below are derived directly.

\hypertarget{p1b:der:g07:s49}{%
\subsubsection{Exact product distance, diameter, and injectivity}\label{p1b:der:g07:s49}}

For unit quaternions \(q_i,q'_i\), put \[
\theta_i=\arccos\operatorname{Re}(q'_i\bar q_i)\in[0,\pi],
\qquad
d_i(q_i,q'_i)=R_i\theta_i.
\] If \(v\in\mathbb R^2\) represents the scalar-angle difference, \[
d_T(0,[v])^2=\min_{n\in\mathbb Z^2}
                    (v+2\pi n)^{\mathsf T}G(v+2\pi n).
\] Then \[
\boxed{\quad d_{g_0}(z,z')^2=d_x^2+d_p^2+d_T^2.\quad}
\tag{G07.3}\label{p1b:der:g07:eq-G07-3}
\] To prove the lower bound, write the factor speeds of a path as \(s_x,s_p,s_T\). Minkowski's integral inequality gives \[
L(\gamma)=\int\|(s_x,s_p,s_T)\|_2
\ge\left\|\left(\int s_x,\int s_p,\int s_T\right)\right\|_2
\ge\sqrt{d_x^2+d_p^2+d_T^2}.
\] For equality choose minimizing factor paths, each with constant speed over the same parameter interval. This proof also applies to every orthogonal Berger product and to the quotient limits.

Define the closed lattice Voronoi cell \[
V_G=\{v:\ v^{\mathsf T}Gn\le\pi n^{\mathsf T}Gn
                         \text{ for every }n\in\mathbb Z^2\},
\] and the lattice quantities \[
\ell_G=\min_{n\ne0}\sqrt{n^{\mathsf T}Gn},\qquad
\rho_G=\max_{v\in V_G}\sqrt{v^{\mathsf T}Gv}.
\tag{G07.4}\label{p1b:der:g07:eq-G07-4}
\] The inequalities for \(n\) and \(-n\) are both included. The torus injectivity radius is \(\pi\ell_G\), and its diameter is \(\rho_G\). Indeed the first boundary of a centered Voronoi cell occurs halfway to a shortest lattice vector; its farthest point gives the greatest distance from the torus origin. Homogeneity makes this independent of the center.

Because factor endpoints can be chosen independently, \[
\boxed{
\begin{aligned}
\operatorname{diam}(M_8,g_0)
 &=\sqrt{\pi^2R_x^2+\pi^2R_p^2+\rho_G^2},\\
\operatorname{inj}(M_8,g_0)
 &=\min(\pi R_x,\pi R_p,\pi\ell_G).
\end{aligned}}
\tag{G07.5}\label{p1b:der:g07:eq-G07-5}
\] For \(P_v\), omit the torus terms. The injectivity formula also follows from the normal domain described in Section~\ref{p1b:der:g07:s50}; it is a minimum, whereas diameter is a Euclidean sum.

For an exact two-dimensional evaluation of \(\rho_G\), choose a unimodularly reduced lattice basis whose Gram matrix is \[
H=\begin{pmatrix}a&c\\c&b\end{pmatrix},
\qquad 0<a\le b,\quad 0\le2c\le a.
\] Then \[
\ell_G=\sqrt a,\qquad
\rho_G=\pi\sqrt{\frac{ab(a+b-2c)}{ab-c^2}}.
\tag{G07.6}\label{p1b:der:g07:eq-G07-6}
\] The reduced basis makes the triangles with vertices \(0,u,v\) and their lattice translates Delaunay triangles. Their common circumradius is \(|u||v||u-v|/(2|u\wedge v|)\). Substituting the lattice Gram matrix \((2\pi)^2H\) gives (G07.6); the rectangular case is the right-triangle limit. Equivalently one can compute the vertices of the half-space description (G07.4), without choosing a reduced basis.

For example, for \(G=\begin{pmatrix}5&3\\3&5\end{pmatrix}\), the basis \((e_1-e_2,e_1)\) has \(a=4,b=5,c=2\). Hence \[
\operatorname{inj}(T^2_G)=2\pi,\qquad
\operatorname{diam}(T^2_G)=\frac52\pi.
\] For \(v=(9\pi/10,9\pi/10)\), coordinatewise rounding selects zero lattice shift, but its squared length is \(324\pi^2/25\). The nearest shifts instead give \(104\pi^2/25\). This exact example shows why the marked torus metric must enter the lattice minimization.

\hypertarget{p1b:der:g07:s50}{%
\subsubsection{Geodesics, cut and conjugate loci, and periodicity}\label{p1b:der:g07:s50}}

The Levi--Civita connection of an orthogonal product is the product connection, as follows immediately from the Koszul formula and vanishing cross-factor brackets. Thus each component of a geodesic is an affinely parametrized factor geodesic. For an ambient sphere point \(p_i\) of length \(R_i\), initial velocity \(v_i\), and speed \(c_i=|v_i|\), \[
\gamma_i(t)=
\cos(c_it/R_i)p_i+\frac{R_i}{c_i}\sin(c_it/R_i)v_i,
\tag{G07.7}\label{p1b:der:g07:eq-G07-7}
\] with the continuous constant-path interpretation when \(c_i=0\). The torus component is \([\theta_0+tw]\). The total speed is \(\sqrt{c_x^2+c_p^2+w^{\mathsf T}Gw}\).

The open injective normal domain at a point is the product \[
\{|v_x|<\pi R_x\}\times\{|v_p|<\pi R_p\}
                   \times\operatorname{int}V_G.
\tag{G07.8}\label{p1b:der:g07:eq-G07-8}
\] Every point represented there has unique minimizing components and a nonsingular exponential differential. On its boundary at least one component ceases to minimize uniquely or becomes conjugate. Consequently \[
\begin{aligned}
\operatorname{Cut}(p_x,p_p,\theta_0)
={}&(\{-p_x\}\times S_p^3\times T^2)\\
&\cup(S_x^3\times\{-p_p\}\times T^2)\\
&\cup(S_x^3\times S_p^3\times[\theta_0+\partial V_G]).
\end{aligned}
\tag{G07.9}\label{p1b:der:g07:eq-G07-9}
\] This is a union of factor cut sets, not a single product antipode.

For a nonzero torus velocity define \[
t_T(w)=\pi\min_{\substack{n\in\mathbb Z^2\\w^{\mathsf T}Gn>0}}
\frac{n^{\mathsf T}Gn}{w^{\mathsf T}Gn},
\qquad t_T(0)=+\infty.
\] This is the first time \(tw\) reaches \(\partial V_G\). The exact product cut time is \[
t_{\rm cut}=
\min\left(\frac{\pi R_x}{c_x},
          \frac{\pi R_p}{c_p},t_T(w)\right),
\tag{G07.10}\label{p1b:der:g07:eq-G07-10}
\] where a zero sphere speed contributes \(+\infty\). A segment is minimizing exactly while all its components are minimizing, by (G07.3).

The torus has no conjugate points. A moving sphere factor has conjugate times \(m\pi R_i/c_i\), \(m=1,2,\ldots\), each with two-dimensional exponential kernel. Product exponential differentials are block diagonal, so coincident multiplicities add. In particular \[
t_{\rm conj,1}
=\min\left(\frac{\pi R_x}{c_x},\frac{\pi R_p}{c_p}\right).
\tag{G07.11}\label{p1b:der:g07:eq-G07-11}
\] The two sphere Jacobi fields in each transverse direction are proportional to \(\sin(c_it/R_i)\); radial variations and flat torus variations have no such zeros. This proves the claim, including the absence of a contribution from a stationary factor. A geodesic moving only in the torus has finite cut time and no conjugate time.

A nonconstant product geodesic closes exactly when some \(T>0\) satisfies \[
\frac{Tc_x}{R_x},\ \frac{Tc_p}{R_p}\in2\pi\mathbb Z
\quad\text{for the moving sphere factors},\qquad
Tw\in2\pi\mathbb Z^2.
\tag{G07.12}\label{p1b:der:g07:eq-G07-12}
\] These are simultaneous commensurability conditions. Without them the orbit can be dense in a torus inside the product of its two great circles and the scalar torus. More precisely its closure dimension is the rational rank of the nonzero frequencies \((c_x/R_x,c_p/R_p,w_1,w_2)\), by the character description of a one-parameter subgroup of a torus. Thus the one-sphere assertion that every geodesic is closed does not transfer to the product.

\hypertarget{p1b:der:g07:s51}{%
\subsubsection{Curvature, volume, and comparison geometry}\label{p1b:der:g07:s51}}

For tangent vectors \(U,V\), the product curvature numerator is \[
\langle R(U,V)V,U\rangle
=\frac{|U_x\wedge V_x|^2}{R_x^2}
 +\frac{|U_p\wedge V_p|^2}{R_p^2}.
\tag{G07.13}\label{p1b:der:g07:eq-G07-13}
\] Divide by \(|U\wedge V|^2\) for a general sectional curvature. For an orthonormal pair the denominator is one. Mixed and torus planes are flat, giving \[
0\le K\le\max(R_x^{-2},R_p^{-2}).
\] The curvature operator on \(\Lambda^2TM_8\) has eigenvalues \(R_x^{-2}\) and \(R_p^{-2}\), each with multiplicity three, and zero with multiplicity 22. The Ricci eigenvalues are \[
\left(\frac2{R_x^2}\right)^{[3]},\qquad
\left(\frac2{R_p^2}\right)^{[3]},\qquad 0^{[2]}.
\] In particular \[
S=\frac6{R_x^2}+\frac6{R_p^2},\qquad
|\mathrm{Rm}|^2=|\mathrm{Ric}|^2
=12(R_x^{-4}+R_p^{-4}),
\] \[
\operatorname{Vol}(M_8,g_0)
=16\pi^6R_x^3R_p^3\sqrt{\det G}.
\tag{G07.14}\label{p1b:der:g07:eq-G07-14}
\] These follow by tracing the curvature blocks and multiplying the factor volume densities. They provide normalization data for the later heat calculations in G08.

The leaf \(P_v\) is Einstein exactly when \(R_x=R_p\), with Einstein constant \(2/R_x^2\). The full product \(M_8\) is never Einstein in this metric family: its torus Ricci eigenvalues vanish while its sphere eigenvalues are positive. Paper~I-D, Section~\ref{p1b:app:g07:3} strengthens this to the absence of a homogeneous Einstein metric on this topology, without making a claim about arbitrary nonhomogeneous Einstein metrics.

The round product is locally symmetric and has parallel curvature. It is not a space form. It is also not conformally flat: with \(a=R_x^{-2}\) and \(b=R_p^{-2}\), the eight-dimensional trace decomposition gives \[
|W|^2
=|\mathrm{Rm}|^2-\frac46|\mathrm{Ric}|^2+\frac2{7\cdot6}S^2
=4(a^2+b^2)+\frac{12}{7}(a+b)^2>0.
\tag{G07.15}\label{p1b:der:g07:eq-G07-15}
\]

Before the cut locus, the exponential volume Jacobian is \[
J(v_x,v_p,v_T)
=\left(\frac{\sin(|v_x|/R_x)}{|v_x|/R_x}\right)^2
 \left(\frac{\sin(|v_p|/R_p)}{|v_p|/R_p}\right)^2.
\tag{G07.16}\label{p1b:der:g07:eq-G07-16}
\] This follows from the transverse sphere Jacobi fields and the flat torus differential. Integrating it over (G07.8) intersected with a tangent ball gives the exact ball volume; the cut locus has measure zero. In particular, with \(\omega_8\) the Euclidean unit-ball volume, \[
\operatorname{Vol}B(p,r)
=\omega_8r^8
\left[1-\frac{r^2}{10}(R_x^{-2}+R_p^{-2})+O(r^4)\right].
\tag{G07.17}\label{p1b:der:g07:eq-G07-17}
\] The coefficient follows by expanding \((\sin z/z)^2=1-z^2/3+O(z^4)\) and using \(\int_{B^8_r}|v_i|^2/\operatorname{Vol}(B^8_r)=3r^2/10\).

The injective normal domain is star-shaped, and each Jacobian factor decreases along a ray until its first cut. Writing the volume integral over the unit tangent ball therefore proves directly that \(\operatorname{Vol}B(p,r)/(\omega_8r^8)\) is nonincreasing and at most one. This is the appropriate nonnegative-Ricci comparison for the full product. The positive three-dimensional Bonnet--Myers bound in Paper I applies to the individual spheres; it is not a diameter bound for \(M_8\), whose torus scale can grow arbitrarily. Likewise a product triangle has component spherical trigonometry and a lattice-minimized flat component, not a single spherical law of cosines.

\hypertarget{p1b:der:g07:s52}{%
\subsubsection{Berger curvature and the complete geodesic formula}\label{p1b:der:g07:s52}}

On one factor set \[
A=R^2,\qquad C=R^2\varepsilon^2,\qquad u=C/A=\varepsilon^2.
\] In the orthonormal frame \(e_1=\xi_1/R,\ e_2=\xi_2/R,\ e_3=\xi_3/(R\varepsilon)\), the brackets are \[
[e_1,e_2]=-\frac{2\varepsilon}{R}e_3,\qquad
[e_2,e_3]=-\frac2{R\varepsilon}e_1,\qquad
[e_3,e_1]=-\frac2{R\varepsilon}e_2.
\] The Koszul formula then gives \[
\begin{array}{lll}
\nabla_{e_1}e_2=-\varepsilon e_3/R,&
\nabla_{e_2}e_1=\varepsilon e_3/R,&
\nabla_{e_1}e_3=\varepsilon e_2/R,\\
\nabla_{e_2}e_3=-\varepsilon e_1/R,&
\nabla_{e_3}e_1=(\varepsilon-2/\varepsilon)e_2/R,&
\nabla_{e_3}e_2=(2/\varepsilon-\varepsilon)e_1/R,
\end{array}
\] and \(\nabla_{e_i}e_i=0\). Substitution in \(R(X,Y)Z=\nabla_X\nabla_YZ-\nabla_Y\nabla_XZ-\nabla_{[X,Y]}Z\) proves \[
\boxed{
K_{12}=\frac{4-3u}{A},\qquad K_{13}=K_{23}=\frac uA,
\qquad
\mathrm{Ric}=\operatorname{diag}
\left(\frac{4-2u}{A},\frac{4-2u}{A},\frac{2u}{A}\right).
}
\tag{G07.18}\label{p1b:der:g07:eq-G07-18}
\] The curvature operator is diagonal in these three bivectors; hence every sectional curvature is between the smallest and largest of the displayed values. Furthermore \[
S_i=\frac{8-2u_i}{A_i},\quad
|\mathrm{Rm}_i|^2
=\frac{4[(4-3u_i)^2+2u_i^2]}{A_i^2},\quad
\operatorname{Vol}(S_i^3,g_{i,\varepsilon_i})
=2\pi^2A_i\sqrt{C_i}.
\tag{G07.19}\label{p1b:der:g07:eq-G07-19}
\] These give the exact one-factor ranges:

\begin{longtable}[]{@{}ll@{}}
\toprule\noalign{}
Property & Condition \\
\midrule\noalign{}
\endhead
\bottomrule\noalign{}
\endlastfoot
Positive sectional curvature & \(0<u<4/3\) \\
Nonnegative sectional curvature & \(0<u\le4/3\) \\
Positive Ricci curvature & \(0<u<2\) \\
Positive scalar curvature & \(0<u<4\) \\
Einstein & \(u=1\) \\
\end{longtable}

For the product, mixed planes remain flat, Ricci is the direct sum with two zero torus eigenvalues, scalar curvature is \(S_x+S_p\), and \[
\operatorname{Vol}(M_8,g_{\boldsymbol\varepsilon})
=16\pi^6A_x\sqrt{C_x}A_p\sqrt{C_p}\sqrt{\det G}.
\tag{G07.20}\label{p1b:der:g07:eq-G07-20}
\] Nonnegative product sectional curvature is equivalent to both \(u_i\le4/3\). Positive product scalar curvature requires the sum \(S_x+S_p>0\), not separate positivity assumptions.

The geodesic equation can be solved explicitly. Write the initial right-trivialized velocity as \(a\,j+b\,k+c\,i\), so its squared speed is \(A(a^2+b^2)+Cc^2\). For the right-trivialized velocity \(v=\dot q q^{-1}\) and momentum \(m=\operatorname{diag}(A,A,C)v\), variation of the energy gives \[
\dot m=2v\times m.
\] Thus \(c\) is constant and \((a(t),b(t))\) rotates with angular speed \(2(1-u)c\). Integrating the group equation gives \[
\boxed{
q(t)=
\exp\!\bigl((1-u)ct\,i\bigr)\,
\exp\!\bigl(t(aj+bk+uc\,i)\bigr)\,q(0).
}
\tag{G07.21}\label{p1b:der:g07:eq-G07-21}
\] Differentiation gives the stated rotating velocity and the initial vector; the energy is constant. This proves the geodesic equation, including the round, horizontal, and purely vertical cases, without assuming a geodesic is a one-parameter subgroup for a general Berger metric.

For unit-speed geodesics put \[
\eta=u^{-1}-1>-1,\qquad
h=\frac{uc}{\sqrt{a^2+b^2+u^2c^2}}\in[-1,1],\qquad
\tau=\frac{t}{R\sqrt{1+\eta h^2}}.
\] With \(\phi\) the initial horizontal angle, (G07.21) becomes \[
q(\tau)=e^{i\eta h\tau}
\left[\cos\tau+\sin\tau
\left(ih+\sqrt{1-h^2}(\cos\phi\,j+\sin\phi\,k)\right)\right].
\tag{G07.22}\label{p1b:der:g07:eq-G07-22}
\] For a different initial point multiply on the right by that point. This formula also supplies an explicit endpoint equation for computing distance: retain the solutions up to their cut times in Section~\ref{p1b:der:g07:s53}. Existence of a minimizing solution follows from compactness and the Riemannian completeness of each positive metric.

\hypertarget{p1b:der:g07:s53}{%
\subsubsection{Global Berger cut data, diameter, and injectivity}\label{p1b:der:g07:s53}}

Define \[
F_{\eta,h}(\tau)
=(1+\eta h^2)\sin\tau+
              \eta(1-h^2)\tau\cos\tau.
\] Differentiating (G07.22), and removing its invertible quaternion rotation, gives the exponential Jacobian up to a nonzero orientation and coordinate factor: \[
\sin\tau\,F_{\eta,h}(\tau).
\tag{G07.23}\label{p1b:der:g07:eq-G07-23}
\] The determinant follows from differentiation of the displayed exponential map. For \(\eta\le0\), \(F/\sin\tau>0\) on \(0<\tau<\pi\): use \(\tau\cot\tau\le1\) and \(1+\eta>0\). The first conjugate parameter is therefore \(\pi\). For \(\eta>0\), \(|h|<1\), the first zero of \(F\) lies uniquely in \((\pi/2,\pi)\), characterized by \[
\tan\tau_c
=-\frac{\eta(1-h^2)}{1+\eta h^2}\tau_c.
\tag{G07.24}\label{p1b:der:g07:eq-G07-24}
\] For \(|h|=1\) it is \(\pi\). These statements follow from monotonicity of \(\tan\tau\) and the linear right-hand side on that interval. At the polar directions, use transverse initial-velocity coordinates or the vertical Jacobi equation to remove the angular-coordinate singularity. Arclength conjugate time is \(R\sqrt{1+\eta h^2}\tau_c\). The later conjugate times are the positive zeros of (G07.23), with the same coordinate qualification at the poles.

For \(\eta>0\) and \(h\ne0\), let \(\tau_3(h)\) be the first positive root of \[
\cos\tau\sin(\eta h\tau)+h\sin\tau\cos(\eta h\tau)=0,
\] and define \(\tau_3(0)=\tau_c(0)\) by continuity. The global SU(2) cut theorem gives \[
\boxed{
t_{\rm cut}(h)=R\sqrt{1+\eta h^2}
\begin{cases}
\pi,&-1<\eta\le0,\\
\tau_3(h),&\eta>0.
\end{cases}}
\tag{G07.25}\label{p1b:der:g07:eq-G07-25}
\] This is the SU(2), rather than SO(3), result in \href{https://arxiv.org/abs/1504.05472}{Podobryaev--Sachkov, §9}. Their metric coefficients translate as \(I_1=I_2=R^2/4,\ I_3=R^2\varepsilon^2/4\). The global optimality conclusion is a theorem input; the Jacobian calculation alone would not prove it.

The cut locus is the image of these cut endpoints. For \(\varepsilon>1\) it is the vertical arc \[
\{-e^{i\pi\eta h}: -1\le h\le1\};
\] at \(\varepsilon=1\) it is the antipode. For \(0<\varepsilon<1\), it is a disk in the two-sphere with zero \(i\)-component, bounded by the conjugate circle \[
\cos\tau_c(0)+\sin\tau_c(0)
                       (\cos\phi\,j+\sin\phi\,k).
\] This is the cut-set description in the same theorem.

The exact diameter is \[
\boxed{
D(R,\varepsilon)=\pi R
\begin{cases}
\dfrac{1}{2\sqrt{1-\varepsilon^2}},
       &0<\varepsilon<1/\sqrt2,\\[4pt]
\varepsilon,&1/\sqrt2\le\varepsilon\le1,\\
1,&\varepsilon\ge1.
\end{cases}}
\tag{G07.26}\label{p1b:der:g07:eq-G07-26}
\] It follows from maximizing (G07.25), and agrees after scaling with \href{https://arxiv.org/abs/1710.02945}{Podobryaev's diameter theorem}. Paper~I-D, Section~\ref{p1b:app:g07:1} gives the derivative calculation and, by taking the minimum instead, obtains \[
\boxed{
\iota(R,\varepsilon)=
\begin{cases}
R\min\{\tau_c(0),\pi\varepsilon\},&0<\varepsilon<1,\\
\pi R/\varepsilon,&\varepsilon\ge1.
\end{cases}}
\tag{G07.27}\label{p1b:der:g07:eq-G07-27}
\] Here \(\tau_c(0)\) uses \(\eta=\varepsilon^{-2}-1\). Thus shrinking and enlarging the fibers can both send injectivity radius to zero, while the diameter remains bounded.

All product global data now follow without an additional cut theorem: \[
\begin{aligned}
\operatorname{diam}(M_8,g_{\boldsymbol\varepsilon})
 &=\sqrt{D(R_x,\varepsilon_x)^2+
          D(R_p,\varepsilon_p)^2+\rho_G^2},\\
\operatorname{inj}(M_8,g_{\boldsymbol\varepsilon})
 &=\min\{\iota(R_x,\varepsilon_x),
          \iota(R_p,\varepsilon_p),\pi\ell_G\}.
\end{aligned}
\tag{G07.28}\label{p1b:der:g07:eq-G07-28}
\] The cut locus is the union of the three pulled-back factor cut loci. For factor speeds \(c_i\), use \(t_{{\rm cut},i}/c_i\) and \(t_{{\rm conj},i}/c_i\), with \(+\infty\) for stationary factors, and take the minimum along with the torus cut time where applicable. This follows from the same distance and exponential-map product proofs as Section~\ref{p1b:der:g07:s50}.

\hypertarget{p1b:der:g07:s54}{%
\subsubsection{Collapse and the reverse control-metric limit}\label{p1b:der:g07:s54}}

The Hopf projection remains a Riemannian submersion for every \(\varepsilon>0\), onto the sphere of radius \(R/2\). The vertical generator has constant length \(R\varepsilon\). Since it is Killing, its integral circles are geodesics and have length \(2\pi R\varepsilon\). For any two endpoints, horizontal lifting of a shortest base path followed by a fiber arc gives \[
d_{S^2}(\pi p,\pi q)
\le d_\varepsilon(p,q)
\le d_{S^2}(\pi p,\pi q)+\pi R\varepsilon.
\tag{G07.29}\label{p1b:der:g07:eq-G07-29}
\]

Collapse any subset \(I\subseteq\{x,p\}\) of the Hopf fibers, keeping \(R_i,G\) and the other sphere metric fixed. Let \(Y_I\) replace each collapsed \(S^3_i\) by \(S^2(R_i/2)\). In the fiber over a base point the intrinsic product diameter is at most \[
\delta_I=\pi\sqrt{\sum_{i\in I}R_i^2\varepsilon_i^2}.
\] The product horizontal lift and the final fiber adjustment prove \[
0\le d_{\boldsymbol\varepsilon}(p,q)
     -d_{Y_I}(\pi_Ip,\pi_Iq)\le\delta_I.
\] The graph of the surjective projection is a correspondence of distortion at most \(\delta_I\), hence in the standard correspondence convention \[
d_{\rm GH}((M_8,g_{\boldsymbol\varepsilon}),Y_I)
\le\frac{\delta_I}{2}.
\tag{G07.30}\label{p1b:der:g07:eq-G07-30}
\] Normalized Haar probability measure pushes forward exactly to the quotient probability measure. This also gives measured collapse with that normalization; the unnormalized volumes instead tend to zero by (G07.20). For the shrinking parameters \(0<\varepsilon_i\le1\), the factor sectional curvatures remain bounded by \(4/R_i^2\).

The opposite limit keeps the same horizontal lengths but makes vertical velocities costly. As both \(\varepsilon_i\to\infty\), \[
d_{\boldsymbol\varepsilon}\ \uparrow\ d_{\mathcal D},
\qquad
\mathcal D=\mathcal H_x\oplus\mathcal H_p\oplus TT^2,
\tag{G07.31}\label{p1b:der:g07:eq-G07-31}
\] uniformly on the compact endpoint space. For a direct proof, horizontal paths provide the upper bound. Minimizers of bounded energy have a weakly convergent subsequence of velocities, while their vertical \(L^2\) norms are \(O(\varepsilon_i^{-1})\). The limit is horizontal, and lower semicontinuity gives the reverse distance inequality. Monotonicity and continuity then give uniform convergence. This is exactly the vertical-penalty construction proved in G05, Paper~I-D, Section~\ref{p1b:app:g05:3}, with the reciprocal parameter. If only one fiber is made long, the same proof constrains just that factor.

\begin{longtable}[]{@{}
  >{\raggedright\arraybackslash}p{(\columnwidth - 4\tabcolsep) * \real{0.3333}}
  >{\raggedright\arraybackslash}p{(\columnwidth - 4\tabcolsep) * \real{0.3333}}
  >{\raggedright\arraybackslash}p{(\columnwidth - 4\tabcolsep) * \real{0.3333}}@{}}
\toprule\noalign{}
\begin{minipage}[b]{\linewidth}\raggedright
Metric limit
\end{minipage} & \begin{minipage}[b]{\linewidth}\raggedright
Limit space or distance
\end{minipage} & \begin{minipage}[b]{\linewidth}\raggedright
Dimension statement
\end{minipage} \\
\midrule\noalign{}
\endhead
\bottomrule\noalign{}
\endlastfoot
Both \(\varepsilon_i\to0\) & \(S_x^2\times S_p^2\times T^2_G\) & Riemannian dimension 6 \\
One \(\varepsilon_i\to0\) & One quotient sphere, one three-sphere, and \(T^2_G\) & Riemannian dimension 7 \\
Both \(\varepsilon_i\to\infty\) & G05's \(d_{\mathcal D}\) on \(M_8\) & Topological dimension 8; homogeneous/Hausdorff dimension 10 \\
\end{longtable}

The last limit has no Riemannian curvature tensor. Indeed the Berger curvature expressions can diverge when fibers are made long, even though the distances converge uniformly.

\hypertarget{p1b:der:g07:s55}{%
\subsubsection{Exact scalar spectra and quantitative spectral limits}\label{p1b:der:g07:s55}}

Use the integer geometric charges of G05--G06, \(Q_i=-iX_{H,i}\). The invariant volume density changes only by a constant, so on normalized Haar \(L^2\) the scalar Laplacian is exactly \[
\Delta_{\boldsymbol\varepsilon}
=\sum_{i=x,p}
\left(L_{H,i}+\frac{Q_i^2}{R_i^2\varepsilon_i^2}\right)
+\Delta_T.
\] On scalar spherical degree \(k_i\) and charge \(q_i=-k_i,-k_i+2,\ldots,k_i\), this gives \[
\boxed{
\lambda=
\sum_{i=x,p}\frac{k_i(k_i+2)+(\varepsilon_i^{-2}-1)q_i^2}{R_i^2}
+n^{\mathsf T}G^{-1}n.
}
\tag{G07.32}\label{p1b:der:g07:eq-G07-32}
\] For fixed \((k_x,q_x,k_p,q_p,n)\), the complex multiplicity is \((k_x+1)(k_p+1)\). Coincident labels contribute additively. This derives the complete scalar spectrum for this family from the already established Casimir/charge decomposition.

The one-factor first positive value is \[
\lambda_1(R,\varepsilon)
=\frac1{R^2}\min\{8,\,2+\varepsilon^{-2}\}.
\tag{G07.33}\label{p1b:der:g07:eq-G07-33}
\] The \(q=0\) modes have even \(k\), with first value \(8/R^2\) and multiplicity three. If \(q\ne0\), then \[
k(k+2)-q^2+\varepsilon^{-2}q^2\ge2+\varepsilon^{-2},
\] with equality precisely at \(k=1,q=\pm1\), of total multiplicity four. At \(\varepsilon=1/\sqrt6\) these eigenspaces coincide in value and their multiplicities add to seven. The full product gap is \[
\min\{\lambda_1(R_x,\varepsilon_x),
       \lambda_1(R_p,\varepsilon_p),\lambda_{T,1}\},
\quad
\lambda_{T,1}=\min_{n\ne0}n^{\mathsf T}G^{-1}n.
\tag{G07.34}\label{p1b:der:g07:eq-G07-34}
\] All these spectra restrict directly to independent scalar cutoffs, because the Casimir bands commute with the operator.

For collapse of the factors in \(I\), let \(\Pi_I\) average over their Hopf circles. Its range is the pullback of \(L^2(Y_I)\), and on that range the Laplacian is exactly \(\Delta_{Y_I}\), at every positive parameter. Every complementary mode has eigenvalue at least \[
\Lambda_I=\min_{i\in I}
\frac{2+\varepsilon_i^{-2}}{R_i^2}.
\] Equality is attained by a charge-one mode in a minimizing factor and constants elsewhere. Therefore for \(z>0\) and \(t>0\), with pullback understood, \[
\boxed{
\begin{aligned}
\|(\Delta_{\boldsymbol\varepsilon}+z)^{-1}
 -(\Delta_{Y_I}+z)^{-1}\Pi_I\|
 &=\frac1{z+\Lambda_I},\\
\|e^{-t\Delta_{\boldsymbol\varepsilon}}
 -e^{-t\Delta_{Y_I}}\Pi_I\|
 &=e^{-t\Lambda_I}.
\end{aligned}}
\tag{G07.35}\label{p1b:der:g07:eq-G07-35}
\] This follows by taking the supremum of the scalar diagonal differences. It proves quantitative spectral decoupling and shows that a fixed spectral window becomes exactly basic as soon as \(\Lambda_I\) exceeds its upper endpoint. For a finite band the same estimates hold as upper bounds; equality requires the relevant charge-one mode to be retained.

There is also a norm-resolvent estimate in the long-fiber direction. Let \[
L=L_{H,x}+L_{H,p}+\Delta_T,\qquad
M=\max_i(R_i^2/\varepsilon_i^2).
\] For \(z>0\), \[
\boxed{
\|(\Delta_{\boldsymbol\varepsilon}+z)^{-1}-(L+z)^{-1}\|
\le\frac M4.
}
\tag{G07.36}\label{p1b:der:g07:eq-G07-36}
\] To prove it, write a horizontal eigenvalue as \(\lambda_H\) and its vertical penalty as \(b=\sum_iq_i^2/(R_i^2\varepsilon_i^2)\). Since \(\lambda_H\ge2\sum_i|q_i|/R_i^2\), \[
b\le\frac M4\lambda_H^2.
\] For nonzero charge, \[
0\le\frac1{\lambda_H+z}
 -\frac1{\lambda_H+b+z}
=\frac b{(\lambda_H+z)(\lambda_H+b+z)}
\le\frac M4.
\] For zero charge the difference vanishes. This proves (G07.36) on the complete basis. The bounded-curvature quotient limit and the control-metric limit therefore have distinct, explicitly derived spectral descriptions.

\hypertarget{p1b:der:g07:s56}{%
\subsubsection{Ricci flow, extinction, and the choice of normalization}\label{p1b:der:g07:s56}}

Use the convention \(\partial_tg=-2\mathrm{Ric}\). For the orthogonal Berger product, the block curvature formulas give the closed system \[
\boxed{
\dot A_i=-8+4u_i,\qquad
\dot C_i=-4u_i^2,\qquad
\dot u_i=\frac{8u_i(1-u_i)}{A_i},\qquad
\dot G=0,\quad u_i=C_i/A_i.
}
\tag{G07.37}\label{p1b:der:g07:eq-G07-37}
\] For example, \(\mathrm{Ric}(\xi_1,\xi_1)=4-2u\) and \(\mathrm{Ric}(\xi_3,\xi_3)=2u^2\), which fixes both scale coefficients without confusing orthonormal and invariant frames. The quotient rule gives the equation for \(u\). The product and axial symmetries are preserved by the equations; uniqueness of the compact Ricci flow identifies this ODE solution with the geometric flow while the coefficients remain positive.

Every positive initial Berger metric rounds in shape before extinction. This includes \(u_0\ge2\); initial positive Ricci curvature is not needed for this homogeneous ODE conclusion. If \(u_0=1\), then \(A=C=A_0-4t\), with extinction time \(T=A_0/4\). If \(u_0\ne1\), division of the first-order equations gives \[
\frac{dC}{du}=-\frac{C}{2(1-u)},\qquad
\frac{C^2}{|1-u|}=\kappa^2,\qquad
\kappa=\frac{C_0}{\sqrt{|1-u_0|}}.
\tag{G07.38}\label{p1b:der:g07:eq-G07-38}
\] Thus \(C=\kappa\sqrt{|1-u|}\) and \(A=C/u\). The sign of \(\dot u\) moves \(u\) monotonically toward one. Writing \(v=\sqrt{|1-u_0|}\) and integrating \(dt=du/\dot u\) gives the exact finite extinction time \[
\boxed{
T=
\begin{cases}
\displaystyle\frac{\kappa}{8}
 \left(\frac{v}{1-v^2}+\operatorname{arctanh}v\right),
 &0<u_0<1,\quad v=\sqrt{1-u_0},\\[6pt]
\displaystyle A_0/4,&u_0=1,\\[4pt]
\displaystyle\frac{\kappa}{8}
 \left(\frac{v}{1+v^2}+\arctan v\right),
 &u_0>1,\quad v=\sqrt{u_0-1}.
\end{cases}}
\tag{G07.39}\label{p1b:der:g07:eq-G07-39}
\] The antiderivatives follow from \(\int(1-v^2)^{-2}dv\) and \(\int(1+v^2)^{-2}dv\). They also show directly that \(u\to1\) as \(t\uparrow T\). Since \(\dot C\to-4\) and \(A=C/u\), \[
A(t)\sim C(t)\sim4(T-t),\qquad
|u(t)-1|\sim\frac{16(T-t)^2}{\kappa^2}
\quad(u_0\ne1).
\tag{G07.40}\label{p1b:der:g07:eq-G07-40}
\] The formula for \(T\) has the continuous round limit \(A_0/4\).

For the product, its maximal unnormalized time is \(T_*=\min(T_x,T_p)\). At that time the first extinct sphere has diameter tending to zero. If exactly one sphere becomes extinct, the compact Gromov--Hausdorff limit is the surviving Berger three-sphere at \(T_*\), times the unchanged flat torus; if both become extinct, it is \(T^2_G\). These follow from the product distance formula and \(A_i,C_i\to0\) in the extinct factors.

There are two distinct normalization operations. First, apply the three-dimensional volume normalization to each sphere separately and keep \(G\) fixed. On a single factor this means \(\dot g=-2\mathrm{Ric}+(2/3)Sg\), and hence \[
\dot A=\frac83(u-1),\qquad
\dot C=\frac{16}{3}u(1-u),\qquad
A^2C=A_0^2C_0.
\tag{G07.41}\label{p1b:der:g07:eq-G07-41}
\] Let \(A_*=(A_0^2C_0)^{1/3}\). Then \[
A=A_*u^{-1/3},\qquad C=A_*u^{2/3},\qquad
\dot u=\frac8{A_*}u^{4/3}(1-u).
\tag{G07.42}\label{p1b:der:g07:eq-G07-42}
\] For every \(u_0>0\), this scalar ODE exists for all positive time, remains between \(u_0\) and one, and converges to one. Its linearization has rate \(-8/A_*\). Thus \(A,C\to A_*\) exponentially, and the limiting round radius is \[
R_*=R_0\varepsilon_0^{1/3}.
\] Doing this independently for \(x,p\) gives convergence of the orthogonal product to its separately volume-matched round product, with no requirement that the two limiting radii agree.

The same factorwise conclusion extends beyond homogeneity for smooth sphere-factor metrics sufficiently close in \(C^2\) to their round metrics: their Ricci curvature remains positive, so Hamilton's three-dimensional normalized convergence theorem, as stated in Paper I at p1:thm:hamilton-ricci-flow-convergence, applies to each factor, modulo diffeomorphisms. This statement concerns product perturbations and the declared factorwise flow. It does not assert stability under arbitrary perturbations coupling the factors.

Second, the full eight-dimensional volume-normalized Ricci flow is \(\dot g=-2\mathrm{Ric}+\tfrac14\bar Sg\). For the homogeneous family \(S=\bar S\); put \(\gamma=S/4\), where \(S=\sum_i(8-2u_i)/A_i\). Then \[
\boxed{
\dot A_i=-8+4u_i+\gamma A_i,\qquad
\dot C_i=-4u_i^2+\gamma C_i,\qquad
\dot G=\gamma G.
}
\tag{G07.43}\label{p1b:der:g07:eq-G07-43}
\] The equations for \(u_i\) are still (G07.37), but the scale equations have changed. Differentiating (G07.20) gives \(\partial_t\log\operatorname{Vol}=-S+4\gamma=0\), as required.

Even initially round spheres of equal squared radius \(a_0\) give the explicit solution \[
\boxed{
A_x(t)=C_x(t)=A_p(t)=C_p(t)=a_0-t,\qquad
G(t)=\left(\frac{a_0}{a_0-t}\right)^3G_0,\quad 0\le t<a_0.
}
\tag{G07.44}\label{p1b:der:g07:eq-G07-44}
\] Indeed \(S=12/(a_0-t)\), so (G07.43) reduces to \(\dot a=-1\), \(\dot G=3G/a\). The volume stays fixed, while sphere curvature and torus diameter diverge at the finite endpoint. Thus the full normalized flow does not have a compact round product attractor even in this symmetric subfamily. Paper~I-D, Section~\ref{p1b:app:g07:3} also derives the instability of relative sphere scales and the homogeneous Einstein obstruction. These are essential distinctions when transferring the one-sphere stability statements from Paper I.

\hypertarget{p1b:der:g07:s57}{%
\subsubsection{Metric comparison and stability of scalar and spin spectra}\label{p1b:der:g07:s57}}

Let \(g,h\) be smooth metrics on the fixed compact \(M_8\), with \[
c\,g\le h\le C\,g,\qquad 0<c\le C.
\] Pointwise comparison of path speeds, followed by minimization, gives \[
\sqrt c\,d_g\le d_h\le\sqrt C\,d_g.
\tag{G07.45}\label{p1b:der:g07:eq-G07-45}
\] The identity correspondence consequently bounds the Gromov--Hausdorff distance by \(\tfrac12\sup_{p,q}|d_g(p,q)-d_h(p,q)|\). This proves metric continuity in any uniformly positive parameter region; collapsing families leave such a region.

For the scalar Laplacian, list eigenvalues with multiplicity, starting at \(\lambda_0=0\). The Dirichlet quotient and min--max principle give \[
\frac{c^4}{C^5}\lambda_j(g)
\le\lambda_j(h)\le
\frac{C^4}{c^5}\lambda_j(g).
\tag{G07.46}\label{p1b:der:g07:eq-G07-46}
\] To see the powers, the covector norm lies between \(1/C\) and \(1/c\) times the old norm, while the volume-density ratio lies between \(c^4\) and \(C^4\). Apply these bounds to numerator and denominator on every finite-dimensional trial subspace. For two invariant metrics the density ratio is constant and cancels, improving the estimate to \[
\boxed{
\frac1C\lambda_j(g)\le\lambda_j(h)\le\frac1c\lambda_j(g).
}
\tag{G07.47}\label{p1b:der:g07:eq-G07-47}
\] Thus all ordered scalar eigenvalues vary continuously with the positive invariant metric matrix, including beyond the orthogonal Berger subfamily.

An isolated finite scalar cluster has locally constant total multiplicity under sufficiently small \(C^2\) changes of a smooth metric. Here is the operator argument, to specify the analytic meaning of that statement. Use the square root of the volume-density ratio to identify the two \(L^2\) spaces unitarily. The transported difference \(\Delta_h-\Delta_g\) tends to zero as an operator \(H^2\to L^2\). Elliptic estimates make \((\Delta_g-z)^{-1}:L^2\to H^2\) uniformly bounded on a compact contour avoiding the spectrum. A Neumann series gives uniform resolvent convergence on that contour. Its integral gives norm convergence of the Riesz projections, whose finite ranks are then equal. Individual eigenvalues inside the cluster may split. An arbitrary perturbation need not preserve the old representation labels; for invariant metrics their Casimir degree bands do remain invariant, as Paper~I-D, Section~\ref{p1b:app:g07:4} proves.

Within the Berger family, each fixed label in (G07.32) is differentiable. For its sphere contribution \(\lambda_i=[k_i(k_i+2)+(u_i^{-1}-1)q_i^2]/A_i\), \[
\dot\lambda_i
=-\frac{\dot A_i}{A_i}\lambda_i
 -\frac{\dot u_i\,q_i^2}{A_i u_i^2},\qquad
\frac d{dt}(n^{\mathsf T}G^{-1}n)
=-n^{\mathsf T}G^{-1}\dot G G^{-1}n.
\tag{G07.48}\label{p1b:der:g07:eq-G07-48}
\] At fixed horizontal radius, charged eigenvalues decrease as the fibers become longer, while basic eigenvalues are fixed. Crossings, including (G07.33), explain why an ordered gap need not be differentiable even though each labelled branch is.

For spin stability fix the product orientation and a spin structure on the scalar torus, encoded by \(\sigma\in\{0,\tfrac12\}^2\). Define \[
\mu_\sigma(G)
=\min_{n\in\mathbb Z^2}(n+\sigma)^{\mathsf T}G^{-1}(n+\sigma).
\] The unique spin structures on the two three-spheres and this torus choice determine the product spin structure. At the round product, the exact squared Dirac gap is \[
\boxed{
\min\operatorname{spec}D_{g_0}^2
=\frac94(R_x^{-2}+R_p^{-2})+\mu_\sigma(G).
}
\tag{G07.49}\label{p1b:der:g07:eq-G07-49}
\] This follows from the product square \(D^2=D_x^2+D_p^2+D_T^2\), interpreted on the graded spin tensor product, and the round sphere values \(\min D_i^2=9/(4R_i^2)\). The Clifford anticommutation cancels cross terms; tensor products of lowest factor modes attain the bound.

There is a useful positive gap bound throughout a larger Berger region. For a three-dimensional spin factor, the twistor decomposition and Lichnerowicz identity give \[
\frac23\|D_i\psi\|^2
=\|P_i\psi\|^2+\int\frac{S_i}{4}|\psi|^2,
\] so \(D_i^2\ge(3/8)\min S_i\). Combining this with nonnegativity of the square and (G07.19) yields the quadratic-form bound \[
\boxed{
D^2\ \ge\
\sum_{i=x,p}\max\left\{0,\frac{3(4-u_i)}{4A_i}\right\}
 +\mu_\sigma(G).
}
\tag{G07.50}\label{p1b:der:g07:eq-G07-50}
\] It is sharp at the round product. Whenever the right side is positive, zero modes are excluded. The formula makes explicit which factor or torus term supplies positivity. It is a bound, not a formula for the full Berger spin spectrum, which belongs to the G08 operator atlas.

For general smooth metric changes, use the positive orthonormal-frame identification of the spin bundles and the volume-density unitary. A \(C^1\)-small metric change makes the transported first-order difference \(D_h-D_g\) small from \(H^1\) to \(L^2\). The same contour/Neumann-series argument as above, now using the first-order elliptic estimate, proves norm-resolvent continuity and preservation of isolated finite spin-cluster ranks. It also preserves a nonzero gap for sufficiently small changes. For \(C^2\)-small changes near the round product, persistence of positive scalar curvature and the Lichnerowicz lower bound give a second direct zero-mode exclusion.

The specific first variation of the lowest three-sphere Dirac pair can be stated without a full higher-spin atlas. For a smooth one-factor family through the round metric put \[
\nu_i=\frac{V_i'(0)}{V_i(0)}
=\frac1{2V_i(0)}\int\operatorname{tr}_{g_i}\dot g_i\,dV_i.
\] Each lowest signed eigenvalue is a quaternionic doublet; its isolated rank-two cluster cannot split under a sufficiently small metric change. The one-factor variation formula is \[
\lambda_{i,\pm}(0)=\pm\frac3{2R_i},\qquad
\lambda_{i,\pm}'(0)
=-\frac13\lambda_{i,\pm}(0)\nu_i.
\tag{G07.51}\label{p1b:der:g07:eq-G07-51}
\] For unit radius this is \(\lambda_\pm'(0)=\mp V'(0)/(4\pi^2)\), as proved by \href{https://arxiv.org/abs/1605.08589}{Fang--Levitin--Vassiliev, Theorem 2.1}; constant rescaling gives (G07.51). It agrees with the direct test \(R_i\mapsto R_i(t)\): \(\nu_i=3R_i'/R_i\), while the eigenvalue varies as \(1/R_i\).

For product perturbations and a fixed torus label \(k=n+\sigma\), the squared branch formed from these lowest sphere pairs therefore satisfies \[
\boxed{
(\Lambda^2)'(0)
=-\frac32\left(\frac{\nu_x}{R_x^2}
                   +\frac{\nu_p}{R_p^2}\right)
 -k^{\mathsf T}G^{-1}\dot G G^{-1}k.
}
\tag{G07.52}\label{p1b:der:g07:eq-G07-52}
\] If several torus labels minimize at the reference metric, this is a statement for each branch; the right derivative of the gap is the minimum of their derivatives. There is no global differentiability claim across crossings.

Finally \(M_8\) is even-dimensional, and its chirality operator satisfies \(\Gamma D=-D\Gamma\). Its full spin Dirac spectrum is therefore exactly symmetric under \(\lambda\mapsto-\lambda\), for every metric. The possible second-order signed asymmetry of a single three-dimensional factor must not be transferred to the eight-dimensional product spectrum. G08 will address the full multiplicities and exchange lifts; the present result establishes the metric-stability bridge.

Paper~I-D, Section~\ref{p1b:app:g07} proves the convention-comparison identities and extended-case formulas used in the results above.


\hypertarget{p1b:sec:g08}{%
\subsection{Full bundle spectra, constraints and analytic invariants}\label{p1b:sec:g08}}


\begin{proposition}[Full spin multiplicities and constrained product spectra]
\label{p1b:prop:g08-product}
On the round product with positive radii, torus metric \(G\), and spin
twist \(\sigma\in\{0,\tfrac12\}^2\), the full spin bundle has complex
rank sixteen. For \(k_x,k_p\ge0\), \(n\in\mathbb Z^2\), put
\[
 E^2=\frac{(k_x+3/2)^2}{R_x^2}
       +\frac{(k_p+3/2)^2}{R_p^2}
       +(n+\sigma)^{\mathsf T}G^{-1}(n+\sigma),\qquad
 d_k=(k+1)(k+2).
\]
Each sign of the Dirac eigenvalue \(\pm E\) contributes multiplicity
\(8d_{k_x}d_{k_p}\) for this label triple; coincident values are added.
The full trace-free symmetric tensor fiber has dimension thirty-five,
and its divergence-free restriction is the TT spectrum calculated below.
The intrinsic gamma-traceless bundle has rank \(112\).
For \(Q=\Pi_\gamma D_T|_{\ker\mu}\) and
\(P=\Pi_\gamma\nabla\), \(Q^2+\tfrac12 PP^*\) has scalar positive
principal symbol. By contrast, \(Q^2\) has squared symbol speeds
\(1\) and \(9/16\), with multiplicities \(96\) and \(16\).
Thus its Weyl coefficient is operator-specific.
\end{proposition}

The Clifford product, isotropy branching and explicit constrained matrices below prove the multiplicities and symbols and specify the full Hodge, TT and intrinsic Rarita spectra, heat traces and counting laws.

\hypertarget{p1b:der:g08:s59}{%
\subsubsection{Scope and the homogeneous spectral construction}\label{p1b:der:g08:s59}}

This section determines spin multiplicities and exchange, higher-spin operators and their constraints, heat kernels and coefficients, and Weyl/counting laws. Paper I already supplies the one-sphere scalar, vector, spinor, TT, and Rarita--Schwinger packages at p1:subsec:s3-harmonic-package and p1:subsec:complete-sphere-spectra. Paper I-C already supplies scalar/Hodge product heat traces at p1b:thm:product-hodge, their exchange traces at p1b:thm:exchange-heat, and the values of the product spin square at p1b:thm:metrics. The derivation below uses and extends these mathematical constructions.

The round metric and spin data are \[
M=S^3(R_x)\times S^3(R_p)\times T_G^2,\quad
G=G^{\mathsf T}>0,\quad
\theta\pmod{2\pi\mathbb Z^2},\quad
\sigma\in\{0,\tfrac12\}^2.
\tag{G08.1}\label{p1b:der:g08:eq-G08-1}
\] Write \(a=R_x^{-2}\), \(b=R_p^{-2}\), \(V=\operatorname{Vol}M\), and \(S=6(a+b)\). All bundle ranks and multiplicities below are complex unless explicitly called real. Squared Dirac operators and the scalar/Hodge/Lichnerowicz Laplacians have the positive convention. The Berger extension is the orthogonal family in G07. Mathematical metric parameters are declared data.

There is one common construction underlying all round bundle spectra. Let \[
\alpha=(j_L,j_R),\qquad
d_\alpha=(2j_L+1)(2j_R+1),\qquad
c_\alpha=2[j_L(j_L+1)+j_R(j_R+1)],
\] and define the branching indicator \[
I_s(\alpha)=
\begin{cases}
1,& |j_L-j_R|\le s\le j_L+j_R,\
                    j_L+j_R-s\in\mathbb Z,\\
0,&\text{otherwise}.
\end{cases}
\tag{G08.2}\label{p1b:der:g08:eq-G08-2}
\] Here \(j_L,j_R,s\) are nonnegative half-integers. The diagonal \(\mathrm{Spin}(3)\) restriction of \(V_{j_L}\otimes V_{j_R}\) contains \(V_s\) once exactly under these conditions, by Clebsch--Gordan. Frobenius reciprocity therefore gives a complete orthogonal decomposition of the homogeneous bundle with fiber \(V_s\).

In the invariant frame, expanding the connection square gives the same Casimir calculation as Paper I's spin-two proof, with the fiber Casimir \(6\) replaced by \(s(s+1)\): \[
\nabla^*\nabla\big|_{\alpha,s}
=\frac{c_\alpha-s(s+1)}{R^2}.
\tag{G08.3}\label{p1b:der:g08:eq-G08-3}
\] The standard curvature Laplacian \(\mathscr L=\nabla^*\nabla+q(\mathrm{Rm})\) restores the fiber Casimir on a round factor, so its eigenvalue is \(c_\alpha/R^2\). On ordinary forms this is the Hodge Laplacian; on symmetric two-tensors it is precisely the Lichnerowicz convention used in Paper I. On spinors \(q(\mathrm{Rm})=S/8\), whereas \(D^2=\nabla^*\nabla+S/4=\mathscr L+S/8\). These operator distinctions fix the curvature shifts.

For a product bundle whose isotropy fiber is \[
W=\bigoplus_{s,t}m_{s,t}\,V_s\otimes V_t,
\] the multiplicity of \((\alpha,\beta,n)\) is \[
m_W(\alpha,\beta)=
\sum_{s,t}m_{s,t}I_s(\alpha)I_t(\beta),
\] and the standard curvature Laplacian is scalar on that entire multiplicity space, with value \[
\ell_{\alpha,\beta,n}
=a c_\alpha+b c_\beta+n^{\mathsf T}G^{-1}n.
\tag{G08.4}\label{p1b:der:g08:eq-G08-4}
\] For spin bundles replace \(n\) by \(n+\sigma\). The full labelled multiplicity is \(m_W(\alpha,\beta)d_\alpha d_\beta\). This applies to the round product, including its unequal radii; no Einstein hypothesis is needed for this Casimir construction. It also proves commutation with every isometry-equivariant differential intertwiner between the specified homogeneous bundles.

An explicit kernel is available at the same level. For a one-factor irreducible representation \(\pi_\alpha\), choose diagonal-isotropy intertwiners \(T_r:V_\alpha\to W\) orthonormal in Hilbert--Schmidt norm. In quotient frames the spectral projector kernel is \[
K_\alpha(gH,hH)
=\frac{d_\alpha}{\operatorname{Vol}S^3}
 \sum_r T_r\pi_\alpha(g^{-1}h)T_r^*.
\tag{G08.5}\label{p1b:der:g08:eq-G08-5}
\] Schur orthogonality proves the normalization: the section \(g\mapsto T_r\pi_\alpha(g^{-1})v\) has squared norm \(\operatorname{Vol}S^3\,\|v\|^2/d_\alpha\). Tensor these projectors, multiply by the torus Fourier kernel, and sum the exponentials of the eigenvalues to obtain the corresponding heat kernel. Finite matrix projectors inside a multiplicity space implement additional invariant constraints. This specifies kernels as well as eigenvalues.

\hypertarget{p1b:der:g08:s60}{%
\subsubsection{Round spin multiplicities and the eight-dimensional square}\label{p1b:der:g08:s60}}

For one sphere the condition \(I_{1/2}=1\) forces \[
\alpha_{k,+}=\left(\frac{k+1}{2},\frac k2\right),
\qquad
\alpha_{k,-}=\left(\frac k2,\frac{k+1}{2}\right),\quad k\ge0.
\] Adding the spin curvature shift in (G08.3) gives \[
D_R\big|_{\alpha_{k,\pm}}
=\pm\frac{k+3/2}{R},\qquad
d_k=(k+1)(k+2).
\tag{G08.6}\label{p1b:der:g08:eq-G08-6}
\] The sign assignment fixes a Clifford/orientation convention; reversing it swaps the two displayed branches. The invariant-frame calculation in Section~\ref{p1b:der:g08:s61} also derives these signs and counts directly at the round parameter.

The spinor bundle of an odd-by-odd product requires the additional Clifford factor. In the present case a convenient realization is \[
\Sigma_8\cong\Sigma_x\otimes\Sigma_p\otimes\mathbb C^4,
\qquad \operatorname{rank}_{\mathbb C}\Sigma_8=2\cdot2\cdot4=16.
\] Let \(\alpha_1,\alpha_2,\beta_1,\beta_2\) be four Hermitian anticommuting matrices on \(\mathbb C^4\), each squaring to one. A torus mode has \(k=n+\sigma\); choose orthonormal components \(\kappa=G^{-1/2}k\). Then \[
D_8=
D_x\otimes\alpha_1+D_p\otimes\alpha_2
 +\kappa_1\beta_1+\kappa_2\beta_2,
\qquad
D_8^2=D_x^2+D_p^2+|\,\kappa\,|^2.
\tag{G08.7}\label{p1b:der:g08:eq-G08-7}
\] Identity factors on the sphere spinors are suppressed. The anticommutation removes every cross term.

For fixed \(k_x,k_p,n\), and fixed signs of the two sphere eigenvalues, (G08.7) is a traceless four-by-four Clifford matrix. It has eigenvalues \(\pm E\), each twice, where \[
E^2=a(k_x+3/2)^2+b(k_p+3/2)^2
                      +(n+\sigma)^{\mathsf T}G^{-1}(n+\sigma).
\] There are four independent choices of sphere signs. It follows that \[
\boxed{
\begin{aligned}
\operatorname{mult}_{(k_x,k_p,n)}(+E)
 &=\operatorname{mult}_{(k_x,k_p,n)}(-E)
   =8d_{k_x}d_{k_p},\\
\operatorname{mult}_{(k_x,k_p,n)}(E^2;D_8^2)
 &=16d_{k_x}d_{k_p}.
\end{aligned}}
\tag{G08.8}\label{p1b:der:g08:eq-G08-8}
\] Coincident numerical eigenvalues from different labels add. Each chirality has half of the squared multiplicity. In particular, for the periodic torus spin structure and the lowest sphere levels, the gap is \(\tfrac32\sqrt{a+b}\), of multiplicity \(32\) per sign; the squared eigenspace has dimension \(64\). Other torus spin structures multiply these lowest counts by the number of minimizing shifted lattice vectors.

Define the one-factor spin heat trace and shifted theta sum \[
\mathcal S_R(t)=
2\sum_{k\ge0}(k+1)(k+2)e^{-t(k+3/2)^2/R^2},
\qquad
\Theta_{G,\sigma}(t)=
\sum_{n\in\mathbb Z^2}e^{-t(n+\sigma)^{\mathsf T}G^{-1}(n+\sigma)}.
\] Then the exact full-spin trace is \[
\boxed{
Z_{\rm spin}(t)=
4\mathcal S_{R_x}(t)\mathcal S_{R_p}(t)\Theta_{G,\sigma}(t).
}
\tag{G08.9}\label{p1b:der:g08:eq-G08-9}
\] The factor four records the auxiliary Clifford factor; each \(\mathcal S_R\) already includes both signs. This reproduces (G08.8) without counting either factor twice. The corresponding heat kernel is the product of the two squared-spin kernels, the shifted torus kernel, and the identity on \(\mathbb C^4\).

\hypertarget{p1b:der:g08:s61}{%
\subsubsection{The exact Berger spin spectrum and product zero modes}\label{p1b:der:g08:s61}}

This section completes the deformation bridge left open in G07. Use its right-invariant frame and \(c_a=-i\sigma_a\) on one sphere. If \(J_a\) are the Hermitian spin-\(j\) angular-momentum matrices, the Koszul spin connection gives \[
\boxed{
D_{R,\varepsilon}\big|_j
=\frac2R\left(\sigma_1J_1+\sigma_2J_2+
                         \varepsilon^{-1}\sigma_3J_3\right)
 +\frac{\varepsilon+2/\varepsilon}{2R}I.
}
\tag{G08.10}\label{p1b:der:g08:eq-G08-10}
\] Each such matrix has the spectator multiplicity \(2j+1\). To derive the constant, insert the G07 connection coefficients in \(\frac14\sum_{a,b,c}\Gamma_{ab}^{\,c}c_ac_bc_c\). The result is \((\varepsilon+2/\varepsilon)/(2R)\); the differentiating part uses \(\xi_a=2iJ_a\).

The paired states \((|\uparrow,m\rangle,|\downarrow,m+1\rangle)\) give a two-by-two matrix with mean diagonal \(\varepsilon/(2R)\), half-difference \((2m+1)/(R\varepsilon)\), and off-diagonal entry \(2\sqrt{(j-m)(j+m+1)}/R\). Writing \(p=j-m\), \(q=j+m+1\), all paired eigenvalues are \[
\lambda_{p,q}^{\pm}
=\frac1R\left[
\frac{\varepsilon}{2}\pm
\sqrt{4pq+\frac{(p-q)^2}{\varepsilon^2}}\right],
\quad p,q\in\mathbb Z_{>0},
\quad\operatorname{mult}=p+q.
\tag{G08.11}\label{p1b:der:g08:eq-G08-11}
\] The two unpaired highest/lowest states give \[
\lambda_r^{\rm u}
=\frac1R\left(\frac{\varepsilon}{2}+\frac r\varepsilon\right),
\quad r\in\mathbb Z_{>0},\quad\operatorname{mult}=2r.
\tag{G08.12}\label{p1b:der:g08:eq-G08-12}
\] Ordered pairs \((p,q)\) label separate subspaces; equal values, including \(p,q\) interchanged, have additive multiplicities. At \(\varepsilon=1\), grouping the pairs by \(p+q\) and adding the unpaired branch recovers (G08.6) exactly. These formulas follow by diagonalizing all invariant blocks, so no extension of a small-deformation theorem is being assumed.

Only the minus paired branch can vanish. Its zero condition is \[
\varepsilon^2
=8pq+2\sqrt{16p^2q^2+(p-q)^2}.
\tag{G08.13}\label{p1b:der:g08:eq-G08-13}
\] It implies \(\varepsilon\ge4\sqrt{pq}\ge4\). The first zero is therefore exactly at \(\varepsilon=4\), \(p=q=1\), and the one-factor kernel has dimension two. At any fixed positive parameter only finitely many pairs can satisfy (G08.13). Let \(h(\varepsilon)\) be the sum of their multiplicities.

For the orthogonal Berger product, (G08.7) still holds with these signed one-factor eigenvalues. For one chosen eigenspace of dimensions \(d_x,d_p\) and one torus mode, the nonzero product values are \(\pm\sqrt{\lambda_x^2+\lambda_p^2+|\kappa|^2}\), each of multiplicity \(2d_xd_p\). At zero the four-dimensional auxiliary matrix is zero. Since a sum of squares vanishes exactly when all three terms vanish, \[
\boxed{
\dim\ker D_8=
\begin{cases}
4h(\varepsilon_x)h(\varepsilon_p),&\sigma=(0,0),\\
0,&\sigma\ne(0,0).
\end{cases}}
\tag{G08.14}\label{p1b:der:g08:eq-G08-14}
\] In particular the first simultaneous Berger crossing \(\varepsilon_x=\varepsilon_p=4\) has a 16-dimensional product kernel, split equally between chiralities. Neither the round gap nor the G07 positive-curvature gap bound was claimed through these distant deformations.

\hypertarget{p1b:der:g08:s62}{%
\subsubsection{The geometric spin lift of exchange}\label{p1b:der:g08:s62}}

Let \(B(q_x,q_p,\theta_t,\theta_E) =(q_p,q_x,\theta_E,\theta_t)\). Assume equal sphere metrics and \(S_2^{\mathsf T}GS_2=G\), where \(S_2\) swaps the scalar coordinates. A spin structure is preserved exactly when \(\sigma_t=\sigma_E\). For the other two spin structures, exchange maps between different spin bundles, so it does not define even/odd sectors inside one of them.

The sphere swap has determinant \((-1)^{3\cdot3}=-1\); the torus swap also has determinant \(-1\). Thus the full eight-dimensional exchange preserves orientation. In matched invariant orthonormal frames it reverses four normal directions: \[
v_i=(e_{x,i}-e_{p,i})/\sqrt2,\quad 1\le i\le3,
\qquad v_4=e_{T,-}.
\] Here \(e_{T,-}\) is the unit anti-diagonal torus vector. The two spin lifts differ by sign and are represented by \[
\widetilde B=\pm c(v_1)c(v_2)c(v_3)c(v_4)
\quad\text{together with pullback of the base variables}.
\] Because four orthogonal Clifford generators anticommute, \[
\widetilde B^{\,2}=(-1)^{4\cdot5/2}=1.
\tag{G08.15}\label{p1b:der:g08:eq-G08-15}
\] This also constructs the global lift in the product trivialization; the condition on \(\sigma\) makes its torus transition signs compatible.

There is an explicit realization matching (G08.7). Choose the torus orthonormal directions as diagonal and anti-diagonal, with auxiliary matrices \(\beta_+,\beta_-\). If \(F\) flips the two sphere-spinor Hilbert factors, then, for one of the two lifts, \[
\widetilde B=F\otimes U,\qquad
U=\frac{i}{\sqrt2}(\alpha_1-\alpha_2)\beta_-,
\quad U^2=1.
\tag{G08.16}\label{p1b:der:g08:eq-G08-16}
\] Include torus pullback in this formula. Conjugation by \(U\) interchanges \(\alpha_1,\alpha_2\), fixes \(\beta_+\), and negates \(\beta_-\). Multiplying the four geometric Clifford normals above gives exactly this operator in the chosen frame. It commutes with \(D_8\) and with chirality. This identifies the smooth spin lift; it is separate from G03's finite protected-module flip.

Both \(U\) and its product with auxiliary chirality have zero trace: each is a nontrivial Clifford monomial. The squared heat operator is scalar on the auxiliary factor. The flip trace identity therefore yields \[
\boxed{
\operatorname{Tr}(\widetilde B e^{-tD_8^2})=0,\qquad
\operatorname{Tr}(\Gamma\widetilde B e^{-tD_8^2})=0.
}
\tag{G08.17}\label{p1b:der:g08:eq-G08-17}
\] This is exact for all \(t>0\), for matched round or matched Berger factors and an exchange-preserved spin structure. Hence the two exchange heat traces are \(Z_{\rm spin}/2\), and each joint chirality/exchange trace is \(Z_{\rm spin}/4\). Uniqueness of the discrete Laplace expansion gives the same equality on every squared eigenspace. Since chirality preserves exchange and reverses the nonzero signed eigenvalue, each signed eigenspace also splits equally into exchange parities. Changing the lift sign merely interchanges the parity names. In particular the first simultaneous Berger kernel in Section~\ref{p1b:der:g08:s61} has four joint chirality/exchange sectors of dimension four.

\hypertarget{p1b:der:g08:s63}{%
\subsubsection{All Hodge degrees, exact/coexact constraints, and kernels}\label{p1b:der:g08:s63}}

Write \(A_R,C_R,\Theta_G\) as in the supplied p1b:thm:product-hodge: \[
\begin{aligned}
A_R(t)&=\sum_{k\ge0}(k+1)^2e^{-tk(k+2)/R^2},\\
C_R(t)&=A_R(t)-1+
       2\sum_{k\ge1}k(k+2)e^{-t(k+1)^2/R^2},\\
\Theta_G(t)&=\sum_{n\in\mathbb Z^2}e^{-tn^{\mathsf T}G^{-1}n}.
\end{aligned}
\] The factor connection and exterior-algebra grading give \[
\boxed{
\mathcal Z(z,t)=
\prod_{i=x,p}\bigl[A_{R_i}(t)(1+z^3)+C_{R_i}(t)(z+z^2)\bigr]
  (1+z)^2\Theta_G(t).
}
\tag{G08.18}\label{p1b:der:g08:eq-G08-18}
\] The coefficient of \(z^p\) is the full \(p\)-form trace. The proof is a tensoring of complete eigenbases: the scalar and volume-form branches coincide; the exact one-form branch is the differential of nonconstant scalars; the coexact branch has values \((k+1)^2/R^2\) and total multiplicity \(2k(k+2)\); Hodge star gives degree two. Eigenvalues add on the product, and the two torus covectors give \((1+z)^2\). All sums converge absolutely for \(t>0\).

This construction also determines the global exact and coexact spectra, which are not obtained by declaring every factor component separately transverse. For \(\lambda>0\), let \(m_p(\lambda)\) be the coefficient multiplicity from (G08.18), and let \(e_p(\lambda)\) be the multiplicity on coexact \(p\)-forms. Hodge decomposition and the isomorphism \(d\) between positive coexact \((p-1)\)-forms and exact \(p\)-forms give \[
e_{-1}(\lambda)=0,\qquad
e_p(\lambda)=m_p(\lambda)-e_{p-1}(\lambda)
 =\sum_{j=0}^p(-1)^{p-j}m_j(\lambda).
\tag{G08.19}\label{p1b:der:g08:eq-G08-19}
\] On the positive \(\lambda\) eigenspace the two orthogonal projections are \[
\Pi_{\rm ex}=\lambda^{-1}dd^*,\qquad
\Pi_{\rm coex}=\lambda^{-1}d^*d.
\tag{G08.20}\label{p1b:der:g08:eq-G08-20}
\] Their sum is the identity. These formulas preserve coincident eigenvalues and implement the full product constraint. Applying them to the spectral kernels in (G08.5) gives the exact/coexact heat kernels.

At zero, (G08.18) gives the Betti polynomial \[
(1+z)^2(1+z^3)^2,
\] of total dimension sixteen. The de Rham supertrace is zero at every positive time. On positive eigenspaces, \((d+d^*)/\sqrt\lambda\) explicitly pairs even and odd form degrees; at zero the alternating Betti sum is zero.

For scalars the zonal projector makes the kernel especially simple: \[
K_{0,R}(t;q,q')
=\frac1{2\pi^2R^3}
 \sum_{k\ge0}(k+1)
 \frac{\sin((k+1)\theta)}{\sin\theta}
 e^{-tk(k+2)/R^2},
\quad
\cos\theta=\operatorname{Re}(q'\bar q).
\tag{G08.21}\label{p1b:der:g08:eq-G08-21}
\] Use the continuous endpoint values at \(\theta=0,\pi\). The torus factor is \[
K_{T,G}(t;\theta,\theta')
=\frac1{(2\pi)^2\sqrt{\det G}}
\sum_{n\in\mathbb Z^2}
 e^{-tn^{\mathsf T}G^{-1}n}e^{in\cdot(\theta-\theta')}.
\] The product of these three kernels is the scalar kernel on \(M\). Orthogonality proves the semigroup law and integral-one normalization directly.

For equal factors the source's ordinary-form exchange trace is recovered as \[
\mathcal Z_B(z,t)=
[A_R(2t)(1-z^6)+C_R(2t)(-z^2+z^4)]
   (1-z^2)\Theta_{\rm diag}(t).
\tag{G08.22}\label{p1b:der:g08:eq-G08-22}
\] Here \(\Theta_{\rm diag}\) sums over the fixed torus labels \((n,n)\). The signs are the Koszul sign on equal sphere-form degrees and \(\det(I+zS_2)=1-z^2\) on torus covectors. The parity traces are \((\mathcal Z\pm\mathcal Z_B)/2\). Equations (G08.19)--(G08.20) also apply within each exchange sector because pullback commutes with \(d,d^*\). The exact zero spin-exchange trace (G08.17) and this generally nonzero form-exchange trace concern different bundles.

\hypertarget{p1b:der:g08:s64}{%
\subsubsection{The complete product TT Lichnerowicz spectrum}\label{p1b:der:g08:s64}}

Define TT to mean both trace-free and divergence-free for the full eight-dimensional metric. Use \[
(\Delta_Lh)_{ab}
=-\nabla^c\nabla_ch_{ab}
 +R_a{}^c h_{cb}+R_b{}^c h_{ac}
 -2R_{acbd}h^{cd}.
\] The trace-free symmetric fiber decomposes under the two sphere isotropy groups as \[
\operatorname{Sym}^2_0T^*M
\cong (2,0)\oplus(0,2)\oplus(1,1)
       \oplus2(1,0)\oplus2(0,1)\oplus4(0,0).
\tag{G08.23}\label{p1b:der:g08:eq-G08-23}
\] The dimensions are \(5+5+9+6+6+4=35\). The four scalar components are the two sphere trace directions and the three torus symmetric directions, subject to one total trace condition. The one-form fiber is \((1,0)\oplus(0,1)\oplus2(0,0)\). The curvature term on the first two summands of (G08.23) is \(6a\), respectively \(6b\); on a cross vector summand it is the sum of the appropriate Ricci values. In each case it cancels the subtracted fiber Casimir in (G08.3). Thus \(\Delta_L\) has the common value \(\ell_{\alpha,\beta,n}\) of (G08.4).

Put \(x_s=I_s(\alpha)\), \(p_s=I_s(\beta)\). The unconstrained trace-free and one-form multiplicities are \[
\begin{aligned}
m_{\rm tf}
 &=x_2p_0+x_0p_2+x_1p_1
    +2x_1p_0+2x_0p_1+4x_0p_0,\\
m_1&=x_1p_0+x_0p_1+2x_0p_0.
\end{aligned}
\] The formal adjoint of divergence is the trace-free symmetrized derivative. Its kernel is the conformal Killing space. On this compact product every conformal Killing field is Killing: for each fixed sphere point, its torus components satisfy the conformal Killing equation on a flat two-torus. Taking a torus divergence shows that these components are harmonic and hence constant in the torus variables; their divergence is zero, so the common conformal factor is zero. The Killing algebra is the sum of the two sphere rotation algebras and the two torus translations. Its dimension is fourteen.

Let \(m_K\) be one on each of the four labels \[
((1,0),(0,0),0),\quad((0,1),(0,0),0),\quad
((0,0),(1,0),0),\quad((0,0),(0,1),0),
\] two on \(((0,0),(0,0),0)\), and zero elsewhere. Since the operators intertwine the complete homogeneous decomposition, rank--nullity and adjointness give \[
\boxed{
m_{\rm TT}
=m_{\rm tf}-m_1+m_K
=x_2p_0+x_0p_2+x_1p_1+x_1p_0+x_0p_1
  +2x_0p_0+m_K.
}
\tag{G08.24}\label{p1b:der:g08:eq-G08-24}
\] The full TT spectrum consists of every value \(\ell_{\alpha,\beta,n}\) with multiplicity \(m_{\rm TT}d_\alpha d_\beta\). This is an explicit all-level formula, including mixed sphere/torus components and the exceptional Killing labels. It is more than the factorwise TT subspace.

Only the trivial homogeneous label has \(\ell=0\). At that label (G08.24) gives four. Consequently \[
\boxed{
\dim\ker(\Delta_L|_{\rm TT})=4,\qquad
\lambda_{\rm TT,1}^{+}
=\min\{3a,3b,\lambda_{T,1}\}.
}
\tag{G08.25}\label{p1b:der:g08:eq-G08-25}
\] The zero modes are precisely parallel trace-free combinations of \(g_x,g_p\), and symmetric torus tensors. They are not automatically infinitesimal Einstein deformations: the full product is not Einstein. The first sphere scalar label, for example, contributes TT multiplicity \(3\cdot4=12\) at \(3a\); a nonzero torus mode with constant sphere labels contributes two. All other allowed positive Casimirs are at least the displayed minimum, which proves the gap formula.

For an exact heat trace, let \[
T_R(t)=2\sum_{L\ge2}(L-1)(L+3)e^{-t[L(L+2)+4]/R^2},
\] the sphere TT trace. The trace-free symmetric sphere trace is \[
B_R(t)=T_R(t)+C_R(t)
              -6e^{-4t/R^2}-4e^{-3t/R^2}.
\] Indeed the trace-free symmetrized gradient on the sphere has the six Killing fields and four first scalar gradients as its kernel; subtracting those shared one-form labels leaves its image. Then (G08.24) gives \[
\boxed{
\begin{aligned}
Z_{\rm TT}(t)
={}&[B_xA_p+A_xB_p+C_xC_p+C_xA_p+A_xC_p+2A_xA_p]
        \Theta_G(t)\\
 &+6e^{-4at}+6e^{-4bt}+2.
\end{aligned}}
\tag{G08.26}\label{p1b:der:g08:eq-G08-26}
\] All factor functions in the first line are evaluated at \(t\). The final finite term restores the Killing kernel in rank--nullity; it must not be discarded when counting zero modes or computing constant heat terms.

\hypertarget{p1b:der:g08:s65}{%
\subsubsection{Rarita--Schwinger operators and the full product constraint}\label{p1b:der:g08:s65}}

Let \(c(v)^2=-|v|^2\) and \[
\mu:T^*M\otimes\Sigma\longrightarrow\Sigma,\qquad
\mu(\psi)=\sum_{\alpha=1}^8c(e_\alpha)\psi_\alpha.
\] Then \(\mu\mu^*=8I\), so \[
\Pi_\gamma=I-\frac18\mu^*\mu,\qquad
\Sigma_{3/2}=\ker\mu,\qquad
\operatorname{rank}_{\mathbb C}\Sigma_{3/2}=112.
\tag{G08.27}\label{p1b:der:g08:eq-G08-27}
\] This is the intrinsic eight-dimensional gamma-traceless bundle. Its isotropy decomposition is \[
4(V_{3/2}\otimes V_{1/2})
\oplus4(V_{1/2}\otimes V_{3/2})
\oplus12(V_{1/2}\otimes V_{1/2}).
\tag{G08.28}\label{p1b:der:g08:eq-G08-28}
\] To obtain it, tensor the spin fiber \(4(V_{1/2}\otimes V_{1/2})\) with \((V_1\otimes V_0)\oplus(V_0\otimes V_1)\oplus2(V_0\otimes V_0)\), apply Clebsch--Gordan, and remove the image of \(\mu^*\). The three dimensions are \(32,32,48\).

Let \(D_T\) be the spin Dirac operator twisted by the tangent bundle, with its Levi--Civita connection. Define \[
Q=\Pi_\gamma D_T|_{\Sigma_{3/2}},\qquad
P=\Pi_\gamma\nabla,\qquad
P^*\psi=-\sum_\alpha(\nabla_{e_\alpha}\psi)_\alpha.
\tag{G08.29}\label{p1b:der:g08:eq-G08-29}
\] Thus \(Q\) is Rarita--Schwinger, \(P\) is twistor, and transversality means \(P^*\psi=0\). The normalization is exactly the gamma projection, with no additional higher-spin scale.

For comparison, on one round three-sphere the fiber is \(V_{3/2}\). Its two branches have spin differences \(3/2\) and \(1/2\). The Einstein Weitzenböck identities give \(Q^2=\mathscr L-5/(4R^2)\) on the transverse branch and \(QP=\frac13PD\) on the twistor image. Inserting their Casimirs yields \[
\begin{array}{c|c|c}
\text{sphere sector}&Q\text{ values}&\text{multiplicity per sign}\\ \hline
\ker P^*&\displaystyle\pm(k+\tfrac52)/R&(k+1)(k+4)\\
\operatorname{Im}P&\displaystyle\pm(k+\tfrac52)/(3R)&(k+2)(k+3)
\end{array}
\qquad k\ge0.
\tag{G08.30}\label{p1b:der:g08:eq-G08-30}
\] The lowest spinor doublets are in \(\ker P\); this explains why the image series starts at the next spinor level. The classical identities and sphere spectrum are stated in \href{https://arxiv.org/abs/2001.06167}{Homma--Tomihisa, Theorems 1--2} and in the supplied p1:thm:rarita-schwinger-spectrum. Equation (G08.30) fixes the distinct full and transverse gaps \(5/(6R)\) and \(5/(2R)\).

For a general \(n\)-dimensional metric, the corresponding commutator retains a trace-free Ricci term: \[
QP-\frac{n-2}{n}PD
=\frac12F,\qquad
(F\varphi)_\alpha
=c\!\left((\mathrm{Ric}-S/n)e_\alpha\right)\varphi.
\tag{G08.31}\label{p1b:der:g08:eq-G08-31}
\] This follows by expanding \(P_\alpha=\nabla_\alpha+ c(e_\alpha)D/n\) in \(Q\), commuting covariant derivatives, and contracting spin curvature. The convention is also recorded in \href{https://arxiv.org/abs/1804.10602}{Homma--Semmelmann, equation (4)}. For the present round product the Ricci eigenvalues are \(2a\) three times, \(2b\) three times, and zero twice. Therefore \(F\ne0\), even when \(a=b\). The sphere's invariant transverse/image split cannot simply be assigned to this intrinsic product operator. The nonzero trace-free Ricci term supplies the stated finite constant-frame coupling between the two spaces.

The twistor projection itself is nevertheless completely determined. The twistor norm identity gives \[
P^*P=\frac78D_8^2-\frac S4
\ \ge\ \frac{15}{32}(a+b)+\frac78\mu_\sigma(G)>0.
\tag{G08.32}\label{p1b:der:g08:eq-G08-32}
\] Hence the product has no twistor spinors, and \[
\Pi_{\operatorname{Im}P}
=P(P^*P)^{-1}P^*,\qquad
\Pi_{\ker P^*}=\Pi_\gamma-\Pi_{\operatorname{Im}P}.
\] These are well-defined orthogonal projections on the appropriate \(L^2\) closures. Failure of \(Q\) to commute with them is an operator property, not missing constraint data.

Here is an exact finite-matrix description of the entire intrinsic \(Q\) spectrum, valid also for the orthogonal Berger family. It avoids an Einstein substitution. In a global invariant orthonormal frame let \(B_\alpha\) be the eight-by-eight tangent connection matrix and \[
\Omega_\alpha^\Sigma
=\frac14\sum_{\beta,\gamma}
 \langle\nabla_{e_\alpha}e_\beta,e_\gamma\rangle
 c(e_\beta)c(e_\gamma).
\] For orbital spins \(j_x,j_p\) use the usual angular matrices \(J_i^{(j_x)},J_i^{(j_p)}\). On that coefficient block the differentiating operators are \[
E_{x,1,2}=2iJ_{1,2}^{(j_x)}/R_x,\quad
E_{x,3}=2iJ_3^{(j_x)}/(R_x\varepsilon_x),
\] with the analogous \(p\) formula, and \(E_{T,\alpha}=i\kappa_\alpha\). The round case sets \(\varepsilon_i=1\). All \(B_\alpha\) are supplied explicitly by the Koszul coefficients in G07. Then form \[
\begin{aligned}
\mathbb D_T(j_x,j_p,n)
&=\sum_{\alpha=1}^8
(I_T\otimes c(e_\alpha))
\bigl[E_\alpha+
       B_\alpha\otimes I_\Sigma+
       I_T\otimes\Omega_\alpha^\Sigma\bigr],\\
\mathbb Q(j_x,j_p,n)
&=Z_\gamma^*\mathbb D_T(j_x,j_p,n)Z_\gamma .
\end{aligned}
\tag{G08.33}\label{p1b:der:g08:eq-G08-33}
\] Identity factors on the orbital spaces are suppressed. Here \(Z_\gamma\) is any orthonormal basis matrix for \(\ker\mu\), tensored with the orbital identity. Thus \(\mathbb Q\) has dimension \(112(2j_x+1)(2j_p+1)\); its eigenvalues are repeated with spectator multiplicity \((2j_x+1)(2j_p+1)\). Take every nonnegative half-integer \(j_x,j_p\) and every \(n\in\mathbb Z^2\). This is an exact spectral specification by explicit Hermitian matrices; it does not leave an unspecified differential problem. Peter--Weyl completeness proves that no eigensections are omitted. The same blocks give \(P^*\) by contraction, so the simultaneous equations \(Q\psi=\lambda\psi\), \(P^*\psi=0\) are exact finite linear systems on each block. A non-invariant transverse space is not silently treated as the domain of a restricted self-adjoint \(Q\).

The coupled positive operator has a simpler geometric meaning: \[
\boxed{
\mathcal A_{\rm RS}
:=
Q^2+\frac12PP^*
=\Pi_\gamma D_T^2|_{\Sigma_{3/2}}.
}
\tag{G08.34}\label{p1b:der:g08:eq-G08-34}
\] The coefficient is \(4/n\) at \(n=8\), obtained by squaring the block decomposition of \(D_T\) into its gamma-trace and gamma-traceless parts. It is a Laplace-type operator; its quadratic form is \(\|Q\psi\|^2+\frac12\|P^*\psi\|^2=\|D_T\psi\|^2\).

On the round product, the tangent twist splits into the two sphere tangent bundles and the flat torus tangent bundle. On a twist in the \(x\) directions the square is \((D_x^{T_x})^2+D_p^2+D_T^{\,2,\mathrm{flat}}\); the \(p\) and torus twists have the corresponding factor sums. Nonnegativity of each twisted square and the other sphere's ordinary Dirac gap give \[
\boxed{
\mathcal A_{\rm RS}\ge
\frac94\min(a,b)+\mu_\sigma(G)>0.
}
\tag{G08.35}\label{p1b:der:g08:eq-G08-35}
\] It follows that the simultaneous equations \(Q\psi=0\), \(P^*\psi=0\) have no nonzero solution. This proves absence of constrained Rarita--Schwinger zero fields on the round product. It is not a claim that gamma-traceless \(Q\) alone has this same gap.

For matched metrics, exchange on \(T^*M\otimes\Sigma\) is the tensor product of the tangent action and the spin lift. Naturality of contraction and connection shows that it commutes with \(Q,\mathcal A_{\rm RS}\) and transports \(P,P^*\). Thus each matrix orbit in (G08.33) can be projected with \((I\pm\widetilde B)/2\). The trace and gamma constraints for higher tensor or tensor-spinor bundles can likewise be formed in the same coefficient spaces: tensor the connection matrices, apply their algebraic orthogonal projectors, and form the chosen natural differential operator. This supplies the higher-rank construction without identifying distinct higher-spin operators by name.

\hypertarget{p1b:der:g08:s66}{%
\subsubsection{Exact round heat profiles and their coefficients}\label{p1b:der:g08:s66}}

Fix the integrated-coefficient convention \[
\operatorname{Tr}e^{-tL}
\sim(4\pi t)^{-4}\sum_{j\ge0}A_{2j}(L)t^j.
\tag{G08.36}\label{p1b:der:g08:eq-G08-36}
\] These \(A_{2j}\) include the spatial integral. The round factor sums determine them directly to every order, so their normalization can be checked without only importing a general heat-coefficient table.

Put \(h=t/R^2\), \(S_0(h)=\sum_{m\ge1}e^{-hm^2}\), and \(S_2(h)=\sum_{m\ge1}m^2e^{-hm^2}=-S_0'(h)\). Poisson summation gives \[
S_0(h)=\frac{\sqrt\pi}{2\sqrt h}-\frac12
             +\text{exponentially small terms}.
\] The exact identities \[
\begin{aligned}
A_R&=e^hS_2,\\
C_R&=e^hS_2-1+2(S_2-S_0),\\
T_R&=2e^{-3h}(S_2-4S_0+3e^{-h})
\end{aligned}
\] follow by replacing the respective levels with \(m=k+1\). For spinors use the half-integer theta sum: \[
\mathcal S_R=
\sum_{m\in\mathbb Z+1/2}(m^2-\tfrac14)e^{-hm^2}.
\] The omitted \(m=\pm1/2\) terms are identically zero. Applying Poisson summation and differentiation therefore gives the following complete local profiles: \[
\begin{array}{c|c}
\text{trace}&\text{profile multiplying }
                    \operatorname{Vol}(S_R^3)(4\pi t)^{-3/2}\\ \hline
A_R&e^h\\
C_R&e^h+2-4h\\
B_R&2e^{-3h}(1-8h)+e^h+2-4h\\
\mathcal S_R&2-h
\end{array}
\tag{G08.37}\label{p1b:der:g08:eq-G08-37}
\] The difference in each row is exponentially small as \(t\downarrow0\), for fixed positive radius. In the \(B_R\) row the finite terms \(4e^{-3h}+6e^{-4h}\) in \(T_R\) cancel against the conformal-Killing subtraction.

The marked torus, including its spin shifts, has \[
\Theta_{G,\sigma}(t)
=\frac{\pi\sqrt{\det G}}t
\sum_{m\in\mathbb Z^2}
 e^{-\pi^2m^{\mathsf T}Gm/t}e^{2\pi i m\cdot\sigma}.
\tag{G08.38}\label{p1b:der:g08:eq-G08-38}
\] Its only local coefficient is the volume term. Thus all four torus spin structures have the same local heat coefficients, although their exact spectra differ.

For the scalar product, multiplication of (G08.37) and (G08.38) yields \[
Z_0(t)\sim\frac{V}{(4\pi t)^4}e^{(a+b)t},
\qquad
\boxed{A_{2j}(\Delta_0)=V\frac{(a+b)^j}{j!}\quad(j\ge0).}
\tag{G08.39}\label{p1b:der:g08:eq-G08-39}
\] In particular \(A_2=VS/6\) and \(A_4=V(a+b)^2/2\), in agreement with the G07 curvature invariants.

For all Hodge degrees define \[
f(z,h)=e^h(1+z^3)+(e^h+2-4h)(z+z^2).
\] The all-order generating function of the coefficients is \[
\sum_{p=0}^8 z^p\sum_{j\ge0}A_{2j}(\Delta_p)t^j
=V(1+z)^2 f(z,at)f(z,bt).
\tag{G08.40}\label{p1b:der:g08:eq-G08-40}
\] The right side here is its Taylor series at zero. With binomial coefficients outside their usual range defined to be zero, coefficient extraction gives \[
\boxed{
\begin{aligned}
\frac{A_0(\Delta_p)}V&=\binom8p,\\
\frac{A_2(\Delta_p)}V
 &=(a+b)\left[\binom8p-6\binom6{p-1}\right],\\
\frac{A_4(\Delta_p)}V
 &=\frac{a^2+b^2}{2}
       \left[\binom8p-2\binom6{p-1}\right]\\
 &\quad+ab\left[\binom8p-12\binom6{p-1}
                          +36\binom4{p-2}\right].
\end{aligned}}
\tag{G08.41}\label{p1b:der:g08:eq-G08-41}
\] This includes degrees zero through eight, their Hodge duals, and the vanishing alternating sums.

For the full spin square the local series terminates: \[
\boxed{
Z_{\rm spin}(t)\sim
\frac{16V}{(4\pi t)^4}
\left[1-\frac{a+b}{2}t+\frac{ab}{4}t^2\right].
}
\tag{G08.42}\label{p1b:der:g08:eq-G08-42}
\] Thus \(A_0=16V\), \(A_2=-8V(a+b)\), \(A_4=4Vab\), and every \(A_{2j}\) for \(j\ge3\) vanishes on this round product. This is an exact statement about local asymptotic coefficients, not about the full heat trace at finite time. The exponentially small terms still contain the global lengths and spin structure.

For TT tensors, substitute the three corresponding profiles from (G08.37) into (G08.26). The first three coefficients are \[
\boxed{
\begin{aligned}
A_0^{\rm TT}&=27V,\\
A_2^{\rm TT}&=-27V(a+b),\\
A_4^{\rm TT}&=V\left[\frac{129}{2}(a^2+b^2)-45ab\right].
\end{aligned}}
\tag{G08.43}\label{p1b:der:g08:eq-G08-43}
\] The exact finite Killing term in (G08.26) starts contributing at the \(t^0\) term of the heat trace, equivalently at \(A_8^{\rm TT}\) in convention (G08.36). Thus the all-order TT expansion is obtained by retaining that term as well as the local factor profiles. The projection onto TT is a global differential constraint; dropping its finite-dimensional correction would give an incorrect constant term.

\hypertarget{p1b:der:g08:s67}{%
\subsubsection{Curvature checks, Berger coefficients, and exchange asymptotics}\label{p1b:der:g08:s67}}

For a Laplace-type operator written \(L=\nabla^*\nabla+E\) on a rank-\(r\) bundle over a closed manifold, the integrated first coefficients are \[
\begin{aligned}
A_0&=\int r,\\
A_2&=\int\operatorname{tr}\left(\frac S6I-E\right),\\
A_4&=\frac1{360}\int\operatorname{tr}\left[
(5S^2-2|\mathrm{Ric}|^2+2|\mathrm{Rm}|^2)I
-60SE+180E^2+30\Omega_{\alpha\beta}\Omega_{\alpha\beta}
\right].
\end{aligned}
\tag{G08.44}\label{p1b:der:g08:eq-G08-44}
\] Integrated divergence terms vanish. Here \(\Omega\) is the skew-Hermitian connection curvature and both ordered curvature indices are summed. This sign convention for \(E\) is opposite to the potential used in \href{https://arxiv.org/abs/hep-th/0306138}{Vassilevich, §4.1}. The formulas apply to the Hodge and spin squares and to \(\mathcal A_{\rm RS}\), but not directly to the non-Laplace-type \(Q^2\).

For the spin square \(E=S/4\), Clifford trace gives \(\operatorname{tr}\Omega_{\alpha\beta} \Omega_{\alpha\beta}=-r|\mathrm{Rm}|^2/8\). Consequently \[
\frac{A_2(D^2)}{rV}=-\frac S{12},\qquad
\frac{A_4(D^2)}{rV}
=\frac{5S^2-8|\mathrm{Ric}|^2-7|\mathrm{Rm}|^2}{1440}
\tag{G08.45}\label{p1b:der:g08:eq-G08-45}
\] for the homogeneous family. For the round product \(|\mathrm{Ric}|^2=|\mathrm{Rm}|^2=12(a^2+b^2)\); the pure \(a^2,b^2\) terms cancel, leaving \(ab/4\). This independently checks (G08.42).

For the Berger product put \(A_i=R_i^2\), \(u_i=\varepsilon_i^2\), and use its actual volume \(V_B\). The G07 curvature calculation gives \[
\begin{aligned}
S&=\sum_i(8-2u_i)/A_i,\\
|\mathrm{Ric}|^2
 &=\sum_i(32-32u_i+12u_i^2)/A_i^2,\\
|\mathrm{Rm}|^2
 &=\sum_i(64-96u_i+44u_i^2)/A_i^2.
\end{aligned}
\tag{G08.46}\label{p1b:der:g08:eq-G08-46}
\] Thus its scalar coefficients are \[
A_0=V_B,\quad A_2=V_BS/6,\quad
A_4=\frac{V_B}{360}(5S^2-2|\mathrm{Ric}|^2+2|\mathrm{Rm}|^2),
\] and its full-spin coefficients are (G08.45) with \(r=16\). No smallness or positive-curvature condition is needed for these fixed positive Berger metrics. Uniformity through a collapsing or infinitely stretched family is a separate limit and is not asserted.

Exchange heat asymptotics also follow from the exact formulas, with no need to infer them from ambient dimension. For an invariant torus metric \[
G=\begin{pmatrix}g&c\\c&g\end{pmatrix},\quad g>|c|,
\] the scalar fixed-label sum is \(\Theta_{\rm diag}(t)= \sum_{n\in\mathbb Z}e^{-2tn^2/(g+c)}\). For round factors its exchange trace is \[
Z_{0,B}(t)=A_R(2t)\Theta_{\rm diag}(t)
\sim \frac{\pi R^3\sqrt{g+c}}{16t^2}e^{2t/R^2}.
\tag{G08.47}\label{p1b:der:g08:eq-G08-47}
\] The fixed set is \(S^3_{\rm diag}\times S^1_{\rm diag}\) with induced volume \(16\pi^3R^3\sqrt{g+c}\). Four normal directions are reversed, so the leading coefficient also equals \(\operatorname{Vol}(\operatorname{Fix}B)/ [16(4\pi t)^2]\), as a direct Gaussian normal integration confirms. This checks both the power \(t^{-2}\) and the factor sixteen. The two scalar parity sectors each have half the ordinary \(t^{-4}\) leading term, with the signed \(t^{-2}\) correction supplied by (G08.47). Hodge parity coefficients follow from (G08.22). For spinors the entire exchange trace vanishes by (G08.17), not merely its first local coefficient.

\hypertarget{p1b:der:g08:s68}{%
\subsubsection{Exact counting, Weyl laws, and the control-heat limit}\label{p1b:der:g08:s68}}

Every atlas above gives an exact counting function by summing its labelled multiplicity over the inequality defining the spectral window. For example, \[
N_{D_8^2}(\Lambda)
=16\sum_{k_x,k_p\ge0}\sum_{n\in\mathbb Z^2}
d_{k_x}d_{k_p}\,\mathbf1_{\{E_{k_x,k_p,n}^2\le\Lambda\}}.
\tag{G08.48}\label{p1b:der:g08:eq-G08-48}
\] The scalar count replaces \(16d_{k_x}d_{k_p}\) by \((k_x+1)^2(k_p+1)^2\) and uses the scalar values. For TT use (G08.24), and for \(Q^2\) use the matrix spectral projectors of (G08.33) with their spectator multiplicities. These sums are finite for fixed \(\Lambda\). They retain the marked torus lattice and all accidental degeneracies.

Set \[
W_8=\frac{V}{(4\pi)^4\Gamma(5)}
    =\frac{V}{6144\pi^4}.
\] The positive heat leading term, followed by the Tauberian theorem for nondecreasing spectral counts, gives \[
\begin{array}{c|c}
L&N_L(\Lambda)\sim C_LW_8\Lambda^4\\ \hline
\Delta_0& C_L=1\\
\Delta_p& C_L=\binom8p\\
D_8^2& C_L=16\\
\Delta_L|_{\rm TT}& C_L=27\\
\mathcal A_{\rm RS}& C_L=112\\
Q^2& C_L=96+16(4/3)^8
\end{array}
\tag{G08.49}\label{p1b:der:g08:eq-G08-49}
\] Finite zero-mode counts do not affect these leading terms. For signed Dirac values, \(N_{D_8,+}(\Lambda)\sim8W_8\Lambda^8\), whereas the absolute-value count has coefficient \(16W_8\). Exchange halves the leading scalar/form counts and exactly halves each squared spin count.

For completeness, the last coefficient in (G08.49) requires its principal symbol, rather than the bundle rank. At a unit covector \(\xi\), the gamma-traceless fiber splits into the image of the twistor symbol (rank 16) and its orthogonal complement (rank 96). Clifford multiplication and the projection (G08.27) give the two squared speeds \[
\sigma(Q)(\xi)^2=
\begin{cases}
\frac9{16}I,&\text{on the rank-16 symbol image},\\
I,&\text{on the rank-96 complement}.
\end{cases}
\tag{G08.50}\label{p1b:der:g08:eq-G08-50}
\] One can verify this by taking \(\xi=e_8\): the projected symbol obeys \((\sigma(Q)^2-I)(\sigma(Q)^2-\frac9{16}I)=0\); its trace is \(96+16\cdot9/16\). Orthogonal covariance gives every \(\xi\). Integrating the two phase-space balls gives the factor \((3/4)^{-8}\) on the smaller-speed subspace. This also proves ellipticity of \(Q\). The term \(\frac12PP^*\) in (G08.34) restores unit principal speed on that subspace, explaining the rank 112 coefficient for \(\mathcal A_{\rm RS}\).

The G05 control operator has a different exact heat asymptotic. On one factor set \(h=t/R^2\) and write the scalar degree and charge as \(k=r+s,\ q=r-s\), where \(r,s\ge0\). Then \[
k(k+2)-q^2=(2r+1)(2s+1)-1,\qquad k+1=r+s+1.
\] Summing one variable, and using symmetry in \(r,s\), gives the exact trace \[
\boxed{
H_R(t):=\operatorname{Tr}e^{-tL_H}
=e^h\sum_{r\ge0}\frac{2r+1}{2\sinh(h(2r+1))}
\sim\frac{\pi^2R^4}{16t^2}.
}
\tag{G08.51}\label{p1b:der:g08:eq-G08-51}
\] For the last step, the midpoint Riemann sum uses \[
\int_0^\infty\frac{x}{2\sinh x}\,dx
=\sum_{m\ge0}\frac1{(2m+1)^2}
=\frac{\pi^2}{8}.
\] The mesh is \(2h\), and the additional prefactor is \(1/h\), giving \(\pi^2/(16h^2)\). The exponential tails justify the sum and the integral.

On \(M_8\), with the torus directions retained, \[
\boxed{
\begin{aligned}
\operatorname{Tr}e^{-t(L_{H,x}+L_{H,p}+\Delta_T)}
&\sim
\frac{\pi^5R_x^4R_p^4\sqrt{\det G}}{256\,t^5},\\
N_H(\Lambda)
&\sim
\frac{\pi^5R_x^4R_p^4\sqrt{\det G}}{30720}\Lambda^5.
\end{aligned}}
\tag{G08.52}\label{p1b:der:g08:eq-G08-52}
\] This derives the homogeneous spectral dimension ten already suggested by G05's tangent geometry. It does not alter the manifold's topological dimension. At each fixed \(t>0\), the Berger eigenvalue formulas and monotone convergence show that long-fiber heat traces tend to this trace. For shrinking fibers, dominated convergence removes every charged mode and leaves the quotient heat trace. Thus the order of the metric limit and small-time limit is material:

\begin{longtable}[]{@{}ll@{}}
\toprule\noalign{}
Parameters before taking \(t\downarrow0\) & Heat power \\
\midrule\noalign{}
\endhead
\bottomrule\noalign{}
\endlastfoot
Fixed positive Berger parameters & \(t^{-4}\) \\
Both fibers collapsed first & \(t^{-3}\) \\
One fiber collapsed first & \(t^{-7/2}\) \\
Both vertical costs sent to infinity first & \(t^{-5}\) \\
One vertical cost sent to infinity first & \(t^{-9/2}\) \\
\end{longtable}

\hypertarget{p1b:der:g08:s69}{%
\subsubsection{Finite levels, zeta functions, and analytic exports}\label{p1b:der:g08:s69}}

An invariant finite spectral band is obtained by retaining a finite set of labels and, where required, a subspace of the finite multiplicity matrix. On scalar, Hodge, spin, and TT bands made from full eigenspaces, all spectral values and kernels restrict exactly. Matrix (G08.33) also preserves every finite orbital coefficient band. An independently imposed total-isometry representation cutoff is implemented by the corresponding equivariant projector. These are different legitimate truncations; their finite dimensions must be computed on the stated space.

If exhaustive spectral cutoffs \(P_N\) are used for a nonnegative operator with compact resolvent, then \[
\|(L+z)^{-1}-(L+z)^{-1}P_N\|
\le\frac1{z+\lambda_{\rm first\ omitted}},
\quad z>0,
\tag{G08.53}\label{p1b:der:g08:eq-G08-53}
\] and the analogous heat-operator norm error is at most \(e^{-t\lambda_{\rm first\ omitted}}\). Both follow by diagonalizing the omitted spectrum. Positive heat traces converge monotonically to the parent trace. For a fixed finite band, however, the small-time trace tends to its dimension. Taking the parent limit first is required for the Weyl powers in (G08.49).

The enriched carrier completion in Paper I-C is not the scalar sphere \(L^2\) completion. If \(K(a,b)=[a(a+2)+b(b+2)]/4\) on its multiplicity-one sectors of dimension \((a+1)(b+1)\), then \[
\operatorname{Tr}e^{-t\gamma K_\infty}
=\left[\sum_{a\ge0}(a+1)e^{-t\gamma a(a+2)/4}\right]^2
\sim\frac4{\gamma^2t^2}.
\tag{G08.54}\label{p1b:der:g08:eq-G08-54}
\] The last step is a one-dimensional Gaussian integral on each factor. In contrast a scalar sphere has heat growth \(\sqrt\pi R^3/(4t^{3/2})\). With a finite scalar-label set \(F\) and momentum carrier \(H_1\) of dimension fifteen, the corresponding position-parent leading terms are respectively \[
\frac{60|F|}{\gamma_x^2t^2},\qquad
\frac{15|F|\sqrt\pi R_x^3}{4t^{3/2}}.
\] With both enriched factors exhaustive the coefficient is \(16|F|/(\gamma_x^2\gamma_p^2t^4)\); with both scalar factors exhaustive it is \(|F|\pi R_x^3R_p^3/(16t^3)\). An unrestricted torus multiplies these by \(\pi\sqrt{\det G}/t\) in place of \(|F|\). This replays the supplied p1b:thm:all-level-limits and keeps its carrier types explicit. An equal power of \(t\) in two constructions does not identify their operators or spaces.

For any of the nonnegative elliptic squares above let \(h_L=\dim\ker L\). Its positive spectral zeta function is \[
\zeta_L(s)=\frac1{\Gamma(s)}
\int_0^\infty t^{s-1}
\bigl[\operatorname{Tr}e^{-tL}-h_L\bigr]\,dt,
\qquad \operatorname{Re}s>4.
\tag{G08.55}\label{p1b:der:g08:eq-G08-55}
\] More generally, use the convergence half-plane dictated by the operator's Weyl law; for the control square it is \(\operatorname{Re}s>5\). Subtracting finitely many small-time terms in (G08.55) continues the function meromorphically wherever the corresponding heat expansion is known. In particular for the scalar round product, \[
\zeta_{\Delta_0}(0)
=\frac{V(a+b)^4}{24(4\pi)^4}-1.
\tag{G08.56}\label{p1b:der:g08:eq-G08-56}
\] For the invertible round full-spin square, the terminating local series (G08.42) gives \[
\boxed{\zeta_{D_8^2}(0)=0.}
\tag{G08.57}\label{p1b:der:g08:eq-G08-57}
\] Moreover the full spin Dirac eta function vanishes: chirality pairs every nonzero \(\lambda\) with \(-\lambda\) at equal multiplicity, first in the absolutely convergent half-plane and then by continuation. This pairing remains true in the Berger family, including when zero modes must be treated separately. It is distinct from a one-sphere eta invariant or from a curl helicity.

The scalar and Hodge Green operators are obtained by integrating their heat kernels after removing harmonic projections. The explicit positive gaps prove convergence at large time; elliptic estimates give the usual Sobolev gains and smoothness away from the diagonal. On positive Hodge modes, \[
d\,\Delta^{-1}d^*+d^*\,\Delta^{-1}d=I-\Pi_{\rm harm},
\] so the atlas also supplies a global contracting homotopy for the nonharmonic de Rham complex. Paper~I-D, Section~\ref{p1b:app:g08} records the product overlap rules, TT exchange trace, and the conformal/Yamabe second variation needed to distinguish sphere conclusions from their product versions.

Paper~I-D, Section~\ref{p1b:app:g08} proves the convention-comparison identities and extended-case formulas used in the results above.


\hypertarget{p1b:sec:g09}{%
\subsection{Reproducing kernels, evaluation measures and exact moments}\label{p1b:sec:g09}}


\begin{proposition}[Product evaluation and exact observable moments]
\label{p1b:prop:g09-product}
Let \(W=\mathcal P_{m_x}\otimes\mathcal P_{m_p}\otimes\mathbb C[F]\),
where \(m_x,m_p\ge0\), \(F\subset\mathbb Z^2\) is nonempty and finite,
and Haar measure is probability measure. Set
\(D_m=(m+1)(m+2)(2m+3)/6\) and \(N=D_{m_x}D_{m_p}|F|\).
The scalar reproducing kernel factorizes, has constant diagonal \(N\),
and its evaluation vectors satisfy
\[
 \int_M |k_z\rangle\langle k_z|\,d\mu(z)=I_W,\qquad
 P_WM_fP_W=\int_M f(z)|k_z\rangle\langle k_z|\,d\mu(z).
\]
For sphere designs of strengths \(t_x,t_p\), polynomial degrees
\(d_x,d_p\) in \(f\), exact sampled moments follow from
\(2m_i+d_i\le t_i\) together with the stated torus no-alias condition.
The enriched form/spin realization instead uses fiber evaluation:
at level one its density has rank nine and trace fifteen, and its
two-factor density has rank eighty-one and trace \(225\).
Its kernel is not a rank-one scalar evaluation kernel.
\end{proposition}

The addition formulas, sphere moments and quadrature identities below prove these assertions and derive the exact boundary corrections, Berezin spectrum and controlled finite-level limits.

\hypertarget{p1b:der:g09:s71}{%
\subsubsection{Scope, scalar bands, and measure}\label{p1b:der:g09:s71}}

This derivation supplies the product version of Paper I's \path{p1:thm:scalar-band-evaluation-povm} and reconciles its exact moments with Paper I-C's finite observable compiler. The source statements used are:

\begin{longtable}[]{@{}
  >{\raggedright\arraybackslash}p{(\columnwidth - 2\tabcolsep) * \real{0.5000}}
  >{\raggedright\arraybackslash}p{(\columnwidth - 2\tabcolsep) * \real{0.5000}}@{}}
\toprule\noalign{}
\begin{minipage}[b]{\linewidth}\raggedright
Source locator
\end{minipage} & \begin{minipage}[b]{\linewidth}\raggedright
Input used here
\end{minipage} \\
\midrule\noalign{}
\endhead
\bottomrule\noalign{}
\endlastfoot
Paper I, \path{p1:thm:scalar-band-evaluation-povm} & Scalar evaluation vectors, probability Haar measure, normalized POVM and first moment \\
Paper I, \path{p1:thm:v180-exceptional-scalar-moment-ranges} & Exceptional strengths \(5,7,11\) and sphere moments \\
Paper I, \path{p1:thm:v203-scalar-observable-compiler} and \path{p1:cor:v203-quadratic-top-block-correction} & The \(2m+\deg f\) condition and degree-two correction \\
Paper I, \path{p1:prop:v203-quadratic-alias-witness} & The specific \(1/48\) witness \\
Paper I-C, result~\ref{p1b:thm:compiler} and result~\ref{p1b:thm:product-cubature} & Product degree/frequency conditions and covariance \\
Paper I-C, result~\ref{p1b:thm:full-sampling} and result~\ref{p1b:thm:all-level-sampling} & Fibre-valued evaluation and all-level realizations \\
Paper I-C, result~\ref{p1b:cor:all-level-design-comparison} & A fixed second factor cannot repair a first-factor moment failure \\
This paper, Section~\ref{p1b:der:g02:s2}--\ref{p1b:der:g02:s8}, \ref{p1b:der:g04:s26}, \ref{p1b:der:g06:s43}, \ref{p1b:der:g08:s59} and \ref{p1b:der:g08:s66} & Typed enriched model, scalar torus, distinct Hardy construction, and round projector kernels \\
\end{longtable}

In particular, scalar bands, the enriched fifteen-dimensional carrier, and the Hardy spaces of G06 remain different Hilbert spaces.

Let \[
M=S^3_{R_x}\times S^3_{R_p}\times\mathbb T^2,\qquad
W=\mathcal P_{m_x}^{x}\otimes\mathcal P_{m_p}^{p}
       \otimes\mathbb C[F],
\quad
\mathbb C[F]=\operatorname{span}\{e^{in\cdot\theta}:n\in F\},
\tag{G09.1}\label{p1b:der:g09:eq-G09-1}
\] where \(R_x,R_p>0\), \(m_x,m_p\in\mathbb Z_{\ge0}\), and \(\varnothing\ne F\subset\mathbb Z^2\) is finite. Here \(\mathcal P_m=\bigoplus_{k=0}^m E_k^{\rm sc}\), with \(\dim E_k^{\rm sc}=(k+1)^2\). Both circle coordinates have period \(2\pi\). Use \[
d\mu=d\mu_x\,d\mu_p\,\frac{d^2\theta}{(2\pi)^2},
\qquad
D_m=\frac{(m+1)(m+2)(2m+3)}6,\qquad
N=\dim W=D_{m_x}D_{m_p}|F|.
\tag{G09.2}\label{p1b:der:g09:eq-G09-2}
\] The sphere measures are probability Haar measures. If a flat torus metric \(G>0\) is specified, its volume is \((2\pi)^2\sqrt{\det G}\); dividing product Riemannian volume by its total volume gives exactly \(d\mu\). Thus no torus metric or dimensionful scale is needed to define this scalar POVM.

The inner product is conjugate-linear in its first argument. Write \(z=(q_x,q_p,\theta)\), with \(q_i\) unit quaternions and radius-scaled embedding coordinates \(x_a=R_xq_{x,a}\), \(p_a=R_pq_{p,a}\). Let \(P\) be the orthogonal projection from \(L^2(M,d\mu)\) onto \(W\).

\hypertarget{p1b:der:g09:s72}{%
\subsubsection{The explicit product reproducing kernel}\label{p1b:der:g09:s72}}

The round degree-\(k\) scalar projector from G08 is \[
\Pi_k(q,q')=(k+1)U_k(q\cdot q'),\qquad
U_k(\cos\vartheta)=\frac{\sin((k+1)\vartheta)}{\sin\vartheta}.
\tag{G09.3}\label{p1b:der:g09:eq-G09-3}
\] Endpoint values are understood by continuity. The coefficient is \(k+1\), whereas the diagonal and the eigenspace dimension are \((k+1)^2\). The addition theorem follows by summing a Haar orthonormal harmonic basis; the invariant zonal polynomial has degree \(k\), and its diagonal is fixed by the trace.

Consequently the complete-band kernel is \[
C_m(s)=\sum_{k=0}^m(k+1)U_k(s)
=\frac{(m+2)U_m(s)-(m+1)U_{m+1}(s)}{2(1-s)},
\quad s\ne1.
\tag{G09.4}\label{p1b:der:g09:eq-G09-4}
\] The second expression follows by telescoping \(2sU_k=U_{k+1}+U_{k-1}\), with \(U_{-1}=0\). At the endpoints, \(C_m(1)=D_m\) and \(C_m(-1)=(-1)^m(m+1)(m+2)/2\). The finite sum, rather than a quotient with a removable singularity, defines the kernel everywhere.

The scalar torus kernel is \(H_F(\theta,\phi)=\sum_{n\in F}e^{in\cdot(\theta-\phi)}\). Tensoring the three orthonormal bases gives \[
\boxed{
K(z,w)=C_{m_x}(q_x\cdot q'_x)\,
       C_{m_p}(q_p\cdot q'_p)\,
       H_F(\theta,\phi).
}
\tag{G09.5}\label{p1b:der:g09:eq-G09-5}
\] For \(K_z(w)=K(w,z)\), one has \[
\langle K_z,u\rangle=u(z),\qquad
K(z,z)=\|K_z\|^2=N,\qquad
\int_M K(z,w)K(w,v)\,d\mu(w)=K(z,v).
\tag{G09.6}\label{p1b:der:g09:eq-G09-6}
\] These equations follow directly from orthonormality, including for unequal sphere cutoffs and nonrectangular \(F\).

With unnormalized product volume \(dV=V\,d\mu\), an orthonormal basis is divided by \(\sqrt V\), and its reproducing kernel is \(K/V\). The two realizations are unitarily equivalent by \(u\mapsto u/\sqrt V\). Factors of \(V\) must be changed in both the inner product and the kernel; they are not extra POVM weights.

\hypertarget{p1b:der:g09:s73}{%
\subsubsection{Normalization, probabilities, and dilation}\label{p1b:der:g09:s73}}

Set \(\kappa_z=K_z/\sqrt N\). For every Borel subset \(A\subset M\), define \[
\boxed{
E(A)=\int_A |K_z\rangle\langle K_z|\,d\mu(z)
=N\int_A|\kappa_z\rangle\langle\kappa_z|\,d\mu(z).
}
\tag{G09.7}\label{p1b:der:g09:eq-G09-7}
\] The density has rank one and trace \(N\). Positivity and countable additivity are immediate in the finite-dimensional weak integral. For \(u,v\in W\), \[
\langle u,E(M)v\rangle
=\int_M\overline{u(z)}v(z)\,d\mu(z)
=\langle u,v\rangle,
\qquad
\boxed{E(M)=I_W}.
\tag{G09.8}\label{p1b:der:g09:eq-G09-8}
\] Thus dividing the density in (G09.7) by \(N\) would destroy normalization. The unit vectors \(\kappa_z\) have measure weight \(N\).

For a density matrix \(\rho\) on \(W\), the outcome density and the normalized trace are \[
p_\rho(z)=\langle K_z,\rho K_z\rangle,\qquad
\int p_\rho\,d\mu=1,\qquad
p_{I/N}(z)=1,\qquad
\frac{\operatorname{tr}E(A)}N=\mu(A).
\tag{G09.9}\label{p1b:der:g09:eq-G09-9}
\] For a pure state \(u\), \(p_{|u\rangle\langle u|}=|u(z)|^2\). The uniform outcome density of the maximally mixed state is a consequence of the constant kernel diagonal, not a choice of state for the other formulas.

There is an explicit dilation. Let \(J:W\hookrightarrow L^2(M,d\mu)\) be inclusion and \(F_{\rm mult}(A)=M_{1_A}\) the multiplication projection. Then \[
E(A)=J^*F_{\rm mult}(A)J,\qquad J^*J=I_W.
\tag{G09.10}\label{p1b:der:g09:eq-G09-10}
\] This proves the dilation identity without an existence assumption. It is minimal: choose any \(n_0\in F\); the retained function \(e^{in_0\cdot\theta}\), constant on the spheres, is nowhere zero, and its multiples by indicator functions span a dense subspace of \(L^2(M,d\mu)\). The compressed measure is generally not projection-valued.

\hypertarget{p1b:der:g09:s74}{%
\subsubsection{Every scalar moment and the product factorization}\label{p1b:der:g09:s74}}

For every bounded measurable scalar symbol \(f\), its exact operator-valued moment is \[
\boxed{
T_f:=\int_M f(z)E(dz)=PM_fP|_W,\qquad
(T_f)_{\alpha\beta}
=\int_M\overline{\phi_\alpha(z)}f(z)\phi_\beta(z)\,d\mu(z).
}
\tag{G09.11}\label{p1b:der:g09:eq-G09-11}
\] In particular \(T_1=I\), \(T_{\bar f}=T_f^*\), \(f\ge0\) implies \(T_f\ge0\), and \(\|T_f\|\le\|f\|_\infty\). For a positive matrix-valued symbol, the same compression is positive on the corresponding tensor product, so the moment map is completely positive as well as unital.

For a separated symbol, Fubini's theorem gives \[
T_{f_xf_ph}
=T^{x}_{f_x}\otimes T^{p}_{f_p}\otimes T^{F}_{h}.
\tag{G09.12}\label{p1b:der:g09:eq-G09-12}
\] Linearity extends this to finite sums of separated symbols. Every polynomial/Fourier moment in the retained coordinates is thereby specified exactly, even when quadrature cannot evaluate it.

The underlying sphere integrals have a closed formula. If \(a_i\ge0\) and \(A=\sum_i a_i\), then \[
\int_{S^3}\prod_{i=0}^3q_i^{2a_i}\,d\mu_H
=\frac{\prod_i(2a_i)!}{4^A(A+1)!\prod_i a_i!}.
\tag{G09.13}\label{p1b:der:g09:eq-G09-13}
\] Any odd coordinate exponent makes the integral zero. For a derivation, write a standard Gaussian in \(\mathbb R^4\) as \(rq\), with \(r\) independent of the uniform direction. Divide the product of the one-dimensional Gaussian even moments by \(\mathbb E r^{2A}=2^A(A+1)!\). Equivalently the result is \(\prod_i(2a_i-1)!!/\prod_{j=0}^{A-1}(4+2j)\). No design assumption enters this continuum identity.

For multi-indices \(\alpha,\beta\) and \(\nu\in\mathbb Z^2\), \[
T_{x^\alpha p^\beta e^{i\nu\cdot\theta}}
=R_x^{|\alpha|}R_p^{|\beta|}
 T^x_{q^\alpha}\otimes T^p_{q^\beta}\otimes C_\nu,
\qquad
\frac{\operatorname{tr}T_f}{N}=\int_M f\,d\mu.
\tag{G09.14}\label{p1b:der:g09:eq-G09-14}
\] The second identity follows from \(K(z,z)=N\). It gives the classical Haar moments only for the normalized trace state; it does not identify the whole operator \(T_f\) with a scalar.

For example, \[
T_{x_ap_b}=T_{x_a}T_{p_b}=T_{p_b}T_{x_a},
\quad
\sum_aT_{x_a^2}=R_x^2I,\quad
\frac{\operatorname{tr}T_{x_ax_b}}N
=\frac{R_x^2}{4}\delta_{ab}.
\tag{G09.15}\label{p1b:der:g09:eq-G09-15}
\] The first equality uses different tensor factors. Within a single sphere, \(T_{x_ax_b}\) generally differs from \(T_{x_a}T_{x_b}\).

In the Fourier basis \(|n\rangle=e^{in\cdot\theta}\), \[
C_\nu:=T^F_{e^{i\nu\cdot\theta}},\qquad
C_\nu|n\rangle=
\begin{cases}
|n+\nu\rangle,&n+\nu\in F,\\
0,&n+\nu\notin F.
\end{cases}
\qquad C_\nu^*=C_{-\nu}.
\tag{G09.16}\label{p1b:der:g09:eq-G09-16}
\] This also covers \(0\notin F\), disconnected frequency sets, and cutoffs coupling the two circle labels.

\hypertarget{p1b:der:g09:s75}{%
\subsubsection{Coordinate boundary terms and exact second moments}\label{p1b:der:g09:s75}}

For a homogeneous harmonic polynomial \(h_k\) on \(\mathbb R^4\), \[
q_a h_k
=\left(q_a h_k-\frac{r^2}{2(k+1)}\partial_a h_k\right)
 +\frac{r^2}{2(k+1)}\partial_a h_k,
\qquad r^2=\sum_b q_b^2.
\tag{G09.17}\label{p1b:der:g09:eq-G09-17}
\] The two terms restrict to degrees \(k+1\) and \(k-1\) on the unit sphere. Indeed, \(\Delta(q_ah_k)=2\partial_ah_k\) and \(\Delta(r^2\partial_ah_k)=4(k+1)\partial_ah_k\). The lowering term vanishes at \(k=0\).

Write \(X_a=P_mM_{q_a}P_m\) and \(P_m^{\rm top}=\Pi_m^{\rm sc}\) for the degree-\(m\) homogeneous projector, rather than the whole-band projector. If \(A_a:E_m^{\rm sc}\to E_{m+1}^{\rm sc}\) is the raising term, then \[
\sum_aA_a^*A_a=\frac{m+2}{2(m+1)}I_{E_m^{\rm sc}}.
\tag{G09.18}\label{p1b:der:g09:eq-G09-18}
\] To determine the coefficient, the squared norm of the lowering part of \(h_k\) is \[
\frac1{4(k+1)^2}\int\sum_a|\partial_ah_k|^2\,d\mu_H
=\frac{k}{2(k+1)}\|h_k\|^2.
\] Here the ambient gradient splits into tangential and radial parts, whose integrated squared norms are \(k(k+2)\|h_k\|^2\) and \(k^2\|h_k\|^2\). Multiplication by the four coordinates has total squared norm \(\|h_k\|^2\); subtraction gives (G09.18).

The only multiplication channel leaving \(\mathcal P_m\) is \(E_m^{\rm sc}\to E_{m+1}^{\rm sc}\). Therefore \[
T_{q_aq_b}-X_aX_b
=P_mM_{q_a}(I-P_m)M_{q_b}P_m
=P_m^{\rm top}A_a^*A_bP_m^{\rm top}.
\tag{G09.19}\label{p1b:der:g09:eq-G09-19}
\] The \(4\)-by-\(4\) matrix of these operator entries is positive. In particular the diagonal differences are positive, and \[
\boxed{
\sum_{a=0}^3X_a^2
=I-\frac{m+2}{2(m+1)}P_m^{\rm top},
\qquad
\frac1{D_m}\operatorname{tr}
 \left(I-\sum_aX_a^2\right)=\frac3{2m+3}.
}
\tag{G09.20}\label{p1b:der:g09:eq-G09-20}
\] At \(m=0\), every \(X_a\) vanishes and the correction is \(I\). At \(m=1\), the constant-to-linear matrix element is \(1/2\). For radius-scaled embedding coordinates, multiply \(X_a\) by \(R\) and every second-moment or commutator identity by \(R^2\).

The antisymmetric part of (G09.19), or direct application of (G09.17), gives \[
[X_a,X_b]
=\frac1{2(m+1)}
P_m^{\rm top}(q_a\partial_b-q_b\partial_a)P_m^{\rm top}.
\tag{G09.21}\label{p1b:der:g09:eq-G09-21}
\] The angular differential operator is anti-Hermitian. For modes below the top degree the two multiplication paths agree, so the commutator vanishes there. On the top degree the lowering-then- raising path gives the displayed coefficient.

Likewise the exact torus second moment is \[
T^F_{|e^{i\nu\cdot\theta}|^2}=I,\qquad
I-C_\nu^*C_\nu
=\sum_{\substack{n\in F\\n+\nu\notin F}}
 |n\rangle\langle n|.
\tag{G09.22}\label{p1b:der:g09:eq-G09-22}
\] These formulas describe the finite compression boundary. They are unrelated to the protected Clifford relations of G02--G04.

\hypertarget{p1b:der:g09:s76}{%
\subsubsection{Exact Berezin transform and the observable symbol space}\label{p1b:der:g09:s76}}

The normalized coherent expectation of \(T_f\) is \[
\mathcal Bf(z)
:=\langle\kappa_z,T_f\kappa_z\rangle
=\frac1N\int_M|K(z,w)|^2f(w)\,d\mu(w).
\tag{G09.23}\label{p1b:der:g09:eq-G09-23}
\] This is a positive, self-adjoint, finite-rank Markov operator on \(L^2(M,d\mu)\). Both row and column integrals of its kernel are one, so it is a contraction.

The elementary character identity \(U_j(s)U_k(s)=\sum_{r=0}^{\min(j,k)}U_{j+k-2r}(s)\) follows by multiplying their finite exponential sums. It is also the \(\lambda=1\) specialization of the \href{https://dlmf.nist.gov/18.18.E22}{ultraspherical linearization formula, DLMF 18.18.22}. It gives \[
C_m(s)^2=\sum_{\ell=0}^{2m}c_{m,\ell}U_\ell(s),\qquad
c_{m,\ell}=
\sum_{\substack{0\le j,k\le m\\
 |j-k|\le\ell\le j+k\\ j+k-\ell\ {\rm even}}}
(j+1)(k+1).
\tag{G09.24}\label{p1b:der:g09:eq-G09-24}
\] Comparing with the projector normalization (G09.3), the one-sphere Berezin eigenvalues are exactly \[
b_{m,\ell}=\frac{c_{m,\ell}}{D_m(\ell+1)}
\quad(0\le\ell\le2m),\qquad b_{m,\ell}=0\quad(\ell>2m).
\tag{G09.25}\label{p1b:der:g09:eq-G09-25}
\] All displayed nonzero eigenvalues are positive. For \(\ell\le m\), take \((j,k)=(\ell,0)\); for \(m<\ell\le2m\), take \((j,k)=(m,\ell-m)\) in (G09.24). The contraction property gives their upper bound \(1\).

For the torus, set \(d_F(\nu)=\#\{(n,k)\in F^2:n-k=\nu\}\). Expansion of \(|H_F|^2\) gives \[
\boxed{
\mathcal B|_{E_{\ell_x}^x\otimes E_{\ell_p}^p\otimes\mathbb C e^{i\nu\cdot\theta}}
=b_{m_x,\ell_x}\,b_{m_p,\ell_p}
  \frac{d_F(\nu)}{|F|}\,I .
}
\tag{G09.26}\label{p1b:der:g09:eq-G09-26}
\] Thus the exact space of observable scalar symbols is \[
\mathcal V=
\mathcal P_{2m_x}^{x}\otimes\mathcal P_{2m_p}^{p}
 \otimes\mathbb C[F-F],
\qquad
\boxed{\operatorname{rank}(f\mapsto T_f)
=D_{2m_x}D_{2m_p}|F-F|.}
\tag{G09.27}\label{p1b:der:g09:eq-G09-27}
\] For a proof of the rank statement, with Hilbert--Schmidt inner product on operators the adjoint of \(T:f\mapsto T_f\) is \(T^*A(z)=\langle K_z,AK_z\rangle\). Hence \(\mathcal B=T^*T/N\). Its strictly positive eigenvalues are exactly those on \(\mathcal V\), and it vanishes on \(\mathcal V^\perp\). Therefore \(T\) is injective on \(\mathcal V\) and zero on \(\mathcal V^\perp\). In particular, (G09.27) specifies every moment without requiring an infinite symbol table.

Useful low-degree values are \[
b_{m,0}=1,\qquad
b_{m,1}=\frac{2m}{2m+3}\quad(m\ge0),\qquad
(b_{2,0},\ldots,b_{2,4})
=\left(1,\frac47,\frac{19}{42},\frac3{14},\frac9{70}\right).
\tag{G09.28}\label{p1b:der:g09:eq-G09-28}
\] At \(m=0\), the \(\ell=1\) expression denotes the zero eigenvalue. For \(\ell=1\), only \(j=k\pm1\) contribute to (G09.24); summing \(2\sum_{k=0}^{m-1}(k+1)(k+2)\) proves the formula.

This POVM resolves \(I_W\) and all moments in its symbol image. It does not determine every density matrix when \(N>1\). Paper~I-D, Section~\ref{p1b:app:g09:2} gives the exact ambiguity and reconstruction of the component that it does determine. This scope distinction also prevents identification with the different Hardy/Berezin quantization in G06.

\hypertarget{p1b:der:g09:s77}{%
\subsubsection{Discrete POVMs, exact cubature, and frequency aliases}\label{p1b:der:g09:s77}}

For nodes \(z_a\) and positive probability weights \(w_a\), define weighted scalar evaluation \(Su=(\sqrt{w_a}u(z_a))_a\). Its candidate discrete effects and Gram operator are \[
E_a=w_a|K_{z_a}\rangle\langle K_{z_a}|,\qquad
\sum_aE_a=S^*S=:\mathsf G.
\tag{G09.29}\label{p1b:der:g09:eq-G09-29}
\] They are a normalized POVM precisely when \(S\) is isometric. For sphere designs of strengths \(t_x,t_p\), and a rectangular uniform torus grid of lengths \(L_t,L_E\), sufficient and explicit conditions are \[
2m_x\le t_x,\qquad2m_p\le t_p,\qquad
(F-F)\cap\Lambda_L=\{0\},
\quad
\Lambda_L=L_t\mathbb Z\times L_E\mathbb Z.
\tag{G09.30}\label{p1b:der:g09:eq-G09-30}
\] The frequency condition is necessary and sufficient for isometric evaluation of \(\mathbb C[F]\) on this grid: the discrete inner product of \(e^{in\cdot\theta}\) and \(e^{ik\cdot\theta}\) is \(1\) exactly when \(n-k\in\Lambda_L\). Here zero belongs to \(F-F\); the condition forbids only its nonzero lattice points.

For a Fourier symbol \(e^{i\nu\cdot\theta}\), the exact matrix entry comparison is \[
(C_\nu^{\rm grid})_{n,k}
=1_{\{k+\nu-n\in\Lambda_L\}},\qquad
(C_\nu)_{n,k}=1_{\{k+\nu-n=0\}}.
\tag{G09.31}\label{p1b:der:g09:eq-G09-31}
\] This is a sharper condition than a rectangular bandwidth estimate. For all symbols supported in \(A\subset\mathbb Z^2\), simultaneous Fourier compiler exactness is equivalent to \((A+F-F)\cap(\Lambda_L\setminus\{0\})=\varnothing\).

If a polynomial/Fourier symbol has sphere degrees bounded by \(d_x,d_p\) and Fourier support \(A\), then \[
\boxed{
S^*\operatorname{diag}(f(z_a))S=T_f
}
\quad\text{provided}\quad
2m_i+d_i\le t_i,\quad
(A+F-F)\cap(\Lambda_L\setminus\{0\})=\varnothing.
\tag{G09.32}\label{p1b:der:g09:eq-G09-32}
\] Indeed each matrix entry integrates \(\bar u f v\); its separate sphere degrees and Fourier frequencies obey exactly these bounds. This proof permits arbitrary finite sums of such symbols, not only separated products.

For the source's rectangular set \(F=\{(n,w):|n|\le N_t,\ |w|\le N_E\}\) and symbol bounds \(|\nu_j|\le D_j\), the convenient sufficient bounds become \[
2N_j+D_j<L_j,\qquad 2m_i+d_i\le t_i.
\tag{G09.33}\label{p1b:der:g09:eq-G09-33}
\] For normalization set \(D_j=d_i=0\). For \(r\)-fold scalar products use \(rN_j<L_j\) and \(rm_i\le t_i\). The strict Fourier inequalities matter at the first multiple of a grid length.

At sphere degree two, the inherited guaranteed symbol degrees are \[
\begin{array}{c|ccc}
\Gamma&2T&2O&2I\\ \hline
t_\Gamma&5&7&11\\
\deg f\ \text{guaranteed at }m=2&1&3&7 .
\end{array}
\tag{G09.34}\label{p1b:der:g09:eq-G09-34}
\] These are separate degree bounds in the two sphere variables. For example \(q_{x,a}q_{p,b}\) is exact on \(2T\times2T\), whereas \(q_{x,a}^2\) can fail there. Both are quadratic polynomials in all eight embedding coordinates.

When the rule is normalized, its positive rank-one effects obey \[
Nw_a\le1,\qquad \#\{a\}\ge N,\qquad
\sum_{a,b}w_aw_b|K(z_a,z_b)|^2=N.
\tag{G09.35}\label{p1b:der:g09:eq-G09-35}
\] The first statements follow from \(0\le E_a\le I\) and the rank of their sum. The last is \(\operatorname{tr}(\sum E_a)^2=N\). On \(\mathcal P_2\) sampled at \(2T\), \(N=14\) and every effect has its single nonzero eigenvalue \(14/24=7/12\).

\hypertarget{p1b:der:g09:s78}{%
\subsubsection{The quadratic correction and its product transfer}\label{p1b:der:g09:s78}}

Let \(T_f^{2T}=S_{2T,2}^*\operatorname{diag}(f(q))S_{2T,2}\) on the fourteen-dimensional scalar band, and define \(D_f=T_f^{2T}-T_f\). For every quadratic symbol, \[
\boxed{
D_f=P_2^{\rm top}D_fP_2^{\rm top},\qquad
\operatorname{rank}D_f\le9.
}
\tag{G09.36}\label{p1b:der:g09:eq-G09-36}
\] If either test vector has degree at most one, \(\bar u f v\) has degree at most five, so its integral is exact. This proves the support assertion without a numerical rank test. The notation \(D_f\) here denotes the discrepancy, not a Dirac operator or the band dimension \(D_m\).

The full correction is explicit. Put \(\delta_6(q^\alpha)=24^{-1}\sum_{q\in2T}q^\alpha -\int q^\alpha\,d\mu_H\) for \(|\alpha|=6\). Independent coordinate sign symmetries annihilate odd-exponent terms. For the three even exponent patterns, \[
\begin{array}{c|ccc}
\text{nonzero exponent pattern}&(6)&(4,2)&(2,2,2)\\ \hline
\delta_6&1/64&-1/192&1/192 .
\end{array}
\tag{G09.37}\label{p1b:der:g09:eq-G09-37}
\] These numbers follow by evaluating the eight axis vertices and sixteen half-coordinate vertices, then applying (G09.13). For \(u,v\in E_2^{\rm sc}\) represented by homogeneous quadratics and a degree-two homogeneous representative \(f_2\) of the symbol, \[
\langle u,D_fv\rangle=\delta_6(\bar u f_2v).
\tag{G09.38}\label{p1b:der:g09:eq-G09-38}
\] Constants and linear parts contribute zero discrepancy. Replacing \(f_2\) by an equivalent sphere representative does not change the result: multiplication by \(r^2=1\) reduces the corresponding ambiguity to a degree-four exact moment. Equations (G09.36)--(G09.38), together with the continuum Gram matrix, specify the correction for every quadratic symbol.

On a product where all remaining evaluation maps are isometric, a symbol depending only on the first sphere satisfies \[
T_{f_x}^{\rm disc}-T_{f_x}
=D_{f_x}\otimes I_{\mathcal P_{m_p}}\otimes I_{\mathbb C[F]}.
\tag{G09.39}\label{p1b:der:g09:eq-G09-39}
\] For separated quadratic symbols on two degree-two \(2T\) bands, the sphere part of the discrepancy is \[
(T_{f_x}+D_{f_x})\otimes(T_{f_p}+D_{f_p})
 -T_{f_x}\otimes T_{f_p}
=D_{f_x}\otimes T_{f_p}
 +T_{f_x}\otimes D_{f_p}
 +D_{f_x}\otimes D_{f_p}.
\tag{G09.40}\label{p1b:der:g09:eq-G09-40}
\] A separately exact torus moment tensors onto every term. Thus the product correction has contributions from either top degree, as well as their intersection. Restricting all corrections to the joint top-by-top subspace would miss terms.

The source witness survives exactly: \[
\boxed{
\langle u,(T^{2T}_{q_0^2}-T_{q_0^2})u\rangle
=\frac1{48},\qquad
u=q_0^2-q_1^2,\quad \|u\|^2=\frac16.
}
\tag{G09.41}\label{p1b:der:g09:eq-G09-41}
\] For the unit vector \(\sqrt6\,u\), the discrepancy is \(1/8\). The same values hold after tensoring with any unit vector in the other retained factors, by (G09.39). In particular, no refinement of the momentum factor removes this position error. Paper~I-D, Section~\ref{p1b:app:g09:1} supplies the direct moment calculation, the full rank-nine spectrum, radius factors, and a covariance diagnostic.

Subtracting the actual linear-in-\(f\) discrepancy gives \(T_f=T_f^{\rm disc}-D_f\) and restores the continuum covariance. This is an exact operator correction on the stated symbol domain. It does not change the normalized effects of the original sampling POVM or make its uncorrected quadratic moments exact.

\hypertarget{p1b:der:g09:s79}{%
\subsubsection{Covariance, exchange, and fibre-valued carriers}\label{p1b:der:g09:s79}}

Let a geometric action \(g\) preserve \(d\mu\) and \(W\), with \((U_gu)(z)=u(g^{-1}z)\). Evaluation gives \(U_gK_z=K_{gz}\), and consequently \[
U_gE(A)U_g^*=E(gA),\qquad
U_gT_fU_g^*=T_{f\circ g^{-1}}.
\tag{G09.42}\label{p1b:der:g09:eq-G09-42}
\] For complete sphere bands this includes \(\operatorname{Spin}(4)_x\times\operatorname{Spin}(4)_p\); for every \(F\) it includes torus translations. The argument also determines covariance of moments of every order, including the corrected quadratic moments.

With matched sphere data and an exchange-preserved frequency set, the scalar flip \(F_{\rm ex}\) and exchange map \(\tau\) obey \[
F_{\rm ex}E(A)F_{\rm ex}^{-1}=E(\tau A),\qquad
F_{\rm ex}T_fF_{\rm ex}^{-1}=T_{f\circ\tau^{-1}}.
\tag{G09.43}\label{p1b:der:g09:eq-G09-43}
\] For a literal geometric isometry the radii and chosen torus metric must also satisfy the G07 exchange conditions. With unequal labelled parameters the formula is instead an intertwiner between the two exchanged Hilbert spaces and parameter sets; it does not define parity sectors in one space. For a sampled rule, geometric node permutations give covariance under its actual finite symmetry group. A full connected amplitude action is transported through the sampling isometry; it need not permute the nodes.

The enriched carrier requires the corresponding typed construction. Let \(H\subset L^2(M,\mathcal E,d\mu)\) be a finite space of smooth sections of a Hermitian bundle, and let \(\mathrm{ev}_z:H\to\mathcal E_z\) be evaluation. Then \[
K_H(z,w)=\mathrm{ev}_z\mathrm{ev}_w^*:
 \mathcal E_w\to\mathcal E_z,\qquad
E_H(dz)=\mathrm{ev}_z^*\mathrm{ev}_z\,d\mu(z),
\quad
\int E_H=I_H.
\tag{G09.44}\label{p1b:der:g09:eq-G09-44}
\] The proof is again the integrated fibre inner product. Its scalar moments are \(P_HM_fP_H\); its matrix-symbol moments are \(P_HM_AP_H\) for bounded measurable symbols \(A(z)\in\operatorname{End}(\mathcal E_z)\). For products, the evaluation maps and effects tensor exactly. Under a bundle isometry, the fibre map is unitary and cancels between evaluation and its adjoint, proving covariance.

For the specific fifteen-dimensional model of G02, convert to Haar-normalized sections by multiplying the volume-normalized ones by \(\sqrt{V_R}\). At every point the effect density has the eigenvalues \[
\operatorname{spec}(\mathrm{ev}_q^*\mathrm{ev}_q)
=\{0^{[6]},\,1^{[1]},\,(4/3)^{[3]},\,2^{[5]}\},
\qquad
\operatorname{rank}=9,\quad\operatorname{tr}=15.
\tag{G09.45}\label{p1b:der:g09:eq-G09-45}
\] To derive this, at a pole the scalar evaluation has squared singular value \(1\); the four retained one-forms span a three-dimensional cotangent fibre with frame operator \((4/3)I_3\); the six retained two-forms span a three-dimensional fibre with frame operator \(2I_3\). The two orthonormal lowest-spinor frames each evaluate unitarily onto the same rank-two spinor fibre, so their combined frame operator is \(2I_2\). The direct sum has rank \(1+3+3+2=9\). Homogeneity propagates these singular values to every point. The unused three-form component of the rank-ten ambient bundle does not contribute.

Thus on \(H_1^x\otimes H_1^p\otimes\mathbb C[F]\), \[
\dim=225|F|,\qquad
\operatorname{rank}(\text{evaluation effect density})=81,
\qquad
\operatorname{tr}(\text{density})=225|F|.
\tag{G09.46}\label{p1b:der:g09:eq-G09-46}
\] The nonzero eigenvalues are products of the one-factor eigenvalues and \(|F|\). On mixed scalar/enriched products, replace the corresponding rank-nine factor by rank one. The same construction applies to the finite smooth realizations at every enriched level, with their actual evaluation ranks.

Scalar multiplication preserves form degree. The intrinsic Clifford symbol \(c(df)\) can change it. Accordingly its exact evaluation formula is the matrix-symbol formula in (G09.44), not the scalar first moment \(T_f\). For instance a scalar constant and a retained one-form have zero scalar-multiplication matrix element, whereas the \(c(dx_a)\) matrix element is nonzero. This preserves the scalar/Clifford distinction already established in G02.

\hypertarget{p1b:der:g09:s80}{%
\subsubsection{Other finite cutoffs and controlled continuum limits}\label{p1b:der:g09:s80}}

The evaluation construction does not require a rectangular sphere cutoff. For a finite label set \(\mathcal I\subset\mathbb Z_{\ge0}^2\times\mathbb Z^2\) retaining each entire sphere representation at its labels, \[
K_{\mathcal I}(z,w)
=\sum_{(k,\ell,n)\in\mathcal I}
(k+1)(\ell+1)U_k(q_x\cdot q'_x)U_\ell(q_p\cdot q'_p)
 e^{in\cdot(\theta-\phi)},
\quad
N_{\mathcal I}
=\sum_{(k,\ell,n)\in\mathcal I}(k+1)^2(\ell+1)^2.
\tag{G09.47}\label{p1b:der:g09:eq-G09-47}
\] Its diagonal is constant, and (G09.7)--(G09.11) apply with \(K_{\mathcal I},N_{\mathcal I}\). This covers joint Casimir cutoffs and complete shells. Their moment matrices follow from the same basis-integral formula; they need not tensor-factor. If only charge subspaces are retained, one instead sums their actual orthonormal basis kernels. Normalization still holds, but a constant diagonal or full Spin(4) covariance must be proved from the symmetries that remain.

For rectangular bands, the finite symbol result has the useful operator consequence \[
T_f=T_{\Pi_{\mathcal V}f}\qquad(f\in L^\infty(M)).
\tag{G09.48}\label{p1b:der:g09:eq-G09-48}
\] Thus a discrete rule exact for the polynomial/Fourier space \(\mathcal V\) with the extra two test-band factors computes all compressed operators once the visible projection \(\Pi_{\mathcal V}f\) is known. It cannot replace that projection by point samples of an arbitrary bounded function.

For an increasing sequence of bands whose harmonic and Fourier labels exhaust the scalar \(L^2\) basis, \(P_j\to I\) strongly. For every bounded \(f\), \[
P_jM_fP_j\longrightarrow M_f
\quad\text{strongly on }L^2(M,d\mu).
\tag{G09.49}\label{p1b:der:g09:eq-G09-49}
\] Indeed \(P_jM_f(P_j-I)u+(P_j-I)M_fu\to0\) for each \(u\). This argument does not assert operator-norm convergence. In fact the one-sphere radial-noise norm in (G09.20) approaches \(1/2\), even as its normalized trace approaches zero.

The Berezin kernels have a quantitative concentration law: \[
\int_{S^3}|q-q'|^2
  \frac{|C_m(q\cdot q')|^2}{D_m}\,d\mu_H(q')
=\frac6{2m+3}.
\tag{G09.50}\label{p1b:der:g09:eq-G09-50}
\] It follows from \(|q-q'|^2=2(1-q\cdot q')\) and (G09.28). For a rectangular torus set, \[
\frac{d_F(\nu)}{|F|}
=\prod_{j=t,E}\left(1-\frac{|\nu_j|}{2N_j+1}\right)_+,
\qquad
\int |e^{i\theta_j}-e^{i\phi_j}|^2
 \frac{|H_F(\theta,\phi)|^2}{|F|}\,d\mu_{\mathbb T^2}(\phi)
=\frac2{2N_j+1}.
\tag{G09.51}\label{p1b:der:g09:eq-G09-51}
\] For a function Lipschitz with constant \(L\) in the Euclidean chord metric of the embedded product, Jensen's inequality gives \[
\|\mathcal Bf-f\|_\infty
\le L\left(
\frac{6R_x^2}{2m_x+3}
+\frac{6R_p^2}{2m_p+3}
+\frac2{2N_t+1}
+\frac2{2N_E+1}
\right)^{1/2}.
\tag{G09.52}\label{p1b:der:g09:eq-G09-52}
\] Here the two scalar-circle embeddings have unit radius for this mathematical estimate. Any separately specified radii multiply their two terms by their squared values. Uniform continuity and concentration extend convergence to all continuous functions when all four cutoffs tend to infinity.

For arbitrary sequences of frequency sets, the precise Fourier condition is \(d_F(\nu)/|F|\to1\) for every fixed \(\nu\). Testing characters proves necessity; density of trigonometric polynomials and the contraction bound prove sufficiency for uniform convergence on continuous functions. Increasing the number of frequencies alone does not imply this condition. Paper~I-D, Section~\ref{p1b:app:g09:3} records the irregular-set example and positive-polar normalization when an evaluation is injective but not isometric.

Paper~I-D, Section~\ref{p1b:app:g09} proves the convention-comparison identities and extended-case formulas used in the results above.


\hypertarget{p1b:sec:g10}{%
\subsection{Finite quotients and marked reference geometry}\label{p1b:sec:g10}}


\begin{proposition}[Action-dependent quotients and invariant transport]
\label{p1b:prop:g10-product}
For a specified finite isometric action with a specified bundle lift,
\[
 P_\Gamma=\frac1{|\Gamma|}\sum_{\gamma\in\Gamma}\rho(\gamma)
\]
projects onto the invariant sections. Commuting operators restrict
to this space. For a free quotient, the kernel normalized with
Riemannian volume is the sum of lifted image kernels; with probability
measure it is their average. Independent, common, correlated, central,
two-sided and exchange actions retain their separately stated quotient
types, invariant ranks and freeness conditions below.
For a left deck action, choose the commuting right Hopf action;
transport from the left common-phase convention requires simultaneous
charge and CR relabelling. For the product of the two degree-twelve
projective reference pencils, let \(h_x,h_p\), respectively \(H_x,H_p\),
be the positive degree-two generators on the domain, respectively
target, product of projective lines. The covering degree is \(144\), the
ramification divisor on the domain has class \(22(h_x+h_p)\), and
the pulled-back line has first Chern class \(-12(h_x+h_p)\) and square
\(288\). The reduced branch locus on the target has class
\(3(H_x+H_p)\).
\end{proposition}

Covering geometry, finite averaging, the integral K"unneth calculation and the explicit invariant polynomials below prove these conclusions. The detailed supporting-polytope and pencil calculations appear in Paper~I-D, Section~\ref{p1b:app:g10}.

\hypertarget{p1b:der:g10:s82}{%
\subsubsection{G10 scope and the finite-action dictionary}\label{p1b:der:g10:s82}}

This section develops finite quotients and reference geometry from the following mathematical inputs:

\begin{longtable}[]{@{}
  >{\raggedright\arraybackslash}p{(\columnwidth - 2\tabcolsep) * \real{0.5000}}
  >{\raggedright\arraybackslash}p{(\columnwidth - 2\tabcolsep) * \real{0.5000}}@{}}
\toprule\noalign{}
\begin{minipage}[b]{\linewidth}\raggedright
Stable source locator
\end{minipage} & \begin{minipage}[b]{\linewidth}\raggedright
Mathematical input
\end{minipage} \\
\midrule\noalign{}
\endhead
\bottomrule\noalign{}
\endlastfoot
Paper I, \path{p1:def:finite-support-atlases} and finite-subgroup classification & Arbitrary finite nets and the separately defined group-closed support family \\
Paper I, \path{p1:prop:exceptional-space-forms}, \path{p1:prop:quotient-invariant-restriction} & Free left quotients, abelianizations, and spectral pullback \\
Paper I, \path{p1:eq:molien-2T}, \path{p1:eq:molien-2O}, \path{p1:eq:molien-2I} & Defining-representation invariant series \\
Paper I, \path{p1:thm:exceptional-normalizer-fork}, \path{p1:thm:2O-2I-generate-SU2} & Common-core normalizers and dense sibling generation \\
Paper I-C, result~\ref{p1b:thm:relative}, result~\ref{p1b:prop:fixed}, result~\ref{p1b:thm:mixed-diagonal} & Right-diagonal quotient, exchange singularity, and diagonal spectral restriction \\
Paper I-C, result~\ref{p1b:thm:all-level-finite-restriction} & Scalar/enriched rank formulas at all independent levels \\
Paper I-C, result~\ref{p1b:prop:five-cells} and result~\ref{p1b:thm:invariant-pencil} & Marked coset geometry and the degree-twelve reference pencil \\
Paper I, \path{p1:cor:v203-geometric-positive-polar} & Dimensionless order comparison of retained LC energy and Hopf charge \\
This paper, G02--G09 & Bundle realizations, product operators, geometric exchange, and evaluation kernels \\
\end{longtable}

Use \(M=S^3_{R_x}\times S^3_{R_p}\times\mathbb T_G^2\), with positive radii, a positive flat metric \(G\), and circle coordinates of period \(2\pi\). The round metric is assumed for the explicit geometric and spectral tables below. Probability measure is used for representation and evaluation formulas; \(dV\) denotes Riemannian volume.

The action is part of every result. In particular:

\begin{longtable}[]{@{}
  >{\raggedright\arraybackslash}p{(\columnwidth - 4\tabcolsep) * \real{0.3333}}
  >{\raggedright\arraybackslash}p{(\columnwidth - 4\tabcolsep) * \real{0.3333}}
  >{\raggedright\arraybackslash}p{(\columnwidth - 4\tabcolsep) * \real{0.3333}}@{}}
\toprule\noalign{}
\begin{minipage}[b]{\linewidth}\raggedright
Construction
\end{minipage} & \begin{minipage}[b]{\linewidth}\raggedright
Action
\end{minipage} & \begin{minipage}[b]{\linewidth}\raggedright
Quotient or retained space
\end{minipage} \\
\midrule\noalign{}
\endhead
\bottomrule\noalign{}
\endlastfoot
Independent finite quotients & \((q_x,q_p)\mapsto(\gamma_xq_x,\gamma_pq_p)\) & \(Q_{\Gamma_x}\times Q_{\Gamma_p}\) \\
Common finite quotient & \((q_x,q_p)\mapsto(q_x\gamma,q_p\gamma)\) & \((S^3\times S^3)/\Delta\Gamma\) \\
Relative continuous quotient & Common right \(K=\mathrm{SU}(2)\times U(1)\) & \(M/\Delta K\cong K\) \\
Finite central identification & A specified subgroup of \(Z(M)\) acting by translations & A smooth quotient with congruence conditions \\
Exchange quotient & Swap the two matched four-dimensional factors & A symmetric square with diagonal singularities \\
Finite reference support & A finite set of points, with or without a subgroup law & A sampling/reference set, not a quotient manifold \\
\end{longtable}

The first two retain dimension eight after the torus is included. The continuous relative quotient has dimension four. A finite support itself has no positive-dimensional manifold topology. None of these constructions silently replaces the original \(M\).

\hypertarget{p1b:der:g10:s83}{%
\subsubsection{Free quotients, correlated actions, and central descent}\label{p1b:der:g10:s83}}

For a finite subgroup \(\Gamma\) of the labelled Lie group \(M\), left translation is free: \(\gamma z=z\) implies \(\gamma=e\). Thus \(X=\Gamma\backslash M\) is a closed smooth manifold and \[
\pi:M\longrightarrow X,\qquad
\deg\pi=|\Gamma|,\qquad
\operatorname{Vol}X=\frac{\operatorname{Vol}M}{|\Gamma|}.
\tag{G10.1}\label{p1b:der:g10:eq-G10-1}
\] It is an orbit quotient; a group structure on \(X\) requires normality of \(\Gamma\). In a connected Lie group a finite normal subgroup is central, since its conjugacy orbits are continuous subsets of a finite set.

The universal cover of \(M\) is \(\widetilde M=\mathrm{SU}(2)^2\times\mathbb R^2\). Let \(\widetilde\Gamma\) be the inverse image of \(\Gamma\). Then \[
\boxed{
\pi_1(X)\cong\widetilde\Gamma,\qquad
1\longrightarrow\mathbb Z^2\longrightarrow
\widetilde\Gamma\longrightarrow\Gamma\longrightarrow1.
}
\tag{G10.2}\label{p1b:der:g10:eq-G10-2}
\] The extension is central because the covering lattice lies in the central real plane. It need not split. For example the identification \((q,e^{i\theta})\sim(-q,-e^{i\theta})\) in one sphere/circle pair has a lift generated by \((-1,\pi)\); its fundamental group is \(\mathbb Z\), rather than \(C_2\times\mathbb Z\).

For independent sphere groups and a finite translation subgroup \(L\subset\mathbb T^2\), let \(\Lambda_L\subset\mathbb R^2\) be the inverse-image lattice. Then \[
X_{\rm ind}\cong Q_{\Gamma_x}\times Q_{\Gamma_p}
                   \times(\mathbb R^2/\Lambda_L),\qquad
\pi_1(X_{\rm ind})\cong\Gamma_x\times\Gamma_p\times\Lambda_L,
\tag{G10.3}\label{p1b:der:g10:eq-G10-3}
\] where \(Q_\Gamma=\Gamma\backslash S^3\) and \(\Lambda_L\cong\mathbb Z^2\).

For a specified central element \(c=((-1)^{\epsilon_x},(-1)^{\epsilon_p},e^{i\phi_t},e^{i\phi_E})\), an irreducible left representation with spins \(j_x,j_p\) and torus weight \(n=(n_t,n_E)\) descends exactly when \[
\boxed{
\exp\!\left(i\pi(2j_x\epsilon_x+2j_p\epsilon_p)
            +in\cdot\phi\right)=1
}
\tag{G10.4}\label{p1b:der:g10:eq-G10-4}
\] for every generator of the central subgroup. This follows by evaluating its scalar action. The one-pair condition is \(2j+n\) even for the example above. Such a congruence belongs to the specified quotient, not to the unquotiented \(M\).

For broader factor-preserving round isometries, write the sphere action as \(q\mapsto a q b^{-1}\). A fixed point exists exactly when \(a\) and \(b\) are conjugate in \(\mathrm{SU}(2)\), equivalently \[
\operatorname{Re}a=\operatorname{Re}b.
\tag{G10.5}\label{p1b:der:g10:eq-G10-5}
\] The pair \((-1,-1)\) is ineffective and must first be divided out. On the two-sphere product with a torus translation \(\phi\), a nonidentity element has a fixed point exactly when both equalities (G10.5) hold and \(\phi=0\) modulo the torus lattice. For a torus affine action \(A\theta+\phi\), replace that last condition by solvability of \((A-I)\theta=-\phi\) modulo \(2\pi\mathbb Z^2\). Checking each nonidentity element gives a complete freeness criterion for these specified actions. If it fails, retain the finite-isotropy orbifold and its local isotropy representations.

This includes the general lens action \((z_1,z_2)\mapsto(\zeta z_1,\zeta^r z_2)\), \(\zeta=e^{2\pi i/n}\). It is free precisely when \(\gcd(n,r)=1\). The general \(L(n,r)\) branch is wider than pure left-cyclic quotients.

\hypertarget{p1b:der:g10:s84}{%
\subsubsection{Integral topology, flat lines, and spin structures}\label{p1b:der:g10:s84}}

For a finite left sphere quotient, \[
H_0(Q_\Gamma;\mathbb Z)=H_3(Q_\Gamma;\mathbb Z)=\mathbb Z,
\quad H_1(Q_\Gamma;\mathbb Z)=\Gamma_{\rm ab},
\quad H_2(Q_\Gamma;\mathbb Z)=0.
\tag{G10.6}\label{p1b:der:g10:eq-G10-6}
\] The first homology is the abelianization of the deck group. Poincaré duality and \(H^1(Q_\Gamma;\mathbb Z)=\operatorname{Hom}(\Gamma_{\rm ab},\mathbb Z)=0\) give \(H_2=0\). Consequently \(H^2(Q_\Gamma;\mathbb Z)\) is its finite torsion group \(\Gamma_{\rm ab}\).

The group presentations and commutators give:

\begin{longtable}[]{@{}
  >{\raggedright\arraybackslash}p{(\columnwidth - 6\tabcolsep) * \real{0.2143}}
  >{\raggedleft\arraybackslash}p{(\columnwidth - 6\tabcolsep) * \real{0.2857}}
  >{\raggedright\arraybackslash}p{(\columnwidth - 6\tabcolsep) * \real{0.2143}}
  >{\raggedleft\arraybackslash}p{(\columnwidth - 6\tabcolsep) * \real{0.2857}}@{}}
\toprule\noalign{}
\begin{minipage}[b]{\linewidth}\raggedright
Group
\end{minipage} & \begin{minipage}[b]{\linewidth}\raggedleft
Order
\end{minipage} & \begin{minipage}[b]{\linewidth}\raggedright
Abelianization
\end{minipage} & \begin{minipage}[b]{\linewidth}\raggedleft
Number of spin structures on \(Q_\Gamma\)
\end{minipage} \\
\midrule\noalign{}
\endhead
\bottomrule\noalign{}
\endlastfoot
\(C_n\) & \(n\) & \(C_n\) & \(1\) for odd \(n\), \(2\) for even \(n\) \\
\(2D_n,\ n\ge2\) & \(4n\) & \(C_4\) if \(n\) odd; \(C_2^2\) if \(n\) even & \(2\) or \(4\), respectively \\
\(2T\) & \(24\) & \(C_3\) & \(1\) \\
\(2O\) & \(48\) & \(C_2\) & \(2\) \\
\(2I\) & \(120\) & \(0\) & \(1\) \\
\end{longtable}

For \(2D_n=\langle a,b:a^{2n}=1,\ b^2=a^n,\ bab^{-1}=a^{-1}\rangle\), abelianization imposes \(a^2=1\): if \(n\) is odd then \(a=b^2\) and \(b^4=1\); if \(n\) is even then both images have order two. The exceptional cases follow from \([2T,2T]=Q_8\), \([2O,2O]=2T\), and the perfectness of \(2I\).

For the product topology put \(A=(\Gamma_x)_{\rm ab}\), \(B=(\Gamma_p)_{\rm ab}\), and \(Z=Q_{\Gamma_x}\times Q_{\Gamma_p}\). Integral Künneth gives \[
\begin{array}{c|ccccccc}
k&0&1&2&3&4&5&6\\ \hline
H_k(Z)&\mathbb Z&A\oplus B&A\otimes B&
\mathbb Z^2\oplus\operatorname{Tor}(A,B)&A\oplus B&0&\mathbb Z .
\end{array}
\tag{G10.7}\label{p1b:der:g10:eq-G10-7}
\] The splitting of the Künneth extension is not asserted to be canonical. For cyclic summands \(C_a,C_b\), both their tensor and Tor groups are \(C_{\gcd(a,b)}\); a two-term resolution \(\mathbb Z\xrightarrow{a}\mathbb Z\) proves the formula. Tensoring with the free circle homology gives all eight-dimensional groups: \[
\boxed{
H_k(X_{\rm ind};\mathbb Z)
\cong H_k(Z)\oplus H_{k-1}(Z)^2\oplus H_{k-2}(Z).
}
\tag{G10.8}\label{p1b:der:g10:eq-G10-8}
\] Groups outside degrees \(0,\ldots,6\) in (G10.7) are zero.

In particular \[
H^2(X_{\rm ind};\mathbb Z)\cong\mathbb Z\oplus A\oplus B,
\quad
H^*(X_{\rm ind};\mathbb Q)
\cong\Lambda_{\mathbb Q}(a_t,a_E,u_x,u_p),
\quad
P(t)=(1+t)^2(1+t^3)^2.
\tag{G10.9}\label{p1b:der:g10:eq-G10-9}
\] The generators have degrees \(1,1,3,3\). The free \(H^2\) class is the scalar-torus class; the finite summands classify additional torsion line bundles. Flat unitary line systems form \(U(1)^2\times\widehat A\times\widehat B\). For example \(Q_{2I}\times Q_{2I}\times T^2\) has the integral homology of \(M\), while its fundamental group contains \(2I\times2I\). Equal homology does not identify the spaces.

Every left-invariant tangent frame on \(M\) descends through a left finite quotient. Therefore \(TX\) is trivial, and the spin structures form a torsor for \[
H^1(X;\mathbb Z_2)=\operatorname{Hom}(\widetilde\Gamma,\{\pm1\}).
\tag{G10.10}\label{p1b:der:g10:eq-G10-10}
\] For independent quotients their number is four times the product of the two sphere entries in the table. The factor four is the torus spin choice. Thus the \(2T/2I\) product has four, and the \(2O/2O\) product has sixteen. These statements concern left quotients; an arbitrary two-sided or exchange action requires its actual bundle lift.

\hypertarget{p1b:der:g10:s85}{%
\subsubsection{Quotient distances, injectivity, and exact diameters}\label{p1b:der:g10:s85}}

For a finite isometric action, path lifting gives \[
\boxed{d_X([z],[w])=\min_{\gamma\in\Gamma}d_M(z,\gamma w).}
\tag{G10.11}\label{p1b:der:g10:eq-G10-11}
\] A path in the quotient lifts to a path ending at one orbit representative; conversely a shortest covering geodesic projects to a path of the same length. This proves both inequalities.

For left translations on the round product, the displacement is constant in \(z\). If \(\gamma=(\gamma_x,\gamma_p,\phi)\), put \[
\delta_\gamma^2=
R_x^2\arccos(\operatorname{Re}\gamma_x)^2+
R_p^2\arccos(\operatorname{Re}\gamma_p)^2+
\min_{\ell\in2\pi\mathbb Z^2}|\phi+\ell|_G^2.
\tag{G10.12}\label{p1b:der:g10:eq-G10-12}
\] Conjugation preserves the quaternion real part, which proves the constant-displacement assertion. The first cut of a quotient exponential map comes either from a covering cut or from two identified radial geodesics. Thus \[
\operatorname{inj}X=
\min\left\{\operatorname{inj}M,\,
\tfrac12\min_{\gamma\ne e}\delta_\gamma\right\},
\quad
\operatorname{inj}M=
\min\left\{\pi R_x,\pi R_p,
\tfrac12\min_{\ell\in2\pi\mathbb Z^2\setminus0}|\ell|_G\right\}.
\tag{G10.13}\label{p1b:der:g10:eq-G10-13}
\] The second minimum over \(\Gamma\) is omitted for the trivial group. Quotient curvature is the covering curvature, so these quotients are locally symmetric and have \(\mathrm{Scal}=6/R_x^2+6/R_p^2\). The product is not a constant-sectional-curvature space: mixed planes still have zero curvature.

Right translations act transitively on \(Q_\Gamma\). Its diameter equals the covering radius of its subgroup support: \[
\operatorname{diam}Q_\Gamma
=\rho(\Gamma):=
\max_{q\in S_R^3}\min_{\gamma\in\Gamma}d(q,\gamma).
\tag{G10.14}\label{p1b:der:g10:eq-G10-14}
\] The exact one-factor values are:

\begin{longtable}[]{@{}
  >{\raggedright\arraybackslash}p{(\columnwidth - 4\tabcolsep) * \real{0.3333}}
  >{\raggedright\arraybackslash}p{(\columnwidth - 4\tabcolsep) * \real{0.3333}}
  >{\raggedright\arraybackslash}p{(\columnwidth - 4\tabcolsep) * \real{0.3333}}@{}}
\toprule\noalign{}
\begin{minipage}[b]{\linewidth}\raggedright
Left subgroup
\end{minipage} & \begin{minipage}[b]{\linewidth}\raggedright
\(\operatorname{inj}Q_\Gamma/R\)
\end{minipage} & \begin{minipage}[b]{\linewidth}\raggedright
\(\operatorname{diam}Q_\Gamma/R\)
\end{minipage} \\
\midrule\noalign{}
\endhead
\bottomrule\noalign{}
\endlastfoot
Trivial & \(\pi\) & \(\pi\) \\
\(C_n,\ n\ge2\) & \(\pi/n\) & \(\pi/2\) \\
\(2D_n,\ n\ge2\) & \(\pi/(2n)\) & \(\arccos(\cos(\pi/(2n))/\sqrt2)\) \\
\(2T\) & \(\pi/6\) & \(\pi/4\) \\
\(2O\) & \(\pi/8\) & \(\arccos((2+\sqrt2)/4)\) \\
\(2I\) & \(\pi/10\) & \(\arccos\sqrt{(7+3\sqrt5)/16}\) \\
\end{longtable}

Injectivity follows from the least nonidentity quaternion angle. For cyclic supports the farthest points are orthogonal to their containing circle. For dicyclic supports, the points lie on two orthogonal circles with \(2n\) equally spaced vertices each. The worst point has equal component norms \(1/\sqrt2\), with both phases halfway between vertices, giving the table. Paper~I-D, Section~\ref{p1b:app:g10:1} gives exact supporting-polytope proofs of all three exceptional covering radii, including a complete 600-facet certificate for \(2I\).

For independent factors, minimizations separate: \[
d_{X_{\rm ind}}^2=d_{Q_{\Gamma_x}}^2+d_{Q_{\Gamma_p}}^2+d_{T/L}^2,
\quad
\operatorname{diam}X_{\rm ind}^2
=\rho(\Gamma_x)^2+\rho(\Gamma_p)^2+
 \operatorname{diam}(T/L)^2.
\tag{G10.15}\label{p1b:der:g10:eq-G10-15}
\] The injectivity radius is the least of the three factor injectivity radii. For correlated actions (G10.11) remains exact, but the common group element must be used in both terms; separate minimization would give a different quotient.

\hypertarget{p1b:der:g10:s86}{%
\subsubsection{Scalar quotient spectra from the actual finite action}\label{p1b:der:g10:s86}}

For a unitary finite-group action, the Reynolds operator \[
\Pi_\Gamma=\frac1{|\Gamma|}\sum_{\gamma\in\Gamma}U(\gamma),
\qquad
\Pi_\Gamma^*=\Pi_\Gamma^2=\Pi_\Gamma
\tag{G10.16}\label{p1b:der:g10:eq-G10-16}
\] projects onto the invariant space. The identities follow by inverting and multiplying the group summation indices. Natural geometric operators commute with it when the action preserves their metric and bundle data.

For a left action on \(S^3\), define \(\mu_k=\dim V_{k/2}^{\Gamma}\). The scalar eigenspace is \(V_{k/2}\otimes V_{k/2}^*\); its untouched factor has dimension \(k+1\). Hence \[
\boxed{\begin{aligned}
\lambda_k&=\frac{k(k+2)}{R^2},\qquad
\operatorname{mult}_{Q_\Gamma}(\lambda_k)=(k+1)\mu_k,\\
M_\Gamma(t)&=\sum_{k\ge0}\mu_kt^k\\
&=\frac1{|\Gamma|}\sum_{\gamma\in\Gamma}
\frac1{1-2t\operatorname{Re}\gamma+t^2}.
\end{aligned}}
\tag{G10.17}\label{p1b:der:g10:eq-G10-17}
\] The last equality follows by summing the symmetric-power character recurrence \(U_{k+1}=2(\operatorname{Re}\gamma)U_k-U_{k-1}\). The scalar multiplicity generating function is \[
\sum_{k\ge0}(k+1)\mu_kt^k
=\frac{d}{dt}\bigl(tM_\Gamma(t)\bigr).
\tag{G10.18}\label{p1b:der:g10:eq-G10-18}
\]

All one-sided finite subgroup families are therefore explicit: \[
\begin{array}{c|c}
\Gamma&M_\Gamma(t)\\ \hline
C_n&(1+t^n)/((1-t^2)(1-t^n))\\
2D_n&(1+t^{2n+2})/((1-t^4)(1-t^{2n}))\\
2T&(1+t^{12})/((1-t^6)(1-t^8))\\
2O&(1+t^{18})/((1-t^8)(1-t^{12}))\\
2I&(1+t^{30})/((1-t^{12})(1-t^{20})).
\end{array}
\tag{G10.19}\label{p1b:der:g10:eq-G10-19}
\] The cyclic formula also reduces correctly to the trivial group at \(n=1\). For the cyclic group, invariant monomials have exponent difference divisible by \(n\), giving the first series. For \(2D_n\), average over the \(2n\) cyclic elements and the \(2n\) trace-zero elements: this gives \(\tfrac12M_{C_{2n}}(t)+1/(2(1+t^2))\), which simplifies to the second row. The exceptional rows are the supplied invariant-ring formulas with their numerators cancelled; the check script independently reproduces each full rational identity from the explicit quaternion classes.

The first nonconstant scalar branches are: \[
\begin{array}{c|ccc}
\Gamma&k_{\min}&R^2\lambda_1&\operatorname{mult}(\lambda_1)\\ \hline
2T&6&48&7\\
2O&8&80&9\\
2I&12&168&13.
\end{array}
\tag{G10.20}\label{p1b:der:g10:eq-G10-20}
\] Thus the first complete invariant bands have dimensions \(8,10,14\), respectively. This is the quotient/invariant filtration, not the degree-two sampling comparison of G09. For context, explicit polynomial eigenmodes for these homogeneous spherical quotients are also developed by \href{https://arxiv.org/abs/math/0502566}{Weeks, Exact Polynomial Eigenmodes}. Here the multiplicity and normalization formulas are derived directly from the group average.

For a general finite orthogonal action on the real four-coordinate sphere, Fischer harmonic decomposition instead gives \[
\sum_{k\ge0}\dim(E_k^{\rm sc})^\Gamma t^k
=\frac{1-t^2}{|\Gamma|}
 \sum_{\gamma\in\Gamma}\frac1{\det_{\mathbb R^4}(I-t\gamma)}.
\tag{G10.21}\label{p1b:der:g10:eq-G10-21}
\] Indeed all homogeneous polynomials decompose into harmonic ones plus \(r^2\) times degree \(k-2\) polynomials, equivariantly. For \(L(n,r)\), the determinant has two factors \(1-2t\cos(2\pi j/n)+t^2\) and \(1-2t\cos(2\pi rj/n)+t^2\), averaged over \(0\le j<n\). This retains the actual lens parameter \(r\); its general spectrum is not obtained by changing the group order in a one-sided exceptional formula.

\hypertarget{p1b:der:g10:s87}{%
\subsubsection{Every level pair and the diagonal comparison}\label{p1b:der:g10:s87}}

For the existing scalar and enriched families, \[
r_\Gamma(\mathcal P_\ell)=\sum_{a=0}^{\ell}(a+1)\mu_a,\qquad
r_\Gamma(H_m)=\sum_{a=0}^{2m}\binom{2m-a+2}{2}\mu_a,
\quad \dim H_m=\binom{2m+4}{4}.
\tag{G10.22}\label{p1b:der:g10:eq-G10-22}
\] The second coefficient counts the untouched spin label in \(H_m=\bigoplus_{a+b\le2m}V_{a/2}\otimes V_{b/2}\). The binomial identity \(\sum_{b=0}^{2m-a}(b+1)=\binom{2m-a+2}{2}\) proves it. Finite averaging commutes with all these level projections and with the Casimir functions used in G03.

Independent actions multiply the invariant ranks. For a common action and arbitrary finite carriers \(A,B\), \[
\boxed{
\dim(A\otimes B)^\Gamma
=\frac1{|\Gamma|}\sum_{\gamma\in\Gamma}\chi_A(\gamma)\chi_B(\gamma),
\quad
\dim\operatorname{Hom}_\Gamma(A,B)
=\frac1{|\Gamma|}\sum_{\gamma\in\Gamma}
\overline{\chi_A(\gamma)}\chi_B(\gamma).
}
\tag{G10.23}\label{p1b:der:g10:eq-G10-23}
\] Taking the trace of the actual average proves the first; applying it to \(A^*\otimes B\) proves the second. They are not interchangeable for general complex characters.

For each exceptional group, low spins \(0,\frac12,1\) remain mutually inequivalent irreducibles. The source branches \(H_1=6V_0\oplus3V_{1/2}\oplus V_1\) and \(\mathcal P_2=V_0\oplus2V_{1/2}\oplus3V_1\) then give \[
\begin{array}{c|cc}
\text{carrier pair}&\text{independent }\Gamma\times\Gamma&
\text{common diagonal }\Gamma\\ \hline
H_1\otimes H_1&36&46\\
\mathcal P_2\otimes H_1&6&15\\
\mathcal P_2\otimes\mathcal P_2&1&14.
\end{array}
\tag{G10.24}\label{p1b:der:g10:eq-G10-24}
\] The low-spin irreducibility is independently checked by the character inner products of the explicit groups. These are concrete examples of why the action must be stated before counting the retained modes.

For \(2T,2O,2I\), the first enriched contribution beyond the acted-spin-zero multiplicity occurs at \(m=3,4,6\), respectively. Its additional dimension is one, and the total ranks at those onsets are \(29,46,92\). At every fixed finite group, \[
\frac{r_\Gamma(\mathcal P_\ell)}{\dim\mathcal P_\ell}\to\frac1{|\Gamma|},
\qquad
\frac{r_\Gamma(H_m)}{\dim H_m}\to\frac1{|\Gamma|}.
\tag{G10.25}\label{p1b:der:g10:eq-G10-25}
\] The identity term contributes the displayed fraction. For noncentral elements the spin character is bounded; its weighted sum is at most \(O(\ell^2)\) or \(O(m^3)\). For the possible element \(-1\), pairing successive alternating terms gives the same bounds. The dimensions grow as \(\ell^3/3\) and \(2m^4/3\). Thus independent averaging on \(\mathcal P_\ell\otimes H_1\), with the same exceptional group on each factor, retains asymptotic fraction \(2/(5|\Gamma|)\). The common action retains instead asymptotic fraction \(1/|\Gamma|\), since only the growing factor supplies the leading identity character term.

For a common finite right subgroup, the explicit smooth identification is \[
(S^3\times S^3)/\Delta\Gamma
\longrightarrow S^3\times(S^3/\Gamma),\qquad
[(q_x,q_p)]\longmapsto(q_xq_p^{-1},[q_p]).
\tag{G10.26}\label{p1b:der:g10:eq-G10-26}
\] Its inverse is \((r,[h])\mapsto[(rh,h)]\). It is a diffeomorphism, but the induced metric generally has cross terms: \[
g=R_x^2|\operatorname{Ad}_{h^{-1}}\theta_r+\theta_h|^2
  +R_p^2|\theta_h|^2.
\tag{G10.27}\label{p1b:der:g10:eq-G10-27}
\] Consequently the common-quotient spectrum is the character restriction (G10.23), not the spectrum of a newly assumed orthogonal product under (G10.26).

Quotienting by the entire right diagonal group gives the source's relative geometry. Its round radius and scalar-circle radius obey \[
R_{\rm rel}^{-2}=R_x^{-2}+R_p^{-2},\qquad
r_{\rm rel}^{-2}=(1,-1)G^{-1}(1,-1)^T.
\tag{G10.28}\label{p1b:der:g10:eq-G10-28}
\] Minimizing the covering norm among velocities with a prescribed relative velocity proves these inverse-metric sums. The finite scalar relative band has \(k\le\min(\ell_x,\ell_p)\) and torus weights \(F_{\rm rel}=\{n:(n,-n)\in F\}\), exactly as in the source.

\hypertarget{p1b:der:g10:s88}{%
\subsubsection{Quotient operators, heat kernels, and scalar POVMs}\label{p1b:der:g10:s88}}

For probability measures, pullback identifies \(L^2(X,d\mu_X)\) unitarily with \(L^2(M,d\mu_M)^\Gamma\). For Riemannian volume the unitary is instead \[
\mathcal U f=|\Gamma|^{-1/2}f\circ\pi,\qquad
\mathcal U A_X=A_M\mathcal U.
\tag{G10.29}\label{p1b:der:g10:eq-G10-29}
\] For a Laplace-type geometric operator the second identity first holds on smooth sections, and then on its self-adjoint domain. Its spectral projectors restrict in the same way. Bundle sections require the specified equivariant fibre lift; spinors require a chosen spin structure.

The scalar heat kernel with respect to quotient Riemannian volume is the image sum \[
\boxed{
K_X(t;[z],[w])=\sum_{\gamma\in\Gamma}K_M(t;z,\gamma w).
}
\tag{G10.30}\label{p1b:der:g10:eq-G10-30}
\] There is no \(1/|\Gamma|\) in this volume-kernel formula. Integrating on a quotient fundamental domain proves the heat equation, initial value and normalization. For probability-normalized kernels the corresponding formula is the average \(|\Gamma|^{-1}\sum K_M^\mu\). For bundles, insert the map from the fibre at \(w\) to that at \(\gamma w\) after \(K_M(z,\gamma w)\).

For independent one-sided sphere quotients, \[
Z_\Gamma(t;R)=\sum_{k\ge0}(k+1)\mu_k e^{-tk(k+2)/R^2},
\qquad
Z_{X_{\rm ind}}(t)=Z_{\Gamma_x}(t;R_x)
                  Z_{\Gamma_p}(t;R_p)\Theta_{G,L}(t).
\tag{G10.31}\label{p1b:der:g10:eq-G10-31}
\] Here \(\Theta_{G,L}\) uses exactly the torus character lattice annihilating \(L\). Its first positive scalar eigenvalue is the minimum of the three factor gaps. Coincident minima add their multiplicities. For a common finite sphere action, the multiplicity for degrees \(k,\ell\) is \[
(k+1)(\ell+1)\frac1{|\Gamma|}
\sum_{\gamma\in\Gamma}\chi_{k/2}(\gamma)\chi_{\ell/2}(\gamma),
\tag{G10.32}\label{p1b:der:g10:eq-G10-32}
\] at eigenvalue \(k(k+2)/R_x^2+\ell(\ell+2)/R_p^2\). This specifies all common-quotient heat sums as well.

Every local heat coefficient for a free finite quotient is the covering coefficient divided by \(|\Gamma|\): all nonidentity images have positive displacement and contribute no term to the local short-time power series. In particular, for a rank-\(r\) Laplace-type principal symbol, \[
Z_X(t)\sim\frac{r\,\operatorname{Vol}X}{(4\pi t)^4}+\cdots,
\qquad
N_X(\Lambda)\sim
\frac{r\,\operatorname{Vol}X}{(4\pi)^4\Gamma(5)}\Lambda^4.
\tag{G10.33}\label{p1b:der:g10:eq-G10-33}
\] For constrained TT or Rarita spaces, use their actual G08 constraint projectors, ranks, and principal-symbol speeds. Formula (G10.33) applies to a \(|\xi|^2I\) principal symbol; otherwise retain the sum of the separate speed contributions. Isometric finite averaging commutes with those constraints, so the quotient matrices are their invariant restrictions. No scalar multiplicity is substituted for a bundle or constrained rank.

The one-sphere Hodge and spin multiplicities can be written explicitly. Exact one-forms have the scalar multiplicity \((k+1)\mu_k\), \(k\ge1\), at \(k(k+2)/R^2\). Coexact one-forms have \[
\lambda_k^{\rm coex}=(k+1)^2/R^2,\qquad
d_k^{\rm coex}=k\mu_{k+1}+(k+2)\mu_{k-1},\quad k\ge1.
\tag{G10.34}\label{p1b:der:g10:eq-G10-34}
\] This is the restriction of \(((k+1)/2,(k-1)/2)\oplus((k-1)/2,(k+1)/2)\). Hodge star supplies degrees two and three.

For the canonical left-frame spin lift twisted by \(\epsilon:\Gamma\to\{\pm1\}\), replace \(\mu_k\) by \(\mu_k^\epsilon=|\Gamma|^{-1}\sum_\gamma \epsilon(\gamma)\chi_{k/2}(\gamma)\). With the sign convention of G08.6, the signed Dirac multiplicities are \[
d_{k,+}^\epsilon=(k+1)\mu_{k+1}^\epsilon,\qquad
d_{k,-}^\epsilon=(k+2)\mu_k^\epsilon,
\qquad
D=\pm(k+3/2)/R.
\tag{G10.35}\label{p1b:der:g10:eq-G10-35}
\] They need not agree on a three-dimensional quotient. The eight-dimensional product uses the G08 Clifford matrix, with these actual multiplicities and the selected torus spin weights. Positivity of the scalar curvature excludes zero Dirac modes for these round products and flat unitary twists.

Finally, a finite scalar band descends by replacing its kernel with \[
K_W^\Gamma(z,w)=\frac1{|\Gamma|}
 \sum_{\gamma\in\Gamma}K_W(z,\gamma w).
\tag{G10.36}\label{p1b:der:g10:eq-G10-36}
\] It is the reproducing kernel of \(W^\Gamma\) in probability measure. Its evaluation POVM and moments are \[
\int |K_z^\Gamma\rangle\langle K_z^\Gamma|\,d\mu_X=I_{W^\Gamma},
\qquad T_f^\Gamma=P_{W^\Gamma}M_{\pi^*f}P_{W^\Gamma}.
\tag{G10.37}\label{p1b:der:g10:eq-G10-37}
\] On left homogeneous quotients the diagonal is the constant \(\dim W^\Gamma\) for the complete invariant bands. Exact sampled evaluation restricts to the invariant subspace. If an invariant weighted node rule is collapsed to node orbits, add their weights and use invariant symbols. This transports G09's actual moment domain; it does not increase a finite design's degree of exactness.

\hypertarget{p1b:der:g10:s89}{%
\subsubsection{The descended Hopf geometry and its orbifold base}\label{p1b:der:g10:s89}}

Choose the right Hopf action \(q\mapsto qe^{itn}\) and its left-invariant connection form \(\eta\), with \(\eta(N)=1\) for the period-\(2\pi\) generator \(N\). Every left deck group commutes with this circle action, so the contact distribution, its coorientation, and its compatible CR structure descend. Use the G05--G06 operator formulas with this right-Hopf marking. Those sections use a left common-phase action: inversion obeys \(\iota(e^{itn}q)=q^{-1}e^{-itn}\), so its transport reverses the angular generator relative to the positive right-circle parameter used here. Relabel charge and CR conventions together before applying their fixed-charge formulas. This statement does not assert descent of the original fixed left-Hopf marking through every left deck group. The base is \[
\Gamma\backslash S^3/U(1)
\cong\bar\Gamma\backslash S^2,
\qquad
\bar\Gamma=\Gamma/(\Gamma\cap\{\pm1\}).
\tag{G10.38}\label{p1b:der:g10:eq-G10-38}
\] It is generally an orbifold base of a smooth Seifert total space. The deck action has no fixed points on \(S^3\); the exceptional fibres arise from stabilizers of the descended circle action.

Let \(c=|\Gamma\cap\{\pm1\}|\), equal to one or two. The kernel of the right circle action on \(Q_\Gamma\) is exactly this central subgroup. Indeed, if a circle element acts trivially at every point, \(qhq^{-1}\in\Gamma\) for every \(q\). Continuity makes it constant, so \(h\) is central. The effective period is \(2\pi/c\), and the normalized effective connection is \(c\eta\). If the original Hopf Chern number in the fixed orientation is \(e_H=\pm1\), then \[
e_{\rm orb}
=\frac1{2\pi}\int_{\bar\Gamma\backslash S^2}^{\rm orb}d(c\eta)
=\frac{c}{|\bar\Gamma|}e_H
=\frac{c^2}{|\Gamma|}e_H.
\tag{G10.39}\label{p1b:der:g10:eq-G10-39}
\] The sign follows the Hopf convention; the magnitude is fixed. For the exceptional groups this yields:

\begin{longtable}[]{@{}llr@{}}
\toprule\noalign{}
Total space & Orbifold base signature & \(\vert e_{\rm orb}\vert\) \\
\midrule\noalign{}
\endhead
\bottomrule\noalign{}
\endlastfoot
\(Q_{2T}\) & \(S^2(2,3,3)\) & \(1/6\) \\
\(Q_{2O}\) & \(S^2(2,3,4)\) & \(1/12\) \\
\(Q_{2I}\) & \(S^2(2,3,5)\) & \(1/30\) \\
\end{longtable}

The cone orders are the stabilizer orders of tetrahedral, octahedral, and icosahedral rotations. There are respectively three branch-value orbits. For the dicyclic branch the base signature is \(S^2(2,2,n)\); a nontrivial projective cyclic branch has two cone points of its projective order. On an untwisted binary quotient, a scalar right-Hopf charge must be even, because the element \(-1\) is already a deck identification. Its effective circle weight is half that original weight.

For independent sphere quotients the bi-Hopf base is the product of the two displayed orbifold bases, with the scalar torus retained. The contact pair, fixed-charge complexes and their tensor constructions are the invariant restrictions of G05--G06. In particular, a charge restriction and finite averaging commute for these compatible left actions.

This compatibility is specific to the marking. For a two-sided deck action \(q\mapsto aqb^{-1}\), the right Hopf axis is transported by \(\operatorname{Ad}_b\). A fixed oriented axis descends when this action fixes \(n\); normalizing its unoriented line suffices only for the unoriented foliation. A global ordered Hopf triple requires trivial action on all three axes. The opposite, left-Hopf marking must be checked against the left deck factors instead. Thus the descended structures follow from their actual equivariance, not merely from the quotient being spherical.

\hypertarget{p1b:der:g10:s90}{%
\subsubsection{Exchange quotients and their exact spectral correction}\label{p1b:der:g10:s90}}

For a matched left group let \(K_\Gamma=Q_\Gamma\times S^1\). The product \(Y=K_\Gamma\times K_\Gamma\) admits exchange \(B(x,y)=(y,x)\). Its fixed set is \(\Delta K_\Gamma\). The tangent action is \(+I\) on a four-dimensional tangent space and \(-I\) on a four-dimensional normal space. Therefore \[
Y/\langle B\rangle=\operatorname{Sym}^2K_\Gamma,\qquad
\text{local model at the diagonal }
\mathbb R^4\times(\mathbb R^4/\{\pm1\}).
\tag{G10.40}\label{p1b:der:g10:eq-G10-40}
\] The normal link is \(\mathbb{RP}^3\), so this is singular along the diagonal and smooth elsewhere. Exchange preserves orientation. A free-quotient formula cannot erase this isotropy or its heat contribution.

There are two relevant fundamental groups: \[
\pi_1^{\rm orb}(Y/\langle B\rangle)
\cong(\Gamma\times\mathbb Z)^2\rtimes C_2,\qquad
\boxed{
\pi_1(|Y/\langle B\rangle|)
\cong\Gamma_{\rm ab}\oplus\mathbb Z.
}
\tag{G10.41}\label{p1b:der:g10:eq-G10-41}
\] The first is the group acting on the simply connected \((S^3\times\mathbb R)^2\). For the underlying space, the stabilizer transposition is killed; its conjugates identify the two copies of each loop, and the copies commuted before identification. The resulting group is the abelianization. The general orbit-space theorem justifying this passage applies to these discontinuous actions on manifolds; the symmetric-square result is \href{https://arxiv.org/pdf/math/0212271}{Brown--Higgins, Proposition 2.12}.

Over \(\mathbb Q\), quotient cohomology is the invariant cohomology. For the graded flip its trace polynomial is \(1-t^2-t^6+t^8\), since the one-factor degrees are \(0,1,3,4\). Hence the quotient Betti polynomial is \[
P_+(t)=\tfrac12\left((1+t)^2(1+t^3)^2+
                         1-t^2-t^6+t^8\right)
=1+t+t^3+2t^4+t^5+t^7+t^8.
\tag{G10.42}\label{p1b:der:g10:eq-G10-42}
\] This also follows from equivariant cellular cochains after inverting the finite stabilizer orders. It is ordinary rational cohomology of the underlying orbifold space.

For an orthogonal matched product metric, scalar heat on the quotient is \[
\boxed{
Z_{\operatorname{Sym}^2K_\Gamma}(t)
=\tfrac12\bigl(Z_{K_\Gamma}(t)^2+Z_{K_\Gamma}(2t)\bigr).
}
\tag{G10.43}\label{p1b:der:g10:eq-G10-43}
\] On a tensor eigenbasis the trace of exchange is supported exactly on equal factors, proving the \(2t\) term. For an exchange-invariant torus metric with a cross term, the more general formula is \[
Z_+(t)=\tfrac12\left(
 Z_\Gamma(t;R)^2\Theta_G(t)
 +Z_\Gamma(2t;R)\!
   \sum_{n\in\mathbb Z}e^{-t(n,n)G^{-1}(n,n)^T}\right).
\tag{G10.44}\label{p1b:der:g10:eq-G10-44}
\] The second term has order \(t^{-2}\) as \(t\downarrow0\); it records the four-dimensional fixed stratum. It is absent from a free two-sheeted quotient.

For a finite band \(A\) on \(K_\Gamma\), with kernel \(k\) and constant diagonal \(D=\dim A\), the symmetric kernel is \[
K_+((z_1,z_2),(w_1,w_2))
=\tfrac12\bigl(k(z_1,w_1)k(z_2,w_2)
              +k(z_1,w_2)k(z_2,w_1)\bigr).
\tag{G10.45}\label{p1b:der:g10:eq-G10-45}
\] Its dimension is \(D(D+1)/2\), but its diagonal is the generally nonconstant \(\tfrac12(D^2+|k(z_1,z_2)|^2)\). The evaluation POVM still integrates to identity with the pushforward probability measure; its normalized coherent vectors therefore carry this varying density. One must not replace that density by the total dimension. For bundles use the G08 geometric exchange lift and exchange-compatible spin structures; scalar symmetrization alone does not specify those lifts.

\hypertarget{p1b:der:g10:s91}{%
\subsubsection{Finite reference resolution and its exact limitation}\label{p1b:der:g10:s91}}

For finite supports in an orthogonal product, nearest-point minimization separates. Thus \[
\rho(Z_x\times Z_p\times Z_T)^2
=\rho(Z_x)^2+\rho(Z_p)^2+\rho(Z_T)^2.
\tag{G10.46}\label{p1b:der:g10:eq-G10-46}
\] If two labelled supports have Hausdorff distances at most \(\epsilon_i\) in their factors, their products have Hausdorff distance at most \((\sum_i\epsilon_i^2)^{1/2}\). This is a reference-accuracy statement, independent of polynomial quadrature accuracy.

On a round sphere a cap of angular radius \(\theta\) has probability volume \[
v_3(\theta)=\frac{\theta-\sin\theta\cos\theta}{\pi},
\qquad 0\le\theta\le\pi.
\tag{G10.47}\label{p1b:der:g10:eq-G10-47}
\] Integrating the radial volume element \(4\pi R^3\sin^2\theta\,d\theta\) proves this formula. A spherical \(\epsilon\)-net of size \(N\) therefore obeys \(Nv_3(\epsilon/R)\ge1\). A maximal \(\epsilon\)-separated set is an \(\epsilon\)-net, and its disjoint \(\epsilon/2\) caps give \(N\le1/v_3(\epsilon/(2R))\). These bounds prove existence of arbitrarily fine finite reference sets, with the expected dimension-three growth.

More generally on the fixed smooth eight-dimensional product, small metric balls have volume \(\omega_8\epsilon^8(1+O(\epsilon^2))\), \(\omega_8=\pi^4/24\). The same covering and packing arguments give \[
\frac{\operatorname{Vol}M}{\omega_8\epsilon^8}(1-C\epsilon^2)
\le N_{\min}(\epsilon)\le
\frac{2^8\operatorname{Vol}M}{\omega_8\epsilon^8}(1+C\epsilon^2).
\tag{G10.48}\label{p1b:der:g10:eq-G10-48}
\] Here \(C>0\) depends on the fixed metric and the estimates hold for all sufficiently small positive \(\epsilon\). In particular the order is \(\epsilon^{-8}\).

Group-closed spherical supports are a different family. The finite-subgroup classification and the diameter table give the sharp universal bound \[
\boxed{
\rho(\Gamma)\ge
R\arccos\sqrt{\frac{7+3\sqrt5}{16}}
\quad\text{for every finite }\Gamma\subset SU(2),
}
\tag{G10.49}\label{p1b:der:g10:eq-G10-49}
\] with equality for \(2I\) up to the unmarked isometry. The cyclic radii are at least \(\pi R/2\), the dicyclic radii exceed \(\pi R/4\), and the other exceptional radii are larger than the displayed \(2I\) value. Thus arbitrarily large finite cyclic or dicyclic groups do not become fine nets of the whole sphere. On a product of group-closed supports, the two positive sphere lower bounds add in quadrature even if the scalar torus grid is refined indefinitely.

This resolves the two reference atlases in the source: all finite \(\epsilon\)-nets provide arbitrary accuracy; finite subgroup supports provide their exact group, incidence, quotient and spectral structures. Neither statement imposes a preferred support or changes the given finite sampling rules.

\hypertarget{p1b:der:g10:s92}{%
\subsubsection{Normalizer forks, coset products, and simultaneous invariants}\label{p1b:der:g10:s92}}

For the marked Hurwitz core \(H=2T\), the source normalizers give the independent and common product normalizers \[
N_{SU(2)^2}(H\times H)=2O\times2O,
\qquad
N_{SU(2)^2}(\Delta H)=\{(a,\epsilon a):a\in2O,\ \epsilon=\pm1\}.
\tag{G10.50}\label{p1b:der:g10:eq-G10-50}
\] For the second identity, both components must induce the same conjugation on \(H\); their ratio lies in \(C_{SU(2)}(H)=\{\pm1\}\). Both normalizer quotients by their respective subgroups are abstractly \(C_2^2\), but the normalizers have different orders, \(2304\) and \(96\), and act on different carriers. For exchange-compatible extensions, append the actual factor-swap semidirect action.

Within \(I=2I\), \(H\) is self-normalizing. Therefore \(I/H\) is the five-element coset space \(X\), not a quotient group. Independent \(I\times I\) acts transitively on \(X\times X\), whereas the common diagonal \(I\) has the two orbits of sizes \(5\) and \(20\). Consequently \[
\dim(\mathbb C^{X\times X})^{I\times I}=1,\qquad
\dim(\mathbb C^{X\times X})^{\Delta I}=2.
\tag{G10.51}\label{p1b:der:g10:eq-G10-51}
\] The one-factor permutation module is \(1\oplus V_4\); under independent actions the product sectors have dimensions \(1,4,4,16\). This is the correct representation decomposition for transporting the source's specified incidence Gram matrix \(120(P\otimes P)+30(Q\otimes Q)\). That Gram matrix has rank \(17\) and kernel dimension \(8\). Those coefficients remain properties of the supplied incidence map; cardinality \(25\) alone does not derive them. Likewise the scalar \(\mathbb F_5^2\) action with orbit sizes \(1,24\) is not this five-cell product action.

The source's dense generation \(\overline{\langle2O,2I\rangle}=SU(2)\) implies, on every one and the same finite continuous representation \(A\), \[
A^{2O}\cap A^{2I}=A^{SU(2)},\qquad
\operatorname{End}(A)^{2O}\cap\operatorname{End}(A)^{2I}
=\operatorname{End}(A)^{SU(2)}.
\tag{G10.52}\label{p1b:der:g10:eq-G10-52}
\] Continuity extends a common fixed vector or commuting operator to the generated closure. The closure is infinite by the finite-subgroup classification, and no proper positive-dimensional closed subgroup can contain \(2I\), so it is \(SU(2)\). Independent actions in both sphere factors give the full \(SU(2)^2\) invariant space; common diagonal actions give only the diagonal \(SU(2)\) invariant space.

There is an explicit finite-dimensional projection procedure. For \(T=\Pi_{2I}\Pi_{2O}\Pi_{2I}\), its \(1\)-eigenspace is the intersection above. All other eigenvalues lie in \([0,1)\); equality in the projection norm inequality proves this assertion. Hence \[
T^n\longrightarrow\Pi_{SU(2)},\qquad
\|T^n-\Pi_{SU(2)}\|=\lambda_*^n,
\tag{G10.53}\label{p1b:der:g10:eq-G10-53}
\] where \(\lambda_*\) is the largest eigenvalue on the orthogonal complement, taken as zero if that complement has no nonzero eigenvalue. The rate is specific to the fixed carrier. This is not an assertion that a single product of the two averages is already the common projector.

\hypertarget{p1b:der:g10:s93}{%
\subsubsection{The product reference pencil and retained order cone}\label{p1b:der:g10:s93}}

The source's absolute degree-twelve \(2T\)-invariant pencil has no base point and gives a holomorphic map \(\varphi:\mathbb{CP}^1\to\mathbb P(I_{12}^*)\cong\mathbb{CP}^1\) of degree \(12\). It is constant on \(A_4=2T/\{\pm1\}\) orbits; a generic orbit has twelve points, so it realizes the projective quotient. Its branching is \((2,3,3)\), matching the orbifold base of \(Q_{2T}\) in Section~\ref{p1b:der:g10:s89}.

For two marked pencils the tensor evaluation and its polar normalization are \[
\epsilon_{L_x,L_p}=\epsilon_{L_x}\otimes\epsilon_{L_p},
\qquad
U_{L_x,L_p}=U_{L_x}\otimes U_{L_p}:
L_x^{12}\otimes L_p^{12}\longrightarrow A_{12,x}\boxtimes A_{12,p}.
\tag{G10.54}\label{p1b:der:g10:eq-G10-54}
\] The tensor product is injective on every product line; positive scalar Gram factors prove the polar identity. Writing \(h_x,h_p\) for the positive generators of the two projective lines, \[
c_1(A_{12,x}\boxtimes A_{12,p})=-12(h_x+h_p),
\qquad
\int_{\mathbb{CP}^1\times\mathbb{CP}^1}c_1^2=288.
\tag{G10.55}\label{p1b:der:g10:eq-G10-55}
\] Projectivization is the product degree-\(12\) quotient map followed by the Segre embedding. It is generically \(144\)-to-one onto the quadric surface in \(\mathbb{CP}^3\). Its ramification divisor on the domain has class \(22(h_x+h_p)\). Paper~I-D, Section~\ref{p1b:app:g10:2} gives the polynomial calculation and distinguishes its projection connection from the connection transported by the unitary polar map.

There is also an exact product of the dimensionless spectral reference in Paper I. On each retained \(H_1\), let \(E_i=R_i\sqrt{K_{{\rm LC},i}}\) and let \(N_i\) be its integer right-Hopf generator. G02 and the source give \([E_i,N_i]=0\) and \(E_i\ge|N_i|\). On the product, the entire separated linear order cone is \[
\boxed{
a_xE_x+a_pE_p+b_xN_x+b_pN_p\ge0
\ \Longleftrightarrow\
a_x\ge|b_x|,\quad a_p\ge|b_p|.
}
\tag{G10.56}\label{p1b:der:g10:eq-G10-56}
\] All coefficients are real. Sufficiency follows on the joint eigenbasis. For necessity, place the other factor in its zero scalar mode and use the states with \(E_i=2\), \(N_i=\pm2\).

In particular, with \(E_\Sigma=E_x+E_p\) and \(N_{\rm rel}=N_x-N_p\), \[
N_\pm=\frac{E_\Sigma\pm N_{\rm rel}}{\sqrt2}\ge0,\qquad
\dim\ker N_+=\dim\ker N_-=4.
\tag{G10.57}\label{p1b:der:g10:eq-G10-57}
\] Each kernel is a tensor product of the two corresponding one-factor two-dimensional kernels. An untouched scalar character factor multiplies these nullities by \(|F|\). Compatible left finite averaging preserves these kernel states and the inequality; a right finite action requires checking commutation with the marked Hopf generator first. Every exact sampling isometry transports these reference operators and their spectra without changing the carrier. The result is a dimensionless spectral order comparison.

Paper~I-D, Section~\ref{p1b:app:g10} proves the convention-comparison identities and extended-case formulas used in the results above.


\subsubsection{Exceptional product invariant rings and McKay--ADE geometry}
\label{p1b:sec:g23}

Let \(\Gamma_x,\Gamma_p\in\{2T,2O,2I\}\) act in their defining
representations \(V_x,V_p\cong\C^2\), and let
\(\Gamma_x\times\Gamma_p\) act on \(V_x\oplus V_p\) factorwise.
The one-factor Molien series are already part of the finite atlas; the
results below close their genuinely reciprocal product geometry.

\begin{theorem}[Product invariant algebra and Molien factorization]
\label{p1b:thm:g23-invariants}
There is a canonical bigraded algebra isomorphism
\[
 \boxed{
 \C[V_x\oplus V_p]^{\Gamma_x\times\Gamma_p}
 \cong
 \C[V_x]^{\Gamma_x}\otimes_\C\C[V_p]^{\Gamma_p}.}
\]
Consequently the bigraded Molien series factorizes,
\[
 \boxed{M_{\Gamma_x\times\Gamma_p}(s,t)
 =M_{\Gamma_x}(s)M_{\Gamma_p}(t).}
\]
Using the one-factor normalizations already fixed in the atlas,
\[
\begin{array}{c|c|c}
\Gamma& M_\Gamma(q)&\text{weighted hypersurface type}\\ \hline
2T&\dfrac{1+q^{12}}{(1-q^6)(1-q^8)}&
X^2+Y^3+Z^4=0\quad(E_6)\\[0.8em]
2O&\dfrac{1+q^{18}}{(1-q^8)(1-q^{12})}&
X^2+Y^3+YZ^3=0\quad(E_7)\\[0.8em]
2I&\dfrac{1+q^{30}}{(1-q^{12})(1-q^{20})}&
X^2+Y^3+Z^5=0\quad(E_8).
\end{array}
\]
Thus each of the nine reciprocal exceptional pairs is obtained by tensoring
the corresponding two rows, not by introducing a new irreducible
four-dimensional invariant-ring type.
\end{theorem}
\begin{proof}
Polynomial functions on a direct sum tensor-factor as
\(\C[V_x]\otimes\C[V_p]\).  Averaging over the direct product group is the
tensor product of the two Reynolds operators, so its image is exactly the
tensor product of the invariant subalgebras.  Taking bigraded dimensions
gives the Molien product.  The displayed one-factor hypersurface forms are
the standard binary-polyhedral/Kleinian normalizations compatible with the
listed Hilbert series.
\end{proof}

\begin{corollary}[Product Kleinian singularity and singular locus]
\label{p1b:cor:g23-kleinian}
Let \(X_\Gamma=\C^2/\Gamma\).  Then
\[
 (V_x\oplus V_p)/(\Gamma_x\times\Gamma_p)
 \cong X_{\Gamma_x}\times X_{\Gamma_p}.
\]
Each factor has its isolated ADE singularity at the origin, while the
four-complex-dimensional product has
\[
 \boxed{\operatorname{Sing}(X_{\Gamma_x}\times X_{\Gamma_p})
 =\bigl(\{0\}\times X_{\Gamma_p}\bigr)
 \cup\bigl(X_{\Gamma_x}\times\{0\}\bigr).}
\]
The two singular components meet at the product origin.  Hence the
reciprocal quotient is not itself a new simple ADE surface singularity.
\end{corollary}
\begin{proof}
The quotient identity is the spectrum of
Theorem~\ref{p1b:thm:g23-invariants}.  Smoothness of a product holds
exactly when both factors are smooth, which gives the singular-locus
formula.
\end{proof}

\begin{theorem}[Product McKay operator]
\label{p1b:thm:g23-mckay}
Let \(A_x,A_p\) be the affine McKay adjacency matrices for the defining
representations of \(\Gamma_x,\Gamma_p\), labelled respectively by affine
\(E_6,E_7,E_8\) as in the McKay correspondence \cite{McKay1980}.  On the
irreducible product representations \(\rho_i\boxtimes\sigma_j\), tensoring
by \(V_x\oplus V_p\) has adjacency matrix
\[
 \boxed{A=A_x\otimes I+I\otimes A_p.}
\]
Thus the reciprocal McKay graph is the Cartesian product of the two affine
ADE graphs.  If \(C_i=2I-A_i\) are the affine Cartan matrices, then the
natural product Cartan operator is
\[
 \boxed{C=4I-A=C_x\otimes I+I\otimes C_p.}
\]
If \(d_x,d_p\) are the one-factor dimension vectors, then
\(d_x\otimes d_p\) spans the nullspace of \(C\).
For a matched pair \(\Gamma_x=\Gamma_p\), Born exchange acts on vertices by
\((i,j)\mapsto(j,i)\); for unequal types it identifies the ordered pair
atlas with the oppositely ordered one rather than defining an internal
symmetry of a single labelled pair.
\end{theorem}
\begin{proof}
Distribute tensor product over the direct sum:
\[
 (V_x\oplus V_p)\otimes(\rho_i\boxtimes\sigma_j)
 =(V_x\otimes\rho_i)\boxtimes\sigma_j
 \oplus\rho_i\boxtimes(V_p\otimes\sigma_j).
\]
This is exactly the Kronecker-sum adjacency formula.  The Cartan identity
follows algebraically.  Affine ADE Cartan matrices are positive semidefinite
with one-dimensional kernels spanned by their dimension vectors; the
Kronecker sum is therefore positive semidefinite and has kernel the tensor
product of those two kernels.
\end{proof}

Generator weights, complete-intersection presentations and the nine-pair
bookkeeping are collected in Paper~I-D, Appendix~\ref{p1b:app:g23}.


\part{Relational classification, reduction and rooted structures}
\section{Relational closure, exhaustion, and regular reductions}
\label{p1b:sec:organized-closure}

\subsection{Cross-structure relational closure}
\label{p1b:sec:relational-closure}
The preceding sections classify the individual structure families.  The
results here close the remaining maps between those families.  No new
geometric datum is introduced: every statement is a consequence of the
already declared Cartan, groupoid, operator, bundle, metric and complex
structures.  In particular, the finite commutator is related to curvature
only to its proved order; the exact statement is the curvature--Frobenius
identity below.

\begin{theorem}[Cartan curvature, Frobenius defect and finite commutator]
\label{p1b:thm:relational-cartan-commutator}
Let \(c(k)=(k,k^{-1})\) be the Cartan embedding and, for
\(X\in\mathfrak k\), put
\[
 U_X=(dc)_eX=(X,-X)\in T_{(e,e)}M_8.
\]
Then \(U_X\) is horizontal for the canonical connection and
\[
 \mathcal A(U_X)=0,\qquad \mathcal E(U_X)=X,\qquad
 F_{\mathcal A}(U_X,U_Y)=-[X,Y].
\]
If \(\widetilde U_X,\widetilde U_Y\) are arbitrary local
\(\mathcal A\)-horizontal extensions of these vectors, then the vertical
Frobenius defect is exactly
\[
 \boxed{\mathcal A([\widetilde U_X,\widetilde U_Y]_{(e,e)})
       =[X,Y]=-F_{\mathcal A}(U_X,U_Y).}
\]
The finite order-swap defect therefore has the canonical infinitesimal
relation
\[
 \boxed{
 \log\mu(e^{\varepsilon X},e^{\varepsilon Y})
   =-\varepsilon^2F_{\mathcal A}(U_X,U_Y)+O(\varepsilon^3)
   =\varepsilon^2[X,Y]+O(\varepsilon^3).}
\]
Thus curvature and horizontal nonintegrability are exactly the same
Lie-algebra-valued defect, while \(\mu\) is its finite group-valued
integration to second order.  By Theorem~\ref{p1b:thm:g12-double}, this
same \(\mu\) is the group-valued moment map of the compact double, so the
Cartan and quasi-Hamiltonian branches share the finite commutator variable
even though their two-forms are not globally identical.  On the scalar
\(\mathfrak u(1)\) direction all three defects vanish identically.
\end{theorem}
\begin{proof}
The Cartan formulas give
\(\mathcal A(X,-X)=0\) and \(\mathcal E(X,-X)=X\).  Hence
\(F_{\mathcal A}=-\mathcal E\wedge\mathcal E\) gives
\(F_{\mathcal A}(U_X,U_Y)=-[X,Y]\).  For horizontal vector fields,
\[
 F_{\mathcal A}(\widetilde U_X,\widetilde U_Y)
 =\dd\mathcal A(\widetilde U_X,\widetilde U_Y)
 =-\mathcal A([\widetilde U_X,\widetilde U_Y]),
\]
since both horizontal values of \(\mathcal A\) vanish. This proves the
exact Frobenius identity.  The Baker--Campbell--Hausdorff expansion of the
group commutator is Proposition~\ref{p1b:prop:commutator}; substitution of
the curvature identity gives the displayed bridge.  The abelian scalar
factor has zero bracket.
\end{proof}

\begin{theorem}[Principal symbols and the tangent-groupoid degeneration]
\label{p1b:thm:relational-tangent-symbol}
On the relative group \(K=\SU(2)\times\U(1)\), let
\(C_{\rm rel}\) be the nonnegative quadratic \(\SU(2)\) Casimir and
\(N_{\rm rel}\) the primitive circle generator.  The complete
bi-invariant scalar differential-operator algebra is
\[
 \mathcal D_{\rm bi}(K)=\C[C_{\rm rel},N_{\rm rel}].
\]
Give it the differential-order filtration
\(\deg C_{\rm rel}=2\), \(\deg N_{\rm rel}=1\).  Under the left
trivialization \(T^*K\cong K\times\mathfrak k^*\), put
\[
 \rho(\xi)=\|\xi_{\mathfrak{su}(2)}\|^2,\qquad
 \nu(\xi)=\xi_{\mathfrak u(1)}.
\]
Then
\[
 \boxed{\operatorname{gr}\mathcal D_{\rm bi}(K)
       \cong S(\mathfrak k_\C)^K=\C[\rho,\nu].}
\]
For an operator \(P\) of order \(m\), its principal symbol is obtained by
taking the weighted degree-\(m\) part of its polynomial in
\((C_{\rm rel},N_{\rm rel})\) and replacing these generators by
\((\rho,\nu)\).  Inversion fixes \(\rho\) and negates \(\nu\), so
\[
 \operatorname{gr}\mathcal D_{\rm bi}(K)^{\rm inv+}
 =\C[\rho,\nu^2].
\]
In the tangent groupoid of Theorem~\ref{p1b:thm:convolution}, the rescaled
order-\(m\) family of \(P\) has at the zero fibre the Fourier multiplier
\(\sigma_m(P)\) on \(T^*K\).  Polynomial principal symbols act as generally
unbounded multiplication operators on the zero-fibre algebra \(C_0(T^*K)\),
with natural domain
\[
 \operatorname{Dom}M_p=\{f\in C_0(T^*K):pf\in C_0(T^*K)\}.
\]
For real symbols these are regular self-adjoint multipliers. Their bounded
transforms belong to the C*-multiplier algebra
\(M(C_0(T^*K))=C_b(T^*K)\); an unbounded polynomial does not itself
belong to that bounded multiplier algebra.  Factorwise on \(M_8\),
\[
 \operatorname{gr}\mathcal D_{\rm bi}(M_8)
 =\C[\rho_x,\rho_p,\nu_t,\nu_E].
\]
\end{theorem}
\begin{proof}
Bi-invariant differential operators are the center of the complex
universal enveloping algebra.  PBW identifies its associated graded with
the invariant polynomial algebra on the dual Lie algebra.  For
\(\mathfrak{su}(2)\), invariant polynomials are generated by the quadratic
norm; the central one-dimensional factor contributes its unrestricted
linear coordinate.  This proves the first three statements and the
inversion parity.  The tangent-groupoid construction identifies the
rescaled highest-order part of a differential operator with its constant
coefficient operator on the tangent fibre; fibrewise Fourier transform
turns that operator into multiplication by its principal symbol
\cite{Connes1994}.  Multiplication by a real continuous symbol has a dense
domain containing \(C_c(T^*K)\); multiplication by \((p\pm i)^{-1}\)
gives its bounded resolvents.  This proves the stated regular
self-adjoint interpretation, with bounded transform
\(p(1+p^2)^{-1/2}\in C_b(T^*K)\).  Taking the product of the four
independent factors gives the final formula.
\end{proof}

\begin{proposition}[The compact-double form and the reciprocal two-form]
\label{p1b:prop:relational-two-forms}
Normalize the invariant inner product used in the compact double by
\[
 \langle A,B\rangle_{\mathfrak{su}(2)}
   =-\frac1{8\pi^2}\Tr(AB),\qquad
 \langle i,i\rangle_{\mathfrak u(1)}=\frac1{4\pi^2},
\]
with the two summands orthogonal.  Let
\(H_v=i\sigma_3\), and on the common maximal torus write
\[
 g_x=e^{\alpha_xH_v},\qquad g_p=e^{\alpha_pH_v}.
\]
Then the compact-double form of Theorem~\ref{p1b:thm:g12-double} and the
reciprocal form of Theorem~\ref{p1b:thm:transgression} agree on this
four-torus:
\[
 \boxed{
 \omega_D\big|_{T^4}=\Omega\big|_{T^4}
 =\frac1{4\pi^2}
   (\dd\alpha_x\wedge\dd\alpha_p
    +\dd\theta_t\wedge\dd\theta_E).}
\]
They are nevertheless globally different structures.  The form \(\Omega\)
has rank eight everywhere, whereas \(\omega_D\) has rank six on the
trace-zero commutator locus.  Moreover
\[
 (\dd\omega_D)_{(e,e)}=0,
 \qquad (\dd\Omega)_{(e,e)}\ne0.
\]
Consequently there is neither a global equality \(\omega_D=\Omega\) nor a
closed two-form \(B_0\) with \(\Omega=\omega_D+B_0\).
\end{proposition}
\begin{proof}
On the common maximal torus all Maurer--Cartan forms commute.  The third
fusion term in \(\omega_D\) wedges a one-form with a scalar multiple of
itself and vanishes, while the first two terms add to
\(\langle H_v,H_v\rangle\dd\alpha_x\wedge\dd\alpha_p\); the scalar factor
is identical.  Since \(H_v^2=-I\) and \(\Tr(H_v^2)=-2\), the displayed
normalization gives \(\langle H_v,H_v\rangle=1/(4\pi^2)\).  The definition
of \(\beta\) in Theorem~\ref{p1b:thm:transgression} gives the same vector
coefficient, and \(\omega_{\rm sc}\) gives the scalar coefficient.

The rank statements are Theorems~\ref{p1b:thm:g12-double} and
\ref{p1b:thm:transgression}.  Since the commutator map has no linear term at
\((e,e)\), \(d\mu_{(e,e)}=0\), so
\(\dd\omega_D=-\mu^*\chi\) vanishes there.  By contrast the proof of
Theorem~\ref{p1b:thm:transgression} evaluates \(\dd\Omega\) nontrivially on
two vector directions in one sphere factor and one in the other; equivalently
choose \(A,B,C\in\mathfrak{su}(2)\) with
\(\Tr([A,B]C)\ne0\).  A closed difference would force equal differentials,
which is impossible.
\end{proof}

\begin{theorem}[Universal spectral evaluation of the scalar operator algebras]
\label{p1b:thm:relational-operator-spectrum}
For
\[
 \pi_{j_x,j_p,n,w}
 =V_{j_x}\otimes V_{j_p}\otimes z_t^n\otimes z_E^w,
\]
every \(P\in\mathcal D_{\rm bi}\) acts by the scalar
\[
 \boxed{
 P\bigl(4j_x(j_x+1),4j_p(j_p+1),n,w\bigr).}
\]
The resulting evaluation map from \(\mathcal D_{\rm bi}\) to functions on
\(\frac12\Z_{\ge0}^2\times\Z^2\) is injective.  On any finite regular
band with joint label set \(S\), its restriction has image exactly
\(\C^S\), the center of the corresponding block algebra
\(\bigoplus_{s\in S}M_{d_s}(\C)\).  In particular every block projector on
that finite carrier is a polynomial in the four commuting generators.
Hence a formally self-adjoint \(P\) is nonnegative on \(L^2(M_8)\) exactly
when every one of these scalar values is nonnegative.

For the Hopf-marked algebra, on a simultaneous Hopf-weight space with
integer charges \(q_x,q_p\),
\[
 \boxed{
 P\in\mathcal D_H\quad\Longrightarrow\quad
 P\mapsto
 P\bigl(4j_x(j_x+1),4j_p(j_p+1),q_x,q_p,n,w\bigr),}
\]
where \(|q_i|\le2j_i\) and \(q_i\equiv2j_i\pmod2\).  These evaluations are
again jointly faithful when all compatible spins and scalar characters are
retained.  Exchange acts on the full joint spectrum by
\((j_x,j_p,q_x,q_p,n,w)\mapsto(j_p,j_x,q_p,q_x,w,n)\).
\end{theorem}
\begin{proof}
The unit-sphere Casimir on spin \(j\) is \(4j(j+1)\), while the circle and
Hopf generators act by their integer weights.  The operator-algebra theorems
show that every operator is a polynomial in these mutually commuting
generators, proving the evaluations.  If a polynomial vanishes on the full
label set, first vary \(n,w\in\Z\), then vary the two infinite Casimir
sequences; repeated one-variable polynomial uniqueness forces every
coefficient to vanish.  On a finite label set, ordinary multivariate
interpolation produces a polynomial taking arbitrary prescribed values on
the distinct joint spectral points; this proves the finite-center statement
and the polynomial block projectors.  The same uniqueness argument works on
any fixed Hopf charge because the compatible spin sequence remains infinite.
Peter--Weyl gives a complete orthogonal joint eigenbasis, so nonnegativity is
equivalent to nonnegativity of all eigenvalues.  The exchange formula is the label action
of Theorem~\ref{p1b:thm:irreducibles} together with the matched Hopf marking.
\end{proof}

\begin{theorem}[Hopf-charge quotient and horizontal operator]
\label{p1b:thm:relational-charge-quotient}
Fix \((q_x,q_p)\in\Z^2\).  Sections of
\(L_{q_x,q_p}\to S_x^2\times S_p^2\) are equivalently functions on
\(P_v\) in the joint Hopf-weight sector
\(Q_x=q_x,Q_p=q_p\).  On the product with the independent scalar torus, the
restriction of the complete marked operator algebra factors faithfully as
\[
 \boxed{
 \mathcal D_H/(Q_x-q_x,Q_p-q_p)
 \cong\C[C_x,C_p,N_t,N_E].}
\]
The horizontal scalar Laplacian is the image of
\[
 \boxed{
 \Delta_H=\frac{C_x-Q_x^2}{R_x^2}
          +\frac{C_p-Q_p^2}{R_p^2}.}
\]
Consequently, on the degree labels \(k_i=2j_i\),
\[
 \lambda_H=
 \frac{k_x(k_x+2)-q_x^2}{R_x^2}
 +\frac{k_p(k_p+2)-q_p^2}{R_p^2},
\]
with \(k_i\ge|q_i|\) and \(k_i\equiv q_i\pmod2\), recovering exactly the
spectral floors of Theorem~\ref{p1b:thm:charge-bundles}.

If \(\pi_H:P_v\to S_x^2\times S_p^2\) is the bi-Hopf bundle, then
\(\pi_H^*L_{q_x,q_p}\) is canonically trivial as a complex line on
\(P_v\), and hence on \(M_8=P_v\times\T^2\).  Its nontrivial information on
the total space is the retained fibre equivariance.  This is distinct from
the scalar Poincar\'e line, whose class \(\eta\) is nonzero on \(M_8\).
Thus Hopf charges \((q_x,q_p)\) and scalar weights \((n,w)\) are independent
until an additional equivariance relation is explicitly imposed.
\end{theorem}
\begin{proof}
The associated-bundle construction identifies a section with an equivariant
function of the corresponding character.  Hence \(Q_i\) acts by \(q_i\), so
polynomial restriction factors through the displayed ideal.  Faithfulness
follows from the joint evaluation theorem because each fixed charge retains
infinitely many compatible spins and all scalar characters.  On a round
sphere the full Casimir splits into horizontal plus vertical squares; with
the present integer normalization the vertical square is \(Q_i^2\).  This
proves the operator and eigenvalue formulas.

The pullback of an associated line to its principal bundle has the canonical
trivialization
\((p,z)\mapsto(p,[p,z])\).  Therefore all base Chern classes disappear on
\(P_v\), in agreement with \(H^2(P_v;\Z)=0\), while the equivariance remains
part of the section model.  The scalar line instead has
\(c_1=\eta\ne0\) on the independent scalar torus factor.
\end{proof}

\begin{proposition}[Placement of scalar Poincar\'e-line operators]
\label{p1b:prop:relational-scalar-line-operators}
For \(L_q=\mathcal P_{\rm sc}^{\otimes q}\), the branch \(q=0\) is a
topologically trivial line; its flat connections give the ordinary Fourier
family, with possible constant holonomy shifts of the scalar generators.
For \(q\ne0\),
\[
 c_1(L_q)=q\eta\ne0,
\]
so no global bundle trivialization can turn its Bochner or twisted Dirac
operators into ordinary scalar differential operators on functions.  The
Landau and Dirac families of Section~\ref{p1b:sec:g13} are therefore a
genuinely twisted operator sector, separate from
\(\mathcal D_{\rm bi}=\C[C_x,C_p,N_t,N_E]\).  Exchange sends this sector
from degree \(q\) to degree \(-q\).
\end{proposition}
\begin{proof}
The zero-degree line is topologically trivial, and a flat connection differs
from the trivial one by a closed constant-holonomy class, producing Fourier
weight shifts.  If \(q\ne0\), a global bundle trivialization would force
\(c_1(L_q)=0\), contradicting primitivity of \(\eta\).  The exchange action
on \(\eta\) is odd by Proposition~\ref{p1b:prop:scalar-moments}, hence it
sends \(q\) to \(-q\).
\end{proof}

\begin{theorem}[Stable blindness, nonflatness and spectral scope of the unstable sector]
\label{p1b:thm:relational-unstable-sector}
Let \(E\to P_v\) be the nontrivial rank-two bundle of
Theorem~\ref{p1b:thm:g15-unstable}, and put
\(\widetilde E=\pi^*E\to M_8\).  Then
\[
 c_j(\widetilde E)=0\quad(j>0),\qquad
 [\widetilde E]-2=0\quad\text{in }K^0(M_8).
\]
Therefore every complex K-homology pairing is blind to the unstable class:
\[
 \boxed{
 \langle[\widetilde E]-2,x\rangle=0
 \qquad\text{for every }x\in K_0(M_8).}
\]
In particular, for every spin structure and Riemannian metric, and for
any unitary connection on \(\widetilde E\),
\[
 \operatorname{ind}D^+_{\widetilde E}=0.
\]
Nevertheless \(\widetilde E\) admits no flat complex connection.
Consequently its isomorphism class is not represented by any of the flat
bundle families of Proposition~\ref{p1b:prop:flat} or
Theorem~\ref{p1b:thm:equivariant-flat}.  Exchange fixes its isomorphism
class, but that fact alone does not produce a chosen involutive exchange
lift.  The independent construction of
Theorem~\ref{p1b:thm:g21-exotic-lift} supplies one for this particular
complex bundle.

A Laplace- or Dirac-type spectrum twisted by \(\widetilde E\) therefore
requires an additional nonflat connection.  The complete spectra in the
G08 family concern the bundles and operators explicitly declared there and
do not, without that extra connection datum, determine a spectrum for
\(\widetilde E\).
\end{theorem}
\begin{proof}
Pullback preserves the stable identity and the vanishing Chern classes from
Theorem~\ref{p1b:thm:g15-unstable}.  The K-homology pairing factors through
K-theory, proving the first boxed identity.  For the spin Dirac class,
\([\widetilde E]=2\) gives
\(\operatorname{ind}D^+_{\widetilde E}=2\operatorname{ind}D^+\); the
untwisted index is zero because the tangent bundle is trivial and the
positive-degree \(\widehat A\)-classes vanish.

Suppose a flat complex connection existed.  The universal cover is
\(S^3\times S^3\times\R^2\), with deck group \(\Z^2\) acting only in the
\(\R^2\) factor.  Its monodromy representation therefore constructs
\(\widetilde E\) as the pullback of a flat rank-two bundle \(F\to\T^2\).
Flatness gives zero real first Chern class; since
\(H^2(\T^2;\Z)\) is torsion-free, \(c_1(F)=0\).  Complex bundles of fixed
rank at least one over the two-dimensional torus are classified by their
first Chern class, so \(F\) is topologically trivial.  Its pullback would
then be trivial, contradicting Theorem~\ref{p1b:thm:g15-unstable}.  The
final spectral statement is a typing consequence: a twisted differential
operator is not defined until its connection is supplied.
\end{proof}

\begin{theorem}[Dolbeault Hodge realization and index bridge]
\label{p1b:thm:relational-dolbeault-hodge}
Let \(J\) be any member of Theorem~\ref{p1b:thm:complex-family} and let
\(g\) be any \(J\)-compatible Hermitian metric, including the compatible
bi-invariant metrics of Section~\ref{p1b:sec:g14}.  For every \((p,q)\),
\[
 \Delta_{\bar\partial}^{p,q}
 =\bar\partial\bar\partial^*+\bar\partial^*\bar\partial
\]
is elliptic and
\[
 \boxed{
 \ker\Delta_{\bar\partial}^{p,q}
 \cong H_{\bar\partial}^{p,q}(M_8,J).}
\]
Hence the zero-mode generating polynomial of the Dolbeault Laplacians is
exactly the polynomial \(P_J(z,w)\) in
Theorem~\ref{p1b:thm:g16-hodge}; in particular it has only the two curvature
rank strata \(\rho=1,2\).

If \(J\) and \(g\) are exchange invariant, pullback by exchange commutes
with \(\Delta_{\bar\partial}\), and the exchange trace on its zero modes is
\[
 \boxed{(1-w^2)(1-z^4w^2).}
\]
Thus the even and odd zero-mode multiplicities are exactly the coefficients
of
\[
 \frac12\{P_J(z,w)\pm(1-w^2)(1-z^4w^2)\}.
\]
Finally, for every fixed holomorphic degree \(p\), the Dolbeault--Dirac
operator graded by \(q\)-parity has index
\[
 \boxed{
 \chi^p(M_8,J)=\sum_q(-1)^q h^{p,q}_{\bar\partial}=0.}
\]
\end{theorem}
\begin{proof}
The first statement is the Hermitian Hodge theorem on the compact complex
manifold.  The Hodge numbers and exchange trace are already computed in
Theorem~\ref{p1b:thm:g16-hodge} and
Proposition~\ref{p1b:prop:g16-frolicher}; an exchange-invariant Hermitian
metric makes exchange a holomorphic isometry, so it commutes with
\(\bar\partial\), its adjoint and the Laplacian.  In both G16 strata the
Hodge polynomial contains the factor \((1+w)^2\).  Setting \(w=-1\)
therefore gives \(P_J(z,-1)=0\), coefficientwise in \(z\), which is exactly
the displayed Dolbeault index.  This also agrees with the smooth triviality
of \((TM_8,J)\), whose positive-degree Chern classes vanish.
\end{proof}

\begin{proposition}[Exact operator-exhaustion boundary]
\label{p1b:prop:relational-operator-scope}
The operator results of this paper have the following exact completeness
scope.
\begin{center}
\begin{tabular}{@{}p{.38\textwidth}p{.53\textwidth}@{}}
\toprule
\textbf{Operator class} & \textbf{Status in this paper}\\
\midrule
Finite-order bi-invariant scalar differential operators on \(M_8\) &
Exhausted by \(\C[C_x,C_p,N_t,N_E]\), with faithful joint spectral evaluation.\\
Finite-order matched Hopf-marked scalar differential operators &
Exhausted by \(\C[C_x,C_p,Q_x,Q_p,N_t,N_E]\); every fixed-charge sector is the quotient of Theorem~\ref{p1b:thm:relational-charge-quotient}.\\
Relative bi-invariant and inversion-even scalar operators &
Exhausted by \(\C[C_{\rm rel},N_{\rm rel}]\) and its even subalgebra \(\C[C_{\rm rel},N_{\rm rel}^2]\), with principal symbols given by Theorem~\ref{p1b:thm:relational-tangent-symbol}.\\
Finite retained carriers &
Their full endomorphism algebras, symmetry commutants, generated coordinate algebras and exact isometric transports are exhausted on the declared carriers.\\
Reciprocal variable-Cartan Hodge--Dirac family on \(S_R^3\) &
For reciprocal-odd scalar multipliers with \(\frac{\sqrt3}{2}\|u\|_\infty<1\), the protected zero line is constructively exhausted by Theorem~\ref{p1b:thm:dirac-zero-resolvent}; its second-order transverse onset and complete dipole--shell generation of the allowed reciprocal-odd coexact vector shells are given by Theorem~\ref{p1b:thm:dirac-transverse-generation}. This does not classify arbitrary bundle Dirac perturbations.\\
Hodge, spin, TT, Rarita, covariant-Casimir, Landau and overlap families &
Their proved spectra, kernels, heat traces, locality or functional-calculus conclusions apply to the particular operators and bundles explicitly defined in their theorems; this is not a classification of every invariant bundle differential operator.\\
Arbitrary twists such as \(\widetilde E\), arbitrary pseudodifferential operators, and nonlocal operators &
Outside the universal operator classifications above unless extra bundle/connection or symbol data are explicitly supplied.\\
\bottomrule
\end{tabular}
\end{center}
Accordingly, the word \emph{complete} applied to an operator statement means
complete in its named operator class.  It does not assert that every natural
operator on every bundle over \(M_8\) belongs to one universal algebra.
\end{proposition}
\begin{proof}
The first three rows are Theorems~\ref{p1b:thm:bi-operators},
\ref{p1b:thm:hopf-operators},
\ref{p1b:thm:relational-operator-spectrum},
\ref{p1b:thm:relational-charge-quotient}, and
\ref{p1b:thm:relational-tangent-symbol}.  The finite-carrier row is the
finite double-commutant and image-transport calculus developed in the finite
and all-level sections.  The named bundle rows each fix their bundle,
connection and differential expression before deriving the spectrum.
Theorem~\ref{p1b:thm:relational-unstable-sector} exhibits a bundle for which
those extra data are not supplied, proving why the broader interpretation
would be false.
\end{proof}


\subsubsection{Relative geometry-saturation basis and completed rank-two classification}
\label{p1b:sec:g25}

No finite paper can assert exhaustion of every geometric structure one
could place on the underlying smooth eight-manifold.  The saturation claim
used here is therefore relative to the explicit search basis below.

\begin{theorem}[Paper-I-B geometry-saturation statement]
\label{p1b:thm:g25-saturation}
Let \(\mathscr B_{\rm geom}\) consist of the following declared families on
the labelled reciprocal product and its canonical hosts:
\begin{enumerate}[label=\textup{(\roman*)}]
\item compact Lie-group, subgroup, quotient, fixed-locus and homogeneous
      geometry;
\item integral topology, cohomology, spin data, complex K-theory and
      characteristic classes;
\item invariant Riemannian metrics, Levi--Civita and metric connections,
      torsion, curvature, holonomy and geodesic operations;
\item Hopf, contact, CR, basic-cohomology, prequantization and harmonic-map
      structures generated by the marked Hopf fibrations;
\item symplectic/LCS obstructions and all left-invariant integrable complex,
      Hermitian, Bismut, Dolbeault and Fr\"olicher structures up to the
      automorphism saturation of Theorem~\ref{p1b:thm:g22-global-complex};
\item smooth complex bundles in the stable range, exact rank-two existence
      for both determinant parities, fixed-Chern principal-\(SU(2)\)
      multiplicities and the complete ordinary complex determinant quotient, rank
      three, the G15 unstable sector, flat equivariant bundles, Born-fixed
      principal-\(SU(2)\) classes and complete ordinary complex Born-fixed counts, the verified
      equivariant-index obstruction, principal-\(SU(2)\) descent fibres and
      the explicit mixed G15 \(U(2)\) lift;
\item finite binary-polyhedral quotients, invariant rings, reciprocal
      McKay/ADE products and their declared finite spectral restrictions;
\item the reciprocal quaternionic cone host, global and independent
      rescalings, full product isometries, and the associated geometric
      spectral/heat scaling laws.
\end{enumerate}
Relative to \(\mathscr B_{\rm geom}\), all branches above are
discharged by explicit owner theorems. The genuine mixed-\(U(2)\)
classification retains its prescribed fixed-locus eigenranks and the
actual boundary model. Its four compatible free-region rows have
evaluated full and auxiliary gauge images, quotient torsors,
fixed-Chern cardinalities and ordinary forgetful fibres in
Theorem~\ref{p1b:thm:complete-mixed-rank-two}.
Theorem~\ref{p1b:thm:natural-rank-two-completion} supplies the original
natural boundary coefficients and every retained named extension test.
This completes the numerical rank-two branch of the displayed search
basis; it does not enlarge that basis.
No claim is made about arbitrary non-left-invariant complex structures,
arbitrary nonhomogeneous metrics or connections, or geometric structures
not generated by this declared basis.
\end{theorem}
\begin{proof}
Items (i)--(iv) are discharged by the group/quotient, topological, metric,
Hopf and G17/G19 atlases; item (v) by the cohomological obstructions,
Theorem~\ref{p1b:thm:g22-global-complex}, G14 and G16; item (vi) by the
stable bundle theorem, G15, Theorems~\ref{p1b:thm:g20-rank2-existence}--
\ref{p1b:thm:g20-rank3-fibre}, and the G21 Born-equivariant theorems; item
(vii) by G10 and G23; and item (viii) by G18, G24 and the spectral sections.
The mixed branch is discharged by the actual coaction and higher-gauge
proofs, the post-Batch-28 image and sphere calculations, the homogeneous
ordinary fibre action, and the complete natural paired-map evaluation.
Those proofs retain nontrivial fixed-locus isotropy rather than replacing
the classification by bundles on the coarse quotient. The stated
exclusions retain the distinction between this relative completion and
absolute exhaustion of arbitrary structures on the smooth manifold.
\end{proof}


\begin{corollary}[Invariant connections on a supplied exchange-equivariant bundle]
\label{p1b:cor:v214-invariant-unstable-connection}
Let $E\to M_8$ carry a unitary lift $\tau$ of exchange with $\tau^2=I$.
Then $E$ admits an exchange-invariant unitary connection. If $E$ is the
nontrivial unstable pullback of Theorem~\ref{p1b:thm:g15-unstable}, every
such connection is nonflat. No lift of the simply transitive left-$M_8$
action exists on this nontrivial bundle.
\end{corollary}
\begin{proof}
Start with any unitary connection $\nabla_0$. Pullback by the lift gives an
affine involution $\mathcal B$ on the affine real space of unitary
connections. Thus $\nabla_B=(\nabla_0+\mathcal B\nabla_0)/2$ is unitary
and fixed. For an exchange-invariant base metric (obtained, if needed,
by averaging any metric with its exchange pullback), its covariant Laplacian
commutes with the induced unitary exchange operator. Flatness would contradict
Theorem~\ref{p1b:thm:relational-unstable-sector}. Finally a lift of the full
left action would trivialize any bundle over the freely transitive group:
fixing a fibre at the identity, $(g,v)\mapsto g\cdot v$ is a global
equivariant trivialization. The nontrivial unstable class excludes this.
\end{proof}

\subsection{The mathematical exhaustion catalogue}
\label{p1b:sec:exhaustion}
The preceding results give finite descriptions of the following complete
families; the parameters range over their entire displayed domains.
This is the reciprocal counterpart of specifying \(\mathcal H_m\),
specializing \(m=1\), and retaining its completion cone in Paper I-B.
\begin{longtable}{@{}p{.42\textwidth}p{.51\textwidth}@{}}
\toprule
Class being exhausted&Classification or exact description\\ \midrule
\endhead
Continuous characters and irreducible representations&
\(\Z^2\), and two sphere spins with two integer charges
(Theorems~\ref{p1b:thm:group}, \ref{p1b:thm:irreducibles})\\
Connected closed subgroup types and hypersubgroups&
Semisimple skeleton plus rational centralizer subtori; every connected
codimension-one subgroup is \(D_6\times S^1_{(m,n)}\)
(Theorem~\ref{p1b:thm:connected-subgroup-skeleton},
Corollary~\ref{p1b:cor:codim-one-normal})\\
Primitive maximal-torus circle quotients&
Complete \(d=0,1,>1\) scalar-projection trichotomy, low homotopy and
exchange-stable primitive slopes
(Theorem~\ref{p1b:thm:primitive-circle-quotients},
Corollary~\ref{p1b:cor:exchange-stable-slopes})\\
Primitive pure-Hopf contact/prequantization geometry&
Residual circle connection and Chern class, sharp contact criterion and
regular prequantization branch
(Theorem~\ref{p1b:thm:g17-contact},
Corollary~\ref{p1b:cor:g17-prequantization})\\
Finite connected covers&
Finite-index sublattices of \(\Z^2\)
(Theorem~\ref{p1b:thm:covers})\\
Integral cohomology and stable complex bundle classes&
Exterior rings, exchange lattices and integral Chern conditions
(Theorems~\ref{p1b:thm:cohomology}, \ref{p1b:thm:k-theory},
\ref{p1b:thm:chern})\\
Spin, string and flat unitary data&
Discrete torsors, joint holonomy spectra, and exchange blocks
(Theorems~\ref{p1b:thm:spin}, \ref{p1b:thm:equivariant-flat})\\
Invariant metrics&
The six full cones of Theorem~\ref{p1b:thm:metric-cones}\\
Left-invariant and bi-invariant metric connections&
Affine tensor spaces of dimensions \(224\) and \(8\), with their
exchange and alternating-torsion subspaces
(Theorems~\ref{p1b:thm:left-metric-connections},
\ref{p1b:thm:bi-metric-connections})\\
Invariant metric-connection operations&
Torsion--contorsion isomorphism, irreducible torsion split, finite curvature
and covariant-derivative calculus, and skew-torsion geodesic invariance
(Proposition~\ref{p1b:prop:g19-torsion},
Theorem~\ref{p1b:thm:g19-curvature},
Corollary~\ref{p1b:cor:g19-geodesics})\\
Principal connections on the relative bundles&
The full affine space of adjoint-valued one-forms;
the homogeneous family and its exchange-fixed member
(Theorems~\ref{p1b:thm:all-principal-connections},
\ref{p1b:thm:homogeneous})\\
Hopf-invariant almost complex structures&
All are integrable; explicit complex-structure parameter spaces and
exchange-fixed parts
(Theorem~\ref{p1b:thm:complex-family})\\
Bi-invariant and Hopf-marked scalar differential operators&
Generator-and-relation algebras, faithful joint spectral evaluation,
charge-sector quotients and principal-symbol degeneration
(Theorems~\ref{p1b:thm:bi-operators}, \ref{p1b:thm:hopf-operators},
\ref{p1b:thm:relational-operator-spectrum}--\ref{p1b:thm:relational-charge-quotient})\\
Cross-structure relational closure&
Curvature--commutator, tangent-symbol, two-form, charge-quotient,
unstable-bundle and Dolbeault-operator bridges, with exact operator scope
(Section~\ref{p1b:sec:relational-closure})\\
Product scalar and form spectra&
All eigenvalues, multiplicities and exchange heat sectors for every
bi-invariant metric
(Theorems~\ref{p1b:thm:product-hodge},
\ref{p1b:thm:exchange-heat})\\
Finite regular scalar bands&
All finite Peter--Weyl index sets, with exact exchange criterion and
charge-restricted dimensions
(Theorem~\ref{p1b:thm:finite-band-classification})\\
Factorwise and full equivariant level-one completions&
\(\R_{\ge0}^6\), exchange-fixed diagonal \(\R_{\ge0}^3\), and
the full 48-edge coordinate atlas with 16 independent cycle phases
(Theorems~\ref{p1b:thm:product-completion},
\ref{p1b:thm:full-product-coordinate-atlas})\\
Arbitrary independent scalar and enriched levels&
All scalar--scalar, enriched--enriched and mixed products, their shells,
parity, algebras, fusion and maps
(Theorems~\ref{p1b:thm:mixed-dimensions}--\ref{p1b:thm:all-level-chiral})\\
Unequal-level exchange and invariant comparisons&
Exact common subspaces, minimal exchange completions, all singlet and
Hopf-charge ranks
(Theorem~\ref{p1b:thm:unequal-exchange},
Theorem~\ref{p1b:thm:mixed-diagonal},
Proposition~\ref{p1b:prop:all-level-hopf})\\
All-level coordinate completions&
Clebsch--Gordan graphs, every support stratum, exchange quotients,
rectangle commutators and generated matrix algebras
(Theorems~\ref{p1b:thm:all-level-graph}--\ref{p1b:thm:all-level-generated-coordinates})\\
Fixed-factor spectra and increasing-level limits&
All specified Casimir towers, exact scalar coordinate boundary identities,
word frontiers and one- or two-factor limits with explicit radius dependence
(Theorems~\ref{p1b:thm:mixed-spectra}--\ref{p1b:thm:all-level-limits})\\
Finite restrictions and sampling at arbitrary levels&
Character formulas, distinct exceptional onsets, homogeneous-bundle
realizations, positive cubature and sharp comparison ranges
(Theorems~\ref{p1b:thm:all-level-finite-restriction},
\ref{p1b:thm:all-level-sampling},
Corollary~\ref{p1b:cor:all-level-design-comparison})\\
Finite sampled presentations&
Exact cubature domains and full algebra, action, tensor and
measure transport
(Theorems~\ref{p1b:thm:product-cubature}--%
\ref{p1b:thm:blocking-transport})\\
\bottomrule
\end{longtable}

Each classification retains its equivalence relation: metric cones are
spaces before quotienting by automorphisms, connection spaces are before
gauge quotient unless stated, completion cones are modulo sector phases,
and flat connections are modulo gauge. Their dimensions therefore measure
different precisely specified objects. Longer computations for particular
members and the additional function and measure families follow below.

At the degree-two action-algebra comparison, the two inherited
normalized traces are related by the unique positive central
density of Corollary~\ref{p1b:cor:degree-two-trace-density}.

The additional product families and their exact scope are as follows.
\begin{longtable}{@{}p{.42\textwidth}p{.51\textwidth}@{}}
\toprule Product family & Result and mathematical domain\\ \midrule\endhead
Coordinate ancestry and retained geometric squares & Two-anchor coordinate grammar, retained geometric square and isometric transport (Proposition~\ref{p1b:prop:g02-product})\\
Product commutants, Clifford algebras and Casimir filtrations & Specified product or diagonal group, coordinate supports and Clifford action (Proposition~\ref{p1b:prop:g03-product})\\
Hurwitz, trimer, Walsh and resolution structures & Actual scalar, trimer and Walsh carriers, resolution spaces/windows, and rooted vector/bivector overlaps and tensor products (Propositions~\ref{p1b:prop:g04-product}, \ref{p1b:prop:trimer-transport-overlap}; Corollary~\ref{p1b:cor:product-trimer-overlap})\\
Contact pairs and product control geometry & Marked orthogonal horizontal distribution, including scalar torus directions (Proposition~\ref{p1b:prop:g05-product})\\
CR complexes, Hardy spaces and bi-Hopf quantization & Specified CR complex, integer twists and compatible graded exchange (Proposition~\ref{p1b:prop:g06-product})\\
Global metric geometry and deformations & Positive round/Berger product metrics and the stated invariant extensions (Proposition~\ref{p1b:prop:g07-product})\\
Specified bundle spectra, constraints and analytic invariants & Complete multiplicities for the declared form, spin, TT and intrinsic Rarita operators; arbitrary twisted bundles require their own connection data (Proposition~\ref{p1b:prop:g08-product}, Theorem~\ref{p1b:thm:relational-unstable-sector})\\
Reproducing kernels, evaluation measures and exact moments & Finite scalar or fiber evaluation, exact moments and stated limit topology (Proposition~\ref{p1b:prop:g09-product})\\
Finite quotients and marked reference geometry & Specified finite actions, smooth/orbifold branches and marked reference maps (Proposition~\ref{p1b:prop:g10-product})\\
\bottomrule\end{longtable}
Completeness refers to these declared constructions and parameter domains. In particular, the product Paneitz definitions do not assert a universal higher-dimensional version of the one-factor rigidity theorems.

The following developments complete six further specified mathematical
families. Their full hypotheses are part of the classifications.
\begin{longtable}{@{}p{.42\textwidth}p{.51\textwidth}@{}}
\toprule Family & Result and scope\\ \midrule\endhead
Cartan cylinder functional & Endpoint transgression, unit-charge minimum
and homogeneous fluctuations on the specified infinite cylinder
(Paper~I-D, Section~\ref{p1b:app:g11})\\
Compact double & Group-valued moment identities, complete pointwise
rank criterion, and the two central-value reductions
(Theorem~\ref{p1b:thm:g12-double},
Proposition~\ref{p1b:prop:g12-reduction})\\
Scalar Poincar\'e line powers & All constant-curvature connections on a
flat scalar torus, theta bases, translations, scalar and spin spectra
(Theorem~\ref{p1b:thm:g13-landau},
Proposition~\ref{p1b:prop:g13-dirac})\\
Closed-torsion Hermitian structures & Bi-invariant compatible metric
family, flat Bismut connection, torsion periods and twisted generalized
Kähler construction (Theorems~\ref{p1b:thm:g14-bismut}
and~\ref{p1b:thm:g14-generalized})\\
Specified unstable bundle sector & The complete \(SU(2)\) bundle family on the
canonical derived core \(D_6=[M_8,M_8]\cong S^3\times S^3\), its pullbacks,
and stable classification for complex rank at least four on \(M_8\)
(Theorem~\ref{p1b:thm:g15-unstable},
Proposition~\ref{p1b:prop:g15-stable-range})\\
Coherent Born lifts in the unstable rank-two sector&
Gauge norm criterion, exact Born-fixed \(SU(2)\) counts, verified quaternionic mod-twenty-four
index obstruction, explicit mixed G15 lift, verified principal-\(SU(2)\)
no-go for the exotic class, and the exact principal descent-fibre reduction (individual fibres unevaluated)
(Proposition~\ref{p1b:prop:g21-norm},
Theorems~\ref{p1b:thm:g21-born-fixed-count}--\ref{p1b:thm:g21-principal-descent})\\
Full marked Dolbeault family & Both curvature-rank strata, complete
Hodge algebras, Frölicher differentials and exchange sectors, together with
the exact derived-subgroup complex locus
(Theorems~\ref{p1b:thm:g16-model}--\ref{p1b:thm:g16-hodge},
Proposition~\ref{p1b:prop:g16-frolicher},
Corollary~\ref{p1b:cor:g16-derived-complex})\\
Aligned exceptional Hopf--Clifford incidence & Exact \(2T,2O,2I\)
Clifford-intersection, Hopf-image, core-fibre and Heegaard counts for the
declared alignment (Proposition~\ref{p1b:prop:exceptional-hopf-clifford})\\
\bottomrule\end{longtable}

The canonical carrier of the G15 unstable sector is \(P_v=D_6=[M_8,M_8]\). Its exhaustive two-class family is indexed by \(\epsilon_{\mathrm{IB15}}\in\mathbb Z/2\) without selecting a member. The restriction bit is a bundle-class index and selects no member.

\begin{longtable}{@{}p{.42\textwidth}p{.51\textwidth}@{}}
\toprule Family & Result and scope\\ \midrule\endhead
Reciprocal double cone & Flat hyper-K\"ahler host, triholomorphic Born exchange, fixed diagonal and compact/noncompact type firewall (Theorem~\ref{p1b:thm:g18-double-cone}--Corollary~\ref{p1b:cor:g18-firewall})\\
Actual low-rank complex bundles & Exact rank-three existence and fixed-Chern fibre; actual rank-two existence and ordinary fixed-Chern multiplicities for both determinant parities; determinant-preserving and ordinary counts retain separate meanings (Theorems~\ref{p1b:thm:g20-rank2-existence}--\ref{p1b:thm:g20-rank3-fibre}, \ref{p1b:thm:v2170-ordinary-counts})\\
Born-equivariant rank-two closure & Exact Born-fixed principal-\(SU(2)\) multiplicities, verified equivariant-index obstruction, mixed G15 lift, verified exotic principal no-go and the exact principal descent-fibre reduction (individual fibres unevaluated) (Theorems~\ref{p1b:thm:g21-born-fixed-count}--\ref{p1b:thm:g21-principal-descent})\\
All left-invariant integrable complex structures & Every invariant complex structure is compact-automorphism conjugate into the marked family; the Dolbeault atlas is therefore automorphism-saturated (Theorem~\ref{p1b:thm:g22-global-complex}, Corollary~\ref{p1b:cor:g22-transfer})\\
Exceptional product McKay--ADE geometry & Tensor-product invariant rings, factorized bigraded Molien series, product Kleinian singular locus and Cartesian-product affine ADE graph (Theorems~\ref{p1b:thm:g23-invariants}, \ref{p1b:thm:g23-mckay})\\
Isometry and scaling atlas & Full round-factor/flat-torus isometry group, exact Born-isometry criterion and dimension-eight scaling weights (Theorems~\ref{p1b:thm:g24-isometry}, \ref{p1b:thm:g24-scaling})\\
Relative geometry saturation & Explicit search basis; ordinary numerical rank-two classification is discharged.  Exact mixed-\(U(2)\) counts, quotient torsors, ordinary fibres and natural extension normalizations are evaluated on the declared rank-two domain (Theorems~\ref{p1b:thm:complete-mixed-rank-two} and \ref{p1b:thm:natural-rank-two-completion}); relative completion retains the displayed search basis (Theorem~\ref{p1b:thm:g25-saturation})\\
\bottomrule\end{longtable}


\subsubsection{The protected native Cartan family}
\label{p1b:sec:native-cartan-premises}
The data here are a round sphere \(S_R^3\), its Hodge operators and
the stated real smooth Cartan twisting profile. They constitute a
mathematical operator family. The Hodge identification, complete
constant-profile spectrum and protected kernel are proved here directly;
later represented realizations may apply these results without being
inputs to their proofs. All operators act on
\(H^1(S_R^3;\Lambda^0\oplus\Lambda^1)_\C\) in the natural
\(L^2\) space, and all stated involutions preserve that domain.
\begin{theorem}[Midpoint Hodge model and constant-parameter spectrum]
\label{p1b:thm:dirac-spectrum}
Couple spinors on \(S_R^3\) to the defining doublet with
\(A_f=f\Theta\). At \(f=\frac12\), the tensor spin bundle identifies
with \((\Lambda^0\oplus\Lambda^1)_\C\), and the Dirac operator is
\[
 B_R=\begin{pmatrix}0&\dd^*\\\dd&-*\dd\end{pmatrix},
 \qquad B_R^2=\Delta_0\oplus\Delta_1,\qquad
 \ker B_R=\C(1,0).
\]
Writing \(u=2f-1\) and \(\delta=f-\frac12\),
\[
 D_f=B_R+\frac{u}{2R}\diag(-3,1,1,1).
\]
Its eigenvalues multiplied by \(R\) are
\[
 \begin{array}{c|c}
 \text{value}&\text{multiplicity}\\\hline
 -3\delta&1\\
 -\delta\pm\sqrt{n(n+2)+4\delta^2}&(n+1)^2\text{ each},\ n\ge1\\
 \delta\pm(n+1)&n(n+2)\text{ each},\ n\ge1 .
 \end{array}
\]
At the midpoint its nonzero gap is \(\sqrt3/R\).
For \(0\le f\le1\), the sole zero crossing is the simple branch
\((3/2-3f)/R\); upward-positive spectral flow is \(-1\).
At the midpoint the eta invariant is zero and reduced eta is \(1/2\).
\end{theorem}
\begin{proof}
Identify the twisting doublet at the midpoint with the dual spin
doublet using its invariant alternating form.  The tensor connection
is then the induced connection on \(\operatorname{End}(\mathbb S)\).
The pointwise unitary identification is
\[
 (\phi,\alpha)\longmapsto
       \frac{\phi I+c(\alpha)}{\sqrt2},
 \qquad c(e^a)=i\sigma_a ,
\]
where the matrix side has its Hilbert--Schmidt metric.  In a normal
orthonormal frame the identity
\(c(e^a)c(e^b)=-\delta_{ab}I-\epsilon_{abc}c(e^c)\)
gives
\[
 D(\phi I+c(\alpha))
   =(d^*\alpha)I+c(d\phi-*d\alpha).
\]
This proves the Hodge matrix without a reference to a later represented
operator.  The identities \(d^2=0\), \(d^**d=0\) and
\((*d)^2=\Delta_1\) on coclosed one-forms give its square.
Hodge theory and \(H^1(S^3;\mathbb C)=0\) give its kernel.

Here is also a direct spectral enumeration.  Restrictions of homogeneous
harmonic polynomials of degree \(n\) on \(\mathbb R^4\) have eigenvalue
\(b_n=n(n+2)/R^2\) and dimension
\(\binom{n+3}{3}-\binom{n+1}{3}=(n+1)^2\).
Their gradients give the exact one-forms.  The coexact vector harmonics
are the tangential divergence-free harmonic vector polynomials of degree
\(n\ge1\); their two irreducible \(SU(2)\times SU(2)\) summands have
spins \(((n+1)/2,(n-1)/2)\) and the reverse spins.
Each has dimension \(n(n+2)\).  The Gauss formula gives rough
eigenvalue \([n(n+2)-1]/R^2\); adding the Ricci term \(2/R^2\)
gives Hodge eigenvalue \((n+1)^2/R^2\).
On these two summands curl has opposite values
\(\pm(n+1)/R\): its square is the Hodge Laplacian, and the
orientation-reversing inversion interchanges the summands.
The scalar, exact and these coexact spaces exhaust one-forms by the
orthogonal Hodge decomposition and the scalar/vector harmonic
polynomial decomposition.

Changing \(f\) from \(1/2\) adds \(4\delta S\cdot I/R\).
The spin--isospin singlet and triplet have \(4S\cdot I=-3,1\),
respectively.  Thus this is the displayed diagonal multiplier.
On the orthonormal pair \((Y,0),(0,dY/\sqrt{b_n})\), after
multiplication by \(R\), the operator is
\[
 \begin{pmatrix}-3\delta&\sqrt{n(n+2)}\\
                 \sqrt{n(n+2)}&\delta\end{pmatrix}.
\]
Its two roots and the coexact roots give the table, including every
multiplicity.  For \(|\delta|\le1/2\), none of these nonconstant
roots vanishes, whereas the constant singlet is \(-3\delta/R\).
It crosses downwards once.  At \(\delta=0\) every nonzero root is
paired with its negative; hence the eta function is zero in its
initial convergence half-plane and by continuation at zero.
The one-dimensional kernel then gives reduced eta \(1/2\).
\end{proof}

\begin{theorem}[Reciprocal involution and a protected simple kernel]
\label{p1b:thm:dirac-reciprocity}
Let \(\iota(q)=q^{-1}\) and
\[
 U(\phi,\alpha)=(\iota^*\phi,-\iota^*\alpha).
\]
Then \(U^2=I\), \(UB_RU=-B_R\), and for constant \(f\),
\(UD_fU=-D_{1-f}\).
If a bounded self-adjoint \(V\) satisfies \(UVU=-V\), the odd
Fredholm block of \(B_R+V\), from the \(U=+1\) space to the
\(U=-1\) space, has index one.
If \(\|V\|<g=\sqrt3/R\), its kernel is exactly one-dimensional
and its other eigenvalues have absolute value at least \(g-\|V\|\).
For a smooth real function \(u\) with \(u(q^{-1})=-u(q)\),
the perturbation above satisfies this hypothesis; if \(\|u\|_\infty\le1\),
its nonzero gap is at least
\[
 \frac{\sqrt3-\frac32\|u\|_\infty}{R}>0.
\]
Its graded heat trace is \(\Tr(Ue^{-tD_u^2})=1\), \(t>0\).
\end{theorem}
\begin{proof}
Inversion is an orientation-reversing round isometry.  It commutes
with \(d,d^*\) and changes the sign of the Hodge star in the
curl expression.  Applying this to the two entries of \(U\)
proves \(UB_RU=-B_R\).  The diagonal singlet/triplet multiplier
commutes with \(U\) when its coefficient is constant, which gives
the constant-\(f\) reflection formula.  An inversion-odd coefficient
instead makes \(M_u\) odd.

With \(\mathcal H_\pm=\ker(U\mp I)\), write the self-adjoint odd
operator as
\(\left(\begin{smallmatrix}0&B_-\\B_+&0\end{smallmatrix}\right)\).
Its only kernel vector is the positive-grade constant.  Thus
\(\ker B_+=\mathbb C\), \(\ker B_-=0\), and
\(\operatorname{ind}B_+=1\).
As a map \(H^1\to L^2\), multiplication by a bounded coefficient
is relatively compact, by the compact Sobolev inclusion.
The odd Fredholm index is therefore unchanged along \(B_R+sV\).

There is a useful stronger estimate than a contour argument:
for \(y\in H^1\cap\mathcal H_-\),
\[
 \|(B_-+V_-)y\|\ge(g-\|V\|)\|y\|.
\]
Thus this adjoint block is injective with closed range.
Its adjoint block is surjective, and its index forces its kernel
to have dimension one.  The nonzero singular values of the two
blocks agree, proving the nonzero gap \(g-\|V\|\) throughout
\(\|V\|<g\), including the range beyond \(g/2\).
Here \(\|M_u\|=3\|u\|_\infty/(2R)\).
Finally \(D\) pairs every positive- and negative-grade eigenspace
of \(D^2\) at a nonzero eigenvalue.  Their graded traces cancel
in the trace-class heat operator, leaving precisely its index one.
\end{proof}

\begin{theorem}[Constructive reciprocal protected zero line]
\label{p1b:thm:dirac-zero-resolvent}
Let \(u\in C^\infty(S_R^3,\R)\) be real and reciprocal odd,
\(u\circ\iota=-u\), and write
\[
 M_u=\frac{u}{2R}\diag(-3,1,1,1),\qquad D_u=B_R+M_u.
\]
Let \(\psi_0=(1,0)\), let \(P\) be the orthogonal projector onto
\(\C\psi_0=\ker B_R\), put \(Q=I-P\), and define the reduced inverse
\[
 G=Q(B_R|_{\operatorname{ran}Q})^{-1}Q.
\]
If
\[
 \rho:=\frac{\sqrt3}{2}\|u\|_\infty<1,
\]
then the unique zero mode normalized by \(P\psi_u=\psi_0\) is
\[
 \boxed{\psi_u=(I+GM_u)^{-1}\psi_0
 =\sum_{n=0}^{\infty}(-GM_u)^n\psi_0.}
\]
The series converges in norm and its order-\(N\) truncation satisfies
\[
 \left\|\psi_u-\sum_{n=0}^{N}(-GM_u)^n\psi_0\right\|
 \le \frac{\rho^{N+1}}{1-\rho}\|\psi_0\|.
\]
Every term has positive reciprocal grade. Thus this formula gives a
constructive representative of the protected zero line throughout the
stated simple-kernel domain.
\end{theorem}
\begin{proof}
The reduced inverse has norm \(R/\sqrt3\), so
\(\|GM_u\|\le\rho<1\).  Its Neumann inverse gives the displayed
convergent series and the geometric remainder bound.
Both \(G\) and \(M_u\) reverse reciprocal grade, so \(GM_u\)
preserves it and every term lies in positive grade.
Also \(PG=0\), hence \(P\psi_u=\psi_0\).
Multiplying \((I+GM_u)\psi_u=\psi_0\) by \(B_R\) gives
\(Q(B_R+M_u)\psi_u=0\).
The residual \(PM_u\psi_u\) vanishes because \(M_u\psi_u\)
has negative grade and \(\operatorname{ran}P\) has positive grade.
Consequently this is a zero mode.  The preceding theorem proves
uniqueness; elliptic regularity makes it smooth for smooth \(u\).
\end{proof}

\begin{theorem}[Graded Feshbach selection and cubic error]
\label{p1b:thm:graded-feshbach-selection}
Let \(D_0\) be self-adjoint with a domain-preserving unitary involution
\(U\), \(UD_0U=-D_0\), and an isolated finite-rank kernel projector
\(P\) entirely in the positive grade, \(UP=P\). Put \(Q=I-P\)
and assume the nonzero gap is \(g>0\). For a bounded self-adjoint
\(M\) with \(UMU=M\), define \(D_\eta=D_0+\eta M\) and its
Schur operator
\[
 K(\eta,z)=P(D_\eta-z)P
 -PD_\eta Q[Q(D_\eta-z)Q]^{-1}QD_\eta P .
\]
On the resolvent domain,
\[
 K(\eta,z)=-K(-\eta,-z),\qquad
 PMGMP=0,\quad G=Q(D_0|_Q)^{-1}Q.
\]
Writing \(v=|\eta|\|M\|<g/2\), for \(|z|\le v\) one has
\[
 \boxed{\|K(\eta,z)-(\eta PMP-zP)\|
 \le\frac{v^2(v+|z|)}{g(g-v-|z|)}
 \le\frac{2v^3}{g(g-2v)}.}
\]
In particular \(K(\eta,0)\) is odd in \(\eta\). A uniquely isolated
simple analytic low eigenvalue is odd under \(\eta\mapsto-\eta\);
for a degenerate cluster the statement concerns the matrix and spectral
multiset unless a compatible branch labeling is established.
\end{theorem}
\begin{proof}[Proof of the graded selection statement]
Conjugating the full Schur expression by the domain-preserving
involution sends \(D_\eta-z\) to \(-(D_{-\eta}+z)\).
Since the kernel is in positive grade, the stated oddness follows.
The reduced inverse changes sign under that involution, so
\(PMGMP=0\). Resolvent expansion about \(QD_0Q\) gives
\[
 \|[Q(D_\eta-z)Q]^{-1}-G\|
 \le\frac{v+|z|}{g(g-v-|z|)}.
\]
Multiplication by the two off-diagonal factors, each of norm at most
\(v\), gives the displayed error. The isolated simple branch is
unique, so its transformed branch has opposite eigenvalue.
\end{proof}

\begin{proposition}[An explicit inhomogeneous zero mode]
\label{p1b:result:an-explicit-inhomogeneous-zero-mode}
In the preceding family let \(s=q_3\) be a unit-sphere coordinate,
\(u=\epsilon s\), \(y=1-s^2\), \(a=\sqrt3\epsilon/2\), and
\[
 F(y)=\frac{\sinh(a\sqrt y)}{a\sqrt y}.
\]
The pair
\[
 \phi=F(y),\qquad
 \alpha=\frac{4R}{\epsilon}F'(y)\dd s
\]
extends smoothly at \(y=0\) and at \(\epsilon=0\), and is a
nonzero zero mode. At \(\epsilon=0\) it equals \((1,0)\).
\end{proposition}
\begin{proof}
On the round sphere \(\Delta s=3s/R^2\) and
\(\|\dd s\|^2=(1-s^2)/R^2\). The entire power series
\[
 F(y)=\sum_{m\ge0}\frac{(a^2y)^m}{(2m+1)!}
\]
gives \(4yF''+6F'=a^2F\) and is smooth at \(y=0\).
For the displayed one-form, \(\dd\alpha=0\), and
\[
 \dd^*\alpha=\frac{4s}{\epsilon R}(3F'+2yF'')
 =\frac{3\epsilon s}{2R}F,
 \qquad
 \dd\phi+\frac{u}{2R}\alpha=0.
\]
These are the two components of \(D_u(\phi,\alpha)=0\).
The series also proves smooth continuation through \(\epsilon=0\),
where the pair is \((1,0)\); its scalar component is nonzero.
\end{proof}


\subsubsection{Uniformly transverse roots of the protected native Cartan family}
\label{p1b:sec:native-source-root-transversality}

Use the round spatial restriction and Hodge identification of
Theorem~\ref{p1b:thm:dirac-spectrum}. For one fundamental twisting
doublet put
\[
 D_u=B_R+M_u,\qquad
 M_u=\frac{u}{2R}\operatorname{diag}(-3,1,1,1),\qquad
 B_R^2=\Delta_0\oplus\Delta_1.
\]
Suppose \(u\in C^\infty(S_R^3,\mathbb R)\), \(u\circ\iota=-u\), \(U=\|u\|_\infty\le1\), and normalize
the unique protected zero as \(s_u=(\phi_u,\alpha_u)\), with
\(\|\phi_u\|^2+\|\alpha_u\|^2=1\).
The constructive zero and its gap are already supplied by
Theorem~\ref{p1b:thm:dirac-zero-resolvent}; the following bound is
independent of an additional choice of its profile.

\begin{theorem}[Scalar weight and a nonvanishing native slope]
\label{p1b:thm:native-source-scalar-weight}
For every such background,
\begin{equation}
 F_u:=\|\phi_u\|^2\ge\frac{12-U^2}{12+8U^2}\ge\frac{11}{20}.
 \label{p1b:eq:native-source-scalar-weight}
\end{equation}
Let \(\lambda_u(\eta)\) be the simple eigenvalue through zero of
\(D_{u+\eta}=D_u+\eta M_1\), for a real constant offset \(\eta\).
It is odd in \(\eta\), and
\begin{equation}
 \kappa_u:=\lambda'_u(0)=\frac{1-4F_u}{2R},\qquad
 -\frac{3}{2R}\le\kappa_u
 \le-\frac{3(3-U^2)}{2R(3+2U^2)}\le-\frac{3}{5R}.
 \label{p1b:eq:native-source-slope}
\end{equation}
In particular this specified signed Dirac readout has a simple zero,
uniformly throughout the admitted protected-profile domain.
\end{theorem}
\begin{proof}
The lowest one-form Hodge eigenvalue is \(3/R^2\): the exact branch
starts at the first scalar shell and the coexact branch at \(4/R^2\).
Since \(B_Rs_u=-M_us_u\), the square of the Hodge operator gives
\[
 \frac{3}{R^2}\|\alpha_u\|^2
 \le\|B_Rs_u\|^2=\|M_us_u\|^2
 \le\frac{U^2}{4R^2}(9\|\phi_u\|^2+\|\alpha_u\|^2).
\]
Thus \((12-U^2)\|\alpha_u\|^2\le9U^2\|\phi_u\|^2\), which
gives the scalar-weight bound after normalization. The Born involution
\(\mathcal U(\phi,\alpha)=(\iota^*\phi,-\iota^*\alpha)\)
anticommutes with \(D_u\) and commutes with \(M_1\).
Theorem~\ref{p1b:thm:graded-feshbach-selection} therefore gives oddness
of the simple branch. Hellmann--Feynman gives the displayed slope;
substitution of the scalar-weight bound gives both uniform inequalities.
\end{proof}

\begin{theorem}[An explicit interval and inverse for the native readout]
\label{p1b:thm:native-source-root-interval}
Put
\[
 \gamma_u=\sqrt3-\tfrac32U>0,\qquad
 g_u=\gamma_u/R,\qquad \delta_u=\gamma_u/9.
\]
On \(|\eta|\le\delta_u\) the continuing low eigenvalue is simple and
\begin{equation}
 -\frac{3}{2R}\le\lambda'_u(\eta)\le-\frac{1}{5R},\qquad
 \frac{|\eta|}{5R}\le|\lambda_u(\eta)|
 \le\frac{3|\eta|}{2R}.
 \label{p1b:eq:native-source-monotone-interval}
\end{equation}
Consequently every supplied real target \(e\) with
\(|e|<\gamma_u/(45R)\) has a unique solution
\(\eta(e)\in(-\delta_u,\delta_u)\) of \(\lambda_u(\eta)=e\).
For \(e\ne0\), its sign is opposite to that of \(e\), and
\begin{equation}
 \frac{2R}{3}|e|\le|\eta(e)|\le5R|e|.
 \label{p1b:eq:native-source-inverse-bound}
\end{equation}
\end{theorem}
\begin{proof}
Write \(v=3|\eta|/(2R)\le g_u/6\). A contour at \(g_u/2\)
separates one low eigenvalue, with \(|\lambda_u(\eta)|\le v\).
Let \(P_u=|s_u\rangle\langle s_u|\), \(Q_u=1-P_u\), and let
\(s_{u,\eta}\) be its normalized continuing eigenvector. The
complementary equation gives
\[
 Q_us_{u,\eta}
 =-(D_u|_{Q_u}-\lambda_u(\eta))^{-1}Q_u\eta M_1s_{u,\eta},
 \qquad
 \|Q_us_{u,\eta}\|\le\frac{v}{g_u-v}\le\frac15.
\]
For normalized rank-one states the difference of their scalar-component
weights is at most this sine of their angle. Indeed their projector
difference has its only two possible nonzero eigenvalues
\(\pm\|Q_us_{u,\eta}\|\); pairing it with any orthogonal projector
has absolute value at most that number. Thus
\(\|\phi_{u,\eta}\|^2\ge11/20-1/5=7/20\).
Hellmann--Feynman proves the derivative bounds. Integrating them and
using strict monotonicity proves the inverse statement.
\end{proof}

These are consequences of the printed real Cartan operator family.
The offset may leave the subrange \(\|u+\eta\|_\infty\le1\);
that subrange is not an assumption of the perturbation theorem. If an
application additionally imposes it, intersect the interval above with
the offsets that satisfy that application's restriction. The target
\(e\) here is a signed spatial Dirac eigenvalue. 


\subsubsection{Midpoint gauge Morse structure and the protected Dirac coordinate}
\label{p1b:sec:native-weak-morse-bridge}
The results below use only the round spatial sphere, the midpoint connection
\(A_{\rm B}=\Theta/2\), the classical pure \(SU(2)\) Yang--Mills constituent,
and the fundamental twisting doublet already fixed in
Theorem~\ref{p1b:thm:dirac-spectrum}. They do not select a Lorentzian history
or a populated topological transition.

\begin{theorem}[Full spatial Morse index of the midpoint connection]
\label{p1b:thm:native-midpoint-gauge-morse}
Let \(A_{\rm B}=\Theta/2\) on the round \(S_R^3\), and consider the
background-gauge-fixed Hessian of the classical pure \(SU(2)\) Yang--Mills
energy on adjoint-valued one-forms.  Let \(L\) denote orbital spin
\(j\in\frac12\mathbb Z_{\ge0}\), let \(S\) be one-form spin one, let
\(T\) be adjoint isospin \(t=1\), and put \(J=L+S+T\).  On the
\((j,J)\) summand,
\begin{equation}
 \boxed{
 H_{\rm gf}=\frac1{R^2}
 \bigl[2J(J+1)+2j(j+1)-2\bigr].}
 \label{p1b:eq:native-midpoint-gauge-hessian}
\end{equation}
The scalar ghost operator on the \(k\subset j\otimes1\) summand has
\begin{equation}
 \boxed{
 \mu_{j}(k)=\frac1{R^2}
 \bigl[2j(j+1)+2k(k+1)-2\bigr].}
 \label{p1b:eq:native-midpoint-ghost}
\end{equation}
After quotienting the positive ghost images, the spatial Hessian on the gauge quotient has
exactly one negative mode,
\[
 (j,J)=(0,0),\qquad \lambda_-=-\frac2{R^2},
\]
no zero mode on the gauge quotient, and smallest positive eigenvalue \(1/R^2\).
Hence the midpoint has spatial Yang--Mills Morse index one.

On the left-invariant \(j=0\) sector,
\[
 \mathbf3_{\rm space}\otimes\mathbf3_{\rm ad}
 =\mathbf1\oplus\mathbf3\oplus\mathbf5.
\]
The three Hessian values are respectively \(-2,2,10\) in units \(R^{-2}\)
before gauge subtraction.  The \(\mathbf3\) is exactly the positive ghost
gauge image, so the homogeneous quotient sector is
\begin{equation}
 \boxed{
 \mathbf1:\ -\frac2{R^2},\qquad
 \mathbf5:\ \frac{10}{R^2}.}
 \label{p1b:eq:native-homogeneous-morse-split}
\end{equation}
\end{theorem}
\begin{proof}
Paper~I-D, Section~\ref{p1b:app:native-weak-morse-bridge} derives
\eqref{p1b:eq:native-midpoint-gauge-hessian} and
\eqref{p1b:eq:native-midpoint-ghost} from the midpoint covariant derivative
and curvature.  For \(j=0\), the decomposition
\(1\otimes1=0\oplus1\oplus2\) gives eigenvalues \(-2,2,10\); the
\(J=1\) summand is the scalar-gauge image.  For \(j=\frac12\), the two
lowest vector copies have \(J=\frac12\), eigenvalue \(1/R^2\); one is
removed by the positive ghost eigenvalue with \(k=\frac12\), leaving one
quotient copy.  All other \(j=\frac12\) values are at least \(7/R^2\).
For \(j\ge1\), even the formal value at \(J=0\) is
\([2j(j+1)-2]/R^2\ge2/R^2\).  Thus no further negative or zero quotient
mode exists, and the stated positive gap is sharp.
\end{proof}

\begin{theorem}[The unique Morse direction is the unique homogeneous source-active direction]
\label{p1b:thm:native-morse-dirac-alignment}
Let \(\xi_0\) be the normalized midpoint fundamental-doublet Dirac zero of
Theorem~\ref{p1b:thm:dirac-spectrum}.  Write a real left-invariant weak
connection perturbation in orthonormal components as
\[
 a_a=i\,b_a{}^c I_c,
 \qquad b=\beta I+A+S,
\]
with \(A^T=-A\), \(S^T=S\), and \(\operatorname{tr}S=0\).  Then:
\begin{enumerate}
\item the antisymmetric \(A\)-sector is exactly the homogeneous infinitesimal
gauge orbit and is removed by the Gauss/background-gauge projection;
\item the symmetric trace-free \(S\)-sector annihilates the protected singlet,
\(c(a_S)\xi_0=0\);
\item the scalar sector acts on the protected line by
\begin{equation}
 \boxed{c(a_{\beta I})\xi_0=\frac32\beta\,\xi_0.}
 \label{p1b:eq:native-source-active-scalar}
\end{equation}
\end{enumerate}
Consequently, among homogeneous adjoint-connection perturbations modulo gauge, the only
mode visible to the protected Dirac source is the \(\mathbf1\) scalar mode,
and it is exactly the unique negative Yang--Mills mode of
Theorem~\ref{p1b:thm:native-midpoint-gauge-morse}.
\end{theorem}
\begin{proof}
For constant adjoint gauge parameter \(\varphi\), the midpoint covariant
derivative \(d_{A_{\rm B}}\varphi=[A_{\rm B},\varphi]\) spans precisely the
antisymmetric matrices \(A\).  For the spin--isospin singlet,
\((\sigma_c/2+I_c)\xi_0=0\); hence
\[
 \sigma_a I_c\xi_0
 =-\frac12\delta_{ac}\xi_0-i\epsilon_{acd}I_d\xi_0
\]
up to the fixed Clifford sign convention already used in the midpoint
Dirac formula.  Symmetric trace-free contraction kills both displayed
terms.  Antisymmetric contraction lies in the triplet complement.  The
remaining scalar coefficient is fixed without sign ambiguity by the existing
projection identity
\(Pc(a)P=\frac12\operatorname{tr}b\,P\), giving
\eqref{p1b:eq:native-source-active-scalar}.  The homogeneous representation
labels then identify \(\beta I\) with the \(J=0\) Morse mode.
\end{proof}

\begin{corollary}[Exact Dirac, Chern--Simons and Yang--Mills barrier coordinate]
\label{p1b:cor:native-dirac-cs-barrier-coordinate}
On the homogeneous Cartan family \(A_f=f\Theta\), \(0\le f\le1\), let
\(\lambda\) be the protected signed Dirac branch.  Then
\begin{equation}
 \boxed{
 \lambda(f)=\frac{3/2-3f}{R},\qquad
 f=\frac12-\frac{R\lambda}{3}.}
 \label{p1b:eq:native-dirac-f-coordinate}
\end{equation}
With \(N_{\rm CS}(f)=3f^2-2f^3\) and
\(\mathcal E_{\rm YM}(f)=4f^2(1-f)^2\),
\begin{equation}
 \boxed{
 N_{\rm CS}(\lambda)
 =\frac12-\frac{R\lambda}{2}+\frac{2R^3\lambda^3}{27},
 \qquad
 \mathcal E_{\rm YM}(\lambda)
 =\frac14\left(1-\frac{4R^2\lambda^2}{9}\right)^2.}
 \label{p1b:eq:native-dirac-cs-potential}
\end{equation}
The flat endpoints and midpoint are therefore
\[
 (f,N_{\rm CS},\lambda)
 =(0,0,3/(2R)),\quad
 (1/2,1/2,0),\quad
 (1,1,-3/(2R)).
\]
Thus the protected signed eigenvalue is an exact coordinate on this specified
unit topological interpolation.  
\end{corollary}
\begin{proof}
The first equality is the simple branch in
Theorem~\ref{p1b:thm:dirac-spectrum}.  Solve it for \(f\) and substitute into
Proposition~\ref{p1b:prop:g11-transgression} and the potential term of
Theorem~\ref{p1b:thm:g11-functional}.  The endpoint values follow directly.
\end{proof}

\subsubsection{An evaluated nonconstant Cartan curve}
\begin{theorem}[Nonconstant Cartan energy and signed low-Dirac response]
\label{p1b:thm:v2217-nonconstant-cartan-dirac}
On $S_R^3$ retain the smooth real connection family
$A_u=(1+u)\Theta/2$ and
$D_u=B_R+(u/(2R))\operatorname{diag}(-3,1,1,1)$ on its common $H^1$
domain. Normalize the spatial Yang--Mills functional to give
$4f^2(1-f)^2$ for constant $f=(1+u)/2$, and let brackets denote unit-sphere
normalized integration. Then
\[
 \mathcal E(u)=\left\langle\frac{R^2|\nabla^{g_R}u|^2}{6}
                      +\frac{(1-u^2)^2}{4}\right\rangle,
 \qquad N_{\rm CS}(u)=\left\langle\frac12+\frac34u-\frac14u^3\right\rangle.
\]
For $h=q_1^2-1/4$, $u=\epsilon h$, one has
\begin{align*}
 \langle h\rangle&=0,&\langle h^2\rangle&=1/16,&
 \langle h^3\rangle&=1/64,&\langle h^4\rangle&=3/256,\\
 \mathcal E(\epsilon h)&=\frac14+\frac5{96}\epsilon^2+
                  \frac3{1024}\epsilon^4,&
 N_{\rm CS}(\epsilon h)&=\frac12-\frac{\epsilon^3}{256}.
\end{align*}
The simple eigenvalue of $D_{\epsilon h}$ continuing the midpoint zero is
analytic for sufficiently small $\epsilon$ and satisfies
\[
 \lambda_{\rm low}(\epsilon)=\frac9{8192R}\epsilon^3+O(\epsilon^5),
 \qquad R\lambda_{\rm low}=-\frac9{32}(N_{\rm CS}-1/2)+O(\epsilon^5).
\]
If $a_n=(h_n/2)\Theta$ for a degree-$n$ scalar harmonic, its orthogonal
background-gauge projection at the midpoint obeys, for $n\ge1$,
\[
 \frac{\|P_Ta_n\|^2}{\|a_n\|^2}
 =1-\frac{n(n+2)}{3[n(n+2)-2]}.
\]
It is zero at $n=1$, $5/9$ at $n=2$, and $8/13$ at $n=3$.
On the nonzero projected image the Yang--Mills Hessian has eigenvalue
$[n(n+2)-2]/R^2$ for $n\ge2$.
\end{theorem}
\begin{proof}
In a left-invariant unit-sphere frame the normalized curvature components are
$(\nabla_a u)\delta_b{}^c-(\nabla_b u)\delta_a{}^c+
(u^2-1)\epsilon_{ab}{}^c$. Their squared sum over $a<b,c$ is
$2|\nabla^{g_1}u|^2+3(1-u^2)^2$; the cross contraction vanishes.
Here $g_1=R^{-2}g_R$ is the pulled-back unit-sphere metric and
$|\nabla^{g_1}u|^2=R^2|\nabla^{g_R}u|^2$.
The constant-profile normalization gives $\mathcal E$. In the
Chern--Simons three-form the possible $du$ term vanishes because
$\operatorname{tr}(\Theta\wedge\Theta)=0$; substitution of $f=(1+u)/2$
in $3f^2-2f^3$ gives $N_{\rm CS}$.
The uniform sphere coordinate has
$\langle q_1^{2k}\rangle=(1/2)_k/(2)_k$. Expanding
$(q_1^2-1/4)^j$, $j=1,2,3,4$, gives the four moments. Also
$\Delta_0h=8h/R^2$ and integration by parts gives
$R^2\langle|\nabla^{g_R}h|^2\rangle=8\langle h^2\rangle=1/2$.
Thus the quadratic energy coefficient is
$\frac16\cdot\frac12-\frac12\cdot\frac1{16}=5/96$;
the energy and CS polynomials follow with no hidden radius factor.

For the eigenvalue, put $s_0=V_R^{-1/2}(1,0)$,
$P=|s_0\rangle\langle s_0|$, $Q=I-P$, $G=Q(B_R|_Q)^{-1}Q$ and
$M_h=(h/(2R))\operatorname{diag}(-3,1,1,1)$. The nonzero gap is
$\sqrt3/R$; bounded perturbation and the finite contour resolvent give
a simple analytic branch for small $\epsilon$. Since
$G(f,0)=(0,d\Delta_0^{-1}f)$ for mean-zero $f$, the linear term is zero
and $PM_hGM_hP=0$. Expanding the complementary inverse yields the cubic
coefficient
\[
 \langle s_0,M_hGM_hGM_hs_0\rangle
 =\frac9{8R^3}\langle h|d\Delta_0^{-1}h|^2\rangle
 =\frac9{8192R}.
\]
Indeed $\langle h|dh|^2\rangle=(4/R^2)\langle h^3\rangle$ by integration
by parts. Group inversion fixes $h$, so the involution of
Theorem~\ref{p1b:thm:dirac-reciprocity} sends $D_{\epsilon h}$ to
$-D_{-\epsilon h}$. Simplicity makes the eigenbranch odd, improving the
remainder to $O(\epsilon^5)$.

For the gauge projection put $j=n/2$. The divergence of $a_n$ is the
equivariant map $v\mapsto\sum_aJ_av\otimes e_a$ into the coupled
$k=j$ part of $V_j\otimes V_1$. The Casimir formula
\eqref{p1b:eq:native-midpoint-ghost} gives its adjoint-scalar eigenvalue
$[n(n+2)-2]/R^2$. Direct contraction gives
$\|d_{A_{\rm B}}^*a_n\|^2=n(n+2)\|a_n\|^2/(3R^2)$.
Subtracting the squared longitudinal norm using this inverse gives the
fraction. The tangent has total diagonal $J=j$; the equivariant projection
preserves it. Formula~\eqref{p1b:eq:native-midpoint-gauge-hessian} then
gives the claimed Hessian eigenvalue.
\end{proof}
The local gauge slice exists on the centralizer complement, where the
adjoint-scalar pivot is invertible. Small gauge transformations implementing
that slice unitarily conjugate $D_u$ and preserve the energy and the CS lift
near the midpoint. The formulas are therefore invariants of this specified
configuration curve. They do not assert nonlinear invariance of its
one-harmonic submanifold.
For the low Schur operator its coefficient of $-z$ is
$1+9\epsilon^2/512+O(\epsilon^4)$, since
$\|GM_hs_0\|^2=9/512$. The isolated even heat contribution begins at order
$\epsilon^6$; the complete heat can already change at order $\epsilon^2$
through its gapped modes. No signed cubic bosonic heat term follows.


\subsubsection{Complete-shell symmetric coexact pair targets}
\begin{theorem}[All-shell pair multiplicities and the mixed spinor target]
\label{p1b:thm:v2217-complete-shell-pair-targets}
For $K\ge1$ put $\mathcal G_{\le K}=\bigoplus_{k=1}^K(A_k\oplus B_k)$,
where $A_k=((k+1)/2,(k-1)/2)$ and $B_k=((k-1)/2,(k+1)/2)$ are the
coexact $\operatorname{Spin}(4)$ representations of
Theorem~\ref{p1b:thm:dirac-spectrum}. In $\operatorname{Sym}^2\mathcal G_{\le K}$
the multiplicities of
\[
 (0,0),\quad(1/2,0),\quad(0,1/2),\quad(1/2,1/2),\quad(1,0),\quad(0,1)
\]
are respectively $(2K,0,0,2(K-1),0,0)$. In particular
\[
 \operatorname{Sym}^2(A_1\oplus B_1)
 =2(0,0)\oplus(2,0)\oplus(0,2)\oplus(1,1).
\]
For $\mathcal S=(1/2,0)\oplus(0,1/2)$ the mixed target instead is
\[
 \mathcal S\otimes\mathcal S^*
 =2(0,0)\oplus(1,0)\oplus(0,1)\oplus2(1/2,1/2).
\]
\end{theorem}
\begin{proof}
Each coexact factor has equal left/right central parity. Tensor products
therefore exclude the two unequal-parity spinor types. For an output
$(1,0)$ in $A_k\otimes A_l$ or $B_k\otimes B_l$, right spin zero forces
$k=l$. The exchange sign of the $(1,0)$ component in that square is
$(-1)^{2k-1}=-1$, so it is absent from the symmetric square.
In $A_k\otimes B_l$, right spin zero requires $k=l+2$; minimum left
spin is then two. The reversed product gives the same exclusion.
Interchanging the two spin factors excludes $(0,1)$.
A scalar occurs once in each same-shell, same-chirality square, with
exchange sign $(-1)^{2k}=1$, and nowhere else. This gives $2K$.
To obtain $(1/2,1/2)$, same-chirality levels must be adjacent and each
adjacent pair contributes once. Opposite chiralities would require both
$|k-l+2|\le1$ and $|k-l-2|\le1$, which is impossible. Thus there are
$2(K-1)$ vector copies. The lowest-level table follows from
$\operatorname{Sym}^2(1)=0\oplus2$ in each factor. Finally the four
summands of $\mathcal S\otimes\mathcal S^*$ follow from
$\tfrac12\otimes\tfrac12=0\oplus1$ and self-duality of $SU(2)$ irreps.
\end{proof}
These are statements about specified representation spaces. A source map
into one of them must be separately specified before Schur's lemma gives
a physical null or nonzero response.

\subsection{Composition of complete regular reductions}
\label{p1b:sec:complete-successive-reductions}
\begin{corollary}[Complete successive Schur and stationary reductions]
\label{p1b:cor:complete-successive-reductions}
Let $H$ be a finite Hermitian Hessian on $x\oplus y\oplus z$.
Assume that the $z$ pivot and the resulting $y$ pivot are invertible on
one common parameter chart. Then the combined $y\oplus z$ pivot is
invertible and
\[
 \operatorname{Schur}_{y}(\operatorname{Schur}_{z}H)
       =\operatorname{Schur}_{y\oplus z}H.
\]
All existing parameter jets agree through the complete chain rule.
For a $C^2$ nonlinear action, the corresponding assertion holds
locally on the same uniquely generated stationary branch, provided the
successive implicit-function pivots are invertible. If an actual positive
bosonic Gaussian integration is licensed, its determinant and measure
factors compose as well:
\[
 \det H_{y\oplus z,y\oplus z}
  =\det H_{zz}\,
       \det\bigl(H_{yy}-H_{yz}H_{zz}^{-1}H_{zy}\bigr).
\]
In Euclidean conventions the effective energy retains
$\frac{\hbar}{2}\log\det C$ from each Gaussian pivot $C$, as well as
its inherited measure normalization and every consumed source derivative.
\end{corollary}
\begin{proof}
Eliminate $z$ from the block stationary equations, then $y$. Block
Gaussian elimination gives the same solution of the combined $y,z$
equations, proving both the remaining Hessian and determinant identities.
This also follows from the residual and jet calculus of
Theorem~\ref{p1b:thm:all-order-schur} and
Theorem~\ref{p1b:thm:schur-jets}. For a nonlinear action the implicit
function theorem gives a unique local solution of these same equations.
Substituting it either way gives the same reduced function. Differentiating
that equality proves equality of all extant jets, including derivatives of
the stationary branch. Successive positive Gaussian integrations use Fubini
and the displayed determinant factorization. Differentiating their complete
normalized result retains the measure and determinant jets.
\end{proof}
This is equality of complete reductions of one parent on one regular
branch. Classical elimination alone introduces no quantum determinant;
singular pivots require their own reduction certificate. Appendix~
\ref{p1b:app:compression-commutator} records a compression identity useful
for diagnosing why compression need not be an algebra restriction.

\subsubsection{Singular Hermitian elimination and its constraint relation}
\label{p1b:sec:v2214-singular-reduction}
\begin{theorem}[Exact elimination at every Hermitian pivot rank]
\label{p1b:thm:v2214-singular-schur}
Let $H=\left(\begin{smallmatrix}A&B\\B^\dagger&C\end{smallmatrix}\right)$
be Hermitian on finite Hilbert spaces $X\oplus Y$. Write $C^+$ for its
Moore--Penrose inverse, $P_0$ for the orthogonal projection onto $\ker C$,
and $S=A-BC^+B^\dagger$. For arbitrary right-hand side $(j_x,j_y)$, the
equation $H(x,y)=(j_x,j_y)$ is equivalent to
\[
 \begin{pmatrix}S&B|_{\ker C}\\P_0B^\dagger&0\end{pmatrix}
 \binom{x}{n}
 =\binom{j_x-BC^+j_y}{P_0j_y},\qquad n\in\ker C,
 \qquad y=C^+(j_y-B^\dagger x)+n.
\]
In particular, eliminating a singular block generally leaves a constraint
and a multiplier. An unconstrained Schur equation on $X$ is valid for all
compatible sources precisely when
\[
 B\ker C=0\quad\Longleftrightarrow\quad
 \operatorname{ran}B^\dagger\subseteq\operatorname{ran}C.
\]
In that case source compatibility is $P_0j_y=0$, $Sx=j_x-BC^+j_y$,
and $n$ is an arbitrary decoupled null coordinate.
\end{theorem}
\begin{proof}
Orthogonal projection of $B^\dagger x+Cy=j_y$ onto $\ker C$ gives the
second displayed row. Projection onto $\operatorname{ran}C$ and inversion
there gives the reconstruction of $y$. Substitution in the first row gives
the remaining equation. These reversible steps prove necessity and
sufficiency at every rank. Removing the constraint for every compatible
source requires $P_0B^\dagger=0$, equivalently $BP_0=0$; the equivalence
with the range condition follows from $C=C^\dagger$.
\end{proof}

\begin{corollary}[Positive pivots, quotient Gaussian and determinant]
\label{p1b:cor:v2214-positive-singular-pivot}
For the preceding block matrix,
\[
 H\succeq0\quad\Longleftrightarrow\quad
 C\succeq0,\quad B\ker C=0,\quad S\succeq0.
\]
Under these conditions the quadratic form is exactly
\[
 \langle(x,y),H(x,y)\rangle
 =\langle x,Sx\rangle+
 \langle y+C^+B^\dagger x,C(y+C^+B^\dagger x)\rangle.
\]
For real symmetric data, on the quotient $Y/\ker C$ with its induced
Euclidean measure, integration of the centered real Gaussian contributes
$(2\pi)^{r/2}(\det{}' C)^{-1/2}$, where $r=\operatorname{rank}C$.
For complex Hermitian data the corresponding quotient Gaussian integral
is $\pi^r/(\det{}'C)$ for the convention $e^{-y^\dagger Cy}$.
If $S$ is positive definite, then
\[
 \ker H=\{0\}\oplus\ker C,
 \qquad \det{}'H=\det S\,\det{}'C.
\]
For singular $S$ the last product is not asserted: the induced metric on
the full null quotient also has to be transported.
\end{corollary}
\begin{proof}
For $n\in\ker C$, positivity on $(tx,n)$ for every real $t$ forces
$\operatorname{Re}\langle x,Bn\rangle=0$; replacing $x$ by $ix$ gives
$Bn=0$. Completing the square then proves the equivalence. On
$\operatorname{ran}C$ the positive matrix is invertible, and the ordinary
Gaussian integral gives the stated factor. If $S>0$, the square completion
identifies the kernel with $\ker C$ and restricts $H$ to the orthogonal
complement $X\oplus\operatorname{ran}C$. Block triangular elimination on
that complement has determinant one and gives the determinant identity.
\end{proof}

\begin{proposition}[Constant-rank differentiation and exact composition]
\label{p1b:prop:v2214-singular-jets}
On a differentiable constant-rank Hermitian family,
\[
 \mathrm dC^+=-C^+(\mathrm dC)C^+
 +(C^+)^2(\mathrm dC)P_0+P_0(\mathrm dC)(C^+)^2,
 \qquad
 \mathrm dP_0=-C^+(\mathrm dC)P_0-P_0(\mathrm dC)C^+.
\]
These formulas, the product rule and the exact constraint relation above
determine all extant fixed-rank parameter jets. Successive eliminations of
one parent linear system give the same relation between the retained
variables and sources whenever every intermediate constraint and null
variable is retained. On a common positive quotient branch their Gaussian
integrals also agree with the transported quotient measures.
\end{proposition}
\begin{proof}
Differentiate $CC^+=C^+C=I-P_0$, $CP_0=0$ and $P_0^2=P_0$, using
$P_0(\mathrm dC)P_0=0$, which holds on a constant-rank stratum.
Solving the resulting range--kernel blocks gives the two formulas.
Repeated differentiation gives higher jets. Every elimination step is the
reversible equivalence in Theorem~\ref{p1b:thm:v2214-singular-schur};
composition therefore preserves exactly the solution relation. Positive
quotient Gaussian composition follows from Fubini on the nondegenerate
directions, with the inherited measures. No assertion identifies derivatives
across a rank-changing point with these stratumwise jets.
\end{proof}

\subsubsection{Exact common positive-parent reconstruction}
\begin{theorem}[Inverse, direct and square compressions of one positive parent]
\label{p1b:thm:v2217-common-positive-parent}
Let $\mathcal I,D,\mathcal G$ be positive definite operators on one finite
Hilbert space $X$, with the same inner product. Put
\[
 S=\mathcal G-D^2,\qquad M=D-\mathcal I^{-1}.
\]
There are a finite Hilbert space $Y$, a positive definite $\mathcal R$ on
$X\oplus Y$, and the coordinate inclusion $V:X\to X\oplus Y$, such that
\[
 V^*\mathcal R^{-1}V=\mathcal I,\quad
 V^*\mathcal RV=D,\quad V^*\mathcal R^2V=\mathcal G,
\]
if and only if $S,M\succeq0$ and $\ker S=\ker M$.
The smallest $\dim Y$ is $r=\operatorname{rank}S=\operatorname{rank}M$.
For $r>0$ let $U:\mathbb C^r\to X$ be an isometry onto their common
support, $S_r=U^*SU$, $M_r=U^*MU$. A minimal realization is
\[
 B=US_r^{1/2},\quad Z=S_r^{1/2}M_r^{-1}S_r^{1/2},\qquad
 \mathcal R=\begin{pmatrix}D&B\\B^*&Z\end{pmatrix}.
\]
It is unique up to a unitary on $Y$ fixing $X$ pointwise, and
\[
 \det\mathcal R=\frac{\det S_r}{\det M_r\det\mathcal I}.
\]
For $r=0$ take $Y=0$, $\mathcal R=D$ and empty support determinants equal to one.
\end{theorem}
\begin{proof}
Write any proposed parent in the displayed block form. Positivity gives
$Z>0$. Squaring and block inversion give
$S=BB^*$ and $M=BZ^{-1}B^*$. Both are positive with kernel $\ker B^*$,
and $\dim Y\ge\operatorname{rank}B=r$. Conversely the stated construction
has $BB^*=S$, $BZ^{-1}B^*=M$ and Schur complement
$D-BZ^{-1}B^*=\mathcal I^{-1}>0$. Thus $\mathcal R>0$ and all three
compressions follow by multiplication and inversion. At minimal dimension
$B$ has full column rank. Any two factors of $BB^*=S$ differ by a right
unitary, after which $M=BZ^{-1}B^*$ uniquely fixes $Z$. Finally
$\det\mathcal R=\det Z\det\mathcal I^{-1}$ gives the formula.
\end{proof}
This is a representability theorem for three specified operators in one
specified metric. It does not determine that tuple, its physical occurrence,
or any identification of three scalar measurements with full matrices.


\section{Special substructures, Clifford triality, and rooted hierarchy}
\label{p1b:sec:organized-rooted}

\subsection{Typed special-substructure atlas of \texorpdfstring{\(M_8\)}{M8}}
\label{p1b:app:m8-substructures}
The subgroup, primitive-circle quotient, exchange-selection,
derived-complex and aligned exceptional-fingerprint results are established
on their stated mathematical domains. This subsection records auxiliary
consequences that do not enter every subsequent reciprocal construction.  Throughout, construction type is part of the datum:
a subgroup, a quotient, an embedded submanifold and an external comparison
space are not identified merely because they have the same dimension or
abstract diffeomorphism type.  The wider incidence and literature comparisons do not add
new occurrence premises to the stated geometric classifications.

\subsubsection{Higher-rank torus quotients}
\begin{proposition}[Smith-normal-form control of connected torus quotients]
\label{p1b:prop:m8-smith-torus-quotients}
Fix the marked maximal torus
\[
 T^4_{\max}=\U(1)_{H,x}\times\U(1)_{H,p}\times\T^2_{\rm sc}.
\]
Let a connected embedded \(r\)-torus \(T^r_\Lambda\subset T^4_{\max}\)
correspond to a saturated rank-\(r\) lattice
\(\Lambda\subset\Z^4\), and let
\[
 P_{\rm sc}:\Lambda\longrightarrow\Z^2
\]
be projection to the two scalar cocharacter coordinates.  For
\(Q_\Lambda=M_8/T^r_\Lambda\),
\[
 \pi_2(Q_\Lambda)\cong\ker P_{\rm sc},\qquad
 \pi_1(Q_\Lambda)\cong\operatorname{coker}P_{\rm sc},\qquad
 \pi_k(Q_\Lambda)\cong\pi_k(M_8)\quad(k\ge3).
\]
Thus the Smith normal form of an integer matrix for \(P_{\rm sc}\)
completely determines the low homotopy groups.  In rank two, scalar
projection of rank zero gives the full bi-Hopf quotient
\(S^2\times S^2\times\T^2\); scalar projection of rank two and determinant
\(\pm1\) makes the torus a graph over \(\T^2_{\rm sc}\), and the quotient is
diffeomorphic to \(S^3\times S^3\).  If the scalar projection has rank one,
then \(\pi_2\cong\Z\) and
\(\pi_1\cong\Z\oplus\Z/d\), where \(d\) is the nonzero Smith invariant.
\end{proposition}
\begin{proof}
The torus acts freely by right translation.  The fibration
\(T^r_\Lambda\to M_8\to Q_\Lambda\), together with
\(\pi_2(M_8)=0\) and \(\pi_1(M_8)=\Z^2\), gives the exact segment
\[
 0\to\pi_2(Q_\Lambda)\to\Lambda
   \xrightarrow{P_{\rm sc}}\Z^2
   \to\pi_1(Q_\Lambda)\to0.
\]
Higher homotopy is unchanged because \(\pi_k(T^r)=0\) for \(k\ge2\).
The three rank-two cases follow from the Smith normal form.  When the
scalar projection is unimodular, projection of the acting torus to the
scalar torus is an isomorphism; gauge fixing both scalar coordinates leaves
\(S^3\times S^3\).  The rank-zero case is
Theorem~\ref{p1b:thm:bihopf} with the spectator scalar torus retained.
\end{proof}
The rank-one case has the displayed homotopy classification, but a full
diffeomorphism classification also depends on the Hopf-slope data and is
not asserted here.

\subsubsection{Hopf--Heegaard--Clifford strata}
Fix
\[
 S^3=\{(z_1,z_2)\in\C^2:|z_1|^2+|z_2|^2=1\}
\]
and the standard Hopf map
\[
 h(z_1,z_2)=
 \bigl(2\Re(z_1\bar z_2),2\Im(z_1\bar z_2),|z_1|^2-|z_2|^2\bigr).
\]
For \(0<\rho<\pi/2\), put
\[
 T^2_\rho=\{|z_1|=\cos\rho,\ |z_2|=\sin\rho\}.
\]

\begin{proposition}[Hopf--Heegaard pullback identity]
\label{p1b:prop:hopf-heegaard-pullback}
Let \(C_\rho\subset S^2\) be the latitude with third coordinate
\(\cos(2\rho)\).  Then
\[
 T^2_\rho=h^{-1}(C_\rho).
\]
In particular the balanced Clifford torus is the Hopf preimage of the
equator, and the two solid tori in the standard genus-one Heegaard splitting
are the inverse images of the two hemispheres.  Their core circles are the
Hopf fibres over the poles.
For the product bi-Hopf map,
\[
 T^4_{\rho,\sigma}
 =(h_x\times h_p)^{-1}(C_\rho\times C_\sigma),
\]
and multiplying by \(\T^2_{\rm sc}\) gives an embedded
six-torus in \(M_8\).
\end{proposition}
\begin{proof}
On \(T^2_\rho\),
\[
 |z_1|^2-|z_2|^2=\cos^2\rho-\sin^2\rho=\cos(2\rho),
\]
so its Hopf image lies in \(C_\rho\).  Conversely the Hopf map fixes this
third coordinate exactly when the two moduli have the displayed values,
giving equality.  Hemisphere and product statements follow by taking
inverse images and products.
\end{proof}

The resulting seven-dimensional hypersurfaces
\[
 \mathcal N_{x,\rho}=T^2_{\rho,x}\times S^3_p\times\T^2_{\rm sc},\qquad
 \mathcal N_{p,\sigma}=S^3_x\times T^2_{\sigma,p}\times\T^2_{\rm sc}
\]
are diffeomorphic to \(T^4\times S^3\).  They are embedded separating
submanifolds, not the connected hypersubgroups
\(S^3\times S^3\times S^1\) of
Corollary~\ref{p1b:cor:codim-one-normal}.  Likewise the bi-Clifford
\(T^6\subset M_8\) is not a Lie subgroup: the maximal Lie torus of
\(M_8\) has rank four.

\subsubsection{The \(\Spin^c(4)\)--\(\Spin(6)\) comparison branch}
\begin{proposition}[Exact central quotient and block embedding]
\label{p1b:prop:spinc4-spin6}
For the abstract seven-dimensional group
\(H_7=\SU(2)\times\SU(2)\times\U(1)\), define
\[
 \Phi(A,B,z)=\operatorname{diag}(zA,z^{-1}B)\in\SU(4).
\]
Then
\[
 \ker\Phi=\langle(-I,-I,-1)\rangle,
 \qquad
 \operatorname{im}\Phi=S(\U(2)\times\U(2)),
\]
and therefore
\[
 \frac{\Spin(4)\times\U(1)}{\Z_2}
 \cong\Spin^c(4)
 \cong S(\U(2)\times\U(2))
 \subset\SU(4)\cong\Spin(6).
\]
For the relative scalar slope in \(M_8\), Born exchange is intertwined by
block exchange with
\((A,B,z)\mapsto(B,A,z^{-1})\).  Under the same embedding,
\[
 \mathbf4\downarrow=(\mathbf2,\mathbf1)_{+1}
   \oplus(\mathbf1,\mathbf2)_{-1},
\]
\[
 \mathbf6\downarrow=(\mathbf2,\mathbf2)_0
   \oplus(\mathbf1,\mathbf1)_{+2}
   \oplus(\mathbf1,\mathbf1)_{-2}.
\]
\end{proposition}
\begin{proof}
The determinant of \(\Phi(A,B,z)\) is
\(z^2z^{-2}\det A\det B=1\).  Its image consists exactly of block
unitaries with determinant one.  The block matrix is the identity precisely
when \(zA=I\) and \(z^{-1}B=I\); since \(A,B\in\SU(2)\), this gives
\((I,I,1)\) and \((-I,-I,-1)\).  Block exchange gives the stated
intertwining relation.  The branching of \(\mathbf4\) is immediate from
the two blocks, and taking \(\Lambda^2\mathbf4\) gives the branching of
\(\mathbf6\).
\end{proof}
\textbf{Status: COMPARISON-ONLY.}  This is an external comparison branch.  It does not identify \(M_8\) with
\(\Spin(6)\), nor does it identify the finite central quotient above with
the scalar-circle quotient \(M_8/\U(1)_{\rm diag}\).

\subsubsection{Complex and metric structures carried by native nodes}
Corollary~\ref{p1b:cor:g16-derived-complex} identifies the rank-one
Dolbeault stratum as exactly the locus where the native
\(D_6=S^3\times S^3\) is a complex submanifold.  After the active Hopf axes
and root-sign conventions are fixed, the induced invariant complex
structures on \(D_6\) are the classical Calabi--Eckmann--Morimoto family;
this is an identification of the invariant complex structure \cite{CalabiEckmann1953,Magnin2008ComplexSU2SU2,Marchidanu2024SasakianProducts}.
The exchange-invariant complex branch remains disjoint from this locus by
Corollary~\ref{p1b:cor:g16-derived-complex}.

\textbf{Status: LITERATURE-EXTRA for the structures named in this paragraph.}\par
Other familiar structures are similarly typed as additional geometric data:
the carrier \(S^3\times S^3\) supports a homogeneous strict nearly-K\"ahler
structure \cite{Butruille2005NearlyKahler}; \(S^3\times S^1\) supports the
standard Hopf/Vaisman constructions after compatible complex and Hermitian
data are chosen \cite{OrneaVerbitsky2003Vaisman}; and the exchange-selected
primitive pure-Hopf slope has carrier
\(T^{1,1}\cong S^2\times S^3\), whose standard homogeneous
Sasaki--Einstein/conifold structure is additional metric/contact data
\cite{KlebanovWitten1998Conifold}.  None of these literature structures is
inferred from carrier topology alone.

\subsubsection{Finite exact-arithmetic certificate and typing firewalls}
The exact finite cardinalities in
Proposition~\ref{p1b:prop:exceptional-hopf-clifford} are independently
reproduced by
\path{Paper1B-M8SubstructureChecks-v1_3.py}; the frozen output is
\path{Paper1B-M8SubstructureChecks-v1_3.json}.  The same checker also
exhausts 5,856 primitive integer weights in a bounded lattice window to test
the exchange eigensublattice criterion and checks the rank-one Smith
invariant for 624 nonzero scalar-weight pairs.  These computations are
independent checks of the finite and integer-lattice statements; the proofs
in the main text and above remain the mathematical justification.

The following distinctions are mandatory throughout the atlas:
\begin{enumerate}[label=(\roman*)]
\item \(S^3\times S^3\times S^1\) is a seven-dimensional carrier, not
      \(\Spin(6)\), which has dimension fifteen.
\item A Heegaard or Clifford \(T^2\subset S^3\) is not a two-dimensional
      Lie subgroup of \(\SU(2)\); the latter has rank one.
\item The embedded bi-Clifford \(T^6\subset M_8\) is not the maximal Lie
      torus, which has dimension four.
\item \(M_8/\Delta K\cong K\) is a homogeneous orbit space, not a quotient
      Lie group, because \(\Delta K\) is not normal.
\item The scalar quotient \(X_7\), an embedded hypersubgroup
      \(H^7_{(m,n)}\), and a primitive-scalar maximal-torus quotient may be
      diffeomorphic while remaining different constructions.
\item The sixteen half-Hurwitz Clifford points are not a subgroup of
      \(2T\).
\item Abstract isomorphisms such as
      \(SL(2,3)\cong2T\) and \(SL(2,5)\cong2I\) do not identify scalar
      torsion points with marked sphere vertices .
\end{enumerate}

\subsubsection{Metric factorization of the primitive bi-Hopf quotients}
\label{p1b:app:g17}

Equip each three-sphere with a round metric for which its marked Hopf map is
its standard Riemannian submersion, and give
\(P_v=S_x^3\times S_p^3\) the product metric.  Quotient spaces below carry
the induced Riemannian-submersion metrics.

\begin{proposition}[Product Hopf harmonic morphism and primitive factorization]
\label{p1b:prop:g17-harmonic}
The product Hopf map
\[
 \Pi_H=\pi_x\times\pi_p:
 S_x^3\times S_p^3\longrightarrow S_x^2\times S_p^2
\]
is a harmonic morphism.  For every primitive pure-Hopf slope \((a,b)\),
its factorization
\[
 P_v\longrightarrow T^{a,b}_{\rm top}
 \longrightarrow S_x^2\times S_p^2
\]
consists of Riemannian submersions with one-dimensional minimal fibres in
the first map and one-dimensional residual torus fibres in the second.
Consequently both arrows are harmonic morphisms.
\end{proposition}
\begin{proof}
Each Hopf map is a Riemannian submersion with geodesic circle fibres.  Their
product is therefore a Riemannian submersion with totally geodesic flat
\(\T^2\)-fibres, hence with minimal fibres.  A Riemannian submersion is a
harmonic morphism precisely when its fibres are minimal.

For the first quotient, the infinitesimal generator
\(aH_x+bH_p\) is a Killing field of constant length: the two Hopf fields
have constant lengths and are tangent to orthogonal product factors.
Integral curves of a constant-length Killing field are geodesics, so the
subcircle fibres are minimal.  The induced metric on
\(T^{a,b}_{\rm top}\) is by definition the quotient metric.  The original
flat torus fibre quotiented by this geodesic subcircle is the residual
circle fibre of the second arrow; its metric and horizontal distribution
are induced from the product principal-torus metric, so that fibre is again
geodesic.  Hence both quotient maps are harmonic Riemannian submersions.
\end{proof}

This metric statement does not alter the quotient classification of
Theorem~\ref{p1b:thm:primitive-circle-quotients}; it only equips the already
classified pure-Hopf branch with its natural quotient metrics.


\subsection{Exact support counterexample and conditional alternative constructors}
\label{p1b:app:v233-support-countermodels}

\begin{proposition}[An anchored 5-design without a Hurwitz subset]
\label{p1b:prop:v233-hopf72}
There is a uniform antipodal spherical 5-design of 72 points which
contains \(1,-1\), has no orthogonal triple and contains no orthogonally
congruent copy of the 24-cell.
\end{proposition}
\begin{proof}
Put \(u_\pm=(1\pm1/\sqrt3)/2\), and for
\(u\in\{u_+,u_-\}\), \(0\le j,k<6\), take
\[
 z(u,j,k)=
 \bigl(\sqrt u\cos(\pi j/3),\sqrt u\sin(\pi j/3),
 \sqrt{1-u}\cos(\pi k/3),\sqrt{1-u}\sin(\pi k/3)\bigr).
\]
These are 72 distinct unit vectors. In these Hopf coordinates the
probability Haar measure is \(du\,d\theta\,d\phi/(2\pi)^2\).
Each six-angle grid integrates Fourier modes with absolute frequency
at most five exactly. After angular averaging, a coordinate monomial
of total degree at most five either vanishes or has even exponents
and radial part proportional to \(u^a(1-u)^b\), \(a+b\le2\).
The equal-weight two-node rule \(u_\pm\) integrates radial degrees
through three: its moments are \(1,1/2,1/3,1/4\).
Thus the 72-node rule is a 5-design.

The shift \((j,k)\mapsto(j+3,k+3)\) gives antipodality.
Two nodes on the same radial layer have inner product
\(u a+(1-u)b\), where \(a,b\in\{1,1/2,-1/2,-1\}\).
It cannot vanish: \(u\) is irrational, so vanishing would force
\(a=b=0\), absent from the list. Any triple contains two nodes in
the same layer; hence no orthogonal triple exists. A 24-cell contains
four mutually orthogonal vertices. Finally, left multiply the entire
set by the inverse of one node. This is an orthogonal transformation,
preserving all properties and sending an antipodal pair to \(1,-1\).
\end{proof}

This counterexample refutes the inference from positive finite exactness,
antipodality and an identity anchor to group closure or a Hurwitz subset.


\begin{proposition}[Frame-normalizer and three-torsion constructions]
\label{p1b:prop:v233-alternative-2t-constructors}
For an oriented orthonormal imaginary-quaternion frame \(F=(i,j,k)\),
put \(Q(F)=\{\pm1,\pm i,\pm j,\pm k\}\). Then
\[
 N_{SU(2)}Q(F)=2O(F),\qquad
 [N_{SU(2)}Q(F),N_{SU(2)}Q(F)]=2T(F).
\]
The full frame family gives the unmarked conjugacy orbit, with
stabilizer the octahedral rotation group. Separately, the symplectic
automorphism group of the scalar three-torsion module is
\[
 \operatorname{Sp}(\mathbb F_3^2)=SL(2,\mathbb F_3)
       \cong Q_8\rtimes C_3\cong2T.
\]
Both constructions retain their specified frame or torsion-module input.
\end{proposition}
\begin{proof}
Under the double cover to \(SO(3)\), \(Q(F)\) projects to the
three orthogonal half-turns and the identity. Its normalizer permutes
their axes, and every orientation-preserving signed permutation does
so. This is the order-24 octahedral group, with binary lift of order
48. Its derived projective group is \([S_4,S_4]=A_4\);
\([i,j]=-1\) supplies the covering kernel, giving the binary group
of order 24. Frame rotation conjugates every construction; its
stabilizer is the octahedral normalizer.

In dimension two an alternating form transforms by the determinant.
At an odd prime \(p\), a determinant-one frame has \(p^2-1\) choices
of nonzero first vector and \(p\) choices of partner, hence the
group has order \(p(p^2-1)\). At \(p=3\), set, over \(\mathbb F_3\),
\[
 A_0=\begin{pmatrix}1&1\\1&2\end{pmatrix},\quad
 A_1=\begin{pmatrix}2&1\\1&1\end{pmatrix},\quad
 A_2=A_0A_1,\quad
 U=\begin{pmatrix}1&1\\0&1\end{pmatrix}.
\]
Multiplication gives \(A_i^2=-I\), \(A_0A_1=-A_1A_0\),
\(U^3=I\), and conjugation by \(U\) cycles \(A_0,A_1,A_2\).
Thus the eight quaternion elements form a normal \(Q_8\), with
the order-three subgroup as complement. The corresponding Hurwitz
half-unit cycles the quaternion axes, proving the isomorphism.
\end{proof}

If a \(2O\) support already occurs and the subgroup-word output is
licensed, its commutator subgroup needs no duplicate occurrence.
The identity does not provide the initial frame support or a normalizer
constructor. Similarly, the torsion calculation does not select \(p=3\)
or identify a scalar-torsion automorphism with an embedded momentum
point. These are the precise domains of the two alternative routes.


\subsection{Primitive-owner selection and odd Schur target shape}
\label{p1b:sec:v125-primitive-owner-odd-schur}

\begin{definition}[Primitive finite owner block]
\label{p1b:def:v125-primitive-finite-owner}
A finite block is \emph{primitive relative to the constituted owner} when all
mandatory projector, involution, quotient and unital refinements already
generated by that owner have been performed, the remaining block carries a
non-scalar owner type constituted upstream, and it is not merely a parent or
composite block whose generated refinement supplies the actual owner pieces.
This is a typed owner notion, not a parity or prime-dimension condition.
\end{definition}

\begin{theorem}[Unique non-scalar \(q=5\) node owner]
\label{p1b:thm:v125-q5-rank2-owner}
Let \(R\) be the regular six-cycle shift, \(J=R^3\), and
\[
 P_1=\frac16\sum_{k=0}^5R^k,\qquad
 P_3=\frac{I-J}{2},\qquad
 P_2=\frac{I+J}{2}-P_1.
\]
Then the ranks are \((1,3,2)\), with
\[
 JP_1=P_1,\qquad JP_2=P_2,\qquad JP_3=-P_3.
\]
Let the target owner be \(E_3(u)=\mathbb R1\oplus u^\perp\), with the
nontrivial action factoring through \(C_6/\{\pm1\}\cong C_3\).  Every
\(C_6\)-equivariant owner map \(A:\mathbb R[C_6]\to E_3(u)\) satisfies
\[
 AP_3=0.
\]
Its \(P_1\) image is scalar, while the unique non-scalar real isotypic owner
is
\[
 \boxed{P_2\simeq u^\perp}.
\]
Consequently the unique primitive non-scalar \(q=5\) node owner has rank two.
\end{theorem}
\begin{proof}
The target action has \(J=I\), hence \(AJ=A\).  Since \(JP_3=-P_3\),
\(AP_3=AJP_3=-AP_3\), so \(AP_3=0\).  The constant projector \(P_1\)
is the trivial character.  The remaining real two-dimensional summand is the
conjugate pair of nontrivial \(C_3\) characters, which is exactly the real
rotation plane \(u^\perp\).
\end{proof}

\begin{corollary}[Paired complement gives a purely odd zero-order Schur descendant]
\label{p1b:cor:v125-paired-complement-odd-schur}
Let
\[
 F_P=F_L\oplus F_R,\qquad F_Q=Q_L\oplus Q_R,
\]
with invertible complement block
\[
 D_Q=\Delta\begin{pmatrix}0&U\\ U^\dagger&0\end{pmatrix},
 \qquad UU^\dagger=U^\dagger U=I,
\]
and chirality-even retained-to-complement coupling
\[
 B=\begin{pmatrix}A_L&0\\0&A_R\end{pmatrix}.
\]
Then the Schur correction \(\Xi=-BD_Q^{-1}B^\dagger\) is
\[
 \boxed{
 \Xi=-\frac1\Delta
 \begin{pmatrix}
 0&A_LUA_R^\dagger\\
 A_RU^\dagger A_L^\dagger&0
 \end{pmatrix}}
\]
and therefore
\[
 P_L\Xi P_L=P_R\Xi P_R=0,
 \qquad
 P_L\Xi P_R=-\Delta^{-1}A_LUA_R^\dagger.
\]
Thus a paired eliminated complement has exactly the zero-order
\(L\leftrightarrow R\) target shape required of a finite odd chiral block.
\end{corollary}
\begin{proof}
The stated involutory off-diagonal form gives
\(D_Q^{-1}=\Delta^{-1}\bigl(\begin{smallmatrix}0&U\\U^\dagger&0\end{smallmatrix}\bigr)\).
Direct block multiplication gives the formula.
\end{proof}

\begin{remark}[Intertwiners between distinct rank-sixteen constructions]
\label{p1b:rem:v125-occurrence-rank16-firewalls}
The preceding statement describes an odd block constructed from $A_L,A_R,U$. The Clifford quotient gives $16=1+9+6=9+7$, but equality of this dimension with another $7+9$ decomposition does not identify the representations. Such an identification requires the acting group and an explicit intertwiner.
\end{remark}


\subsection{Spin(8) triality restriction to the Clifford channel}
\label{p1b:sec:v126-d4-triality-clifford-restriction}

Paper~I-B supplies the finite ordered triality-incidence parent.  The
following branching statement identifies its protected four-dimensional
restriction with the Clifford channel already present in the framework.

\begin{theorem}[Spin(8) triality restriction]
\label{p1b:thm:v126-spin8-triality-restriction}
Use
\[
\Spin(8)\supset
 \bigl(\Spin(4)_A\times\Spin(4)_B\bigr)/\langle(-1,-1)\rangle.
\]
The following branching labels are pulled back to its covering product
\[
\Spin(4)_A\times\Spin(4)_B
\cong SU(2)_{A+}\times SU(2)_{A-}\times SU(2)_{B+}\times SU(2)_{B-}.
\]
With a standard triality convention the three eight-dimensional modules branch
as
\begin{align*}
 8_v&=(2,2,1,1)\oplus(1,1,2,2),\\
 8_s&=(2,1,2,1)\oplus(1,2,1,2),\\
 8_c&=(2,1,1,2)\oplus(1,2,2,1).
\end{align*}
The unique Spin(8)-invariant triality trilinear contains the component
\[
 (2,2,1,1)\otimes(2,1,2,1)\otimes(1,2,2,1)\longrightarrow\mathbb C.
\]
Contracting the spectator \(SU(2)_{B+}\)-doublets by their invariant
alternating form yields a nonzero \(\Spin(4)_A\)-equivariant map
\[
 \boxed{V_A\otimes S_{A+}\longrightarrow S_{A-}.}
\]
The opposite chirality is obtained by the companion summand.  Each retained
channel has multiplicity one.
\end{theorem}

\begin{proof}
The displayed branchings are the standard restriction of the vector and two
half-spin representations.  In the indicated tensor product, the
\(SU(2)_{B+}\) factors form \(2\otimes2\), whose invariant alternating line is
one-dimensional; the remaining \(A\)-factors are
\((2,2)\otimes(2,1)\to(1,2)\), again with multiplicity one.  Nonvanishing
follows from nonvanishing of the parent triality form on its weight-zero
support.  The conjugate component gives the reverse chirality.
\end{proof}

\begin{corollary}[The normalized triality restriction is Clifford multiplication]
\label{p1b:cor:v126-triality-restriction-is-clifford}
Use the signed Fock tensor and the real orthonormal vector normalization of
Paper~I-B, Theorem~\ref{p1:thm:v214-normalized-spin8-triality}. On the
spectator doublet use its alternating form $\epsilon_B$, with
$\|\epsilon_B\|_{\rm HS}^2=2$, and define the coefficient extraction
\[
 \operatorname{Ext}_B(F)=\frac12
       \sum_{\alpha,\beta=1}^2
          \overline{(\epsilon_B)_{\alpha\beta}}F_{\alpha\beta}.
\]
Then the retained four-dimensional triality coefficient is exactly
\[
 \boxed{\operatorname{Ext}_B
   \bigl(\Theta_{\rm Spin(8)}|_{V_A,S_{A+},S_{A-}}\bigr)(v)=c_A(v).}
\]
All signs are those of the stated Fock bases. This equation uses the signed
coefficient tensor, not the unsigned $192$-incidence counting vector.
\end{corollary}
\begin{proof}
Write $U=U_A\oplus U_B$ and identify
$\Lambda^*U=\Lambda^*U_A\widehat\otimes\Lambda^*U_B$. For $v\in V_A$,
the explicit wedge/contraction matrices give
$c_8(v)=c_A(v)\otimes I$. Chevalley pairing restricts to the product
of the $A$ pairing and $\epsilon_B$, with the signs already evaluated in
the cited one-factor theorem. Thus the retained trilinear is
$\Theta_A\otimes\epsilon_B$, coefficient for coefficient.
The extraction gives
$\frac12\sum_{\alpha,\beta}|(\epsilon_B)_{\alpha\beta}|^2\Theta_A
=\Theta_A$, proving the equation and its magnitude one. Evaluating a
specified spectator pair with $\epsilon_B(\xi,\eta)=1$ gives the same
coefficient. Full contraction with $\epsilon_B^*$ instead gives
$2\Theta_A$, and contraction with the unit tensor
$\epsilon_B^*/\sqrt2$ gives $\sqrt2\Theta_A$; those are distinct
specified contractions and are not silently identified here.
\end{proof}

\begin{corollary}[Primitive owner descendants of one rooted triality parent]
\label{p1b:cor:v126-triality-primitive-owner-descendants}
After rooting an oriented Hurwitz frame, the presently used primitive
non-scalar owner ranks
\[
 \boxed{2,\ 4,\ 6}
\]
are descendants of the same triality parent: rank four is the vector slot
\(V\), rank six is the Clifford/exterior-square descendant \(\Lambda^2V\), and
rank two is the augmentation plane \(u^\perp\) of the residual cyclic
\(C_3\)-trimer
\[
 \mathbb R[C_3]=\mathbb R\mathbf1\oplus u^\perp.
\]
This is a statement about the displayed representation-theoretic descendants.
\end{corollary}


\subsection{Axial triality restriction and half-step compatibility}
\label{p1b:sec:v127-axial-triality-restriction}

Paper~I-B computes the finite rooted axial triality contraction before
passing to the protected chiral representation.  The result sharpens the
v1.26 restriction theorem: the primitive axial contraction itself identifies
the two complementary trimers by the identity map, while the orientation-odd
half-step transports remain separate marked descendants.

\begin{corollary}[Axial finite restriction matches the protected Clifford root]
\label{p1b:cor:v127-axial-finite-clifford-root}
Under the branching of
Theorem~\ref{p1b:thm:v126-spin8-triality-restriction}, choose the rooted axial
vector corresponding to \(-e_1\).  The finite weight contraction on the
complementary trimer is the identity by Paper~I-B, while its protected
chiral factor is the already normalized Clifford multiplication
\(c(-e_1)\).  Thus the finite trimer identity and the protected Clifford map
are two typed restrictions of one rooted triality contraction; no additional
family transport coefficient is introduced.
\end{corollary}

\begin{corollary}[Half-step transport remains a separate rooted descendant]
\label{p1b:cor:v127-half-step-remains-marked}
The triality/Clifford parent does not by itself distinguish the two antiprism
matchings \(X_G\) and \(X_G^\dagger\).  Equivalently, the sign line
\(C_{2,G}=-i(X_G-X_G^\dagger)\) is not generated by the unmarked axial
contraction alone.  A theorem using that line must exhibit the rooted
orientation/half-step constructor that consumes it.
\end{corollary}


\subsection{Rooted q=5/trimer isometry and oriented complex structure}
\label{p1b:sec:v128-q5-trimer-isometry}

The unique primitive non-scalar $q=5$ owner of
Theorem~\ref{p1b:thm:v125-q5-rank2-owner} is the real rotation plane
$u^\perp$.  The rooted Hurwitz trimer supplies another real two-plane: the
augmentation-zero family plane.  For the same oriented Hurwitz diagonal these
planes are canonically and isometrically identified.

\begin{theorem}[Exact augmentation-plane/\texorpdfstring{$q=5$}{q=5} isometry]
\label{p1b:thm:v128-q5-trimer-isometry}
Let
\[
 u=\frac1{\sqrt3}(1,1,1)
\]
and let the three imaginary Hurwitz-link vectors belonging to the rooted
trimer be
\[
 y_0=\frac12(-1,1,1),\qquad
 y_1=\frac12(1,-1,1),\qquad
 y_2=\frac12(1,1,-1).
\]
Define
\[
 G_\perp^{\mathbb R}:=
 \{(a_0,a_1,a_2)\in\mathbb R^3:a_0+a_1+a_2=0\}
\]
and
\[
 \boxed{
 \Psi_u(a_0,a_1,a_2):=a_0y_0+a_1y_1+a_2y_2.}
\]
Then
\[
 \boxed{
 \Psi_u:G_\perp^{\mathbb R}\xrightarrow{\;\cong\;}u^\perp}
\]
is an exact real isometry.
\end{theorem}

\begin{proof}
The map is the imaginary-coordinate restriction of \(L_+\)
to the augmentation plane in
Proposition~\ref{p1b:prop:trimer-transport-overlap}.
There it equals \(a\mapsto-a\), with image \(u^\perp\),
so it is an exact real isometry.
\end{proof}

\begin{theorem}[Rooted \texorpdfstring{$C_3$}{C3} intertwiner]
\label{p1b:thm:v128-q5-c3-intertwiner}
Let $X_G$ be the rooted cyclic shift of the trimer and let $R_u$ be the
already generated $2\pi/3$ rotation of the $q=5$ plane $u^\perp$.  With the
compatible rooted cyclic orientation,
\[
 \boxed{\Psi_uX_G=R_u\Psi_u.}
\]
Reversing the rooted cyclic orientation replaces both $X_G$ and $R_u$ by
their inverses, so the statement is presentation covariant.
\end{theorem}

\begin{proof}
The ambient cyclic coordinate permutation
$(x_1,x_2,x_3)\mapsto(x_3,x_1,x_2)$ fixes $u$ and sends
$y_0\mapsto y_1\mapsto y_2\mapsto y_0$.  Its restriction to $u^\perp$ is
the nontrivial real $C_3$ rotation $R_u$.  Applying both sides to the trimer
basis and extending linearly gives the identity.
\end{proof}

\begin{corollary}[The trimer sign line is the \texorpdfstring{$q=5$}{q=5} complex structure]
\label{p1b:cor:v128-q5-complex-structure-signline}
On $u^\perp$ define
\[
 J_u:=\frac{R_u-R_u^{-1}}{\sqrt3}.
\]
Because $I+R_u+R_u^2=0$ on the nontrivial $C_3$ plane,
\[
 J_u^2=-I_{u^\perp}.
\]
The intertwiner theorem gives
\[
 \boxed{
 \Psi_u^{-1}J_u\Psi_u
 =\frac{X_G-X_G^\dagger}{\sqrt3}.}
\]
Consequently
\[
 \boxed{
 C_{2,G}:=-i(X_G-X_G^\dagger)
 =-i\sqrt3\,\Psi_u^{-1}J_u\Psi_u.}
\]
Thus the primitive D4/Hurwitz sign line is the Hermitian representative of the
canonical oriented complex structure of the same-root $q=5$ owner; no
independent family-orientation datum is required.
\end{corollary}

\begin{remark}[Marked half-step normalization]
\label{p1b:rem:v128-occurrence-firewall}
The corollary identifies the geometry, parity and normalization of the half-step/sign-line direction relative to the stated root and orientation. The full rooted relation cell consists of $\{I_G,X_G,X_G^{-1}\}$; an individual member is specified by its marked edge orbit.
\end{remark}


\subsection{Rooted q=5 local system and associated-owner descent}
\label{p1b:sec:v130-q5-associated-owner}

The preceding same-root isometry identifies the nontrivial real $C_3$ module
but does not yet package the two rooted half-step transports as one generated
finite geometric object.  The Hurwitz-link six-cycle already supplies exactly
that package.  This section makes the transport object explicit and then
records the general scalar-owner retraction carried by any finite unitary
associated module.  The latter is a mathematical descent map of a different
type from the ordinary action trace.

\begin{definition}[Rooted two-edge quotient of the Hurwitz six-cycle]
\label{p1b:def:v130-q5-quotient-graph}
Write the alternating Hurwitz trimers as
\[
 T_+=\{a_k:k\in\mathbb Z_3\},\qquad
 T_-=\{b_k:k\in\mathbb Z_3\},
\]
with the rooted cyclic action $r(a_k)=a_{k+1}$ and
$r(b_k)=b_{k+1}$.  The induced six-cycle of Paper~I-B has edge set
\[
 a_k\sim b_{k+1},\qquad a_k\sim b_{k-1},\qquad k\in\mathbb Z_3.
\]
Let $B_{q5}$ be the multigraph with vertices $+$ and $-$ and two distinct
edges $e_+,e_-:+\to-$.  Define
\[
 \pi(a_k)=+,
 \qquad
 \pi(b_k)=-,
\]
and send the edge $a_kb_{k+1}$ to $e_+$ and $a_kb_{k-1}$ to $e_-$.
\end{definition}

\begin{theorem}[The Hurwitz six-cycle is a regular \texorpdfstring{$C_3$}{C3} covering]
\label{p1b:thm:v130-q5-c3-cover}
The map
\[
 \boxed{\pi:C_6^{\rm link}\longrightarrow B_{q5}}
\]
is a regular three-sheeted graph covering.  Its deck group is the rooted
cyclic group $\langle r\rangle\cong C_3$.
\end{theorem}

\begin{proof}
At every $a_k$ the two incident edges $a_kb_{k+1}$ and $a_kb_{k-1}$ map
bijectively to the two incident edge germs $e_+$ and $e_-$ at $+$.  The same
holds at every $b_k$ after reversing the corresponding edge germs.  Thus
$\pi$ is a graph covering of degree three.  The common shift $r$ preserves
both edge classes and acts freely and transitively on every vertex fibre, so
it is a deck transformation group of order three.  A deck transformation is
determined by its action on one fibre point, hence the full deck group is
exactly $C_3$.
\end{proof}

\begin{definition}[Rooted finite transport constructor]
\label{p1b:def:v130-q5-transport-constructor}
Let
\[
 E_+:=\mathbb C^{T_+},\qquad E_-:=\mathbb C^{T_-},
\]
with basis vectors $e_{a_k},e_{b_k}$.  The aligned axial triality contraction
supplies the rooted identification
\[
 I_0:E_+\longrightarrow E_- ,\qquad I_0e_{a_k}=e_{b_k}.
\]
The two edge classes of $B_{q5}$ lift to the unitary maps
\[
 U_+e_{a_k}=e_{b_{k+1}},\qquad
 U_-e_{a_k}=e_{b_{k-1}}.
\]
In the aligned gauge put
\[
 A_+:=I_0^{-1}U_+,
 \qquad
 A_-:=I_0^{-1}U_-.
\]
For the chosen root orientation,
\[
 A_+=X_G,\qquad A_-=X_G^{-1}=X_G^\dagger.
\]
Define the single rooted transport constructor
\[
 \boxed{
 \mathfrak C_{q5}^{\rm geom}
 :=\mathfrak C_{q5}(C_6^{\rm link},\mathfrak r)
 :=(I_G,C_{1,G},C_{2,G})}
\]
with
\[
 C_{1,G}:=A_++A_-,
 \qquad
 C_{2,G}:=-i(A_+-A_-).
\]
\end{definition}

\begin{theorem}[Single-constructor status and cyclic reconstruction]
\label{p1b:thm:v130-q5-generated-cell}
The three ports of $\mathfrak C_{q5}^{\rm geom}$ are simultaneous outputs of
one rooted Hurwitz-link constructor, not a tuple assembled from independent
ambient endomorphisms.  Moreover
\[
 \boxed{
 X_G=\frac12(C_{1,G}+iC_{2,G}),\qquad
 X_G^{-1}=\frac12(C_{1,G}-iC_{2,G}).}
\]
Thus the Hermitian and cyclic presentations
\[
 (I_G,C_{1,G},C_{2,G})
 \quad\longleftrightarrow\quad
 \{I_G,X_G,X_G^{-1}\}
\]
are two coordinate presentations of the same generated transport object.
Root reversal interchanges $X_G$ and $X_G^{-1}$ and sends
$C_{2,G}\mapsto-C_{2,G}$ while fixing $I_G$ and $C_{1,G}$.
\end{theorem}

\begin{proof}
The rooted graph and its two edge orbits determine $U_+$ and $U_-$, while the
aligned axial contraction determines $I_0$.  Hence $A_+$, $A_-$ and their
Hermitian sum and difference are all fixed before any scalar evaluation.  The two displayed reconstruction formulas follow by adding and
subtracting the definitions of $C_{1,G}$ and $C_{2,G}$.  Reversing the root
orientation exchanges the two edge orbits, so it exchanges $A_+$ and $A_-$.
\end{proof}

In the aligned realization of
Proposition~\ref{p1b:prop:trimer-transport-overlap}, the two
ambient metric overlaps are \(C_{1,G}^{-1}\) and
\(C_{1,G}/2\). They are fixed descendants of the existing
even port of this rooted constructor. The axial identification,
the edge incidence and these metric overlaps retain their
respective definitions.

\begin{theorem}[Augmentation local system equals the primitive \texorpdfstring{$q=5$}{q=5} owner]
\label{p1b:thm:v130-q5-augmentation-local-system}
The regular trimer module decomposes as
\[
 \mathbb R^3=\mathbb R\mathbf 1\oplus G_\perp^{\mathbb R}.
\]
The scalar line is the trivial local subsystem of the covering, while the
augmentation plane carries the nontrivial real $C_3$ monodromy.  Under the
same-root isometry $\Psi_u$ of
Theorem~\ref{p1b:thm:v128-q5-trimer-isometry}, this monodromy is exactly the
primitive $q=5$ rotation:
\[
 \boxed{\Psi_uX_G=R_u\Psi_u.}
\]
Writing
\[
 J_u=\frac{R_u-R_u^{-1}}{\sqrt3},
\]
one equivalently has on $u^\perp$
\[
 \boxed{R_u=-\frac12I+\frac{\sqrt3}{2}J_u.}
\]
Hence the primitive $q=5$ owner is the nontrivial associated real local
system of the rooted Hurwitz $C_3$ covering.
\end{theorem}

\begin{proof}
The decomposition into the invariant scalar line and its zero-sum orthogonal
complement is the usual decomposition of the real permutation module of
$C_3$.  The intertwining statement is
Theorem~\ref{p1b:thm:v128-q5-c3-intertwiner}.  On the nontrivial real plane,
$I+R_u+R_u^2=0$ and $R_u^{-1}=R_u^2$, so
$R_u-R_u^{-1}=2R_u+I=\sqrt3J_u$, which gives the displayed expression.
\end{proof}

\begin{definition}[Canonical graded odd \texorpdfstring{$q=5$}{q=5} geometry]
\label{p1b:def:v130-q5-graded-odd-geometry}
On $\mathcal E_{q5}:=E_+\oplus E_-$ define
\[
 \Gamma_{q5}:=
 \begin{pmatrix}I&0\\0&-I\end{pmatrix},
\]
and the two root-related normalized odd operators
\[
 \mathscr D_+:=
 \begin{pmatrix}0&X_G\\X_G^{-1}&0\end{pmatrix},
 \qquad
 \mathscr D_-:=
 \begin{pmatrix}0&X_G^{-1}\\X_G&0\end{pmatrix}.
\]
\end{definition}

\begin{theorem}[Odd involution and \texorpdfstring{$A_2$}{A2} Coxeter lift]
\label{p1b:thm:v130-q5-coxeter-lift}
The two operators satisfy
\[
 \mathscr D_\pm^\dagger=\mathscr D_\pm,
 \qquad
 \mathscr D_\pm^2=I,
 \qquad
 \Gamma_{q5}\mathscr D_\pm=-\mathscr D_\pm\Gamma_{q5}.
\]
Moreover
\[
 \mathscr D_-\mathscr D_+
 =\begin{pmatrix}X_G&0\\0&X_G^{-1}\end{pmatrix},
 \qquad
 (\mathscr D_-\mathscr D_+)^3=I.
\]
Thus $\mathscr D_+$ and $\mathscr D_-$ obey the Coxeter relations of
$W(A_2)\cong S_3$: the two rooted half-steps are reflection lifts and their
product is the $q=5$ cyclic rotation.
\end{theorem}

\begin{proof}
Self-adjointness, involutivity and oddness follow by block multiplication and
$X_G^{-1}=X_G^\dagger$.  The product formula is direct; its cube is the
identity because $X_G^3=I$.
\end{proof}

\begin{remark}[The complete rooted transport object]
\label{p1b:rem:v130-g15-synchronization}
The object $\mathfrak C_{q5}^{\rm geom}$ and the relation cell $\{I_G,X_G,X_G^{-1}\}$ are constructed on the same marked Hurwitz/trimer data. They specify operator capacity and relations; they do not select coefficients of an arbitrary operator in their span.
\end{remark}

\subsubsection{Finite associated-owner geometry and scalar descent}

\begin{definition}[Finite associated-owner geometry]
\label{p1b:def:v130-associated-owner}
Let $\mathcal G\rightrightarrows B$ be a finite groupoid and let
$E\to B$ be a constant-rank $n$ Hermitian representation with unitary arrow
maps
\[
 \rho(g):E_x\longrightarrow E_y,
 \qquad g:x\to y.
\]
The associated-owner datum is
\[
 \mathfrak A_E:=(\mathcal G\rightrightarrows B,E,\rho,\iota_E,\epsilon_E),
\]
where the scalar inclusion and scalar-owner retraction are
\[
 \iota_E(f)_x:=f(x)I_{E_x},
 \qquad
 \boxed{\epsilon_E(A)(x):=\frac1n\Tr_{E_x}(A_x).}
\]
\end{definition}

\begin{theorem}[Canonical scalar-owner retraction]
\label{p1b:thm:v130-canonical-owner-retraction}
For the datum above,
\[
 \boxed{\epsilon_E\circ\iota_E=\id.}
\]
The retraction is natural under every groupoid arrow:
\[
 \boxed{
 \epsilon_E\!\left(\rho(g)A\rho(g)^{-1}\right)
 =\epsilon_E(A).}
\]
Fibrewise, $\epsilon_E$ is the unique complex-linear functional on
$\End(E_x)$ that is unital and invariant under all unitary conjugations.
\end{theorem}

\begin{proof}
The first identity is
$\epsilon_E(fI)=n^{-1}\Tr(fI)=f$.  Naturality is ordinary trace invariance
under conjugation.  For uniqueness choose an orthonormal basis of a fibre and
let $\lambda:M_n(\mathbb C)\to\mathbb C$ be complex-linear, unital and
unitary-conjugation invariant.  Conjugation by diagonal unitary matrices forces
$\lambda(E_{ij})=0$ for $i\ne j$.  Conjugation by permutation matrices gives
the same value $c$ on every diagonal matrix unit $E_{ii}$.  Since
$1=\lambda(I)=nc$, one has $c=1/n$, hence
$\lambda(A)=n^{-1}\Tr(A)$ for every $A$.
\end{proof}

\begin{corollary}[Amplification invariance of primitive-owner descent]
\label{p1b:cor:v130-owner-amplification-invariance}
Let $F$ be a rank-$m$ Hermitian spectator factor.  Then for every
$A\in\End(E)$,
\[
 \boxed{
 \epsilon_{E\otimes F}(A\otimes I_F)=\epsilon_E(A).}
\]
By contrast the ordinary represented trace scales as
\[
 \Tr_{E\otimes F}(A\otimes I_F)=m\Tr_E(A).
\]
\end{corollary}

\begin{proof}
Using $\rank(E\otimes F)=nm$ and
$\Tr(A\otimes I_F)=m\Tr(A)$ gives
\[
 \epsilon_{E\otimes F}(A\otimes I_F)
 =\frac1{nm}m\Tr(A)=\frac1n\Tr(A).
\]
\end{proof}

\begin{remark}[Ordinary trace and central retraction]
\label{p1b:rem:v130-trace-owner-firewall}
The maps $\operatorname{Tr}_E$ and $\epsilon_E$ have different domains and normalizations. The former is the ordinary matrix trace and counts represented multiplicity; the latter retracts a represented central scalar to the corresponding scalar. Neither is substituted for the other in a quadratic form or determinant without an explicit transport of that form or determinant.
\end{remark}

\subsection{Protected-edge zero face and source normalization}
\label{p1b:sec:v132-fermion-source-zero-face}

The factorwise completion theorem classifies the protected chiral-spinor
edge by the nonnegative coordinate $\varrho_{\rm sp}$. The corollary below
identifies its exact zero face within the same completion cone.

\begin{corollary}[Exact protected chiral zero face]
\label{p1b:cor:v132-protected-chiral-zero-face}
In the level-one completion cone
\[
 \mathfrak M_1\cong\R_{\ge0}^3,
 \qquad (r_{1,0},r_{0,1},\varrho_{\rm sp}),
\]
the face
\[
 \boxed{\varrho_{\rm sp}=0}
\]
is an exact admissible boundary stratum.  On this face the unnormalized
protected chiral-spinor coordinate quadruple vanishes, while the remaining
retained completion data stay inside the same classified completion family.
Consequently the existence of the protected Clifford module on the
nondegenerate stratum does not imply that its unnormalized protected edge is
nonzero in every completion.
\end{corollary}

\begin{proof}
This is the $\varrho_{\rm sp}=0$ face already included in
Theorem~\ref{p1b:thm:product-completion}.  The last sentence of
Corollary~\ref{p1b:cor:product-clifford} states that when a protected magnitude
vanishes, the corresponding unnormalized four coordinates vanish and that
boundary member remains in the completion cone.
\end{proof}

\begin{remark}[Normalized Clifford port versus source actuality]
\label{p1b:rem:v132-normalized-port-not-actuality}
On the nondegenerate stratum $\varrho_{\rm sp}>0$, division by the positive
magnitude produces normalized Clifford generators satisfying the usual
Clifford identities.  Such a normalized nonzero port must not be read
backwards as a theorem that $\varrho_{\rm sp}>0$ is actual.  The normalization
is conditional on the nondegenerate stratum; the zero face above remains an
exact member of the completion space.
\end{remark}


\subsection{Canonical phase-twisted antiprism family primitive}
\label{p1b:sec:v133-antiprism-family-primitive}

The following construction uses the stated one-sphere and exceptional
geometric inputs from Papers~I-A and I-B. Let
\[
 C:=J-I,
 \qquad
 R_G:=\operatorname{diag}(1,-i,i),
\]
let
\[
 \mathcal A_6:=\begin{pmatrix}0&C\\ C&0\end{pmatrix},
 \qquad
 \widehat R:=\operatorname{diag}(R_G,R_G),
\]
and use the already-generated $q=5$ grading
\[
 \Gamma_{q5}:=\operatorname{diag}(I,-I),
 \qquad
 P_A:=\Gamma_{q5}\widehat R
      =\operatorname{diag}(R_G,-R_G).
\]


\begin{theorem}[Canonical phase-twisted antiprism odd primitive]
\label{p1b:thm:v133-antiprism-odd-primitive}
Define
\[
 \boxed{\mathcal D_A:=-i\bigl(\mathcal A_6P_A-P_A^\dagger\mathcal A_6\bigr)}.
\]
Then $\mathcal D_A$ is self-adjoint and $\Gamma_{q5}$-odd, with block form
\[
 \mathcal D_A=\begin{pmatrix}0&F_A\\F_A^\dagger&0\end{pmatrix},
 \qquad F_A=iQ_A,
\]
where
\[
 \boxed{Q_A=CR_G+R_G^\dagger C
 =\begin{pmatrix}
 0&1-i&1+i\\
 1+i&0&2i\\
 1-i&-2i&0
 \end{pmatrix}.}
\]
Its characteristic polynomial factorizes exactly as
\[
 \boxed{\chi_{Q_A}(\lambda)
   =(\lambda+2)(\lambda^2-2\lambda-4)},
\]
so
\[
 \operatorname{spec}Q_A=\{-2,1-\sqrt5,1+\sqrt5\},
 \qquad
 \boxed{\sigma(F_A)=\{\sqrt5-1,2,\sqrt5+1\}}.
\]
In particular $|\det F_A|=8$, hence the chiral block is full rank.
\end{theorem}

\begin{proposition}[Full-rank quadrature selection]
\label{p1b:prop:v133-antiprism-quadrature-selection}
The competing canonical Hermitian quadrature has chiral block
\[
 F_{A,+}:=-CR_G+R_G^\dagger C.
\]
It obeys
\[
 \boxed{\chi_{F_{A,+}^\dagger F_{A,+}}(\lambda)
       =\lambda(\lambda-4)^2},
\]
so $\sigma(F_{A,+})=\{0,2,2\}$.  Therefore, among the two canonical phase
quadratures generated by the same antiprism data, the $-i$ quadrature of
Theorem~\ref{p1b:thm:v133-antiprism-odd-primitive} is the unique one up to
overall sign/orientation that preserves full rank.
\end{proposition}

\begin{proposition}[Canonical polar data and exact paired-Feshbach self-realization]
\label{p1b:prop:v133-antiprism-feshbach-self-realization}
Let
\[
 F_A=U_AP_A^{(+)}
\]
be the polar decomposition.  Since $F_A=iQ_A$ with $Q_A$ Hermitian and
invertible, both $U_A$ and $P_A^{(+)}=|F_A|$ are functions of $Q_A$ and commute.
In Corollary~\ref{p1b:cor:v125-paired-complement-odd-schur}, choose
\[
 A_L=A_R=(P_A^{(+)})^{1/2},
 \qquad U=U_A.
\]
Then
\[
 A_LUA_R^\dagger=F_A,
\]
and hence the Schur descendant is exactly
\[
 \boxed{-BD_Q^{-1}B^\dagger=-\Delta^{-1}\mathcal D_A.}
\]
Thus the phase-twisted antiprism primitive is an exact mathematical
paired-Feshbach descendant of its own canonical polar data.
\end{proposition}

\begin{remark}[The antiprism operator family]
\label{p1b:rem:v133-antiprism-occurrence-firewall}
The statements specify $\mathcal D_A$, its polar factors and its Schur block on their stated finite Hilbert spaces. A particular quadratic form using them must retain its displayed coefficients and metric.
\end{remark}


\subsection{Native-owner transfers and the preterminal family trimer}
\label{p1b:sec:v137-native-owner-preterminal}
The primitive ranks $2,4,6$ of
Corollary~\ref{p1b:cor:v126-triality-primitive-owner-descendants} are actual
owner ranks, not merely integers used by a later spectrum table.  Their
canonical coevaluation must therefore be retained at owner level until the specified operator contraction has been formed.

\begin{theorem}[Native unit transfer on an associated owner]
\label{p1b:thm:v137-native-unit-transfer}
Let $E$ be a generated finite Hermitian owner of rank $n$ with the inherited
Hilbert--Schmidt metric on $\operatorname{End}(E)$.  The already constituted
associated-owner inclusion and retraction are
\[
 \iota_E:\mathbb C\to\operatorname{End}(E),\qquad z\mapsto zI_E,
 \qquad
 \epsilon_E(A)=\frac1n\operatorname{Tr}_E A.
\]
They obey
\[
 \boxed{
 \iota_E^\dagger=\operatorname{Tr}_E,\qquad
 \iota_E^\dagger\iota_E=nI,\qquad
 \iota_E^+=\epsilon_E,
 }
\]
and therefore have unique nonzero singular magnitudes
\[
 \boxed{\sigma(\iota_E)=\sqrt n,\qquad
        \sigma(\epsilon_E)=\frac1{\sqrt n}.}
\]
For the generated owners $E_2=(u^\perp)_\mathbb C$, $E_4=V_\mathbb C$ and
$E_6=(\Lambda^2V)_\mathbb C$, this gives
$\sqrt2,2,\sqrt6$ and their reciprocals on the actual owners.
\end{theorem}

\begin{proof}
For every $A\in\operatorname{End}(E)$,
$\langle zI_E,A\rangle_{\rm HS}=\bar z\operatorname{Tr}A$ and
$\|I_E\|_{\rm HS}^2=n$.  The adjoint and Gram identities follow, and the
Moore--Penrose inverse of the injective rank-one map is
$(\iota_E^\dagger\iota_E)^{-1}\iota_E^\dagger=\operatorname{Tr}/n$.
\end{proof}

\begin{theorem}[The rooted $q=5$ owner carries the complete family trimer]
\label{p1b:thm:v137-q5-operator-trimer}
Let $E=u^\perp$ be the real $q=5$ owner identified with the rooted family
augmentation plane by
Theorem~\ref{p1b:thm:v128-q5-trimer-isometry}.  Let $r_0,r_1,r_2$ be the three
rooted unit rays in $E$ and put $P_j=|r_j\rangle\langle r_j|$.  For
$a=(a_0,a_1,a_2)\in G_{\mathbb R}$ define
\[
 \mathcal F(a)=\sum_{j=0}^2a_jP_j,
 \qquad
 M:=\mathcal F^\dagger\mathcal F
   =\frac34I+\frac14\mathbf1\mathbf1^T,
 \qquad
 \mathcal U:=\mathcal F M^{-1/2}.
\]
Then
\[
 \boxed{\mathcal U:G_{\mathbb R}\xrightarrow{\ \cong\ }
              \mathcal S(E)}
\]
is an isometry onto the three-dimensional real symmetric-operator space,
intertwines the rooted $C_3$ action, and satisfies
\[
 \boxed{
 \mathcal U\frac{\mathbf1}{\sqrt3}=\frac{I_E}{\sqrt2},\qquad
 \mathcal U(G_\perp)=\mathcal S_0(E).}
\]
Thus the normalized rank-two identity transfer is exactly the scalar family
line, while the two traceless symmetric channels are the remaining family
directions.
\end{theorem}

\begin{proof}
Three unit rays separated by $120^\circ$ obey
$\operatorname{Tr}(P_jP_j)=1$ and
$\operatorname{Tr}(P_iP_j)=1/4$ for $i\ne j$, yielding the displayed Gram
$M$.  It is positive with eigenvalues $3/2,3/4,3/4$, so $\mathcal U$ is
isometric.  Since $\sum_jP_j=(3/2)I_E$, the normalized constant vector maps to
$I_E/\sqrt2$; orthogonality gives the traceless complement.  Rooted cyclic
transport permutes the $P_j$, proving equivariance.
\end{proof}

\begin{corollary}[Early scalar closure loses the augmentation channels]
\label{p1b:cor:v137-early-trace-loss}
The scalar retraction $\epsilon_E\circ\mathcal F$ has rank one and vanishes on
$G_\perp$, whereas $\mathcal F$ has rank three.  Hence tracing the owner before
forming any required common-parent incidence discards exactly two family
channels.
\end{corollary}

\begin{remark}[Transfer domains]
\label{p1b:rem:v137-preterminal-occurrence-firewall}
The maps above act on their stated owner and trimer spaces. Compositions retain the intermediate Hilbert metric and are formed before any trace or scalar evaluation.
\end{remark}


\subsection{Rooted hierarchy--family block bridge}
\label{p1b:sec:v138-hierarchy-family-bridge}
The rank-$1\oplus2$ Clifford corridor used by the rooted hierarchy and the
rank-$1\oplus2$ rooted family decomposition are descendants of the same rooted
Hurwitz data.  This section records the exact intertwiner before any terminal
scalar terminal evaluation is taken.

\begin{theorem}[Rooted hierarchy--family isometry]
\label{p1b:thm:v138-hierarchy-family-isometry}
Let
\[
 c:=\frac1{\sqrt3}(1,1,1)\in G_{\mathbb R},\qquad
 Q_{\rm fam,0}:=|c\rangle\langle c|,\qquad
 Q_{\rm fam,\perp}:=I_G-Q_{\rm fam,0}.
\]
Let $u$ be the rooted axial imaginary quaternion and let
$\Psi_u:G_{\perp}^{\mathbb R}\to u^\perp$ be the exact rooted $q=5$
isometry of Theorem~\ref{p1b:thm:v128-q5-trimer-isometry}.  Define
\[
 \boxed{
 J_{\rm HG}:=|u\rangle\langle c|+\Psi_uQ_{\rm fam,\perp}:
 G_{\mathbb R}\longrightarrow \mathbb Ru\oplus u^\perp.}
\]
Then $J_{\rm HG}$ is an exact real isometry, extends unitarily after
complexification, and obeys
\[
 \boxed{
 J_{\rm HG}Q_{\rm fam,0}J_{\rm HG}^\dagger=Q_4,\qquad
 J_{\rm HG}Q_{\rm fam,\perp}J_{\rm HG}^\dagger=Q_5.}
\]
The rooted $C_3$ action on $G_\perp$ is carried to the already generated
$q=5$ rotation on $u^\perp$.
\end{theorem}

\begin{proof}
The two summands have orthogonal domains and orthogonal images.  The axial
summand is unitary on the one-dimensional line and $\Psi_u$ is an exact
isometry on the augmentation plane.  Hence $J_{\rm HG}^\dagger J_{\rm HG}=I_G$.
The projector identities follow from its images.  The $C_3$ statement is
Theorem~\ref{p1b:thm:v128-q5-c3-intertwiner} on the augmentation plane and is
trivial on the fixed line.
\end{proof}

\begin{theorem}[Hierarchy blocks are the rooted family-operator blocks]
\label{p1b:thm:v138-hierarchy-family-block-isomorphism}
Let
\[
 \Phi_{\rm HG}(A):=J_{\rm HG}^\dagger A J_{\rm HG}.
\]
On the $Q_4\oplus Q_5$ Clifford corridor of
Proposition~\ref{p1b:prop:block-edges}, $\Phi_{\rm HG}$ is a $*$-isomorphism
onto $\operatorname{End}(G)$ and maps the hierarchy blocks as
\[
 \boxed{
 \begin{aligned}
 Q_4\operatorname{End}Q_4
   &\xrightarrow{\ \Phi_{\rm HG}\ }
     Q_{\rm fam,0}\operatorname{End}(G)Q_{\rm fam,0},\\
 Q_4\operatorname{End}Q_5\oplus Q_5\operatorname{End}Q_4
   &\xrightarrow{\ \Phi_{\rm HG}\ }
     Q_{\rm fam,0}\operatorname{End}(G)Q_{\rm fam,\perp}
     \oplus Q_{\rm fam,\perp}\operatorname{End}(G)Q_{\rm fam,0},\\
 Q_5\operatorname{End}Q_5
   &\xrightarrow{\ \Phi_{\rm HG}\ }
     Q_{\rm fam,\perp}\operatorname{End}(G)Q_{\rm fam,\perp}.
 \end{aligned}}
\]
Their dimensions are $1,4,4$ and their balanced hierarchy labels are
respectively $32,36,40$.  Thus the $E_{32}/E_{36}/E_{40}$ corridor is an
explicit rooted decomposition of the full $9$-dimensional family-operator
algebra; no external three-space embedding is required.
\end{theorem}

\begin{proof}
Conjugation by the unitary $J_{\rm HG}$ is a $*$-isomorphism and carries
$Q_4,Q_5$ to $Q_{\rm fam,0},Q_{\rm fam,\perp}$.  Apply it to each matrix
corner.  The hierarchy labels and multiplicities are those of
Proposition~\ref{p1b:prop:block-edges}.
\end{proof}

\begin{corollary}[The $E_{36}$ block is the native singlet--augmentation incidence edge]
\label{p1b:cor:v138-e36-family-incidence}
For $v\in u^\perp$ let
$T(v)=Q_5L_vQ_4=|vu\rangle\langle u|$ and put
$g_v:=J_{\rm HG}^\dagger(vu)\in G_\perp$.  Then
\[
 \boxed{
 \Phi_{\rm HG}(T(v))=|g_v\rangle\langle c|,\qquad
 \Phi_{\rm HG}(T(v)^\dagger)=|c\rangle\langle g_v|.}
\]
Moreover
\[
 \Phi_{\rm HG}(T(v))^\dagger\Phi_{\rm HG}(T(w))
 =\langle v,w\rangle Q_{\rm fam,0}.
\]
Hence the midpoint hierarchy scale $E_{36}^2=E_{32}E_{40}$ belongs to the
native off-diagonal family-incidence block and not to a third independent
carrier.
\end{corollary}

\begin{remark}[Incidence and spectral data]
\label{p1b:rem:v138-mass-incidence-firewall}
The rooted carrier/host isomorphism transports incidence and orientation. It does not identify an independently specified list of scalar coefficients with eigenvalues of a family operator.
\end{remark}

\subsection{Complete rooted C3 relation cell}
\label{p1b:sec:v139-rooted-c3-relation-cell}
The two alternating rooted trimers are regular $C_3$ modules.  The root
involution and the two oriented short-edge half-steps therefore form the
complete equivariant relation space between them, not three unrelated family
choices.

\begin{theorem}[Complete $C_3$-equivariant trimer intertwiners]
\label{p1b:thm:v139-complete-c3-intertwiners}
Let $E_+$ and $E_-$ be the two ordered rooted trimer carriers with their common
cyclic action.  Let
\[
 I_0e_{p_j}=e_{m_j},\qquad
 U_+e_{p_j}=e_{m_{j+1}},\qquad
 U_-e_{p_j}=e_{m_{j-1}}
\]
(indices modulo $3$).  Then
\[
 \boxed{
 \operatorname{Hom}_{C_3}(E_+,E_-)
 =\operatorname{span}_{\mathbb C}\{I_0,U_+,U_-\}.}
\]
After alignment by $I_0$, this is
\[
 \boxed{\operatorname{span}\{I,X,X^{-1}\}\cong\mathbb C[C_3].}
\]
\end{theorem}

\begin{proof}
In aligned cyclic coordinates an intertwiner $T$ obeys $XT=TX$.  Hence its
matrix is circulant and is uniquely of the form
$aI+bX+cX^{-1}$.  Multiplying on the left by $I_0$ gives the stated basis
between the two labelled trimer copies.
\end{proof}

\begin{corollary}[Orbit and character presentations are Fourier dual]
\label{p1b:cor:v139-c3-orbit-character-duality}
With $\omega=e^{2\pi i/3}$, the Fourier idempotents
\[
 \Pi_0=\frac13(I+X+X^{-1}),\quad
 \Pi_+=\frac13(I+\omega^2X+\omega X^{-1}),\quad
 \Pi_-=\frac13(I+\omega X+\omega^2X^{-1})
\]
are the three rank-one character lines of the same relation cell.  Thus the
transport basis $\{I,X,X^{-1}\}$ and the rooted family flag
$\{Q_0,E_+,E_-\}$ are two presentations of one generated three-state relation,
up to the already retained common cyclic relabelling and complete rooted
reversal.
\end{corollary}

\begin{corollary}[Unique channel representative of every hierarchy incidence]
\label{p1b:cor:v139-hierarchy-incidence-channel-representative}
Combine Theorem~\ref{p1b:thm:v138-hierarchy-family-block-isomorphism} with
Paper~I-B's signed triality isomorphism
$\chi(W)=D_\eta W D_\eta$.  Every hierarchy-owned incidence
$\iota\in\operatorname{End}(G)$ therefore has the unique rooted channel
representative
\[
 \boxed{W_\iota=\chi^{-1}(\iota)=D_\eta\iota D_\eta.}
\]
No additional three-space matching is needed to compare an owner-resolved operator with the rooted family carrier.
\end{corollary}

\begin{remark}[Incidence versus singular values]
\label{p1b:rem:v139-c3-mass-semantics}
The relation cell transports incidence and orientation. Eigenvalues and singular values are computed from the particular operator and Hilbert metric supplied in each construction.
\end{remark}


\part{Reciprocal response, operators and principal descent}
\section{Primitive transfers, Born compatibility, and boundary response}
\label{p1b:sec:organized-boundary}

\subsection{Full primitive owner support and inverse-transfer transport}
\label{p1b:sec:v2168-primitive-owner-support}
\subsubsection{The full primitive lift, including its Moore--Penrose inverse}
Let $E_n$ be a native rank-$n$ Hermitian owner. On $\End(E_n)$ retain the
Hilbert--Schmidt pairing. Let $\Omega_n$ be the counting vector with $n$ unit
entries. Define the two active unit directions
\[
 \omega_n=\frac{\Omega_n}{\sqrt n},\qquad
 s_n=\frac{I_{E_n}}{\sqrt n},
\]
and the support-preserving partial isometry
\[
 \Lambda_n=|s_n\rangle\langle\omega_n|:
 \mathbb C^n\longrightarrow\End(E_n).
\]
The native unit transfer and its retraction are
$\iota_{E_n}(z)=zI_{E_n}$ and $\epsilon_{E_n}(X)=\Tr(X)/n$.

\begin{theorem}[Operator-level native primitive transport]
\label{p1b:thm:v2168-primitive-owner-support}
For the full counting transfer $N_n(z)=z\Omega_n$,
\[
 \boxed{\begin{gathered}
 \Lambda_n^\dagger\Lambda_n=|\omega_n\rangle\langle\omega_n|,
 \qquad\Lambda_n\Lambda_n^\dagger=|s_n\rangle\langle s_n|,\\
 \Lambda_nN_n=\iota_{E_n},\qquad
 \epsilon_{E_n}\Lambda_n=N_n^+,\qquad
 N_n^+\Lambda_n^\dagger=\epsilon_{E_n}.
 \end{gathered}}
\]
Thus both directions of the primitive transfer commute with the owner lift
as full-domain rank-one maps, with their inherited active-support metrics.
The complementary kernels have different dimensions and are not identified
by an isometry. In particular
\[
 \boxed{s_4=I_{E_4}/2,\qquad\epsilon_{E_4}(s_4)=1/2.}
\]
\end{theorem}
\begin{proof}
Both active directions have norm one. Their outer product gives the two
support projections. Since $N_n=\sqrt n\,|\omega_n\rangle$,
$\Lambda_nN_n=\sqrt n\,|s_n\rangle=\iota_{E_n}$.
Also $\epsilon_{E_n}=\langle s_n|/\sqrt n$ and
$N_n^+=\langle\omega_n|/\sqrt n$. Multiplying by $\Lambda_n$ or its adjoint
gives the last two identities on the entire respective domains.
For $n=4$, $\Tr(I_4/2)/4=1/2$.
\end{proof}

The active primitive direction is the displayed coevaluation/support map. The calculation retains the complete operator transfer, including its phase and intermediate metric.


\subsection{Native local harmonic boundary map and the IR family bridge}
\label{p1b:sec:v2208-neutral-boundary-geometry}
This section integrates geometric results on their stated carriers.
The six-point port belongs to the existing rooted Hurwitz family; the port coordinates are transported together with the rooted family. On the full
degree-twelve round scalar harmonic space with probability Haar metric,
\[
 \dim\mathcal H_{12}=169,\qquad
 \Delta_{S^3_R}|_{\mathcal H_{12}}=(168/R^2)I.
\]
The harmonic kernel is the existing normalized kernel of
\ref{p1b:der:g09:eq-G09-3}. The harmonic carrier and the boundary/trimer carriers retain their distinct definitions and explicit maps.

\subsubsection{Native six-point support}
For \(\varepsilon\in\{\pm1\}^3\), the generated Hurwitz cube-link points are
\[
 h_\varepsilon=\frac12
 (1+\varepsilon_1 i+\varepsilon_2 j+\varepsilon_3 k).
\]
Choose a calculation representative with body diagonal
\((+,+,+)\leftrightarrow(-,-,-)\). The parent decomposition gives
\[
 T_+=\{(-,+,+),(+,-,+),(+,+,-)\},\qquad T_-=-T_+,
\]
and the local support
\[
 \Omega=\{h_\varepsilon:\varepsilon\in T_+\sqcup T_-\}.
\]
All six points have real part \(1/2\). Opposite imaginary trimer points
are conjugate quaternions, not antipodal quaternions.

\begin{proposition}[Native regular action]
\label{p1b:proposition:atlas-native-regular-action-2b29373}% stable atlas label; provenance v2_277

The unoriented diagonal stabilizer \(D_3\cong S_3\) acts freely and
transitively on \(\Omega\). Thus \(\mathbb C^\Omega\cong\mathbb C[S_3]\).
\end{proposition}
\begin{proof}
For a permutation \(p\) of the three imaginary coordinates, the action is
\(v\mapsto\operatorname{sgn}(p)P_pv\). It is orientation-preserving on
\(\mathbb R^3\) because
\(\det(\operatorname{sgn}(p)P_p)=\operatorname{sgn}(p)^4=1\).
The even three-cycle moves through one trimer; odd permutations reverse
the diagonal and move to the other trimer. The orbit has six points.
Its stabilizer in a group of order six is therefore trivial.
The binary lift is \(2D_3\); its central sign acts trivially by quaternion
conjugation, so this is the projective \(S_3\) action, not an embedding
of the ordinary symmetric group into \(SU(2)\).
\end{proof}

\subsubsection{Evaluation and its exact Gram}
With probability Haar normalization, the degree-\(n\) kernel is
\[
 K_n(q,h)=(n+1)U_n(q\cdot h),\qquad
 U_n(\cos\theta)=\frac{\sin((n+1)\theta)}{\sin\theta}.
\]
Endpoint values are by continuity. In particular
\(K_{12}(q,q)=169\).
Define the raw evaluation map and its adjoint by
\begin{equation}\label{p1b:eq:v2208-boundary-evaluation}
 E:\mathcal H_{12}\longrightarrow\mathbb C^\Omega,\qquad
 (Ef)_a=f(h_a),\qquad
 (E^\dagger c)(q)=\sum_{a=1}^{6}c_aK_{12}(q,h_a).
\end{equation}
The codomain carries its standard counting inner product.

\begin{theorem}[Exact local evaluation Gram]\label{p1b:thm:v2208-boundary-gram}
\[
 G:=EE^\dagger=156I_6+13J_6.
\]
Consequently \(E\) has rank six.
\end{theorem}
\begin{proof}
The native inner products are
\[
 h_\varepsilon\cdot h_\delta
 =\frac{1+\varepsilon\cdot\delta}{4}.
\]
For distinct non-axial cube points these are \(-1/2,0,1/2\).
The Chebyshev recurrence \(U_{n+1}(t)=2tU_n(t)-U_{n-1}(t)\) gives
\[
 U_{12}(-1/2)=U_{12}(0)=U_{12}(1/2)=1,\qquad U_{12}(1)=13.
\]
Thus the diagonal entries of \(G\) are \(169\), and every off-diagonal
entry is \(13\). Writing \(P_{\Omega,0}=J_6/6\), \(P_{\Omega,a}=I_6-P_{\Omega,0}\), one obtains
\[
 G=234P_{\Omega,0}+156P_{\Omega,a}>0.
\]
Both eigenvalues are positive, hence the rank is six.
\end{proof}

\subsubsection{The derived map \texorpdfstring{\(B\)}{B}}
\begin{theorem}[Unique positive normalized restriction]\label{p1b:thm:v2208-boundary-coisometry}
The positive codomain normalization of \(E\) is
\begin{equation}\label{p1b:eq:v2208-boundary-map}
 \boxed{B=G^{-1/2}E
 =\left(\frac{P_{\Omega,a}}{\sqrt{156}}+
         \frac{P_{\Omega,0}}{\sqrt{234}}\right)E.}
\end{equation}
It satisfies \(BB^\dagger=I_6\), and \(\ker B=\ker E\).
Writing \(\bar f=\frac16\sum_af(h_a)\), its coordinate formula is
\begin{equation}\label{p1b:eq:v2208-boundary-coordinates}
 \boxed{(Bf)_a=\frac{f(h_a)-\bar f}{\sqrt{156}}
                    +\frac{\bar f}{\sqrt{234}}.}
\end{equation}
\end{theorem}
\begin{proof}
The spectral resolution of Theorem \ref{p1b:thm:v2208-boundary-gram} gives
\[
 G^{-1/2}=P_{\Omega,a}/\sqrt{156}+P_{\Omega,0}/\sqrt{234}.
\]
Hence \(BB^\dagger=G^{-1/2}GG^{-1/2}=I_6\).
Invertibility of \(G^{-1/2}\) preserves the kernel.
For uniqueness, let \(C>0\) satisfy \(CGC=I\).
Then the positive operator \(G^{1/2}CG^{1/2}\) has square \(G\).
Uniqueness of the positive square root gives
\(G^{1/2}CG^{1/2}=G^{1/2}\), hence \(C=G^{-1/2}\).
The coordinate formula follows from \(P_{\Omega,0}Ef=\bar f\,\mathbf1\).
\end{proof}

\begin{remark}[Counting versus probability coordinates]
Using uniform probability weights on \(\Omega\) replaces the coordinate
evaluation by \(E/\sqrt6\), with Gram \(G/6\). Its positive normalization
gives the same \(B\) under the canonical unitary coordinate identification.
Similarly, consistent sphere-metric and kernel changes preserve
the normalized map under the corresponding unitary coordinate identifications. Both conventions transport the same stated Hilbert metric.
\end{remark}

\subsubsection{Kernel and the exact \texorpdfstring{\(163+5+1\)}{163+5+1} decomposition}
\begin{theorem}[Explicit orthogonal kernel projector]\label{p1b:thm:v2208-boundary-projectors}
The operators
\[
 P_{\rm sing}=\frac1{234}E^\dagger P_{\Omega,0}E,\qquad
 P_{\rm aug}=\frac1{156}E^\dagger P_{\Omega,a}E,\qquad
 Q=I_{\mathcal H_{12}}-P_{\rm sing}-P_{\rm aug}
\]
are mutually orthogonal projectors, with
\[
 \boxed{\rank Q=163,\quad\rank P_{\rm aug}=5,\quad
        \rank P_{\rm sing}=1.}
\]
Their ranges give
\[
 \boxed{\mathcal H_{12}=\ker B\oplus
 B^\dagger\ker\varepsilon\oplus B^\dagger\mathbb C\mathbf1,}
 \qquad \varepsilon(c)=\sum_ac_a.
\]
The kernel projector is explicitly
\begin{equation}\label{p1b:eq:v2208-kernel-coordinates}
 \boxed{(Qf)(q)=f(q)-\sum_{a=1}^6K_{12}(q,h_a)
 \left(\frac{f(h_a)-\bar f}{156}+\frac{\bar f}{234}\right).}
\end{equation}
\end{theorem}
\begin{proof}
Since \(BB^\dagger=I_6\), the adjoint \(B^\dagger\) is an isometry.
It transports \(P_{\Omega,0},P_{\Omega,a}\) to orthogonal projectors
\(B^\dagger P_{\Omega,0}B=P_{\rm sing}\) and
\(B^\dagger P_{\Omega,a}B=P_{\rm aug}\), of ranks one and five.
Their sum is \(B^\dagger B=E^\dagger G^{-1}E\), the visible projector.
Its orthogonal complement projects onto \(\ker B\) and has rank
\(169-6=163\). Expanding \(G^{-1}=P_{\Omega,a}/156+P_{\Omega,0}/234\) gives
\eqref{p1b:eq:v2208-kernel-coordinates}.
\end{proof}

\begin{corollary}[Exact Born norm decomposition]
\label{p1b:corollary:atlas-exact-born-norm-decomposition-f1b6c68}% stable atlas label; provenance v2_277

The map \(f\mapsto(Bf,Qf)\) is unitary from \(\mathcal H_{12}\) to
\(\mathbb C^\Omega\oplus\ker B\), and
\[
 \|f\|^2=\|P_{\rm sing}f\|^2+\|P_{\rm aug}f\|^2+\|Qf\|^2.
\]
These ranks become probabilities \(1/169,5/169,163/169\) only for the
additional maximally mixed state \(I/169\), not for an unspecified source.
\end{corollary}

\subsubsection{Covariance and character certificate}
For \(g\in S_3\), let \(U_gf(q)=f(g^{-1}\cdot q)\) and let \(R_g\)
permute the boundary coordinates. Isometric kernel covariance gives
\[
 EU_g=R_gE,\quad R_gG=GR_g,\quad BU_g=R_gB.
\]
Thus each lifted projector commutes with this stabilizer.
For any ambient isometry \(F\), with the entire parent local port transported,
\[
 B_{F\Omega}U_F=B_\Omega,\qquad
 Q_{F\Omega}=U_FQ_\Omega U_F^\dagger.
\]
This constructs the equivariant family on the existing parent orbit.
The whole isometry orbit, rather than one preferred representative, is retained.

For the conjugation action, \(\mathcal H_{12}\) is
\(V_6\otimes V_6^\ast\). Its character on the identity, transpositions,
and three-cycles is \((169,1,1)\); the two nonidentity values follow by
squaring the spin-six characters at rotation angles \(\pi\) and \(2\pi/3\).
Character averaging therefore yields
\[
 \mathcal H_{12}\cong29\,\mathbf1\oplus28\,\mathrm{sgn}
                  \oplus56\,\mathrm{std}
              \cong\mathbf1\oplus28\,\mathbb C[S_3].
\]
The visible range is one regular representation, so
\[
 \ker B\cong\mathbf1\oplus27\,\mathbb C[S_3],
 \qquad \dim(\mathcal H_{12})^{S_3}=29.
\]
The \(163\) kernel is a restriction kernel, not the invariant quotient.
The boundary augmentation has character \((5,-1,-1)\) and is
\(\mathrm{sgn}\oplus2\,\mathrm{std}\).
The spin-two quadrupole instead has character \((5,1,-1)\) and is
\(\mathbf1\oplus2\,\mathrm{std}\) under this \(S_3\) action. Equal dimension
does not identify these carriers.


\begin{theorem}[Complete nonsinglet owner, canonical Green Gram and Schur residual]
\label{p1b:thm:v239-nonsinglet-green-schur}
On the existing degree-twelve harmonic owner, define only from the
already established boundary projectors
\[
 P_{\rm ns}:=I_{\mathcal H_{12}}-P_{\rm sing}=Q+P_{\rm aug},
 \qquad E_{168}:=\operatorname{ran}P_{\rm ns}.
\]
Then $\operatorname{rank}P_{\rm ns}=168$ and, for the rooted stabilizer
representation of this section,
\[
 \boxed{E_{168}\cong 28\,\mathbb C[S_3],\qquad
        E_{163}=\ker B\cong\mathbf1\oplus27\,\mathbb C[S_3].}
\]
For the \emph{dimensionless round harmonic kinetic operator}
$K_{12}:=R^2\Delta_{S^3_R}|_{\mathcal H_{12}}=168I$, put
\[
 \mathcal G^{(K)}_{ab}=\operatorname{Tr}_{\mathcal H_{12}}
    (P_aK_{12}^{-1}P_b),\qquad
 (P_1,P_2)=(P_{\rm ns},P_{\rm aug}).
\]
Its exact two-channel return is
\begin{equation}\label{p1b:eq:v239-nonsinglet-green-gram}
 \boxed{\mathcal G^{(K)}=
 \begin{pmatrix}1&5/168\\[2pt]5/168&5/168\end{pmatrix},
 \quad \mathcal G^{(K)}_{11}
       -\mathcal G^{(K)}_{12}(\mathcal G^{(K)}_{22})^{-1}
                      \mathcal G^{(K)}_{21}=\frac{163}{168}.}
\end{equation}
Thus the mixed trace equals $5/168$ \emph{before} the augmentation
channel is Schur-eliminated, while the reduced channel returns $163/168$.
The equal number $168=12(12+2)=(12+1)^2-1$ follows from the round
harmonic eigenvalue and the nonsinglet rank; its additional factorization
$168=7\cdot24$ does not, by itself, construct a Clifford--Hurwitz
intertwiner.
\end{theorem}
\begin{proof}
The rank and decomposition follow from
Theorem~\ref{p1b:thm:v2208-boundary-projectors}. The projector
$P_{\rm sing}$ carries the single trivial summand of the visible regular
$S_3$ module. Removing it from
$\mathcal H_{12}\cong\mathbf1\oplus28\,\mathbb C[S_3]$ gives the
claimed nonsinglet module. Since $K_{12}^{-1}=I/168$,
$\mathcal G^{(K)}_{ab}=\Tr(P_aP_b)/168$. Orthogonality of $Q$
and $P_{\rm aug}$ yields the three entries, and direct scalar Schur
reduction gives $1-(5/168)^2/(5/168)=163/168$. Equivalently,
$P_{\rm ns}-P_{\rm aug}=Q$ is the Gram-orthogonal residual.
\end{proof}
\begin{remark}[Source, mass and reduction-order firewall]
\label{p1b:rem:v239-nonsinglet-action-scope}
The matrix \eqref{p1b:eq:v239-nonsinglet-green-gram} is an exact
geometric quadratic return \emph{for the stated round kinetic inverse}.
It is not a licensed K56 matter action, neutral pole, mass-squared
splitting, or probability for an unspecified state. In particular the
complete $168$-dimensional owner and the reduced $163$-dimensional
kernel have different cross relations with $P_{\rm aug}$:
$P_{\rm ns}P_{\rm aug}=P_{\rm aug}$ but $QP_{\rm aug}=0$.
Replacing one by the other before an action's common contraction
would discard a generated mixed return unless an action theorem proves
that elimination occurs there. No rank-to-terminal-role selector is
inferred from the numeric value $5/168$.
\end{remark}

\begin{proposition}[Orthogonal harmonic Green reduction and equivariant nonuniqueness]
\label{p1b:prop:v240-k56-orthogonal-green-firewall}
On the precise round harmonic owner of
Theorem~\ref{p1b:thm:v239-nonsinglet-green-schur}, the two traces
$\mathcal G^{(K)}_{\rm ns,aug}=5/168$ and
$\mathcal G^{(K)}_{\rm Schur}=163/168$ are traces of overlapping
\emph{observable projectors}. Their orthogonal carrier counterpart satisfies
\[
 \boxed{Q_{163}K_{12}^{-1}P_{\rm aug}=0,\qquad
        P_{\rm ns}-P_{\rm aug}=Q_{163}.}
\]
In particular the geometric $5/168$ overlap does not furnish an
off-diagonal neutral matter Hessian between $E_{163}$ and $E_5$.
For the inherited complex $S_3$ representations,
\[
 E_{163}\cong28\,\mathbf1\oplus27\,\mathrm{sgn}
                 \oplus54\,\mathrm{std},\qquad
 E_5\cong\mathrm{sgn}\oplus2\,\mathrm{std},
\]
and hence
\[
 \boxed{\dim_{\mathbb C}\operatorname{Hom}_{S_3}(E_{163},E_5)
          =27+54\cdot2=135.}
\]
Neither these ranks nor $S_3$-equivariance select an action-owned
kernel--augmentation contraction, a family operator, or a neutral mass.
\end{proposition}
\begin{proof}
Here $K_{12}^{-1}=I/168$ and $Q_{163}P_{\rm aug}=0$; the
first claims follow by projector multiplication. The full harmonic
representation is $\mathbf1\oplus28\mathbb C[S_3]$, and the visible
augmentation is $\mathrm{sgn}\oplus2\mathrm{std}$; subtract their
multiplicities to obtain the displayed kernel type. The dimension of
an intertwiner space is the sum of products of like irreducible
multiplicities. The physical source requires its separately typed
same-action contraction.
\end{proof}


\subsubsection{A native geometric bridge to the family carrier}
For the existing rooted Clifford geometry put
\[
 E_3(u)=\mathbb R\mathbf1\oplus u^\perp,\quad
 \Pi_3=Q_7+Q_5,\quad Q_7=P_0,\quad
 (\rank Q_7,\rank Q_5)=(1,2).
\]
The endpoint and the following bridge are defined on these displayed geometric carriers. The inherited plane map and rooted bridge are
\ref{p1b:thm:v128-q5-trimer-isometry} and
\ref{p1b:thm:v138-hierarchy-family-isometry}.

Use exactly the family decomposition and rooted plane isometry already
present in this paper:
\[
 c=\frac1{\sqrt3}(1,1,1)^T,\qquad Q_c=cc^\dagger,
 \qquad Q_\perp=I_G-Q_c,
 \qquad\Psi_u:G_\perp^{\mathbb R}\longrightarrow u^\perp.
\]
The existing $\Psi_u$ is isometric and intertwines the rooted cyclic shift
$X_G$ with the $2\pi/3$ plane rotation induced by $\operatorname{Ad}_{c_u}$.
With the printed coordinates, $\Psi_u(a)=-a$ for $\sum_i a_i=0$.

\begin{theorem}[Exact IR--family geometric isometry]
\label{p1b:thm:v2208-ir-family-bridge}
Given the unit-line identification $c\mapsto\mathbf1$ and the already fixed
rooted plane map $\Psi_u$, the unique linear extension is
\[
 \boxed{J_{\rm IR}=|\mathbf1\rangle\langle c|+\Psi_uQ_\perp:
 G_{\mathbb R}\longrightarrow E_3(u).}
\]
It is isometric, extends unitarily onto the complexified corridor, and
obeys
\[
 J_{\rm IR}^\dagger J_{\rm IR}=I_G,\qquad
 J_{\rm IR}J_{\rm IR}^\dagger=\Pi_3,\qquad
 J_{\rm IR}X_G=\operatorname{Ad}_{c_u}J_{\rm IR},
\]
\[
 \boxed{J_{\rm IR}^\dagger P_0J_{\rm IR}=Q_c
 =\frac{I_G+X_G+X_G^2}{3},\qquad
 J_{\rm IR}^\dagger Q_5J_{\rm IR}=Q_\perp.}
\]
\end{theorem}
\begin{proof}
The two domain summands are orthogonal, as are their images. The first
summand maps one normalized vector to the normalized Clifford unit, and
$\Psi_u$ is the existing plane isometry. Their direct sum gives the two
isometry identities and fixes the map on a direct-sum basis, proving
uniqueness with the stated restrictions. The scalar line is fixed by both
cyclic actions; on the plane the intertwiner is the existing theorem.
$P_0$ and $Q_5$ project onto the respective images. Finally the cyclic
average of the rooted three-cycle is the constant-line projector $Q_c$.
\end{proof}

\begin{theorem}[Full operator bridge and its relative partial isometry]
\label{p1b:thm:v2208-ir-operator-bridge}
The map
\[
 \Phi_{\rm IR}(A)=J_{\rm IR}^\dagger A J_{\rm IR}
\]
is a $*$-isomorphism from
$\Pi_3\operatorname{End}(\mathcal A_{\mathbb C})\Pi_3$ onto
$\operatorname{End}(G_{\mathbb C})$. The corridor identity $\Pi_3$ maps to
$I_G$. With the already generated rooted bridge
$J_{\rm HG}=|u\rangle\langle c|+\Psi_uQ_\perp$,
\[
 \boxed{J_{\rm IR}J_{\rm HG}^\dagger
 =|\mathbf1\rangle\langle u|+P_{u^\perp}.}
\]
Thus the relative partial isometry replaces the axial line by the scalar
line and preserves the same rooted plane.
\end{theorem}
\begin{proof}
On the corner, $J_{\rm IR}J_{\rm IR}^\dagger=\Pi_3$ shows that conjugation
preserves products and adjoints. Its inverse is
$B\mapsto J_{\rm IR}B J_{\rm IR}^\dagger$, so it is onto all nine matrix
units. Multiplying the two bridge formulas cancels the cross terms and
uses $\Psi_u\Psi_u^\dagger=P_{u^\perp}$, giving the displayed relative map.
\end{proof}
The relative partial isometry compares the two stated decompositions of the same finite carrier.

\subsubsection{Conditional harmonic amplitude control}
This control compares a particular linear harmonic coupling with the
earlier rank trace proposal. It does not identify $\mathcal H_{12}$ as the
actual K56 input. Three maps must be kept separate: Clifford $O(3)$ Haar
projection, $O(3)$ conjugation averaging on harmonics, and translation-Haar
averaging of functions on $S^3$.

\begin{proposition}[Scalar harmonic amplitude is not a rank trace]
\label{p1b:prop:v2208-h12-scalar-control}
In the full degree-twelve round harmonic carrier, the normalized
$O(3)$ conjugation-invariant line is generated by
\[
 Z_{12}(q)=U_{12}(\operatorname{Re}q).
\]
It has probability-Haar norm one, translation-Haar mean zero, and value one
at each of the six native Hurwitz points. For the earlier boundary
singlet, augmentation and kernel projectors,
\[
 \boxed{\bigl(\langle Z_{12},P_{\rm sing}Z_{12}\rangle,
 \langle Z_{12},P_{\rm aug}Z_{12}\rangle,
 \langle Z_{12},Q_{163}Z_{12}\rangle\bigr)
       =(1/39,0,38/39).}
\]
The rotational scalar line differs from the local boundary singlet line;
with $\Pi_Z=|Z_{12}\rangle\langle Z_{12}|$,
\[
 \|[\Pi_Z,Q_{163}]\|_{\rm HS}^2=76/1521.
\]
\end{proposition}
\begin{proof}
The radial harmonic equation on $S^3$, restricted to functions of
$t=\operatorname{Re}q$, is the Gegenbauer equation of parameter one;
its degree-twelve polynomial solution is the one-dimensional line
$\mathbb R U_{12}(t)$. Under probability Haar measure the $t$ density is
$(2/\pi)\sqrt{1-t^2}$. Its moments are
$\mathbb E[t^{2k}]=\binom{2k}{k}/[4^k(k+1)]$ and odd moments vanish.
Inserting the exact polynomial
\[
 U_{12}(t)=\sum_{k=0}^{6}(-1)^k\binom{12-k}{k}(2t)^{12-2k}
\]
gives mean zero, norm one, and $U_{12}(1/2)=1$. Thus $EZ_{12}=\mathbf1_6$.
The exact boundary Gram is $EE^\dagger=234P_{\rm sing}^{\Omega}
+156P_{\rm aug}^{\Omega}$. The polar-normalized boundary map gives
$\|P_{\rm sing}Z_{12}\|^2=6/234=1/39$ and zero augmentation weight.
Orthogonality leaves kernel weight $38/39$. For a rank-one projector onto
a unit vector and any orthogonal projector $Q$,
$\|[\Pi_Z,Q]\|_{\rm HS}^2=2p(1-p)$ with
$p=\langle Z,QZ\rangle$. Substitution gives $76/1521$.
\end{proof}
A linear coupling through that one normalized rotational scalar therefore
does not recover the rank counts $163$ and $5$. A collective trace on
$\operatorname{End}(\ker B)$ or on the augmentation owner has a different
domain and can retain those counts before core terminalization. Collective traces and linear amplitudes retain these distinct domains.


\begin{remark}[Domains of the harmonic and trimer maps]
\label{p1b:rem:v2208-neutral-origin-scope}
The map $B$, its $163+5+1$ orthogonal decomposition, the isometry $J_{\rm IR}$ and the induced operator isomorphism are assertions on their displayed harmonic, boundary and trimer Hilbert spaces. A map to any other carrier requires a specified intertwiner.
\end{remark}

\subsection{Reciprocal Born compatibility and finite boundary returns}
\label{p1b:sec:v2210-lit-reciprocal-transfers}
This section extends the supplied reciprocal atlas by consequences of its
actual mixed pairing, finite reduction and rooted intertwiners.  The positive
Born metric below is mathematical geometry relative to a specified factor
metric.  

\begin{theorem}[Compatible positive Born triple from the reciprocal pairing]
\label{p1b:thm:v2210-lit-born-triple}
Use the nondegenerate two-form $\Omega$ of
Theorem~\ref{p1b:thm:transgression} and the product splitting
$TM_8=L_x\oplus L_p$.  Let $K$ be $+I$ on $L_x$ and $-I$ on $L_p$,
and set
\[
 \eta(X,Y):=\Omega(KX,Y).
\]
Then $\eta$ is a symmetric nondegenerate metric of signature $(4,4)$,
$\eta(KX,KY)=-\eta(X,Y)$, and both eigendistributions of $K$ are
integrable.  For a specified positive metric $h$ on $L_x$, define the
isomorphism $C:L_p\to L_x$ by
\[
 h(Cv,u)=\eta(v,u),\qquad u\in L_x,\quad v\in L_p,
\]
and the positive metric
\[
 \mathcal H(u+v,u'+v')=h(u,u')+h(Cv,Cv').
\]
As maps $TM_8\to T^*M_8$, these tensors obey
\[
 \eta^{-1}\mathcal H=\mathcal H^{-1}\eta,
 \qquad \Omega^{-1}\mathcal H=-\mathcal H^{-1}\Omega.
\]
Thus they form a positive Born triple on the supplied product.  For fixed
$\Omega$, splitting and $h$, this is the unique compatible positive metric
which restricts to $h$ on $L_x$ and makes $L_x,L_p$ orthogonal.
\end{theorem}
\begin{proof}
The transgression form pairs the two three-dimensional vector factors,
while $\omega_{\rm sc}$ pairs the two scalar directions.  It vanishes on
each factor separately and its mixed pairing is nondegenerate.  Hence
$\Omega(KX,KY)=-\Omega(X,Y)$, which proves symmetry of $\eta$; an
$\eta$-dual frame for the two factors gives its split signature.
The distributions are tangent to the product factors, so their brackets
remain in the same factor.  Nondegeneracy of the mixed pairing makes $C$
invertible.  The endomorphism $J=\eta^{-1}\mathcal H$ exchanges the factors:
\[
 Ju=C^{-1}u,\qquad Jv=Cv.
\]
Consequently $J^2=I$ and $JK=-KJ$.  Since
$\eta^{-1}\Omega=K$, the two stated identities follow from $J^2=K^2=I$
and anticommutation.  Conversely those identities force the metric on
$L_p$, after identifying it with $L_x^*$ by $\eta$, to be the dual of $h$;
this gives the displayed expression and uniqueness.
\end{proof}

\begin{corollary}[The available Born connection has a positive-metric owner]
\label{p1b:cor:v2210-lit-born-connection-scope}
For the triple of Theorem~\ref{p1b:thm:v2210-lit-born-triple}, the canonical
$D$-bracket and the unique compatible connection of zero generalized torsion
are supplied by the Born-connection theorem
\cite[Theorem~1]{FreidelRudolphSvoboda2019Born}.  Their existence requires no
closedness of $\Omega$.  The common structure group for this positive triple is $O(4)$.  The independently
specified product metric agrees with $\mathcal H$ only when its second
factor metric is the dual metric prescribed above.
\end{corollary}
\begin{proof}
The preceding theorem checks the compatibility hypotheses of the cited
theorem, rather than inferring them from a shared name.  Its structure-group
intersection is $\operatorname{Sp}(8)\cap O(4,4)\cap O(8)=O(4)$.
Nonclosedness of the native $\Omega$ is unchanged.  The metric-equality
condition is the uniqueness clause of the preceding theorem.
\end{proof}
The ordinary symplectic/K\"unneth integrable specialization cannot be
applied globally to this unchanged compact $M_8$: the obstruction in
Proposition~\ref{p1b:prop:lcs} remains.  Equality with another connection requires an explicit comparison of their covariant derivatives.

\begin{theorem}[Finite boundary Weyl function from a complete Hermitian owner]
\label{p1b:thm:v2210-lit-finite-weyl}
Let a supplied complete finite Hermitian operator split as
\[
 H=\begin{pmatrix}A&B\\B^\dagger&C\end{pmatrix}
 \quad\hbox{on}\quad\mathcal K\oplus\mathcal E.
\]
On $z\notin\operatorname{spec}C$, define its retained response
\[
 M_H(z)=zI_{\mathcal K}-A+B(C-zI_{\mathcal E})^{-1}B^\dagger.
\]
For $\operatorname{Im}z>0$, this response is analytic and obeys
\[
 \operatorname{Im}M_H(z)
 =(\operatorname{Im}z)
 \left[I_{\mathcal K}+B(C-\bar z)^{-1}(C-z)^{-1}B^\dagger\right]>0.
\]
Where the inverses exist, the complete retained resolvent is exactly
\[
 P_{\mathcal K}(zI-H)^{-1}|_{\mathcal K}=M_H(z)^{-1}.
\]
Successive eliminations give the same function on their common regular
domain.  At a real regular value $E$ with a simple visible zero
$M_H(E)v=0$, its retained pole residue is
\[
 \underset{z=E}{\operatorname{Res}}\,M_H(z)^{-1}
 =\frac{vv^\dagger}{v^\dagger M_H'(E)v},\qquad
 M_H'(E)=I+B(C-E)^{-2}B^\dagger>0.
\]
\end{theorem}
\begin{proof}
For $R(z)=(C-z)^{-1}$, the resolvent identity gives
$\operatorname{Im}R(z)=(\operatorname{Im}z)R(z)^\dagger R(z)$.
This proves positivity.  The Schur complement of the lower block in
$zI-H$ is $zI-A-B(zI-C)^{-1}B^\dagger=M_H(z)$; block inversion gives
the retained resolvent, and uniqueness of that inverse gives complete
elimination invariance.  Real differentiation gives the stated $M_H'$.
For a simple zero, expansion on the null line and inversion on its
complement give the displayed rank-one residue.
\end{proof}
This is the finite same-owner version of the Weyl/Nevanlinna and resolvent
tools used in boundary-triple theory
\cite{LitGEODerkachHassiMalamud2017Boundary}.  It constructs the actual finite
response once $H$ and its retained/environment split are supplied; the split is part of the input data.


\begin{theorem}[Minimal visible Hermitian realization of a finite response]
\label{p1b:thm:response-minimal-visible-realization}
For the complete matrix response of
Theorem~\ref{p1b:thm:v2210-lit-finite-weyl}, write
\[
 M_H(z)=zI-A+\sum_{\lambda}\frac{R_\lambda}{\lambda-z},
 \qquad R_\lambda=BP_\lambda B^\dagger\succeq0,
\]
where only nonzero coefficient matrices are included.  Exact knowledge of
this matrix-valued rational function determines $A$, the distinct visible
pole positions $\lambda$, and all $R_\lambda$.  Its least Hermitian
environment, relative to these retained tests, is
\[
 \mathcal E_{\min}=\bigoplus_\lambda\operatorname{ran}R_\lambda,
 \qquad C_{\min}=\bigoplus_\lambda\lambda I,
 \qquad B_{\min}((v_\lambda))=\sum_\lambda R_\lambda^{1/2}v_\lambda,
\]
and obeys
\[
 \boxed{\dim\mathcal E_{\min}
       =\sum_\lambda\operatorname{rank}R_\lambda.}
\]
It reproduces $M_H$ and is unique up to a unitary on its environment which
intertwines $C$ and the coupling $B$, with the retained frame fixed.
In every realization the source-visible reducing subspace and its
orthogonal complement are
\[
 \mathcal E_{\rm vis}
  =\operatorname{span}\{C^kB^\dagger v:k\ge0,\ v\in\mathcal K\},
 \qquad \mathcal E_{\rm dark}=\mathcal E_{\rm vis}^{\perp}.
\]
The response determines the visible restriction, while
$B\mathcal E_{\rm dark}=0$ and $C|_{\mathcal E_{\rm dark}}$ is
unconstrained by this response.
\end{theorem}
\begin{proof}
The polynomial part at infinity recovers $A$; uniqueness of partial
fractions recovers the distinct poles and their coefficient matrices.
The stated square-root factors have Gram matrices $R_\lambda$, so they
reproduce every summand.  In any realization
$R_\lambda=(BP_\lambda)(BP_\lambda)^\dagger$ requires at least
$\operatorname{rank}R_\lambda$ dimensions at that eigenvalue.
Equality is attained by the construction.  Any two full-column-rank
factors with the same Gram matrix differ by a unitary on their domains;
these polewise unitaries intertwine the scalar $C$ blocks and combine to
prove uniqueness.  Spectral interpolation shows that the cyclic visible
space is the orthogonal sum of the ranges of $P_\lambda B^\dagger$.
It is $C$-invariant, and Hermiticity makes its complement reducing.
Orthogonality to $\operatorname{ran}B^\dagger$ then gives
$B\mathcal E_{\rm dark}=0$.
\end{proof}
The positive matrices $R_\lambda$ are the coefficients of
$(\lambda-z)^{-1}$, hence the complex residues of $M_H$ itself have the
opposite sign.  A degenerate complementary eigenvalue contributes one
matrix $R_\lambda$, whose rank counts its visible multiplicity; it is
not counted once for each arbitrarily chosen eigenvector.  Partial pole
data or scalar source contractions leave an inverse family.  Minimality
for these tests does not remove a mode consumed by another declared test,
.

\begin{corollary}[Nonzero visible response retains exact energy dependence]
\label{p1b:cor:response-exact-memory}
On a real regular interval,
\[
 \frac{\mathrm d}{\mathrm dE}\bigl[M_H(E)-EI\bigr]
 =\sum_\lambda\frac{R_\lambda}{(\lambda-E)^2}\succeq0.
\]
The environmental term is constant on a nonempty open interval only if
all $R_\lambda$ vanish, equivalently $B=0$ in that realization.
\end{corollary}
\begin{proof}
Differentiate the partial fractions.  A zero sum of these positive
semidefinite matrices forces each nonzero coefficient to vanish.
Finally $\sum_\lambda R_\lambda=BB^\dagger$.
\end{proof}
An approximate local replacement still needs its actual energy band and
error bound; a nonzero exact return cannot be replaced by a constant for
free.

\begin{theorem}[Positive complementary gap couples lowering and retained weight]
\label{p1b:thm:response-positive-gap-weight}
Retain the same energy-independent Hermitian $H$, split and canonical
Hilbert norm as in Theorem~\ref{p1b:thm:v2210-lit-finite-weyl}.
At a real eigenvalue $E$ suppose
\[
 0<\delta I\preceq C-EI\preceq DI,\qquad
 R=(C-EI)^{-1},\qquad M_H(E)u=0,\quad u^\dagger u=1.
\]
Put $s=u^\dagger BRB^\dagger u\ge0$.  The reconstructed eigenvector
$(u,-RB^\dagger u)$ has normalized retained population
\[
 w_{\rm ret}=\frac{1}{1+u^\dagger BR^2B^\dagger u},\qquad
 \boxed{\frac{\delta}{\delta+s}\le w_{\rm ret}
                    \le\frac{D}{D+s}.}
\]
Here $s=u^\dagger Au-E$ is the lowering relative to the direct-block
Rayleigh value in that same direction.  At a simple visible pole the
retained resolvent residue is exactly $w_{\rm ret}uu^\dagger$.

For a possibly degenerate zero let $U$ be an isometry onto $\ker M_H(E)$
and define
\[
 S=U^\dagger BRB^\dagger U,\qquad
 T=U^\dagger BR^2B^\dagger U.
\]
Then the entire retained matrix residue is
\[
 \mathop{\rm Res}_{z=E}M_H(z)^{-1}=U(I+T)^{-1}U^\dagger,
\]
with the basis-independent bounds
\[
 (I+S/\delta)^{-1}\preceq(I+T)^{-1}
                         \preceq(I+S/D)^{-1}.
\]
\end{theorem}
\begin{proof}
The full eigenvector equation gives the complementary component
$-RB^\dagger u$ and $u^\dagger Au-E=s$.  Its squared norm is
$1+u^\dagger BR^2B^\dagger u$.  Spectral calculus on $C-EI$ gives
$R/D\preceq R^2\preceq R/\delta$, proving the scalar inequalities.
For degeneracy the lifted null frame has Gram matrix $I+T$.
The full spectral projector is therefore
$LU(I+T)^{-1}U^\dagger L^\dagger$, where
$L=\left(\begin{smallmatrix}I\\-RB^\dagger\end{smallmatrix}\right)$.  Compressing it proves the residue
formula.  Congruence gives $S/D\preceq T\preceq S/\delta$;
order reversal under inversion proves the matrix inequalities without
assuming that $S$ and $T$ commute.
\end{proof}
At a degenerate pole a chosen lifted vector's population is not by itself
the complete matrix residue.  Retained-source residues are contractions
of that complete residue; a source with complementary support also needs
the inherited effective source.  The bound constrains one fixed eliminated
sector and does not apply to an indefinite complement, an arbitrary
energy-dependent inverse propagator, or a continued resonance without
additional hypotheses.

\begin{corollary}[Visible zero modes cannot be inverted away]
\label{p1b:cor:v2210-lit-boundary-zero-mode}
If $P_\lambda$ is a spectral projector of $C$, the pole of $M_H$ at
$z=\lambda$ is removable exactly when $BP_\lambda=0$.  In particular a
constant-mode zero eigenvalue of an assigned boundary operator can be
eliminated at $z=0$ only when its actual source coupling annihilates that
mode or when the mode is retained in the retained block.
\end{corollary}
\begin{proof}
The coefficient of $(\lambda-z)^{-1}$ is $BP_\lambda B^\dagger$.
It vanishes if and only if $BP_\lambda=0$, since it is the Gram matrix of
$BP_\lambda$.  Otherwise the original regular-domain inverse does not
exist at the eigenvalue.
\end{proof}
The continuum boundary-triple or determinant-gluing transfer still requires
the actual Green identity, operator domain, traces and boundary conditions.
A rational finite response alone neither proves a smooth product collar
nor replaces its geometric zero-mode and normal-orientation data.

\begin{theorem}[Rooted unitary transport preserves complete source jets]
\label{p1b:thm:v2210-lit-unitary-source-jets}
Let $s$ range over a regular source/metric parameter domain and let $U(s)$
be a smooth unitary intertwiner between the same finite owner presentations.
For every consumed operator $A(s)$, set
$\widetilde A=UAU^\dagger$ and $\Theta=\mathrm dU\,U^\dagger$.
Then
\[
 \mathrm d\widetilde A-[\Theta,\widetilde A]
 =U(\mathrm dA)U^\dagger.
\]
The transported derivative extends by the product rule and intertwines
every ordered mixed jet.  Finite traces of words and source contractions
$j^\dagger F(A_1,\ldots,A_k)j$, with $\widetilde j=Uj$, have identical
values and ordinary mixed derivatives in both presentations.
\end{theorem}
\begin{proof}
Differentiate $U^\dagger U=I$ to obtain
$\mathrm dU^\dagger=-U^\dagger\mathrm dU\,U^\dagger$.  Differentiating
$UAU^\dagger$ gives the displayed commutator identity.  Applying it
recursively proves the jet statement.  Cyclicity of the trace and
$U^\dagger U=I$ give scalar equality at every parameter; ordinary
derivatives of that scalar equality agree as well.
\end{proof}
For an isometric embedding rather than a unitary owner identification, use
the inherited image and its trace and include the omitted-image defects of
Lemma~\ref{p1b:lem:v2167-compression-multiplication}.  A rank-fifteen
fibre-sampling map is not an identification with the scalar left-regular
$\mathbb C^{2T}$ target.  These transport results preserve the already
generated ranks, rooted trimer relations and finite split outputs.


\subsection{Born pencils, conjugate boundary energies and source transport}
\label{p1b:sec:v223-born-boundary}
The following results are finite geometric constructions.  They extend the
Born triple of Theorem~\ref{p1b:thm:v2210-lit-born-triple}, the collar
construction of Theorem~\ref{p1b:thm:v124-return-coordinate-collar}, and
complete source transport beyond the unitary case.  Their physical use
requires the actual operator, source and boundary maps; a scalar fit does
not establish those maps.

\begin{theorem}[Weighted Born pencil and positive reciprocal chart]
\label{p1b:thm:v223-born-pencil}
On the positive compatible real eight-dimensional fibre let
$J^2=K^2=I$, $JK=-KJ$, with $J,K$ self-adjoint.  For $x>0$ put
\[
 A_x=J+xK,\quad G(x)=\sqrt{1+x^2},\quad
 U_c=(J+K)/\sqrt2,\quad t(x)=\frac{x}{1+G(x)}.
\]
Then $A_x^2=G(x)^2I$, $U_c$ is an isometric involution, and
\[
 U_cA_xU_c=xA_{1/x},\qquad G(1/x)=G(x)/x,\qquad
 t(1/x)=\frac{1-t(x)}{1+t(x)}.
\]
The positive reciprocal fixed point is $x=1$, $t=\sqrt2-1$.
Within the real affine class $A_x=(1+ax)J+bxK$, normalized by $A_0=J$,
the displayed weighted covariance for every $x>0$ forces $a=0,b=1$.
\end{theorem}
\begin{proof}
Anticommutation cancels the cross term in $A_x^2$ and gives
$U_c^2=I$, $U_cJU_c=K$, $U_cKU_c=J$.  Self-adjointness then proves
isometry.  Rationalizing $t$ proves the reciprocal formula.  For the
general affine member the covariance defect is
\[
 [(b-1)x-a]J+[ax+1-b]K.
\]
Linear independence of $J,K$ and its vanishing for all $x$ give the
claimed coefficients.  Solving the reciprocal fixed-point equation
gives the stated positive root.
\end{proof}
This is uniqueness in a specified affine operator class.  It does not
derive the affine physical port or identify $U_c$ with every inherited
marked Born lift.  The eight-dimensional central identity does not by
itself identify a physical scalar constituent or its trace weight.

\begin{theorem}[Square-root extension and conjugate-flux cap]
\label{p1b:thm:v223-extension-flux}
Let $\mathcal A\succeq g_0^2I>0$ on a finite positive fibre.  For
$v(0)=w$, minimize
\[
 \mathcal E[v]=\frac12\int_0^\infty
 (\|v_y\|^2+\langle v,\mathcal A v\rangle)\,dy
\]
on the finite-energy decaying domain.  The unique minimizer is
$v(y)=e^{-y\sqrt{\mathcal A}}w$, its outward response at zero is
$-v_y(0)=\sqrt{\mathcal A}w$, and
\[
 E_D(w)=\tfrac12\langle w,\sqrt{\mathcal A}w\rangle,
 \qquad E_D^*(p)=\tfrac12\langle p,\mathcal A^{-1/2}p\rangle.
\]
For $\mathcal A=A_x^2$, a scalar conjugate-flux port with $\Delta>0$
and prescribed total flux $w=2\sqrt\Delta\,r$ has
\[
 \min_{a+b=w}\tfrac12(a^2+G^{-1}b^2)
 =\frac{2\Delta r^2}{1+G},\qquad
 a=\frac{w}{1+G},\quad b=\frac{Gw}{1+G}.
\]
\end{theorem}
\begin{proof}
Diagonalize $\mathcal A$.  On each eigenspace the Euler equation is
$-v_{yy}+\mathcal A v=0$; finite energy removes its growing solution.
Integration by parts gives $E_D$ and positivity gives uniqueness.
Completing the square in $\langle p,w\rangle-E_D(w)$ proves the
convex dual.  A Lagrange multiplier for $a+b=w$ gives $b=Ga$ and the
displayed value.  Displacing $(a,b)$ by $(z,-z)$ increases the value by
$z^2(1+G^{-1})/2$, proving it is the unique minimum.
\end{proof}
The stationary value of $E_D(w)-\langle p,w\rangle$ is $-E_D^*(p)$;
the boundary work fixes this sign.  Joining the raw Dirichlet weights
$1,G$ instead gives the normalized coefficient $2G/(1+G)$, rather
than $2/(1+G)$.  Thus the boundary role and flux incidence matter.
The extension parameter is an analytic normal coordinate, not an added
physical spacetime dimension.  On the spectral domain the identity
\[
 e^{-y\sqrt{\mathcal A}}=\frac{y}{2\sqrt\pi}
 \int_0^\infty s^{-3/2}e^{-y^2/(4s)}e^{-s\mathcal A}\,ds
\]
follows eigenvalue by eigenvalue from the Gaussian integral.  It is
Poisson subordination of the exponential heat semigroup, consistent
with the extension results of Caffarelli--Silvestre and
Stinga--Torrea\cite{CaffarelliSilvestre2007,StingaTorrea2010Extension}.
It does not identify the native gravitational heat generator with
$A_x^2$.  The operator, form domain, normal orientation, endpoint
conditions and curvature-to-flux map remain part of any such identity,
as in Remark~\ref{p1b:rem:v124-boundary-gluing-owner-firewall}.

\begin{theorem}[Positive metric comparison of complete source words]
\label{p1b:thm:v223-positive-source-transport}
Let $P,Q$ be smooth positive definite real $m$-dimensional metrics in a
fixed reference presentation.  There is a unique symmetric positive
matrix $C$ with $C^{\mathsf T}QC=P$, namely
\[
 C=Q^{-1/2}(Q^{1/2}PQ^{1/2})^{1/2}Q^{-1/2}.
\]
More generally let $T_s=A^{\mathsf T}T_rA$, where $A$ is smooth and
invertible.  For $m$ by $n$ matrices set
$Y_r=AY_s$, $j_r=A^{-\mathsf T}j_s$.  Then the complete quadratic
source word
\[
 W_T(Y,j)=\operatorname{Tr}(Y^{\mathsf T}TY)
                   +2\operatorname{Tr}(j^{\mathsf T}Y)
\]
agrees off shell in the two presentations.  Its stationary value is
$-\operatorname{Tr}(j^{\mathsf T}T^{-1}j)$, and all ordinary mixed
parameter derivatives of the pulled-back word and value agree.
The induced finite volumes agree as well:
\[
 (\det T_r)^{n/2}|\det A|^n\,d^{mn}Y_s
       =(\det T_s)^{n/2}\,d^{mn}Y_s.
\]
Consequently the complete convergent Gaussian integrals coincide.
\end{theorem}
\begin{proof}
Writing $D=Q^{1/2}CQ^{1/2}$ reduces the first assertion to the unique
positive square root $D^2=Q^{1/2}PQ^{1/2}$.  Substitution in the word
proves its equality, including the linear source.  Completing its
square gives the stationary value; differentiating the identity proves
every jet.  The matrix-entry metric is $I_n\otimes T$, so its volume
is $(\det T)^{n/2}dY$.  Each of the $n$ columns contributes $|\det A|$
to the Jacobian, proving the measure identity and the integral claim.
\end{proof}
The canonical positive comparison is relative to the specified reference
metric and presentations.  Existence of an isometry neither selects an
incoming physical section nor proves that this comparison is the native
source map.  For differentiated source fields its connection, derivative
and boundary terms must also be transported.  Tensoring with a spectator
preserves the identities while retaining its full multiplicity.

\begin{theorem}[Noncommuting positive tangential perturbation]
\label{p1b:thm:v223-root-perturbation}
On a finite positive spectral band let $\mathcal A\succeq g_0^2I>0$,
$V\succeq0$, $\|V\|\le\nu$.  Set $S_0=\sqrt{\mathcal A}$ and
$S_1=\sqrt{\mathcal A+V}$.  Without a commutation hypothesis,
\[
 \|S_1-S_0\|\le\frac{\nu}{2g_0},\qquad
 \|2(I+S_1)^{-1}-2(I+S_0)^{-1}\|
       \le\frac{\nu}{g_0(1+g_0)^2}.
\]
\end{theorem}
\begin{proof}
For $X=S_1-S_0$, $S_1X+XS_0=V$.  Its unique solution is
$X=\int_0^\infty e^{-tS_1}Ve^{-tS_0}\,dt$, yielding the first bound.
The resolvent identity and $S_i\succeq g_0I$ yield the second.
\end{proof}
These are Euclidean finite-band operator estimates.  Actual metric and
source jets, unbounded domains and Lorentzian continuation require their
own comparison; this norm estimate does not supply them.


\subsection{Reference Poisson defects and complete retained responses}
\label{p1b:sec:v224-reference-poisson}
The positive extension of Theorem~\ref{p1b:thm:v223-extension-flux}
admits the following source-complete translation. All matrices below
act on a finite positive fibre; an unbounded application additionally
requires its common form domain and convergence. The normal coordinate
is an analytic extension coordinate.

\begin{theorem}[Noncommuting reference-Poisson defect]
\label{p1b:thm:v224-poisson-defect}
Let $\mathcal A_r\succeq g_0^2I>0$ be a smooth family, $K_r=
\sqrt{\mathcal A_r}$, and hold $r_0$ and a boundary matrix $B_0$ fixed.
On the decaying finite-energy domain set
\[
 \mathscr E_r[Q]=\frac12\int_0^\infty
 \operatorname{Tr}(Q_y^\dagger Q_y+Q^\dagger\mathcal A_rQ)\,dy,
 \quad Q(0)=B_0,\quad Q_0=e^{-yK_0}B_0,
 \quad K_0=K_{r_0}.
\]
For $Q=Q_0-F$, $F(0)=0$, the exact off-shell identity is
\begin{align}
 \mathscr E_r[Q_0-F]&=\mathscr E_r[Q_0]
 +\tfrac12\|F\|_{L_r}^2-\operatorname{Re}\langle j_r,F\rangle,
 &L_r&=-\partial_y^2+\mathcal A_r,\nonumber\\
 j_r&=(\mathcal A_r-\mathcal A_0)Q_0.
 \label{p1b:eq:v224-poisson-defect}
\end{align}
The unique stationary fields and complete energy are
\[
 Q_*=e^{-yK_r}B_0,\quad F_*=(e^{-yK_0}-e^{-yK_r})B_0,
 \qquad E_* =\tfrac12\operatorname{Tr}(B_0^\dagger K_rB_0).
\]
If $K_0S_r+S_rK_0=\mathcal A_0+\mathcal A_r$, then
$\mathscr E_r[Q_0]=\tfrac12\operatorname{Tr}(B_0^\dagger S_rB_0)$.
No pair of the displayed matrices is required to commute.
\end{theorem}
\begin{proof}
Expand the full quadratic form. Integrate the cross derivative onto
$Q_0$; the boundary term vanishes because $F(0)=0$ and the fields decay.
Since $L_0Q_0=0$, the cross term is exactly
$-\operatorname{Re}\langle(\mathcal A_r-\mathcal A_0)Q_0,F\rangle$.
The positive gap proves strict convexity. Solving the normal equation
gives the decaying $Q_*$ and integration by parts gives its energy.
For the reference energy, integrate
$e^{-yK_0}(\mathcal A_0+\mathcal A_r)e^{-yK_0}$; the resulting
matrix is the unique positive Sylvester solution $S_r$. This order of
the factors proves the assertion also for noncommuting data.
\end{proof}
Differentiating this functional identity retains the source, reference
energy and every moving-coordinate contact together. The affine
translation has absolute Jacobian one at a finite regulator; it preserves
the original quadratic determinant rather than cancelling it. A zero
boundary datum produces a zero defect source. Thus the theorem does
not select a nonzero physical incoming datum from an operator alone.

\begin{proposition}[Noncommuting root jets and sourced support]
\label{p1b:prop:v224-root-source-jets}
Put $\mathcal L_K(X)=KX+XK$. For smooth positive $\mathcal A=K^2$,
\[
 \mathcal L_K^{-1}(V)=\int_0^\infty e^{-yK}Ve^{-yK}\,dy,\qquad
 K_a=\mathcal L_K^{-1}(\mathcal A_a),\qquad
 K_{ab}=\mathcal L_K^{-1}(\mathcal A_{ab}-K_aK_b-K_bK_a).
\]
At fixed boundary datum, $(E_*)_a=\tfrac12\operatorname{Tr}
(B_0^\dagger K_aB_0)$, with the analogous mixed second formula.
If the baseline is $k^2I$ and $V\succeq0$, then
\[
 0\le \tfrac12\operatorname{Tr}
 [B_0^\dagger(\sqrt{k^2I+V}-kI)B_0]
 \le\frac1{4k}\operatorname{Tr}(B_0^\dagger VB_0).
\]
Its vanishing is equivalent to $VB_0=0$.
\end{proposition}
\begin{proof}
Differentiate $K^2=\mathcal A$ once and twice and invert its positive
Sylvester map by the displayed convergent integral. Differentiate the
fixed-datum energy. For the scalar baseline diagonalize $V$ and use
$0\le\sqrt{k^2+v}-k\le v/(2k)$, with equality zero precisely at $v=0$.
Summing the nonnegative source-weighted terms proves the support test.
\end{proof}

\begin{proposition}[Canonical auxiliary trace and full Schur source]
\label{p1b:prop:v224-canonical-heat-source}
For a positive finite Euclidean word $\operatorname{Tr}(Le^{-A})$,
let $T=e^A$ and $J=L^{1/2}$. Then
\[
 \operatorname{Tr}(Le^{-A})=
 \sup_X\{2\operatorname{ReTr}(J^\dagger X)
                       -\operatorname{Tr}(X^\dagger TX)\}.
\]
For $T=\left(\begin{smallmatrix}T_{PP}&B\\B^\dagger&C\end{smallmatrix}\right)$,
elimination of the $Q$ rows yields
\[
 T_{\rm eff}=T_{PP}-BC^{-1}B^\dagger,\qquad
 J_{\rm eff}=J_P-BC^{-1}J_Q,
 \qquad +\operatorname{Tr}(J_Q^\dagger C^{-1}J_Q)
\]
in the supremum word. The last term is retained in its value and jets.
\end{proposition}
\begin{proof}
Complete the square around $X=T^{-1}J$ and use cyclicity of the finite
trace. For fixed $X_P$, complete the square in
$X_Q-C^{-1}(J_Q-B^\dagger X_P)$; the remaining terms give precisely
the displayed kernel, source and centre.
\end{proof}
This is canonical in this specified auxiliary encoding, not a uniqueness
principle for physical sources. On an infinite fibre the same statement
requires $T^{-1/2}J$ Hilbert--Schmidt and the corresponding closed form;
an infinite identity coevaluation vector is not assumed.

\begin{proposition}[Retained primitive coefficient comparison]
\label{p1b:prop:v224-retained-primitive}
Let $\dim_{\mathbb C}\mathbb S=4$, $\dim_{\mathbb C}F=48$,
$\Gamma_{\rm rel}=\gamma_x\otimes\Gamma_W$ with
$\operatorname{Tr}_F\Gamma_W=0$, and $P_+=(I+\Gamma_{\rm rel})/2$.
The primitive retraction $\epsilon_F=\operatorname{Tr}_F/48$ of
Theorem~\ref{p1b:thm:v130-canonical-owner-retraction} gives
\[
 Z_{\rm rp}=P_+/48,\qquad \operatorname{rank}Z_{\rm rp}=96,
 \quad\operatorname{Tr}Z_{\rm rp}=2,\qquad
 \frac1{48}\operatorname{Tr}[P_+(Q_g\otimes I_F)]
       =\frac12\operatorname{Tr}_{\mathbb S}Q_g.
\]
Relative to the ordinary pure-metric trace, this is exactly $1/96$.
\end{proposition}
\begin{proof}
Expand $P_+$. Its graded internal trace vanishes, while the identity
term contributes $48\operatorname{Tr}_{\mathbb S}Q_g/2$ before
the retraction. The relative grading has equally many signs; its
projector has rank 96 and its nonzero rescaled eigenvalue is $1/48$.
\end{proof}
If $D$ is relative odd, $D^2$ reduces $P_+$. This licenses a retained
squared-operator readout, not a rank-two Clifford-stable projection of
the first-order operator. The complete marked readout is
\[
 \mathcal I_{\rm rp}=\frac1{96}
 \left[\mathcal I_{\rm raw}
 +\operatorname{Tr}(\Gamma_{\rm rel}Le^{-D^2/\Lambda_g^2})\right].
\]
The graded term need not vanish for an interacting or boundary word.
The ordinary action's adoption of this marked readout remains an
action comparison, with all source, gauge and boundary data transported.

\begin{proposition}[Spectral and positive heat representations of the cap]
\label{p1b:prop:v224-heat-cap}
For the normalized two-dimensional Born pencil, its retained response is
$M_x(z)=z-1-x^2/(1+z)$. With $G=\sqrt{1+x^2}$,
\[
 M_x(G)=0,\qquad \frac{M_x'(G)}G=\frac2{1+G}.
\]
For $s\ge0$ the same scalar function has the positive representation
\[
 H(s)=\frac2{1+\sqrt{1+s}}=\int_0^\infty\rho(t)e^{-ts}\,dt,
 \quad \rho(t)=\frac{2e^{-t}}{\sqrt{\pi t}}
                       -2\operatorname{erfc}\sqrt t.
\]
Its moments are $\int\rho=1$, $\int t\rho=1/4$,
$\int t^2\rho=1/4$.
\end{proposition}
\begin{proof}
Differentiate $M_x$ and use $x^2=(G-1)(G+1)$.
Integration by parts gives
$\rho(t)=\pi^{-1/2}\int_t^\infty e^{-v}v^{-3/2}\,dv>0$.
Tonelli's theorem evaluates its Laplace transform as
$2(\sqrt{1+s}-1)/s$, continuously equal to one at $s=0$.
The same integral gives
$\int t^n\rho(t)dt=\Gamma(n+1/2)/[\sqrt\pi(n+1)]$ for
$n=0,1,2$, proving the moments.
\end{proof}
A single heat time with the same first moment has second moment $1/16$,
so cannot produce this function. Positive normalized finite mixtures
$\sum_iw_iH(\lambda_i s)=H(\lambda_0s)$ near zero, with
$\lambda_i\ge0$, force $\sum_iw_i\lambda_i=\lambda_0$ and
$\sum_iw_i\lambda_i^2=\lambda_0^2$; hence all visible rates coincide.
These are functional-calculus identities, not a regulator replacement.
More generally, for $K\succeq0$, $W\succeq0$, $\operatorname{Tr}W=M_0>0$,
the equality $\operatorname{Tr}[WH(sK)]=M_0H(cs)$ near zero forces
$\operatorname{Tr}(WK)=cM_0$ and $\operatorname{Tr}(WK^2)=c^2M_0$.
Consequently
$\|(K-cI)W^{1/2}\|_{\rm HS}^2=\operatorname{Tr}[W(K-cI)^2]=0$.
This is equivalent to source support in the $c$ eigenspace and also
suffices for the full response equality; no $[W,K]=0$ premise is needed.
Indeed $H(s)\sim2s^{-1/2}$, so its zeroth and first positive-power
profile integrals diverge. The four-dimensional Weyl counting law
$N_{D^2}(L)\sim cL^2$ makes $H(D^2)$ non-trace-class on the full
compact elliptic spectrum. Finite retained windows and the boundary
response role remain valid. Neither the spectral parameter $x$ nor an
analytic normal coordinate is thereby identified with physical curvature.

\begin{proposition}[Affine relative-odd normal completion]
\label{p1b:prop:v224-normal-completion}
For self-adjoint anticommuting involutions $J,K$ and $x\ge0$,
\[
 D_{\rm comp}=\ell[\sqrt{2x}\,J+(x-1)K],\qquad
 D_{\rm comp}^2=\ell^2(1+x^2)I.
\]
Within an affine channel $b(x)=a_0+a_1x$, the condition
$xb(1/x)=-b(x)$ forces $b(x)=a_1(x-1)$; unit endpoint norm fixes
$|a_1|=1$.
\end{proposition}
\begin{proof}
Anticommutation eliminates the cross term, leaving
$\ell^2[2x+(x-1)^2]I$. Comparing the constant and linear
coefficients of the relative-odd condition gives $a_0=-a_1$.
\end{proof}
This is uniqueness only in the stated completion class. Its real
self-adjoint expression has $x\ge0$ domain although its square extends.
The native round-collar scale and source comparisons remain separate
in Paper~III-F's native tensor and reference-boundary analysis.


\subsection{Unitary singlet corners and complete cyclic transfer tests}
\label{p1b:sec:v228-unitary-cyclic-tests}
All statements in this section are finite geometric operator identities.
Their application to a source consumes that source's actual constructor.

\begin{theorem}[A unit singular channel forced by two marked lines]
\label{p1b:thm:v228-unitary-singlet}
Let $W:G_R\to G_L$ be unitary between complex three-spaces, with unit
vectors $c_L,c_R$. Put $Q_a=I-c_ac_a^\dagger$, $F=Q_LWQ_R$,
$\rho=c_L^\dagger Wc_R$ and $v=Q_LWc_R$. Then
\[
 FF^\dagger=Q_L-vv^\dagger,\qquad \|v\|^2=1-|\rho|^2.
\]
If $|\rho|<1$, the plane corner has squared singular values $1,|\rho|^2$.
The line orthogonal to $c_L,Wc_R$ is unique. For its unit vector $g$,
$h=W^\dagger g$ satisfies $h\perp c_R$, $Fh=g$, $F^\dagger g=h$.
Every scalar-weighted four-corner word of the same edge therefore maps
$h\mapsto a_{00}g$ and $g\mapsto\bar a_{00}h$ under its adjoint,
where $a_{00}$ is its plane-to-plane coefficient.
\end{theorem}
\begin{proof}
Insert $Q_R=I-c_Rc_R^\dagger$ and $WW^\dagger=I$.
On $Q_LG_L$, the rank-one subtraction has eigenvalues
$1-\|v\|^2=|\rho|^2$ and $1$. The two root constraints leave one
line. Root-touching corners vanish on the displayed pair.
\end{proof}

\begin{proposition}[Maximal paired subspace avoiding the roots]
\label{p1b:prop:v228-root-orbit}
For unitary edges $W_j:G_R\to G_L$, choose $W_0$ and let $K_R$ be the
smallest subspace invariant under $W_0^\dagger W_j$ and their adjoints
that contains $c_R,W_0^\dagger c_L$. Then $D_R=K_R^\perp$,
$D_L=W_0D_R$ is the maximal paired reducing subspace annihilated by
both root projectors. Every path visiting a root vanishes there.
\end{proposition}
\begin{proof}
The invariant subspace is reducing because adjoints are included.
Writing $W_j=W_0(W_0^\dagger W_j)$ proves the edge property.
Any other such pair is orthogonal to both seeds and all their loop
images, proving maximality.
\end{proof}

For the rooted family data set $c=\mathbf1/\sqrt3$, $Q_0=cc^\dagger$,
$Q=I-Q_0$, $R=\operatorname{diag}(1,i,-i)$ and $Xe_j=e_{j+1}$.
Define the following geometric vectors and arrows:
\[
 g_e=\tfrac12(-1-i,i,1)^T,\quad
 g_\mu=\tfrac1{2\sqrt6}(2,-1+3i,-1-3i)^T,\quad
 h_e=R^\dagger g_e,\quad h_\tau=\bar g_\mu,
\]
\[
 B_e=g_eh_e^\dagger,\quad B_\mu=g_\mu c^\dagger,\quad
 B_\tau=ch_\tau^\dagger,\quad
 F_\gamma=B_e+\gamma g_\mu h_\tau^\dagger.
\]
Put $P_e=g_eg_e^\dagger$, $P_\mu=g_\mu g_\mu^\dagger$ and $P_\tau=Q_0$.
Both $(g_e,g_\mu,c)$ and $(h_e,c,h_\tau)$ are unitary frames;
direct substitution gives $F_{-1/3}=QRQ$.
Put $C_1=X+X^\dagger$, $C_2=-i(X-X^\dagger)$.
The supplied returns obey $K_7=4I+2C_1$, $K_9=I$, so
$C^*(K_7,K_9,X)=\mathbb C[X]\cong\mathbb C^3$.
With $\omega=(-1+i\sqrt3)/2$, its plane character projectors are
\[
 E_+=\tfrac13(I+\bar\omega X+\omega X^2),\qquad
 E_-=\tfrac13(I+\omega X+\bar\omega X^2).
\]
The two seed vectors $c,r=R^\dagger c$, together with $Xr$, have
Gram determinant $4/27$. Thus an independent cycle makes their
root-generated orbit full; it removes universal root-avoiding protection
for the enlarged path algebra, without deciding a particular word.

\begin{theorem}[Complete cyclic functional-calculus defect]
\label{p1b:thm:v228-character-defect}
Let $C_e,C_\tau\in\mathbb C[X]$ have plane character values
$a_\pm,d_\pm$. For real $e,m,t>0$ put
$T=eF_\gamma C_e+mB_\mu+tB_\tau C_\tau$ and
\[
 w_\pm=\frac{2\mp\sqrt3}{4},\quad
 n_\pm=|a_\pm|^2,\quad q_\pm=d_\pm\bar a_\pm,\quad
 S=w_+n_++w_-n_-.
\]
Then the exact positive marked defect is
\[
 \boxed{\|(TT^\dagger-e^2I)g_e\|^2
 =e^4(S-1)^2+\frac{e^4|\gamma|^2}{16}(n_+-n_-)^2
       +\frac{e^2t^2}{16}|q_+-q_-|^2.}
\]
For $\gamma\ne0$ it vanishes exactly when
$|a_+|=|a_-|=1$ and $q_+=q_-$. For $\gamma=0$ the exact
conditions are $S=1$ and $q_+=q_-$. Root character values and
the muon coefficient do not enter this marked test.
\end{theorem}
\begin{proof}
Direct evaluation gives
$h_e^\dagger E_\pm h_e=w_\pm$ and
$h_\tau^\dagger E_\pm h_e=\pm\kappa$,
where $\kappa=-\sqrt2(1+i)/8$ and $|\kappa|^2=1/16$.
The other two row adjoints annihilate $g_e$, whereas
$F_\gamma^\dagger g_e=h_e$. Consequently
\[
 (TT^\dagger-e^2I)g_e=e^2(S-1)g_e
  +e^2\gamma\kappa(n_+-n_-)g_\mu+et\kappa(q_+-q_-)c.
\]
The three left vectors are orthonormal. Squaring gives the formula;
its nonnegative summands give both equivalences.
\end{proof}

\begin{corollary}[Noncommuting finite outer coefficients]
\label{p1b:cor:v228-operator-character-defect}
On a finite common space $K$, replace the character values by matrices
$A_\pm,D_\pm$, and take
$T=e(I\otimes F_\gamma)C_e+mM\otimes B_\mu+
t(I\otimes B_\tau)C_\tau$.
For $J_Lv=v\otimes g_e$, put
$N_\pm=A_\pm A_\pm^\dagger$, $S=w_+N_++w_-N_-$,
$\Delta=N_+-N_-$ and $L=D_+A_+^\dagger-D_-A_-^\dagger$.
No commutation of these outer matrices is assumed. Then
\[
 J_L^\dagger(TT^\dagger-e^2I)^2J_L
 =e^4(S-I)^2+\frac{e^4|\gamma|^2}{16}\Delta^2
                 +\frac{e^2t^2}{16}L^\dagger L.
\]
For $\gamma\ne0$, vanishing is equivalent to $N_+=N_-=I$ and $L=0$.
For $\gamma=0$, it is equivalent to $S=I,L=0$.
\end{corollary}
\begin{proof}
Expand in the orthogonal character projectors. Cross-character products
vanish before any outer coefficients are multiplied. The proof of the
theorem therefore gives the same three orthogonal output components,
with the displayed matrix coefficients and order retained.
\end{proof}


\begin{proposition}[Cyclic equivariance does not identify two character coefficients]
\label{p1b:prop:v240-c3-character-no-go}
Let $E_\pm$ be the two nontrivial $C_3$-character projections of
Corollary~\ref{p1b:cor:v228-operator-character-defect}. An exterior
operator $D_+E_++D_-E_-$ commutes with the cyclic generator for
arbitrary $D_+,D_-$ acting in the common exterior coefficient space.
Consequently $C_3$-equivariance alone does not imply $D_+=D_-$. In
particular the allowed algebraic coefficients
\[
 A_+=A_-=I,\quad D_+=2I,\quad D_-=I,\quad \gamma=-1/3
\]
give the strictly nonzero marked defect
\[
 \boxed{J_L^\dagger(TT^\dagger-e^2I)^2J_L
           =\frac{e^2t^2}{16}I.}
\]
If $\gamma\ne0$ and $A_+,A_-$ are positive Hermitian, exact
protection within this operator-character class is equivalent to
\[
 \boxed{A_+=A_-=I,\qquad D_+=D_-.}
\]
These are tests of the \emph{actual} mixed source word, not
consequences of a path being cyclically equivariant.
\end{proposition}
\begin{proof}
The cyclic generator acts by scalar characters on each $E_\pm$,
so it commutes with independently valued outer coefficients.
For the displayed example $N_+=N_-=I$ and $L=I$, whence the
nonnegative defect formula gives $e^2t^2I/16$.
For positive Hermitian $A_\pm$, $A_\pm A_\pm^\dagger=I$ implies
$A_\pm^2=I$ and thus $A_\pm=I$. The remaining defect equation is
then exactly $D_+=D_-$. The converse is immediate.
\end{proof}

\begin{proposition}[A constant-return positive fibre]
\label{p1b:prop:v228-opposite-odd-fibre}
Let real $A,B$ satisfy $A^2-3AB+3B^2=1$, and put
$C_e=Q_0+AQ+BC_2$, $C_\tau=Q_0+AQ-BC_2$.
The arrows $\iota_e=B_eC_e$, $\iota_\mu=B_\mu$,
$\iota_\tau=B_\tau C_\tau$ obey
$\iota_r\iota_s^\dagger=\delta_{rs}P_r$.
Thus $T=\sum_r\varepsilon_r\iota_r$ has
$TT^\dagger=\sum_r\varepsilon_r^2P_r$ for every point of the fibre.
Both dressings are positive if $A>\sqrt3|B|$.
\end{proposition}
\begin{proof}
The electron row norm is $A^2+3B^2-3AB$; exchanging the character
weights and the odd sign gives the same tau norm. Their plane product
is $(A^2-3B^2)Q$, so the cross return vanishes by $h_e\perp h_\tau$.
The root and plane supports kill the remaining cross returns.
The eigenvalues $1,A\pm\sqrt3B$ give positivity.
For example $A=\sqrt3,B=1/\sqrt3$ gives two positive, nonunitary
dressings with plane product $2I$ and the same complete left return.
\end{proof}
These results construct finite operators and common readouts on their
specified domains. They do not place an arbitrary such operator in a
physical action. A complete co-realized transport cell is a tuple:
its unweighted coherent cyclic sum is $I+X+X^{-1}=3Q_0$, which
annihilates the plane. Co-realization alone does not select that sum.


\section{Native operator interfaces, principal descent, and readout controls}
\label{p1b:sec:organized-interfaces}

\subsection{Compression commutators and omitted-image terms}
\label{p1b:app:compression-commutator}
This appendix is a mathematical reduction diagnostic. It assumes a supplied orthogonal projector and compares compression with restriction on the stated operator algebra.
For $P^2=P=P^\dagger$, $Q=I-P$, and finite operators $A,B$, the retained
range obeys
\[
 [PAP,PBP]=P[A,B]P+PBQAP-PAQBP.
\]
Indeed, insert $I=P+Q$ into $PABP-PBAP$ and rearrange. Even commuting
parent operators can therefore acquire a compression commutator when
their omitted images contribute. If the retained range reduces both
operators, both image terms vanish and compression is a true restriction.
Thus a compression is an algebra restriction precisely under the stated reducing condition.


\subsection{Pair faithfulness and normal-current retention}
\label{p1b:app:v2167-pair-current-retention}
This appendix tests a generic field-slot implication. Its mode chart is
comparison data, not a licensed alternative K16 realization. The outcome is
that faithful pair transport on one full-rank canonical chart need not retain
the normal-current algebra on that same chart.

\begin{lemma}[Multiplication defect of an isometric compression]
\label{p1b:lem:v2167-compression-multiplication}
Let $Z$ be an isometry, $P=ZZ^\dagger$, and
$\pi_Z(A)=Z^\dagger AZ$. Then
\[
 \pi_Z(A)\pi_Z(B)-\pi_Z(AB)
 =-Z^\dagger A(I-P)BZ.
\]
For a unital star algebra, this compression preserves multiplication
exactly if and only if $P$ reduces the algebra.
\end{lemma}
\begin{proof}
Insert $ZZ^\dagger=P$ between the two compressed factors. If the defect
vanishes, take $B=A^\dagger$: the positive matrix
$Z^\dagger A(I-P)A^\dagger Z$ vanishes, so
$(I-P)A^\dagger Z=0$. Applying this to every $A$ and its adjoint proves
reduction. The converse follows from $[P,A]=0$.
\end{proof}

\begin{proposition}[One canonical chart with a faithful pair and a lossy current]
\label{p1b:prop:v2167-faithful-pair-lossy-current}
Let $\mathcal M=\mathbb C^3$ be a mode multiplicity distinct from the family
carrier $G=\mathbb C^3$. Put
\[
 Q_m=\tfrac13\mathbf1\mathbf1^{\mathsf T},\qquad
 U=I+(i-1)Q_m,\qquad Je_i=(Ue_i)\otimes e_i.
\]
With parent Hilbert metric $I$, $J$ is an isometric, full-rank, constant
field chart with zero embedding connection. For every family normal coefficient,
\[
 \pi_J(A):=J^\dagger(I_{\mathcal M}\otimes A)J
 =\operatorname{diag}(A_{00},A_{11},A_{22}).
\]
For every symmetric family pair $M'$, the same creation-slot pullback,
extended by the identity on a two-dimensional antisymmetric spin slot, is
\[
 \Gamma_J(M')=[(I-2Q_m)\circ M']\otimes\epsilon,
 \qquad \epsilon=\begin{pmatrix}0&1\\-1&0\end{pmatrix}.
\]
This pair map is injective on the entire symmetric family coefficient
space, while $\pi_J(E_{ij})=0$ for every $i\ne j$.
\end{proposition}
\begin{proof}
$U$ acts by $i$ on $Q_m\mathcal M$ and by $1$ on its complement. Hence
$U^\dagger U=I$, and its columns $\phi_i=Ue_i$ are orthonormal.
The normal matrix entries are
$\langle\phi_i,\phi_j\rangle A_{ij}=\delta_{ij}A_{ij}$.
The creation slots instead contract by
$\phi_i^\dagger\overline{\phi_j}$; their matrix is
$U^\dagger\bar U=I-2Q_m$. Its diagonal entries are $1/3$ and its
off-diagonal entries $-2/3$, all nonzero. Entrywise division therefore
recovers every $M'$. Symmetry of this matrix preserves pair antisymmetry
after multiplication by $\epsilon$.
\end{proof}
For the unchanged Born coefficient
$M=\frac{-1+i}{2}\operatorname{diag}(1,-i,i)$, the comparison returns
$(M/3)\otimes\epsilon\ne0$. The factor $1/3$ is an evaluated comparison
mode contraction, not a fitted native efficiency. Although all three
diagonal role tests survive, their off-diagonal products do not:
$\pi_J(E_{21})\pi_J(E_{12})=0$ whereas $\pi_J(E_{22})=E_{22}$.

\begin{proposition}[Complete regular reduction does not supply intertwining]
\label{p1b:prop:v2167-complete-return-diagonal}
On this fixed comparison multiplicity let
$h_p=I_{\mathcal M}\otimes h_f$, with $h_f$ a complete Hermitian
port/family normal member, and retain
$J_Re_i=\phi_i\otimes|R,i\rangle$. On every common regular domain,
\[
 [J_R^\dagger(\zeta I-h_p)^{-1}J_R]_{ij}
 =\delta_{ij}[(\zeta I-h_f)^{-1}]_{Ri,Rj}.
\]
The inverse complete Schur kernel has this same diagonal return.
Source variations of $h_f$ at the fixed chart retain that property.
\end{proposition}
\begin{proof}
The complete resolvent factors by $I_{\mathcal M}$, giving the displayed
orthogonality factor. The Schur inverse equals the full retained resolvent;
successive complete eliminations agree on their common regular branch.
Differentiation at this fixed comparison chart preserves the mode factor.
\end{proof}
Thus rank, faithful pair insertion, canonical normalization, and complete
Schur identities alone do not imply normal-current intertwining. This
refutes a shortcut to a retention certificate; it is not a native
non-derivability theorem. The required mixed return can be proved structurally or evaluated from a specified embedding.


\setlength{\emergencystretch}{2em}
\subsection{Equivariant embeddings and matrix-unit retention}
\label{p1b:sec:v2201-geometry-first-consequences}
The reciprocal $M_8$ rooted family algebra supplies the specified inherited
field embedding. Canonical compression and current retention are compared
under their stated domain, metric and intertwining assumptions.

\subsubsection{Canonical chart normalization and family-current retention}
\label{p1b:sec:v2201-derived-d08}

For a family-equivariant field embedding $J=J_m\otimes I_G$ and its
canonical isometry $Z$, the compression preserves every family matrix
unit. Consequently for any two unitary frame matrices $W_1,W_2$, their comparison matrix remains
$W_1^\dagger W_2$ after canonical normalization.
\begin{proof}
The inherited Gram and its positive square root factor by
$I_G$. The normalized isometry therefore factors as well. Compression
of $I\otimes A$ returns $I\otimes A$ for every $A$.
\end{proof}
The probability functional can change with its density matrix, positive
effect, unitary family or nonreducing mixed return. It is
\[
 P_D(t)=\operatorname{Tr}\{E_D\,V(t)\rho_PV(t)^\dagger\}.
\]
For diagonal role propagation its phase coefficients are
$(\rho_P)_{rs}(E_D)_{sr}$. These are explicit density/effect coefficients. Faithful pair transport alone
does not establish the compressed-operator identity.
\par\noindent\textit{Comparison:} \ref{p1b:prop:v2167-faithful-pair-lossy-current}. The factorization argument above is independent of any later-series application.


\subsection{Connection, reduction and coordinate provider identities}
\label{p1b:sec:v215-geometric-provider-completion}
\subsubsection{Projection, flat dilation and the exact connection-to-operator maps}
\label{p1b:sec:providers-connections}

The connection families above have different underlying bundles.  The
following results give explicit maps between some of them, and identify
the information lost when a connection is contracted to a spin Dirac
operator.  In particular, equality of two connection matrices in one
representation is distinguished from equality of differential operators
on different tensor factors.

\begin{proposition}[Projected connections and their curvature]
\label{p1b:prop:providers-projection}
Let \((\mathcal V,D)\) be a Hermitian bundle with unitary connection over
a manifold \(N\), and let \(P=P^*=P^2\) be a smooth constant-rank
orthogonal projector.  The image \(E=\operatorname{im}P\) has the
unitary connection \(\nabla^E=P D|_E\), with curvature
\[
 \boxed{R^E=P R^D P+P(DP\wedge DP)P.}
\]
If \(\mathcal V=N\times\mathbb C^k\) and \(D=d+A\), and if
\(V:N\times\mathbb C^r\to\mathcal V\) is an orthonormal frame of
\(E\), its connection matrix and curvature are
\[
 A_E=V^*AV+V^*dV,\qquad
 F_E=dA_E+A_E\wedge A_E.
\]
Under \(V\mapsto Vg\), \(g:N\to U(r)\), this matrix changes to
\(g^{-1}A_Eg+g^{-1}dg\).  These constructions commute with pullback.
Flatness of \(D\) therefore does not imply flatness of \(\nabla^E\).
\end{proposition}
\begin{proof}
For a section \(s=Ps\), apply \(PD\) twice and use
\(P(DP)P=0\), obtained by differentiating \(P^2=P\).  The result is
\(P R^D s+P(DP\wedge DP)s\), proving the first formula.  In an
orthonormal frame a section is \(V\xi\), so
\(V^*D(V\xi)=d\xi+(V^*AV+V^*dV)\xi\).  This proves the local
formula, metric compatibility and the stated change of frame.  Pullback
commutes with multiplication, covariant differentiation and exterior
products, proving naturality.  The curvature formula exhibits the
additional term even when \(R^D=0\).
\end{proof}

\begin{theorem}[Complementary projector realization of a Cartan profile]
\label{p1b:thm:providers-profile-projectors}
Let \(U:N\to SU(2)\), put \(\Theta=U^*dU\), and let
\(f:N\to(0,1)\) be smooth.  In the trivial Hermitian rank-four bundle
define
\[
 V_f^+=\begin{pmatrix}\sqrt f\,U\\\sqrt{1-f}\,I\end{pmatrix},
 \qquad
 V_f^-=\begin{pmatrix}-\sqrt{1-f}\,U\\\sqrt f\,I\end{pmatrix},
 \qquad P_f^\pm=V_f^\pm(V_f^\pm)^*.
\]
The two projectors are orthogonal and complementary.  Their projected
flat connections, expressed in these frames, are exactly
\[
 \boxed{A_f^+=f\Theta,\qquad A_f^-=(1-f)\Theta.}
\]
In particular no derivatives of \(f\) are omitted from this identity.
Their full curvatures are
\[
 F_f^+=df\wedge\Theta+f(f-1)\Theta\wedge\Theta,\qquad
 F_f^-=-df\wedge\Theta+f(f-1)\Theta\wedge\Theta.
\]
The frames and projectors extend across endpoint values whenever the
displayed square roots extend smoothly; a smooth angular
parametrization \(f=\cos^2\vartheta\), with the corresponding signed
amplitudes \(\cos\vartheta,\sin\vartheta\), supplies such an extension.

At the constant midpoint the two projected connections are unitarily
isomorphic, by the map \(V_{1/2}^-(V_{1/2}^+)^*\) between their images.
For \(N=S^3\) and \(U\) its marked quaternionic parametrization, both
have restricted holonomy \(SU(2)\), although the ambient connection
is flat.
\end{theorem}
\begin{proof}
Matrix multiplication gives \((V_f^\pm)^*V_f^\pm=I\) and
\((V_f^+)^*V_f^-=0\); the ranks sum to four.  In each differentiated
frame the two real-amplitude terms sum to
\(\tfrac12df+\tfrac12d(1-f)=0\).  The remaining term is
\(fU^*dU\) or \((1-f)U^*dU\), respectively.  Maurer--Cartan,
\(d\Theta=-\Theta\wedge\Theta\), then proves both curvature
formulas.  The endpoint assertion follows directly from the frames;
smoothness of \(f\) alone is not being used to infer smoothness of an
arbitrary square root.  At the midpoint the two frame matrices of the
connections coincide, so their frame-induced map intertwines covariant
derivatives.  Finally the values of
\(-\Theta\wedge\Theta/4\) on a three-sphere span
\(\mathfrak{su}(2)\).  The restricted holonomy algebra contains these
curvature values and lies in \(\mathfrak{su}(2)\); its connected group
is therefore \(SU(2)\).
\end{proof}

\begin{corollary}[Clifford latitude and its four-sphere extension]
\label{p1b:cor:providers-projector-four-sphere}
In a chiral quaternionic frame write
\[
 \Gamma(n)=\begin{pmatrix}0&U(n)\\U(n)^*&0\end{pmatrix},\quad
 \Gamma_* =\begin{pmatrix}I&0\\0&-I\end{pmatrix},\quad u=2f-1.
\]
Then
\[
 2P_f^+-I=\sqrt{1-u^2}\,\Gamma(n)+u\Gamma_*.
\]
For \(q=rU(n)\in\mathbb H\), \(\rho>0\), and
\(f=r^2/(r^2+\rho^2)\), the projector is
\[
 P_\rho^+(q)=\frac1{r^2+\rho^2}
 \begin{pmatrix}r^2I&\rho q\\\rho q^*&\rho^2I\end{pmatrix}.
\]
It extends smoothly to the one-point compactification \(S^4\).
Its image and complementary image are rank-two bundles with trivial
determinants and opposite second Chern numbers of absolute value one.
On \(\mathbb H\setminus\{0\}\), the first connection is precisely
\(f\Theta\).  This is an identity of bundles and connections, using
the same profile \(f\).
\end{corollary}
\begin{proof}
Multiplication of the rows of \(V_f^+\) gives the latitude formula.
The explicit rational projector is smooth at \(q=0\).  In the chart
\(q'=q/|q|^2\) around infinity its entries are
\[
 \frac1{1+\rho^2|q'|^2}
 \begin{pmatrix}I&\rho q'\\\rho q'^*&\rho^2|q'|^2I\end{pmatrix},
\]
so it is smooth there as well.  A smooth frame around zero is
\((q,\rho I)^t/(r^2+\rho^2)^{1/2}\).  Around infinity use that frame
multiplied on the right by \(U^*\); its two blocks become
\((I,\rho q'^*)^t/(1+\rho^2|q'|^2)^{1/2}\).  The transition is
\(U^*:S^3\to SU(2)\), whose degree has absolute value one, so the
clutching second Chern number has absolute value one and the determinant
is trivial.  Its sign is fixed by the chosen chart and orientation
conventions.  The complementary bundle has trivial direct sum with it,
so its first Chern class is zero and its second Chern class is the
negative.  On the punctured chart the zero frame is exactly \(V_f^+\),
and the preceding theorem gives its connection.
\end{proof}

\begin{theorem}[A globally flat dilation of the sphere Levi--Civita connections]
\label{p1b:thm:providers-flat-dilation}
On \(M_8=S^3(R_x)\times S^3(R_p)\times\mathbb T^2\) give the
central plane its declared flat metric and choose an orthonormal
constant central frame \(W_x,W_p\).  The bi-invariant metric connection
with parameters
\[
 s_x=s_p=\tfrac12,\qquad
 \ell_i(z)=g(W_i,z)/R_i,\qquad\kappa=0
\]
is globally flat and has trivial tangent holonomy.  Its orthogonal
projection to \(TS_i^3\) is exactly \(\nabla_i^{\rm LC}\).
The spin connection of that projected rank-three connection is
\(\Theta_i/2\) in the convention below.

When the product metric is exchange invariant, the frame can be chosen
so that this dilation is exchange invariant: for unit central exchange
eigenvectors \(E_+,E_-\), take
\(W_x=(E_++E_-)/\sqrt2\) and
\(W_p=(E_+-E_-)/\sqrt2\).  All these frame choices give the same
projected sphere connections.

More generally, among the flat members of the entire eight-parameter
bi-invariant connection family, a projected sphere connection has only
the values \(s_i=0,\tfrac12,1\).  Its unique curved possibility is
\(s_i=\tfrac12\), realized exactly when
\(\ell_i\ne0\); then \(|\ell_i^\sharp|=1/R_i\).
\end{theorem}
\begin{proof}
On \(TS_i^3\oplus\mathbb RW_i\) the stated connection reads
\[
 D_X(Y+aW_i)=\nabla_X^{\rm LC}Y+\frac a{R_i}X+
       \left(Xa-\frac{g(X,Y)}{R_i}\right)W_i.
\]
Identify this rank-four bundle isometrically with the trivial
\(\mathbb R^4\) bundle restricted to the sphere by
\(Y+aW_i\mapsto Y+a n_i\), where \(n_i=q_i/R_i\) is its Euclidean
unit normal.  The identities
\(\partial_Xn_i=X/R_i\) and
\((\partial_XY)^\perp=-g(X,Y)n_i/R_i\) show that \(D\) becomes
ordinary Euclidean differentiation.  Taking the two independent
rank-four blocks gives a global isometric trivialization over all of
\(M_8\); it is constant along the scalar torus.  The flat tangent
holonomy is therefore trivial, including around the scalar cycles.
Projection removes the normal terms and leaves exactly the displayed
Levi--Civita connection.  This projection occurs before taking the
rank-three spin lift; it is not a claim about contracting the full
rank-eight flat connection.

Exchange invariance of the metric gives \(R_x=R_p\) and makes the
two central exchange eigenspaces orthogonal.  The proposed frame obeys
\(SW_x=W_p\), hence \(\ell_p=\ell_x\circ S\); the other exchange
conditions are immediate.  For the last assertion the flatness equations
of Theorem~\ref{p1b:thm:v214-flat-bi-connections}, with
\(c_i=2/R_i\) and \(v_i=\ell_i^\sharp\), are
\[
 (2s_i-1)v_i=0,\qquad
 s_i(s_i-1)+|v_i|^2/c_i^2=0.
\]
If \(v_i=0\), they force \(s_i=0\) or \(1\); otherwise they force
\(s_i=1/2\) and \(|v_i|=c_i/2=1/R_i\).  The sphere projection is
\(s_i[X,Y]\), whose curvature is \(s_i(s_i-1)\operatorname{ad}_{[X,Y]}\).
This proves exactly the three possibilities and identifies the sole
curved one.  The rank-three spin formula follows from the next
proposition.  For a fixed spin structure on \(M_8\), lifting this
globally flat tangent connection can still retain its two central spin
transition signs; trivial tangent holonomy does not remove those signs.
\end{proof}

\begin{theorem}[The exact spin-Dirac contraction of the full metric-connection family]
\label{p1b:thm:providers-dirac-quotient}
Fix any left-invariant metric on \(M_8\), any of its spin structures,
and its full complex spinor bundle.  Write a metric connection as
\(\nabla=\nabla^{\rm LC}+K\), with
\(K_{abc}=g(K_{e_a}e_b,e_c)=-K_{acb}\), and write its torsion as
\(T_{abc}=K_{abc}-K_{bac}\).  Put
\[
 T^{\rm alt}_{abc}=\tfrac13(T_{abc}+T_{bca}+T_{cab}),\qquad
 t_a=\sum_bT_{abb}.
\]
With \(c(e_a)c(e_b)+c(e_b)c(e_a)=-2\delta_{ab}\), the induced
spin Dirac operator is exactly
\[
 \boxed{D^\nabla=D^{\rm LC}+\tfrac34c(T^{\rm alt})-
                                      \tfrac12c(t).}
\]
Consequently the affine map from the \(224\)-dimensional invariant
metric-connection space to contracted spin Dirac operators has
\(160\)-dimensional kernel, precisely the traceless Cartan torsion
part.  Its image has real dimension \(56+8=64\).  It is formally
self-adjoint in the metric volume exactly when \(t=0\); that subfamily
has dimension \(216\), and its operator image has dimension \(56\).
For an exchange-invariant metric and exchange-invariant connections,
the corresponding dimensions are kernel \(80\), image \(28+4=32\),
and formally self-adjoint connection subfamily \(108\).

The symmetric part of the spin contraction of an arbitrary metric
connection is \(D^{\rm LC}+3c(T^{\rm alt})/4\).  Thus even this
symmetrization forgets both trace and Cartan torsion; it does not retain
the full connection.
\end{theorem}
\begin{proof}
The lifted connection difference on spinors is
\(\tfrac14\sum_{b,c}K_{abc}c(e_b)c(e_c)\).  Contracting once more gives
\(\tfrac14\sum_{a,b,c}K_{abc}c(e_a)c(e_b)c(e_c)\).  Use
\[
 c_ac_bc_c=c_{[a}c_bc_{c]}-\delta_{ab}c_c+
                                      \delta_{ac}c_b-\delta_{bc}c_a.
\]
Full alternation of \(T=K-K^{(12)}\) is twice full alternation of
\(K\), and \(c(T^{\rm alt})=\tfrac16\sum T^{\rm alt}_{abc}c_ac_bc_c\).
The fully alternating contribution is therefore
\(3c(T^{\rm alt})/4\).  The repeated-index terms are
\(-\tfrac12\sum_{a,c}K_{aac}c_c\).  Since
\(t_c=\sum_aT_{caa}=\sum_aK_{aac}\), they give exactly
\(-c(t)/2\), proving the operator formula.

In real dimension eight the full complex Clifford representation is
faithful: the Clifford algebra is the matrix algebra on the full spin
module.  Its distinct exterior degrees, in particular one and three,
are linearly independent.  Thus the contraction vanishes precisely
when \(t=0=T^{\rm alt}\).  The torsion decomposition of
Proposition~\ref{p1b:prop:g19-torsion} identifies this kernel with the
\(160\)-dimensional Cartan summand and gives all image dimensions.
Real one-form Clifford multiplication is skew-adjoint, whereas real
three-form Clifford multiplication is self-adjoint.  The Levi--Civita
Dirac operator is formally self-adjoint; the displayed expression is
therefore formally self-adjoint exactly when \(c(t)=0\), equivalently
\(t=0\).  Removing this eight-dimensional trace summand leaves
\(216\) dimensions.  Restricting the same decomposition to exchange
fixed vectors uses \(112=4+28+80\), giving the stated restricted
counts.  Taking the symmetric part removes precisely the skew-adjoint
one-form term.
\end{proof}

\begin{corollary}[Domains and the untwisted product index]
\label{p1b:cor:providers-domains-index}
For smooth metric connections with \(t=0\), the operators in the
preceding theorem are self-adjoint on the same Sobolev domain
\(H^1(M_8;\Sigma)\) and have compact resolvent.  They are odd for the
fixed eight-dimensional chirality.  Their chiral index is zero, so their
positive and negative chiral kernel dimensions are equal.  A continuous
path of these operators with invertible endpoints has zero full
self-adjoint spectral flow.
\end{corollary}
\begin{proof}
The connection change is a smooth self-adjoint bundle endomorphism and
therefore bounded on \(L^2\), since \(M_8\) is compact.  Bounded
self-adjoint perturbation preserves self-adjointness and its \(H^1\)
domain.  Ellipticity and compactness of the Sobolev inclusion give compact
resolvent.  Clifford multiplication by a three-form is odd, so chirality
is unchanged.  The principal symbol is that of the untwisted spin
operator; hence the chiral index is its spin index.  Parallelizability
of \(TM_8\) makes the positive-degree Pontryagin classes zero, and the
degree-eight term of \(\widehat A(TM_8)\) is consequently zero.  The
index is zero.  Formal adjointness identifies the two chiral kernels
with the index kernel and cokernel.  Finally fixed chirality conjugates
every operator in the path to its negative.  Unitary invariance and
sign reversal of spectral flow give \(\operatorname{sf}(D)=-
\operatorname{sf}(D)\), proving the last assertion.
\end{proof}

\begin{proposition}[Sphere spin connection and its two inequivalent contractions]
\label{p1b:prop:providers-sphere-dictionary}
On \(S_R^3\) choose an oriented orthonormal left frame with
\([e_a,e_b]=2\epsilon_{abc}e_c/R\), and set
\(c(e_a)=i\sigma_a\), \(\Theta(e_a)=-i\sigma_a/R\).
The tangent connection \(\nabla^s_XY=s[X,Y]\) has spin connection
matrix \(s\Theta\).  Its spin Dirac contraction \(T_s\), on the
rank-two spinor bundle \(\Sigma\), is
\[
 T_s=D^{\rm LC}+3(s-1/2)I_\Sigma/R.
\]
For constant \(s\), its complete values are
\[
 \frac{\ell+3s}{R},\qquad\frac{-\ell-3+3s}{R},\qquad
 \ell\ge0,
\]
each with labelled multiplicity \((\ell+1)(\ell+2)\).

Instead keep the first spin connection Levi--Civita and let
\(A_f=f\Theta\) act on a separate fundamental twisting doublet
\(E\).  The resulting rank-four twisted spin Dirac operator obeys
\[
 \boxed{\mathcal D_f=B_R+\frac{f-1/2}{R}
                    \operatorname{diag}(-3,1,1,1),\qquad
 B_R=\begin{pmatrix}0&d^*\\d&-*d\end{pmatrix}.}
\]
Here the unitary identification is
\(\Sigma\otimes E\simeq\operatorname{End}(\Sigma)
\simeq(\Lambda^0\oplus\Lambda^1)_\mathbb C\), with
\((\phi,\alpha)\mapsto(\phi I+c(\alpha))/\sqrt2\).
It gives \(B_R^2=\Delta_0\oplus\Delta_1\) and
\(\ker B_R=\mathbb C(1,0)\).  The formulas for the operators remain
valid for smooth profiles \(s,f\); the spectral formula assumes
constant \(s\).  Reversing the volume-Clifford convention reverses
the signed sphere-Dirac convention, not the kernel multiplicities.
\end{proposition}
\begin{proof}
The spin lift is
\[
 \Omega_s(e_a)=\tfrac14\sum_{b,c}(2s/R)\epsilon_{abc}c_bc_c
             =-s c_a/R=s\Theta(e_a),
\]
using \(c_bc_c=-\delta_{bc}I-\epsilon_{bcd}c_d\).
Contraction of \((s-1/2)\Theta\) gives
\(3(s-1/2)I/R\).  Adding this scalar to the complete
Levi--Civita spectrum \(\pm(\ell+3/2)/R\) proves the two labelled
branches and their multiplicities.  This convention has
\(c(e_1)c(e_2)c(e_3)=I\); it is the opposite volume convention
to the signed choice in Theorem~\ref{p1b:thm:v214-torsion-dirac}.

The invariant alternating form on a fundamental \(SU(2)\) doublet
identifies it with its dual, so the midpoint tensor product identifies
with endomorphisms carrying the induced Levi--Civita connection.
The trace and traceless endomorphisms are exactly the scalar and vector
summands in the stated unitary map.  Contracting the induced connection
and using the same Clifford product gives scalar part \(d^*\alpha\)
and vector part \(d\phi-*d\alpha\), proving the displayed \(B_R\).
The identities \(d^2=0\), \((d^*)^2=0\) and the Hodge identities
give its square.  On a three-sphere the only harmonic zero- or one-forms
are the constants, proving the kernel.  The change on the twisting
factor is
\[
 \sum_a c_a\otimes(f-1/2)\Theta(e_a)
       =(f-1/2)\sum_a\sigma_a\otimes\sigma_a/R.
\]
The last sum has eigenvalue \(-3\) on the singlet and \(+1\) on
the triplet: it is twice the interchange operator minus the identity.
This proves the twisted formula, including nonconstant multipliers.
No derivative of a connection profile is introduced by a first-order
Dirac contraction.
\end{proof}

\begin{theorem}[Exact obstruction to replacing a twisting profile by tangent torsion]
\label{p1b:thm:providers-torsion-obstruction}
Let \(s,f:S_R^3\to\mathbb R\) be smooth and compare equal-rank
operators \(T_s\otimes I_2\) and \(\mathcal D_f\).  A smooth unitary
bundle map covering the identity and commuting with the fixed Clifford
action intertwines them if and only if
\[
 \boxed{s\equiv1/2,\qquad f\equiv0\ \text{or}\ f\equiv1.}
\]
In particular the rank-four midpoint \(B_R\) is not such a doubled
torsion-spin operator.  Even allowing an abstract unitary of Hilbert
spaces, \(B_R\) is not equivalent to any constant-\(s\)
\(T_s\otimes I_2\): its kernel has dimension one, whereas the latter
has kernel dimension zero or a positive multiple of four.

If \(Q_s\) instead contracts the connection induced by the same
\(\nabla^s\) on both factors of \(\Sigma\otimes\Sigma^*\), then
\[
 Q_s=B_R+4(s-1/2)\Pi_{\Lambda^1}/R,\qquad Q_s(1,0)=0.
\]
In the displayed fixed identification \(Q_s=\mathcal D_f\) holds
only for \(s=f=1/2\).
\end{theorem}
\begin{proof}
Irreducibility of the complex three-dimensional Clifford doublet
implies that every Clifford-commuting fibre unitary is \(I\otimes g\)
for a smooth \(g:S^3\to U(2)\).  It acts only by a change of gauge
on the twisting connection.  Subtracting the common untwisted
Levi--Civita spin operator, the proposed equality is
\[
 \sum_a c_a\otimes A_f^g(e_a)=3(s-1/2)I\otimes I/R.
\]
Partial trace over the spin doublet is zero on the left and forces
\(s=1/2\).  Linear independence of the three Clifford matrices then
forces \(A_f^g=0\).  Thus its curvature vanishes.  In
\(F_f=df\wedge\Theta+f(f-1)\Theta\wedge\Theta\), the coefficient
of the third independent Lie direction in \(F_f(e_a,e_b)\),
\(a\ne b\), is a nonzero multiple of \(f(f-1)\); the derivative
terms have only the first two directions.  Therefore \(f(f-1)=0\)
pointwise.  Connectedness makes \(f\) a constant endpoint.
Conversely \(A_0=0\) and \(A_1=U^*dU\) is globally pure gauge,
proving both sufficiency statements.  A zero of a constant torsion
branch has multiplicity \((\ell+1)(\ell+2)\), an even integer;
the two zero arithmetic families are disjoint.  Doubling therefore
gives a multiple of four, whereas \(\ker B_R\) is one-dimensional.

Changing the first spin factor contributes \(3(s-1/2)I/R\).
Changing the second dual factor contributes
\((s-1/2)\operatorname{diag}(-3,1,1,1)/R\), by the same alternating
identification used above.  Their sum gives the formula for \(Q_s\).
Equivalently the identity endomorphism is parallel for every induced
endomorphism connection \(d+[s\Theta,\,\cdot\,]\).  Equality of the
scalar entries of \(Q_s-B_R\) and \(\mathcal D_f-B_R\) first forces
\(f=1/2\); equality of their vector entries then forces \(s=1/2\).
\end{proof}

\begin{corollary}[Flat Hermitian connection versus the curved midpoint]
\label{p1b:cor:providers-bismut-midpoint}
On the compatible bi-invariant Hermitian branch, the left-frame Bismut
connection of Theorem~\ref{p1b:thm:g14-bismut} is the endpoint
\(s_x=s_p=0\), with three-form torsion \(-H_g\) and trivial tangent
holonomy.  Its projected sphere spin connections are zero in their
global left frames.  They are not gauge-equivalent to the midpoint
\(\Theta_i/2\), whose curvature is
\(-\Theta_i\wedge\Theta_i/4\) and whose restricted holonomy is
\(SU(2)\).  The other flat endpoint \(s_i=1\) is pure gauge on a
sphere.  The closed Cartan torsion also cannot be changed into an exact
three-form by a global two-form shift.
\end{corollary}
\begin{proof}
The Bismut assertion and its torsion are precisely the global
parallel-frame connection of the cited theorem.  The sphere spin
dictionary gives its matrix and the midpoint matrix.  Curvature is
conjugated by gauge transformations; zero curvature cannot be conjugate
to the nonzero midpoint curvature.  Their restricted holonomies give
the same obstruction.  Finally
\(\int_{S_i^3}H_g=4\pi^2R_i^2\ne0\), whereas every exact form has
zero integral on a closed three-cycle.  Adding an exact two-form
differential therefore cannot erase this distinction.
\end{proof}

\begin{proposition}[The exact data required by a connection consumer]
\label{p1b:prop:providers-consumer-map}
Let \(F:N\to M_8\) be a supplied smooth map, let
\(j:E\hookrightarrow F^*TM_8\otimes\mathbb C\) be a supplied
isometric bundle inclusion, and fix a connection from the invariant
metric-connection family.  Then
\[
 \nabla^E=j^*F^*\nabla j
\]
is a defined unitary connection, and a change \(\nabla\mapsto\nabla+K\)
changes it by exactly \(j^*F^*Kj\).  A supplied spin structure and
metric on \(N\) then define its twisted spin operator, with change
\[
 \delta D=\sum_\alpha c_N(\varepsilon^\alpha)
                       \otimes j^*(F^*K)_{\varepsilon_\alpha}j.
\]
Alternatively, a supplied principal-bundle reduction and unitary
representation define an associated connection by the differential of
that representation.  These maps, or another explicit bundle map with
the same role, are required to turn a parameter of one connection
family into a parameter of an operator on another bundle.  A parameter
list by itself defines none of these maps.

For a fixed metric, Clifford bundle and twisting bundle, two twisting
connections give the same twisted spin Dirac differential operator
exactly when their connection one-forms agree.  This last assertion
concerns twisting connections; the \(160\)-dimensional kernel in
Theorem~\ref{p1b:thm:providers-dirac-quotient} shows why it would be
false for tangent connections contracted through their spin lifts.
\end{proposition}
\begin{proof}
Pullback and orthogonal projection are the operations in
Proposition~\ref{p1b:prop:providers-projection}; subtracting two such
connections gives the stated tensorial variation.  The definition of a
twisted spin Dirac operator is Clifford contraction of the tensor-product
connection, giving the second formula.  For an associated bundle, changing
the principal matrix by \(a\) changes the associated matrix by
\(d\rho(a)\), directly from its equivariant section definition.
These constructions have specified bases, fibres and transformation laws,
so their outputs are well-defined.  Without the supplied map, reduction
or inclusion, no expression identifies the two different connection
bundles.  Finally the matrices \(c_N(\varepsilon^\alpha)\) are linearly
independent and Hilbert--Schmidt orthogonal: their trace pairing is a
nonzero multiple of \(\delta_{\alpha\beta}\).  Taking the spinor
partial trace of a zero connection difference after multiplication by
each Clifford matrix forces every twisting coefficient to vanish.
\end{proof}

\begin{proposition}[The complex-family operator and its comparison data]
\label{p1b:prop:providers-complex-operator}
Every member \(J\) of the invariant complex family and every supplied
Hermitian metric \(h\) define the self-adjoint elliptic operator
\[
 \mathscr D_{J,h}=\sqrt2(\bar\partial_J+\bar\partial_{J,h}^*)
       \quad\hbox{on }\Lambda^{0,*}_JT^*M_8.
\]
Its square is \(2\Delta_{\bar\partial}\), its kernel is the full
Dolbeault cohomology in holomorphic degree zero, and its graded index
is zero.  Thus the complex-family parameters have an exact mathematical
operator map, without identifying this operator with a different
previously fixed spin or Hodge operator.

The canonical line has the holomorphic square root
\(L=\pi^*\mathcal O(-1,-1)\), where \(\pi\) is the holomorphic
torus fibration in Corollary~\ref{p1b:cor:v214-canonical-kodaira}.
For the corresponding spin structure there is a Clifford-module
identification
\[
 \Sigma\cong\Lambda^{0,*}_J\otimes L,
 \qquad\Lambda^{0,*}_J\cong\Sigma\otimes L^{-1}.
\]
The other spin structures differ by the two independent scalar-cycle
sign bits.
These are module identifications; an equality of differential operators
also requires the specific line, spin and Hermitian connections and all
their zero-order terms.  Smooth triviality of \(L\) does not supply a
holomorphic trivialization: it has infinite holomorphic order, as does
the canonical line.
\end{proposition}
\begin{proof}
Integrability gives \(\bar\partial_J^2=0\), so squaring the displayed
operator gives \(2\Delta_{\bar\partial}\).  Its principal symbol is
the standard Clifford action by exterior multiplication and contraction;
it is elliptic.  The Hermitian Hodge theorem on a compact manifold gives
its self-adjoint realization and its Dolbeault kernel.  The index is the
alternating sum of the Hodge numbers, which vanishes by
Theorem~\ref{p1b:thm:relational-dolbeault-hodge} (or the factor
\((1+w)^2\) in the full Hodge polynomial).

The canonical-bundle formula gives \(L^{\otimes2}=K_J\) directly.
The restriction of the complex spin representation to the unitary
frame group is the exterior \((0,*)\) representation tensored with
the half-canonical character; its Clifford exterior/contraction action
is unaffected by this scalar character.  Gluing with the square-root
line therefore gives exactly the two module identifications above.
Changing a spin lift multiplies its transition functions by a flat
real sign character of \(\pi_1(M_8)=\mathbb Z^2\), giving the stated
remaining choices.  Finally \(c_1(K_J)=0\) and torsion-freeness of
\(H^2(M_8;\mathbb Z)\) give \(c_1(L)=0\), hence smooth triviality.
If any positive power of \(L\) were holomorphically trivial, the same
would hold for a positive power of \(K_J=L^2\), contrary to the
canonical-bundle theorem.  Thus a smooth module isomorphism by itself
cannot set its holomorphic connection data to zero.  Connections and
zeroth-order terms must be compared separately as asserted.
\end{proof}

\subsubsection{Rank changes and meromorphic Schur reduction}
\label{p1b:sec:geo-response-providers}
All spaces in this subsection are finite-dimensional Hilbert spaces.
Adjoints, orthogonal projections and Moore--Penrose inverses use their
specified Hilbert metrics. Real quadratic functions of complex variables
are differentiated on the realification. The results supplement the
singular Hermitian relation without extending constant-rank inverse
derivatives across a change of rank.

\begin{lemma}[Reduced quadratic function and its null constraint]
\label{p1b:lem:geo-constrained-quadratic}
Let
\[
 H=\begin{pmatrix}A&B\\B^*&C\end{pmatrix}=H^*,\quad
 N=\ker C,\quad P_N=P_{\ker C},\quad S=A-BC^+B^*.
\]
For \(Q_j(x,y)=\langle(x,y),H(x,y)\rangle/2
-\operatorname{Re}\langle(j_x,j_y),(x,y)\rangle\), stationary elimination
of the \(N^\perp\)-component of \(y\) gives exactly
\begin{align*}
 Q_j^{\rm red}(x,n)
 &=\tfrac12\langle x,Sx\rangle
 -\operatorname{Re}\langle j_x-BC^+j_y,x\rangle\\
 &\quad+\operatorname{Re}\langle n,B^*x-j_y\rangle
 -\tfrac12\langle j_y,C^+j_y\rangle,\qquad n\in N,\\
 y&=C^+(j_y-B^*x)+n.
\end{align*}
Thus the remaining equations are
\(Sx+Bn=j_x-BC^+j_y\) and \(P_NB^*x=P_Nj_y\).
In particular the remaining equation is unconstrained on all of \(X\)
if and only if \(BN=0\), and then it still requires \(P_Nj_y=0\).

If \(H\succeq0\), then \(C\succeq0\), \(BN=0\), \(S\succeq0\), and
\[
 \ker H=\{(x,-C^+B^*x+n):x\in\ker S,\ n\in N\}.
\]
\end{lemma}
\begin{proof}
Write \(y=y_r+n\), \(y_r\in N^\perp\), and solve the projected
stationary equation as \(y_r=C^+(j_y-B^*x)\). Substitution, or square
completion on the invertible restriction of \(C\) to \(N^\perp\), gives
the displayed function. Differentiating it gives both remaining equations.
Its constraint has zero linear part in \(x\) precisely when
\(P_NB^*=0\), equivalently \(BN=0\).

If \(H\succeq0\), its restriction gives \(C\succeq0\). For \(n\in N\),
positivity on \((tx,n)\) for both signs of small real \(t\), and then
for \(ix\), gives \(Bn=0\). Consequently
\[
 \langle(x,y),H(x,y)\rangle
 =\langle x,Sx\rangle+
 \langle y+C^+B^*x,C(y+C^+B^*x)\rangle.
\]
Taking \(y=-C^+B^*x\) proves \(S\succeq0\). The two nonnegative
terms vanish exactly under the stated kernel conditions.
\end{proof}

\begin{theorem}[Input compatibility and output independence]
\label{p1b:thm:geo-null-output}
Let \(H=H^*\) on \(V\), and let \(L:V\to Z\) be linear. The following
are equivalent:
\begin{enumerate}
\item \(Hq=L^*h\) is solvable for every \(h\in Z\);
\item \(L\ker H=0\);
\item for every compatible right-hand side, \(Lq\) is independent of
the choice of solution \(q\).
\end{enumerate}
Under these conditions the input-output map is \(LH^+L^*\). It is
positive semidefinite when \(H\succeq0\). In the preceding lemma, if
\(H\succeq0\) and \(S\succ0\), then \(\ker H=0\oplus N\). A map
\(L=(L_X,0)\) therefore has input-output map \(L_XS^{-1}L_X^*\).
\end{theorem}
\begin{proof}
Hermiticity gives \(\operatorname{Ran}H=(\ker H)^\perp\), so the first
condition is equivalent to \(P_{\ker H}L^*=0\), or the second. Every
compatible solution is \(H^+j+k\), \(k\in\ker H\). Applying \(L\)
proves the equivalence with the third condition and the stated map.
Positivity follows from the spectral decomposition of \(H^+\). The
kernel formula follows from the lemma; for the special map its first
reduced equation is \(Sx=L_X^*h\), which proves the final formula.
\end{proof}

\begin{theorem}[Laurent coefficient of an interior crossing]
\label{p1b:thm:geo-interior-laurent}
Let \(C(z)\) be holomorphic near a real \(z_0\), Hermitian on the real
axis, and write \(\delta=z-z_0\). Suppose
\[
 C(z)=C_0-\delta M+O(\delta^2),\quad C_0=C_0^*,\quad
 N=\ker C_0,\quad M_N=P_NM|_N\succ0.
\]
Let \(A(z),B(z),j_x(z),j_y(z)\) be holomorphic, with the full matrix
Hermitian on the real axis. Put
\(B^\sharp(z)=\overline{B(\overline z)}^{\,T}\),
\(F=B(z_0)|_N\), and \(j_N=P_Nj_y(z_0)\). Then, on a punctured
neighborhood,
\begin{align*}
 C(z)^{-1}&=-\delta^{-1}P_NM_N^{-1}P_N+O(1),\\
 A(z)-B(z)C(z)^{-1}B^\sharp(z)
   &=\delta^{-1}FM_N^{-1}F^*+O(1),\\
 j_x(z)-B(z)C(z)^{-1}j_y(z)
   &=\delta^{-1}FM_N^{-1}j_N+O(1).
\end{align*}
The exact constraint at \(z_0\) is \(F^*x=j_N\). Its equivalence to the
residue equation is
\[
 F^*x=j_N\quad\Longleftrightarrow\quad
 \begin{cases}
 P_{\ker F}j_N=0,\\
 FM_N^{-1}(F^*x-j_N)=0.
 \end{cases}
\]
In particular a vanishing Schur residue does not remove the component
of the source in \(\ker F\). The constraint is only one part of the
full singular block equations; their finite remaining row is retained.
\end{theorem}
\begin{proof}
In \(N\oplus N^\perp\), the lower-right block is
\(C_1+O(\delta)\), with \(C_1\) invertible, the off-diagonal blocks
are \(O(\delta)\), and the upper-left block is
\(-\delta M_N+O(\delta^2)\). Its Schur complement is therefore also
\(-\delta M_N+O(\delta^2)\). The block inverse has upper-left block
\(-\delta^{-1}M_N^{-1}+O(1)\) and bounded other blocks. This proves
the first expansion; multiplication gives the other two.

Projecting \(B(z_0)^*x+C_0y=j_y(z_0)\) onto \(N\) gives the constraint.
If the two conditions on the right hold, put \(d=F^*x-j_N\). The first
makes \(d\perp\ker F\). The second makes \(k=M_N^{-1}d\in\ker F\),
whence \(0=\langle k,d\rangle=\langle k,M_Nk\rangle\), so \(d=0\).
The converse follows by substitution. No invariance of \(\ker F\)
under \(M_N\) was assumed.
\end{proof}

\begin{theorem}[Full pencil residues and projected residues]
\label{p1b:thm:geo-full-pencil-residues}
Let \(D(z)=K-zM\), with \(K=K^*\), \(M\succ0\), and let \(L:V\to Z\)
be linear. If \(\Pi_\lambda\) are the spectral projections of
\(T=M^{-1/2}KM^{-1/2}\), then
\[
 LD(z)^{-1}L^*=
 \sum_\lambda\frac{W_\lambda}{\lambda-z},\qquad
 W_\lambda=LM^{-1/2}\Pi_\lambda M^{-1/2}L^*\succeq0.
\]
The \(\lambda\)-eigenspace is absent from this projected rational
function exactly when \(LM^{-1/2}\Pi_\lambda=0\). Equivalently,
\(L\ker(K-\lambda M)=0\). A right-hand side \(j\) admits a solution
of \((K-\lambda M)q=j\) exactly when
\(j\perp\ker(K-\lambda M)\).

An interior block inverse can have a pole where \(D(z)^{-1}\) is
holomorphic. In particular, if \(D(z_0)\) is invertible, no singularity
of its chosen Schur coordinates creates a pole of the full inverse or
of the projected inverse.
\end{theorem}
\begin{proof}
Factor \(D(z)=M^{1/2}(T-zI)M^{1/2}\) and apply the finite spectral
theorem. Each \(W_\lambda=UU^*\),
\(U=LM^{-1/2}\Pi_\lambda\), is zero exactly when \(U=0\).
The map \(M^{-1/2}\) takes the \(\lambda\)-eigenspace of \(T\)
bijectively to \(\ker(K-\lambda M)\), proving the equivalence.
The range statement follows from Hermiticity. The adjugate formula makes
the full inverse holomorphic at every invertible matrix. For an explicit
example, \(\left(\begin{smallmatrix}a&b\\b&-\delta\end{smallmatrix}
\right)\), with real \(b\ne0\), is invertible at \(\delta=0\), although
its lower-right inverse and its Schur complement have poles there.
\end{proof}

\begin{theorem}[Continuity criterion across a positive rank change]
\label{p1b:thm:geo-positive-rank-change}
For a real parameter \(t\), let
\(H(t)=\left(\begin{smallmatrix}A(t)&B(t)\\B(t)^*&C(t)
\end{smallmatrix}\right)\succeq0\) be continuous at zero. Suppose
\(C(t)\succ0\) for \(t\ne0\) and let \(P(t)\) be the small spectral
projection converging to \(P_{\ker C(0)}\); its orthogonal complement
stays uniformly gapped. Then the regular Schur complements
\(S(t)=A(t)-B(t)C(t)^{-1}B(t)^*\) converge to
\(S(0)=A(0)-B(0)C(0)^+B(0)^*\) if and only if
\[
 \left\|B(t)P(t)
       \bigl(C(t)|_{\operatorname{Ran}P(t)}\bigr)^{-1/2}\right\|
       \longrightarrow0.
\]
Positivity alone gives only
\(\|B(t)v\|^2\le\|A(t)\|\gamma\) for each unit eigenvector
\(C(t)v=\gamma v\), \(\gamma>0\), and does not imply this continuity.
Indeed
\[
 H(t)=\begin{pmatrix}2&t\\t&t^2\end{pmatrix}\succeq0
 \quad\text{has}\quad S(t)=1\ (t\ne0),\qquad S(0)=2.
\]
\end{theorem}
\begin{proof}
Choose the small band by a fixed spectral interval separating zero from
the positive eigenvalues of \(C(0)\). On its complement the inverse and
its Schur term converge continuously to the inverse term of \(C(0)^+\).
The remaining term is
\[
 E(t)=B(t)P(t)
       \bigl(C(t)|_{\operatorname{Ran}P(t)}\bigr)^{-1}P(t)B(t)^*,
\]
which is positive and has norm equal to the square of the stated norm.
It is absent from the pseudoinverse term at zero. This proves necessity
and sufficiency. Positivity on the two-dimensional form spanned by
\((x,0),(0,v)\) gives
\(|\langle x,Bv\rangle|^2\le\langle x,Ax\rangle\gamma\); maximize
over unit \(x\). In the example \(\det H(t)=t^2\) and direct
substitution gives the asserted Schur complements.
\end{proof}

\begin{proposition}[Necessary jets of a singular stationary branch]
\label{p1b:prop:geo-singular-stationary-jets}
Let \(f(q,\eta)\) be a real \(C^3\) function and let \(\bar q(\eta)\)
be a \(C^2\) stationary branch. At a singular point put
\(H=f_{qq}\), \(J=f_{q\eta}\), \(P=P_{\ker H}\). Necessarily
\[
 PJ=0,\qquad \dot{\bar q}=-H^+J+k,\quad k\in\ker H,
\]
and
\[
 P\bigl(f_{qqq}[\dot{\bar q},\dot{\bar q}]
             +2f_{qq\eta}\dot{\bar q}+f_{q\eta\eta}\bigr)=0.
\]
The first compatibility condition is sufficient for a solution of the
first linearized equation, but does not establish a nonlinear stationary
branch or select its kernel tangent.
\end{proposition}
\begin{proof}
Differentiate \(f_q(\bar q(\eta),\eta)=0\) once and twice:
\[
 H\dot{\bar q}+J=0,\qquad
 H\ddot{\bar q}+f_{qqq}[\dot{\bar q},\dot{\bar q}]
       +2f_{qq\eta}\dot{\bar q}+f_{q\eta\eta}=0.
\]
Orthogonal projection to \(\ker H\), and inversion on its complement,
give the assertions. For example \(f(q,\eta)=q^3/3+\eta^2q\) has
\(H=J=0\) at the origin, so its first equation is compatible, while its
stationary equation \(q^2+\eta^2=0\) has no real branch for
\(\eta\ne0\). Its second compatibility would require
\(2(\dot{\bar q})^2+2=0\).
\end{proof}

\begin{proposition}[Positive rational gap sums and closing-gap scaling]
\label{p1b:prop:geo-positive-gap-scaling}
Let \(A_r\succeq0\), \(\Delta_r>0\), for a fixed finite index set, and
define the operator-valued rational function
\[
 R(z)=2\sum_r\frac{\Delta_r A_r}{\Delta_r^2-z^2}
     =S_0+z^2T+O(z^4),\qquad
 S_0=2\sum_r\frac{A_r}{\Delta_r},\quad
 T=2\sum_r\frac{A_r}{\Delta_r^3}.
\]
If every nonzero term has \(\delta\le\Delta_r\le D\), then
\[
 D^{-2}S_0\preceq T\preceq\delta^{-2}S_0,
 \qquad \ker S_0=\ker T=\bigcap_{A_r\ne0}\ker A_r.
\]
For a parameter family with fixed finite index set, bounded \(S_0\)
forces \(\|A_r\|=O(\Delta_r)\) for every closing gap, while bounded
\(T\) forces \(\|A_r\|=O(\Delta_r^3)\). If the other gaps stay
uniformly positive and their matrices stay bounded, these respective
conditions on each closing gap are also sufficient. They are distinct:
\(A_r=\Delta_r A\), \(A\ne0\), has bounded contribution to \(S_0\)
and unbounded contribution to \(T\).
If a nonzero coupled gap tends to zero, these rational functions have no
common nonzero disk of holomorphy about zero.
\end{proposition}
\begin{proof}
Expand each denominator for \(|z|<\min_{A_r\ne0}\Delta_r\). Multiplying
the positive term \(2A_r/\Delta_r\) by
\(D^{-2}\le\Delta_r^{-2}\le\delta^{-2}\) and summing gives the bounds.
A positive finite sum vanishes on a vector precisely when every positive
summand does, giving the kernels. The inequalities
\(0\preceq2A_r/\Delta_r\preceq S_0\) and
\(0\preceq2A_r/\Delta_r^3\preceq T\) prove the necessary norm rates.
With finitely many summands the corresponding termwise bounds prove
sufficiency under the stated bounded-remainder condition. At any
\(\Delta>0\) shared by one or more gaps the residue is minus the sum
of their \(A_r\); positivity prevents its cancellation if any is nonzero.
A closing coupled gap therefore brings a pole arbitrarily close to zero.
\end{proof}

\begin{proposition}[Bipartite first-order realization and internal pole cancellation]
\label{p1b:prop:geo-bipartite-resolvent}
Let \(X=X_0\oplus E\oplus X_s\), and let \(G:Y\to E\) be invertible.
Let \(A_s=A_s^*\) on \(X_s\), let \(A=0\oplus0\oplus A_s\) on \(X\),
and extend \(G\) by zero into \(X\). Put
\[
 D=\begin{pmatrix}A&G\\G^*&0\end{pmatrix},\qquad
 \Lambda=GG^*|_E\succ0.
\]
For every regular \(z\ne0\),
\[
 P_X(D-zI)^{-1}|_X
 =-\frac{P_{X_0}}z+(A_s-zI)^{-1}P_{X_s}
             +z(\Lambda-z^2I)^{-1}P_E.
\]
The \(E\)-component is holomorphic and zero at \(z=0\). Each
\(\Lambda\)-eigenspace of eigenvalue \(s^2>0\) has projected spectral
weight \(1/2\) at both \(D\)-eigenvalues \(\pm s\). At zero the
equation \(D(x,y)=(j_x,0)\) requires \(x_E=0\) and \(j_{X_0}=0\),
as well as its \(X_s\)-equation. Thus replacing the zero interior block
by its zero pseudoinverse without its constraint loses \(\dim E\)
independent equations.
\end{proposition}
\begin{proof}
For \(z\ne0\), eliminating the interior block \(-zI\) gives the Schur
complement \(A-zI+GG^*/z\). Invert it on the three orthogonal
summands. On \(E\), its inverse is \(z(\Lambda-z^2I)^{-1}\), proving
the formula and regularity at zero. On an \(s^2\)-eigenspace,
\[
 \frac{z}{s^2-z^2}
  =\frac12\left(\frac1{s-z}+\frac1{-s-z}\right),
\]
which gives the weights. At zero the second block equation is
\(G^*x_E=0\), and invertibility of \(G\) gives \(x_E=0\); the
\(X_0\)-part of the first is \(0=j_{X_0}\). These are exact equations
of the full Hermitian matrix, independent of any choice of coordinates
for a Schur inverse.
\end{proof}

\subsubsection{A prescribed weight in the Cartan profile functional}
\label{p1b:sec:v215-weighted-profile}
\begin{proposition}[Weighted profile equation and persistence test]
\label{p1b:prop:v215-weighted-profile}
Let $Z$ be a positive $C^1$ function on an interval and consider
$I_Z(f)=\int Z(f'^2+V(f))\,dt$, where
$V(f)=4f^2(1-f)^2$. Under compactly supported variations its Euler
equation is
\[
 (Zf')'-\tfrac12ZV'(f)=0.
\]
If $f_*$ is an unweighted critical profile, then it is also critical
for $I_Z$ if and only if $Z'f_*'=0$ pointwise. In particular the
nonconstant logistic profile $f_*'=2f_*(1-f_*)$ remains critical on a
connected open interval precisely when $Z$ is constant there.
The constant profiles $0,1,1/2$ remain critical for every such weight.
\end{proposition}
\begin{proof}
Integration by parts gives the Euler expression
$2(Zf')'-ZV'(f)$. On an unweighted solution
$f_*''=V'(f_*)/2$, its half is exactly $Z'f_*'$. The logistic derivative
is strictly positive at every finite point, which gives the constant-
weight conclusion. Each stated constant is a critical point of $V$.
\end{proof}

\subsubsection{The logarithmic coordinate of a one-variable Legendre function}
\label{p1b:sec:v215-log-legendre}
\begin{proposition}[Exact logarithmic Hessian and its derivative]
\label{p1b:prop:v215-log-legendre-hessian}
Let $f\in C^3(I;\R)$, with $F=f'>0$ and $\chi=f''>0$ on an
interval $I$, and let $a,b>0$. Put
\[
 \phi=a\log F(R),\qquad P=RF-f,\qquad B=2f-RF,
 \qquad W(\phi(R))=bP/F^2.
\]
This coordinate is strictly increasing, and on its image
\begin{align*}
 W_\phi&=\frac ba\frac B{F^2},\\
 W_{\phi\phi}
 &=\frac b{a^2}\left(\frac1\chi-\frac R F-\frac{2B}{F^2}\right),\\
 \frac{d}{dR}W_{\phi\phi}
 &=\frac b{a^2}\left(-\frac{\chi_R}{\chi^2}
       -\frac{3(F-R\chi)}{F^2}+\frac{4B\chi}{F^3}\right).
\end{align*}
In particular the Hessian is not generally the rescaled expression
$(b/a^2)(F-R\chi)/(\chi F)$: their difference is
$-2bB/(a^2F^2)$. They agree exactly where $B=0$.
\end{proposition}
\begin{proof}
Here $\phi_R=a\chi/F>0$, $P_R=R\chi$ and $B_R=F-R\chi$.
Thus $d/d\phi=F(a\chi)^{-1}d/dR$. Applying this differential
operator first to $bP/F^2$ and then to its first derivative gives the
first two identities. Differentiation in $R$ of the second identity,
using $B_R=F-R\chi$, gives the third. Subtracting the indicated
rescaled expression gives the final statement. All formulas are local
on the displayed regular coordinate interval.
\end{proof}

\subsection{A compact-pulse scattering estimate for the canonical scalar equation}
\label{p1b:sec:v215-pulse-provider}
The following one-dimensional estimate is independent of any field
interpretation. Its frequency is a positive real number; a later spectral
application may supply a discrete frequency set.

\begin{theorem}[Uniform first coefficient and an explicit Volterra remainder]
\label{p1b:thm:v215-compact-pulse-scattering}
Let $h\in C_c^\infty((a,b);\R)$, $\lambda\in\R$ and
$|\lambda|\|h\|_\infty\le\delta<1$. Put
\[
 Z=1+\lambda h,\qquad s=\sqrt Z,\qquad V_\lambda=s''/s.
\]
For $\omega>0$ let $u$ be the unique solution of
$u''+(\omega^2-V_\lambda)u=0$ equal to
$e^{-i\omega t}/\sqrt{2\omega}$ for $t\le a$. For $t\ge b$ write
\[
 u(t)=\frac{\alpha_\omega e^{-i\omega t}
                    +\beta_\omega e^{i\omega t}}{\sqrt{2\omega}}.
\]
Then $|\alpha_\omega|^2-|\beta_\omega|^2=1$ and, with
$\widehat h(\xi)=\int h(t)e^{-i\xi t}\,dt$,
\[
 \beta_\omega=i\lambda\omega\widehat h(2\omega)+r_\omega.
\]
Define
\[
 A=\tfrac12\|h''\|_1,\qquad
 B_\delta=\frac{\|hh''\|_1}{2(1-\delta)}
            +\frac{\|h'^2\|_1}{4(1-\delta)^2},\qquad
 L=|\lambda|A+\lambda^2B_\delta.
\]
The exact remainder satisfies
\[
 \boxed{|r_\omega|\le e_\omega
 :=\frac{\lambda^2B_\delta}{2\omega}
  +\frac12\left(e^{L/\omega}-1-\frac L\omega\right).}
\]
In particular this is $O(\lambda^2)$ uniformly for $\omega\ge\omega_0>0$.
For $b_\omega=i\lambda\omega\widehat h(2\omega)$,
\[
 \bigl||\beta_\omega|^2-|b_\omega|^2\bigr|
       \le2|b_\omega|e_\omega+e_\omega^2.
\]
Every finite positively weighted sum of these squared coefficients has
the sum of the corresponding displayed error bounds.
\end{theorem}
\begin{proof}
The real smooth compactly supported potential gives a unique global
solution of the linear initial-value problem. Its constant Wronskian
$u\bar u'-u'\bar u=i$ proves the asserted coefficient identity.
Direct differentiation gives
\[
 V_\lambda=\frac{\lambda h''}{2(1+\lambda h)}
       -\frac{\lambda^2h'^2}{4(1+\lambda h)^2}
       =\frac\lambda2h''+R_\lambda,
 \qquad \|R_\lambda\|_1\le\lambda^2B_\delta.
\]
Thus $\|V_\lambda\|_1\le L$. For $y=\sqrt{2\omega}\,u$ the
exact Volterra equation is
\[
 y(t)=e^{-i\omega t}+\int_a^t
       \frac{\sin\omega(t-v)}\omega V_\lambda(v)y(v)\,dv.
\]
Iteration of this triangular integral, or its integral Gronwall bound,
gives $|y(t)-e^{-i\omega t}|\le e^{M(t)/\omega}-1$, where
$M(t)=\int_a^t|V_\lambda|$. Reading off the outgoing exponential gives
\[
 \beta_\omega=\frac1{2i\omega}\int_a^b
                 e^{-i\omega t}V_\lambda(t)y(t)\,dt.
\]
Two integrations by parts, with zero endpoint terms, make its linear
term $\lambda\widehat{h''}(2\omega)/(4i\omega)
=i\lambda\omega\widehat h(2\omega)$.
The $R_\lambda$ term against the incoming exponential is bounded by
$\lambda^2B_\delta/(2\omega)$. The remaining term is bounded by
\[
 \frac1{2\omega}\int_a^b |V_\lambda(t)|
            (e^{M(t)/\omega}-1)\,dt
 =\frac12(e^{M(b)/\omega}-1-M(b)/\omega),
\]
which proves the estimate. The squared-modulus inequality follows by
expanding $|b+r|^2$ and applying the triangle inequality. Uniformity
uses the Taylor expansion of the exponential with
$\omega\ge\omega_0$ and bounded $|\lambda|$. Finite weighted sums
preserve the bounds term by term.
\end{proof}

\begin{remark}[A first jet is not a second-jet estimate]
\label{p1b:rem:v215-pulse-jet-domain}
The stated bound uses the exact family $Z=1+\lambda h$. Knowing only
$Z(\varepsilon,t)=1+\varepsilon h_1(t)+O(\varepsilon^2)$ pointwise
does not control $(\sqrt Z)''/\sqrt Z$. A parameter-family application
must retain a uniform $C^2$ remainder (and its support or integrable
derivative bounds); its contribution to $\|R\|_1$ is then added to
the same Volterra estimate. No interchange of an unbounded-frequency
sum and a parameter or long-interval limit is asserted here.
\end{remark}

\subsection{End gauges and the actual intermediate principal descent fibre}
\label{p1b:v215:sec:principal-fibre}
Put \(A=\operatorname{SP}^2S^3\), \(P=S^3\times S^3\),
\(Y=BSU(2)\), and retain the actual symmetric-square cofibration
\[
 S^3\xrightarrow{i}A\xrightarrow{c}M_5
 \xrightarrow{\partial}S^4,\qquad M_k=S^k\cup_2e^{k+1}.
\]
The bottom restriction of \(\partial\) is \(\eta_4\), and the
already proved values are \([A,Y]=0\) and
\(\operatorname{im}(i^*:[A,S^3]\to\mathbb Z)=4\mathbb Z\).
These are the cofibration and degree assertions of
Lemma~\ref{p1b:lem:pure-descent-symmetric-loop-image}; the actual
null map \(cq_Y\) used below is the suspended double-cover map in
the symmetric-square no-go proof of
Theorem~\ref{p1b:thm:g21-exotic-nogo}.  The specified double cylinder
and mapping-space pullback are those of
Theorem~\ref{p1b:thm:principal-descent-mapping-torus} and
Corollary~\ref{p1b:cor:principal-descent-source-cohomology}.
Write \(q_Y:P\to A\) for the quotient. Its actual gauge image is
\[
 H=\langle\ell\rangle\subset G=[P,S^3],\quad
 \ell=(4,4;z_0),\quad z_0\in\{2,8\},
\]
where
\((a,b;z)(a',b';z')=(a+a',b+b';z+z'+ab')\), \(z\in C_{12}\).
No argument below selects \(z_0\) without evaluating the actual map.

\begin{proposition}[The complete end-gauge groups in the needed degrees]
\label{p1b:v215:prop:end-gauges}
Abstractly,
\[
 [A,S^3]=\mathbb Z\oplus C_2,\qquad
 [\Sigma A,S^3]=C_2.
\]
The first group's degree image is \(4\mathbb Z\); its order-two
kernel is killed by \(q_Y^*\). The map
\[
 (\Sigma q_Y)^*:[\Sigma A,S^3]\longrightarrow
                    [\Sigma P,S^3]=C_2^3
\]
is zero. Equivalently the actual map
\(\pi_2\operatorname{Map}(A,Y)\to\pi_2\operatorname{Map}(P,Y)\)
is zero.
\end{proposition}
\begin{proof}
The Moore exact sequences have end groups \(C_2,C_2\):
\[
 0\to C_{12}/2C_{12}\to[M_5,S^3]\to(\pi_5S^3)[2]\to0,
\quad
 0\to\pi_7S^3/2\pi_7S^3\to[M_6,S^3]\to(\pi_6S^3)[2]\to0.
\]
The identity \(2\operatorname{id}_{M_k}=i\eta_kc\) shows that
twice a lift of \(\eta_3\eta_4\) in the first is \((6\nu')c=0\)
in the injected \(C_{12}/2C_{12}\). Twice a lift of \(6\nu'\)
in the second is \((6\nu'\eta_6)c=0\). Thus both middle groups
are \(C_2^2\), rather than \(C_4\).

The Puppe map \(\pi_4S^3\to[M_5,S^3]\) is injective: the
bottom restriction of \(\eta_3\partial\) is
\(\eta_3\eta_4\ne0\). Its quotient is a single \(C_2\), the
degree-zero kernel of \([A,S^3]\to\pi_3S^3\). The established
degree image is infinite cyclic \(4\mathbb Z\), so choosing one
degree-four map splits this extension. Its order-two kernel is central
because \(\operatorname{Aut}(C_2)=1\).

The next exact segment is
\[
 \pi_5S^3\longrightarrow[M_6,S^3]\longrightarrow
 [\Sigma A,S^3]\longrightarrow\pi_4S^3
                   \xrightarrow{\partial^*}[M_5,S^3].
\]
Its first map sends \(\eta_3\eta_4\) to a class with bottom value
\(\eta_3\eta_4\eta_5=6\nu'\ne0\). Its last map is injective,
so \([\Sigma A,S^3]=C_2\), represented by the surviving top Moore
class; every such map factors through \(\Sigma c\).
The actual composite \(cq_Y\) factors, after the coordinate wedge
is collapsed, through the already proved null map
\(\Sigma^4(S^2\to\mathbb{RP}^2):S^6\to M_5\).
It and its suspension are null. This kills the first kernel and the
whole suspended group on pullback. Evaluation changes none of these
mapping groups because \(Y\) is three-connected.
\end{proof}

\begin{theorem}[Every principal bundle component on the double cylinder]
\label{p1b:v215:thm:F7-components}
For \(F_7=A_-\cup_{q_Y}(P\times I)\cup_{q_Y}A_+\),
\[
 [F_7,Y]=H\backslash G/H.
\]
Put \(d=a-b\) and \(h_d=\gcd(12,4d)=4\gcd(3,d)\), where
\(\gcd(3,0)=3\). Its unique normal forms are
\[
 \boxed{(u,u-d;\bar z),\quad u\in\{0,1,2,3\},\quad
                                  \bar z\in C_{h_d}.}
\]
There are \(16\gcd(3,d)\) classes at each \(d\). Orienting the
Mayer--Vietoris transgression by \((a,b)\mapsto a-b\), this is
the second Chern coefficient in \(H^4(F_7;\mathbb Z)=\mathbb Z\).
\end{theorem}
\begin{proof}
Both end mapping spaces are connected. Homotope their two maps to the
constant map. The intervening path then becomes an element of
\(\pi_1\operatorname{Map}(P,Y)_0=G\). Changing the two end homotopies
acts on its two sides by the actual image \(H\), proving the double
coset description including its gauge equivalence.
The actual powers are
\(\ell^r=(4r,4r;rz_0+8r(r-1))\). For \(k=r+s\), direct multiplication gives
\[
 \ell^r(a,b;z)\ell^s=
 (a+4k,b+4k;z+kz_0+8k(k-1)+4ka+4r(b-a)).
\]
Normalize \(a\) modulo four; then \(b=u-d\). A transformation
preserving that normalized pair has \(k=0\) and translates the
central coordinate by \(\langle4d\rangle\). The resulting quotient
is \(C_{h_d}\), proving uniqueness and the count independently of
\(z_0\). The Chern class is the clutching transgression into
\(\mathbb Z^2/\mathbb Z(1,1)\); its chosen coordinate is \(a-b\).
\end{proof}

\begin{proposition}[Exact component monodromy]
\label{p1b:v215:prop:F7-monodromy}
The specified involution \(R\) acts on raw representatives by
\[
 \boxed{R(a,b;z)=(-b,-a;z).}
\]
For a normalized row, choose \(u'\in\{0,1,2,3\}\) congruent to
\(d-u\pmod4\) and put \(k=(u'+u-d)/4\). Its exact formula is
\[
 (u,u-d;\bar z)\mapsto
 (u',u'-d;\overline{z+kz_0+8k(k-1)+4k(d-u)}).
\]
Odd \(d\) has no invariant component. For even \(d\), the actual
candidate \(z_0=8\) gives \(8\gcd(3,d)\) invariant components,
whereas \(z_0=2\) gives \(4\gcd(3,d)\). These are components
of the seven-dimensional fibre, not numerical descent fibres on \(Q\).
\end{proposition}
\begin{proof}
The involution exchanges endpoints, reverses the interpolating path and
swaps the two sphere factors. Thus its raw action is \(\sigma(g)^{-1}\).
The known \(\sigma(a,b;z)=(b,a;ab-z)\) proves the first formula;
right multiplication by \(\ell^k\) normalizes it and proves the second.
A fixed row requires \(2u=d\pmod4\), giving two rows exactly when
\(d\) is even. The central correction is automatically zero modulo
three when \(3\mid d\). Modulo four it is \(kz_0\). Both rows
survive for \(z_0=8\); for \(z_0=2\), the additional condition
\(2u=d\pmod8\) selects one row. Each has \(h_d\) central choices.
\end{proof}

\begin{proposition}[The exact remaining fundamental-group extension]
\label{p1b:v215:prop:F7-gauge-extension}
At \(g=(a,b;z)\), put \(k_g=3/\gcd(3,a-b)\). There is a central
extension
\[
 \boxed{0\to C_2^3\to\pi_1\operatorname{Map}(F_7,Y)_g
                 \to k_g\mathbb Z\oplus C_2^2\to0.}
\]
The infinite coordinate is the common degree-four endpoint power; the
two finite coordinates are the two separate invisible end gauges.
No splitting or extension class is asserted by this sequence alone.
\end{proposition}
\begin{proof}
The homotopy-pullback exact sequence and the zero maps on \(\pi_2\)
just proved inject the entire target \(C_2^3\). The quotient is the
stabilizer under the end loop groups \((\mathbb Z\oplus C_2)^2\).
Their sphere degrees force a common infinite power \(r\), and then
\[
 g\ell^rg^{-1}=\ell^r
 \ \Longleftrightarrow\ 12\mid4r(a-b)
 \ \Longleftrightarrow\ k_g\mid r.
\]
The two invisible end-kernel coordinates remain independent. Their
action on the injected module is the established G20 action of their
common target image, whose sphere degrees \((4r,4r)\) are even.
The formula \((r_1,r_2,e)\mapsto(r_1,r_2,e+br_1+ar_2)\)
from Paper~I-D, Section~\ref{p1b:app:g20} is therefore the identity.
This proves centrality.
\end{proof}

\subparagraph{The remaining origin-sensitive calculation.}
For the actual mapping torus \(Q\), choose a point \(x\) in an
\(R\)-invariant component and one path \(\kappa:x\to Rx\).
Let \(\Gamma=\pi_1\operatorname{Map}(F_7,Y)_x\). Its contribution
to \([Q,Y]\) is exactly the twisted-conjugacy set
\[
 \Gamma/\{u\sim h^{-1}u\Theta_\kappa(h)\},\qquad
 \Theta_\kappa(h)=\kappa(Rh)\kappa^{-1}.
\]
Changing \(\kappa\) to \(v\kappa\) conjugates the automorphism by
\(v\) and translates the coordinates, preserving the underlying set.
Its square is \(\operatorname{Ad}_{\kappa R\kappa}\), which is the identity for the abelian group established below in
Theorem~\ref{p1b:v215:thm:F7-gauge-split}. The actual involution on its
generators and the ordered-cover pullback still require evaluation. A
canonical zero for the twisted-loop coordinates does not follow from the
abstract group list. The completed mixed \(U(2)\) classification does
not supply those distinct maps.

\subsection{Principal stabilization and the geometric carrier selection}
\label{p1b:v215:sec:principal-stabilization}
Write \(M=T^2\times S^3\times S^3\),
\(c_1=0\), \(c_2=Aa_tu_x+Ba_tu_p+Ca_Eu_x+Da_Eu_p\), and
\(\Delta=BC-AD\).  Assume \(12\mid\Delta\), and set
\[
 N=\begin{pmatrix}B&-A\\D&-C\end{pmatrix},\qquad
 T_m=\mathbb Z^2/(m\mathbb Z^2+N\mathbb Z^2),\qquad Q_m=|T_m|.
\]
All source objects below carry a specified determinant trivialization,
and all source isomorphisms preserve it.

\begin{theorem}[Uniform principal stabilization fibres]
\label{p1b:v215:thm:principal-stabilization}
The actual inclusion \(SU(2)\hookrightarrow SU(3)\) induces a
surjection from the principal rank-two fixed-\(c_2\) fibre, separately
in each restriction sector \(\varepsilon\in C_2\), onto the principal
rank-three fibre with \(c_2\) as prescribed and \(c_3=0\).  The latter
is an affine copy of \(T_6\).  Every target has exactly
\[
 \boxed{\nu_{\rm pr}\,Q_{12}/Q_6}
\]
preimages in \emph{each} restriction sector, where
\[
 \nu_{\rm pr}=\begin{cases}
 4,&\text{some of }A,B,C,D\text{ is odd},\\
 4+2^{\operatorname{rank}_{\mathbb F_2}(N/2)},
       &\text{all four are even}.
 \end{cases}
\]
At \(c_2=0\), there are thirty-six target classes and forty principal
rank-two preimages of each: twenty in each sector.  This differs from
the ordinary complex-bundle fibre of twenty-five because forgetting
the source determinant trivialization is a further quotient.
\end{theorem}
\begin{proof}
The determinant-preserving calculation of
Theorem~\ref{p1b:thm:g20-rank2-fibre-even}, specifically the
``Orbit counts'' proof in Paper~I-D, Section~\ref{p1b:app:g20}, gives, at
each primary in \(T_{12}\), exactly \(\nu_{\rm pr}\) secondary
orbits in either restriction sector.  The displayed formula includes
its zero and rank-one cases: if \(N\) has integral rank one and
nonzero Smith invariant \(d\), then \(Q_{12}=12\gcd(12,d)\),
and \(\nu_{\rm pr}\) is respectively \(4,6,5\) when \(d\) is
odd, congruent to two modulo four, or divisible by four.

For the target put \(P=S^3\times S^3\).  The condition \(c_3=0\)
makes its restriction to \(P\) trivial.  The target mapping-space
component has
\[
 \pi_1\operatorname{Map}(P,BSU(3))_0
   =\{(a,b;z):a,b\in\mathbb Z,\ z\in C_6\},
 \qquad \pi_2\operatorname{Map}(P,BSU(3))_0=0.
\]
The first multiplication is the level-six Heisenberg multiplication;
the second vanishing follows from
\(\Sigma P\simeq S^4\vee S^4\vee S^7\),
\(\pi_4SU(3)=0\), and \(\pi_7SU(3)=0\).
Maps of the two-torus are therefore commuting edge pairs modulo
simultaneous conjugation, without a retained finite face module.
The fixed free edge degrees are the prescribed Chern coefficients;
their central commutator vanishes since \(6\mid\Delta\).
Conjugation changes their two central coefficients by \(N\mathbb Z^2\),
so their target fibre is an affine \(T_6\).

Naturality of the block inclusion gives the identity on \(\pi_3\),
reduction \(C_{12}\to C_6\) on \(\pi_6\), and zero on the
\(\pi_4\) and \(\pi_7\) source differences.  It kills the slice
restriction bit as well.  Choose a target origin and, separately in
each source sector, an origin mapping to it.  The stabilization map
in these charts is exactly \(T_{12}\to T_6\), with every secondary
orbit sent to its primary.  This quotient homomorphism is onto and
has constant kernel size \(Q_{12}/Q_6\), proving both surjectivity
and the fibre formula.  When \(N=0\), these values are
\(Q_{12}=144\), \(Q_6=36\), and \(\nu_{\rm pr}=5\).
\end{proof}

\begin{remark}[Why the target determinant quotient changes nothing here]
\label{p1b:v215:rem:target-determinant}
For this \(c_3=0\) target fibre, forgetting its determinant
trivialization is bijective.  Its determinant boundary is conjugation
by a full unitary determinant loop.  On the two sphere-degree
generators this operation lands in \(\pi_4SU(3)=0\), and on the
central generator it lands in \(\pi_7SU(3)=0\).  The slice map is
literally constant, and the face module is zero.  Thus it induces the
identity on all the target torus charts above.  This argument uses the
actual target degrees; it supplies no such conclusion for the
rank-two source.
\end{remark}

\begin{corollary}[Stabilization cannot settle the principal descent origin]
\label{p1b:v215:cor:stabilization-loses-descent}
The principal stabilization calculation does not distinguish the
two remaining degree-four endpoint images
\(\ell=(4,4;2)\) and \(\ell=(4,4;8)\): both have image
\((4,4;2)\) in the level-six target.  Nor does stable triviality
imply the existence of a principal Born lift.
\end{corollary}
\begin{proof}
Their difference is the central element six, exactly the kernel of
\(\pi_6SU(2)\to\pi_6SU(3)\).  Further stabilization loses the
remaining fixed-Chern rank-three differences.  More concretely the
nonzero rank-two restriction bundle on \(S^3\times S^3\), pulled
back to \(M\), has zero stable class difference and admits no
principal exchange lift by the actual symmetric-square no-go.  It is
therefore already a counterexample to inferring descent from stable
data.  Compatible affine origins in a stabilization chart are choices,
not a replacement for the actual descent map or its homotopies.
\end{proof}

The following proposition concerns the specified mathematical
carrier.  Its application to the native weak role is conditional
on \(\mathsf H_{\rm vert}\) of
Definition~\ref{p1b:def:conditional-residual-datum}, discharged
later in Paper~III-F.

\begin{proposition}[The specified residual doublet selects an actual point]
\label{p1b:v215:prop:carrier-selection}
In the compact residual presentation
\[
 SU(2)_{\rm diag}\longrightarrow
 SU(2)_x\times SU(2)_p\times U(1)_{\rm rel}
 \longrightarrow S^3_{\rm rel}\times S^1_{\rm rel},
\]
the section \(s(r,z)=(r,e,z)\) makes the principal bundle and its
fundamental associated doublet trivial.  Its pullback to \(M\)
is one specified principal rank-two class, with \(c_2=0\) and
\(\varepsilon=0\); the descended exchange action has fixed
isotropy \(+I\).  It therefore selects one point, rather than
the entire stabilization fibre, in the native bundle classification.
\end{proposition}
\begin{proof}
Right diagonal multiplication preserves \(r=g_xg_p^{-1}\), and
every point has the unique form \(s(r,z)h\).  The displayed map is
therefore a global principal section.  Pullback preserves this
trivialization.  Exchange sends the section representative to
\((e,r,z^{-1})=s(r^{-1},z^{-1})\cdot r\).  On the exchange
fixed locus of the original two factors, \(r=e\), so the vertical
factor acts by the identity.  These are statements about the bundle
and its equivariant structure; they require no dynamics or state
selection.
\end{proof}

\subsubsection{The intermediate gauge extension splits}
\label{p1b:v215:sec:intermediate-gauge-splitting}

\begin{lemma}[Skeletal commutator filtration for a group bundle]
\label{p1b:v215:lem:twisted-skeletal-commutator}
Let \(\mathcal G\to X\) be a locally trivial bundle of topological
groups over a finite CW complex, with transition functions preserving
multiplication and identity.  In the group of homotopy classes of
sections, let \(F^p\) consist of classes represented by the identity
section on \(X^{(p-1)}\).  Then
\[
 [F^p,F^q]\subset F^{p+q}.
\]
This applies to the adjoint group bundle of any principal \(SU(2)\)
bundle, including a nontrivial one.
\end{lemma}
\begin{proof}
On \(X^{(p+q-1)}\), cellularly approximate the diagonal by
\(D=(D_1,D_2):X^{(p+q-1)}\to X\times X\).  The image is in the
union of product cells with dimensions \(i+j\le p+q-1\).
Thus at each point either \(D_1\) lands in \(X^{(p-1)}\)
or \(D_2\) lands in \(X^{(q-1)}\).  The two coordinate
homotopies from the identity to \(D_i\) give group-bundle
isomorphisms between \(\mathcal G\) and \(D_i^*\mathcal G\);
one may construct them by trivializing the pulled-back bundle along
the homotopy interval.  Pull a representative of \(F^p\) and one
of \(F^q\) back into the original fibre using these isomorphisms.
At every point at least one is the identity, so their commutator is
the identity.  The homotopies identify this commutator with the
original one on that skeleton.  Extend this homotopy of sections to
\(X\) by the relative homotopy lifting property.  It gives a
representative equal to the identity on \(X^{(p+q-1)}\), as required.
\end{proof}

\begin{theorem}[The actual abstract component gauge group]
\label{p1b:v215:thm:F7-gauge-split}
For every bundle component \(g\) on \(F_7\),
\[
 \boxed{\pi_1\operatorname{Map}(F_7,BSU(2))_g
                  \cong\mathbb Z\oplus C_2^5.}
\]
The isomorphism is not a distinguished coordinate chart.  In
particular the common endpoint degree is still restricted to
\((3/\gcd(3,d))\mathbb Z\) in the exact endpoint sequence.
\end{theorem}
\begin{proof}
Retain \(c:A\to M_5\), and choose an actual based nullhomotopy
\[
 H:P\times I\longrightarrow M_5,\qquad
 H(p,0)=cq_Y(p),\quad H(p,1)=*.
\]
The source has its double-cylinder CW structure.  Since \(M_5\)
is four-connected, choose \(H\) cellular relative to its ends;
it is then constant on the four-skeleton, as is \(c\).
Put \(\mathcal U=\operatorname{Map}_*(M_5,BSU(2))_0\), and
represent \(g\) by a loop of maps \(g(p,t)\) with both endpoints
the constant map.  There is an actual continuous map
\[
 \Phi_g:\mathcal U\times\mathcal U
          \longrightarrow\operatorname{Map}(F_7,BSU(2))_g.
\]
It assigns \(f_-c\) and \(f_+c\) to the two ends, and assigns
to the three consecutive thirds of the cylinder
\[
 f_-H(p,3t),\qquad g(p,3t-1),\qquad f_+H(p,3-3t),
\]
respectively.  These formulas agree at every seam.  At the basepoint
pair this is merely the chosen clutch loop with constant collars.

The already computed group
\(\pi_1\mathcal U=[M_5,S^3]=C_2^2\) maps onto the invisible
endpoint \(C_2\subset[A,S^3]\).  Choose an order-two element
mapping to its generator.  Apply it in the first and second source
factors of \(\Phi_g\).  Because this is a product map, its two
images \(e_-,e_+\) commute and have order two.  Their endpoint
projections are the two separate generators, so they give an
embedded lift of the endpoint \(C_2^2\).

A loop family in either source factor is constant on \(F_7^{(4)}\)
by the cellular choice of \(c,H\); the middle clutch loop stays
fixed throughout that family.  Under the relative correspondence
between loops in the classifying mapping space and components of
the gauge group, \(e_\pm\) therefore have representatives equal
to the identity on \(F_7^{(4)}\).  They lie in filtration \(F^5\)
of the adjoint group-bundle sections.  Every \(SU(2)\) gauge
class lies in \(F^3\), since its fibre is two-connected.  The
preceding lemma and \(\dim F_7=7\) imply that \(e_\pm\)
commute with every gauge class.

The exact sequence of
Proposition~\ref{p1b:v215:prop:F7-gauge-extension} already makes its
kernel \(C_2^3\) central.  Lift one generator of its infinite
cyclic quotient.  Together with this kernel and \(e_\pm\), it
generates the entire group, and all these generators commute.
The five order-two generators are independent, by the kernel
injection and the independent endpoint projections.  The infinite
generator has no relation with them, by its nonzero infinite
projection.  This proves the stated abstract splitting.
\end{proof}

\subparagraph{Precisely reduced descent frontier.}
The central extension class is consequently no longer a missing input.
Moreover, for a joining path \(\kappa:x\to Rx\), the automorphism
\(\Theta_\kappa\) of this abelian group is independent of that
path, because changing it acts by inner conjugation.  The identity
\(\Theta_\kappa^2=\operatorname{Ad}_{\kappa R\kappa}\)
now gives \(\Theta_\kappa^2=1\).  This does not provide an
origin for the twisted loops: changing \(\kappa\) still translates
their coordinates.  Their set in this component is an affine copy of
\[
 \operatorname{coker}(\Theta_\kappa-1:
      \mathbb Z\oplus C_2^5\longrightarrow\mathbb Z\oplus C_2^5).
\]
The actual involution on these generators and the ordered-cover
pullback of these affine classes still require evaluation before
individual numerical \(\operatorname{Desc}(E)\) fibres can be claimed.

\begin{proposition}[The specified residual lift has the trivial quotient origin]
\label{p1b:v215:prop:carrier-equivariant-origin}
For the residual doublet of
Proposition~\ref{p1b:v215:prop:carrier-selection}, its pulled-back
\(+I\) lift is equivariantly isomorphic to the canonical exchange
lift on the trivial bundle.  Its quotient is the trivial principal
bundle on \(Q\).
\end{proposition}
\begin{proof}
In the section \(s(r)\) the fibre cocycle is
\(b(a,b)=ab^{-1}\).  The global determinant-one gauge
\(u(a,b)=b^{-1}\) changes it to
\[
 u(b,a)b(a,b)u(a,b)^{-1}
       =a^{-1}(ab^{-1})b=I.
\]
The relative scalar coordinate is carried to its inverse and does
not enter this gauge.  Equivalently the pulled-back principal bundle
has the tautological equivariant section \(t(a,b)=(a,b)\), which
differs from \(s(ab^{-1})\) by the vertical factor \(b\).
This proves the actual quotient origin, without a count or uniqueness
assertion about the other quotient bundles pulling back trivially.
\end{proof}


\subsection{Hurwitz-cell quadratic frame and comparison principal angles}
\label{p1b:sec:v221-hurwitz-comparison-frame}
The following is a geometric statement about a finite evaluation frame.
It creates neither extra propagating scalar degrees of freedom nor a new
fermion-source owner.  Write $I=2I$, $H=2T$ and choose five right coset
representatives $g_i$ for $I/H$.  The twenty-five ordered cells
$\mathcal C_{ij}=g_iHg_j^{-1}$ contain twenty-four vertices each and are
isometric images of the Hurwitz spherical $5$-design.  Distinct cells can
overlap; the cell index below is an \emph{orthogonal evaluation-register}
label, not a declaration of twenty-five independent fields.

\begin{theorem}[Exact quadratic Hurwitz evaluation frame]
\label{p1b:thm:v221-hurwitz-quadratic-frame}
Let $\mathcal K_{14}$ be the real round-$S^3$ harmonic band of degrees
$0,1,2$, of dimensions $1+4+9=14$, with the normalized Haar $L^2$
metric.  For every ordered cell $r=(i,j)$ define
\[
 E_r:\mathcal K_{14}\longrightarrow\mathbb R^{24},\qquad
 (E_rf)(q)=24^{-1/2}f(q),\quad q\in\mathcal C_r.
\]
Then $E_r^*E_r=I_{14}$ for each $r$.  With the isometric
multi-cell evaluation
\[
 F:=\frac15\begin{pmatrix}E_1\\[-1mm]\vdots\\E_{25}\end{pmatrix}
 :\mathcal K_{14}\longrightarrow\mathbb R^{600},
\qquad F^*F=I_{14},
\]
let $Q_{r_0}$ be the orthogonal projector deleting one \emph{entire}
$24$-component rooted-cell register.  Then
\[
 \boxed{F^*Q_{r_0}F=\eta_H I_{14},\qquad
 \eta_H=1-\frac1{25}=\frac{24}{25}.}
\]
\end{theorem}
\begin{proof}
For harmonics $f,g$ of degree at most two, $\overline f g$ is a
polynomial restricted to $S^3$ of total degree at most four.  The
spherical $5$-design rule on each rotated Hurwitz cell gives
$\langle E_rf,E_rg\rangle=\langle f,g\rangle_{S^3}$.
Consequently each of the twenty-five blocks contributes $I_{14}$
to the unscaled Gram, giving $F^*F=(25/25)I_{14}$.
Removing one block leaves $(24/25)I_{14}$.  This uses the
quadratic band, not the interacting quartic Hessian at a general
background, whose integrands can have degree eight.
\end{proof}

\begin{corollary}[Canonical two-frame principal angle]
\label{p1b:cor:v221-hurwitz-principal-angle}
Put $\eta_H=24/25$, $t=\sqrt{\eta_H}=2\sqrt6/5$ and
\[
 J_0=F,\qquad J_1=\eta_H^{-1/2}Q_{r_0}F.
\]
Both maps are isometries and
\[
 \boxed{J_0^*J_1=tI_{14},\qquad
 \cos\theta=t,\quad\sin\theta=\frac15.}
\]
Every one of the fourteen principal angles between their images is
identical.  Moreover, for every isometric complex-linear inclusion
$U:G\hookrightarrow\mathcal K_{14}\otimes\mathbb C$ of a Hermitian
three-dimensional space $G$, the maps $j_a=(J_a\otimes1)U$ satisfy
$j_0^*j_1=tI_G$.  Independence of the chosen $U$ is a numerical
identity, not a proof that a physical source carries a particular $U$.
\end{corollary}
\begin{proof}
Since $Q_{r_0}$ is an orthogonal projector,
$J_1^*J_1=\eta_H^{-1}F^*Q_{r_0}F=I_{14}$ and
$J_0^*J_1=\eta_H^{-1/2}F^*Q_{r_0}F=tI_{14}$.
The principal-angle identity and the isometric trimer pullback follow.
\end{proof}

\begin{proposition}[Exterior comparison and oriented-volume character]
\label{p1b:prop:v221-exterior-comparison-character}
For the three-dimensional $G$ of the preceding corollary, put
$L=j_0^*j_1=tI_G$.  The determinant inner products on its exterior
powers give
\[
 \boxed{(\Lambda^rj_0)^*(\Lambda^rj_1)
 =\Lambda^rL=t^r I_{\Lambda^rG},\quad r=1,2,3.}
\]
Thus the available comparison amplitudes include $t$, $t^2$,
$t^{-2}$ on the reciprocal two-area sort, and $t^{-3}$ on the
reciprocal oriented-volume sort.  With the standard three-dimensional
Hodge identification,
\[
 (\Lambda^2L)^{-1}=(\det L)^{-1}L^{T},\qquad
 t^{-2}=t^{-3}t.
\]
The source/target, owner selection and physical readout of these
algebraically available sorts are \emph{not} consequences of this
proposition.  In particular neither $\Lambda^3G$ nor its dual is
identified here with a charged particle.
\end{proposition}
\begin{proof}
The cross inner product of wedges is the determinant of the matrix of
pairwise cross inner products.  Every such matrix is $tI_r$, so its
determinant is $t^r$.  The inverse and Hodge statements are the usual
nondegenerate three-dimensional determinant identities.
\end{proof}

\begin{remark}[Quadratic versus nonlinear comparison]
\label{p1b:rem:v221-quadratic-vs-quartic-comparison}
The scalar frame theorem concerns the quadratic Haar pairing.  It is
not a universal $24/25$ multiplier on nonlinear vertices, the
full interacting Hessian, or the physical fermion mass operator.
The rank-twenty-four register block deletion differs from the
rank-one exclusion in the abstract twenty-five-label carrier.  These
are distinct but related realizations of the same numerical ratio.
\end{remark}

\subsection{Explicit family embedding and its finite nonlinear comparison}
\label{p1b:sec:v242-family-quadratic-embedding}
\begin{theorem}[Unital quadratic-band embedding from the operator trimer]
\label{p1b:thm:v242-family-quadratic-embedding}
Let \(E=u^\perp\subset\operatorname{Im}\mathbb H\), with
\(u=(1,1,1)/\sqrt3\), and keep the rooted rays
\[
 r_0=(-2,1,1)/\sqrt6,\quad r_1=(1,-2,1)/\sqrt6,\quad
 r_2=(1,1,-2)/\sqrt6.
\]
Let \(\mathcal U:G\to\mathcal S(E)_{\mathbb C}\) be the complexified
isometry of Theorem~\ref{p1b:thm:v137-q5-operator-trimer}.
For \(A=\mathcal Ua\), set
\(A_0=A-\frac12\Tr(A)P_E\), where \(P_E\) projects onto \(E\).
Then
\[
 \boxed{(\mathcal Ja)(q)=\langle c,a\rangle
       +\sqrt{12}(\operatorname{Im}q)^TA_0(\operatorname{Im}q),
       \qquad c=(1,1,1)^T/\sqrt3,}
\]
defines a complex-linear isometry
\(\mathcal J:G\hookrightarrow\mathcal H_0\oplus\mathcal H_2
\subset\mathcal K_{14}\otimes\mathbb C\), with \(\mathcal Jc=1\).
The three basis images are \(f_j=1/\sqrt3+b_j\), where
\[
 b_j(q)=4\langle\operatorname{Im}q,r_j\rangle^2
             -2\|P_E\operatorname{Im}q\|^2,
 \qquad \sum_jb_j=0.
\]
On each of the twenty-five Hurwitz cells,
\[
 \langle b_j\rangle=0,\qquad
 \langle b_jb_k\rangle=\delta_{jk}-1/3,\qquad
 \langle f_jf_k\rangle=\delta_{jk}.
\]
\end{theorem}
\begin{proof}
The operator trimer carries \(c\) to \(P_E/\sqrt2\) and
\(G_\perp\) isometrically to trace-free symmetric operators on \(E\).
For trace-free complex symmetric \(A_0\), the normalized round
\(S^3\) fourth moment gives
\[
 \int_{S^3}|q^TA_0q|^2\,d\sigma
       =\frac1{12}\Tr(A_0^\dagger A_0).
\]
Its mean is zero. The constant and trace-free terms are orthogonal,
proving isometry and unitality. The Gram normalization of the operator
trimer gives the displayed \(f_j\) by substitution. Their first and
second products have degree at most four, so the Hurwitz quadratic-frame
theorem proves the same identities on every cell. The expanded moment
calculation is given in
Proposition~\ref{p1c:prop:v242-native-family-cubic}.
\end{proof}

\begin{corollary}[Exact cubic comparison for the embedded augmentation modes]
\label{p1b:cor:v242-family-cubic-comparison}
For \(a\in G_\perp\), the rooted Hurwitz and full icosian averages obey
\[
 \boxed{\langle(\mathcal Ja)^3\rangle_H=6a_0a_1a_2,
 \qquad \langle(\mathcal Ja)^3\rangle_{2I}=0.}
\]
With the original \(1/25\) action/register normalization, exclusion of
the root Hurwitz register therefore gives
\[
 \boxed{C_{\rm raw}=-C_H/25,\qquad Z_{\rm raw}=(24/25)I_G.}
\]
For the normalized average over retained registers the cubic is instead
\(C_{24}=-C_H/24\). These identities extend by polarization to the
corresponding trilinear tensors on \(G_\perp\).
\end{corollary}
\begin{proof}
Use the finite sixth-moment computation and fivefold vertex multiplicity
in Proposition~\ref{p1c:prop:v242-native-family-cubic}.
Every register has the same quadratic Gram, while the full cubic average
on these modes vanishes. Removing one register subtracts its complete
cubic with the declared normalization.
\end{proof}

\begin{remark}[Coefficient modes and physical field profiles]
\label{p1b:rem:v242-family-band-scope}
This supplies an explicit embedding of the family coefficient space, not a
new identification with physical fermion spinor wavefunctions or an internal
Higgs profile. Three canonically normalized copies of these scalar modes
give \(C_{\rm can}=(24/25)^{-3/2}C_{\rm raw}\), with the raw vertex
retained. This computation does not assign that factor to an electron mass
or license replacing a finite cubic by its Haar average. A physical
source or nonlinear action vertex must inherit its actual field embedding,
chiral contraction, kinetic pairing and comparison frames.
\end{remark}

\subsection{One rooted icosian comparison: finite moments and harmonic boundary carriers}
\label{p1b:sec:v243-one-root-geometric-closures}
The statements in this subsection concern the already defined finite
comparison registers and round scalar harmonic carriers.  The embedded
point sets are geometric evaluation supports; the operations below do
not, by themselves, constitute a physical source, action or state.
Keep $I=2I$, $H=2T<I$, $X=I/H$ and the $25$ ordered cell registers
$\mathcal C_{ij}=g_iHg_j^{-1}$ of
Theorem~\ref{p1b:thm:v221-hurwitz-quadratic-frame}.  For a chosen
rooted coset choose $g_0=1$, so the deleted diagonal register is
$\mathcal C_{00}=H$.

\begin{theorem}[One rooted register functional on all finite observables]
\label{p1b:thm:v243-rooted-register-functional}
For any function $F:I\to\mathbb C$, write
$\tau_I(F)=120^{-1}\sum_{q\in I}F(q)$ and
$\tau_H(F)=24^{-1}\sum_{q\in H}F(q)$.  The unnormalized average
of the original $25$ registers, with the entire root register deleted,
is precisely
\begin{equation}\label{p1b:eq:v243-rooted-register-functional}
 \boxed{\tau_{\rm root}(F)=\frac1{600}
 \sum_{(i,j)\ne(0,0)}\ \sum_{q\in\mathcal C_{ij}}F(q)
 =\tau_I(F)-\frac1{25}\tau_H(F).}
\end{equation}
In particular $\tau_{\rm root}(1)=24/25$.  On the labelled Hilbert
carrier $\ell^2(X\times X)\otimes\mathbb C^{24}$, deletion is the
operator $Q_{\rm rel}\otimes I_{24}$, where
$Q_{\rm rel}=I-|x_0,x_0\rangle\langle x_0,x_0|$.  This is a
block realization of the relative-label operator, not an equality of
operators on two different underlying carriers.  A retained-register
probability average instead divides $\tau_{\rm root}$ by $24/25$.
\end{theorem}
\begin{proof}
For fixed $j$, the five left cosets $g_iH$, translated by
$g_j^{-1}$, partition $I$.  Hence every $q\in I$ is counted once
for each of the five values of $j$ and five times in all $25$
registers.  Subtracting the single rooted $H$ register gives the
identity.  The Hilbert-carrier statement follows by the orthogonal
labelled direct-sum realization, with the internal order of each
$24$-point cell immaterial.
\end{proof}

\begin{theorem}[Nested Fourier projectors and two distinct harmonic onsets]
\label{p1b:thm:v243-nested-reynolds-onsets}
For the spin-$n/2$ representation $V_n$ of $SU(2)$ let
$P_I^{(n)}$ and $P_H^{(n)}$ be its finite-group Reynolds projectors,
and let $\mu_n^I,\mu_n^H$ denote their ranks.  On matrix
coefficients, the single functional
\eqref{p1b:eq:v243-rooted-register-functional} is represented by
\[
 \boxed{T_n=P_I^{(n)}-\frac1{25}P_H^{(n)}.}
\]
Its eigenvalues and multiplicities are respectively
\[
 \begin{array}{c|ccc}
  \text{eigenvalue}&24/25&-1/25&0\\\hline
  \text{multiplicity}&\mu_n^I&\mu_n^H-\mu_n^I&(n+1)-\mu_n^H.
 \end{array}
\]
For $1\le n\le5$ it vanishes.  At $n=6$ and $n=8$ it has one
$-1/25$ channel and no $+24/25$ channel.  At $n=12$ it has one
channel of each nonzero sign.  In contrast, the first nonconstant
\emph{unrooted} $I$-invariant scalar harmonic occurs at degree $12$.
The rooted first response at degree $6$ must not replace the
preterminal unrooted degree-$12$ output.
\end{theorem}
\begin{proof}
Since $H\subset I$, one has $P_HP_I=P_IP_H=P_I$.  Decompose
$V_n=\operatorname{ran}P_I\oplus(\operatorname{ran}P_H\ominus
\operatorname{ran}P_I)\oplus\ker P_H$ and evaluate $T_n$ on its
three summands.  The ranks through $n=12$ are the exact group-character
sums of Proposition~\ref{p1c:prop:v243-reynolds-character-replay}.
\end{proof}

\begin{theorem}[Exact four-atom Hurwitz law and every rooted family moment]
\label{p1b:thm:v243-four-atom-family-law}
Let $a=(a_0,a_1,a_2)$ be real, $\sum_ja_j=0$, and let
$f_a=\mathcal Ja$ be the \emph{actual} quadratic embedding of
Theorem~\ref{p1b:thm:v242-family-quadratic-embedding}.
Its pushforward of the uniform $H$ counting measure is
\[
 \boxed{(f_a)_*\mu_H=\tfrac14(\delta_0+\delta_{2a_0}
 +\delta_{2a_1}+\delta_{2a_2}).}
\]
Consequently $\tau_H(f_a^k)=2^{k-2}\sum_ja_j^k$ for $k\ge1$.
Writing $r^2=\sum_ja_j^2$ and $c=a_0a_1a_2$, the first moments of
$\tau_{\rm root}=\tau_I-\tau_H/25$ are
\[
 \begin{array}{c|rrrrrr}
 k&0&1&2&3&4&5\\\hline
 \tau_{\rm root}(f_a^k)&24/25&0&24r^2/25&-6c/25&
 131r^4/50&-4r^2c/5.
 \end{array}
\]
In particular the $I$ and $H$ fourth moments are respectively
$27r^4/10$ and $2r^4$.  For the raw rooted quadratic metric
$Z=24I/25$, canonical $a_{\rm can}=\sqrt{24/25}\,a$ gives the
rooted quartic coefficient $3275\|a_{\rm can}\|^4/1152$.
For arbitrary functions $\mathcal F$ of these finite profile values,
\[
 \boxed{\tau_{\rm root}(\mathcal F(f_a))=
 \frac1{120}\sum_{q\in I}\mathcal F(f_a(q))-
 \frac1{100}\bigl[\mathcal F(0)+\sum_j\mathcal F(2a_j)\bigr].}
\]
The same marked functional determines all its nonlinear coefficients;
it does not impose a universal $24/25$ multiplier on distinct vertices.
\end{theorem}
\begin{proof}
The four-atom enumeration, fourth-moment calculation and Newton
identities on $\sum_ja_j=0$ are given in
Proposition~\ref{p1c:prop:v243-four-atom-moment-proof}.  Apply
Theorem~\ref{p1b:thm:v243-rooted-register-functional} and rescale by
the displayed quadratic metric.
\end{proof}

\begin{theorem}[Canonical first icosian harmonic from the sextic defect]
\label{p1b:thm:v243-icosian-h12-origin}
For $a=(1,-1,0)$ the same embedded quadratic satisfies the exact
sampling difference
\[
 \boxed{\tau_I(f_a^6)-\int_{S^3}f_a^6\,d\sigma
 =\frac{1241}{16}-\frac{540}{7}=\frac{47}{112}\ne0.}
\]
The $I$ counting functional on the complex degree-$12$ scalar harmonic
space $\mathcal H_{12}$ has a unique Riesz representative
$h_{12}$ (normalized Haar inner product), explicitly
\[
 h_{12}(q)=\frac{13}{120}\sum_{g\in I}\chi_{12}(g^{-1}q),
 \quad \boxed{\|h_{12}\|^2=13,\qquad h_{12}(g)=13\ (g\in I)}.
\]
It is nonzero and its $SU(2)_L\times SU(2)_R$ \emph{mathematical}
orbit spans $\mathcal H_{12}$, of dimension $169$.  The orbit-span
statement does not license arbitrary orbit evaluations as physical
finite-carrier constructors.
\end{theorem}
\begin{proof}
The $I$ design is exact through polynomial degree $11$, while
$f_a$ is degree two; hence the first possible nonzero sampling defect
of its powers is the degree-$12$ contribution to $f_a^6$.  The
independent exact sixth-moment and reproducing-kernel calculations
are in Proposition~\ref{p1c:prop:v243-icosian-sextic-riesz-proof}.
Finally $\mathcal H_{12}\simeq V_{12}\otimes V_{12}^*$ is
irreducible under the product of the full left and right round group
actions; a nonzero orbit therefore spans it.
\end{proof}

\begin{theorem}[Complete rooted boundary augmentation from degree two plus degree one]
\label{p1b:thm:v243-neutral-five-isometry}
Let $\Omega$ be the six Hurwitz half-units of
Section~\ref{p1b:sec:v2208-neutral-boundary-geometry}, chosen on
the rooted diagonal line, and let
$\operatorname{ev}_\Omega$ be pointwise scalar evaluation and
$E_\Omega=\operatorname{ev}_\Omega|_{\mathcal H_{12}}$ its restriction.
Retain the canonical positive coisometry
$B=(E_\Omega E_\Omega^\dagger)^{-1/2}E_\Omega$ of
Theorem~\ref{p1b:thm:v2208-boundary-coisometry}, with
$E_5=B^\dagger\mathbf1^\perp$.  The intrinsic opposite-pair
parity splits $\mathbf1^\perp=W_{2,+}\oplus W_{3,-}$ with
dimensions $2$ and $3$.  On the calculation representative of the
rooted line $u=(1,1,1)/\sqrt3$, put
\[
 V(w)(h_\varepsilon)=w\cdot\operatorname{Im}h_\varepsilon,
 \qquad M=V^\dagger V=2I_3-\tfrac12J_3\succ0.
\]
Then the explicit map
\begin{equation}\label{p1b:eq:v243-five-augmentation-isometry}
 \boxed{\mathcal I_5(a,w)=B^\dagger\left(
 \frac{\operatorname{ev}_\Omega(\mathcal Ja)}{\sqrt8}+VM^{-1/2}w\right):
 G_\perp\oplus\mathbb C^3\longrightarrow E_5}
\end{equation}
is a surjective isometry.  Here
$M^{-1/2}=2^{-1/2}P_{u^\perp}+\sqrt2 P_u$.
For the $h_{12}$ of the preceding theorem,
\[
 \boxed{P_{\rm aug}h_{12}=0,\quad
 \|P_{\rm sing}h_{12}\|^2=13/3,\quad
 \|Q_{163}h_{12}\|^2=26/3.}
\]
Thus the unrooted group-average harmonic alone cannot source the
augmentation; the full unsummed six-port incidence carries the
missing geometric data.  Neither statement is a physical K56 mass
operator or neutral state probability.
\end{theorem}
\begin{proof}
The six points form three opposite-imaginary pairs.  The family
values on the pairs are $2a_0,2a_1,2a_2$, yielding
$\|\operatorname{ev}_\Omega(\mathcal Ja)\|^2=8\|a\|^2$, whereas the imaginary
linear coordinates have Gram $M$ and are odd under pair exchange.
These two mutually orthogonal isometries exhaust the mean-zero
boundary image.  Apply $BB^\dagger=I_6$ and the exact projection
of $h_{12}$ calculated in
Proposition~\ref{p1c:prop:v243-neutral-boundary-source-proof}.
\end{proof}

\begin{theorem}[Finite uniform quartic-exact subgroup discriminator]
\label{p1b:thm:v243-quartic-exact-subgroup}
Let $\Gamma<SU(2)$ be a finite subgroup with its uniform counting
state, and let $\mathcal P_2(S^3)$ be the \emph{complete} real scalar
polynomial band of degrees at most two.  If
\[
 \tau_\Gamma(f^4)=\int_{S^3}f^4d\sigma
 \quad\text{for every }f\in\mathcal P_2(S^3),
\]
then $\Gamma$ is conjugate to $2I$.  Conversely $2I$ satisfies this
identity.  In particular a physically occurring \emph{finite,
subgroup-closed, uniform and quartic-exact} comparison role would
force the $2I$ group class, but neither occurrence nor that full
comparison requirement follows from this classification theorem.
\end{theorem}
\begin{proof}
By polarization with the constant function, exactness of all fourth
powers implies exactness of the real polynomial band through degree
$8$.  Finite subgroups of $SU(2)$ are cyclic, binary dihedral,
$2T$, $2O$ or $2I$.  Cyclic groups have a transverse degree-one
coordinate vanishing on them; binary dihedral groups have a nonzero
quadratic transverse product vanishing on both of their great circles.
For $2T$ and $2O$ use $f(q)=q_0^2-q_1^2$; their fourth moments are
$1/6$ and $5/48$ respectively, whereas the round moment is $3/40$.
The $2I$ counting state agrees with Haar through degree $11$.
The exact arithmetic and the hypotheses of the subgroup exclusion
are supplied by Proposition~\ref{p1c:prop:v243-reynolds-character-replay}.
\end{proof}

\begin{remark}[Geometry completed; physical action insertion remains typed]
\label{p1b:rem:v243-one-root-scope}
Theorems~\ref{p1b:thm:v243-rooted-register-functional}--
\ref{p1b:thm:v243-quartic-exact-subgroup} close geometric register,
family-moment and neutral-boundary consequences from one root.  The
comparison used by a physical action must still be an occurring
licensed constructor with an actual retained input/output and source
embedding.  The preterminal unrooted degree-$12$ output and the
rooted degree-$6$ sibling are not interchangeable.  The separately
open individual principal-$SU(2)$ descent fibres and the mathematical
validity of the preceding theorems are unaffected.
\end{remark}


\subsection{Principal descent: established maps and unevaluated fibres}
\label{p1b:sec:v214-mathematical-status}
The Hodge identification, reciprocal index and gap, variable protected
kernel, finite response calculus and the positive appendix results used by
the main proofs are proved in this manuscript from the displayed mathematical data. Their proofs do
not import results from later papers in the series. The explicit triality
restriction uses the signed coefficient tensor and spectator extraction of
Paper~I-B, Theorem~\ref{p1:thm:v214-normalized-spin8-triality}, as calculated
in Corollary~\ref{p1b:cor:v126-triality-restriction-is-clifford}.

The numerical mixed-$U(2)$ classification and its retained natural tests
are evaluated in Theorems~\ref{p1b:thm:complete-mixed-rank-two} and
\ref{p1b:thm:natural-rank-two-completion}. Ordinary determinant-preserving
Born-fixed classes use only the single-map intertwining of
Proposition~\ref{p1b:prop:ordinary-born-component-translation}; they do
not supply coherent principal lifts. The square obstruction
of Proposition~\ref{p1b:prop:exotic-born-square-obstruction} remains nonzero.

The principal-$SU(2)$ descent reduction of
Theorem~\ref{p1b:thm:g21-principal-descent} is exact, but its individual
numerical fibres are not evaluated. Theorems~\ref{p1b:thm:principal-descent-mapping-torus}
and~\ref{p1b:cor:principal-descent-source-cohomology} give its actual
mapping-torus source and mapping-space pullback. Lemma~\ref{p1b:lem:pure-descent-symmetric-loop-image}
computes the symmetric-sphere loop image to a unique actual-map
normalization $z_0\in\{2,8\}$ in the fixed Heisenberg chart. Choosing
between those candidates still requires that actual map; the full
twisted-loop gauge action is also retained in the descent fibre.
The zero-Chern ordinary image is evaluated in Theorem~\ref{p1b:thm:pure-descent-zero-chern-image}: exactly thirty-six classes admit pure lifts, each with at least the two fixed-fibre signs. Their individual descent fibres remain unevaluated. Every pure complex lift reduces equivariantly to a principal lift by Proposition~\ref{p1b:prop:pure-u2-determinant-reduction}.
The restriction and equivariant-index congruences are necessary tests,
and are not asserted sufficient. Thus the remaining task is precisely to
evaluate $q^*:[\operatorname{SP}^2(S^1\times S^3),BSU(2)]\to[M_8,BSU(2)]$
with its individual fibres. The quotient's cohomology alone does not
evaluate this unstable classifying-map problem.


The additional projection, dilation, torsion-contraction and complex-module
identities in Section~\ref{p1b:sec:v215-geometric-provider-completion}
are mathematical maps on their displayed bundles. The rank-changing
Hermitian reduction, compact-pulse estimate and logarithmic Legendre
calculus likewise have no later-series premise. The principal
stabilization and seven-dimensional double-cylinder calculations of
Sections~\ref{p1b:v215:sec:principal-fibre}
and~\ref{p1b:v215:sec:principal-stabilization} evaluate their stated
intermediate components and maps. They do not evaluate the final
individual fibres of the symmetric-square descent map. In particular
the abstract central extension is split in
Theorem~\ref{p1b:v215:thm:F7-gauge-split}, while Theorem~\ref{p1b:thm:v226-principal-two-types} reduces the monodromy type to one carry bit. Theorem~\ref{p1b:thm:trivial-component-zero-carry} below evaluates
the trivial-component carry using an invariant primitive map.
Other-component carries, endpoint normalization and the
ordered-cover pullback remain separate classification questions.

\subsection{The principal-descent group extension and remaining fibres}
\label{p1b:sec:v226-principal}

This section concerns only the existing principal mapping-space reduction.
Its result is an abstract group of differences for twisted loops, not the
individual fibres of $q^*:[\operatorname{SP}^2(S^1\times S^3),BSU(2)]
\to[M_8,BSU(2)]$.

Retain the double mapping cylinder $F_7$, its reflection/exchange $R$, and
an $R$-invariant bundle component $x$. Paper I-C has already proved
\[
 \Gamma=\pi_1\operatorname{Map}(F_7,BSU(2))_x\cong\mathbb{Z}\oplus C_2^5.
\]
Its exact endpoint sequence is
\begin{equation}\label{p1b:eq:v226-extension}
 0\longrightarrow K=C_2^3\longrightarrow\Gamma
 \longrightarrow k_x\mathbb{Z}\oplus C_2^2\longrightarrow0,
 \qquad k_x=3/\gcd(3,d).
\end{equation}
The two quotient $C_2$ factors are the separate invisible end gauges.
The kernel is $[\Sigma(S^3\times S^3),S^3]=C_2^3$; its two sphere
primaries form a natural $C_2^2$ quotient, with a remaining top $C_2$.
The existing gauge action on the three kernel coordinates changes only
the top entry:
\begin{equation}\label{p1b:eq:v226-shear}
 (r_x,r_p,e)\longmapsto(r_x,r_p,e+b r_x+a r_p).
\end{equation}
These are the actual endpoint and primary maps, not an arbitrary splitting
of an abstract abelian group.

\begin{proposition}[The free loop direction survives monodromy]
\label{p1b:prop:v226-free}
For any path $\kappa:x\to Rx$, the involution
$\Theta_\kappa$ acts as $+1$ on the free quotient of $\Gamma$.
\end{proposition}
\begin{proof}
The free endpoint coordinate in \eqref{p1b:eq:v226-extension} records equal positive
sphere degree at both ends. The geometric involution exchanges those ends,
so sends $(r,r)$ to $(r,r)$. A joining-path transport conjugates endpoint
loops; their degree homomorphism to $\mathbb{Z}$ is unchanged by conjugation.
Thus the free quotient action is $+1$. There is no reversal of the
parameter of a gauge loop at an endpoint: the cylinder reflection reverses
the gluing path, not that loop parameter.
\end{proof}

\begin{theorem}[Only two abstract twisted-loop difference groups]
\label{p1b:thm:v226-principal-two-types}
In each $R$-invariant component the affine twisted-loop set has difference
group of one of the two types
\begin{equation}\label{p1b:eq:v226-two-types}
 \boxed{\operatorname{coker}(\Theta_\kappa-I)
       \cong\mathbb{Z}\oplus C_2^3\quad\text{or}\quad\mathbb{Z}\oplus C_2^2.}
\end{equation}
The remaining abstract-type distinction is one bit.  This generic
reduction does not assign it.  The trivial component is evaluated
separately in Theorem~\ref{p1b:thm:trivial-component-zero-carry};
other-component values and the ordered-cover pullback remain open.
\end{theorem}
\begin{proof}
Let $T=\operatorname{Tor}\Gamma=C_2^5$. The induced action on $T/K=C_2^2$
exchanges the two invisible endpoint gauges. On the two primary quotient
coordinates of $K$ it exchanges the two sphere factors. The cylinder
reversal has no effect on a $C_2$ sign, and joining-path transport may
introduce the shear \eqref{p1b:eq:v226-shear} but leaves those primaries unchanged.
Hence $\Theta-I$ has rank at least one on $K$. Its induced map on $T/K$
also has rank one, so
\[
 \operatorname{rank}_{\mathbb{F}_2}(A-I)\ge2,
 \qquad A=\Theta|_T.
\]
Since $A^2=I$ on a five-dimensional $\mathbb{F}_2$ space,
$(A-I)^2=0$ and its rank is at most two. Its rank is therefore exactly two.
In particular
\[
 \dim\ker(A-I)=3,\qquad \dim\operatorname{im}(A-I)=2.
\]
By Proposition~\ref{p1b:prop:v226-free}, choose any free generator $z$ and write
$\Theta z=z+v$, $v\in T$. The identity $\Theta^2=I$ forces $Av=v$.
Changing $z$ by a torsion vector changes $v$ by an element of
$\operatorname{im}(A-I)$. Thus precisely one invariant bit remains:
\[
 [v]\in\ker(A-I)/\operatorname{im}(A-I)\cong C_2.
\]
In coordinates $(n,t)\in\mathbb{Z}\oplus T$,
$(\Theta-I)(n,t)=(0,nv+(A-I)t)$. If $[v]=0$, its image has dimension two;
otherwise it has dimension three. Taking its cokernel proves
\eqref{p1b:eq:v226-two-types}. No choice of zero in the affine twisted-loop set
has been made.
\end{proof}

\begin{remark}[What still has to be calculated]
This generic reduction does not assign the actual carry bits.
Theorem~\ref{p1b:thm:trivial-component-zero-carry} proves zero
in the trivial component by an additional invariant-map witness;
the other components remain open. The actual endpoint generator
$(4,4;z_0)$, $z_0\in\{2,8\}$ in the fixed Heisenberg chart, still needs
its specified geometric normalization. The actual affine origin/transport
and ordered-cover pullback are still required to enumerate descent fibres.
The surviving $\mathbb{Z}$ direction ranges over additional Chern data on the
mapping torus; it is not evidence of infinitely many lifts of a fixed
zero-Chern bundle. In particular the established thirty-six-element pure
ordinary image and the two fixed-isotropy signs are unchanged.
\end{remark}

\subsection{Complete native coefficient tests and source derivatives}
\label{p1b:sec:v226-finite-coefficients}
The following finite linear-algebra provider consumes the actual native
block algebra of Paper I-B; it has no premise from a later dynamical paper.

\begin{theorem}[Explicit dual coefficients for the whole native algebra]
\label{p1b:thm:v226-dual-coefficients}
Use the 137 effects $E_j=(I+\eta B_j)/137$, $\eta=1/2$, of
Theorem~\ref{p1:thm:v226-tomography}, and put
\[
 D_j=I/15+(136/\eta)B_j,\qquad
 q_j(A)=\operatorname{Tr}(E_jA).
\]
For every complex $A\in\mathcal A_D$,
\[
 \operatorname{Tr}(E_iD_j)=\delta_{ij},\qquad
 \boxed{A=\sum_{j=1}^{137}q_j(A)D_j.}
\]
For Hermitian $A$, these coefficients obey
\[
 \|A\|\le\|A\|_{\rm HS}
 \le\frac{\sqrt{137\cdot136}}{\eta}
             \left(\sum_j|q_j(A)|^2\right)^{1/2}.
\]
Consequently equality of native matrices, including any specified native
response or source residual, is equivalent to 137 scalar coefficient
identities. This is reconstruction, not a positive inverse channel.
\end{theorem}
\begin{proof}
The simplex pairing is $\operatorname{Tr}(B_iB_j)=1$ for $i=j$ and
$-1/136$ otherwise. Substitution gives the biorthogonality and hence the
inverse on the 137-dimensional algebra. On $I/\sqrt{15}$ the effect-frame
eigenvalue is $15/137$, and on the traceless space it is
$\eta^2/(137\cdot136)$. The latter is the smaller, proving the norm bound.
Complex linearity extends the inverse to every complex matrix in the algebra.
\end{proof}

\begin{corollary}[Fixed-frame source jets and operator-valued coefficients]
\label{p1b:cor:v226-source-jets}
If $A(s)$ is $C^k$ in a fixed native frame, every multi-index $\nu$ with
$|\nu|\le k$ satisfies
\[
 \partial^\nu A=\sum_j\partial^\nu q_j(A)D_j.
\]
For a jointly transported unitary frame $U(s)$ use
$\nabla_\mu A=\partial_\mu A-[\Theta_\mu,A]$,
$\Theta_\mu=(\partial_\mu U)U^*$; transport effects and duals by the same
$U$. The same coefficient identities hold with covariant derivatives.
For any finite outer carrier $\mathcal K$ and
$A\in\operatorname{End}(\mathcal K)\otimes\mathcal A_D$ they hold with
operator coefficients
\[
 q_j(A)=(\operatorname{id}\otimes\operatorname{Tr}_{\mathcal H})
          [(I\otimes E_j)A],\quad
 A=\sum_jq_j(A)\otimes D_j,
 \quad \|A\|\le\sum_j\|q_j(A)\|\,\|D_j\|.
\]
\end{corollary}
\begin{proof}
Differentiate the fixed finite linear inverse. For moving unitary frames,
$\nabla_mu(UAU^*)=U(\partial_\mu A)U^*$, so the fixed identities transport
exactly. The tensor assertion follows on elementary tensors by the inverse
and then by linearity; the norm bound is the triangle inequality.
Nonreducing isometries still retain their omitted-image terms.
\end{proof}

\begin{remark}[An exact test does not choose its input]
The theorem evaluates equality once the two actual matrix words are supplied.
It supplies neither an action matrix, a boundary datum, a gauge coupling nor
its physical normalization. A common carrier or a matching rank does not
force the 137 coefficients to agree. Derivative claims consume the actual
jets, including the complete moving-frame and source terms.
\end{remark}


\subsection{Finite response comparison from complete sampled data}
\label{p1b:sec:v227-response-comparison}
The complete Hermitian response and minimal visible realization already
proved in Theorems~\ref{p1b:thm:v2210-lit-finite-weyl} and
\ref{p1b:thm:response-minimal-visible-realization} give the following finite
comparison certificate. The energy variable below is the common normalized
spectral parameter, not a curvature or hierarchy coordinate.

\begin{theorem}[Finite upper-half-plane certificate for the whole response]
\label{p1b:thm:v227-finite-response}
Let two actual Hermitian finite parents have the same retained target and
canonical kinetic normalization, with responses
\[
 M_i(z)=zI-A_i+B_i(C_i-zI)^{-1}B_i^*,\qquad
 A_i=A_i^*,\quad C_i=C_i^*,\quad\dim C_i\le m_i.
\]
Put $N=m_1+m_2$. Equality at
$s=\lfloor N/2\rfloor+1$ distinct points in the open upper half-plane
is necessary and sufficient for $M_1=M_2$ as matrix-valued meromorphic
functions. If $m_1,m_2\le m$, $m+1$ such points suffice.
Alternatively $N+1$ distinct regular real points suffice.
If $A_1=A_2$ is independently known, the upper-half-plane count improves
to $\lceil N/2\rceil$ (with zero tests needed when $N=0$).

If both responses belong to the native algebra $\mathcal A_D$,
full matrix equality at each point is equivalent to equality of all 137
coefficients $\operatorname{Tr}(E_jM_i(z))$ of
Theorem~\ref{p1b:thm:v226-dual-coefficients}. A larger retained/source
target consumes its own complete test frame, rather than a universal 137 count.
\end{theorem}
\begin{proof}
Let $d_i(z)=\det(C_i-zI)$, with the determinant one for a zero-dimensional
environment. Then
$P(z)=d_1(z)d_2(z)[M_1(z)-M_2(z)]$ is a matrix polynomial of degree
at most $N$, since the common $zI$ term cancels. The determinants have
real coefficients and $M_i(\bar z)^*=M_i(z)$, so
$P(\bar z)^*=P(z)$. A matrix zero at each upper-half-plane sample gives
another matrix zero at its conjugate. Thus $2s>N$ distinct zeros force
every entry polynomial to vanish. There are no complementary poles at
these points because $C_i$ are Hermitian. The real-node assertion follows
from the usual $N+1$ polynomial-zero count. If the direct blocks agree,
the numerator degree is at most $N-1$, giving the improved count.
The converse is immediate. The coefficient version follows from the
exact complex-linear native inverse.
\end{proof}

\begin{proposition}[The count can be necessary even for positive parents]
\label{p1b:prop:v227-sharp-sample-count}
For $m_1=m_2=2$, equality at two upper-half-plane points need not imply
response equality. Let
\[
\begin{split}
 M_1(w)&=w+\frac{13/7}{-3-w}+\frac{4/3}{2-w},\\
 M_2(w)&=w-1+\frac{1/3}{-1-w}+\frac{34/7}{4-w}.
\end{split}
\]
Then $M_1=M_2$ at $w=i,2i$, but not identically. Both are Hermitian-parent
responses with positive coefficient residues. A common energy shift by ten
makes both parent matrices positive definite while preserving the example
at $z=10+i,10+2i$.
\end{proposition}
\begin{proof}
Their difference is exactly
\[
 \frac{(w^2+1)(w^2+4)}{(-3-w)(-1-w)(2-w)(4-w)}.
\]
The partial-fraction residues are $13/7,-1/3,4/3,-34/7$, with the positive
and negative terms assigned to the respective parents. Use these positive
weights as squared couplings to their two real complementary eigenvalues,
and direct blocks zero and one. Adding ten to the full parents shifts
$w=z-10$; strict positive diagonal dominance then proves positivity.
The displayed nonzero rational function proves both assertions about samples.
\end{proof}

\begin{corollary}[Finite source-jet certificate and visible realization]
\label{p1b:cor:v227-response-jets}
Let the preceding parents be $C^k$ families in real source/metric parameters
$t$, on fixed finite carriers and in one common kinetic/energy frame.
At $t_*$, equality of all consumed source jets through order $k$ at
\[
 s_k=\left\lfloor\frac{(k+1)(m_1+m_2)}2\right\rfloor+1
\]
fixed distinct upper-half-plane energies implies equality of those jets
as meromorphic functions of energy. If the complete compared response
includes all consumed source ports, equality fixes every admitted retained
Green response, pole and residue and the minimal visible environment up
to its unitary equivalence. Unused dark complements may differ.
\end{corollary}
\begin{proof}
An order-$\nu$ derivative of $(C_i-zI)^{-1}$ is a sum of words containing
at most $|\nu|+1$ resolvents. Multiplication by
$(d_1d_2)^{|\nu|+1}$ therefore leaves a matrix polynomial of degree at
most $(|\nu|+1)(m_1+m_2)$ with the same Hermitian reality relation.
Apply the preceding zero-count proof to every consumed jet. Fixed energies
avoid derivatives of the sampling map; moving unitary frames use
Corollary~\ref{p1b:cor:v226-source-jets}. Matrix inversion on its regular
domain and contour residues preserve equality. The minimal visible
realization statement is
Theorem~\ref{p1b:thm:response-minimal-visible-realization}.
\end{proof}

\begin{remark}[Exact comparison, complete ports and variable typing]
These tests compare independently licensed actual parent/source words.
One may include already owned source columns in an augmented response to
compare their incidence as well; equality of an unaugmented kernel alone
does not identify an omitted source. A known finite dimension bound is
part of the certificate. Approximate sample agreement consumes its own
conditioning and spectral-separation estimates. The theorem gives all
energy dependence at a fixed source point, and finite source jets there;
it does not turn finitely many curvature samples into a global identity
for an arbitrary running curvature profile.
\end{remark}

\begin{corollary}[Complementary zero modes do not obstruct nonreal comparison]
\label{p1b:cor:v227-zero-mode-comparison}
No positive lower bound on $C_i$ is needed for these upper-half-plane
tests. A zero complementary eigenvalue is retained as part of the same
meromorphic response. At real zero it is removable exactly when its actual
coupling is zero, as in
Corollary~\ref{p1b:cor:v2210-lit-boundary-zero-mode}. Otherwise the coupled
mode must remain in a real-zero stationary calculation.
\end{corollary}
\begin{proof}
A Hermitian matrix has real spectrum, so $C_i-zI$ is invertible for every
nonreal $z$. Its coefficient at a real pole is $B_iP_\lambda B_i^*$,
which vanishes exactly when $B_iP_\lambda=0$.
\end{proof}

\subsection{Action symmetry controls canonical finite-state covariance}
\label{p1b:sec:v227-state-symmetry}
\begin{theorem}[Quantitative full-ground and Gibbs covariance]
\label{p1b:thm:v227-state-symmetry}
Let $H=H^*$ be finite, $U$ unitary and $\epsilon=\|[H,U]\|$.
If $P$ is the full ground projector of rank $r<\dim H$, with gap
$\Delta>0$ to its complement, set $\rho_0=P/r$. Then
\[
 \boxed{\|U\rho_0U^*-\rho_0\|_1
       \le2\min\{1,\epsilon/\Delta\}.}
\]
For the normalized Gibbs density
$\rho_\beta=e^{-\beta H}/\operatorname{Tr}e^{-\beta H}$, $\beta\ge0$,
\[
 \boxed{\|U\rho_\beta U^*-\rho_\beta\|_1
       \le2\min\{1,\beta\epsilon\}.}
\]
Both states are exactly invariant when $[H,U]=0$.
If $H$ is scalar, its normalized full-ground state is $I/\dim H$ and
is already invariant. The ground statement does not apply to an arbitrary
chosen orientation inside a degenerate ground space.
\end{theorem}
\begin{proof}
Let $Q=I-P$ and $E_0$ be the ground energy. Then
\[
 QUP=[Q(H-E_0)Q]^{-1}Q[H,U]P,\qquad
 \|QUP\|\le\epsilon/\Delta.
\]
The equal-rank projectors $P$ and $UPU^*$ have paired principal-angle
singular values, so their trace-norm difference is at most $2r\|QUP\|$.
Divide by $r$ and use the universal trace distance bound two.

For Gibbs states put $H_s=H+s(UHU^*-H)$. Finite matrix exponential
differentiation gives
\[
 \partial_s\rho_s=-\beta\int_0^1
 \rho_s^{1-v}\bigl(\partial_sH_s-
        \operatorname{Tr}(\rho_s\partial_sH_s)I\bigr)\rho_s^v\,dv.
\]
Schatten H\"older bounds the trace norm of each sandwiched bounded matrix
by its operator norm. The centered matrix has norm at most $2\epsilon$,
so integration in $s$ gives $2\beta\epsilon$. At $s=1$ the density is
$U\rho_\beta U^*$ by unitary functional calculus. Again clamp by two.
The assertions at zero defect follow. A proper substate of a degenerate
$P$ need not commute with a unitary acting inside that eigenspace.
\end{proof}


\begin{theorem}[Sharp finite Gibbs covariance under approximate unitary symmetry]
\label{p1b:thm:v230-sharp-gibbs-covariance}
For a finite Hermitian $H$, unitary $U$, inverse temperature $\beta\ge0$,
$\epsilon=\|[H,U]\|$, and the normalized canonical state
$\rho_\beta=e^{-\beta H}/\operatorname{Tr}(e^{-\beta H})$, one has
\begin{equation}\label{p1b:eq:v230-gibbs-sharp}
 \boxed{\|U\rho_\beta U^*-\rho_\beta\|_1
 \le\min\{2,\beta\epsilon\}.}
\end{equation}
The linear coefficient one cannot be decreased uniformly, even on a
single two-level canonical branch.  This strengthens, but does not alter,
the valid earlier estimate in Theorem~\ref{p1b:thm:v227-state-symmetry}.
\end{theorem}
\begin{proof}
Put $K=UHU^*$, $\rho=\rho_\beta(H)$, and
$\sigma=\rho_\beta(K)=U\rho U^*$, so
$\|K-H\|=\|[H,U]\|=\epsilon$.  Both partition functions agree.
For $\beta>0$, the quantum relative entropies obey the exact identity
\[
 D(\rho\|\sigma)+D(\sigma\|\rho)
 =\beta\operatorname{Tr}\bigl[(\rho-\sigma)(K-H)\bigr].
\]
Quantum Pinsker's inequality, applied to both ordered pairs, gives
$\|\rho-\sigma\|_1^2\le D(\rho\|\sigma)+D(\sigma\|\rho)$.
Trace H\"older then yields
$\|\rho-\sigma\|_1^2\le\beta\epsilon\|\rho-\sigma\|_1$.
Division when the distance is nonzero, together with the universal bound
two, proves the estimate.  At $\beta=0$ the state is maximally mixed.
For optimality, take a two-level Hamiltonian with energy separation $d$
and let $U$ exchange its two eigenspaces.  The distance is
$2\tanh(\beta d/2)$ and $\epsilon=d$; their small-$\beta d$
ratio tends to one.
\end{proof}

\begin{corollary}[Sharpened finite Gibbs representative comparison]
\label{p1b:cor:v230-sharp-gibbs-chart}
Under the representative hypotheses and notation of
Corollary~\ref{p1b:cor:v227-chart-covariance}, on a Gibbs branch
\[
 \bigl|\operatorname{Tr}(\rho_\beta A^{\rm cmp})
       -\operatorname{Tr}(\rho_\beta A)\bigr|
 \le\delta_A+\|A\|\min\{2,\beta\epsilon\}.
\]
No change is made to the full-ground estimate, to the requirement of an
actual observable representative, or to the independent continuum
covariance and Osterwalder--Schrader obligations.
\end{corollary}
\begin{proof}
Use the same transported-representative decomposition as in
Corollary~\ref{p1b:cor:v227-chart-covariance}, with the improved state
trace-distance bound of Theorem~\ref{p1b:thm:v230-sharp-gibbs-covariance}.
\end{proof}

\begin{corollary}[The finite state defect is not a separate selector]
\label{p1b:cor:v227-chart-covariance}
For any bounded observable $A$ and compared representative $A^{\rm cmp}$,
put $\delta_A=\|A^{\rm cmp}-UAU^*\|$.
On the normalized full-ground branch,
\[
 |\operatorname{Tr}(\rho_0 A^{\rm cmp})-
   \operatorname{Tr}(\rho_0 A)|
 \le\delta_A+2\|A\|\min\{1,\epsilon/\Delta\}.
\]
On the Gibbs branch replace $\epsilon/\Delta$ by $\beta\epsilon$.
Thus an admitted refinement has state covariance once its representative
error and the printed action-defect ratio vanish. Uniform positive ground
gap is sufficient but not necessary. Exact equivariance removes the state
term entirely. This does not establish the remaining chart or continuum
comparison, reflection positivity, clustering or a physical invariant measure.
\end{corollary}
\begin{proof}
Add and subtract the transported expectation, use the norm-one state bound
for $\delta_A$, then apply the trace duality bound and the theorem.
If $H^{\rm ref}$ commutes with $U$, an actual native residual
$D=H-H^{\rm ref}$ gives $\epsilon\le2\|D\|$; the norm of $D$ may be
bounded by the 137 coefficients of
Theorem~\ref{p1b:thm:v226-dual-coefficients}, or by its operator-valued
tensor bound. The limiting statements follow from these actual inequalities.
\end{proof}

\subsection{The exact centre quotient for a fixed neutral readout}
\label{p1b:sec:v227-central-readouts}
\begin{theorem}[Block-central invariance of native even dynamics]
\label{p1b:thm:v227-central-readouts}
Let $\mathcal A=\bigoplus_a M_{n_a}(\mathbb C)$, with central projections
$P_a$, and let $H\in\mathcal A$ be Hermitian. Put
$C=\sum_a c_aP_a$ with real $c_a$. Then
\[
 \operatorname{ad}_{H+C}|_{\mathcal A}
       =\operatorname{ad}_H|_{\mathcal A}.
\]
Every even Heisenberg operator word and its fixed-state stationary source
response formed from $\operatorname{ad}_H$ is unchanged when the same
actual density and support are retained.
For fixed probabilities $p_a$, the sectorwise Gibbs family
\[
 \rho_{\beta,p}=\bigoplus_a
       p_a\frac{e^{-\beta H_a}}{\operatorname{Tr}e^{-\beta H_a}}
\]
is also unchanged by $H_a\mapsto H_a+c_aI$. The same holds for normalized
full-ground projectors chosen within each occupied block, with fixed block
weights. These conclusions do not choose the probabilities.
\end{theorem}
\begin{proof}
$C$ commutes with every element of $\mathcal A$ and with $H$, proving the
commutator and exponential-conjugation identities. Each block Gibbs
numerator and denominator have the same factor $e^{-\beta c_a}$.
Within a fixed block, a scalar shift changes no eigenspace or ordering.
Products, stationary expectations and source resolvents are built from
these unchanged operators and the retained actual state/support.
\end{proof}

\begin{proposition}[Odd sources and absolute energy retain the centre]
\label{p1b:prop:v227-central-boundary}
For a source block $q_{ab}=P_aqP_b$,
\[
 [H+C,q_{ab}]=[H,q_{ab}]+(c_a-c_b)q_{ab}.
\]
Hence a source crossing central blocks can read their relative energy;
its frequencies shift by $c_a-c_b$ in the common energy convention.
The absolute heat trace instead obeys
\[
 \operatorname{Tr}e^{-t(H+C)}
 =\sum_a e^{-tc_a}\operatorname{Tr}e^{-tH_a}.
\]
Globally normalized Gibbs weights are
$p_a=e^{-\beta c_a}Z_a/\sum_b e^{-\beta c_b}Z_b$ and can change.
The global ground sector can also change. Metric/source derivatives of
$C$ remain in absolute energy and stress; the even-dynamics quotient does
not remove them.
\end{proposition}
\begin{proof}
Multiply $q_{ab}$ on its two ends by the central projectors to obtain
the commutator. Blockwise functional calculus proves both exponential
identities and the global weights. Absolute energy changes by
$\operatorname{Tr}(\rho C)$; its actual differentiated word includes the
central derivatives and any state/weight response.
\end{proof}

\begin{corollary}[Occupied-block covariance needs only occupied-block gaps]
\label{p1b:cor:v227-block-covariance}
If $U$ preserves the central blocks and the actual state is a fixed mixture
of normalized full-ground or Gibbs block densities, its trace-norm covariance
defect is bounded by the weighted sum of
Theorem~\ref{p1b:thm:v227-state-symmetry}'s block bounds, over occupied
blocks only. A relative crossing of their scalar ground energies creates
no additional within-block state-covariance defect at fixed weights.
\end{corollary}
\begin{proof}
The trace norm is additive on direct sums. Apply the theorem separately
on each block with positive weight and multiply by that weight.
Unoccupied blocks contribute zero. Reordering global block energies changes
a globally selected ensemble, a different state from the stated fixed mixture.
\end{proof}


\subsection{Physical weak-carrier descent and the generic mapping-torus carry}
\label{p1b:sec:v229-physical-weak-vs-carry}
\begin{remark}[Native descent and the separate invariant carry witness]
\label{p1b:rem:v229-physical-triviality-vs-carry}
The explicit equivariant section and determinant-one gauge of
Proposition~\ref{p1b:v215:prop:carrier-equivariant-origin}
show that the specified native residual carrier descends to the
trivial principal $SU(2)$ bundle with pure $+I_2$ fixed-fibre
isotropy.  Under $\mathsf H_{\rm vert}$ this is the weak-role
carrier; the later proof discharges that role premise.
This physical quotient-origin certificate is independent of the
abstract mapping-torus twisted-loop difference group.  For a free
gauge-loop lift $z\in\Gamma\cong\mathbb Z\oplus C_2^5$, its carry
\[
 v=\Theta_\kappa(z)-z,\qquad
 [v]\in\ker(A-I)/\operatorname{im}(A-I)\cong C_2
\]
requires evaluating the actual monodromy on $z$.  A fixed identity,
zero bundle class, or canonical physical section does \emph{not} force
$[v]=0$ by itself.  Theorem~\ref{p1b:thm:trivial-component-zero-carry}
supplies that separate certificate in the trivial component.
The endpoint normalization, other-component carries and
ordered-cover pullbacks remain classification questions.
Triviality of the underlying weak principal bundle does not assert a
flat inherited gauge connection or zero gauge curvature.
\end{remark}


\section{Conditional native \texorpdfstring{\(SU(2)\)}{SU(2)} atlas and its geometric consequences}
\label{p1b:sec:v259-conditional-native-atlas}

\subsection{Minimal native handoff and the derived residual provider}
\label{p1b:sec:conditional-residual-descent}

The conditional atlas has the same logical organization as the level-one
atlas of Paper~I-A.  In this paper the native weak-role premise is
provisional: every weak-sector application below is conditional on
\(\mathsf H_{\rm vert}\).  Later Papers~III-B and III-F prove
that the generated compact weak sector satisfies this premise; its
explicit discharge is Corollary~\ref{p3f:cor:v259-native-su2-discharge}.
It is therefore not an additional axiom of the completed framework.
The labelled geometry supplies mathematical constructors; the later
papers establish that their generated sector realizes those constructors.  A hypothesis identifying an
entire bundle, exchange lift and connection is stronger than needed here.
The carrier constructor and its inherited connection are already determined
by the relative geometry of this paper.

\begin{definition}[Minimal native vertical-role handoff]
\label{p1b:def:conditional-residual-datum}
Retain the labelled reciprocal product, its actual relative map
\(r(a,b)=ab^{-1}\), its scalar factor, and exchange.  The assertion
\(\mathsf H_{\rm vert}\) means that a named weak principal role uses the
connected derived right-translation action inside the actual fibres of
\(r\), after the stated diagonal scalar quotient, with exchange induced
from the source.  It asserts the role of this actual action, rather than
only an abstract Lie-algebra isomorphism with \(\mathfrak{su}(2)\).
It does not prescribe a connection, metric, action, representation
multiplicity, boundary condition or source.

The mathematical residual provider will be denoted
\[
 P_{\rm res}=SU(2)_x\times SU(2)_p\times U(1)_{\rm rel},\qquad
 \pi(a,b,z)=(ab^{-1},z),\qquad (a,b,z)h=(ah,bh,z).
\]
Write \(r_v=ab^{-1}\), \(z=z_tz_E^{-1}\),
\(K_{\rm rel}=SU(2)\times U(1)\), and
\(\beta(a,b,z)=(b,a,z^{-1})\).
These displayed objects are derived below; their physical role is the only
content of \(\mathsf H_{\rm vert}\).

In particular, this premise does not assert that every principal
\(SU(2)\) bundle on \(M_8\) admits descent.  The independent
mathematical results below hold on their displayed input domains;
\(\mathsf H_{\rm vert}\) is used only when applying them to the
named weak role.  A result using the inherited connection constructs it
from the parent Cartan datum, rather than adding it to this premise.
\end{definition}

\begin{theorem}[The native residual carrier and connection are derived data]
\label{p1b:thm:minimal-native-provider}
The right-translation fibres of \(r\) determine \(\Delta K\) uniquely.
Its connected derived group is \(\Delta SU(2)\).  Quotienting by the
already typed \(Z_{\rm sc}=U(1)_{\rm diag}\) gives precisely
\(P_{\rm res}\to K_{\rm rel}\), with the displayed induced exchange.
The fundamental associated module is the unique nontrivial irreducible
complex two-dimensional \(SU(2)\) module, up to intertwining; no assertion
of physical multiplicity is included.

Let \(q_{\rm sc}:M_8\to P_{\rm res}\), and project the parent Cartan
connection to its semisimple quotient.  The result is scalar-basic and has
one unique descent \(\omega_{\rm B}\), with
\[
 q_{\rm sc}^*\omega_{\rm B}=\pi_{\mathfrak{su}}\mathcal A_{M_8},
 \qquad \omega_{\rm B}=\tfrac12(\theta_x+\theta_p).
\]
It is the unique connection invariant under both the full left homogeneous
action and exchange.  Neither its formula nor the choice of its homogeneous
midpoint is an extra hypothesis of \(\mathsf H_{\rm vert}\).
This constructs an inherited connection; a statement that a separately
supplied physical operator consumes it retains that operator's own origin
and domain conditions.
\end{theorem}
\begin{proof}
Theorems~\ref{p1b:thm:relative} and
\ref{p1b:cor:relative-carrier-rigidity} give the actual fibres and their
unique right-translation group.  Its derived subgroup is
\(\Delta SU(2)\); the scalar quotient and exchange formulas follow by
evaluation.  The section is \(s(r,z)=(r,e,z)\).

By Theorem~\ref{p1b:thm:cartan}, the parent connection reproduces scalar
vertical vectors.  Semisimple projection kills those values and is invariant
under the central scalar action.  It is therefore basic.  Descent through
the surjective submersion is unique; residual reproduction and equivariance
make it a principal connection.  Projecting the displayed parent average
gives its formula and exchange invariance.

Independently, a homogeneous connection is an equivariant map
\(W(X,Y,t)=MX+NY+vt\).  Principal reproduction gives \(M+N=I\).
Commuting with every adjoint rotation forces \(M=fI\),
\(N=(1-f)I\), and \(v=0\).  Exchange forces \(f=1/2\).
This replays Theorem~\ref{p1b:thm:homogeneous} without assuming the
connection as part of a later physical identification.
\end{proof}

The later proof of the provisional premise has two explicit stages.
Paper~III-B, Theorems~\ref{p3b:thm:canonical-compact-seven-quotient} and
\ref{p3b:thm:parent-cartan-scalar-diagonal-descent}, generates this same
quotient diagram, its induced exchange and its inherited connection from
the already constituted source.  Paper~III-F,
Theorem~\ref{p3f:thm:v2183-native-weak-carrier-class} and
Corollary~\ref{p3f:cor:v259-native-su2-discharge}, proves that the
generated weak carrier satisfies exactly \(\mathsf H_{\rm vert}\),
with the trivial ordinary class and pure fixed isotropy asserted below.
Those proofs apply the independent mathematical constructions in this
paper.  No later physical-occurrence conclusion is a premise of the
constructions' proofs, so this order of exposition is not circular.

The generic antecedent
\(\mathsf D_E:\operatorname{Desc}(E)\ne\varnothing\) is a separate
existence statement for an arbitrary bundle.  Discharging the native
premise neither asserts \(\mathsf D_E\) for every bundle nor evaluates
every point of a generic descent fibre.

\begin{theorem}[Specified descent, global gauge and the inherited connection]
\label{p1b:thm:conditional-residual-descent}
Let \(E_{\rm w}=P_{\rm res}\times_{SU(2)}\C^2\), and pull it back to
\(M_8\) by \((r_v,z)\).  In the section
\(s(r,z)=(r,e,z)\), its exchange cocycle is \(c(y)=ab^{-1}\).
The determinant-one gauge \(u(y)=b^{-1}\) satisfies
\[
 u(By)c(y)u(y)^{-1}=I.
\]
Thus this particular pullback is equivariantly trivial with pure \(+I_2\)
fixed-fibre isotropy.  Its quotient bundle is \(Q\times SU(2)\), where
\(Q=\operatorname{SP}^2(S^1\times S^3)\).  In particular it has
\(c_2=0\), \(\epsilon_{\rm IB15}=0\), and a specified trivial quotient
origin, independently of the other points of
\(\operatorname{Desc}(M_8\times SU(2))\).

For the section-frame connection
\(A_{\rm rel}=\tfrac12r_v^{-1}dr_v\), the same gauge gives
\[
 \boxed{
 A_{\rm taut}=uA_{\rm rel}u^{-1}-du\,u^{-1}
             =\tfrac12(\theta_x+\theta_p),\qquad
 F_{\rm taut}=-\tfrac14(\theta_x-\theta_p)
                         \wedge(\theta_x-\theta_p).}
\]
This connection is exchange invariant and nonflat, with restricted holonomy
\(SU(2)\).  Nevertheless \(\operatorname{Tr}(F_{\rm taut}\wedge
F_{\rm taut})=0\) pointwise.
\end{theorem}
\begin{proof}
Every point of \(P_{\rm res}\) is uniquely \(s(r,z)b\), proving ordinary
triviality.  Exchange sends
\(s(r,z)\) to \(s(r^{-1},z^{-1})r\), giving the cocycle.  Direct
multiplication yields \(a^{-1}(ab^{-1})b=I\); the transformed lift is
therefore \((y,h)\mapsto(By,h)\), whose quotient is the displayed product.
This proves the actual quotient origin, rather than inferring it from Chern
classes.

The identity
\(r_v^{-1}dr_v=b(\theta_x-\theta_p)b^{-1}\) and
\(du\,u^{-1}=-\theta_p\) give the transformed connection.  The two
Maurer--Cartan equations give its curvature.  With
\(\mathcal E=\tfrac12(\theta_x-\theta_p)\), the horizontal curvature is
\(-[X,Y]\); these values span \(\mathfrak{su}(2)\), so the restricted
holonomy is the full connected group.  Finally the curvature is a gauge
transform of the pullback of a two-form on the relative three-sphere.  Its
square is therefore zero.  Neither vanishing characteristic classes nor this
vanishing four-form implies vanishing curvature.
\end{proof}

\begin{corollary}[Associated representations and the two central signs]
\label{p1b:cor:conditional-residual-associated}
For every finite-dimensional representation \(\rho:SU(2)\to U(V)\), the
associated residual bundle and its pullback are ordinarily trivial.  For the
specified \(+I\) lift their quotient is the trivial associated bundle on
\(Q\).  Tensor products, duals, exterior and symmetric powers, and equivariant
intertwiners commute with these pullbacks and descents.  This supplies no
physical occurrence assertion for additional representations.

For spin \(\iota\), the alternative central lift acts on sections by
\[
 U_{\iota,-}=(-1)^{2\iota}U_{\iota,+}.
\]
Thus the two lifts give the same associated exchange action in integer-spin
representations and interchange the two exchange sectors in half-integer-spin
representations.  In particular their adjoint connections and local
adjoint-valued curvature words agree.  A half-integer representation of the
negative lift has nontrivial stabilizer action and is not an ordinary
coarse-quotient vector-bundle descent until the central sign is changed.

For the midpoint connection, the parallel associated vectors are precisely
the \(SU(2)\)-invariants; in particular a nontrivial irreducible has none.
Parallel intertwiners have exactly the multiplicity-space form in
Proposition~\ref{p1b:result:parallel-maps-under-the-midpoint-connection}.
Every associated bundle of this specified branch has vanishing positive
Chern classes.  On the closed spin manifold \(M_8\), its twisted chiral
Dirac index is zero.
\end{corollary}
\begin{proof}
Apply the associated-bundle functor to the explicit equivariant product
trivialization.  The central element \(-I\) acts on spin \(\iota\) as
\((-1)^{2\iota}I\), which proves the sector statements; it acts trivially
in the adjoint representation.  Full restricted holonomy gives the parallel
section and intertwiner assertions.  Triviality gives the Chern classes, and
\(\widehat A(TM_8)=1\) makes the top-degree twisted index density zero.
This index statement does not exclude pairs of chiral zero modes.
\end{proof}

\subsubsection{The minimal spatial bit and its complete invariant consequences}
\label{p1b:sec:minimal-born-spatial}

Spatial assertions in this subsection concern either invariant scalar slice
\(z=1\) or \(z=-1\) of \(K_{\rm rel}\), or a connection pulled back
from its sphere.  A general connection on the four-dimensional product need
not be independent of \(z\).  Exchange otherwise relates the slices
\(z\) and \(z^{-1}\).

\begin{theorem}[Four spatial lifts and the minimal odd Born class]
\label{p1b:thm:minimal-spatial-lifts}
Let \(\jmath(r)=r^{-1}\) on \(S^3=SU(2)\).  Every principal
\(SU(2)\) bundle there is trivial.  In a section, a coherent lift is
\[
 (r,h)\longmapsto(r^{-1},c(r)h),\qquad c(r^{-1})c(r)=I.
\]
Its gauge-isomorphism class is determined exactly by the two central signs
\(c(1)=\sigma_+I\), \(c(-1)=\sigma_-I\).
There are exactly four classes, represented by \(I,-I,r,-r\).
Put \(\nu=(1-\sigma_+\sigma_-)/2\).  Then
\[
 \deg c\equiv\nu\pmod2,
 \qquad \deg\bigl(g(r^{-1})c(r)g(r)^{-1}\bigr)
                  =\deg c-2\deg g.
\]
The native residual lift is \(c(r)=r\), so has signs \((+,-)\),
\(\nu=1\), and degree one.  Thus the minimal spatial topological
condition for the odd branch is one bit, \(\nu=1\).  The sign at \(r=1\)
fixes the sign of its reciprocal index.

An ordinary coarse-quotient principal descent through
\(S^3\to S^3/\jmath\) exists exactly for the class \((+,+)\).
The native odd spatial lift remains an equivariant bundle there.  This
does not contradict its pulled-back pure \(+I\) descent on \(M_8/B\):
the latter has fixed set \(\Delta K\), where \(r=1\).
\end{theorem}
\begin{proof}
Ordinary triviality follows from \(\pi_2SU(2)=0\).  At either fixed
point coherence forces \(c^2=I\), hence the two displayed central
values.  Gauge changes conjugate these central values, so preserve them.
The four representatives realize all sign pairs.

For completeness, compare two lifts with the same signs.  Their principal
isomorphism bundle has fibre \(SU(2)\).  Use the equivariant CW structure
on \(S^3\) with two fixed zero-cells and free cells elsewhere.  Choose
an isomorphism at the fixed cells; matching central signs make this
equivariant.  Extend over each free cell and its mate.  Obstructions are
in \(\pi_0,\pi_1,\pi_2\) of \(SU(2)\), all zero, because the
base has dimension three.  This gives an actual equivariant isomorphism,
proving completeness without an unstable eight-dimensional calculation.

Multiplication adds degrees in \(H^3(SU(2);\Z)\); inversion has degree
minus one.  The gauge-degree formula follows.  Comparing with the four
representatives proves the parity formula.  Alternatively a regular fixed
target value has one fixed preimage modulo two exactly for opposite signs;
all other preimages occur in pairs.  Trivial stabilizer action is necessary
for ordinary coarse descent.  Only \((+,+)\) has it, and its representative
\(c=I\) descends to the product bundle, proving sufficiency.
\end{proof}

\begin{corollary}[Spatial equivariant gauge components]
\label{p1b:cor:minimal-spatial-gauges}
For each of the four lift classes, every equivariant gauge has ordinary
winding zero, while its equivariant homotopy components form \(C_2\).
These are three-dimensional spatial gauge components, not the full gauge
group on \(M_8\) or the intermediate \(F_7\) loop group.
\end{corollary}
\begin{proof}
The degree formula for a gauge preserving the same lift gives
\(2\deg g=0\).  Its adjoint group bundle has trivial stabilizer action
at the two fixed points and descends as a group bundle over
\(S^3/\jmath\cong\Sigma\mathbb {RP}^2\).  The fibre is two-connected
and inner conjugation acts trivially on \(\pi_3SU(2)=\Z\).
Obstruction theory for its sections in dimension three, with the identity
section as origin, gives
\(H^3(\Sigma\mathbb {RP}^2;\Z)=C_2\).  Pointwise multiplication gives
addition of these obstruction classes.
\end{proof}

\begin{theorem}[Half Chern--Simons class and forced full holonomy]
\label{p1b:thm:minimal-born-cs-holonomy}
Let \(A\) be any smooth connection invariant under a specified spatial
lift.  Use the source normalization
\(N(A)=(8\pi^2)^{-1}\int\operatorname{tr}(A\wedge dA+2A^3/3)\)
and \(\deg r=1\).  Then
\[
 \boxed{2N(A)=\deg c,\qquad N(A)\equiv\nu/2\pmod\Z.}
\]
Every derivative of \(N\) along a smooth family invariant under this
same lift is zero.  For a globally invariant connection on
\(S^3\times S^1\), its slice values instead satisfy
\(N(z)+N(z^{-1})=1\) in the native section, and equal \(1/2\) at
\(z=\pm1\).

If \(\nu=1\), every invariant connection has full restricted
\(SU(2)\) holonomy; in particular none is flat.  Conversely, when
\(\nu=0\), an invariant flat connection exists.  Thus opposite fixed
signs are necessary and sufficient to obstruct invariant flatness.
For every nontrivial irreducible associated representation, the covariant
scalar Laplacian on the compact sphere has zero kernel and a positive
connection-dependent first eigenvalue.  Parallel intertwiners have the
usual Schur multiplicity-space form.
\end{theorem}
\begin{proof}
Invariance in the section gives
\(\jmath^*A=cAc^{-1}-dc\,c^{-1}=A^c\).
Orientation reversal and the gauge-transgression identity give
\(-N(A)=N(A)-\deg c\).  This proves the formula, its gauge-invariant
fractional part, and constancy along invariant families.  Applying the same
identity to the two exchanged scalar slices gives their paired formula.

On simply connected \(S^3\), proper holonomy would be contained in a
circle or be trivial.  It would give a nonzero parallel adjoint direction
\(n\).  Exchange preserves parallel directions and is the identity on
the adjoint fibre at a fixed point, so fixes \(n\) everywhere.  Since
\(H^2(S^3;\Z)=0\), lift this direction through the Hopf fibration to
a gauge in which it is constant.  The transformed exchange cocycle must
then lie in its centralizer \(U(1)\).  Its degree is zero.  But the
gauge-degree formula preserves the odd degree parity, a contradiction.
Hence holonomy is the full group.  Equal signs have the constant lift
\(c=I\) or \(-I\), for which the zero connection is invariant.
Finally the Bochner identity identifies the scalar kernel with parallel
vectors.  Full holonomy and Schur's lemma give the remaining assertions;
compact ellipticity gives positivity of the nonzero spectral floor.
\end{proof}

\begin{theorem}[Reciprocal Dirac index for every invariant connection]
\label{p1b:thm:minimal-born-dirac-index}
Supply a \(\jmath\)-invariant Riemannian metric on \(S^3\), the standard
spin structure, and the orientation-reversing involutive spin transport
continued from the round reference.  Choose its sign positive at \(r=1\).
For spin \(\iota\) twisting, let \(D_{A,\iota}\) use any smooth
connection invariant under the given lift.  The combined transport
\(U_\iota\) preserves its \(H^1\) domain and obeys
\[
 U_\iota^2=I,\qquad U_\iota D_{A,\iota}U_\iota=-D_{A,\iota}.
\]
The index of its odd Fredholm block, or graded heat trace, is
\[
 \boxed{\mathcal I_\iota=
 \frac{2\iota+1}{4}(\sigma_+^{2\iota}-\sigma_-^{2\iota}),\qquad
 \operatorname{Tr}(U_\iota e^{-tD_{A,\iota}^2})=\mathcal I_\iota.}
\]
For the native \((+,-)\) lift this is zero for integral \(\iota\), and
\(\iota+1/2\) for half-integral \(\iota\).  If
\(h_\iota=\dim\ker D_{A,\iota}\), then
\[
 h_\iota\ge|\mathcal I_\iota|,\qquad
 h_\iota\equiv\mathcal I_\iota\pmod2,\qquad
 \eta(D_{A,\iota})=0,\qquad
 \bar\eta(D_{A,\iota})\equiv\mathcal I_\iota/2\pmod\Z.
\]
For a general associated module the indices add.  In particular the native
index is half the complex dimension of its \(-I\)-odd isotypic part.
The fundamental zero is forced without assuming a Cartan profile or a round
metric.  Simplicity and a numerical gap require additional control.
Among the four lift classes, \(\nu=1\) is also necessary for a condition
on the lift alone to force a fundamental zero for every invariant metric
and connection.
\end{theorem}
\begin{proof}
For the round left spin frame the spin transport is
\(r^{-1}\) composed with pullback by inversion.  The twisting transport
is \(\rho_\iota(c(r^{-1}))\) composed with that pullback.  Coherence
gives an involution, and the spin transport changes the Clifford sign under
orientation reversal.  Naturality of the spin and twisting connections
therefore gives the anticommutation relation.  It persists for invariant
metrics by equivariant orthonormalization and continuation from the round
metric; the two fixed-point spin signs cannot change along this path.

Nonzero eigenvalues pair with their negatives.  The graded heat trace is
constant and is the difference of the two kernel grades.  Its small-time
limit can be calculated directly.  The only fixed points are \(r=1,-1\),
with derivative \(-I_3\).  The off-diagonal Gaussian heat kernel integrates
to \(|\det(I_3-(-I_3))|^{-1}=1/8\) at each point.  The combined fibre
traces are \(2(2\iota+1)\sigma_+^{2\iota}\) and
\(-2(2\iota+1)\sigma_-^{2\iota}\).  Their sum divided by eight is
the displayed index.  This gives an independent local calculation of the
fundamental index one.  The kernel inequality and parity follow from
\(h=h_++h_-\), \(\mathcal I=h_+-h_-\).  Spectral pairing makes the
eta function zero in its initial convergence half-plane and by continuation;
the reduced-eta assertion follows.
For equal fixed signs the constant representative admits the flat connection
on the round sphere; its fundamental twist is two untwisted spin copies,
with no zero and gap \(3/(2R)\).  This proves the stated necessity.
\end{proof}

\begin{corollary}[Representation trace and phase consequences]
\label{p1b:cor:minimal-born-representation-phase}
In spin \(\iota\), set
\(I_\iota=\frac23\iota(\iota+1)(2\iota+1)\), with fundamental
trace normalized to \(I_{1/2}=1\).  Then
\[
 N_\iota(A)=I_\iota N_{1/2}(A),\qquad
 e^{2\pi iN_\iota(A)}
   =e^{2\pi i\bar\eta(D_{A,\iota})}=(-1)^{\mathcal I_\iota}
\]
for an odd lift.  These are geometric trace and spectral signs; no physical
determinant action or anomaly adoption is asserted.
\end{corollary}
\begin{proof}
The invariant quadratic trace form scales by the stated Dynkin index,
obtained from \(\sum_{m=-\iota}^{\iota}m^2\); the cubic Chern--Simons
form scales by the same factor.  For integral \(\iota\), \(I_\iota\)
is even.  If \(d=\iota+1/2\) is integral, then
\(I_\iota=d(4d^2-1)/3\equiv d\pmod2\).  Combine this with the
half Chern--Simons and reduced-eta formulas.  Reversing the odd fixed-sign
pair changes the index sign and preserves this parity.
\end{proof}

\subsubsection{Complete midpoint Dirac provider in every isospin}
\label{p1b:sec:minimal-all-isospin-dirac}
\begin{theorem}[Complete signed midpoint spectrum]
\label{p1b:thm:minimal-all-isospin-dirac}
Supply the round radius \(R\) and the inherited midpoint connection.
Let \(j,k\in\tfrac12\Z_{\ge0}\), \(L=j+k\),
\(\delta=j-k\), \(\tau=\iota-1/2\), and \(p=\iota+1/2\).
Only labels with \(L-\tau\in\Z\) occur.  The following disjoint cases
give the complete signed spectrum of \(D_{B,\iota}\):
\[
\begin{array}{c|c|c}
\text{label domain}&R\lambda&\text{multiplicity per displayed value}\\\hline
\iota\ge1/2,\ L=\tau,\ |\delta|\le\tau
   &\delta&(2j+1)(2k+1)\\
|\delta|=p,\ L\ge p
   &-\operatorname{sgn}(\delta)(L+1)&(2j+1)(2k+1)\\
\iota\ge1/2,\ L\ge p,\ |\delta|\le\tau
   &\pm\sqrt{L(L+2)+\delta^2-\tau(\tau+2)}&(2j+1)(2k+1)
\end{array}
\]
The last row has one copy of each sign.  Accidental equal numerical values
add their labelled multiplicities.  In particular
\[
 h_{B,\iota}=\begin{cases}0,&\iota\in\Z_{\ge0},\\
 (\iota+1/2)^2,&\iota\in1/2+\Z_{\ge0},\end{cases}
\]
and the nonzero absolute gap and its multiplicity per sign are
\[
\begin{array}{c|c|c}
\iota&Rg_\iota&m_\iota\\\hline
0&3/2&2\\
1/2&\sqrt3&4\\
\Z_{\ge1}&1/2&\iota(\iota+1)\\
3/2+\Z_{\ge0}&1&(\iota-1/2)(\iota+3/2).
\end{array}
\]
For half-integral \(\iota\), with \(d=\iota+1/2\), the zero grades
have dimensions \(d(d+1)/2\) and \(d(d-1)/2\).
\end{theorem}
\begin{proof}
Use the angular normalization of the fundamental theorem.  On
\(V_{1/2}\otimes V_j\otimes V_\iota\), with spectator dimension
\(2j+1\), direct Clifford contraction gives
\[
 RD=\tfrac32+2\sigma\cdot J+\sigma\cdot I.
\]
The Pauli identity \((\sigma\cdot J)^2=J^2-\sigma\cdot J\),
and its isospin version, give the exact matrix identity
\[
 R^2D^2=2K^2+2j(j+1)+\tfrac34-\iota(\iota+1),
 \qquad K=J+I+\sigma/2.
\]
Its scalar on a total-\(k\) block is
\(L(L+2)+\delta^2-\tau(\tau+2)\).
Clebsch--Gordan gives one copy on each boundary row and two in the
interior.  In the latter, couple \(J+I=u=k\pm1/2\).
The two diagonal expectations of \(2S\cdot J\), \(S=\sigma/2\),
sum to minus one: use the projection coefficient
\([u(u+1)+j(j+1)-\iota(\iota+1)]/[2u(u+1)]\).
The two values of \(2S\cdot(J+I)\) sum to minus two.
Consequently the multiplicity matrix of \(RD\) has trace zero, so its
roots are the displayed opposite square roots.  A zero \(u\) term is
evaluated by continuity, or directly as zero.
For \(L=\tau\), the same one-dimensional calculation gives
\(R\lambda=j-k\).  For \(|j-k|=p\), it gives
\(-\operatorname{sgn}(j-k)(L+1)\).

Only the lower boundary can contain zero, and it does so exactly at
\(j=k=\tau/2\), which is an allowed spin precisely when \(\iota\)
is half-integral.  Its dimension is \((\tau+1)^2\).
For integral positive \(\iota\), the balanced lower-boundary pairs have
\(|\delta|=1/2\).  For \(\tau\ge1\), the first nonzero lower pairs
have \(|\delta|=1\).  For \(\tau=0\), the first interior value is
\(\sqrt3\).  The untwisted boundary gives \(3/2\).
All higher \(L\) increase the square.  These observations and the
representation dimensions prove every gap and multiplicity in the table.
The two zero grades follow from the previously computed index \(d\).
\end{proof}

\begin{corollary}[All-isospin compact circle products]
\label{p1b:cor:minimal-all-isospin-circle}
Supply a product metric, a connection pulled back from the sphere with no
circle component, a circle radius \(\ell\), and spin structure
\(s\in\{0,1/2\}\).  For every spatial signed eigenvalue \(\lambda\)
and Fourier label \(n\), the product roots are
\[
 \pm\sqrt{\lambda^2+(n+s)^2/\ell^2},
\]
each with that spatial multiplicity, or twice that multiplicity at a zero
label.  If the spatial kernel has dimension \(h>0\), the periodic circle
gives \(2h\) zeros, \(h\) in each chirality, and nonzero gap
\(\min(g,1/\ell)\).  The antiperiodic circle gives exact gap
\(1/(2\ell)\), multiplicity \(2h\) per sign.  If \(h=0\), the gap is
\(\sqrt{g^2+s^2/\ell^2}\).  The general labelled rule also specifies
all ties and multiplicities.  Circle-dependent connections require their
actual coupled operator and are outside this Fourier-product assertion.
\end{corollary}
\begin{proof}
The product operator is \(D\otimes\sigma_1+(n+s)I\otimes\sigma_2/\ell\).
The auxiliary matrices anticommute, giving the stated square.  Minimize
over the exact circle frequencies and use the spatial zero and gap data.
\end{proof}

\subsubsection{General invariant zero responses and necessary extra premises}
\begin{proposition}[The exact analytic condition for a single protected line]
\label{p1b:prop:minimal-born-simple-line}
For the native fundamental \((+,-)\) lift, let \(D_{A,-}\) be the Dirac block from
negative to positive reciprocal grade.  Then
\[
 \dim\ker D_A=1\quad\Longleftrightarrow\quad\ker D_{A,-}=0.
\]
Thus injectivity of this block is the exact analytic condition for a single
line.  A small-connection norm bound is only a sufficient certificate.
At any simple point, compact ellipticity supplies a positive nonzero gap
and the local inverse and jet formulas below apply with that point as
reference.  For smooth \(SU(2)\) connection families at fixed metric,
the simple-kernel line is real, its Berry curvature is zero, and its global
Berry holonomy lies in \(\{+1,-1\}\).  A global positive orientation
on an arbitrary parameter loop does not follow from the local statement.
\end{proposition}
\begin{proof}
The index gives \(h_+-h_-=1\), hence total kernel dimension \(1+2h_-\).
At a simple point zero is isolated by ellipticity on the compact sphere.
The spin and fundamental quaternionic structures tensor to a real structure
commuting with the Dirac operator.  Local normalized real eigenvectors have
zero Berry one-form; their transition functions are signs.  This proves
flatness and the stated holonomy restriction.  The resolvent construction
uses the reference reduced inverse and a perturbation smaller than its gap;
this produces a local atlas covering every simple point.
\end{proof}

\begin{theorem}[Complete local response of the fundamental protected line]
\label{p1b:thm:minimal-born-zero-jets}
On the round fundamental bundle, fix the native lift and write
\(D_A=B_R+V\), where \(V=c(A-A_B)\) is self-adjoint and odd under \(U\).
Put \(g=\sqrt3/R\), \(P=|\psi_0\rangle\langle\psi_0|\), \(Q=I-P\),
and \(G=Q(B_R|_Q)^{-1}Q\).  If \(v=\|V\|<g\), there is exactly
one zero, all other absolute eigenvalues are at least \(g-v\), and
\[
 \psi=(I+GV)^{-1}\psi_0
       =\sum_{n\ge0}(-GV)^n\psi_0,\qquad P\psi=\psi_0.
\]
Writing \(\rho=v/g<1\), the order-\(N\) tail is at most
\(\rho^{N+1}\|\psi_0\|/(1-\rho)\).
For a smooth parameter family let \(\mathcal R=(I+GV)^{-1}\).
Every existing parameter jet is determined recursively by
\[
 \partial^\alpha\psi=-\mathcal R
 \sum_{0<\beta\le\alpha}{\alpha\choose\beta}
 G(\partial^\beta V)(\partial^{\alpha-\beta}\psi).
\]
The normalized line has zero Berry curvature on this small-connection
ball, and admits a globally zero Berry connection in its real normalized
frame.  Physical sources and parameter-dependent gauge frames still retain
their complete source and contact transport.
\end{theorem}
\begin{proof}
Naturality in the reciprocal-index theorem makes every such connection
difference odd.  The odd-block singular-value estimate of
Theorem~\ref{p1b:thm:dirac-reciprocity} gives the simple zero and gap
throughout \(v<g\).  Since \(\|G\|=g^{-1}\), the Neumann series
converges.  Both \(G\) and \(V\) are odd, so the constructed vector is
positive-grade; \(PV\psi=0\), which makes the complementary equation
the full zero equation.  Differentiating it gives the recursion.

The three-dimensional spin doublet and fundamental twisting doublet each
have their parallel quaternionic structure.  Their tensor product is a
parallel real structure \(C\) with \(C^2=1\), commuting with \(D_A\)
and \(U\).  Choose \(C\psi_0=\psi_0\); \(G,V,\mathcal R\) preserve
reality.  The normalized real vector has
\(\langle\widehat\psi,d\widehat\psi\rangle=0\), proving the Berry
statement on this contractible ball.  The standard transported-operator
contact \((dU)U^{-1}\) is retained when the frame itself varies.
\end{proof}

\begin{proposition}[Sharp limits of the minimal condition]
\label{p1b:prop:minimal-born-countermodels}
The odd spatial bit forces neither the midpoint connection nor its numerical
scalar gap, circle spin, physical action, source normalization or boundary
domain.  In particular there is no positive scalar-gap lower bound uniform
over all invariant connections in the odd class, despite full holonomy of
every member.
\end{proposition}
\begin{proof}
The profile \(u(r)=\epsilon q_1\), \(r=(q_0,q_1,q_2,q_3)\), is
inversion odd and gives a nonhomogeneous invariant connection.  The energy
formula of Theorem~\ref{p1b:thm:v2217-nonconstant-cartan-dirac} gives
\(\mathcal E=1/4+\epsilon^4/32\), using
\(\langle q_1^2\rangle=1/4\),
\(\langle q_1^4\rangle=1/8\) and
\(\langle|\nabla q_1|^2\rangle=3/4\).  It is thus not even gauge
equivalent to the midpoint for nonzero \(\epsilon\).

For the stronger scalar-gap countermodel, the square root
\(w(r)=r^{1/2}\) is smooth off \(-1\) and trivializes the native lift
there.  The connection \(A_{\rm flat}=w^{-1}dw\) is flat and invariant
on that punctured sphere.  Interpolate invariantly between \(A_B\) near
\(-1\) and \(A_{\rm flat}\) outside a ball of radius \(2\varepsilon\),
using an inversion-even cutoff of angular radius.  This gives a smooth invariant \(A_\varepsilon\).
A unit associated section parallel outside the ball, cut to zero inside
radius \(2\varepsilon\) and to one outside \(3\varepsilon\), has
Rayleigh quotient \(O(\varepsilon/R^2)\).  Its derivative is supported
where the connection is flat; a three-dimensional shell contributes
\(O(\varepsilon^3)\) volume and \(O(\varepsilon^{-2})\) squared
derivative.  Its norm tends to a nonzero constant.  Hence the first scalar
eigenvalues tend to zero, while no individual member has a scalar kernel.

Rescaling \(R\), changing the circle spin, multiplying the selected action
by a positive constant, or replacing a source \(B\) by \(2B\) all preserve
the carrier and lift while changing the respective output.  The last
replacement multiplies \(B^*(z-H)^{-1}B\) by four.  Distinct closed boundary
conditions likewise give different boundary spectra on the same interior
bundle.  These supply concrete independence witnesses for the extra
premises; none is a refutation of a source theorem retaining those premises.
\end{proof}

\begin{corollary}[The Born-fixed Yang--Mills tangent excludes the Morse escape]
\label{p1b:cor:minimal-born-morse-parity}
Under the selected round pure Yang--Mills Hessian and gauge quotient of
Theorem~\ref{p1b:thm:native-midpoint-gauge-morse}, its unique negative
direction is odd under exchange.  On the exchange-fixed physical tangent,
the Hessian is positive and bounded below by \(1/R^2\).
The full fluctuation space still has the stated one negative mode.
\end{corollary}
\begin{proof}
The negative line is the constant Cartan-profile variation.  Exchange sends
that variation to its negative.  The Hessian and gauge projection commute
with exchange, so the fixed tangent excludes this line.  The source theorem
has no quotient zero and all remaining eigenvalues at least \(1/R^2\).
This restricts a tangent space; it does not impose Born-fixed fields on the
complete physical action or select a dynamical history.
\end{proof}

\begin{theorem}[Closure under all typed mathematical consumers]
\label{p1b:thm:minimal-provider-consumer-closure}
Take any well-typed mathematical constructor word formed from the residual
provider and explicitly supplied additional data using pullbacks,
restrictions, associated representations, tensor operations, equivariant
intertwiners, projected connections, covariant differentiation, compatible
elliptic operators, functional calculus, supplied source contractions and
regular reductions.  An isomorphism of its consumed data transports the
entire word.  Every defined spectral quantity, trace word, response,
residue, parameter jet and source pairing is preserved with its own domain
and convergence hypotheses.  Consequently all later mathematical consumers
of this provider can use this Paper-I-B construction without re-assuming
an isomorphism of an already completed datum.

This is a universal naturality statement, not permission to discard a
consumer's other inputs or introduce a physical constructor.  If a later
word additionally consumes a metric, local realization, quantization member,
functional, source, pairing or boundary domain, that input remains necessary
until its actual framework owner supplies it.
\end{theorem}
\begin{proof}
Induct on the finite typed constructor word.  Each listed operation commutes
with the stated isomorphisms; for projected connections use
Proposition~\ref{p1b:prop:providers-projection}, and for closed operators
transport their graphs and domains.  Functional calculus and regular inverse
operations commute with unitary conjugation on their declared domains.
Retain full source columns and insertion operators before contractions.
Differentiation includes the frame-contact commutator, so induction also
holds for every existing jet.  This proves closure for all finite typed
compositions, without claiming that a topology-only hypothesis supplies the
extra arguments of those compositions.
\end{proof}


\begin{theorem}[Sharp midpoint covariant scalar gaps in every isospin]
\label{p1b:thm:conditional-residual-scalar-gaps}
Supply a round relative sphere of radius \(R>0\).  On the spin-\(\iota\)
associated bundle with the midpoint connection, the covariant scalar
Laplacian has complete labelled values
\[
 R^2\lambda_{jk}=2j(j+1)+2k(k+1)-\iota(\iota+1),
 \quad |j-\iota|\le k\le j+\iota,
\]
with multiplicity \((2j+1)(2k+1)\).  Its exact floor is
\[
 \boxed{\lambda_{\min,\iota}=\frac{\iota+\delta_\iota}{R^2},\qquad
 \delta_\iota=\begin{cases}0,&\iota\in\Z_{\ge0},\\
                           \tfrac14,&\iota\in\tfrac12+\Z_{\ge0}.
             \end{cases}}
\]
The floor multiplicity is \((\iota+1)^2\) for integral \(\iota\), and
\(2(\iota+\tfrac12)(\iota+\tfrac32)\) for half-integral \(\iota\).
In particular the fundamental and adjoint floors are
\(3/(4R^2)\) and \(1/R^2\), both of multiplicity four.
On \(S_R^3\times S_\ell^1\), with the stated product metric and no extra
circle holonomy, add \(n^2/\ell^2\); the floor is unchanged.
\end{theorem}
\begin{proof}
Set \(f=1/2\) in Theorem~\ref{p1b:thm:covariant-spectrum}.  For a second
minimization argument put \(L=j+k\), \(d=j-k\).  Clebsch--Gordan implies
\(L-\iota\in\Z_{\ge0}\), and
\[
 R^2\lambda_{jk}
  =\iota+(L-\iota)(L+\iota+2)+d^2.
\]
For integral \(\iota\), the minimum is attained exactly at
\(j=k=\iota/2\).  For half-integral \(\iota\), it is attained at the
two balanced pairs with \(L=\iota\) and \(d=\pm1/2\).  Any larger
\(L\) adds a positive term.  Multiplying the two representation dimensions
gives the stated multiplicities.  The circle contributes its independent
nonnegative Fourier square.
\end{proof}

\begin{corollary}[Compact residual Dirac problem and its circle-spin dependence]
\label{p1b:cor:conditional-residual-dirac}
Supply the orthogonal product \(S_R^3\times S_\ell^1\), one fundamental
twisting doublet, a connection pulled back from the sphere with zero circle
component, and a scalar-circle spin structure
\(\sigma\in\{0,1/2\}\).  Let \(D_u\) be the spatial Hodge/Cartan
operator of Theorem~\ref{p1b:thm:dirac-reciprocity}, with smooth reciprocal-odd
\(u\), \(\|u\|_\infty\le1\).  The compact product operator has blocks
\[
 \mathscr D_{u,n}=D_u\otimes\sigma_1+
                 \frac{n+\sigma}{\ell}I\otimes\sigma_2,
 \qquad
 \mathscr D_{u,n}^2=D_u^2\otimes I+
                        \frac{(n+\sigma)^2}{\ell^2}I.
\]
For every signed spatial eigenvalue \(\lambda\) of multiplicity
\(d_\lambda\), the nonzero product values are
\(\pm\sqrt{\lambda^2+(n+\sigma)^2/\ell^2}\), each of multiplicity
\(d_\lambda\) for that full label.  A zero label has multiplicity
\(2d_\lambda\).

For \(\sigma=0\), there are exactly two zero modes, one of each product
chirality.  At \(u=0\) the first positive absolute value is
\(\min\{\sqrt3/R,1/\ell\}\).  For \(\sigma=1/2\), the operator is
invertible and its exact gap is \(1/(2\ell)\), with multiplicity two per
sign.  These are compact geometric Dirac statements; they do not identify
the gap with a physical rest mass, choose a clock continuation, or infer a
physical scalar-circle spin structure.
\end{corollary}
\begin{proof}
The two auxiliary Pauli matrices anticommute, so the square has no mixed
term.  Tensor the complete spatial eigenbasis with the circle Fourier basis;
each remaining two-by-two matrix has the stated roots.  The spatial zero is
unique by Theorem~\ref{p1b:thm:dirac-reciprocity}.  The periodic circle has
one zero frequency and the antiperiodic circle has two minimizing frequencies
\(\pm1/(2\ell)\).  The auxiliary chirality \(\sigma_3\) splits a zero
label into one vector of each chirality.  At the midpoint the spatial
nonzero gap is exactly \(\sqrt3/R\).
\end{proof}

\begin{corollary}[The conditional Yang--Mills and protected-source branch]
\label{p1b:cor:conditional-residual-morse}
Supply, in addition to the residual datum, the round \(S_R^3\), the
classical pure \(SU(2)\) Yang--Mills energy and one fundamental twisting
doublet.  The mathematical conclusions of
Theorems~\ref{p1b:thm:native-midpoint-gauge-morse} and
\ref{p1b:thm:native-morse-dirac-alignment} apply to this specified datum:
the gauge quotient has one negative mode \(-2/R^2\), no zero mode and
positive spectrum beginning at \(1/R^2\).  Its homogeneous scalar is
the unique homogeneous gauge-quotient direction acting nontrivially on the
protected Dirac line.  On the Cartan family, with the source's normalized
energy and topological orientation,
\[
 \lambda=\frac{3/2-3f}{R},\qquad
 N_{\rm CS}=\frac12-\frac{R\lambda}{2}
                         +\frac{2R^3\lambda^3}{27},\qquad
 \mathcal E_{\rm YM}=\frac14
                    \left(1-\frac{4R^2\lambda^2}{9}\right)^2.
\]
The upward-positive spectral flow along \(0\le f\le1\) is \(-1\).
These conclusions retain the extra functional and metric premises; they
do not supply a physical action normalization or a Lorentzian history.
\end{corollary}
\begin{proof}
The explicit gauge in
Theorem~\ref{p1b:thm:conditional-residual-descent} identifies the actual
connection and fundamental doublet with the inputs of the named Morse
and source theorems.  Their closed operator and gauge-quotient problems
are transported together.  Substitute
\(f=1/2-R\lambda/3\) into \(3f^2-2f^3\) and
\(4f^2(1-f)^2\), as in
Corollary~\ref{p1b:cor:native-dirac-cs-barrier-coordinate}.  The simple
branch decreases through zero once; no other branch crosses in that interval.
\end{proof}

\begin{proposition}[Connection families, source transport and the readout condition]
\label{p1b:prop:conditional-residual-consumers}
The same global gauge transports every specified profile connection as
\[
 u\,r_v^*(f\theta_r)u^{-1}-du\,u^{-1}
       =f\theta_x+(1-f)\theta_p.
\]
It is exchange invariant exactly when \(f+f\circ B=1\); among constant
profiles this is the midpoint.  The space of exchange-invariant connections
on a specified coherent bundle is nonempty and affine, obtained by averaging
any connection with its exchange transport.  The two central signs induce
the same action on this connection space.

For a supplied operator/source/boundary problem, transport the entire
consumed tuple by the induced unitary \(U=\rho(u)\), including the closed
domain, sources and any grading or insertion.  Then its spectra, resolvents,
trace words and source contractions agree in the two frames.  In a smooth
parameter family, with \(\Theta=(dU)U^{-1}\),
\[
 d(UTU^{-1})-[\Theta,UTU^{-1}]=U(dT)U^{-1}.
\]
Consequently all existing transported jets agree.  An observable descends
through an unspecified descent fibre whenever it factors through the same
complete consumed tuple.  Dependence only on the underlying topological
bundle is insufficient to identify connections, source couplings, boundary
domains, or parity-resolved states.
\end{proposition}
\begin{proof}
The connection formula follows from the same Maurer--Cartan calculation as
Theorem~\ref{p1b:thm:conditional-residual-descent}.  Exchanging the two forms
gives the profile condition.  Connections form an affine space, so finite
group averaging gives the invariant subspace; central multiplication acts
trivially on adjoint-valued differences.  Unitary graph transport preserves
the full operator problem.  Differentiation gives the displayed identity,
and equality of the complete scalar words gives equality of their ordinary
derivatives.  Finally, if an observable is \(F\circ\operatorname{For}\),
it is constant on each fibre of that same forgetting map by substitution.
The map must forget only data that the observable does not consume.
\end{proof}

\subsection{Conditional principal fibres: finiteness and the trivial-component carry}
\label{p1b:sec:conditional-principal-fibres}

For a general principal bundle \(E\to M_8\), the minimal descent
antecedent is
\[
 \mathsf D_E:\qquad \operatorname{Desc}(E)\ne\varnothing.
\]
It is not replaced by \(c_2(E)=0\), by \(\epsilon_{\rm IB15}=0\), or by
stable triviality.  Under \(\mathsf D_E\), the two central signs give
coherent principal lifts as in Theorem~\ref{p1b:thm:g21-principal-descent}.
The all-Chern exotic restriction and quaternionic-index obstructions remain
necessary.  The following conditional calculation uses the exact principal
mapping-space reduction already stated in this paper.  Its topological input
theorems are retained as explicit premises; no value of an unevaluated map is
chosen.

\begin{proposition}[Necessary filters on the conditional principal branch]
\label{p1b:prop:conditional-descent-filters}
Suppose \(\mathsf D_E\) holds, with the specified determinant trivialization.
Its underlying second Chern class has the exchange-invariant form
\[
 c_2(E)=A(a_tu_x+a_Eu_p)+B(a_tu_p+a_Eu_x).
\]
The restriction bit is \(\epsilon_{\rm IB15}=0\), and
\[
 12\mid B^2-A^2,\qquad 4\mid A+B,\qquad
 24\mid B^2-A^2+3(A+B).
\]
Equivalently, on the ordinary principal existence domain, both pure
equivariant-index multiplicities
\[
 m_\pm=\frac{B^2-A^2}{12}\pm\frac{A+B}{4}
\]
are even.  These are necessary filters, not a sufficient descent test.
The mixed \(U(2)\) counts \(72/120\) do not count this branch.  For the
specified residual datum, \(A=B=0\), and all these filters hold; its actual
global gauge, rather than these filters alone, establishes its descent.
\end{proposition}
\begin{proof}
A quotient pullback is exchange invariant.  The ordinary rank-two
existence constraint, the principal restriction no-go and
Theorem~\ref{p1b:thm:g21-index-obstruction} give the first filters.
The invariant quaternionic structure in
Corollary~\ref{p1b:cor:g21-quaternionic-index-obstruction} makes both
virtual multiplicities even.  Their two numerators differ by \(6(A+B)\);
on the ordinary existence domain, the displayed mod-twenty-four condition
also forces \(4\mid A+B\).  This proves the arithmetic equivalence and
retains the source's necessary-versus-sufficient distinction.
\end{proof}

\begin{proposition}[Finiteness with weaker topological premises]
\label{p1b:prop:conditional-primary-finiteness}
Finiteness of a fixed-bundle descent fibre does not require the complete
intermediate gauge-group calculation.  It already follows from the finite
eight-dimensional quotient, the torsion-free group
\(H^4(Q;\Z)=\Z^2\), and the standard finite homotopy groups
\[
 \pi_5(BSU(2))=C_2,\quad \pi_6(BSU(2))=C_2,\quad
 \pi_7(BSU(2))=C_{12},\quad \pi_8(BSU(2))=C_2.
\]
These are the low homotopy inputs retained throughout the principal
calculation; no value of \(z_0\) or carry enters this argument.
\end{proposition}
\begin{proof}
An equivariant triangulation and barycentric subdivision of the compact
product make its finite-group quotient a finite eight-dimensional CW
space.  Rational cochains of the quotient identify with the invariant
cochains; averaging over the finite group is exact over \(\Q\).
Consequently
\(q^*:H^4(Q;\Q)\to H^4(M_8;\Q)\) is injective.  The displayed
torsion-free integral group makes the integral pullback injective too.
Thus a fixed \(c_2(E)\) permits at most one quotient second Chern class.

For a fixed representative of that primary class, lift its map to
\(K(\Z,4)\) through \(c_2:BSU(2)\to K(\Z,4)\).  The homotopy
fibre is four-connected, and its homotopy groups in degrees five through
eight are the finite groups displayed above.  Obstruction theory on a
finite eight-dimensional CW complex has finitely many choices at each
remaining stage, since its cohomology with each finite coefficient group
is finite.  It therefore gives finitely many lifts, even before quotienting
by homotopies of the primary map.  Empty lift sets cause no problem.
Hence every fixed-\(c_2\) quotient class set, and its subset
\(\operatorname{Desc}(E)\), is finite.  The stronger numerical bounds
below use the additional exact intermediate reduction.
\end{proof}

\begin{lemma}[The nonzero Chern increment of the free loop direction]
\label{p1b:lem:conditional-free-chern-increment}
Let \(A=\operatorname{SP}^2S^3\) and let
\(p:Q\to A\) forget the two scalar coordinates.  If
\(a\in H^1(S^1;\Z)\) is the scalar-product circle class and
\(\gamma\in H^3(A;\Z)\) is normalized by
\(q_Y^*\gamma=u_x+u_p\), put
\[
 \zeta=\mu^*a\,p^*\gamma\in H^4(Q;\Z).
\]
Then
\[
 \boxed{q^*\zeta=(a_t+a_E)(u_x+u_p)\ne0.}
\]
Retain the endpoint sequence and abstract splitting of
Proposition~\ref{p1b:v215:prop:F7-gauge-extension} and
Theorem~\ref{p1b:v215:thm:F7-gauge-split}.  In a component of second-Chern
coefficient \(d\), let \(g_d=\gcd(3,d)\) and \(k_d=3/g_d\).
A primitive free loop changes the quotient second Chern class by
\(4k_d\zeta\), up to a consistent overall orientation.  Torsion loops
do not change that class.  Therefore a prescribed pullback second Chern
class fixes at most one free integer in that component's affine twisted-loop
set, independently of its carry bit.
\end{lemma}
\begin{proof}
Naturality gives \(q^*\mu^*a=a_t+a_E\) and
\(q^*p^*\gamma=u_x+u_p\).  Their product is nonzero in the exterior
cohomology ring.  The restriction of \(p^*\gamma\) to \(F_7\) is the
primitive degree-three class: on either end \(A\) the forgetting map is the
identity.  Thus \(\zeta\) is its Wang transgression, with the scalar-product
circle as base.

For a loop in the classifying mapping space, the change of \(c_2\) under
clutching is the transgression of its degree-three class.  Restricting this
class to either end detects its coefficient, since
\(H^3(F_7;\Z)\to H^3(A;\Z)\) is an isomorphism.  The actual endpoint
sequence sends a primitive free generator to the common endpoint power
\(k_d\); a degree-one endpoint power has sphere degree four by
Lemma~\ref{p1b:lem:pure-descent-symmetric-loop-image}.  Its transgression
is therefore \(4k_d\zeta\).  A torsion loop maps to zero in the free group
\(H^3(F_7;\Z)\).  Pullback preserves the nonzero free increment, so two
different free integers cannot have the same pullback \(c_2\).
\end{proof}

\begin{theorem}[Finite fixed-bundle descent fibres under the principal reduction]
\label{p1b:thm:conditional-principal-finiteness}
Retain the actual double-cylinder component classification, endpoint
sequence, splitting and monodromy reduction of
Theorems~\ref{p1b:v215:thm:F7-components},
\ref{p1b:v215:thm:F7-gauge-split} and
\ref{p1b:thm:v226-principal-two-types}.  Every fixed-bundle descent fibre
\(\operatorname{Desc}(E)\) is finite.  If it is nonempty, write
\[
 c_2(E)=A(a_tu_x+a_Eu_p)+B(a_tu_p+a_Eu_x),\qquad
 d=A-B,\qquad g_d=\gcd(3,d).
\]
Then \(d\) is even.  In each invariant \(F_7\) component \(x\), let
\(c_x\in\{0,1\}\) be its carry bit.  Its contribution at prescribed
pullback \(c_2\), and hence at a prescribed full bundle \(E\), has at most
\(2^{3-c_x}\) elements.  In particular
\[
 \begin{array}{c|c|c}
 \text{actual endpoint normalization}&|\operatorname{Desc}(E)|
                  &|\mathscr L_{SU(2)}(E)|\\\hline
 z_0=2&\le32g_d&\le64g_d\\
 z_0=8&\le64g_d&\le128g_d
 \end{array}
\]
and, under \(\mathsf D_E\),
\[
 \boxed{1\le|\operatorname{Desc}(E)|\le192,\qquad
 2\le|\mathscr L_{SU(2)}(E)|
             =2|\operatorname{Desc}(E)|\le384.}
\]
These are bounds, not evaluations of the individual fibres or a choice
between \(z_0=2\) and \(8\).
\end{theorem}
\begin{proof}
The labelled cover of a scalar-product fibre is
\(q_0:S^1\times(S^3\times S^3)\to F_7\), represented by
\((e^{i\phi},x;e^{-i\phi},y)\).  For the Mayer--Vietoris generator
\(\eta=\delta(u_x)\in H^4(F_7;\Z)\), the two half-circles give
\(q_0^*\eta=a_\phi(u_x-u_p)\): the second half reverses the cylinder
and exchanges the two sphere factors.  Restricting \(c_2(E)\) to this
cover gives \((A-B)a_\phi(u_x-u_p)\), so its fibre coefficient is
\(d=A-B\), with the stated orientation.

The component monodromy of
Proposition~\ref{p1b:v215:prop:F7-monodromy} has no fixed component when
\(d\) is odd.  When \(d\) is even it has \(4g_d\) fixed components for
\(z_0=2\), and \(8g_d\) for \(z_0=8\).  In one fixed component, the
affine twisted-loop difference group is
\(\Z\oplus C_2^{3-c_x}\).  The preceding lemma shows that fixing the
pullback \(c_2\) fixes at most one free integer.  Hence at most
\(2^{3-c_x}\le8\) quotient classes can contribute there, and imposing
the full bundle class can only reduce that set.  Sum these bounds over the
fixed components and use \(g_d\le3\).  Finally
Theorem~\ref{p1b:thm:g21-principal-descent} gives the two inequivalent
central signs.  Nonemptiness gives the lower bounds.
\end{proof}

\begin{theorem}[An invariant primitive free generator in the trivial component]
\label{p1b:thm:trivial-component-zero-carry}
Let \(x_0\) be the constant-map component of
\(\operatorname{Map}(F_7,BSU(2))\), with the constant joining path at its
basepoint.  The actual involution on
\(\Gamma_0=\pi_1\operatorname{Map}(F_7,BSU(2))_{x_0}\) has zero carry.
Consequently
\[
 \boxed{\operatorname{coker}(\Theta-I)
                    \cong\Z\oplus C_2^3.}
\]
This calculation does not require choosing the value of \(z_0\).
\end{theorem}
\begin{proof}
Define the continuous forgetting map
\[
 p_F:F_7\longrightarrow A,\qquad
 p_F[\phi,x,y]=\{x,y\},\qquad p_F|_{A_\pm}=\mathrm{id}_A.
\]
The formulas agree at both seams of the double mapping cylinder, and
\(p_FR=p_F\) exactly.  Choose any based map
\(f_4:A\to SU(2)\) whose coordinate-sphere degree is four; its existence
is Lemma~\ref{p1b:lem:pure-descent-symmetric-loop-image}.  Under
\[
 \Gamma_0\cong[F_7,SU(2)],
\]
the map \(f_4p_F\) defines a loop \(z\).  Its two endpoint restrictions
are both \(f_4\).  Since \(d=0\) gives \(k_d=1\), the endpoint sequence
sends \(z\) to a primitive common free endpoint power.  Its kernel is
finite, so \(z\) is a primitive free generator of \(\Gamma_0\).
At the constant map and constant joining path, \(\Theta\) is
precomposition by \(R\).  The exact identity \(f_4p_FR=f_4p_F\) gives
\(\Theta z=z\).  Thus the carry vanishes for this explicit primitive
free generator.  Replacing it by another free lift changes the carry only
by \(\operatorname{im}(\Theta-I)\) on torsion, so its invariant bit is
zero.  The rank-two torsion-monodromy theorem then gives the displayed
cokernel.  The proof evaluates an actual invariant free map; it does not
infer the carry from the identity bundle alone.
\end{proof}

\begin{corollary}[Trivial-component module, cohomology and fixed-Chern bound]
\label{p1b:cor:trivial-component-descent-consequences}
Retain the intermediate gauge-group input of the preceding theorem.  As a
\(C_2\)-module, \(\Gamma_0\) is isomorphic to a trivial free \(\Z\)
summand, two length-two unipotent blocks over \(\F_2\), and one trivial
\(C_2\) summand.  In particular
\[
 \Gamma_0^{\Theta}\cong\Z\oplus C_2^3,\qquad
 H^{2r+1}(C_2;\Gamma_0)\cong C_2\quad(r\ge0),\qquad
 H^{2r}(C_2;\Gamma_0)\cong C_2^2\quad(r\ge1).
\]
These are cohomology groups of the intermediate loop group, not the
nonabelian \(H^1(C_2;\mathcal G(E)_\Phi)\) of the full gauge group on
\(M_8\).

At fixed pullback second Chern class, this trivial \(F_7\) component
contributes at most eight quotient bundles and at most sixteen principal
coherent lifts when both central signs are counted.  Under \(\mathsf H_{\rm vert}\), the specified weak-role
residual datum chooses the trivial quotient and the \(+I\) sign, rather
than this whole mathematical family.  Other \(F_7\) components may still
contribute to the descent fibre of the same underlying bundle.
\end{corollary}
\begin{proof}
The invariant primitive free generator gives an equivariant free/torsion
splitting.  On \(T=C_2^5\), the map \(N=\Theta-I\) has
\(N^2=0\) and rank two.  Its Jordan blocks therefore have sizes
\(2,2,1\); \(\ker N\) and \(T/\operatorname{im}N\) have dimension
three.  The periodic cochain complex alternates \(\Theta-I\) and
\(I+\Theta\).  On the free summand these maps are zero and multiplication
by two; on torsion both equal \(N\).  Thus positive odd cohomology has
one \(C_2\), while positive even cohomology has that \(C_2\) together
with the free summand modulo two.  The fixed-Chern bound follows from
Lemma~\ref{p1b:lem:conditional-free-chern-increment} and the eight-element
torsion quotient.  The actual residual quotient origin is the explicit one
in Theorem~\ref{p1b:thm:conditional-residual-descent}.
\end{proof}

\begin{remark}[Scope of the remaining calculation]
\label{p1b:rem:conditional-descent-frontier}
The native residual \(+I\) descent is already proved by its actual
constructor and global gauge, and later Papers~III-B and III-F supply its
physical role.  The trivial-component carry is now evaluated by an explicit
invariant primitive map.  The generic endpoint normalization, carries in
other components, and the actual ordered-cover pullback of torsion classes
remain separate classification questions.  The finiteness and upper bounds
above do not evaluate those maps.  No scalar trace, family mass, global
boundary gap, \(1/96\) action insertion, or cap source is inferred merely
from the bundle descent.  Such a consumer must retain its actual operator,
source, pairing and domain data.
\end{remark}


\part{Reciprocal operations and exact boundary constructions}
\section{\texorpdfstring{$M_8$}{M8}-inherited dynamic operations on the finite \texorpdfstring{$S^3_1$}{S3 level-one} package}
\label{p1:sec:v188-m8-s31-dynamic-operations}

The notation $S^3_1$ continues to mean the round momentum sphere together
with the declared level-one finite carrier package.  It is not a second
manifold factor inside $M_8$.  The following statements concern operations
obtained by restricting group words of the $SU(2)_p$ factor.

\begin{theorem}[Exceptional support defect residues]
\label{p1:thm:v188-exceptional-defect-residues}
For every retained exceptional support $\Gamma\in\{2T,2O,2I\}$ the $M_8$
order-swap word restricts to $\Gamma\times\Gamma\to\Gamma$, and its generated
image is the derived subgroup.  Consequently
\[
 [2T,2T]=Q_8,\qquad [2O,2O]=2T,\qquad [2I,2I]=2I,
\]
with
\[
 [Q_8,Q_8]=\{\pm1\},\qquad [\{\pm1\},\{\pm1\}]=1.
\]
The residual abelian quotients are therefore
\[
 \boxed{(2T)_{\rm ab}=2T/Q_8\cong C_3,\qquad
        (2O)_{\rm ab}=2O/2T\cong C_2,\qquad
        (2I)_{\rm ab}=1.}
\]
Hence $2I$ is the unique exceptional support whose full finite support is
regenerated by order-swap defects, whereas the nonabelian defect of $2O$
falls immediately into the common tetrahedral core.
\end{theorem}

\begin{proof}
The order-swap word is the group commutator, so its generated image is the
derived subgroup.  Proposition~\ref{p1:prop:exceptional-abelianizations}
already proves $[2T,2T]=Q_8$, $[2O,2O]=2T$ and $[2I,2I]=2I$, including the
central quaternion witness.  The standard quaternion presentation of $Q_8$
gives $[Q_8,Q_8]=\{\pm1\}$, and the latter is central abelian.  Taking
successive quotients gives the stated residual groups.
\end{proof}

\begin{corollary}[The trimer clock is the abelian residue of the Hurwitz core]
\label{p1:cor:v188-trimer-as-2t-abelianization}
For an oriented body diagonal $u$, the Hurwitz unit
\[
 c_u=\cos\frac\pi3+u\sin\frac\pi3
\]
has order six, $c_u^3=-1$, and
$C_6(u)\cap Q_8=\{\pm1\}$.  The quotient map $2T\to2T/Q_8$ therefore restricts
to an isomorphism
\[
 \boxed{C_6(u)/\{\pm1\}\ \cong\ 2T/Q_8\ \cong C_3.}
\]
Thus the cyclic trimer action used in the level-one package is not an
unrelated $C_3$: after an oriented diagonal has been generated it is a concrete
representative of the abelianized residue of the $2T$ order-swap defect.
Changing the diagonal orientation exchanges the two generators.
\end{corollary}

\begin{proof}
The quotient homomorphism $C_6(u)\to2T/Q_8$ has kernel
$C_6(u)\cap Q_8=\{\pm1\}$.  Its image has order three.  Since
$|2T/Q_8|=3$, the induced map is an isomorphism.  Reversing the oriented
axis replaces the order-three generator by its inverse.
\end{proof}

\begin{corollary}[Cartan square pairs the two oriented tetrahedral presentations]
\label{p1:cor:v188-square-pairs-2t-orientations}
The parent Cartan relation $g\mapsto g^2$ descends to the $2T$ abelianization.
On $C_3\cong2T/Q_8$ it is
\[
 \boxed{x\mapsto x^2=x^{-1}.}
\]
Hence it exchanges the two oriented $C_3$ generators appearing in
$Q_8\rtimes_\varphi C_3$ and $Q_8\rtimes_{\varphi^{-1}}C_3$.  It pairs the
two orientations but does not select one of them.
\end{corollary}

\begin{proof}
For the quotient homomorphism $q:2T\to2T/Q_8$, normality of
$Q_8$ gives $q(g^2)=q(g)^2$, so the square is representative-independent
in the quotient and descends to $2T/Q_8$.  In a cyclic group of order three every
nontrivial element satisfies $x^2=x^{-1}$, which exchanges the two oriented
generators.
\end{proof}

\begin{corollary}[Unique exceptional real sign residue]
\label{p1:cor:v188-unique-exceptional-sign-residue}
Among $2T,2O,2I$, only $2O$ has a nontrivial real one-dimensional sign
character:
\[
 \Hom(2T,\{\pm1\})=0,\qquad
 \Hom(2O,\{\pm1\})\cong C_2,\qquad
 \Hom(2I,\{\pm1\})=0.
\]
Indeed every such character factors through the abelianization, and only
$(2O)_{\rm ab}\cong C_2$ admits a nontrivial map to $\{\pm1\}$.  This gives a
dynamic group-theoretic reason that the normalizer branch, rather than the
$2T$ core or $2I$ selector, owns the exceptional sign/parity residue.
\end{corollary}

\begin{proof}
Any homomorphism to $\{\pm1\}$ factors through the abelianization.  The
preceding theorem gives $C_3$, $C_2$ and the trivial group respectively, and
only $C_2$ has a nontrivial homomorphism to $\{\pm1\}$.
\end{proof}

\begin{remark}[Carrier firewall]
\label{p1:rem:v188-operation-carrier-firewall}
Closure of a retained support under inverse, square, multiplication or
commutator does not enlarge the declared level-one representation carrier
$\mathcal H_1$ and does not replace the condition $m_p=1$.  These group words
are parent operations/relations.  Products of amplitudes may require a larger
comparison algebra, but they are not new momentum states merely because the
parent group operation exists.
\end{remark}

\section{Normal-jet realization of the oriented boundary quotient}
\label{p1:sec:v357-normal-jet}
\begin{theorem}[Exact two-channel normal-jet sequence]
\label{p1:thm:v357-normal-jet-exact}
Adopt the previously declared, orientation-compatible splitting
$X=\CC\eta_\sigma\oplus F_\sigma$ of
Definition~\ref{p1:def:oriented-boundary-face}.  Treat $\eta_\sigma$
as the formal normal coordinate in the corresponding symmetric algebra.
For every $r\ge3$ define the linear \emph{comparison} map
\[
 J_r:\Sym^r(X)\longrightarrow
 \Sym^r(F_\sigma)\oplus\Sym^{r-2}(X),\qquad
 J_r(p)=\left(p\big|_{\eta_\sigma=0},
                \partial_{\eta_\sigma}^{\,2}p\right).
\]
Then $J_r$ is surjective and
\[
 \ker J_r=\eta_\sigma\Sym^{r-1}(F_\sigma)=\mathcal L_r^\sigma.
\]
In particular the following sequence is exact:
\[
 0\longrightarrow\mathcal L_r^\sigma
 \longrightarrow\mathcal B_r\xrightarrow{J_r}
 \Sym^r(F_\sigma)\oplus\Sym^{r-2}(X)
 \longrightarrow0,
\]
so $\mathcal N_r\cong\Sym^r(F_\sigma)\oplus\Sym^{r-2}(X)$
as complex vector spaces.  The normal-coordinate basis fixes the
particular formula, while its kernel is invariant under rescaling that
basis. The source-defined degree-three onset and its zero boundary in
degrees $0,1,2$ are retained unchanged.
\end{theorem}
\begin{proof}
Every $p\in\Sym^r(X)$ has the unique homogeneous expansion
$p=\sum_{k=0}^r\eta_\sigma^kp_k$ with
$p_k\in\Sym^{r-k}F_\sigma$.  Restriction keeps precisely $p_0$.
Formal second differentiation keeps precisely the independent
components $k\ge2$, with nonzero coefficients $k(k-1)$.
Thus $J_rp=0$ holds exactly when $p=\eta_\sigma p_1$ with
$p_1\in\Sym^{r-1}F_\sigma$.  Conversely, a prescribed tangential
polynomial $a$ and homogeneous $(r-2)$-polynomial $b$ have preimage
$a+\int_0^{\eta_\sigma}\int_0^{u}b(v,F_\sigma)\,dv\,du$,
where the integrals are formal polynomial antiderivatives. This
proves surjectivity and exactness. The quotient assertion follows
from the first isomorphism theorem.
\end{proof}
\begin{remark}[No physical source or algebra multiplication is generated]
\label{p1:rem:v357-normal-jet-scope}
$J_r$ is a mathematical witness for the already defined degreewise
quotient. It is not the QCD-to-nuclear source map $S_r$, a nuclear
Hamiltonian, a normal-response constitutive law, or a proof that the
onset convention is selected by baryon number.  By
Proposition~\ref{p1:prop:boundary-quotient-not-algebra}, the
$\mathcal L_r^\sigma$ do not constitute a graded ideal; exactness of
each $J_r$ supplies no multiplication on the direct sum of quotients.
\end{remark}


\section{Rooted Dirac boundaries and exact conormal responses}
\label{p1b:sec:v229-rooted-boundary}
The results in this section concern the displayed geometric operators,
source maps and closed domains. Their use by a physical action consumes
those same data. Normal Clifford multiplication obeys
$c_n^\dagger=-c_n$, $c_n^2=-I$. All unbounded identities are understood
on the specified invariant core and then on the associated closed forms.

\begin{theorem}[Transport of the full source and boundary domain]
\label{p1b:thm:v229-boundary-transport}
Let $\gamma$ and $\Gamma$ be parallel Hermitian gradings, with
$\{D_M,\gamma\}=0$ and $\{\Gamma,\Phi\}=0$ for a Hermitian finite field.
Set
\[
 D_1=D_M\otimes I+\gamma\otimes\Phi,\qquad
 D_2=D_M\otimes\Gamma+I\otimes\Phi,\qquad
 U=P_\gamma^+\otimes I+P_\gamma^-\otimes\Gamma.
\]
Then $U=U^\dagger=U^{-1}$, $UD_1U=D_2$ and
$U(c_n\otimes I)U=c_n\otimes\Gamma$.
For $\chi_a=c_n\gamma\otimes I$ and
$\chi_b=c_n\gamma\otimes\Gamma$, one has
$U\chi_aU=\chi_b$ and $U\chi_bU=\chi_a$.
Consequently these presentations define equivalent closed boundary
problems when the whole domains are related by $U$.
\end{theorem}
\begin{proof}
Use $D_MP_\gamma^\pm=P_\gamma^\mp D_M$ and
$\Gamma\Phi\Gamma=-\Phi$ in the two external chiral blocks.
The same multiplication proves the normal-symbol and projector identities.
Unitary transport of the graph of an operator transports its closed domain.
Keeping the same untransported projector in both presentations instead
specifies a different boundary problem.
\end{proof}

\begin{proposition}[Source conormal isometry]
\label{p1b:prop:v229-conormal-isometry}
For $(D_1,\chi_b)$ write $\Psi=\gamma\otimes\Phi$ and
$J_n=c_n\Psi=(c_n\gamma)\otimes\Phi$. Then
$J_n^\dagger=J_n$, $J_n^2=I\otimes\Phi^2$ and
$\{\chi_b,J_n\}=0$. If $\Phi^2=m^2I$, $m>0$, and
$\Pi_\pm=(I\pm\chi_b)/2$, the map $V=J_n/m$ satisfies
\[
 (V\Pi_+)^\dagger V\Pi_+=\Pi_+,\qquad
 V\Pi_+(V\Pi_+)^\dagger=\Pi_-.
\]
For noncentral $\Phi^2$, its polar analogue holds on the nonzero support;
the original positive source Gram must still be retained.
\end{proposition}
\begin{proof}
The factors $c_n\gamma$ and $\Phi$ are commuting Hermitian matrices,
$(c_n\gamma)^2=I$, and $\Gamma$ anticommutes with $\Phi$.
Thus $V$ is a Hermitian unitary exchanging the two projectors.
The polar assertion follows by spectral calculus on the nonzero support.
\end{proof}

\begin{proposition}[The bulk square does not determine a boundary gap]
\label{p1b:prop:v229-robin-zero-mode}
On a frozen flat product collar, the mass part of the Robin operator
inherited from $\Pi_-\psi=\Pi_-D\psi=0$ is
$S_\Phi=-\Pi_+c_n\Psi\Pi_+$. It vanishes for $(D_1,\chi_b)$,
whereas for $(D_1,\chi_a)$ it equals $-\Phi\Pi_+$.
On the half line $y\ge0$, with $\partial_n=-\partial_y$ and
$\Phi=mK$, $K^2=I$, the latter problem has a normalizable zero mode
in a zero tangential channel although its bulk square is
$-\partial_y^2+m^2$.
\end{proposition}
\begin{proof}
Multiply the second inherited boundary condition by $c_n$.
For the example choose $\chi_av=v$, $Kv=v$. Then $c_nv=-\gamma v$,
so $\psi=e^{-my}v$ obeys $D_1\psi=0$ and
$(\partial_n-\Phi)\psi=0$. A periodic tangential torus supplies an
actual constant tangential section. This is a counterexample for the
stated product problem, not a claim about every global boundary.
\end{proof}

\subsection{Parent matching, leakage and a surviving gap}
\begin{proposition}[Matched quaternionic source overlap]
\label{p1b:prop:v229-parent-overlap}
Suppose two midpoint doublets are identified by the parallel intertwiner
$c(u)$ for a unit quaternion $u$. Their matched unit singlets obey
\[
 \langle\xi_L,(I\otimes c(h))\xi_R\rangle
 =\tfrac12\operatorname{Tr}(c(u)^\dagger c(h))=\langle u,h\rangle.
\]
If this overlap is spatially constant, the squared complementary weight
of the unit source is $1-|\langle u,h\rangle|^2$.
For $u=(1,1,1,1)/2$, $h=-e_1$, these numbers are $-1/2$ and $3/4$.
\end{proposition}
\begin{proof}
The singlet reduced doublet density is $I/2$. Take the partial trace
and use the quaternionic Clifford Gram. Unitarity gives the complementary
weight by subtraction. The common-frame invariance is also
Corollary~\ref{p1:cor:v229-common-frame-overlap}.
\end{proof}
A different actual identification first requires its parallel intertwiner
$T_{\rm par}$; its overlap is $\operatorname{Tr}(T_{\rm par}^\dagger
c(h))/2$. The numerical fraction above is not asserted independently
of this matching.

\begin{theorem}[Rooted gap with its domain hypothesis]
\label{p1b:thm:v229-rooted-gap}
In the matched realization, decompose the odd completion as
$K_h=rK_s+K_v$, where $r=-1/2$, $K_s^2=I$,
$K_v^2=3I/4$ and $\{K_s,K_v\}=0$. Let
$S_s=\gamma\otimes K_s$, $S_v=\gamma\otimes K_v$ and
$\mathscr D=\mathscr B+m(rS_s+S_v)$ for constant $m>0$.
If $\mathscr D_\perp=\mathscr B+mS_v$ is self-adjoint on a common
domain preserved by $S_s$, and $\{\mathscr D_\perp,S_s\}=0$ there, then
\[
 \mathscr D^2=\mathscr D_\perp^2+m^2I/4,
 \qquad |\mathscr D|\ge mI/2.
\]
The compatible total-graded product bag preserves $S_s$ when $K_s$
is parallel. The result is not extended to arbitrary bags or varying $m$.
\end{theorem}
\begin{proof}
Expand the square on the common core. The mixed term vanishes.
Closure gives $\|\mathscr D\psi\|^2=
\|\mathscr D_\perp\psi\|^2+m^2\|\psi\|^2/4$ and the spectral bound.
\end{proof}

\begin{proposition}[Complete constant singlet--triplet return]
\label{p1b:prop:v229-rooted-return}
On the frozen source-active constant round-$S^3_a$ sector, $g=2/a$,
the inherited two-partner square and eigenvalues are
\[
 A(m,g)=\begin{pmatrix}m^2&\sqrt3gm/2\\\sqrt3gm/2&m^2+g^2\end{pmatrix},
 \quad \lambda_\pm=m^2+g^2/2\pm\tfrac g2\sqrt{g^2+3m^2}.
\]
Although its singlet compression is $m^2$, its square-root compression is
\[
 \frac{m^2+\sqrt{m^4+g^2m^2/4}}
 {\sqrt{g^2+2m^2+2\sqrt{m^4+g^2m^2/4}}}.
\]
It tends to $m/2$ as $g/m\to\infty$ and is $\sqrt3$ for $m=2,g=3$.
\end{proposition}
\begin{proof}
The off-diagonal source weight is $\sqrt3/2$ by the preceding overlap.
The characteristic polynomial has trace $2m^2+g^2$ and determinant
$m^4+g^2m^2/4$. For a positive two-by-two matrix,
$\sqrt A=(A+\sqrt{\det A}I)/\sqrt{\operatorname{Tr}A+2\sqrt{\det A}}$,
as follows by squaring and using its characteristic polynomial.
\end{proof}

\begin{theorem}[Exact conormal spectral response]
\label{p1b:thm:v229-conormal-response}
On a frozen tangential realization let $\mathscr D$ be self-adjoint,
$A=\mathscr D^2>0$, $\mathscr BP=0$ and
$\mathscr DP=mS_0P$. Suppose $c_n$ is unitary and anticommutes with
$\mathscr D$ as an operator on its tangential domain. For
$J_0=c_nS_0P$ and every spectral function with the required source
and quadratic-form domains,
\[
 J_0^\dagger f(A)J_0=m^{-2}PAf(A)P.
\]
For the product extension $-\partial_y^2+A$ on an interval of length $L$,
with respectively Neumann and Dirichlet far termination, this implies
\[
 K_N=\sqrt A\tanh(L\sqrt A),\quad K_D=\sqrt A\coth(L\sqrt A),
 \qquad J_0^\dagger K_N^{-1}J_0=m^{-2}PK_DP.
\]
\end{theorem}
\begin{proof}
The assumed operator anticommutation gives $[c_n,f(A)]=0$.
Substitute $S_0P=\mathscr DP/m$ on both sides of the source form and
commute $\mathscr D$ with $f(A)$. Solving the scalar normal ODE gives
the two impedances, and $AK_N^{-1}=K_D$ proves the last formula.
\end{proof}
Here $A$ is the tangential operator defining the extension. The normal
symbol exchanges bag sectors; its domain invariance is not inferred
for an arbitrary full normal boundary Dirac operator. The complete
source-column trace retains the actual $1/4,3/4$ partner covariance.

\section{Geometry forcing a homothetic boundary response}
\label{p1b:sec:v229-response-scaling}
\begin{proposition}[Common-return Robin homothety]
\label{p1b:prop:v229-robin-homothety}
For $-\partial_y^2+k^2$, $k>0$, let $kL=\ell$ be fixed and impose
at the far end $q'(L)+k\zeta q(L)=0$, with fixed $\zeta\ge0$.
The near-end impedance is
\[
 \kappa=k\frac{\zeta+\tanh\ell}{1+\zeta\tanh\ell}.
\]
Thus $k_R=Gk_0$, $L_R=L_0/G$ gives $\kappa_R=G\kappa_0$.
\end{proposition}
\begin{proof}
Write $q(y)=a\cosh(ky)+b\sinh(ky)$, impose the far condition and
evaluate $-q'(0)/q(0)$. Equivalently the entire positive interval form
scales by $G$ after $z=Gy$. A fixed dimensionful Robin constant is
different data and need not have this scaling.
\end{proof}

\begin{theorem}[Exact support criterion with fixed tangential geometry]
\label{p1b:thm:v229-fixed-tangent-support}
Let $F_R=-\partial_y^2+\mu^2G^2+T$ on $[0,L_0/G]$, where
$\mu>0$, $G>1$, and $T\ge0$ is fixed self-adjoint, with the split
product Neumann termination. Set $\mathcal R=K_0^{-1/2}K_RK_0^{-1/2}$.
For complete source columns $M$ with $\operatorname{Tr}M^\dagger M=1$
and the required inverse and trace domains,
\[
 E_* =\frac{w^2}{2}\operatorname{Tr}
 [M^\dagger(I+\mathcal R)^{-1}M]
 =\frac{w^2}{2(1+G)}
 \quad\Longleftrightarrow\quad \operatorname{Ran}M\subseteq\ker T
\]
when $w\ne0$.
\end{theorem}
\begin{proof}
For $T$ eigenvalue $\lambda$, $K_R/G=f(\mu^2+\lambda/G^2)$ and
$K_0=f(\mu^2+\lambda)$, with $f(s)=\sqrt s\tanh(L_0\sqrt s)$
strictly increasing. Hence $\mathcal R\le GI$, with equality space
$\ker T$. The positive difference
$(I+\mathcal R)^{-1}-(1+G)^{-1}I$ has the same kernel.
Its complete source trace vanishes exactly on that kernel.
The displayed energy follows by minimizing the two flux forms under
their sum constraint. No commutation of $M$ with $T$ is required.
\end{proof}
This split theorem does not discard the $D\Phi$ cross terms of a
different actual Dirac square.

\begin{theorem}[Necessary and sufficient coupled scaling]
\label{p1b:thm:v229-coupled-homothety}
For $A_R=A(\mu G,g_R)$ from
Proposition~\ref{p1b:prop:v229-rooted-return}, $L_R=L_0/G$ and
$g_R,g_0>0$, the body-frame identity satisfies
\[
 K_{N,R}=GK_{N,0}\quad\Longleftrightarrow\quad g_R=Gg_0.
\]
On a round factor this requires $a_R=a_0/G$. If all remaining metric,
connection and mass constituents obey the corresponding frozen
homothety, the whole tangential operator scales, not only this block.
\end{theorem}
\begin{proof}
$K_{N,R}/G=f(A_R/G^2)$ for the strictly increasing function in the
preceding proof. Its scalar inverse defines an inverse by spectral
calculus, so equality implies $A_R/G^2=A_0$. Compare the second
diagonal entry, then the off-diagonal entry. The converse is immediate.
Ordinary round Dirac homothety scales the entire spectrum under its
unitary density transport. A whole product requires all its factors.
\end{proof}

\begin{proposition}[Affine conormal covariance]
\label{p1b:prop:v229-affine-conormal}
Suppose $\Gamma_W$ commutes with $c_n$, $\mathscr B$ and $P$, and
anticommutes with $S_0$. Put $S_1=i\Gamma_WS_0$,
$\Psi_R=\mu S_0+\gamma_NR S_1$, $\gamma_N=\mu/R_s$,
$G=\sqrt{1+(R/R_s)^2}$ and $U_R=e^{i\theta\Gamma_W/2}$,
$\theta=\arctan(R/R_s)$. Then
\[
 \delta\mathcal J_R=c_n(\Psi_R-\Psi_0)P=\gamma_NR J_1,
 \quad J_1=-iJ_0\Gamma_W,
 \quad J_1J_1^\dagger=J_0J_0^\dagger=:\mathcal C,
\]
$J_0J_1^\dagger+J_1J_0^\dagger=0$ and $[\mathcal C,\Gamma_W]=0$.
Consequently $\delta\mathcal J_R\delta\mathcal J_R^\dagger
=\gamma_N^2R^2\mathcal C$. Under whole frozen homothety,
$\mathscr D_R=GU_R\mathscr D_0U_R^\dagger$ and
$K_{N,R}=GU_RK_{N,0}U_R^\dagger$.
\end{proposition}
\begin{proof}
Anticommutation gives $S_1=-iS_0\Gamma_W$. Multiply the resulting
source identities and use $\Gamma_W^2=I$. The two cross covariances
cancel, and $\mathcal C$ is invariant under $U_R$. Finally
$\Psi_R=\mu G U_RS_0U_R^\dagger$; the assumed whole background
scaling gives the operator and impedance statements.
\end{proof}
The source statement uses every original column. It constructs an
operator map; its interpretation as a scalar action source is a
separate application. For nonconstant $R$, the connection contact
$U_R^\dagger\nabla U_R=i\Gamma_Wd\theta/2$ remains present.

\section{Exact tests separating cap and fixed positive heat profiles}
\label{p1b:sec:v229-heat-cap-tests}
\begin{theorem}[Gapped fixed-cutoff heat profiles cannot equal the cap]
\label{p1b:thm:v229-heat-tail}
Let $A_0\ge aI>0$ be self-adjoint with compact resolvent, $W\ge0$
bounded and $0<\operatorname{Tr}(We^{-\tau A_0})<\infty$ for fixed
$\tau>0$. Assume the heat trace is analytic for positive times.
Then
\[
 F_\tau(s)=\frac{\operatorname{Tr}(We^{-\tau(1+s)A_0})}
 {\operatorname{Tr}(We^{-\tau A_0})}
 \le e^{-\tau as},\qquad s\ge0,
\]
and $F_\tau$ cannot equal $H(s)=2/(1+\sqrt{1+s})$ on any open
positive interval. A compact visible zero eigenspace instead gives a
positive constant limit; absent that weight, the first visible positive
eigenvalue gives the exponential bound again.
\end{theorem}
\begin{proof}
The spectral resolution gives a positive probability measure with rates
$t=\tau\lambda\ge\tau a$, even if $W$ does not commute with $A_0$.
Integrating $e^{-ts}$ proves the bound. The cap tends to $2s^{-1/2}$.
Real analyticity extends any open-interval equality to the connected
positive domain, contradicting the tails. Compact resolvent gives the
stated alternatives when zero modes are admitted.
\end{proof}
This theorem concerns the specified positive heat word. It does not
exclude arbitrary physical field reductions or indefinite causal forms.

\begin{proposition}[Finite heat rank test for the cap]
\label{p1b:prop:v229-cap-moment-rank}
The positive cap density of Proposition~\ref{p1b:prop:v224-heat-cap}
has moments $m_n=\Gamma(n+1/2)/(\sqrt\pi(n+1))$. In particular,
\[
 \det\begin{pmatrix}1&1/4&1/4\\1/4&1/4&15/32\\
 1/4&15/32&21/16\end{pmatrix}=71/1024>0.
\]
No positive two-rate heat mixture equals the cap. More generally no
finite-rate heat mixture equals it.
\end{proposition}
\begin{proof}
Tonelli applied to the strictly positive density gives the moments in
all degrees. A two-atom measure has moment matrix rank at most two;
direct determinant expansion gives the displayed nonzero value.
For any nonzero polynomial $p$, $\int p(t)^2\rho(t)dt>0$, so the cap
has positive moment matrices of every size, whereas a finite atomic
measure has finite rank. This proves the general assertion.
\end{proof}


\subsection{Balanced finite heat normalization and source ports}
\label{p1b:sec:v250-balanced-heat-source-boundary}
\begin{proposition}[Exact retained normalization of a balanced closed finite odd heat]
\label{p1b:prop:v250-balanced-odd-heat-normalization}
Let $\mathcal H=\mathcal H_+\oplus\mathcal H_-$ be an actually retained
fixed finite carrier, let $\Gamma=I_+\oplus(-I_-)$, and suppose
$D=D^\dagger$ satisfies $\{D,\Gamma\}=0$. For any scalar function $f$
defined on its finite squared spectrum,
\[
 \boxed{\operatorname{Tr}(\Gamma f(D^2))
       =(\dim\mathcal H_+-\dim\mathcal H_-)f(0).}
\]
If the actual carrier is balanced and $P_+=(I+\Gamma)/2$, then the
primitive trace normalization of
Proposition~\ref{p1b:prop:v224-retained-primitive} therefore obeys
\[
 \boxed{\frac1{48}\operatorname{Tr}(P_+e^{-D^2/\Lambda_g^2})
       =\frac1{96}\operatorname{Tr}(e^{-D^2/\Lambda_g^2}).}
\]
The identity holds throughout every smooth fixed balanced relative-odd
finite family, and hence for all of its existing parameter/source jets,
including their actual common cutoff and scalar coefficient derivatives.
It covers the uninserted heat constituent only.
\end{proposition}
\begin{proof}
Write $D=\left(\begin{smallmatrix}0&A^\dagger\\A&0\end{smallmatrix}\right)$.
Every positive eigenvalue of $A^\dagger A$ is paired with the same
multiplicity in $AA^\dagger$ by its normalized $A$ map. Only the zero
eigenspaces contribute to the graded trace, and their dimension difference
is $\dim\mathcal H_+-\dim\mathcal H_-$. This gives the first formula.
In the balanced case its right side is zero; expand $P_+$ to obtain the
second formula. Differentiating an identity holding throughout the stated
fixed family gives its jets; no insertion is added or omitted.
\end{proof}

\begin{proposition}[The same marked trace is not a uniform normalization on all ports]
\label{p1b:prop:v250-marked-trace-not-whole-action-scaling}
On a nonzero balanced carrier of the preceding proposition define
\[
 \tau_{\rm rp}(L)=\frac1{48}\operatorname{Tr}(P_+Le^{-D^2/\Lambda_g^2}),
 \qquad
 \tau_{\rm unif}(L)=\frac1{96}\operatorname{Tr}(Le^{-D^2/\Lambda_g^2}).
\]
They agree at $L=I$ but are different linear functionals on the already
available relative-projector algebra:
\[
 \boxed{\tau_{\rm rp}(P_+)=2\tau_{\rm unif}(P_+)>0,
 \qquad \tau_{\rm rp}(P_-)=0<\tau_{\rm unif}(P_-).}
\]
Thus the uninserted all-jet heat identity cannot be extended to a common
$1/96$ rescaling of every source/port insertion by retaining this same
marked trace. Any physical equality on a smaller actually consumed port
family must be proved on that family with its actual field and boundary
maps.
\end{proposition}
\begin{proof}
$D^2$ reduces both projectors. Their heat traces agree by the preceding
pairing theorem and are strictly positive on this nonzero balanced carrier.
Use $P_+^2=P_+$ and $P_+P_-=0$ in the two definitions. This uses the same
$D$ on both sides, not a changed source or a candidate action.
\end{proof}

The balanced finite hypothesis must be checked on the actual action window.
Equal dimensions of the pointwise spin/internal fibres do not alone prove
this assertion for a global untruncated elliptic operator, whose graded heat
trace may contain an index. An actual finite compression must preserve the
graded operator domain and its displayed balanced carrier. The local joint
formula $D_A+\gamma_g\otimes\Phi$ is relative odd when its complete gauge
connection preserves $\Gamma_W$ and its actual scalar obeys
$\{\Phi,\Gamma_W\}=0$; the finite compression must retain these relations.
Boundary conditions which are not preserved by the grading are outside the
proposition. Neither result constructs the physical adoption of the marked
readout, a first-order fermionic pairing, or the cap's normal-source domain.
In particular the positive theorem does not remove the outstanding
whole-action condition in
Proposition~\ref{p3e:prop:v2211-whole-action-newton-certificate}.

\begin{proposition}[Hermitian marked Dirac form and the strongest closed finite candidate]
\label{p1b:prop:v250-marked-real-dirac-normalization}
On the same fixed balanced finite relative-odd carrier, with its fixed
ordinary complex Hilbert pairing, put $Z_{\rm rp}=P_+/48$. Then
\[
 \boxed{\operatorname{Herm}(Z_{\rm rp}D)
 :=\tfrac12(Z_{\rm rp}D+DZ_{\rm rp})=D/96,\qquad
 \operatorname{Re}\langle\psi,Z_{\rm rp}D\psi\rangle
                 =\langle\psi,D\psi\rangle/96.}
\]
Consequently, for any inherited real scalar weights $w,c$, the declared
closed finite real functional
\[
 \mathscr S_{\rm rp}^{\rm cand}(D,\psi)
 ={w\over48}\operatorname{Tr}(P_+e^{-D^2/\Lambda_g^2})
       +c\operatorname{Re}\langle\psi,Z_{\rm rp}D\psi\rangle
\]
is exactly
\[
 \boxed{\mathscr S_{\rm rp}^{\rm cand}(D,\psi)
 ={1\over96}\left[w\operatorname{Tr}(e^{-D^2/\Lambda_g^2})
                          +c\langle\psi,D\psi\rangle\right].}
\]
This equality includes all existing parameter and real field/source jets
on that shared fixed-pairing family. It is a finite real-form identity,
not an assertion of physical adoption of this candidate or of a Grassmann
reality convention.
\end{proposition}
\begin{proof}
Relative oddness gives $P_+D+DP_+=D$. Since $D$ and $Z_{\rm rp}$ are
Hermitian, the adjoint of $Z_{\rm rp}D$ is $DZ_{\rm rp}$; its Hermitian
part is therefore $D/96$. The real part of a diagonal ordinary Hilbert
quadratic form is the quadratic form of that Hermitian part. Combine this
identity with Proposition~\ref{p1b:prop:v250-balanced-odd-heat-normalization}
and differentiate the exact functional equality where its jets exist.
\end{proof}

The same-action consumer has not been inferred from this calculation.
The actual local fermionic theorem
\ref{p3f:thm:v177-local-clocked-fermion-jet} uses the inherited
\emph{uninserted} pairing and fixes $c_F$ jointly on its kinetic and
Yukawa terms; its insertion firewall is
Remark~\ref{p3f:rem:v177-fermion-insertion-quantum-firewall}.
Neither that theorem nor the primitive trace retraction asserts that
$Z_{\rm rp}$ acts on the first-order fermionic action and is then
Hermitian-symmetrized in precisely the real form used here. That typed
adoption remains the missing action map. A Lorentzian clock-hermitizing
form, independent ordered Grassmann variables, a Weyl restriction or an
endpoint pairing must use its actual real structure and domains before
this identity can be consumed there. No fermionic measure, determinant
phase or quantum equivalence at fixed operational $\hbar$ follows.
The candidate also contains no new cap normal source or boundary action;
its closed finite equality does not establish those missing incidences.


\subsection{Explicit rooted operator coordinates for forward action tests}
\label{p1b:sec:v251-rooted-action-coordinates}

The following maps are geometric operator coordinates, not a selection of
a physical odd-source coefficient. Retain the actual rooted family pair
\[
 R=\operatorname{diag}(1,i,-i),\quad
 c=(1,1,1)^T/\sqrt3,\quad P=cc^\dagger,\quad Q=I-P,
\]
and the existing hierarchy-family unitary $J_{HG}$.

\begin{proposition}[A complete coordinate map from the rooted pair]
\label{p1b:prop:v251-rooted-matrix-units}
For $\rho=(1,i,-i)$ put
\[
 \Pi_j=\prod_{k\ne j}\frac{R-\rho_kI}{\rho_j-\rho_k},
 \qquad E_{jk}=3\Pi_jP\Pi_k.
\]
These are matrix units: $E_{jk}^\dagger=E_{kj}$,
$E_{jk}E_{\ell m}=\delta_{k\ell}E_{jm}$ and $\sum_jE_{jj}=I$.
Consequently the rooted geometric algebra is all of $\operatorname{End}(G)$,
with the explicit inverse and hierarchy transport
\[
 M=\sum_{jk}\operatorname{Tr}(E_{kj}M)E_{jk},\qquad
 \widehat E_{jk}=J_{HG}E_{jk}J_{HG}^\dagger.
\]
\end{proposition}
\begin{proof}
The simple-spectrum polynomial projectors satisfy
$\operatorname{Tr}(P\Pi_j)=1/3$ and $P\Pi_jP=P/3$.
Insert these relations in the product of two displayed units. Their
diagonal sum is the spectral resolution of the identity. Trace duality
gives the coefficient inverse, and unitary conjugation preserves every
product and adjoint.
\end{proof}

\begin{proposition}[One polar word and its exact geometric return]
\label{p1b:prop:v251-single-polar-transfer}
Put $s^2=8/9$, $P_\mu=s^{-2}QRPR^\dagger Q$, $P_e=Q-P_\mu$, and
\[
 A_0=P_eRQ+QRP+PRQ,\qquad V_0=(A_0A_0^\dagger)^{-1/2}A_0.
\]
The positive factor is
$A_0A_0^\dagger=P_e+s^2(P_\mu+P)$, so $V_0$ is unitary. With
$B_e=P_eRQ$, $B_\mu=s^{-1}QRP$, $B_\tau=s^{-1}PRQ$,
\[
 V_0=B_e+B_\mu+B_\tau,\qquad
 M_{\rm geom}=(eP_e+mP_\mu+tP)V_0,
\]
\[
 M_{\rm geom}M_{\rm geom}^\dagger=e^2P_e+m^2P_\mu+t^2P.
\]
\end{proposition}
\begin{proof}
The rooted polar-Hopf identities give the three orthogonal left returns
and zero cross returns. Positive normalization of their respective
coefficients $1,s,s$ gives $V_0$. The weighted return follows by multiplying
the unitary out. This is an explicit geometric free-response comparison,
not an evaluation of the actual mixed fermionic action coefficient.
\end{proof}

An actual mass/current source can now be extracted in a complete finite
rooted coordinate system. The target-free construction of that source
still has to precede its comparison with $M_{\rm geom}$. In particular,
equivariance alone cannot supply the missing source word: the stabilizer
of $(R,P)$ is the scalar unitary group, so every fixed matrix $M_0$ defines
an equivariant map on its presentation orbit by simultaneous conjugation.
This is not a proof of multiple physically licensed source members.


\subsection{Forward neutral response from the full harmonic packet}
\label{p1b:sec:v251-neutral-forward-packet}
Let $\mathcal H=\mathcal H_{12}$ have normalized round Haar metric,
$K_{12}=168I$, and $E:\mathcal H\to\mathbb C^\Omega$ be the native
six-point evaluation. Write
\[
 EE^\dagger=234p+156p^\perp,\quad
 B=(EE^\dagger)^{-1/2}E,\quad
 P_1=B^\dagger pB,\quad P_5=B^\dagger p^\perp B,\quad
 Q=I-P_1-P_5.
\]
Here $p=|s\rangle\langle s|$, $s=\mathbf1_6/\sqrt6$, and the three
projector ranks are $1,5,163$. Let $h=h_{12}$ be the actual icosian
seed, with $\|h\|^2=13$, $Eh=13\mathbf1_6$ and $P_5h=0$.
Retain both generated derivative isometries
\[
 J_Lg=\frac{L_{J_{HG}g}h}{\sqrt{728}},\qquad
 J_Rg=\frac{R_{J_{HG}g}h}{\sqrt{728}}.
\]
The derivatives use the dimensionless curves $\exp(tw)q$ and
$q\exp(tw)$, for unit imaginary quaternions $w$. They are admitted
finite isometry-derivative words on the invariant harmonic band,
as in Paper~II-C's canonical transvection and all-band covariance
theorems. Their group labels do not yet identify physical fermion
chirality.

\begin{lemma}[The complete seed derivative packet]
\label{p1b:lem:v251-neutral-seed-derivative-packet}
The two derivative maps above satisfy
\[
 QJ_a=J_a,\quad J_a^\dagger J_a=I_G,\quad
 J_L^\dagger J_R=0,\quad J_a^\dagger h=0\qquad(a=L,R).
\]
\end{lemma}
\begin{proof}
In the thirteen-dimensional spin-six representation write
$h(q)=13\operatorname{Tr}(P_I\rho_{12}(q))$, with $P_I$ the icosian
Reynolds projector. Peter--Weyl orthogonality and $\|h\|^2=13$ imply
$\operatorname{rank}P_I=1$. With $T_w=d\rho_{12}(w)$, conjugation
invariance of $h$ under the Hurwitz subgroup makes the derivative
Gram scalar on the imaginary three-space: the half-turns about the
coordinate axes eliminate its off-diagonal entries and the cyclic
Hurwitz rotation makes its diagonal entries equal. The round
Casimir gives the trace of that Gram as
$168\|h\|^2=2184$, so each entry is $728$.

The linear functional $w\mapsto\operatorname{Tr}(P_IT_w)$ is invariant
under the same rotations, hence zero. Rank one then gives
$P_IT_wP_I=0$. For $g\in I$, $\rho_{12}(g)P_I=P_I$ and cyclicity
show that both derivatives of $h$ at $g$ equal
$13\operatorname{Tr}(P_IT_w)=0$. Since $\Omega\subset I$, both
derivative ranges lie in $\ker E$.

For $f_A(q)=13\operatorname{Tr}(A\rho_{12}(q))$, the coefficient
pairing is $13\operatorname{Tr}(A^\dagger B)$. The left and right
seed derivatives have coefficients $P_IT_w$ and $T_zP_I$.
Their cross inner product is
$13\operatorname{Tr}(T_w^\dagger P_IT_zP_I)=0$. Their inner products
with $h$ vanish by the same identity. Division by $\sqrt{728}$ and
transport by the existing unitary $J_{HG}$ prove the assertions.
\end{proof}

\begin{theorem}[Complete joint source Gram and spectral response]
\label{p1b:thm:v251-neutral-joint-source-response}
Define the source carrier and full joint column
\[
 \mathscr S=\mathbb C\oplus\mathbb C^\Omega\oplus G\oplus G,
 \qquad
 \mathcal J=(h/\sqrt{13},B^\dagger,J_L,J_R):\mathscr S\to\mathcal H.
\]
For every $F=f_1P_1+f_5P_5+f_0Q$, with complex scalars $f_a$, its
exact complete joint return is
\begin{equation}
\label{p1b:eq:v251-neutral-joint-spectral-response}
\boxed{
 \mathcal J^\dagger F\mathcal J=
 \begin{pmatrix}
 (f_1+2f_0)/3&f_1s^\dagger/\sqrt3&0&0\\
 f_1s/\sqrt3&f_1p+f_5p^\perp&0&0\\
 0&0&f_0I_3&0\\
 0&0&0&f_0I_3
 \end{pmatrix}.}
\end{equation}
In particular its source Gram $\mathsf G=\mathcal J^\dagger\mathcal J$
has spectrum $1+1/\sqrt3,1-1/\sqrt3,1$ with multiplicities $1,1,11$.
The joint source has rank thirteen, including seven bulk directions;
the orthogonal unseen complement has dimension $156$ and lies in $Q$.
\end{theorem}
\begin{proof}
The seed decomposition is
$h/\sqrt{13}=e_s/\sqrt3+\sqrt{2/3}\,e_b$, where
$e_s=B^\dagger s$ and $e_b=Qh/\sqrt{26/3}$ are orthonormal. The
five augmentation columns, the two derivative triples and $e_b$ are
mutually orthogonal, and each is orthogonal to $e_s$. Inserting the
three spectral projectors gives the matrix. Its only nondiagonal
subspace is the seed/boundary-singlet plane, with Gram
$\left(\begin{smallmatrix}1&1/\sqrt3\\1/\sqrt3&1\end{smallmatrix}\right)>0$.
The boundary contributes six directions and the bulk contributes
$1+3+3=7$. Since $\operatorname{rank}Q=163$, the unseen complement
has dimension $163-7=156$.
\end{proof}

This retains the unsummed boundary ports and both bulk derivative
packets. It does not replace the $168$-owner by the $163$-owner before
performing the common contraction. Raw ports are retained exactly by
\[
 \mathcal J_{\rm raw}=(h,E^\dagger,D_L,D_R)
 =\mathcal J\operatorname{diag}
 (\sqrt{13},(EE^\dagger)^{1/2},\sqrt{728}I_3,\sqrt{728}I_3),
\]
where $D_a=\sqrt{728}J_a$. Thus normalization has not discarded their
native source weights.

\begin{corollary}[The round forward response has no rank-mass splitting]
\label{p1b:cor:v251-neutral-round-response}
For the actual round kinetic, the entire joint Green response is
$\mathsf G/168$, and the joint resolvent is $\mathsf G/(z-168)$.
On the correctly normalized source space it is $I_{13}/168$ and
$I_{13}/(z-168)$, respectively. The geometric Gram eigenvalues are
source overlaps, not distinct pole energies. No squared spectrum
$\{1/163,5,163\}$ follows from this full round response.
\end{corollary}
\begin{proof}
Use the theorem with constant spectral function and normalize both
source legs by $\mathsf G^{-1/2}$. The range of $\mathcal J$ reduces
the round scalar operator. Every source direction has the same pole.
\end{proof}

\subsubsection{One complete quadratic parent distinguishes vector and trace ports}
\label{p1b:sec:v251-neutral-complete-quadratic-parent}
For positive numbers $k_1,k_5,k_0$ retain the full positive operator
\[
 K=k_1P_1+k_5P_5+k_0Q
\]
and the mathematically defined complex quadratic parent
\[
 S_K(\xi,j)=\langle\xi,K\xi\rangle
                  -2\operatorname{Re}\langle\xi,\mathcal Jj\rangle,
 \qquad \xi\in\mathcal H.
\]
Stationary elimination gives $\xi_*=K^{-1}\mathcal Jj$ and
$S_{\rm red}=-j^\dagger\mathcal J^\dagger K^{-1}\mathcal Jj$.
This is a geometric quadratic comparison parent, not an assertion of
new physical harmonic fields or a native neutrino action.

The exact finite complex Gaussian with its round Lebesgue measure is
\begin{equation}
\label{p1b:eq:v251-neutral-complete-gaussian}
 Z_K(j)=\int_{\mathcal H}e^{-S_K(\xi,j)}\frac{d^{338}\xi}{\pi^{169}}
 =\frac{\exp(j^\dagger\mathcal J^\dagger K^{-1}\mathcal Jj)}
 {k_1k_5^5k_0^{163}}.
\end{equation}
It derives vector Green responses and collective trace responses from
one supplied parent. For example, at $j=0$,
\[
 -\partial_{k_a}\log Z_K=\frac{n_a}{k_a},\qquad
 \partial_{k_a}\partial_{k_b}\log Z_K
 =\delta_{ab}\frac{n_a}{k_a^2},\quad (n_1,n_5,n_0)=(1,5,163).
\]
The vector inverse, trace inverse and connected quadratic-source
Hessian are differently typed objects. In particular the last contains
two inverse factors, rather than one. None is automatically an odd
mass coefficient.

The native nonsinglet/augmentation Green trace is also recovered from
this same $K$:
\[
 \begin{pmatrix}
 \operatorname{Tr}((Q+P_5)K^{-1}(Q+P_5))&
 \operatorname{Tr}((Q+P_5)K^{-1}P_5)\\
 \operatorname{Tr}(P_5K^{-1}(Q+P_5))&\operatorname{Tr}(P_5K^{-1}P_5)
 \end{pmatrix}
 =\begin{pmatrix}163/k_0+5/k_5&5/k_5\\5/k_5&5/k_5\end{pmatrix}.
\]
Its scalar Schur residual is $163/k_0$, while the orthogonal carrier
block $QK^{-1}P_5$ is zero. This preserves the full $168$-owner
cross return without declaring it an off-diagonal matter Hessian.

There is a further integration consequence. The thirteen-dimensional
joint visible determinant is $k_1k_5^5k_0^7$. The unseen complement
contributes $k_0^{156}$ to the full determinant. It is invisible to
the retained vector response but not to the complete generating
functional or its derivatives. At the Gaussian level, dropping that
complement would lose the term $-156\log k_0$ in $\log Z_K$.

\subsubsection{A concrete mixed derivative map supplies all family matrix units}
\begin{proposition}[Mixed left/right harmonic operator extraction]
\label{p1b:prop:v251-neutral-mixed-derivative-extraction}
In an orthonormal rooted imaginary frame, let $L_i,R_j$ be the
dimensionless Killing derivatives above. Then
\begin{equation}
\label{p1b:eq:v251-neutral-mixed-matrix-unit}
 J_L^\dagger L_iR_jJ_L=J_R^\dagger L_iR_jJ_R=0,\qquad
 \boxed{J_R^\dagger L_iR_jJ_L=-56|j\rangle\langle i|.}
\end{equation}
The right side is transported to the native family frame by $J_{HG}$.
Thus the already available mixed second-derivative words provide
explicit directed family matrix-unit maps. The isotropic word
$\sum_iL_iR_i$ gives the central cross return $-56I_3$.
\end{proposition}
\begin{proof}
Write $h(q)=13\operatorname{Tr}(P_I\rho_{12}(q))$ and
$T_i=d\rho_{12}(e_i)$, so $T_i^\dagger=-T_i$.
The seed/derivative identities imply
\[
 P_IT_iP_I=0,\quad
 P_IT_b^\dagger T_jP_I=56\delta_{bj}P_I,\quad
 \operatorname{Tr}(P_IT_aT_i)=-56\delta_{ai}.
\]
For $f_A(q)=13\operatorname{Tr}(A\rho_{12}(q))$, left differentiation
takes $A$ to $AT_i$ and right differentiation takes $A$ to $T_jA$.
Peter--Weyl pairing is $\langle f_A,f_B\rangle=13\operatorname{Tr}(A^\dagger B)$.
The required matrix element is therefore
\[
 \frac{13}{728}\operatorname{Tr}
 (P_IT_b^\dagger T_jP_IT_aT_i)
 =-56\delta_{bj}\delta_{ai}.
\]
For either same-port block one instead obtains a factor $P_IT_jP_I$
or $P_IT_iP_I$, which vanishes. This proves all identities.
\end{proof}

For any actual consumed operator $O$, the finite cross target is
$M=J_R^\dagger OJ_L$. Existence of the derivative map does not choose
the coefficients of $O$ in the physical action. It does show exactly
where an oriented, noncentral coefficient can enter. An actual word
containing only scalar spectral functions of $E^\dagger E$ has zero
cross return, and the isotropic mixed second-derivative word has a
central one.

For a supplied family coefficient $M$, there is a unique self-adjoint
extension of minimal Hilbert--Schmidt norm with only the required
off-diagonal support:
\[
 O_{\min}(M)=J_RMJ_L^\dagger+J_LM^\dagger J_R^\dagger,
 \qquad \|O_{\min}(M)\|_{HS}^2=2\|M\|_{HS}^2.
\]
This follows by orthogonal block decomposition relative to
$\operatorname{ran} J_L\oplus\operatorname{ran} J_R\oplus(\operatorname{ran} J_L\oplus\operatorname{ran} J_R)^\perp$.
It is a mathematical minimal extension, not a GP3 principle selecting $M$.
For this specifically projected extension,
\[
 J_R^\dagger(\kappa I+\varepsilon O_{\min}(M))^{-1}J_L
 =-\varepsilon M(\kappa^2I-\varepsilon^2M^\dagger M)^{-1}
\]
when $|\varepsilon|\|M\|<\kappa$. Full unprojected words $L_iR_j$
can have additional boundary and unseen couplings; their higher
returns must not be replaced by this projected model.

\subsubsection{Complete source jets and the remaining action calculation}
\label{p1b:sec:v251-neutral-complete-source-jets}
For a fixed retained source map, $R=K^{-1}$ obeys
\[
 \partial_i(\mathcal J^\dagger R\mathcal J)
 =-\mathcal J^\dagger RK_iR\mathcal J,
\]
\[
 \partial_i\partial_j(\mathcal J^\dagger R\mathcal J)
 =\mathcal J^\dagger(RK_iRK_jR+RK_jRK_iR-RK_{ij}R)\mathcal J.
\]
Moving maps and metrics additionally contribute their ordinary
product-rule/frame contacts. In a fixed transported Hilbert frame,
with $\mathcal J_i=\partial_i\mathcal J$, the full first jet is
\[
 \partial_i(\mathcal J^\dagger R\mathcal J)
 =\mathcal J_i^\dagger R\mathcal J+\mathcal J^\dagger R\mathcal J_i
      -\mathcal J^\dagger RK_iR\mathcal J.
\]
Its next derivative retains $\mathcal J_{ij}$ and every mixed
$\mathcal J_i,K_j$ product. At the round point, a perturbation
$K\mapsto\kappa I+\varepsilon O$ has
\[
 \left.\partial_\varepsilon
 [J_R^\dagger K^{-1}J_L]\right|_0=-M/\kappa^2.
\]
This gives a direct finite extraction test for the actual mixed operator.
Even if its first visible compression vanishes, the second jet can
contain $2\mathcal J^\dagger O\Pi_{\rm unseen}O\mathcal J/\kappa^3$.
All required jets therefore use the full retained operator or a proved
jet-complete reduction.

The positive outcome is an evaluated common harmonic response,
retaining the seed, six boundary ports, both derivative triples,
the $168/163$ distinction and the complete trace contribution.
The remaining physical task is to derive the native neutral source
operator and its action-owned relation to these ports, including
the current/kinetic/pole readout. The calculation does not install
the inverse/direct rank diagram by declaration and does not identify
spatial left/right derivatives with fermionic chiral inputs.


\begin{theorem}[Exact positive finite-word cubature from the two dense icosahedral completions]
\label{p1b:thm:v275-dense-twin-word-cubature}
Let $H\cong2T\le G=SU(2)$ be supplied, let $O=N_G(H)$,
let $I_+,I_-$ be its unordered icosahedral completions, and put
$D_H=\langle I_+,I_-\rangle$. For every finite-dimensional real
subspace $V\subset C(G,\mathbb R)$ containing $1$, with $d=\dim V$,
there is a positive probability
\[
 \nu_V=\sum_{j=1}^{s}w_j\delta_{q_j},\qquad
 1\le s\le d,\quad w_j>0,\quad q_j\in D_H,
\]
such that $\int f\,d\nu_V=\int f\,d\sigma$ for every $f\in V$.
The nodes are therefore finite words in the \emph{already specified}
unordered pair of completions. This is a mathematical existence
statement relative to $H$ and not a canonical node/weight selection
or an occurring finite measurement process.

If $V=\mathcal P_L$ is the real scalar polynomial band on $S^3$,
there exists in addition a finite positive $\mathcal P_L$-exact
measure $\widetilde\nu_{H,L}$ supported inside $D_H$ that is
left- and right-$H$-invariant and invariant under conjugation by
$O$. This construction makes no choice between $I_+$ and $I_-$.
It can be obtained from an initial rule with at most
\[
 d_L=\dim\mathcal P_L=(L+1)(L+2)(2L+3)/6
\]
nodes by finite $H$ bi-averaging and $O$-conjugation averaging.
\end{theorem}
\begin{proof}
Take a real basis $1,f_1,\ldots,f_{d-1}$ of $V$ and set
$F(q)=(f_j(q)-\int f_j\,d\sigma)_j$.
The origin lies in the \emph{interior} of
$\operatorname{conv}F(G)$ in $\mathbb R^{d-1}$. Indeed if it
lay on a supporting hyperplane, a nonzero linear combination
$\ell\cdot F$ would be nonnegative on $G$ with zero Haar integral.
Full support of Haar and continuity would force this combination
to vanish everywhere, contradicting linear independence modulo
constants. By
Theorem~\ref{p1:thm:v273-two-icosahedral-completions},
$D_H$ is dense in $G$. Approximate a finite collection of vertices
whose convex hull surrounds the origin by points from $F(D_H)$;
a sufficiently small perturbation retains the origin in their
convex hull. Carath\'eodory's theorem then yields at most $d$
points in $D_H$ and strictly positive weights realizing the origin.
The case $d=1$ is immediate.

Apply the first assertion to $\mathcal P_L$, then average the
resulting probability over left and right multiplication by $H$
and conjugation by $O$. Since $H\subset D_H$ and conjugation by
$O$ interchanges $I_+$ and $I_-$, all nodes remain in $D_H$.
The polynomial band is stable under these orthogonal actions and
Haar is invariant, so exactness is preserved. $H$ is normal in $O$,
so the two symmetrizations commute at the level needed to retain
all three invariances. Positivity and finiteness are preserved.
\end{proof}


\begin{corollary}[Finite observable data do not select unique exceptional-word supports]
\label{p1b:cor:v275-no-finite-observable-support-selection}
For every finite-dimensional real $V\subset C(SU(2),\mathbb R)$
containing constants, there exist infinitely many distinct positive
finite probabilities on $D_H$ that agree \emph{exactly} with Haar
on $V$ and whose supports are pairwise disjoint. For
$V=\mathcal P_L$, all these measures may simultaneously be
chosen $H$-bi-invariant and $O$-conjugation invariant.
Thus no fixed finite list of continuous observable expectations,
without an independently supplied support/occurrence constraint,
can identify a unique finite word sampler, even among such
symmetry-preserving positive presentations.
\end{corollary}
\begin{proof}
Removing any finite set from the dense subgroup $D_H$ leaves
a dense subset of the perfect compact group $SU(2)$.
Apply the interior-convex-hull proof of
Theorem~\ref{p1b:thm:v275-dense-twin-word-cubature}
to $D_H$ minus the union of the finitely many previously used
supports. Induction produces pairwise disjoint finite exact
positive rules. In the polynomial case, first replace the
forbidden set by its finite orbit under left and right $H$
multiplication and $O$ conjugation. Removing this orbit still
leaves a dense subset; averaging the new rule over the finite
symmetries preserves exactness and cannot enter the forbidden
invariant set. Induction proves the symmetric assertion.
\end{proof}

\begin{theorem}[Positive finite correction of the branch-symmetric twin convolution]
\label{p1b:thm:v275-positive-exactification-twin}
With $H,D_H$ as above, let
$\xi_{H,n}=\xi_H^{*n}$ be the branch-symmetric positive finite
word measures of
Theorem~\ref{p1b:thm:v274-branch-symmetric-twin-convolution}.
For every fixed degree $L$ there are positive finite word measures
$\eta_{H,L,n}$ and numbers $0\le\varepsilon_{L,n}<1$ such that
\[
 \nu_{H,L,n}=(1-\varepsilon_{L,n})\xi_{H,n}
                  +\varepsilon_{L,n}\eta_{H,L,n}
\]
is exactly Haar-cubature on $\mathcal P_L$ and
$\varepsilon_{L,n}=O(q_L^n)$ for some $q_L<1$.
Moreover $\nu_{H,L,n}$ may be chosen left/right $H$-invariant
and $O$-conjugation invariant, retaining exactness and the bound
\[
 \|\nu_{H,L,n}-\xi_{H,n}\|_{\rm TV,var}
       \le 2\varepsilon_{L,n}=O(q_L^n).
\]
All these finite supports are \emph{not} subgroups: they contain
$I_+$ and $I_-$, which generate a dense proper subgroup.

For $L\ge12$, no nontrivial convex mixture of finitely many
$\xi_{H,n}$ alone is exact even on $\mathcal P_{12}$: its
$\rho_{12}$ Fourier coefficient is strictly positive on $V^H$.
In particular the exactifying $\eta$ provides genuinely different
positive word weights, rather than another mixture of the same
symmetric convolution powers. For $L=19$ one may take
$q_{19}=525/8192$ in the stated asymptotic bound.
\end{theorem}
\begin{proof}
Using the centered coordinate map $F$ for $V=\mathcal P_L$
from the previous proof, the interior argument supplies a finite
set $T_L\subset D_H$ and $r_L>0$ with
$B_{r_L}(0)\subset\operatorname{conv}F(T_L)$.
Set $e_n=\int F\,d\xi_{H,n}$. The finite-band Fourier
contraction in the cited theorem gives $\|e_n\|=O(q_L^n)$.
For $e_n\ne0$, set
\[
 \varepsilon_{L,n}=\frac{2\|e_n\|}{r_L+2\|e_n\|},
 \qquad y_n=-\frac{1-\varepsilon_{L,n}}{\varepsilon_{L,n}}e_n.
\]
Then $\|y_n\|=r_L/2$ and $y_n$ is realized by a positive
finite probability $\eta_{H,L,n}$ on $T_L$. Therefore
$(1-\varepsilon)e_n+\varepsilon y_n=0$.
If $e_n=0$, take $\varepsilon=0$. Convolve the resulting
exact measure on left and right by $\mu_H$ and average its
$O$-conjugates. This preserves $\mathcal P_L$ exactness, retains
support in $D_H$, and contracts total variation. The reference
$\xi_{H,n}$ is already invariant under these operations; hence
its stated total-variation bound follows. Since
$\varepsilon<1$, the support still contains that of
$\xi_{H,n}$, in particular $I_+\cup I_-$. A finite subgroup
cannot contain both. At degree 12 the Fourier coefficients
of the $\xi_{H,n}$ have strictly positive eigenvalues
$(275/8192)^n$ and $(525/8192)^n$ on $V^H$; positive
convex combinations cannot make them zero.
For $L=19$, the only nonzero nonconstant Fourier coefficient
of $\xi_{H,n}$ is the displayed $k=12$ component, hence its
moment-vector error is $O((525/8192)^n)$.
\end{proof}

\begin{theorem}[Nested positive word resolutions and unequal-resolution partial isometries]
\label{p1b:thm:v275-nested-positive-resolution-tower}
For every fixed $H$ above, there is a sequence of positive finite
probabilities $\nu_m$ on $D_H$ such that
\[
 \operatorname{supp}\nu_m\subseteq\operatorname{supp}\nu_{m+1},
 \qquad \nu_m|_{\mathcal P_{2m}}=\sigma|_{\mathcal P_{2m}}
 \quad(m\ge0).
\]
They may all be chosen $H$-bi-invariant and $O$-conjugation
invariant. Their finite supports need not be subgroups, while
$\nu_m\rightharpoonup\sigma$ weakly, and the union of their
supports is dense in $G$.

For arbitrary positive rules $\alpha,\beta$ exact on
$\mathcal P_{2m}$ and $\mathcal P_{2n}$, respectively, with
$m\le n$, their weighted evaluation maps
$S_{\alpha}^m:\mathcal P_m\to\ell^2(\operatorname{supp}\alpha)$
and $S_\beta^n$ are isometries. Let
$J_{nm}:\mathcal P_m\hookrightarrow\mathcal P_n$ be the
canonical harmonic inclusion. Then
\[
 U_{\beta,n\leftarrow\alpha,m}
   =S_\beta^n J_{nm}(S_\alpha^m)^*
\]
is a partial isometry, with initial projection
$S_\alpha^m(S_\alpha^m)^*$ and final projection
$S_\beta^nJ_{nm}J_{nm}^*(S_\beta^n)^*$.
For $m\le n\le\ell$, with any third such rule $\gamma$,
\[
 U_{\gamma,\ell\leftarrow\beta,n}
 U_{\beta,n\leftarrow\alpha,m}
 =U_{\gamma,\ell\leftarrow\alpha,m}.
\]
These are comparison-frame identities, not physical evolution maps.
\end{theorem}
\begin{proof}
Assume $\nu_m$ constructed. For the centered finite moment
coordinates of $\mathcal P_{2m+2}$, the Haar moment is interior
to the convex hull of values on $D_H$ by the first theorem.
Choose $0<\theta<1$ sufficiently small that the residual target
$-\theta(1-\theta)^{-1}\int F\,d\nu_m$ lies in that convex
hull. A positive finite rule $\eta$ on $D_H$ realizes the
residual. Then $\nu_{m+1}=\theta\nu_m+(1-\theta)\eta$ is
exact at degree $2m+2$, and every old support atom retains
positive weight. Start at $m=0$ with $\nu_0=\mu_H$.
Symmetrizing the newly added rule by $H$ and $O$ preserves
exactness and includes the previous invariant support.
Polynomial density in $C(G)$ plus positivity yields weak
convergence. Since Haar has full support, a weakly convergent
nested finite-support sequence must have dense union.

Degree-$2m$ exactness gives
$\langle S_\alpha^mf,S_\alpha^mg\rangle
=\int\overline f g\,d\sigma$, and similarly for $\beta$.
Thus $S_\alpha^{m*}S_\alpha^m=I$ and
$S_\beta^{n*}S_\beta^n=I$. The projection and composition
identities follow by multiplication and
$J_{\ell n}J_{nm}=J_{\ell m}$.
\end{proof}

\begin{proposition}[Sharp finite-support obstruction and Parseval node lower bound]
\label{p1b:prop:v275-no-fixed-finite-universal-sampler}
If a probability $\nu$ has finite support $S\subset S^3$ with
$N=|S|$, then
\[
 p_S(q)=\prod_{s\in S}(1-\langle q,s\rangle)\ge0,
 \qquad \int p_S\,d\nu=0<\int p_S\,d\sigma.
\]
Hence no fixed finite support reproduces Haar on all polynomial
bands, and in fact it cannot be exact even on $\mathcal P_N$.
If a positive $N$-node rule is exact on $\mathcal P_L$, then
\[
 \boxed{N\ge d_{\lfloor L/2\rfloor}}
\]
because its weighted evaluation is an isometry on
$\mathcal P_{\lfloor L/2\rfloor}$. For $L=5,7,11,19$ the
corresponding minimal frame ranks are $14,30,91,385$.
\end{proposition}
\begin{proof}
Each factor in $p_S$ is nonnegative on the unit sphere and
vanishes only at the corresponding point. Its product therefore
vanishes precisely on $S$ and is positive off $S$, a set of
full Haar measure. The degree is $N$. The lower rank bound is
the injectivity of the Parseval evaluation map on the indicated
polynomial band, since products of two degree-$m$ scalar
polynomials have degree at most $2m\le L$.
\end{proof}


\part{Marked reciprocal comparisons and finite spectral calculus}
\section{Reciprocal lift of rank-two marked polygon projectors}
\label{p1b:sec:v340-reciprocal-ranktwo}
\begin{proposition}[Factorwise incidence projectors and Born exchange]
\label{p1b:prop:v340-factorwise-ranktwo-projectors}
Suppose, purely as mathematical markings, that binary icosahedral
subgroups have been supplied in both $SU(2)_x$ and $SU(2)_p$.
For $\xi\in\{x,p\}$ let $G_\xi\cong A_5$, let
$M_{k,\xi}$ be the order-$k$ cyclic-coset point Gram matrices on
$\mathbb R[G_\xi]$, and set
\[
 \mathscr P_{5,\xi}=(M_{2,\xi}-M_{3,\xi})/12,
 \qquad \mathscr P_{4,\xi}=(M_{3,\xi}-M_{5,\xi})/15.
\]
They are orthogonal idempotents of ranks $25$ and $16$, respectively.
Their product projectors on the \emph{sampled} reciprocal group-function
carrier have ranks $r_a r_b$ for $r_4=16,r_5=25$. Under an
explicit factor-identification and Born flip $\tau$,
\[
 \tau(\mathscr P_{a,x}\otimes\mathscr P_{b,p})\tau^{-1}
 =\mathscr P_{b,x}\otimes\mathscr P_{a,p}.
\]
The rank-$(16,25)$ character sectors are \emph{not} the
real oriented fifteen-edge Petersen module $W_P$ in each factor,
and are not the complex fifteen-dimensional momentum module $H_1$.
For a separately identified $W_{P,x}\simeq W_{P,p}$, the previously
established exchange dimensions stay
$\dim\operatorname{Sym}^2W_P=120$ and
$\dim\Lambda^2 W_P=105$. Those dimensions do not identify either
space with the $120$ binary-icosahedral vertices, nor with any
simplex count in $\mathbb H P^2$.
\end{proposition}
\begin{proof}
The factorwise identities are the character idempotents in
Theorem~\ref{p1:thm:v340-incidence-character-projectors}.
Tensor products of orthogonal projections have multiplicative ranks,
while the flip exchanges the two factors after an identification of
the source carriers. The dimensions $120$ and $105$ follow from
$15(15\pm1)/2$. An independent $A_{5,x}\times A_{5,p}$ action
need not preserve either exchange eigenspace; the latter carries
in particular the \emph{diagonal} group action after matching frames.
\end{proof}


\section{Reciprocal realization of the marked coset--Hodge edge module}
\label{p1b:sec:v341-coset-coefficient-exchange}
\begin{proposition}[Finite coefficient realization of the exchange carrier]
\label{p1b:prop:v341-coset-pi15-exchange}
Given independently marked $2I$ packets in the two $SU(2)$ factors,
let $B_\xi:\mathbb R^{25}\to\ell^2(2I_\xi)$ be their
transporter incidence matrices, $Q=I_5-J_5/5$, and
$\chi_{4,\xi}$ the standard four-dimensional character vectors.
The factorwise operators
\[
 \Pi_{15,\xi}
 =\frac1{30}B_\xi(Q\otimes Q)B_\xi^{\mathsf T}
 -\frac1{120}\chi_{4,\xi}\chi_{4,\xi}^{\mathsf T},
 \qquad\xi\in\{x,p\},
\]
are orthogonal projectors of rank fifteen.  Their ranges
realize $\mathbf4\oplus\mathbf5\oplus\mathbf3\oplus\mathbf3'$
by explicit finite group-function coefficients, equivalently
the marked oriented Petersen edge module $W_P$.
The two Hodge triplets arise from skew transporter indicators,
whereas the nine-dimensional $\mathbf4\oplus\mathbf5$ sector
arises from the non-invariant part of the symmetric indicators.
Under a chosen common marking and inner-product identification,
Born exchange interchanges the factors and the matched product has
\[
 \dim\operatorname{im}(\Pi_{15,x}\otimes\Pi_{15,p})=225,
 \qquad(\dim E_+,\dim E_-)=(120,105).
\]
No physical selection is implied by the coefficient realization,
and without matched markings the independent $A_{5,x}\times A_{5,p}$
action does not preserve a proposed diagonal exchange sector.
\end{proposition}
\begin{proof}
The factorwise projector identity and module decomposition are
Corollary~\ref{p1:cor:v341-coset-pi15}; the existing edge-to-
traceless-matrix isometry is
Theorem~\ref{p1:thm:v339-petersen-sl4-isometry}.
A specified factor identification conjugates the ordinary tensor
flip to the matched coefficient realization; its symmetric and
antisymmetric ranks are $15(15\pm1)/2$, exactly as in
Corollary~\ref{p1b:cor:petersen-born-even-odd}.
The coefficient projector is mathematical and does not identify
$W_P$ with the complex momentum module $H_1$.
\end{proof}


\section{Matched reciprocal fivefold contrast and its exchange-resolved carrier}
\label{p1b:sec:v343-reciprocal-signed-contrast}
\begin{theorem}[Unique double-odd factor-central contrast and exchange spectra]
\label{p1b:thm:v343-doubleodd-contrast}
Fix two mathematically marked real Petersen modules $W_x,W_p\cong W_P$
and a specified orthogonal factor identification.  Let
$L_\xi=\Pi_{\mathbf3,\xi}-\Pi_{\mathbf3',\xi}$ for $\xi=x,p$,
as normalized in
Proposition~\ref{p1:prop:v343-fivefold-signed-contrast}, and define
\[
 \boxed{\mathcal C_{xp}=L_x\otimes L_p
 =\frac{K_x\otimes K_p}{80}
 =\frac{45}{16}\,
 (\Phi_x^{-1}D_{5,x}\Phi_x)\otimes
 (\Phi_p^{-1}D_{5,p}\Phi_p).}
\]
With matched orientations and Born exchange $\tau$, one has
$\tau\mathcal C_{xp}\tau^{-1}=\mathcal C_{xp}$ and
$\mathcal C_{xp}^3=\mathcal C_{xp}$.  Among factorwise central
operators odd under the abstract five-cycle-class interchange
\emph{in each factor}, $\mathcal C_{xp}$ is the unique normalized
double-odd, exchange-even direction, up to the joint naming of
$5A$ and $5B$.  The full independent $A_{5,x}\times A_{5,p}$
commutant has dimension $16$; imposing matched Born exchange
leaves dimension $10$, while also imposing simultaneous outer
class interchange leaves dimension $7$.
On $W_x\otimes W_p$ and its matched exchange sectors the exact spectra are
\[
\begin{array}{c|rrr|r}
 & +1&-1&0&\dim\\\hline
 W_x\otimes W_p&18&18&189&225\\
 \operatorname{Sym}^2 W_P&12&9&99&120\\
 \Lambda^2W_P&6&9&90&105
\end{array}.
\]
In particular $\operatorname{tr}_{120}\mathcal C_{xp}=3$,
$\operatorname{tr}_{105}\mathcal C_{xp}=-3$ and the
corresponding quadratic traces are $21$ and $15$.
\end{theorem}
\begin{proof}
Each single-factor central algebra is the real span of the four
irreducible projectors, and its odd subspace under the interchange
$\mathbf3\leftrightarrow\mathbf3'$ is $\mathbb R L$.
The symmetric square of the resulting four-dimensional coefficient
space has dimension $10$; with simultaneous outer interchange its
even part is $\operatorname{Sym}^2\mathbb R^3_{\rm even}
\oplus\operatorname{Sym}^2\mathbb R_{\rm odd}$, of dimension $6+1=7$.
Born exchange sends $L_x\otimes L_p$ to $L_p\otimes L_x$.
The $L$-eigenspace dimensions are $3,3,9$; tensor products and the
standard symmetric/exterior square dimensions yield the table.
The equality with $D_5$ uses
Proposition~\ref{p1:prop:v343-fivefold-signed-contrast} in both factors.
\end{proof}

\begin{theorem}[Canonical fifteen-dimensional active summand of the exchange-odd carrier]
\label{p1b:thm:v343-canonical-selected15}
Set $U=\operatorname{im}L^2\cong\mathbf3\oplus\mathbf3'$ and
$V=\ker L\cong\mathbf4\oplus\mathbf5$, of dimensions $6$ and $9$.
The exchange-odd carrier admits the \emph{orthogonal, operator-selected}
splitting
\[
\begin{split}
\Lambda^2 W_P&=\underbrace{\Lambda^2 U}_{15}
\oplus\underbrace{(U\otimes V)\oplus\Lambda^2V}_{54+36=90},\\
\boxed{P_{15}=\mathcal C_{xp}^2\big|_{\Lambda^2W_P}},&\qquad
\operatorname{im}P_{15}=\Lambda^2U,\quad
\operatorname{rank}P_{15}=15.
\end{split}
\]
On the selected image, $\mathcal C_{xp}$ has eigenvalues $+1$
on $\Lambda^2\mathbf3\oplus\Lambda^2\mathbf3'\cong
\mathbf3\oplus\mathbf3'$ (dimension $6$), and $-1$ on
$\mathbf3\otimes\mathbf3'\cong\mathbf4\oplus\mathbf5$
(dimension $9$).  Thus, as $A_5$ modules,
\[
\boxed{\Lambda^2W_P\cong W_P^{\oplus7},\qquad
 \Lambda^2U\cong W_P,\qquad
 \ker\mathcal C_{xp}\big|_{105}\cong W_P^{\oplus6}.}
\]
The first $15$ is distinguished by $P_{15}$; the remaining six
abstract copies of $W_P$ are not individually distinguished by it.
\end{theorem}
\begin{proof}
Since $L^2$ is the orthogonal projector onto $U$,
$\mathcal C_{xp}^2=L^2\otimes L^2$ projects onto $U\otimes U$.
Its exchange-odd restriction has range $\Lambda^2U$ and rank
$\binom62=15$.  The indicated $6+9$ split follows from
$\Lambda^2\mathbf3=\mathbf3$,
$\Lambda^2\mathbf3'=\mathbf3'$ and
$\mathbf3\otimes\mathbf3'=\mathbf4\oplus\mathbf5$.
The exterior-square character
$\chi_{\Lambda^2W}(g)=(\chi_W(g)^2-\chi_W(g^2))/2$
equals $7\chi_W(g)$ on $A_5$; subtracting the selected copy
identifies the residual six $W_P$ copies.
\end{proof}

\begin{theorem}[$S_5$-sign refinement and a compact symmetric pair]
\label{p1b:thm:v343-s5-twist-so6}
For the oriented fifteen-edge representation of the full relabeling
group $S_5$, the prior sevenfold $A_5$ identity refines to
\[
 \boxed{\Lambda^2W_P\cong_{S_5}3W_P\oplus
 4(\operatorname{sgn}\otimes W_P),\qquad
 \Lambda^2U\cong_{S_5}\operatorname{sgn}\otimes W_P.}
\]
The invariant inner product identifies the \emph{selected}
$\Lambda^2U$ with $\mathfrak{so}(U)\cong\mathfrak{so}(6)$.
Conjugation by $L|_U=\operatorname{diag}(I_3,-I_3)$
induces precisely $\mathcal C_{xp}$ on this subalgebra and
gives the $\mathbb Z_2$-graded symmetric pair
\[
 \boxed{\mathfrak{so}(6)=
 [\mathfrak{so}(3)\oplus\mathfrak{so}(3)]_{+}
 \oplus[\mathbb R^3\otimes\mathbb R^3]_{-},\qquad 15=6+9.}
\]
The standard skew-generator quadratic adjoint Casimir of the
$\mathfrak{so}(3)\oplus\mathfrak{so}(3)$ subalgebra, extended by
zero on the inactive $90$ modes, is
\[
 \boxed{\Delta_{\mathrm{rec}}=3\mathcal C_{xp}^{\,2}-\mathcal C_{xp},
 \qquad\operatorname{Spec}\Delta_{\mathrm{rec}}|_{105}
 =\{2^{(6)},4^{(9)},0^{(90)}\}.}
\]
Neither $\mathfrak{so}(6)$ nor this auxiliary Casimir is thereby a
physical gauge algebra, mass or kinetic operator.
\end{theorem}
\begin{proof}
The real signed-edge character of $W_P$, on the cycle types
$1^5,2\,1^3,2^2 1,3\,1^2,3\,2,4\,1,5$, is
$(15,3,-1,0,0,-1,0)$.  Its exterior-square character equals
$(105,-3,-7,0,0,1,0)=(3+4\operatorname{sgn})\chi_W$.
The six-dimensional $U$ is invariant under $S_5$ even though $L$
changes sign under odd relabeling; the same character calculation
for $\Lambda^2U$ yields $\operatorname{sgn}\chi_W$.
Writing $U=U_+\oplus U_-$ gives the block-skew realization
$\bigl(\begin{smallmatrix}A&B\\-B^{\mathsf T}&D\end{smallmatrix}\bigr)$.
Conjugation by $L$ fixes $A,D$ and reverses $B$; the
block-matrix bracket proves the symmetric-pair rules.
The two adjoint $\mathfrak{so}(3)$ blocks carry Casimir $2$;
the mixed $(\mathbf3,\mathbf3)$ block carries $2+2=4$.
The polynomial has eigenvalues $2,4,0$ at $\mathcal C_{xp}=1,-1,0$.
\end{proof}

\begin{remark}[Identification and symmetry firewall]
\label{p1b:rem:v343-reciprocal-nonactualization}
The decomposition into the $120$ and $105$ exchange eigenspaces
requires an explicit matching $W_x\simeq W_p$; independent
$A_{5,x}\times A_{5,p}$ actions do not preserve that split.
An independent abstract outer odd relabeling reverses
$\mathcal C_{xp}$, while the selected projector $P_{15}$ remains
unchanged; simultaneous odd relabeling of both factors preserves
$\mathcal C_{xp}$.  The Lie algebra $\mathfrak{so}(6)$ is compact
and is \emph{not isomorphic as a real Lie algebra} to the
$\mathfrak{sl}(4,\mathbb R)$ algebra reconstructed in Paper I-B,
despite both having dimension $15$.
No physical Born-exchange sector, fixed coupling constant,
rest-mass ratio or chirality effect is inferred from these
finite mathematical statements alone.
\end{remark}

\clearpage
\section{The degree-twelve icosahedral alias and its exact positive cubature repair}
\label{p1b:sec:v349-twisted-harmonic-quadrature}
Let $\mu$ be probability Haar measure on $S^3(R)$ and
$\nu_I$ the uniform measure on the $120$ vertices of $2I$.
The comparison coefficient shell is the degree-four scalar
harmonic space $\mathcal H_4$, not the physically selected
momentum shell.

\begin{theorem}[Unique degree-twelve sampling defect through degree nineteen]
\label{p1b:thm:v349-single-degree12-alias}
For $\mathcal P_{19}=\bigoplus_{n=0}^{19}\mathcal H_n(S^3)$,
the functional $\nu_I-\mu$ is supported on a single
harmonic direction in $\mathcal H_{12}$.  If $v$ is a unit
$2I$-fixed vector in the 13-dimensional spin-six
$SU(2)$ representation, its Riesz representer is
\[
 h_{12}(q)=13\langle v,\rho_6(q)v\rangle,
 \quad \|h_{12}\|_{L^2}^2=13,
 \quad (\nu_I-\mu)(p)=\langle p,h_{12}\rangle_{L^2}
 \quad(p\in\mathcal P_{19}).
\]
The normalized spin-two character
$f(q)=16(\operatorname{Re}q)^4-12(\operatorname{Re}q)^2+1$
obeys
\[
 \int f^2d\mu=\int f^2d\nu_I=1,\qquad
 \int f^3d\mu=1,\qquad\int f^3d\nu_I=2.
\]
Thus raw single-$2I$ sampling is quadratically exact on the
25-dimensional coefficient shell but has an additional cubic
invariant, of exact sample/continuum discrepancy one on $f^3$.
\end{theorem}
\begin{proof}
Finite character averaging on $2I$ finds a fixed-vector rank
one in spherical harmonic degrees $n=0,12$ and rank zero for
all other $1\leq n\leq19$.  Degree twelve has the single
spin-six fixed line, so Schur orthogonality gives its
Riesz representer and squared norm $13$.  The spin-two
character takes $A_5$ values $(5,1,-1,0,0)$;
continuous Clebsch--Gordan coupling gives one cubic singlet,
whereas finite class averaging gives
$(5^3+15-20)/60=2$.  The common quadratic value is one.
\end{proof}

\begin{corollary}[Existing positive $IJI$ cubature is exact through quartic order]
\label{p1b:cor:v349-twisted25-quartic-cubature}
The previously proved positive triple-icosahedral convolution
$\nu_{IJI}=\mu_I*\mu_J*\mu_I$ of
Theorem~\ref{p1b:thm:v273-triple-icosahedral-nineteen-design}
is exact for every scalar polynomial of total degree at most
nineteen.  Therefore for any $f_1,\ldots,f_k\in\mathcal H_4$,
$k\leq4$, both
\[
 \int\prod_{a=1}^k f_a\,d\mu
 =\int\prod_{a=1}^k f_a\,d\nu_{IJI}
\]
and, for $2\leq k\leq4$,
\[
 \int(\nabla f_1\cdot\nabla f_2)\prod_{a=3}^k f_a\,d\mu
 =\int(\nabla f_1\cdot\nabla f_2)\prod_{a=3}^k f_a\,d\nu_{IJI}.
\]
\end{corollary}
\begin{proof}
A degree-four spherical harmonic is the restriction of a
homogeneous degree-four harmonic polynomial $F$ on $\mathbb R^4$.
Products of four such fields have degree at most sixteen.
For unit radius,
$\nabla_{S^3}f\cdot\nabla_{S^3}g=
\nabla_{\mathbb R^4}F\cdot\nabla_{\mathbb R^4}G-16FG$
on the sphere, again of degree at most eight.
The $R^{-2}$ gradient factor is common to both measures.
Thus the already constructed positive degree-nineteen cubature
integrates every displayed integrand exactly.
\end{proof}

\begin{remark}[No physical measure substitution]
\label{p1b:rem:v349-cubature-action-guard}
This correction is an exact \emph{mathematical comparison}
for polynomial vertices and their first-gradient contractions.
It does not decree that the actual common action is a spin-two
field theory, that a 120-site sum is its measure, or that the
$IJI$ word measure is the physical integration prescription.
Nonpolynomial action functions require independent control.
\end{remark}


\section{Grassmannian and topology of the reciprocal nine-mode comparison}
\label{p1b:sec:cc14-grassmann}
\begin{theorem}[Exact oriented Grassmannian realization]
\label{p1b:thm:cc14-grassmann}
Fix the selected matched $U=U_x\oplus U_p\cong\mathbb R^3\oplus\mathbb R^3$
of Theorem~\ref{p1b:thm:v343-s5-twist-so6}.  Its nine-dimensional
$\mathcal C_{xp}$-odd summand of $\mathfrak{so}(U)$ is
$\mathfrak p=\{X(B):B\in\mathrm{Mat}_{3\times3}(\mathbb R)\}$, with
\[
 X(B)=\begin{pmatrix}0&B\\-B^{\mathsf T}&0\end{pmatrix}.
\]
It is the tangent space at the marked plane $U_x$ of the compact
oriented Grassmannian
\[
 \boxed{\mathcal G=\mathrm{Gr}^+_3(\mathbb R^6)
 \cong SO(6)/(SO(3)\times SO(3)),\qquad\dim\mathcal G=9.}
\]
With the invariant metric $\langle X,Y\rangle=-\frac12\operatorname{tr}(XY)$,
$\langle X(B),X(C)\rangle=\operatorname{tr}(B^{\mathsf T}C)$.
The origin-based geodesic $\exp(tX(B))U_x$ has infinitesimal principal
angle speeds equal to the three singular values of $B$.
The largest flat tangent abelian subspace has rank three, represented
by diagonal $B$ after independent orthogonal frame choices.
For linearly independent $B,C$, the sectional curvature is
\[
 \boxed{K(B,C)=\frac{\|CB^{\mathsf T}-BC^{\mathsf T}\|_F^2
 +\|C^{\mathsf T}B-B^{\mathsf T}C\|_F^2}
 {2\bigl(\|B\|_F^2\|C\|_F^2-
 (\operatorname{tr}B^{\mathsf T}C)^2\bigr)}\ge0.}
\]
It vanishes exactly when both numerator commutators vanish.
\end{theorem}
\begin{proof}
The $SO(6)$-stabilizer of the marked oriented three-plane is
$SO(3)\times SO(3)$, giving the homogeneous space and tangent blocks.
For $B=U\operatorname{diag}(\sigma_i)V^{\mathsf T}$, constant frame
rotations conjugate $X(B)$ into three independent two-plane rotation
generators, proving the principal-angle description.  Their diagonal
span commutes and has rank three, which is maximal for the
Grassmannian $SO(3+3)/(SO(3)SO(3))$ by the standard singular-value
normal form.
For the compact symmetric-space metric,
$K(X,Y)=\|[X,Y]\|^2/(\|X\|^2\|Y\|^2-\langle X,Y\rangle^2)$.
The commutator is block diagonal with entries
$CB^{\mathsf T}-BC^{\mathsf T}$ and
$C^{\mathsf T}B-B^{\mathsf T}C$; taking its squared
$-\frac12\operatorname{tr}$ norm yields the displayed formula.
\end{proof}

\begin{corollary}[Simply connected nine-mode space with a torsion sphere obstruction]
\label{p1b:cor:cc14-grassmann-topology}
The comparison manifold $\mathcal G$ has
\[
 \boxed{\pi_1(\mathcal G)=0,\qquad\pi_2(\mathcal G)\cong\mathbb Z_2.}
\]
Thus the nine-dimensional linear representation extends to a
nontrivial compact curved homogeneous space whose second homotopy
contains two-torsion; its tangent dimensions alone do not recover
this global topology.
\end{corollary}
\begin{proof}
The fibration $SO(3)\times SO(3)\to SO(6)\to\mathcal G$
has $\pi_2SO(6)=0$, $\pi_1(SO(3)\times SO(3))=\mathbb Z_2\oplus
\mathbb Z_2$, and $\pi_1SO(6)=\mathbb Z_2$.  The induced map of
fundamental groups is $(a,b)\mapsto a+b$, since a basic $2\pi$
rotation in either three-plane represents the generating loop in
$SO(6)$.  The long exact sequence therefore gives
$\pi_2\mathcal G=\ker(+)=\mathbb Z_2$ and
$\pi_1\mathcal G=\operatorname{coker}(+)=0$.
\end{proof}
\begin{remark}[No physical Born-exchange metric claim]
This is the canonical Riemannian geometry of the \emph{selected
compact symmetric pair}, not a derived physical target-space metric,
Lorentzian spacetime, actual gauge manifold, new bosonic particle, or
an identification with the unrelated $\mathfrak{sl}(4,\mathbb R)$
comparison Lie algebra of the same dimension.
\end{remark}

\begin{theorem}[Non-spin reciprocal frames with a spin nine-mode tangent]
\label{p1b:thm:cc15-grassmann-spin-cancellation}
For the selected Grassmannian $\mathcal G=\mathrm{Gr}^+_3(\mathbb R^6)$
of Theorem~\ref{p1b:thm:cc14-grassmann}, let $\mathcal E$ be the
universal oriented rank-three subbundle and $\mathcal F=\mathcal E^\perp$
its oriented rank-three complement. Then
\[
 \mathcal E\oplus\mathcal F\cong\underline{\mathbb R}^{6},\qquad
 T\mathcal G\cong\mathcal E^*\otimes\mathcal F.
\]
The mod-two degree-two cohomology is
$H^2(\mathcal G;\mathbb Z_2)\cong\mathbb Z_2$, and
\[
 \boxed{w_2(\mathcal E)=w_2(\mathcal F)\ne0,
 \qquad w_2(T\mathcal G)=0.}
\]
Thus the two canonical reciprocal three-plane bundles have no global spin
lift, while the curved nine-dimensional tangent bundle is spin.
Since $\pi_1(\mathcal G)=0$, the spin structure of the tangent
bundle is unique up to isomorphism.
\end{theorem}
\begin{proof}
In the principal bundle
$SO(3)\times SO(3)\to SO(6)\to\mathcal G$, the map
on fundamental groups is $(a,b)\mapsto a+b$ in $\mathbb Z_2$.
The generator of $\pi_2\mathcal G=\mathbb Z_2$ therefore has boundary
$(1,1)$; both tautological oriented rank-three bundles restrict to this
clutching sphere with the nontrivial $SO(3)$ loop. Their $w_2$
classes evaluate to one, and since $\mathcal G$ is simply connected,
Hurewicz and universal coefficients give the displayed nonzero
class in $H^2(\mathcal G;\mathbb Z_2)$.
The direct-sum identity gives
$w_2(\mathcal E)+w_2(\mathcal F)=0$ because both first
Stiefel--Whitney classes vanish. For oriented bundles of ranks $r,s$,
$w_2(E\otimes F)=s\,w_2(E)+r\,w_2(F)$ modulo two.
Here $r=s=3$ and $\mathcal E^*\cong\mathcal E$ at the level of
Stiefel--Whitney classes, so
$w_2(T\mathcal G)=3w_2(\mathcal E)+3w_2(\mathcal F)=0$.
An oriented tangent bundle is spin precisely when $w_2=0$;
its spin structures form an $H^1(\mathcal G;\mathbb Z_2)$-torsor,
which is trivial by simple connectedness.
\end{proof}
\begin{remark}[Two mod-two obstructions cancel, not two new particles]
The cancellation occurs in the tangent tensor product of the selected
auxiliary compact symmetric space. It neither introduces physical
reciprocal spinors nor fixes a spin coupling, QCD flavor, gravitational
spin structure or interaction coefficient. Any physical use requires the
existing common-owner and operator-incidence maps.
\end{remark}


\begin{theorem}[Positive curvature and the canonical Grassmannian spin-Dirac gap]
\label{p1b:thm:cc16-grassmann-spin-gap}
On the selected compact reciprocal nine-mode comparison manifold
$\mathcal G=SO(6)/(SO(3)\times SO(3))$ of
Theorem~\ref{p1b:thm:cc14-grassmann}, use exactly the invariant
metric $\langle X,Y\rangle=-\frac12\operatorname{tr}(XY)$.
Its Ricci tensor and scalar curvature are
\[
 \boxed{\operatorname{Ric}_{\mathcal G}=4g,\qquad
        \operatorname{Scal}_{\mathcal G}=36.}
\]
With the unique tangent spin structure established in
Theorem~\ref{p1b:thm:cc15-grassmann-spin-cancellation}, the
\emph{untwisted geometric} spin-Dirac operator satisfies
\[
 \boxed{D_{\mathcal G}^{\,2}\ge9I,\qquad
 \operatorname{ker}D_{\mathcal G}=\{0\}.}
\]
After uniformly scaling the metric by $\ell^2$, the squared gap is
$9/\ell^2$. This is an auxiliary compact Riemannian result,
not a physical fermion-mass lower bound or a Lorentzian Dirac spectrum.
\end{theorem}
\begin{proof}
Choose the orthonormal tangent basis
$X(E_{ia})$, $1\le i,a\le3$, of the $3\times3$ off-diagonal
matrices. The compact symmetric-space curvature formula gives
$K(X,Y)=\|[X,Y]\|^2/\|X\wedge Y\|^2$. For a fixed
$X(E_{ia})$, exactly four other basis elements share its row or
column and have sectional curvature one; the remaining four have
sectional curvature zero. Thus every diagonal Ricci entry equals
four. Invariance of the metric and irreducibility of the isotropy
representation $\mathbb R^3\otimes\mathbb R^3$ force the same
multiple of the metric on all tangent directions, whence
$\operatorname{Ric}=4g$ and $\operatorname{Scal}=9\cdot4=36$.
On the spin manifold, the Lichnerowicz identity is
$D^2=\nabla^*\nabla+\operatorname{Scal}/4$; the nonnegative
connection Laplacian gives $D^2\ge9I$ and the trivial kernel.
Scaling the metric gives the indicated reciprocal-square law.
\end{proof}
\begin{remark}[No invented physical chiral zero modes]
The spin cancellation proves existence of the \emph{geometric}
untwisted spin bundle, but the positive-curvature comparison has no
untwisted harmonic spinors. Twisted operator couplings and physical
fermion sources are separate typed claims requiring their own action
and source certificates.
\end{remark}

\subsection{Exact 120-vertex scalar spectral and mixed-degree Galerkin calculus}
\label{p1b:sec:cc16-120-spectral-compiler}
\begin{theorem}[Explicit degree-shell projectors and exact scalar functional calculus]
\label{p1b:thm:cc16-120-spectral-compiler}
Let $q_1,\ldots,q_{120}$ be the quaternionic vertices of one marked
$2I\subset S^3(R)$, identified by unit quaternion coordinates. With
$U_n$ the degree-$n$ Chebyshev polynomial of the second kind, put
\[
 (E_n)_{ij}=\frac{n+1}{120}U_n(q_i\cdot q_j),\qquad
 0\leq n\leq5,\qquad P_5^{\rm sam}=\sum_{n=0}^5E_n.
\]
Then the matrices are real symmetric orthogonal projectors, and
\[
 \boxed{E_nE_m=\delta_{nm}E_n,\qquad
 \operatorname{rank}E_n=(n+1)^2,\qquad
 \operatorname{rank}P_5^{\rm sam}=91,\quad
 \dim\ker P_5^{\rm sam}=29.}
\]
The explicit 120-by-120 scalar Laplace representative
\[
 \boxed{L_5=R^{-2}\sum_{n=0}^5 n(n+2)E_n}
\]
intertwines the positive round scalar Laplacian with the normalized
vertex evaluation $S$ on
$\mathcal P_5=\oplus_{n=0}^5\mathcal H_n(S^3(R))$:
$L_5 Sf=S\Delta_+f$.
For any function $F$ on the six-point eigenvalue set,
\[
 \boxed{\operatorname{tr}\big(P_5^{\rm sam}F(L_5)\big)
 =\sum_{n=0}^5(n+1)^2F(n(n+2)/R^2).}
\]
Taking the ambient trace $\operatorname{tr}F(L_5)$ without the sample
projector would add the spurious term $29F(0)$.
All statements hold in the 91-dimensional scalar harmonic comparison
band, not on the distinct physical $15$-dimensional momentum carrier
or the physical spinorial Dirac operator.
\end{theorem}
\begin{proof}
The addition theorem for the Haar probability-normalized scalar
$S^3$ harmonics states that the degree-$n$ reproducing kernel is
$(n+1)U_n(x\cdot y)$. Hence $E_n=S_nS_n^*$ where
$S_n:\mathcal H_n\to\mathbb C^{120}$ is the evaluation divided
by $\sqrt{120}$. The binary-icosahedral degree-eleven design of Theorem~\ref{p1:thm:v180-exceptional-scalar-moment-ranges}, with the frame of Theorem~\ref{p1b:thm:v273-cosetwise-resolution-groupoid}, makes
$S_n^*S_m=\delta_{nm}I$ for all $n,m\leq5$, since the integrand
is a polynomial of degree at most $n+m\leq10$. Thus the $E_n$
are orthogonal projectors with rank $\dim\mathcal H_n=(n+1)^2$.
Summing ranks gives 91 and a 29-dimensional sample complement.
On each $S_n\mathcal H_n$, the indicated $L_5$ is multiplication by
$n(n+2)/R^2$, proving the intertwiner and trace claims.
On the sample complement $L_5$ is zero, explaining exactly the
additional $29F(0)$ term of the unprojected trace.
\end{proof}
\begin{corollary}[Exact finite scalar heat and massive resolvent operators]
\label{p1b:cor:cc16-120-spectral-green}
On the sample image, for $t\geq0$ and $\mu^2>0$,
\[
 H_5(t)=\sum_{n=0}^5e^{-tn(n+2)/R^2}E_n,\qquad
 G_5(\mu)=\sum_{n=0}^5\frac{E_n}{\mu^2+n(n+2)/R^2},
\]
obey $H_5(t)H_5(s)=H_5(t+s)$ and
$(L_5+\mu^2P_5^{\rm sam})G_5(\mu)=P_5^{\rm sam}$.
For a sampled scalar field $Sf$ in the image,
$\langle Sf,L_5Sf\rangle=\langle f,\Delta_+f\rangle_{L^2}$.
These are finite scalar comparisons, not automatically stochastic
kernels or interacting propagators.
\end{corollary}
\begin{proposition}[Exact mixed-band Galerkin product criterion]
\label{p1b:prop:cc16-120-mixed-galerkin}
For nonnegative integers $a,b,c\leq5$ satisfying $a+b+c\leq11$,
let $f\in\mathcal P_a$, $g\in\mathcal P_b$. Write $P_{\le c}$
for the Haar-orthogonal projection onto $\mathcal P_c$ and
$\mathsf E_{\le c}=\sum_{n=0}^cE_n$.
Then the sampled Galerkin product is exactly
\[
 \boxed{S_cP_{\le c}(fg)=\mathsf E_{\le c}
 \bigl(\sqrt{120}\,(S_af)\odot(S_bg)\bigr).}
\]
In particular, two degree-three fields have their complete
$\mathcal P_5$ Galerkin product computed exactly by the 120-point
samples because $3+3+5=11$, even though $fg$ may contain degree-six
content. Similarly all mixed scalar triple overlaps
$\int\overline hfg$ with $h\in\mathcal P_c$ are exactly given by
uniform 120-point summation under the same degree condition.
\end{proposition}
\begin{proof}
For every test $h\in\mathcal P_c$ the integrand
$\overline hfg$ is a restricted polynomial of coordinate degree
at most $a+b+c\leq11$. The binary-icosahedral eleven-design
therefore equates its Haar integral with the 120-node average.
This is precisely the weak defining identity of the orthogonal
Galerkin projection. The isometric evaluation $S_c$ then gives
the displayed finite-matrix formula. The projection is essential:
finite pointwise products are not asserted to reproduce all
untruncated higher harmonics.
\end{proof}


\subsection{A 720-node scalar spectral compiler and chamber-harmonic leakage}
\label{p1c:sec:cc17-chamber-spectral-compilers}
\begin{theorem}[Explicit spectral projectors through scalar harmonic degree nine]
\label{p1c:thm:cc17-720-shell-projectors}
Use the 720 distinct nodes $z_\alpha\in V\cup F$ and per-node weights
$w_\alpha=1/504$ ($z\in V$) and $w_\alpha=2/1575$ ($z\in F$)
from Theorem~\ref{p1b:thm:cc17-positive-vertex-facet-19}.
For $0\le n\le9$ put
\[
 (E_n)_{\alpha\beta}
 =\sqrt{w_\alpha w_\beta}(n+1)
 U_n(z_\alpha\cdot z_\beta),\qquad P_9=\sum_{n=0}^9E_n.
\]
Then the $E_n$ are pairwise orthogonal self-adjoint projections,
\[
\boxed{E_nE_m=\delta_{nm}E_n,\qquad
\operatorname{rank}E_n=(n+1)^2,\qquad
\operatorname{rank}P_9=385,\quad\dim\ker P_9=335.}
\]
For the round scalar Laplacian on $S^3(R)$ the exact finite band operator is
\[
 L_9=R^{-2}\sum_{n=0}^9n(n+2)E_n,
\qquad\operatorname{Tr}(P_9F(L_9))
 =\sum_{n=0}^9(n+1)^2F(n(n+2)/R^2).
\]
This is a 385-dimensional scalar comparison; an unprojected 720-by-720
trace includes the spurious factor $335F(0)$.
\end{theorem}
\begin{proof}
The scalar harmonic degree-$n$ reproducing kernel is $(n+1)U_n$.
The positive degree-nineteen rule integrates products of two
harmonics of degrees $n,m\le9$ exactly since $n+m\le18$.
Thus weighted evaluation $S_n$ satisfies $S_n^*S_m=\delta_{nm}I$,
so $E_n=S_nS_n^*$ has the stated projection and dimension properties.
The Laplacian is scalar $n(n+2)/R^2$ on $\mathcal H_n$.
Restriction of the trace to $P_9$ yields the displayed formula.
\end{proof}

\begin{theorem}[The tetrahedral chamber walk is a non-invariant harmonic Galerkin operator]
\label{p1c:thm:cc17-face-graph-harmonic-Galerkin}
Let $F$ denote the uniform 600 facet centres, and let $A$ denote
the source's symmetric four-regular face-adjacency matrix.
Set for $0\le n\le5$
\[
 (E^{F}_n)_{ij}=\frac{n+1}{600}U_n(c_i\cdot c_j),
\quad c_*:=\frac{1+3\sqrt5}{8},\quad d_*:=\frac{5+\sqrt5}{8}.
\]
These $E_n^F$ are rank-$(n+1)^2$ mutually orthogonal projections.
For $n\le5$ the following are exact:
\[
\boxed{E_n^F A E_n^F=\lambda_n E_n^F,\qquad
\lambda_n=\frac4{n+1}U_n(c_*).}
\]
Define the off-shell leakage $K_n=(I-E_n^F)A E_n^F$.
Then
\[
\boxed{K_n^*K_n=\delta_n E_n^F,\qquad
\delta_n=4+\frac{12}{n+1}U_n(d_*)-\lambda_n^2.}
\]
In degrees $0,1,2$ the leakage vanishes, whereas in degrees $3,4,5$
it is \emph{strictly positive}:
\[
\begin{array}{c|c|c}
 n&\lambda_n&\delta_n\\\hline
 3&13/8+3\sqrt5/4&27/64-3\sqrt5/16\\
 4&(43+33\sqrt5)/40&(1173-519\sqrt5)/800\\
 5&(55+11\sqrt5)/32&(1273-557\sqrt5)/512.
\end{array}
\]
Thus the compression is exact but the source's discrete chamber walk
does not preserve $\mathcal H_3,\mathcal H_4,\mathcal H_5$ as
embedded sample subspaces.
\end{theorem}
\begin{proof}
The 600 facet centres are individually an eleven-design as the
$I$-bi-invariant double coset; hence the $E_n^F$ are projections for
$n\le5$. The $I\times I$ isometric action makes both the face graph
and these harmonic projections equivariant. The restrictions
$\operatorname{Sym}^n\mathbb C^2|_{2I}$ for $0\le n\le5$ are irreducible:
this follows from the source's 120-element character sums
$120^{-1}\sum_{g\in I}|U_n(\operatorname{Re}g)|^2=1$.
Schur's lemma makes the compression $E_n^FAE_n^F$ scalar.
Each chamber centre has four nearest chamber neighbours at dot
product $c_*$, giving
$\lambda_n=4U_n(c_*)/(n+1)$ by the reproducing-kernel trace.
Each centre has four closed two-step returns and twelve distinct
two-step neighbours at dot product $d_*$, so the normalized trace of
$(E_n^FA^2E_n^F)$ is
$4+12U_n(d_*)/(n+1)$. Apply Schur again to the positive operator
$K_n^*K_n=E_n^F A^2 E_n^F-\lambda_n^2E_n^F$.
Inserting $c_*,d_*$ and simplifying Chebyshev polynomials yields
all three exact radical coefficients. The nearest/two-step geometry
and every matrix equality admit a finite exact quaternionic enumeration.
\end{proof}

\begin{corollary}[No common smooth-heat clock and a rigorous walk error bound]
\label{p1c:cor:cc17-chamber-heat-obstruction}
For the symmetric stochastic chamber generator $Q=A/4-I$, the
exact degree-one and degree-two decay rates are
\[
 \gamma_1=\frac{7-3\sqrt5}{8},\qquad
 \gamma_2=\frac{3-\sqrt5}{8},\qquad
 \boxed{\gamma_2/\gamma_1=\frac{3+\sqrt5}{2}\ne\frac83.}
\]
Therefore no single round-sphere diffusion-time conversion can
identify both with scalar Laplace rates $n(n+2)/R^2$.
For unit $x\in\operatorname{Ran}E_n^F$, $1\le n\le5$,
\[
\boxed{\|e^{tQ}x-e^{-\gamma_n t}x\|_2
 \le\frac{\sqrt{\delta_n}}{4\gamma_n}
  (1-e^{-\gamma_nt}),\qquad\gamma_n=1-\lambda_n/4.}
\]
The bound is zero for the preserved $n=1,2$ sectors; it controls the
nontrivial leakage for $n=3,4,5$ without identifying the graph walk
with continuous heat flow.
\end{corollary}
\begin{proof}
The ratio follows by substitution. Since $A$ is symmetric and
four-regular, $Q\preceq0$ and $\|e^{tQ}\|\le1$.
For $x\in\operatorname{Ran}E_n^F$,
$\|(Q+\gamma_n)x\|=\sqrt{\delta_n}/4\,\|x\|$.
Duhamel's formula and contraction of $e^{tQ}$ give
$\|e^{tQ}x-e^{-\gamma_nt}x\|
 \le(\sqrt{\delta_n}/4)\int_0^t e^{-\gamma_ns}\,ds$.
\end{proof}

\begin{proposition}[First register-visible scalar harmonics]
\label{p1c:prop:cc17-register-degree-six-eight}
Let $S:\mathbb R^{25}\to\mathbb R^{600}$ be the exact
$24$-points-per-register isometry and let
$P_c=(I_5-J_5/5)\otimes(I_5-J_5/5)$ be the centered rank-16
register projector. Put $K_n(i,j)=(n+1)U_n(c_i\cdot c_j)/600$
for all $n$, without asserting $K_n$ is a projection outside $n\le5$.
Then
\[
\boxed{S^TK_nS=0\ (1\le n\le5),\qquad
 S^TK_6S=\tfrac7{16}P_c,\quad S^TK_8S=\tfrac9{16}P_c.}
\]
The last two are \emph{sampled harmonic Gram coefficients}, not
orthogonal spectral weights or a decomposition of the whole chamber
state. They identify degree six as the first register-visible
nonconstant smooth scalar band.
\end{proposition}
\begin{proof}
Each of the 25 blocks is a 24-element bi-translate of $H=2T$.
For any degree $n$,
$\operatorname{Tr}(S^TK_nS)=(n+1)m_H(n)$, where
$m_H(n)$ is the coefficient of the canonical $2T$ Molien series.
That coefficient vanishes for $1\le n\le5$, so positivity of the
sampled Gram forces $S^TK_nS=0$.
At $n=6,8$ the $H$-fixed rank is one while the $I$-fixed rank
is zero. Frobenius reciprocity for the five-point permutation module
$\operatorname{Ind}_H^I\mathbf1=\mathbf1\oplus\mathbf4$
places both left and right slots entirely in the centered
$\mathbf4\otimes\mathbf4$ register sector $P_c$.
Equivariance and Schur's lemma make $S^TK_nS$ a scalar on that
sector. Its trace is $n+1$, its rank-16 factor has trace 16,
thus the coefficients are $7/16$ and $9/16$.
\end{proof}

\begin{remark}[Distinct operators stay distinct]
The first two harmonic shells happen to be invariant under the
chamber walk, but their decay-rate ratio differs from round scalar
heat. At degree three and above the off-band leakage is positive.
The chamber coarse-graining, comparison scalar Laplacian, and native
physical Dirac/Einstein/QCD operators are not identified by this module.
\end{remark}


\section{Amplitude versus probability information in the 720-point frame}
\label{p1c:sec:cc18-amplitude-probability}
\begin{proposition}[Exact finite-frame amplitude reconstruction and strict POVM-tomography bound]
\label{p1c:prop:cc18-720-tomography-bound}
For the scalar $S^3$ polynomial band $\mathcal P_9$ of dimension
$d=\sum_{n=0}^9(n+1)^2=385$, the 720-point frame of
Theorem~\ref{p1c:thm:cc17-720-shell-projectors} is an isometry
$S:\mathbb C^{385}\to\mathbb C^{720}$.  Its complex coherent
amplitudes $y=S\psi$ determine $\psi=S^*y$ exactly.  The associated
rank-one effects $E_z=|s_z\rangle\langle s_z|$ with
$\sum_zE_z=I_{385}$, however, are \emph{not} informationally complete
on arbitrary mixed states: the map from traceless Hermitian
operators to outcome-probability deviations has a real kernel of
dimension at least
\[
\boxed{(385^2-1)-(720-1)=147505.}
\]
No physical 720-outcome detector or phase-sensitive amplitude
instrument is claimed.
\end{proposition}
\begin{proof}
The source's degree-nineteen positive integration formula exactly
integrates products of two harmonics through degree nine; hence
$S^*S=I_d$.  Its weighted rows $s_z$ define positive rank-one
operators summing to $I_d$.  The probability vector of any density
matrix has 720 entries with one trace-normalization relation, so its
traceless linear output rank is at most 719.  The real vector space
of traceless Hermitian $385\times385$ matrices has dimension
$385^2-1=148224$.  Rank-nullity gives the bound; phase-sensitive
vector reconstruction is a different data modality.
\end{proof}


\section{Full register harmonic inversion and complementary finite readouts}
\label{p1c:sec:cc20-register-inversion}
\begin{theorem}[Degree twelve detects the complete twenty-five-register space]
\label{p1c:thm:cc20-degree12-register-frame}
Let $S=U/\sqrt{24}:\mathbb R^{25}\to\mathbb R^{600}$ be the
isometric register embedding, and put
$K_n(c,d)=(n+1)U_n(c\cdot d)/600$ and $B_n=S^{\mathsf T}K_nS$.
With the established orthogonal projectors $P_0,P_a,P_c$ of ranks
$1,8,16$, respectively, the exact geometric register Gram values are
\[
 B_6=\frac7{16}P_c,\quad B_8=\frac9{16}P_c,\quad B_{10}=0,
 \qquad
 \boxed{B_{12}=\frac{325}{256}P_0+\frac{3003}{1024}P_a+
 \frac{325}{4096}P_c.}
\]
In particular the degree-twelve readout alone is injective on all
$25$ register directions, with inverse
\[
 B_{12}^{-1}=\frac{256}{325}P_0+\frac{1024}{3003}P_a+
 \frac{4096}{325}P_c.
\]
For $T_{12}^{\mathsf T}T_{12}=B_{12}$, the map
$x\mapsto T_{12}x$ has sharp inverse norm $64/\sqrt{325}$.
\end{theorem}
\begin{proof}
The established double-Hurwitz register matrix is constant on the
three pair types: equal labels, exactly one equal label, and neither
equal.  The exact degree-twelve zonal block sums are $26/25$, $39/100$,
and $-18499/102400$, respectively.  Diagonalizing this real symmetric
matrix on the $1\oplus8\oplus16$ register decomposition gives the
three displayed positive eigenvalues.  The degrees six, eight and ten
follow from the source's normalized facet harmonic projectors.  The
inverse and sharp stability factor follow by spectral calculus.
\end{proof}
\begin{corollary}[Complementary Hopf and degree-six register readout]
\label{p1c:cor:cc20-hybrid-hopf-degree6}
Combine the two fibre-average maps of
Theorem~\ref{p1b:thm:cc20-dual-hopf-incidence} with the degree-six
harmonic map $T_6$. Then
\[
 \boxed{C_L^{\mathsf T}C_L+C_R^{\mathsf T}C_R+T_6^{\mathsf T}T_6
 =\frac{24}{5}P_0+\frac{12}{5}P_a+\frac7{16}P_c>0.}
\]
This is a complete mathematical register readout using a degree-six
harmonic complement to the rank-nine phase-blind fibre marginals; its
condition number is $\sqrt{384/35}$.  On the registered response
$M(t)=P_0+e^{-5t/4}P_a+f_c(t)P_c$ one also has the exact transfer traces
\[
 \operatorname{Tr}(B_6M(t))=7f_c(t),\quad
 \operatorname{Tr}(B_8M(t))=9f_c(t),\quad
 \operatorname{Tr}(B_{12}M(t))=\frac{325}{256}(1+f_c(t))+
 \frac{3003}{128}e^{-5t/4}.
\]
\end{corollary}
\begin{proof}
The two Hopf maps have the rank-nine Gram of the imported finite
incidence theorem, while the degree-six Gram is $(7/16)P_c$.
These orthogonal contributions give the positive block-diagonal sum.
The transfer traces are the sums of products of the projector eigenvalues
with their multiplicities.  No continuous-time physical generator is
selected by this mathematical readout.
\end{proof}


\subsection{Exceptional invariant synthesis from the correlated chamber module}
\label{cc19:sec:correlation-invariant-synthesis}
Let $I=2I\subset SU(2)$, $F=IuI$ for $u=(1+i)/\sqrt2$,
$B_{ab}$ the twenty-five equal 24-facet registers, and $P_c$ the
rank-sixteen centred-register projector.  Define
\[
 b_n(q)_{ab}=\frac{n+1}{120}\sum_{c\in B_{ab}}U_n(q\cdot c),
 \qquad T_n^*T_n\big|_{P_c}=\lambda_n I_{16},
 \quad(\lambda_6,\lambda_8,\lambda_{12})
 =\left(\frac7{16},\frac9{16},\frac{325}{4096}\right).
\]
The third Gram scalar follows from Theorem~\ref{p1c:thm:cc20-degree12-register-frame}; its mathematical proof domain is retained. Set $\mathsf U_n=\lambda_n^{-1/2}T_nP_c$.
For a \emph{common} orthonormal basis $(e_\alpha)_{\alpha=1}^{16}$ of
$P_c\mathbb R^{25}$, put $f_{n,\alpha}=\mathsf U_ne_\alpha$ and
\[
\rho_n=\sum_\alpha |f_{n,\alpha}|^2,\quad
D_{nm}=\sum_\alpha f_{n,\alpha}\overline{f_{m,\alpha}},\quad
C_n(q)=\frac1{120}\sum_{g\in I}U_n\bigl(\mathrm{Re}(g^{-1}q)\bigr).
\]
\begin{theorem}[Exact first and second exceptional harmonic synthesis]
\label{cc19:thm:two-exceptional-generators}
Under these fixed scalar comparison premises, pointwise on $S^3$,
\[
 \boxed{\rho_6=16+12C_{12},\qquad\rho_8=16+20C_{12},}
\]
\[
 \boxed{D_{6,8}=-4\sqrt{63}C_{12},\qquad
 D_{6,12}=-4\sqrt{91}C_{12},}
\]
\[
 \boxed{D_{8,12}=\sqrt{13}\left(\frac{84}{17}C_{12}
                                      +\frac{120}{17}C_{20}\right).}
\]
Hence $C_{12}=(\rho_6-16)/12$ and
$C_{20}=17D_{8,12}/(120\sqrt{13})-7C_{12}/10$.
\end{theorem}
\begin{proof}
Every displayed quadratic/cross contraction is bi-$I$ invariant,
since the shared register module is the irreducible real
$\mathbf4_L\otimes\mathbf4_R$; its polynomial degree is at most
$2n$ or $n+m$.  The source Molien series supplies only the constant
and the unique degree-twelve invariant below degree twenty, and adds
one degree-twenty invariant at that degree.  The cross terms have
zero Haar mean by orthogonality of distinct $S^3$ scalar harmonics.
The exact $\mathbb Q(\sqrt5)$ Hurwitz-coset sums at the identity give
$(b_6)_{aa}=-7/5$, $(b_6)_{a\ne b}=7/20$,
$(b_8)_{aa}=9/5$, $(b_8)_{a\ne b}=-9/20$,
$(b_{12})_{aa}=0$, and $(b_{12})_{a\ne b}=-65/64$.
These yield $\rho_6(1)=28$, $\rho_8(1)=36$,
$D_{6,8}(1)=-4\sqrt{63}$, $D_{6,12}(1)=-4\sqrt{91}$,
and $D_{8,12}(1)=12\sqrt{13}$.
The further exact facet-centre contraction
$b_8(u)^TP_cb_{12}(u)=2925/2048$ gives
$D_{8,12}(u)=15\sqrt{13}/8$.
Since $C_{12}(1)=C_{20}(1)=1$ and
$(C_{12}(u),C_{20}(u))=(-5/16,31/64)$,
the unique invariant coefficients are the stated ones.
\end{proof}
\begin{corollary}[Exact alias repair and a selected degree-twelve frame]
\label{cc19:cor:720-mixed-band-frame}
On $W_6\oplus W_8$, normalized 120-vertex and 600-facet sampling have
respective two-by-two Gram blocks, each tensored with $I_{16}$,
\[
G_V=\frac14\begin{pmatrix}7&-\sqrt{63}\\-\sqrt{63}&9\end{pmatrix},
\quad G_F=\frac1{64}\begin{pmatrix}49&5\sqrt{63}\\5\sqrt{63}&39\end{pmatrix}.
\]
Their spectra are $\{4,0\}$ and $\{21/16,1/16\}$, while
$(5G_V+16G_F)/21=I_2$.  In particular all 32 modes are
720-node isometrically sampled, though 120 vertices alone retain an
exact sixteen-dimensional alias.
For the selected 48-dimensional $W_6\oplus W_8\oplus W_{12,c}$,
the positive 720-node Gram is
\[
\boxed{\begin{pmatrix}1&0&0\\0&1&15\sqrt{13}/56\\
0&15\sqrt{13}/56&2795/1792\end{pmatrix}\otimes I_{16}.}
\]
The nontrivial two-by-two determinant is $7865/12544>0$,
so these 48 selected modes are also reconstructible.
\end{corollary}
\begin{proof}
Each orbit-averaged sampling form commutes with the $I_L\times I_R$
action on the irreducible $\mathbf4\otimes\mathbf4$ and so is scalar
in each matched pair of harmonic degrees.  Taking traces of the
pointwise identities in the theorem, using $C_{12}=1$ on vertices,
$C_{12}=-5/16$ on facet centres, and the exact degree-twelve
correlation intensities $52$ and $4225/256$, gives the matrices.
The degree-nineteen design removes $6+8$ and $6+12$ cross terms,
while the degree-twenty $8+12$ term remains $15\sqrt{13}/56$.
\end{proof}
\begin{corollary}[A source-geometric invariant bound]
\label{cc19:cor:invariant-positivity-bound}
$-1/3\le C_{12}(q)\le1$ for every $q\in S^3$.
\end{corollary}
\begin{proof}
The matched $2\times2$ pointwise Gram of the sextic and octic
16-vectors is positive semidefinite, with determinant
$256(1-C_{12})(1+3C_{12})$. The two diagonal entries are
nonnegative. The determinant inequality gives the stated interval.
\end{proof}


\subsection{An exact higher-band scalar spectral and nonlinear compiler (mathematical comparison)}
\label{p1c:sec:cc110-orbit-spectral-propagation}

\begin{theorem}[Positive orbit cubatures induce exact finite scalar projectors]
\label{p1c:thm:cc110-orbit-spectral-projectors}
Let $\nu$ be one of the four positive orbit cubatures in
Theorem~\ref{p1b:thm:cc110-positive-orbit-cubatures}, with
$N$ nodes $z_i$ and positive per-node weights $\omega_i$.
For its degree $d\in\{23,29,31,35\}$, let
$B=\lfloor d/2\rfloor$, and define the weighted kernel matrices
\[
(E_n)_{ij}=\sqrt{\omega_i\omega_j}(n+1)
 U_n(z_i\cdot z_j),\qquad 0\le n\le B.
\]
Then $E_n E_m=\delta_{nm}E_n$, $\operatorname{rank} E_n=(n+1)^2$,
and the exact sample image has dimensions
\[
\boxed{\operatorname{rank}\sum_{n=0}^{B}E_n
=\sum_{n=0}^B(n+1)^2=(650,1240,1496,2109).}
\]
The finite round scalar Laplacian and spectral functional calculus on
that image are exactly
\[
L_B=R^{-2}\sum_{n=0}^B n(n+2)E_n,
\quad \operatorname{Tr}_{\operatorname{im}S}F(L_B)
=\sum_{n=0}^B(n+1)^2F(n(n+2)/R^2).
\]
All complementary sample directions are excluded from the scalar band
trace; their count for the degree-35 rule is $8880-2109=6771$.
\end{theorem}
\begin{proof}
The round scalar kernel is $(n+1)U_n(z_i\cdot z_j)$.
Since $n+m\le2B\le d$, the positive cubature exactly integrates
products of any two scalar harmonics in the designated band.
Weighted evaluation is therefore an isometry, so its shell
projectors are mutually orthogonal with their full harmonic ranks.
The Laplacian eigenvalues are $n(n+2)/R^2$ on the $n$th shell.
\end{proof}

\begin{corollary}[Exact finite nonlinear Galerkin coefficient budget]
\label{p1c:cor:cc110-nonlinear-galerkin-budget}
For $f\in\mathcal P_a$, $g\in\mathcal P_b$ and
$h\in\mathcal P_c$, a cubature of degree $d$ satisfies
\[
\int_{S^3}\overline hfg\,d\sigma
=\sum_i\omega_i\overline{h(z_i)}f(z_i)g(z_i)
\quad\text{whenever }a+b+c\le d.
\]
In particular, for the 8880-node degree-35 rule, projecting products of
two band-nine scalar fields into band seventeen is exact
($9+9+17=35$). Quartic scalar actions on band-eight fields,
including their band-projected first and second variations, are exact
because $4\times8=32\le35$. This is a finite mathematical Galerkin
calculus; no physical interaction coefficient is introduced.
\end{corollary}

\begin{corollary}[Sharp first surviving heat error and invariant tomography]
\label{p1c:cor:cc110-heat-error-invariant-tomography}
For an orbit-cubature measure $\nu$ exact through degree $d$ with first
uncancelled invariant degree $n_*=d+1$, write
$r_{n_*}=\nu(C_{n_*})\ne0$.
Then, for the positive round scalar Laplacian $\Delta_+$,
\[
\boxed{\|\nu e^{-t\Delta_+}-\sigma e^{-t\Delta_+}\|_{L^2\to\mathbb C}
\sim\sqrt{n_*+1}\,|r_{n_*}|
 e^{-n_*(n_*+2)t/R^2},\quad t/R^2\to\infty.}
\]
The four values of $n_*$ are $24,30,32,36$, respectively.
Furthermore, for each $d$ above, the $k$ uniform orbit means of
$f\in\mathcal P_d$ uniquely determine the $k$ coefficients of its
$I_L\times I_R$-invariant projection in the normalized basis
$\{1,C_{12},C_{20},\ldots,C_{d_{\mathrm{last}}}\}$.
The coefficients are recovered by the inverse of the exact finite
moment matrix of Theorem~\ref{p1b:thm:cc110-positive-orbit-cubatures}.
\end{corollary}
\begin{proof}
Peter--Weyl diagonalizes the round heat operator in degrees $n$.
The first nonzero cubature Fourier transform is rank one because the
corresponding $2I$ invariant multiplicity is one; its eigenvalue is the
scalar residual $r_{n_*}$. The squared dual functional norm on shell
$n$ is $(n+1)\|\widehat\nu(n)\|_{HS}^2$, so the first term gives the
asymptotic formula. Orbit means annihilate all non-bi-invariant
harmonics; on the invariant subspace they act by the displayed finite
moment matrix, already shown to be invertible.
\end{proof}

\begin{remark}[Downstream mathematical role]
Potential consumers include Paper III-B's scalar spectral/Galerkin
comparisons, Paper IV-A's finite evaluation calculus, and conditional
Paper VI-B/VII-C sampling tests. Neither the complete physical Dirac
operator nor the whole Born--Hopf action is identified with these
scalar comparison matrices. Exact geometric quadrature does not prove
physical source occurrence or continuum QFT reconstruction.
\end{remark}


\subsection{An exact 3311-dimensional scalar spectral window}
\label{p1c:sec:cc111-band20}
\begin{theorem}[Positive 12720-node band-twenty scalar sampling]
\label{p1c:thm:cc111-band20}
Assume the mathematical positive degree-41 rule of the proposed
Paper I-B theorem and let $z_j$ be its $12720$ distinct nodes with
per-node weights $\omega_j>0$. For $0\le n\le20$ define
\[
(E_n)_{ij}=\sqrt{\omega_i\omega_j}(n+1)
U_n(\operatorname{Re}(z_i^{-1}z_j)),\qquad P_{20}=\sum_{n=0}^{20}E_n.
\]
Then $E_nE_m=\delta_{nm}E_n$,
$\operatorname{rank}E_n=(n+1)^2$ and
\[
\boxed{\operatorname{rank}P_{20}=3311,\quad
\dim\ker P_{20}=9409.}
\]
The exact round-sphere scalar band Laplacian and its band trace are
\[
L_{20}=R^{-2}\sum_{n=0}^{20}n(n+2)E_n,\qquad
\operatorname{Tr}_{P_{20}}F(L_{20})
=\sum_{n=0}^{20}(n+1)^2F(n(n+2)/R^2).
\]
The rule integrates exactly quartic polynomial potentials and their
projected first and second variations for scalar fields of harmonic
degree at most ten, when all tested variations lie in the same band.
\end{theorem}
\begin{proof}
Products of two degree-$\le20$ harmonics have coordinate degree
at most $40<41$, so weighted evaluation is an exact isometry.
The scalar reproducing kernel is $(n+1)U_n$; hence the $E_n$ are
mutually orthogonal projections with the stated ranks.
The scalar Laplacian has eigenvalue $n(n+2)/R^2$ on the $n$th
harmonic shell. Quartic integrands and their projected variations
contain no more than four degree-ten factors (total degree $40$),
so the same cubature evaluates them exactly.
\end{proof}


\section{Finite scalar spectral and variational calculus through degree thirty}
\label{p1c:sec:cc112-band30}
\begin{theorem}[Positive band-thirty scalar sampling; mathematical comparison]
\label{p1c:thm:cc112-band30-spectral-projectors}
Suppose the exact positive degree-$61$ sixty-thousand-node cubature
of the proposed Paper I-B theorem is used, with point weights
$\omega_j=w_{O(j)}/|O(j)|$. Its weighted evaluation operator
$S:\mathcal P_{30}(S^3)\to\mathbb C^{60000}$ is an isometry on the
complete scalar harmonic band
\[
 \dim \mathcal P_{30}(S^3)=\sum_{n=0}^{30}(n+1)^2=10416.
\]
For each harmonic degree $0\le n\le30$, the matrix-free kernel
\[
 (E_n)_{ij}=\sqrt{\omega_i\omega_j}\,(n+1)
 U_n\bigl(\operatorname{Re}(z_i^{-1}z_j)\bigr)
\]
represents an orthogonal projector of rank $(n+1)^2$ on the sample
image: $E_nE_m=\delta_{nm}E_n$. In particular
\[
 L_{30}=R^{-2}\sum_{n=0}^{30}n(n+2)E_n
\]
acts as the exact scalar Laplacian on the sampled band, and for any
function $F$ defined on these $31$ eigenvalues
\[
 \operatorname{Tr}_{\mathcal P_{30}}F(-\Delta)
 =\sum_{n=0}^{30}(n+1)^2F\bigl(n(n+2)/R^2\bigr).
\]
The $49584$ sample-image complement directions are not additional
physical zero modes.
\end{theorem}
\begin{proof}
The product of two band-$30$ harmonics has scalar polynomial degree
at most $60$, so degree-$61$ exact integration reproduces their
continuous $L^2$ inner product. The standard scalar $S^3$ addition
kernel then provides the mutually orthogonal degree projectors.
The spectral calculus follows on the isometric image.
\end{proof}
\begin{corollary}[Polynomial Galerkin degree budget]
\label{p1c:cor:cc112-exact-nonlinear-degree-budget}
A finite sampled polynomial interaction whose total field-and-test
degree is at most $61$ has exactly its continuous band-projected
scalar Galerkin integral. Examples include quartic actions on fields
in $\mathcal P_{15}$, quintic expressions on $\mathcal P_{12}$, and
two $\mathcal P_{30}$ factors tested against $\mathcal P_1$.
\end{corollary}
\noindent\emph{Scope:} This scalar Galerkin identity neither
identifies the framework's interacting physical action nor supplies
an action-owned operator/source matching certificate.


% BEGIN RECIPROCAL BERGER ROLE FIREWALL v1.17
\section{Physical localization versus mathematical reciprocal Berger metrics}
\label{bh:berger:v117:p1c:reciprocal-scope}
\begin{remark}[Independent metric parameters do not imply independent physical variation]
\label{bh:berger:v117:p1c:fixed-momentum}
The full mathematical $S_x^3\times S_p^3\times T^2$ metric atlas admits
separate Berger parameters $s_x,s_p$, and Born exchange compares
particular mathematical metric choices.  On the physically localized
branch owned by Paper III-F, however, the position metric is varied
against a source while the already constituted compact momentum seed
is retained as fixed spectral data. Thus $\delta g_x\ne0$ does not
imply $\delta g_p\ne0$; no new value of $m_p$ is generated.
The equality $s_x=s_p$ is an isometry condition in a chosen mathematical
comparison family, not a derived physical evolution equation.
The exact orthogonal-product Dirac square retains its own factor
conditions, so a spatial Dirac crossing is not automatically a zero
of the full eight-dimensional product operator.
\end{remark}
% END RECIPROCAL BERGER ROLE FIREWALL v1.17


\begingroup
\small
\printbibliography
\endgroup


% BEGIN FRAMEWORK v3.068 GEOMETRY BRIDGE SOURCE INTEGRATION
\section{Quadratic Born--Petersen Cartan geometry: first non-gauge test family}
\noindent\emph{Scope.} The statements below retain their declared finite geometric, source-chart, and action-relative hypotheses. They do not assert physical particle decay or unconditional source occupation.
\begin{theorem}[First non-gauge Born--Petersen Cartan harmonic and exact static response]
\label{p1c:thm:v368-first-born-petersen-harmonic}
On the generated round relative $S_R^3$, let $\theta=q^{-1}dq$,
$A_B=\theta/2$, and fix a marked unit imaginary quaternion $u_0$.
Let $K:\mathbb R^{15}\to\mathbb R^3$ be the existing oriented Petersen flux,
$KK^T=5I_3$, and $v_e=-(2/3)P_{u_0}Ke$.
Within the pre-existing Cartan family $A_w=(1+w(q))\theta/2$,
the degree-one inversion-odd source $w_e^{(1)}=v_e\cdot\operatorname{Im}q$
has zero transverse background-gauge projection at $A_B$.
Under the diagonal quaternion-conjugation symmetry the unique
inversion-odd degree-two scalar triplet is
$w_e^{(2)}(q)=q_0(v_e\cdot\operatorname{Im}q)$.
For $a_e=w_e^{(2)}\theta/2$, its gauge-transverse projection has exact
normalized-Haar $L^2$ Gram
\[
 \langle P_Ta_e,P_Ta_f\rangle_{S_R^3}
  =\frac{5}{648R^2}\,e^TK^TP_{u_0}Kf,
\]
with two nonzero Gram eigenvalues $25/(648R^2)$ and thirteen zero eigenvalues.
The existing normalized static Yang--Mills and Chern--Simons functionals
along $A_{t w_e^{(2)}}$ evaluate for every real $t$ to
\[
 \mathcal E_e(t)=\frac14+\frac5{324}t^2\|P_{u_0}Ke\|^2
                     +\frac1{4320}t^4\|P_{u_0}Ke\|^4,
 \qquad N_{\mathrm{CS},e}(t)=\frac12.
\]
Thus the normalized energy Hessian on the fifteen-edge source is
$(5/162)K^TP_{u_0}K$; its two nonzero eigenvalues are $25/162$.
The first degree-two transverse curvature mode has native spatial
Yang--Mills Hessian eigenvalue $6/R^2$ on its image.
These are static connection-geometry and test-space results; they do not
assert an action-selected Petersen insertion, a proper weak vertex, or
particle decay.
\end{theorem}
\begin{proof}
Theorem~\ref{p1b:thm:v2217-nonconstant-cartan-dirac} gives the exact
background-gauge survival fraction
$1-n(n+2)/[3(n(n+2)-2)]$, equal to $0$ for $n=1$ and $5/9$ for $n=2$,
and the native energy and Chern--Simons functionals.
The diagonal-conjugation action on quaternion coordinates is
$\mathbb R^4=\mathbf1\oplus\mathbf3$, and
$\mathcal H_2=\mathrm{Sym}^2(\mathbb R^4)_0=\mathbf1\oplus\mathbf3\oplus\mathbf5$;
its inversion-odd triplet is spanned by $q_0q_i$.
The exact Haar moments are
$\langle q_0^2q_i^2\rangle=1/24$ and
$\langle q_0^4q_i^4\rangle=3/640$, with scalar Laplace eigenvalue
$8/R^2$, giving the energy polynomial and constant Chern--Simons value.
Theorem~\ref{p1:thm:petersen-signed-and-flux} gives $KK^T=5I_3$,
and Proposition~\ref{p1:prop:v371-petersen-quadrupole-flux}
gives the fixed factor $-2/3$. The normalized unit source
$a_v=(q_0v\cdot\operatorname{Im}q)\theta/2$ has
$\|a_v\|^2=|v|^2/(32R^2)$;
its projected norm is $5|v|^2/(288R^2)$.
Polarization and $v_e=-(2/3)P_{u_0}Ke$ prove the source Gram.
Differentiation of the exact energy curve proves its Hessian.
No physical insertion is used in any step.
\end{proof}
% END FRAMEWORK v3.068 GEOMETRY BRIDGE SOURCE INTEGRATION

\end{document}
